\documentclass[11pt]{article}
\usepackage[margin=1in]{geometry}
\usepackage{amsmath,amssymb,amsthm}
\usepackage{enumitem}
\usepackage[hidelinks]{hyperref}

\newtheorem{theorem}{Theorem}[section]
\newtheorem{lemma}[theorem]{Lemma}
\newtheorem{proposition}[theorem]{Proposition}
\newtheorem{corollary}[theorem]{Corollary}
\newtheorem{definition}[theorem]{Definition}
\newtheorem{assessment}[theorem]{Assessment}
\newtheorem{firewall}[theorem]{Firewall}

\title{Renewal Standard Model--Einstein Continuum Research Notes\\
Fifth Series, V1}
\author{}
\date{\today}

\begin{document}
\maketitle
\tableofcontents

\section{Context and purpose}
\label{sec:v5v1-context}

These notes continue a long constructive programme aimed at a full
Standard Model--Einstein continuum theory.  The earlier files became too
large to track efficiently.  In accordance with the version rule, the
fourth series has been frozen at
\texttt{Renewal\_SM\_Einstein\_Continuum\_Research\_Notes\_V100\_f4.tex}.
This new V1 is a compact restart.  It records the major proved results,
states the present obstruction boundary without claiming that it has been
crossed, and fixes the next acquisition order.

The goal remains the full continuum: a cutoff-independent, anomaly-free,
gauge- and diffeomorphism-covariant Schwinger functional with controlled
fermion determinants and zero modes, reflection positivity, clustering,
and Osterwalder--Schrader reconstruction.  Conditional certificates are
useful only when their hypotheses are attacked explicitly.

\section{Major results inherited from the archived series}
\label{sec:v5v1-results}

\subsection{Fermion topology, determinant lines, and zero modes}

The archived work constructs the anomaly-normalized three-generation
Standard Model determinant line for the faithful gauge group, separates
hard and soft fermion sectors by exact Schur complements, and transports
their determinant lines without assigning independent anomaly phases to
the factors.  It records the chiral index and the electroweak and QCD
zero-mode selection rules, including exact finite Grassmann saturation.
A determinant zero is retained as a zero of a line section; it is never
removed by an inverse soft gap.

\subsection{Global spectral reduction}

For a fixed principal symbol, an intrinsic relative perturbation
\(K\in\mathcal S_{4,\infty}\) is split into a norm-small tail and a
finite-rank large part.  The exact factorization
\[
 1+K=T(1+AB),\qquad
 \det_F(1+AB)=\det_E(1+BA)
\]
captures every hard kernel and cokernel in a finite Sylvester matrix.
Weak-Schatten counting pays its rank by the critical \(L^4\) load.
Frozen-centre charts, quantitative singular-value collars, and
determinant-line transition maps sew these blocks globally without a
Cech-index shortcut.

A partition-flat reference connection is constructed from the Uhlenbeck
transition cocycle:
\[
 A_{\sharp,i}=-\sum_j\rho_j\,dg_{ij}g_{ij}^{-1},\qquad
 A_i-A_{\sharp,i}=\sum_j\rho_jg_{ij}A_jg_{ij}^{-1}.
\]
This identity retains the bundle topology in one owner and gives the
action-paid global \(L^4\) comparison.

\subsection{Metric normalization and renormalization}

A metric variation changes the Dirac principal symbol, so its direct
relative operator is order zero and generally noncompact.  The archived
work corrects this by an elliptic order-zero microlocal normalizer \(Q\):
\[
 Q_{g,\circ}D_g-D_\circ\in\Psi^0,\qquad
 (Q_{g,\circ}D_g-D_\circ)C_\circ\in\Psi^{-1}
 \subset\mathcal S_{4,\infty}.
\]
The order-zero normalizer is not discarded.  Its ultraviolet defect is
assigned to local gravitational counterterms.

The four-dimensional heat coefficients identify the cosmological,
Einstein--Hilbert, gauge, Higgs, nonminimal, curvature-squared, and boundary
owners.  Half-integer boundary coefficients produce cubic and linear
cutoff powers on artificial core boundaries.  The finite determinant has
the exact form
\[
 \exp(\mathcal B_{\le4}^{\rm ren}+\mathcal T_{\rm loc})
 \det{}_5(1+K)
\]
as a determinant-line section.  The first four renormalized Born
amplitudes retain their finite nonlocal parts; \(\det_5\) controls the
high-order tail and remains entire at zeros.

For \(K\in\Psi^{-1}\), the marginal logarithm is the local Wodzicki
residue
\[
 \operatorname{res}(K^4)
 =\frac1{(2\pi)^4}\int_{S^*X}
   \operatorname{tr}\sigma_{-1}(K)^4.
\]
Its scale change is exact.  A fixed-scale \(L^4\) bound fails under
concentration when this residue is nonzero, while evaluation of one running
coupling trajectory at each terminal-core scale cancels the logarithm
without introducing core-dependent theories.

\subsection{Wilsonian compactness and positivity}

Finite-depth compatible ancestors in compact admissible shell spaces have
an all-scale inverse-limit trajectory.  The irrelevant interaction sector
is compact from a strong rooted kernel norm to a weaker one when there is
one derivative of smoothing, a strict exponential localization gap, and
an arity gap.  Covariant Taylor forests remove local distributions, and
literal heat-shell spanning trees give graphwise strong bounds.  A
Banach-valued Kotecky--Preiss condition sums the Ward-complete small-field
activities.

Reflection positivity is formulated directly on reflected cylinder
matrices:
\[
 \sum_{i,j}\overline c_i c_j
 S((\Theta F_i)F_j)\ge0.
\]
It is closed under entrywise convergence.  A positive transfer
factorization or an exact contour deformation preserving every reflected
matrix element is sufficient; a formal conformal rotation is not.

\section{Established obstruction boundary}
\label{sec:v5v1-obstructions}

\begin{theorem}[Present-route boundary]
\label{thm:v5v1-boundary}
The archived results do not yield a full ultraviolet continuum while one
simultaneously keeps:
\begin{enumerate}
 \item the unmodified \(U(1)_Y\) Gaussian-boundary trajectory;
 \item a finite Einstein coupling list; and
 \item the real Euclidean Einstein conformal contour.
\end{enumerate}
\end{theorem}
\begin{proof}
The inherited one-loop coefficients are
\[
 b_3=7,\qquad b_2=\frac{19}{6},\qquad b_Y=-\frac{41}{6}.
\]
Under the proved Gaussian-dominance hypothesis, positive inverse
hypercharge coupling fails after finitely many ultraviolet shells.
Independently, the nonzero two-loop curvature-cubic divergence shows that a
finite weight-at-most-four Einstein list is not closed.  Finally, for
\(g=e^{2\varphi}\bar g\),
\[
 S_{\rm EH}(g)\supset
 -\frac{6}{16\pi G}\int e^{2\varphi}|\bar\nabla\varphi|^2,
\]
which is unbounded below along fixed-amplitude high-frequency conformal
modes.  Thus the real Einstein factor supplies no large-field suppression.
\end{proof}

\begin{firewall}[Meaning of the boundary]
The theorem does not exclude every ultraviolet completion compatible with
low-energy SM--Einstein physics.  It excludes treating the three listed
choices as though the existing estimates had already made them into a
compact, positive all-scale construction.
\end{firewall}

\section{Six live gates}
\label{sec:v5v1-gates}

\begin{enumerate}
 \item Construct an ultraviolet completion or fixed point whose trajectory
 matches the Standard Model and exists at every terminal-core scale.
 \item Control the generated gravitational operator tower rather than
 truncating it at curvature squared.
 \item Construct a metric integration prescription satisfying the metric
 OS certificate: reflected positivity, domination, and convergence.
 \item Prove the renormalized one-through-four insertion bounds on the
 actual rough critical profile space.
 \item Prove the Ward-complete polymer and large-field estimates for the
 joint gauge--Higgs--fermion--metric shell.
 \item Establish tightness, clustering, uniform integrability, and the
 remaining OS axioms for the cofinal Schwinger functions.
\end{enumerate}

\section{Fifth-series acquisition order}
\label{sec:v5v1-next}

A proposed ultraviolet mechanism must provide compact admissible sets
\(K_j\), exact shell maps
\(\mathcal R_j:K_{j+1}\to K_j\), and finite-depth ancestors of the matched
low-energy datum.  The first new task is to study a mechanism-independent
fixed-point entrance theorem.

\begin{definition}[Fixed-point entrance data]
\label{def:v5v1-entrance}
A shell fixed-point candidate consists of a compact invariant set
\(K_*\), a fixed point \(x_*\), a local splitting into finitely many
relevant/marginal coordinates and a rooted irrelevant kernel space, a
literal contraction on the latter, and Ward-, anomaly-, and
reflection-compatible coordinate charts.
\end{definition}

\begin{assessment}[V2 target]
V2 will prove a graph-transform theorem for the irrelevant kernel
coordinate over a prescribed bounded marginal trajectory.  The theorem
will state the exact contraction and Lipschitz constants needed for a
unique ultraviolet stable graph.  It will then test the result against the
hypercharge direction: the theorem may organize a genuine interacting
fixed point, but cannot manufacture one from the Gaussian trajectory.
\end{assessment}

The full SM--Einstein continuum remains unproved.  This restart preserves
the proved mathematical tools and directs the next iteration at the first
ultraviolet mechanism that can change that status.

\clearpage
\section{V2: nonautonomous stable graph over a marginal trajectory}
\label{sec:v5v2}

The fifth-series V1 asks for a fixed-point entrance theorem that does not
pretend to construct the marginal ultraviolet trajectory.  This note proves
the fibre theorem.  Given a bounded compatible marginal trajectory, a
literal contraction in the rooted irrelevant kernel norm produces a unique
bounded ultraviolet irrelevant ancestry, with two marked derivatives.  The
hypercharge and gravitational directions remain base-space gates.

\subsection{Sequence-space formulation}

Let \(Y=\mathfrak X_{\rm w}\) be the rooted weak kernel space from the
fourth-series V96.  Let
\[
 u=(u_j)_{j\ge0}
\]
be a prescribed compatible trajectory in the local relevant/marginal
coordinates.  The irrelevant part of the shell equation has the inverse
orientation
\begin{equation}
 v_j=F_j(u_{j+1},v_{j+1}),\qquad j\ge0.
\label{eq:v5v2-shell}
\end{equation}
Put
\[
 \mathscr Y_\infty:=\ell^\infty(\mathbb N_0;Y),\qquad
 \|v\|_\infty:=\sup_j\|v_j\|_Y,
\]
and define
\begin{equation}
 (\mathcal T_uv)_j:=F_j(u_{j+1},v_{j+1}).
\label{eq:v5v2-T}
\end{equation}

\begin{theorem}[Uniform nonautonomous graph transform]
\label{thm:v5v2-graph}
Assume there are \(R>0\), \(r_*\ge0\), and \(0\le c<1\) such that,
uniformly in \(j\),
\begin{align}
 \|F_j(u_{j+1},0)\|_Y&\le r_*,
\label{eq:v5v2-source}\\
 \|F_j(u_{j+1},y)-F_j(u_{j+1},\widetilde y)\|_Y
 &\le c\|y-\widetilde y\|_Y
\label{eq:v5v2-contraction}
\end{align}
for \(\|y\|,\|\widetilde y\|\le R\), and
\begin{equation}
 r_*+cR\le R.
\label{eq:v5v2-invariant-ball}
\end{equation}
Then \(\mathcal T_u\) has a unique fixed point
\(v(u)\) in the closed radius-\(R\) ball of \(\mathscr Y_\infty\).
It is the unique bounded irrelevant trajectory in that ball satisfying
\eqref{eq:v5v2-shell}, and
\begin{equation}
 \|v(u)\|_\infty\le\frac{r_*}{1-c}.
\label{eq:v5v2-bound}
\end{equation}
\end{theorem}
\begin{proof}
For \(\|v\|_\infty\le R\),
\[
 \|(\mathcal T_uv)_j\|
 \le r_*+c\|v_{j+1}\|\le r_*+cR\le R.
\]
Thus the closed ball is invariant.  Moreover,
\[
 \|\mathcal T_uv-\mathcal T_u\widetilde v\|_\infty
 \le c\|v-\widetilde v\|_\infty.
\]
Banach's fixed-point theorem gives existence and uniqueness.  At the fixed
point,
\(\|v\|_\infty\le r_*+c\|v\|_\infty\), which proves
\eqref{eq:v5v2-bound}.
\end{proof}

\begin{corollary}[Exact stable ancestry series]
\label{cor:v5v2-series}
If
\[
 F_j(u,y)=A_j(u)y+r_j(u),\qquad \|A_j(u)\|\le c<1,
\]
then
\begin{equation}
 v_j=
 \sum_{k=j}^{\infty}
 A_jA_{j+1}\cdots A_{k-1}r_k,
\label{eq:v5v2-series}
\end{equation}
with the empty product equal to the identity.
\end{corollary}
\begin{proof}
Iterate the equation to level \(N+1\).  The terminal term has norm at most
\(c^{N-j+1}R\) and tends to zero.  The series is dominated by
\(\sum_{k\ge j}c^{k-j}r_*\).
\end{proof}

\subsection{Entrance into the fixed point}

Let \(r_j:=\|F_j(u_{j+1},0)\|_Y\).

\begin{proposition}[Vanishing source gives ultraviolet entrance]
\label{prop:v5v2-entrance}
Under the hypotheses of Theorem~\ref{thm:v5v2-graph}, if \(r_j\to0\), then
\[
 \|v_j(u)\|_Y\longrightarrow0.
\]
More quantitatively,
\begin{equation}
 \|v_j\|_Y\le\sum_{k=j}^{\infty}c^{k-j}r_k.
\label{eq:v5v2-tail}
\end{equation}
\end{proposition}
\begin{proof}
Iteration gives the finite version of \eqref{eq:v5v2-tail} plus a terminal
term bounded by \(c^{N-j+1}R\).  Let \(N\to\infty\).  To prove convergence
to zero, split the series after a fixed number of terms.  The finite head
tends to zero because \(r_k\to0\), and the geometric tail is uniformly
small.
\end{proof}

\begin{firewall}[Bounded ancestry is not fixed-point entrance]
\label{fw:v5v2-bounded}
A bounded nonzero source produces a bounded irrelevant graph but need not
give \(v_j\to0\).  Ultraviolet entrance requires decay of the source or a
nonzero fixed irrelevant coordinate included in the definition of
\(x_*\).
\end{firewall}

\subsection{Dependence on the marginal path}

Let \(\mathscr U\) be a metric space of bounded marginal trajectories with
metric \(d_{\mathscr U}\).

\begin{theorem}[Lipschitz stable graph]
\label{thm:v5v2-Lipschitz}
Assume, in addition,
\[
 \sup_j\sup_{\|y\|\le R}
 \|F_j(u_{j+1},y)-F_j(\widetilde u_{j+1},y)\|_Y
 \le L_u d_{\mathscr U}(u,\widetilde u).
\]
Then
\begin{equation}
 \|v(u)-v(\widetilde u)\|_\infty
 \le\frac{L_u}{1-c}d_{\mathscr U}(u,\widetilde u).
\label{eq:v5v2-Lipschitz}
\end{equation}
\end{theorem}
\begin{proof}
Insert and subtract
\(\mathcal T_uv(\widetilde u)\):
\[
 \|v(u)-v(\widetilde u)\|_\infty
 \le c\|v(u)-v(\widetilde u)\|_\infty
     +L_ud_{\mathscr U}(u,\widetilde u).
\]
Rearrange.
\end{proof}

\subsection{Two marked derivatives}

Let \(z\) range over a finite-dimensional parameter chart and write
\(\mathcal T(z,v)\) for \(\mathcal T_{u(z)}v\).

\begin{theorem}[Twice-marked graph]
\label{thm:v5v2-marks}
Suppose \(\mathcal T\) is \(C^2\) on the invariant ball,
\[
 \|D_v\mathcal T\|\le c<1,
\]
and all first \(z\)-derivatives are bounded by \(B_1\), while all second
partial derivatives in \((z,v)\) are bounded by \(B_2\).  Then the fixed
point \(v(z)\) is \(C^2\).  With
\(A=D_v\mathcal T(z,v(z))\),
\begin{align}
 \partial_av
 &=(1-A)^{-1}\partial_a\mathcal T,
\label{eq:v5v2-D1}\\
 \partial_a\partial_bv
 &=(1-A)^{-1}\Bigl(
 \partial_a\partial_b\mathcal T+
 D_{\partial_av}\partial_b\mathcal T+
 D_{\partial_bv}\partial_a\mathcal T+
 D_v^2\mathcal T[\partial_av,\partial_bv]\Bigr).
\label{eq:v5v2-D2}
\end{align}
In particular,
\[
 \|\partial_av\|_\infty\le\frac{B_1}{1-c},
\]
and
\[
 \|\partial_a\partial_bv\|_\infty
 \le\frac{B_2}{1-c}
 \left(1+\frac{2B_1}{1-c}
             +\frac{B_1^2}{(1-c)^2}\right).
\]
\end{theorem}
\begin{proof}
The Neumann series gives
\(\|(1-A)^{-1}\|\le(1-c)^{-1}\).  Differentiate
\(v=\mathcal T(z,v)\) once and solve for \(\partial_av\), giving
\eqref{eq:v5v2-D1}.  Differentiate again, retaining the two ordered mixed
terms and the quadratic fibre term, to obtain \eqref{eq:v5v2-D2}.
The implicit-function theorem on the sequence Banach space gives \(C^2\)
dependence, and the displayed bounds follow.
\end{proof}

The marked theorem applies to the rooted kernel sectors only after the
Ward-complete shell map proves the literal contraction and its marked
bounds.  It does not differentiate a gauge choice; frozen-centre
determinant charts remain in force.

\subsection{Sharpness of strict contraction}

\begin{proposition}[The boundary \(c=1\)]
\label{prop:v5v2-sharp}
At \(c=1\), existence or uniqueness can fail even in \(Y=\mathbb R\).
\end{proposition}
\begin{proof}
For \(F_j(y)=y+1\), the equation \(v_j=v_{j+1}+1\) has no bounded solution.
For \(F_j(y)=y\), every constant sequence is a bounded solution.  Thus
strict contraction is needed for the general theorem.
\end{proof}

\subsection{The marginal obstruction remains upstream}

Let the prescribed base contain the positive inverse hypercharge coupling.
The fourth-series V6/V94 result shows that, in the Gaussian-dominance wedge,
no such base trajectory reaches every ultraviolet depth.

\begin{corollary}[The stable graph cannot create hypercharge ancestry]
\label{cor:v5v2-hypercharge}
Theorem~\ref{thm:v5v2-graph} cannot repair the missing hypercharge
trajectory or close the gravitational local-operator tower.
\end{corollary}
\begin{proof}
The theorem assumes the marginal path \(u\) before defining
\(\mathcal T_u\).  If no admissible all-scale \(u\) exists, there is no base
over which to construct the graph.  Local gravitational operators also
belong to \(u\), not to the compact rooted fibre \(Y\).
\end{proof}

\begin{theorem}[Fixed-point entrance certificate]
\label{thm:v5v2-certificate}
A candidate ultraviolet fixed point supplies the irrelevant part of the
RILC if it provides:
\begin{enumerate}
 \item an all-scale bounded marginal trajectory \(u\);
 \item a Ward-complete shell map satisfying
 \eqref{eq:v5v2-source}--\eqref{eq:v5v2-invariant-ball} in the rooted
 kernel norm, through two marks; and
 \item \(r_j\to0\) after subtracting the candidate fixed interaction.
\end{enumerate}
Then there is a unique twice-marked irrelevant trajectory converging to the
fixed interaction.
\end{theorem}
\begin{proof}
Apply Theorem~\ref{thm:v5v2-graph},
Theorem~\ref{thm:v5v2-marks}, and
Proposition~\ref{prop:v5v2-entrance}.
\end{proof}

V2 proves the stable-fibre graph theorem and its exact constants.  The
remaining first gate is the marginal base.  V3 will formulate a
mechanism-independent marginal fixed-point theorem using a compact
coupling domain, inward-pointing beta field, and a hyperbolic zero.  It will
separate what follows from degree/topology from the genuinely unproved
existence of a joint SM--gravity beta zero.



\clearpage
\section{V3: a degree certificate and a local ultraviolet stable manifold}
\label{sec:v5v3}

V2 reduced the irrelevant reconstruction to a graph over a prescribed
marginal trajectory.  This note now asks when such a marginal trajectory
exists.  The answer is deliberately split into two statements.  A
finite-dimensional degree condition forces at least one zero of a proposed
beta field, while a quantitative Lyapunov--Perron argument constructs the
trajectories that converge to a hyperbolic zero.  Neither statement proves
that the exact Standard Model--Einstein beta field has the required
properties.  Their purpose is to turn that missing assertion into explicit,
checkable inequalities.

Throughout this note, ultraviolet time is \(t=\log(\mu/\mu_0)\), so increasing
\(t\) means increasing energy.  Let \(G=\mathbb R^d\) be a finite marginal
chart and consider
\begin{equation}
 \dot g(t)=\beta(g(t)).
 \label{eq:v5v3-flow}
\end{equation}
A finite chart is legitimate only after closure of the local operator basis
has been proved.  In particular, the Einstein truncation does not acquire
that property merely by being written in finitely many coordinates.

\subsection{Existence of a zero from an inward boundary field}

Let \(D\Subset G\) be a bounded, open, convex set with \(C^1\) boundary.
Write \(n(g)\) for the outward unit normal at \(g\in\partial D\).

\begin{theorem}[Inward degree certificate]
\label{thm:v5v3-degree}
Suppose \(\beta:\overline D\to G\) is continuous and
\begin{equation}
 \beta(g)\mathbin{\cdot}n(g)<0
 \qquad(g\in\partial D).
 \label{eq:v5v3-inward}
\end{equation}
Then \(\beta\) has a zero in \(D\).  More precisely, for every \(g_0\in D\),
\begin{equation}
 \deg(\beta,D,0)=(-1)^d.
 \label{eq:v5v3-degree-value}
\end{equation}
\end{theorem}

\begin{proof}
Because \(D\) is convex and \(g_0\) is interior, the supporting hyperplane at
\(g\in\partial D\) separates \(g_0\) strictly from \(g\).  Hence
\[
 (g-g_0)\mathbin{\cdot}n(g)>0.
\]
For \(s\in[0,1]\), define
\[
 H_s(g)=(1-s)\beta(g)-s(g-g_0).
\]
Both summands have strictly negative scalar product with \(n(g)\) on
\(\partial D\).  Thus \(H_s(g)\ne0\) there for every \(s\).  Homotopy
invariance of Brouwer degree gives
\[
 \deg(\beta,D,0)
 =\deg\bigl(-(g-g_0),D,0\bigr)
 =\operatorname{sgn}\det(-I)=(-1)^d.
\]
The degree is nonzero, so the existence property of degree supplies
\(g_*\in D\) with \(\beta(g_*)=0\).
\end{proof}

\begin{remark}[What must be estimated]
\label{rem:v5v3-degree-data}
The theorem needs the beta field on an entire codimension-one boundary, not
only a perturbative expansion at one point.  A constructive application
therefore needs a compact coupling domain, a common renormalisation chart,
and error bars strong enough to preserve the strict sign in
\eqref{eq:v5v3-inward}.  Scheme changes that are \(C^1\) diffeomorphisms
transport the zero and the degree, but a truncation error is not such a
diffeomorphism.
\end{remark}

\begin{proposition}[Robustness under a certified error]
\label{prop:v5v3-robust-degree}
Assume that a computable field \(\beta_0\) satisfies
\[
 \beta_0(g)\mathbin{\cdot}n(g)\le-\delta
 \quad(g\in\partial D)
\]
for some \(\delta>0\).  If
\[
 \sup_{g\in\partial D}\|\beta(g)-\beta_0(g)\|<\delta,
 \label{eq:v5v3-degree-error}
\]
then \(\beta\) obeys \eqref{eq:v5v3-inward} and has a zero in \(D\).
\end{proposition}

\begin{proof}
Cauchy--Schwarz and \(\|n(g)\|=1\) give
\[
 \beta(g)\mathbin{\cdot}n(g)
 \le -\delta+\|\beta(g)-\beta_0(g)\|<0.
\]
Apply Theorem~\ref{thm:v5v3-degree}.
\end{proof}

\subsection{A quantitative stable manifold at a certified zero}

Assume now that \(\beta\) is \(C^1\) near a zero \(g_*\).  Put
\[
 x=g-g_*,
 \qquad
 A=D\beta(g_*),
 \qquad
 N(x)=\beta(g_*+x)-Ax.
\]
Suppose that \(G=E^s\oplus E^u\) is an \(A\)-invariant splitting with
projections \(P_s,P_u\), and that for some \(M\ge1\) and \(\alpha>0\),
\begin{align}
 \|e^{tA}P_s\|&\le M e^{-\alpha t}
 &&(t\ge0),
 \label{eq:v5v3-semigroup-s}\\
 \|e^{tA}P_u\|&\le M e^{\alpha t}
 &&(t\le0).
 \label{eq:v5v3-semigroup-u}
\end{align}
Choose \(0<\eta<\alpha\).  After a standard local cut-off, suppose that
\(N(0)=0\) and
\begin{equation}
 \|N(x)-N(y)\|\le L\|x-y\|
 \label{eq:v5v3-N-Lipschitz}
\end{equation}
globally, with
\begin{equation}
 q:=ML\left(\frac1{\alpha-\eta}+\frac1{\alpha+\eta}\right)<1.
 \label{eq:v5v3-q}
\end{equation}
The cut-off changes no trajectory that remains in the chosen local ball.

Let
\[
 \mathcal X_\eta
 =\left\{x\in C([0,\infty);G):
   \|x\|_\eta:=\sup_{t\ge0}e^{\eta t}\|x(t)\|<\infty\right\}.
\]

\begin{theorem}[Lyapunov--Perron ultraviolet entrance]
\label{thm:v5v3-LP}
Under \eqref{eq:v5v3-semigroup-s}--\eqref{eq:v5v3-q}, for every sufficiently
small \(\xi\in E^s\) there is a unique \(x_\xi\in\mathcal X_\eta\) satisfying
\(P_sx_\xi(0)=\xi\) and \(\dot x_\xi=Ax_\xi+N(x_\xi)\).  It is the unique
decaying solution with that stable initial coordinate and obeys
\begin{equation}
 \|x_\xi\|_\eta\le\frac{M}{1-q}\|\xi\|.
 \label{eq:v5v3-LP-bound}
\end{equation}
Its initial value is a graph
\begin{equation}
 x_\xi(0)=\xi+h(\xi),
 \qquad
 h(\xi)=-\int_0^\infty e^{-\tau A}P_uN(x_\xi(\tau))\,d\tau,
 \label{eq:v5v3-stable-graph}
\end{equation}
and the corresponding coupling trajectory converges to \(g_*\) at least as
\(e^{-\eta t}\).
\end{theorem}

\begin{proof}
For fixed \(\xi\), define on \(\mathcal X_\eta\)
\begin{align}
 (\mathcal L_\xi x)(t)
 := {}&e^{tA}\xi
 +\int_0^t e^{(t-\tau)A}P_sN(x(\tau))\,d\tau
 \notag\\
 &-\int_t^\infty e^{(t-\tau)A}P_uN(x(\tau))\,d\tau.
 \label{eq:v5v3-LP-map}
\end{align}
The two integrals converge by
\eqref{eq:v5v3-semigroup-s}--\eqref{eq:v5v3-N-Lipschitz}.
For \(x,y\in\mathcal X_\eta\), the stable integral is bounded in the weighted
norm by
\[
 \frac{ML}{\alpha-\eta}\|x-y\|_\eta,
\]
and the unstable integral by
\[
 \frac{ML}{\alpha+\eta}\|x-y\|_\eta.
\]
Thus \(\mathcal L_\xi\) is a contraction with constant \(q\).  Moreover,
\(\|\mathcal L_\xi0\|_\eta\le M\|\xi\|\).  Banach's theorem gives a unique
fixed point and the bound
\[
 \|x_\xi\|_\eta
 \le M\|\xi\|+q\|x_\xi\|_\eta,
\]
which is \eqref{eq:v5v3-LP-bound}.

Differentiating the two variation-of-constants integrals shows that the fixed
point solves the differential equation.  Applying \(P_s\) at \(t=0\) gives
\(\xi\), while applying \(P_u\) gives
\eqref{eq:v5v3-stable-graph}.  Conversely, any solution in
\(\mathcal X_\eta\) with stable coordinate \(\xi\) has the stable
variation-of-constants formula.  Its unstable formula has a terminal term
\(e^{(t-T)A}P_ux(T)\); the weighted decay and
\eqref{eq:v5v3-semigroup-u} make this term tend to zero as \(T\to\infty\).
It therefore satisfies \eqref{eq:v5v3-LP-map} and equals the fixed point.
Finally, \(\|x_\xi(t)\|\le e^{-\eta t}\|x_\xi\|_\eta\).
For sufficiently small \(\xi\), \eqref{eq:v5v3-LP-bound} keeps the trajectory
inside the ball on which the cut-off agrees with the original field.
\end{proof}

\begin{corollary}[Dimension of the ultraviolet critical surface]
\label{cor:v5v3-critical-dimension}
If \(A\) is hyperbolic, the local set of trajectories converging to \(g_*\)
in ultraviolet time is a graph over \(E^s\).  Its local dimension is
\(\dim E^s\); the complementary \(\dim E^u\) coordinates must be tuned.
If \(N\) is \(C^k\), the graph is \(C^k\).
\end{corollary}

\begin{proof}
The graph and uniqueness follow from
Theorem~\ref{thm:v5v3-LP}.  The usual parameter-dependent contraction
argument differentiates \eqref{eq:v5v3-LP-map} through order \(k\); the
inverse \(1-D_x\mathcal L_\xi\) is the convergent Neumann series because its
norm is at most \(q<1\).  Since \(DN(0)=0\), differentiation at \(\xi=0\)
gives \(Dh(0)=0\).  Hence the graph has tangent space \(E^s\) and the stated
dimension.
\end{proof}

\begin{firewall}[Zero, entrance, and physical matching are distinct]
\label{fw:v5v3-three-gates}
Theorem~\ref{thm:v5v3-degree} proves a zero from a global boundary sign.
Theorem~\ref{thm:v5v3-LP} proves local ultraviolet entrance from hyperbolicity
and a small nonlinear remainder.  Neither theorem proves that the critical
surface intersects the observed infrared coupling surface.  That is a
separate matching problem.
\end{firewall}

\subsection{The inherited hypercharge obstruction in degree form}

The fourth-series analysis already proved, in its stated
Gaussian-dominance wedge, monotone decrease of the positive inverse
hypercharge coupling \(z_Y=g_Y^{-2}\) and finite-depth loss of positivity.
We do not repeat that calculation.  We only record its consequence for the
new degree certificate.

\begin{proposition}[Outward hypercharge face]
\label{prop:v5v3-hypercharge-face}
Let \(D\Subset\{g:z_Y(g)>0\}\) be a bounded convex \(C^1\) domain.  Suppose
\(z_Y\) is one of its affine coordinates and
\[
 \beta_{z_Y}(g)\le-\kappa<0
 \qquad(g\in\overline D).
 \label{eq:v5v3-hypercharge-monotone}
\]
Then \(\beta\) cannot satisfy the inward condition
\eqref{eq:v5v3-inward} on \(\partial D\).
\end{proposition}

\begin{proof}
Let \(a=\min_{\overline D}z_Y\), attained at
\(g_a\in\partial D\).  The supporting hyperplane
\(z_Y=a\) has outward normal \(n(g_a)=-e_Y\), where \(e_Y\) is the unit
vector in the increasing \(z_Y\) direction.  Therefore
\[
 \beta(g_a)\mathbin{\cdot}n(g_a)
 =-\beta_{z_Y}(g_a)\ge\kappa>0,
\]
contradicting \eqref{eq:v5v3-inward}.
\end{proof}

\begin{corollary}[What a joint fixed point must change]
\label{cor:v5v3-joint-change}
Within the wedge of Proposition~\ref{prop:v5v3-hypercharge-face}, neither the
degree certificate nor an all-scale positive hypercharge trajectory is
available.  A joint SM--gravity ultraviolet construction must leave that
wedge, change the exact hypercharge beta field enough to create a zero, or
use a rigorously equivalent chart in which the same zero is visible.
\end{corollary}

\begin{proof}
The first assertion is Proposition~\ref{prop:v5v3-hypercharge-face}; the
finite-depth obstruction is inherited from the fourth-series result.  A
regular change of coordinates transports zeros and therefore cannot turn a
nowhere-zero vector field into one with a zero.  The listed alternatives
exhaust the ways to invalidate the hypotheses of the obstruction.
\end{proof}

\subsection{A checkable marginal fixed-point package}

\begin{definition}[Marginal fixed-point certificate]
\label{def:v5v3-MFPC}
An MFPC consists of:
\begin{enumerate}
 \item a proved closed marginal chart \(G\) and a convex \(C^1\) domain
 \(D\Subset G\);
 \item a beta field \(\beta\) with a normed error estimate on
 \(\overline D\);
 \item either the exact inward sign \eqref{eq:v5v3-inward}, or the robust
 version \eqref{eq:v5v3-degree-error};
 \item an isolated zero \(g_*\), a verified hyperbolic splitting, and
 constants satisfying \eqref{eq:v5v3-semigroup-s}--\eqref{eq:v5v3-q};
 \item a physical infrared surface and a proof that it intersects the local
 stable graph.
\end{enumerate}
\end{definition}

\begin{theorem}[Conditional marginal ancestry]
\label{thm:v5v3-conditional-ancestry}
An MFPC supplies an all-scale marginal trajectory converging to \(g_*\).
If, along that trajectory, the three hypotheses of the V2 fixed-point
entrance certificate hold, then it also supplies the unique twice-marked
irrelevant trajectory converging to the fixed interaction.
\end{theorem}

\begin{proof}
Item 5 selects initial data on the graph constructed by
Theorem~\ref{thm:v5v3-LP}; that theorem gives the all-scale marginal
trajectory and its convergence.  Apply
Theorem~\ref{thm:v5v2-certificate} to reconstruct the irrelevant fibre.
\end{proof}

No MFPC is presently proved for the exact Standard Model--Einstein system.
The finite-dimensional closure in item 1 is itself obstructed by the
inherited curvature-counterterm tower unless a nonperturbative construction
controls that tower.  V3 has nevertheless removed an ambiguity: a proposed
joint fixed point must come with a boundary degree estimate, local spectral
bounds, and an infrared intersection argument.

V4 will address the last of these three tasks.  It will formulate the
infrared matching problem as a finite-dimensional shooting map, prove a
quantitative transversality theorem with certified errors, and identify the
dimension balance required for a discrete predictive match.


\clearpage
\section{V4: certified infrared shooting and transversality}
\label{sec:v5v4}

V3 supplied a conditional ultraviolet critical surface.  A continuum theory
that converges to the ultraviolet fixed point must still reproduce the
infrared data.  This note makes that requirement finite and quantitative.
The ultraviolet critical surface is transported through a finite RG
interval, composed with an observation map, and tested against a physical
constraint surface.  A contraction estimate then certifies a nearby match
from approximate data.  The derivative rank also gives the precise local
dimension count.

Let \(g_*\) be a hyperbolic zero satisfying the hypotheses of
Theorem~\ref{thm:v5v3-LP}.  Write \(s=\dim E^s\), and let
\begin{equation}
 \iota(\xi)=g_*+\xi+h(\xi)
 \qquad(\xi\in U\subset E^s)
 \label{eq:v5v4-iota}
\end{equation}
be its local ultraviolet stable graph.  Fix a finite RG time \(T>0\).
Assume every trajectory under consideration remains in a compact tube
\(K_T\) on the interval \([-T,0]\).  Denote the flow of
\(\dot g=\beta(g)\) by \(\varphi_t\).

Let \(\mathcal O\) be a \(C^1\) observation map from couplings to
\(\mathbb R^m\).  Let the acceptable infrared data form, locally, the
regular level set
\[
 \mathcal M=\{y:C(y)=0\},
 \qquad
 C:\mathbb R^m\to\mathbb R^r,
 \qquad
 \operatorname{rank}DC=r
 \quad\hbox{on }\mathcal M.
\]
The shooting map is
\begin{equation}
 S_T(\xi)
 :=C\!\left(\mathcal O\!\left(\varphi_{-T}(\iota(\xi))\right)\right)
 \in\mathbb R^r.
 \label{eq:v5v4-shooting}
\end{equation}
A physical match is exactly a zero of \(S_T\).

\subsection{Transport through a finite RG interval}

\begin{proposition}[Finite-time transport of the critical surface]
\label{prop:v5v4-transport}
Suppose \(\beta\) is \(C^1\) on a neighbourhood of \(K_T\) and all indicated
solutions remain in that neighbourhood.  Then
\[
 F_T:=\varphi_{-T}\circ\iota
\]
is a \(C^1\) immersion after \(U\) is made sufficiently small.  Its
derivative is
\begin{equation}
 DF_T(\xi)
 =X(-T;\iota(\xi))\bigl(I+Dh(\xi)\bigr),
 \label{eq:v5v4-transport-derivative}
\end{equation}
where the fundamental matrix solves
\[
 \partial_tX(t;g_0)
 =D\beta(\varphi_t(g_0))X(t;g_0),
 \qquad
 X(0;g_0)=I.
\]
\end{proposition}

\begin{proof}
Standard local existence and differentiability for a \(C^1\) vector field
give the derivative formula for the flow.  The matrix \(X(t;g_0)\) is
invertible, because uniqueness gives
\[
 X(-t;\varphi_t(g_0))X(t;g_0)=I.
\]
At \(\xi=0\), Corollary~\ref{cor:v5v3-critical-dimension} gives
\(Dh(0)=0\), so \(I+Dh(0):E^s\to G\) is injective.  Injectivity persists for
small \(\xi\), and multiplication by the invertible fundamental matrix
preserves it.
\end{proof}

\begin{corollary}[Exact shooting derivative]
\label{cor:v5v4-shooting-derivative}
Under Proposition~\ref{prop:v5v4-transport},
\begin{equation}
 DS_T(\xi)
 =DC_{\mathcal O(F_T(\xi))}
  D\mathcal O_{F_T(\xi)}
  X(-T;\iota(\xi))
  \bigl(I+Dh(\xi)\bigr).
 \label{eq:v5v4-DS}
\end{equation}
\end{corollary}

\begin{proof}
Apply the chain rule to \eqref{eq:v5v4-shooting} and use
\eqref{eq:v5v4-transport-derivative}.
\end{proof}

\begin{remark}[Conditioning is part of the certificate]
\label{rem:v5v4-conditioning}
Even when the stable graph is well conditioned near \(g_*\), the norm of
\(X(-T)\) can grow rapidly.  A numerical shooting result without a rigorous
bound on the fundamental matrix does not control the error at the infrared
end.
\end{remark}

\subsection{A posteriori certification of a square match}

We first treat the predictive square case \(s=r\), after choosing linear
coordinates on \(E^s\).  Let
\[
 B_\rho(\xi_0)=\{\xi\in\mathbb R^r:\|\xi-\xi_0\|\le\rho\}.
\]

\begin{theorem}[Quantitative shooting certificate]
\label{thm:v5v4-shooting-certificate}
Let \(S:B_\rho(\xi_0)\to\mathbb R^r\) be \(C^1\), and let
\(A:\mathbb R^r\to\mathbb R^r\) be invertible.  Define
\begin{equation}
 \eta=\|A^{-1}S(\xi_0)\|,
 \qquad
 q=\sup_{\xi\in B_\rho(\xi_0)}
       \|I-A^{-1}DS(\xi)\|.
 \label{eq:v5v4-eta-q}
\end{equation}
If
\begin{equation}
 q<1,
 \qquad
 \eta\le(1-q)\rho,
 \label{eq:v5v4-radii}
\end{equation}
then \(S\) has a unique zero \(\xi_*\) in \(B_\rho(\xi_0)\), and
\begin{equation}
 \|\xi_*-\xi_0\|\le\frac{\eta}{1-q}.
 \label{eq:v5v4-root-bound}
\end{equation}
\end{theorem}

\begin{proof}
Define
\[
 \mathcal N(\xi)=\xi-A^{-1}S(\xi).
\]
The mean-value formula and \eqref{eq:v5v4-eta-q} give
\[
 \|\mathcal N(\xi)-\mathcal N(\zeta)\|
 \le q\|\xi-\zeta\|
\]
on the convex ball.  Moreover,
\[
 \|\mathcal N(\xi)-\xi_0\|
 \le q\|\xi-\xi_0\|+\|\mathcal N(\xi_0)-\xi_0\|
 \le q\rho+\eta\le\rho.
\]
Thus \(\mathcal N\) maps the complete ball into itself and is a strict
contraction.  Its unique fixed point is the unique zero of \(S\) in the
ball.  At the fixed point,
\[
 \|\xi_*-\xi_0\|
 \le q\|\xi_*-\xi_0\|+\eta,
\]
which proves \eqref{eq:v5v4-root-bound}.
\end{proof}

\begin{corollary}[Certificate from an approximate shooting computation]
\label{cor:v5v4-approximate}
Let \(\widehat S_0\in\mathbb R^r\) and an invertible matrix
\(\widehat A\) be computed approximations at \(\xi_0\).  Suppose rigorous
bounds establish
\begin{align}
 \|S(\xi_0)-\widehat S_0\|&\le\varepsilon_0,
 \label{eq:v5v4-value-error}\\
 \sup_{\xi\in B_\rho(\xi_0)}
 \|DS(\xi)-\widehat A\|&\le\varepsilon_1.
 \label{eq:v5v4-derivative-error}
\end{align}
Put
\begin{equation}
 \widehat\eta
 =\|\widehat A^{-1}\|
   \bigl(\|\widehat S_0\|+\varepsilon_0\bigr),
 \qquad
 \widehat q=\|\widehat A^{-1}\|\varepsilon_1.
 \label{eq:v5v4-hat-bounds}
\end{equation}
If
\[
 \widehat q<1,
 \qquad
 \widehat\eta\le(1-\widehat q)\rho,
\]
then \(S\) has exactly one zero in \(B_\rho(\xi_0)\), with
\[
 \|\xi_*-\xi_0\|
 \le\frac{\widehat\eta}{1-\widehat q}.
\]
\end{corollary}

\begin{proof}
With \(A=\widehat A\),
\[
 \|A^{-1}S(\xi_0)\|
 \le\|A^{-1}\|
      \bigl(\|\widehat S_0\|+\varepsilon_0\bigr)
 =\widehat\eta
\]
and
\[
 \|I-A^{-1}DS(\xi)\|
 =\|A^{-1}(A-DS(\xi))\|
 \le\widehat q.
\]
Apply Theorem~\ref{thm:v5v4-shooting-certificate}.
\end{proof}

\subsection{Controlling flow error}

The previous corollary is useful only if the errors in the RG transport are
propagated to \(\varepsilon_0\) and \(\varepsilon_1\).  The value error has a
simple universal estimate.

\begin{proposition}[Finite-time beta-field error]
\label{prop:v5v4-flow-error}
Let \(\beta,\widehat\beta\) be vector fields on the common tube \(K_T\).
Assume
\[
 \|\beta(x)-\beta(y)\|\le L\|x-y\|,
 \qquad
 \sup_{x\in K_T}\|\beta(x)-\widehat\beta(x)\|\le\varepsilon_\beta.
\]
Let \(g,\widehat g\) solve the two flow equations from the same initial
point and remain in \(K_T\) for \(0\le t\le T\).  Then
\begin{equation}
 \|g(t)-\widehat g(t)\|
 \le
 \begin{cases}
 \displaystyle
 \varepsilon_\beta\frac{e^{Lt}-1}{L},&L>0,\\
 \varepsilon_\beta t,&L=0.
 \end{cases}
 \label{eq:v5v4-flow-error}
\end{equation}
The same estimate holds for backward transport after replacing the vector
fields by \(-\beta\) and \(-\widehat\beta\).
\end{proposition}

\begin{proof}
The integral equations imply
\[
 \|g(t)-\widehat g(t)\|
 \le L\int_0^t\|g(\tau)-\widehat g(\tau)\|\,d\tau
     +\varepsilon_\beta t.
\]
Equivalently, the upper right derivative of the distance is at most
\(L\|g-\widehat g\|+\varepsilon_\beta\).  Gronwall's inequality gives
\eqref{eq:v5v4-flow-error}; the \(L=0\) case follows directly.
\end{proof}

\begin{corollary}[Observation error]
\label{cor:v5v4-observation-error}
If \(\mathcal O\) and \(C\) are Lipschitz on the relevant images with
constants \(L_{\mathcal O}\) and \(L_C\), then the contribution of the
beta-field error to the shooting value error is at most
\[
 L_CL_{\mathcal O}\,
 \varepsilon_\beta\frac{e^{LT}-1}{L}
\]
for \(L>0\), with the evident \(L=0\) replacement.
\end{corollary}

\begin{proof}
Compose \eqref{eq:v5v4-flow-error} with the two Lipschitz estimates.
\end{proof}

The derivative error in \eqref{eq:v5v4-derivative-error} additionally
requires interval bounds for \(D\beta\), the fundamental matrix, \(Dh\),
\(D\mathcal O\), and \(DC\).  Equation~\eqref{eq:v5v4-DS} identifies every
factor; omitting any one of them leaves the transversality uncertified.

\subsection{Dimension balance}

\begin{theorem}[Local dimension of the matched set]
\label{thm:v5v4-dimension}
Let \(S_T:U\subset\mathbb R^s\to\mathbb R^r\) be \(C^1\), and suppose
\(DS_T(\xi_*)\) is surjective at every zero under consideration.  Then the
matched set \(S_T^{-1}(0)\) is locally a \(C^1\) manifold of dimension
\begin{equation}
 s-r.
 \label{eq:v5v4-dimension}
\end{equation}
Consequently:
\begin{enumerate}
 \item if \(s=r\), transverse matches are locally isolated;
 \item if \(s>r\), \(s-r\) local parameters remain;
 \item if \(s<r\), transversality to the point \(0\) is impossible at a
 zero, and a generic perturbation of the \(r\) constraints has no local
 match.
\end{enumerate}
\end{theorem}

\begin{proof}
Surjectivity makes \(0\) a regular value.  The regular-value theorem gives
a local manifold whose tangent space is \(\ker DS_T(\xi_*)\).  Rank-nullity
then gives
\[
 \dim\ker DS_T=s-r.
\]
The first two conclusions follow.  If \(s<r\), a linear map from
\(\mathbb R^s\) to \(\mathbb R^r\) cannot be surjective.  The generic
statement is the standard finite-dimensional transversality consequence:
after a sufficiently small generic translation of the target value, the
translated value is regular; its putative inverse image would have negative
dimension and is therefore empty.
\end{proof}

\begin{firewall}[Parameter counting does not prove a match]
\label{fw:v5v4-counting}
The equality \(s=r\) only permits isolated transverse solutions.  Existence
requires a residual-and-derivative estimate such as
\eqref{eq:v5v4-value-error}--\eqref{eq:v5v4-hat-bounds}.  Conversely, a
numerically small residual without a lower singular-value bound for the
shooting derivative can lie near a tangent miss rather than a true root.
\end{firewall}

\subsection{The matching package for the continuum programme}

\begin{definition}[Infrared matching certificate]
\label{def:v5v4-IMC}
An IMC at RG time \(T\) consists of:
\begin{enumerate}
 \item the local stable graph and constants from
 Theorem~\ref{thm:v5v3-LP};
 \item a compact tube proving existence of the backward flow to time
 \(-T\);
 \item a renormalised observation map \(\mathcal O\) and a regular physical
 constraint map \(C\);
 \item interval bounds for every factor of \eqref{eq:v5v4-DS};
 \item a ball, an approximate residual, and an invertible approximate
 derivative satisfying Corollary~\ref{cor:v5v4-approximate}.
\end{enumerate}
\end{definition}

\begin{theorem}[Conditional matched marginal trajectory]
\label{thm:v5v4-matched-trajectory}
In the square case \(s=r\), an IMC supplies a unique marginal trajectory in
the certified ball which converges to \(g_*\) in ultraviolet time and meets
the infrared constraints at time \(-T\).  If the V2 fixed-point entrance
certificate holds along it, there is also a unique twice-marked irrelevant
trajectory above it.
\end{theorem}

\begin{proof}
Corollary~\ref{cor:v5v4-approximate} gives the unique shooting zero.
Proposition~\ref{prop:v5v4-transport} identifies it with a trajectory on the
ultraviolet critical surface, and
Theorem~\ref{thm:v5v3-LP} gives ultraviolet convergence.  Finally apply
Theorem~\ref{thm:v5v2-certificate}.
\end{proof}

No IMC is presently available for the exact Standard Model--Einstein
system, because the joint marginal fixed point, the closed gravitational
chart, and the required nonperturbative error bounds remain open.  V4
nevertheless closes the logical gap between the claim that relevant
parameters can be fitted and a proof: the fit must be a zero of a specified
shooting map with a certified inverse bound.

V5 will perform the required companion-programme audit before choosing the
next bottleneck.  It will inspect the newest Yang--Mills and Navier--Stokes
notes in the shared folder, import only proved reusable mechanisms, and
then decide whether the most productive next step is operator-tower closure
or a construction of the exact joint beta field.



\clearpage
\section{V5: companion audit and essential operator-tower closure}
\label{sec:v5v5}

This is the required five-version companion checkpoint.  Its mathematical
decision is to attack the gravitational local-operator tower before trying
to compute a finite joint beta field.  A beta function on a truncation is
not the beta function of a continuum theory unless omitted essential
coordinates are controlled.  The note first records the audit, then gives
a scheme-invariant single-operator obstruction and an infinite-tail
elimination theorem.

\subsection{Compulsory audit of the shared folder}

At the time of this audit, the newest live companion files were
\begin{align*}
 &\texttt{Renewal\_Yang\_Mills\_Mass\_Gap\_Research\_Notes\_V23.tex},\\
 &\hspace{1em}\text{SHA--256 }
 \texttt{1E8528A10785E0C8E0B0197854ED03B90A1E602992BB002154B66D16B9DEE9A9},
 \\
 &\texttt{Renewal\_NS\_Stability\_Research\_Notes\_V31.tex},\\
 &\hspace{1em}\text{SHA--256 }
 \texttt{F1121B62238AD5491BB7CF03635EA9934942EDF573ED8FAAA9EE79E1D38A6696}.
\end{align*}

Yang--Mills V23 isolates a legal coefficient-one resolved fibre, computes
its complete exchange normalization, and proves that any rescue must occur
inside the final physical output block with a specified
\(d_j^{-1}\) square loss.  It does not infer cancellation from the presence
of other fibres.  Navier--Stokes V31 proves an exact dyadic first-moment
identity and shows that removing a probability mark costs the missing
critical shell weight.  Its high-parent lift returns the original
continuation stock, and a balanced-helicity solution gives an exact
nullspace for the signed hinge mechanism.

\begin{theorem}[V5 audit decision]
\label{thm:v5v5-audit}
The companion mechanisms justify the following rules for the present
programme.
\begin{enumerate}
 \item Test a gravitational truncation with a single essential quotient
 coordinate after all redundant directions have been removed.
 \item Evaluate the complete beta coefficient in that coordinate; an
 individual graph is not a certificate.
 \item Do not replace uniform tail decay by a bounded or marked estimate.
 A truncation-removal theorem needs a compact-image or equivalent weighted
 tail bound.
 \item Keep the resulting obstruction separate from the construction of a
 convergent infinite operator tower.
\end{enumerate}
\end{theorem}

\begin{proof}
Item 1 is the abstract form of the Yang--Mills coefficient-one fibre test,
and item 2 preserves its no-inter-fibre-cancellation logic at the level of
the complete coefficient.  Item 3 is the exact lesson of the
Navier--Stokes mark-removal calculation.  Item 4 distinguishes a no-go for
one finite ansatz from a no-go for every infinite-dimensional completion.
\end{proof}

The audit therefore changes the acquisition order.  We first construct the
essential quotient and state the exact tail-tightness condition.  Only
after that condition is available can the V3 degree and V4 shooting
certificates be applied to finite approximations.

\subsection{Essential local theory space}

Let \(\mathcal L\) be a Banach space of local diffeomorphism- and
gauge-invariant action densities in a fixed regularity and weight class.
Let \(\mathcal R\subset\mathcal L\) be the closed subspace generated by
total divergences and infinitesimal local field redefinitions.  Define the
essential theory space
\begin{equation}
 \mathcal E=\mathcal L/\mathcal R.
 \label{eq:v5v5-essential-space}
\end{equation}
All statements below concern the induced beta vector field
\[
 \mathcal B:\mathcal U\subset\mathcal E\longrightarrow\mathcal E.
\]
Thus an equation-of-motion counterterm is already zero in the space being
tested.

Let \(V_N\subset\mathcal E\) be a finite-dimensional linear ansatz and
\(a_0+V_N\) its affine translate.  An essential read outside the ansatz is
a functional \(\ell\in\mathcal E^*\) satisfying
\begin{equation}
 \ell|_{V_N}=0.
 \label{eq:v5v5-annihilator}
\end{equation}

\begin{theorem}[Single essential-coordinate exit]
\label{thm:v5v5-single-read}
Let \(a(t)\) be a differentiable RG trajectory satisfying
\(\dot a=\mathcal B(a)\) and \(a(0)=a_0\).  If an essential read
\(\ell\) obeys \eqref{eq:v5v5-annihilator} and
\begin{equation}
 b_\ell:=\ell(\mathcal B(a_0))\ne0,
 \label{eq:v5v5-nonzero-read}
\end{equation}
then the trajectory cannot remain in \(a_0+V_N\) on any interval about
zero.  Quantitatively,
\begin{equation}
 \operatorname{dist}\bigl(a(t),a_0+V_N\bigr)
 \ge \frac{|b_\ell|}{2\|\ell\|}|t|
 \label{eq:v5v5-exit-bound}
\end{equation}
for all sufficiently small nonzero \(t\).
\end{theorem}

\begin{proof}
Differentiability gives
\[
 a(t)-a_0=t\mathcal B(a_0)+o(t)
 \quad\text{in }\mathcal E.
\]
For every \(v\in V_N\), the annihilator property gives
\[
 |\ell(a(t)-a_0-v)|
 =|tb_\ell+o(t)|.
\]
The dual norm inequality and an infimum over \(v\) yield
\[
 \operatorname{dist}(a(t),a_0+V_N)
 \ge\frac{|tb_\ell+o(t)|}{\|\ell\|}.
\]
For sufficiently small nonzero \(t\), the numerator is at least
\(\frac12|t b_\ell|\), proving \eqref{eq:v5v5-exit-bound}.
\end{proof}

\begin{proposition}[Scheme covariance of the read]
\label{prop:v5v5-scheme}
Let \(\Psi:\mathcal U\to\mathcal U'\) be a \(C^1\) diffeomorphic change of
renormalisation coordinates.  The transformed beta field is
\[
 \mathcal B'(\Psi(a))=D\Psi(a)\mathcal B(a).
\]
At \(a_0\), define
\[
 \ell'=\ell\circ D\Psi(a_0)^{-1}.
\]
Then
\[
 \ell'(\mathcal B'(\Psi(a_0)))
 =\ell(\mathcal B(a_0)).
 \label{eq:v5v5-scheme-invariant-read}
\]
Moreover, \(\ell'\) annihilates
\(D\Psi(a_0)V_N\).
\end{proposition}

\begin{proof}
Both claims follow by substitution:
\[
 \ell'D\Psi(a_0)=\ell.
\]
Thus the nonzero transverse beta component cannot be removed by a regular
scheme change; only its coordinate representation changes.
\end{proof}

\begin{firewall}[The read must see the complete coefficient]
\label{fw:v5v5-complete-coefficient}
A nonzero contribution from one Feynman graph does not establish
\eqref{eq:v5v5-nonzero-read}, because other graphs may cancel it.  The
single-read certificate applies only to the summed, renormalised beta
coefficient in the essential quotient.
\end{firewall}

\subsection{Distance from a truncated fixed point}

Suppose
\(\mathcal E=H_N\oplus T_N\) with bounded projections \(P_N,Q_N\).
A projected fixed point \(a_N\in H_N\) satisfies
\[
 P_N\mathcal B(a_N)=0,
\]
but its full residual is
\[
 r_N=Q_N\mathcal B(a_N).
 \label{eq:v5v5-tail-residual}
\]

\begin{theorem}[Residual lower bound on fixed-point distance]
\label{thm:v5v5-residual-distance}
Assume \(\mathcal B\) is \(L\)-Lipschitz on a convex neighbourhood
containing \(a_N\) and an exact fixed point \(a_*\).  Then
\begin{equation}
 \|a_*-a_N\|
 \ge\frac{\|r_N\|}{\|Q_N\|L}.
 \label{eq:v5v5-distance-lower}
\end{equation}
For every \(\ell\in\mathcal E^*\) annihilating \(H_N\),
\begin{equation}
 \|a_*-a_N\|
 \ge
 \frac{|\ell(\mathcal B(a_N))|}{\|\ell\|L}.
 \label{eq:v5v5-dual-distance-lower}
\end{equation}
\end{theorem}

\begin{proof}
Since \(\mathcal B(a_*)=0\),
\[
 \|r_N\|
 =\|Q_N(\mathcal B(a_N)-\mathcal B(a_*))\|
 \le\|Q_N\|L\|a_N-a_*\|.
\]
This proves \eqref{eq:v5v5-distance-lower}.  The dual version follows from
\[
 |\ell(\mathcal B(a_N))|
 =|\ell(\mathcal B(a_N)-\mathcal B(a_*))|
 \le\|\ell\|L\|a_N-a_*\|.
\]
\end{proof}

This result reverses a common inference: a zero of the projected beta field
is evidence for a nearby full fixed point only after the omitted residual
is small.  A large single-coordinate read gives a rigorous lower bound
against that interpretation.

\subsection{Nonperturbative elimination of an infinite tail}

A finite truncation can still converge if the entire omitted tower is
solved rather than discarded.  Let
\[
 \mathcal E=H\oplus T,
\]
where \(H\) is finite dimensional and \(T\) is a Banach space of essential
higher operators.  Let \(D\Subset H\) be a compact parameter domain.
Assume the exact tail fixed-point equation is equivalent to
\begin{equation}
 h=\mathcal K(g,h),
 \qquad
 (g,h)\in D\times\overline B_R^T.
 \label{eq:v5v5-tail-equation}
\end{equation}

\begin{theorem}[Exact tail elimination]
\label{thm:v5v5-tail-elimination}
Suppose
\begin{align}
 \|\mathcal K(g,h)\|_T&\le R,
 \label{eq:v5v5-tail-invariant}\\
 \|\mathcal K(g,h)-\mathcal K(g,\widetilde h)\|_T
 &\le q\|h-\widetilde h\|_T
 \quad\text{with }q<1
 \label{eq:v5v5-tail-contraction}
\end{align}
uniformly for \(g\in D\).  Then there is a unique tail
\(h=H(g)\in\overline B_R^T\) solving
\eqref{eq:v5v5-tail-equation}.  If \(\mathcal K\) is continuous, then
\(H\) is continuous.  If it is \(C^1\) and
\(\|D_g\mathcal K\|\le B_g\), then
\begin{equation}
 DH(g)
 =\bigl(I-D_h\mathcal K(g,H(g))\bigr)^{-1}
   D_g\mathcal K(g,H(g)),
 \qquad
 \|DH(g)\|\le\frac{B_g}{1-q}.
 \label{eq:v5v5-tail-derivative}
\end{equation}
\end{theorem}

\begin{proof}
For each \(g\), Banach's fixed-point theorem on the invariant ball gives a
unique \(H(g)\).  For \(g,\widetilde g\),
\begin{align*}
 \|H(g)-H(\widetilde g)\|
 \le{}&q\|H(g)-H(\widetilde g)\|\\
 &+\|\mathcal K(g,H(\widetilde g))
       -\mathcal K(\widetilde g,H(\widetilde g))\|.
\end{align*}
Continuity of \(\mathcal K\), uniformly on the compact parameter set and
the invariant ball where needed, gives continuity of \(H\).
For \(C^1\) dependence, differentiate
\(H(g)=\mathcal K(g,H(g))\).  The inverse in
\eqref{eq:v5v5-tail-derivative} is the Neumann series, whose norm is at
most \((1-q)^{-1}\).
\end{proof}

Let \(P_H\mathcal B\) be the head component.  Define the exact reduced field
\begin{equation}
 \beta_{\mathrm{red}}(g)
 :=P_H\mathcal B(g,H(g)).
 \label{eq:v5v5-reduced-beta}
\end{equation}

\begin{corollary}[Full and reduced zeros]
\label{cor:v5v5-full-reduced}
Under Theorem~\ref{thm:v5v5-tail-elimination}, and assuming
\eqref{eq:v5v5-tail-equation} is equivalent to the vanishing of the tail
beta component,
\[
 \mathcal B(g,h)=0
 \quad\Longleftrightarrow\quad
 h=H(g)\ \text{ and }\ \beta_{\mathrm{red}}(g)=0.
\]
Thus the degree theorem of V3 may be applied to
\(\beta_{\mathrm{red}}\), but not to a field obtained by setting \(h=0\)
without an error estimate.
\end{corollary}

\begin{proof}
The tail equation and uniqueness force \(h=H(g)\).  The remaining equation
is exactly \(\beta_{\mathrm{red}}(g)=0\), and the converse is immediate.
\end{proof}

\subsection{When finite tail projections converge}

Let \(T\) have uniformly bounded finite-rank projections
\(\Pi_N:T\to T\) with \(\Pi_Nh\to h\) for every \(h\in T\).  Define
\(H_N(g)\in\operatorname{Ran}\Pi_N\) by
\begin{equation}
 H_N(g)=\Pi_N\mathcal K(g,H_N(g)).
 \label{eq:v5v5-finite-tail}
\end{equation}
The same contraction constant gives a unique \(H_N(g)\).

\begin{theorem}[Uniform truncation-removal bound]
\label{thm:v5v5-truncation-removal}
Set
\begin{equation}
 \varepsilon_N
 :=\sup_{\substack{g\in D\\ \|h\|_T\le R}}
 \|(I-\Pi_N)\mathcal K(g,h)\|_T.
 \label{eq:v5v5-tail-tightness}
\end{equation}
Then
\begin{equation}
 \sup_{g\in D}\|H(g)-H_N(g)\|_T
 \le\frac{\varepsilon_N}{1-q}.
 \label{eq:v5v5-tail-convergence}
\end{equation}
In particular, \(\varepsilon_N\to0\) proves uniform removal of the
operator truncation.
\end{theorem}

\begin{proof}
For fixed \(g\), insert and subtract
\(\mathcal K(g,H_N(g))\):
\begin{align*}
 \|H(g)-H_N(g)\|
 &\le
 \|\mathcal K(g,H(g))-\mathcal K(g,H_N(g))\|\\
 &\quad+
 \|(I-\Pi_N)\mathcal K(g,H_N(g))\|\\
 &\le q\|H(g)-H_N(g)\|+\varepsilon_N.
\end{align*}
Rearrange and take the supremum over \(g\).
\end{proof}

\begin{proposition}[Compact image supplies tail tightness]
\label{prop:v5v5-compact-image}
If
\[
 \mathscr K
 :=\{\mathcal K(g,h):g\in D,\ \|h\|_T\le R\}
\]
is relatively compact in \(T\), then \(\varepsilon_N\to0\).
\end{proposition}

\begin{proof}
Strong convergence \(\Pi_Ny\to y\) and the uniform boundedness of
\(\Pi_N\) imply uniform convergence on compact subsets.  Explicitly, cover
the compact closure of \(\mathscr K\) by a finite \(\delta\)-net.  Choose
\(N\) so that \(\|(I-\Pi_N)y_i\|<\delta\) for every net point \(y_i\).
Uniform boundedness of \(I-\Pi_N\) controls the displacement from each
\(y\) to its net point.  Letting \(\delta\downarrow0\) proves the claim.
\end{proof}

\begin{proposition}[Bounded contraction is not tail tightness]
\label{prop:v5v5-shift-counterexample}
Let \(T=\ell^1(\mathbb N)\), let \(S\) be the right shift, and set
\(\mathcal K(h)=qSh\) with \(0<q<1\).  Then
\(\mathcal K\) is a strict contraction of every centred ball and has the
unique fixed point \(0\), but
\[
 \sup_{\|h\|_1\le R}\|(I-\Pi_N)\mathcal K(h)\|_1=qR
\]
for every coordinate projection \(\Pi_N\).
\end{proposition}

\begin{proof}
The shift is an isometry, so the contraction and fixed-point claims are
immediate.  For each \(N\), take \(h=Re_N\); then
\((I-\Pi_N)\mathcal K(h)=qRe_{N+1}\).
\end{proof}

The counterexample is the operator-tower version of the NS mark warning:
bounded control may prove existence in the full space while providing no
uniform finite-truncation certificate.  Compactness or an explicit
weighted tail estimate is an independent input.

\subsection{Application boundary for Einstein gravity}

The fourth-series V94 theorem already used the complete
Goroff--Sagnotti coefficient to prove that the cosmological,
Einstein--Hilbert, and curvature-squared ansatz is not closed under
graviton loops.  V5 does not repeat that computation.  In the present
language, the independent on-shell curvature-cubic class supplies an
essential read \(\ell_6\) annihilating the at-most-four-derivative ansatz,
while the known nonzero coefficient gives
\[
 \ell_6(\mathcal B(a_{\mathrm{EH}}))\ne0.
\]
Theorem~\ref{thm:v5v5-single-read} makes the exit explicitly invariant
under regular scheme changes and local field redefinitions.

\begin{firewall}[No no-go for an infinite gravitational completion]
\label{fw:v5v5-infinite-completion}
The single-read theorem excludes the stated finite ansatz.  It does not
exclude a fixed point in a weighted infinite-dimensional essential theory
space.  Such a construction must prove
\eqref{eq:v5v5-tail-invariant},
\eqref{eq:v5v5-tail-contraction}, and
\eqref{eq:v5v5-tail-tightness}, or a stronger substitute, for the exact
gravity--matter shell map.
\end{firewall}

\begin{theorem}[Operator-tower certificate for the MFPC]
\label{thm:v5v5-OTC}
Suppose an essential SM--gravity theory space admits:
\begin{enumerate}
 \item a finite head \(H\), an infinite tail \(T\), and an exact tail map
 satisfying Theorem~\ref{thm:v5v5-tail-elimination};
 \item finite-rank projections satisfying
 \(\varepsilon_N\to0\) in
 \eqref{eq:v5v5-tail-tightness};
 \item a reduced beta field \eqref{eq:v5v5-reduced-beta} satisfying the V3
 inward-degree and hyperbolicity conditions; and
 \item the V4 infrared matching certificate.
\end{enumerate}
Then the marginal fixed-point and matching portions of the MFPC hold in
the full essential theory space, and their finite-tail approximations
converge uniformly at the rate
\(\varepsilon_N/(1-q)\).
\end{theorem}

\begin{proof}
Theorem~\ref{thm:v5v5-tail-elimination} solves the exact infinite tail.
Corollary~\ref{cor:v5v5-full-reduced} lifts every reduced zero to a full
zero.  Theorem~\ref{thm:v5v5-truncation-removal} gives the approximation
rate.  The V3 stable-manifold construction and V4 shooting theorem then
supply entrance and matching under the stated certificates.
\end{proof}

No operator-tower certificate is yet proved for quantum
Standard Model--Einstein shells.  V5 has identified its missing analytic
core without adding another unverified finite truncation: one needs a
weighted essential curvature-word space in which the exact tail map has a
compact image and a strict contraction.

V6 will construct a candidate factorially weighted curvature-word space
and prove an abstract heat-kernel compactness criterion for its tail.  It
will state exactly which uniform graph coefficients are needed and test
whether the inherited graphwise smoothing estimates provide them.


\clearpage
\section{V6: a Gevrey curvature-word space and quantitative tail compactness}
\label{sec:v5v6}

V5 showed that a finite gravitational truncation can be removed only after
the exact higher-operator map has a strict contraction and a uniformly
small projected tail.  This note constructs a candidate Banach scale for
that purpose.  The scale allows factorial graph growth, makes local word
multiplication continuous, and converts one stronger analytic-radius bound
into an explicit compactness and truncation rate.  The final subsection
tests the inherited V96--V97 smoothing theorem against the resulting
hypotheses.

\subsection{The graded essential local space}

Fix spacetime dimension four and a finite list of field species.  For each
grade \(n\ge0\), let \(E_n\) be the normed space spanned by essential local
gauge- and diffeomorphism-invariant scalar words of grade \(n\), modulo
algebraic identities, total divergences, and equation-of-motion
directions.  The grade may count curvature factors together with half the
number of additional covariant derivatives; only additivity under word
multiplication will be used.

\begin{lemma}[Finite multiplicity at fixed grade]
\label{lem:v5v6-finite-grade}
If the field list and the maximal jet order at grade \(n\) are finite, then
\(\dim E_n<\infty\).
\end{lemma}

\begin{proof}
At fixed jet order, a local polynomial is a polynomial of bounded degree in
finitely many jet coordinates and finitely many tensor components.
Therefore the ambient polynomial space is finite dimensional.  Imposing
linear symmetry identities and quotienting by the subspace of total
divergences and equation-of-motion directions produces a quotient of a
finite-dimensional space.
\end{proof}

For \(\sigma\ge0\) and \(\rho>0\), define
\begin{equation}
 \mathcal E_{\rho,\sigma}
 =
 \left\{a=(a_n)_{n\ge0}:
 \|a\|_{\rho,\sigma}
 :=\sum_{n\ge0}\frac{\rho^n}{(n!)^\sigma}
       \|a_n\|_{E_n}<\infty\right\}.
 \label{eq:v5v6-Gevrey-space}
\end{equation}
The case \(\sigma=0\) is analytic in the grade variable; positive
\(\sigma\) permits Gevrey factorial growth of the coefficients.

\begin{proposition}[Completeness]
\label{prop:v5v6-complete}
If every \(E_n\) is finite dimensional, then
\(\mathcal E_{\rho,\sigma}\) is a Banach space.
\end{proposition}

\begin{proof}
The map
\[
 a\longmapsto
 \left(\frac{\rho^n}{(n!)^\sigma}a_n\right)_{n\ge0}
\]
is an isometry onto the \(\ell^1\) direct sum of the finite-dimensional
Banach spaces \(E_n\).  An \(\ell^1\) direct sum of Banach spaces is
complete.
\end{proof}

Assume there are bilinear graded contractions
\[
 \star:E_p\times E_q\longrightarrow E_{p+q}
\]
with
\begin{equation}
 \|u\star v\|_{E_{p+q}}
 \le C_\star\|u\|_{E_p}\|v\|_{E_q}.
 \label{eq:v5v6-graded-product}
\end{equation}
For sequences, set
\[
 (a\star b)_n=\sum_{p+q=n}a_p\star b_q.
\]

\begin{theorem}[Gevrey word algebra]
\label{thm:v5v6-algebra}
Under \eqref{eq:v5v6-graded-product},
\begin{equation}
 \|a\star b\|_{\rho,\sigma}
 \le C_\star\|a\|_{\rho,\sigma}\|b\|_{\rho,\sigma}.
 \label{eq:v5v6-algebra-bound}
\end{equation}
\end{theorem}

\begin{proof}
The factorial inequality
\[
 (p+q)!\ge p!\,q!
\]
implies
\[
 \frac{\rho^{p+q}}{((p+q)!)^\sigma}
 \le
 \frac{\rho^p}{(p!)^\sigma}
 \frac{\rho^q}{(q!)^\sigma}.
\]
Use \eqref{eq:v5v6-graded-product}, sum first over \(p+q=n\), and then
over all \(p,q\).  Tonelli's theorem applies because every term is
nonnegative, and gives \eqref{eq:v5v6-algebra-bound}.
\end{proof}

This algebra estimate is only a bookkeeping theorem.  Covariant
commutators can lower the displayed derivative grade while adding
curvature factors; a concrete shell construction must choose a filtration
for which every such operation is bounded by
\eqref{eq:v5v6-graded-product}.

\subsection{An analytic-radius compact embedding}

Let \(0<\rho_{\rm w}<\rho_{\rm s}\), and let
\(\Pi_Na=(a_0,\ldots,a_N,0,\ldots)\).

\begin{theorem}[Compact radius-gap embedding]
\label{thm:v5v6-compact-embedding}
The inclusion
\begin{equation}
 \mathcal E_{\rho_{\rm s},\sigma}
 \hookrightarrow
 \mathcal E_{\rho_{\rm w},\sigma}
 \label{eq:v5v6-radius-embedding}
\end{equation}
is compact.  More precisely,
\begin{equation}
 \|(I-\Pi_N)a\|_{\rho_{\rm w},\sigma}
 \le
 \left(\frac{\rho_{\rm w}}{\rho_{\rm s}}\right)^{N+1}
 \|a\|_{\rho_{\rm s},\sigma}.
 \label{eq:v5v6-geometric-tail}
\end{equation}
\end{theorem}

\begin{proof}
For \(n>N\),
\[
 \frac{\rho_{\rm w}^n}{(n!)^\sigma}
 =
 \left(\frac{\rho_{\rm w}}{\rho_{\rm s}}\right)^n
 \frac{\rho_{\rm s}^n}{(n!)^\sigma}
 \le
 \left(\frac{\rho_{\rm w}}{\rho_{\rm s}}\right)^{N+1}
 \frac{\rho_{\rm s}^n}{(n!)^\sigma}.
\]
Summing proves \eqref{eq:v5v6-geometric-tail}.  On a bounded strong ball,
the discarded weak tail is therefore uniformly small.  The retained range
of \(\Pi_N\) is finite dimensional by
Lemma~\ref{lem:v5v6-finite-grade}, so every bounded sequence has a
subsequence converging in that range.  A diagonal choice followed by the
uniform tail estimate gives convergence in the weak-radius norm.
\end{proof}

\begin{corollary}[Explicit V5 tail tightness]
\label{cor:v5v6-V5-tail}
If a tail map \(\mathcal K\) satisfies
\begin{equation}
 \sup_{\substack{g\in D\\
       \|h\|_{\rho_{\rm w},\sigma}\le R}}
 \|\mathcal K(g,h)\|_{\rho_{\rm s},\sigma}
 \le M_{\rm s},
 \label{eq:v5v6-strong-image}
\end{equation}
then the V5 error obeys
\begin{equation}
 \varepsilon_N
 \le M_{\rm s}
 \left(\frac{\rho_{\rm w}}{\rho_{\rm s}}\right)^{N+1}.
 \label{eq:v5v6-epsilon}
\end{equation}
If the contraction constant in the weak-radius space is \(q<1\), then
\begin{equation}
 \sup_{g\in D}\|H(g)-H_N(g)\|_{\rho_{\rm w},\sigma}
 \le
 \frac{M_{\rm s}}{1-q}
 \left(\frac{\rho_{\rm w}}{\rho_{\rm s}}\right)^{N+1}.
 \label{eq:v5v6-fixed-tail-rate}
\end{equation}
\end{corollary}

\begin{proof}
Apply \eqref{eq:v5v6-geometric-tail} to every value of \(\mathcal K\), take
the supremum, and then use
Theorem~\ref{thm:v5v5-truncation-removal}.
\end{proof}

\subsection{A componentwise graph-coefficient criterion}

Write
\[
 \mathcal K(g,h)=\bigl(\mathcal K_n(g,h)\bigr)_{n\ge0}.
\]

\begin{theorem}[Gevrey coefficient bound implies compact image]
\label{thm:v5v6-GCB}
Suppose there are constants \(A>0\) and \(R_{\rm c}>\rho_{\rm s}\) such
that
\begin{equation}
 \sup_{\substack{g\in D\\
       \|h\|_{\rho_{\rm w},\sigma}\le R}}
 \|\mathcal K_n(g,h)\|_{E_n}
 \le A\,R_{\rm c}^{-n}(n!)^\sigma
 \qquad(n\ge0).
 \label{eq:v5v6-component-majorant}
\end{equation}
Then \eqref{eq:v5v6-strong-image} holds with
\begin{equation}
 M_{\rm s}
 \le\frac{A}{1-\rho_{\rm s}/R_{\rm c}}.
 \label{eq:v5v6-strong-majorant}
\end{equation}
Consequently the image is relatively compact in
\(\mathcal E_{\rho_{\rm w},\sigma}\), with the explicit tail bound
\eqref{eq:v5v6-epsilon}.
\end{theorem}

\begin{proof}
Multiply \eqref{eq:v5v6-component-majorant} by
\(\rho_{\rm s}^n/(n!)^\sigma\) and sum:
\[
 \|\mathcal K(g,h)\|_{\rho_{\rm s},\sigma}
 \le A\sum_{n\ge0}
       \left(\frac{\rho_{\rm s}}{R_{\rm c}}\right)^n
 =\frac{A}{1-\rho_{\rm s}/R_{\rm c}}.
\]
The compactness and tail conclusions follow from
Theorem~\ref{thm:v5v6-compact-embedding}.
\end{proof}

For a graph expansion
\[
 \mathcal K_n=\sum_{\Gamma\in\mathfrak G_n}
               \mathcal K_{n,\Gamma},
\]
a sufficient, cancellation-free hypothesis is
\begin{equation}
 \sum_{\Gamma\in\mathfrak G_n}
 \sup_{g,h}\|\mathcal K_{n,\Gamma}(g,h)\|_{E_n}
 \le A\,R_{\rm c}^{-n}(n!)^\sigma.
 \label{eq:v5v6-absolute-graph-majorant}
\end{equation}
This is intentionally an absolute sum.  A conditional resummation may be
used instead only if its summation procedure defines the same exact shell
map and preserves the Ward and locality identities.

\begin{firewall}[A fixed-graph constant is not a Gevrey bound]
\label{fw:v5v6-fixed-graph}
For each fixed \(\Gamma\), an estimate with a constant \(C_\Gamma\) gives
no control of the sum in
\eqref{eq:v5v6-absolute-graph-majorant}.  One must control simultaneously
the number of graphs, their symmetry factors, subdivergence forests,
tensor contractions, and the growth of differentiated local jets.
\end{firewall}

\subsection{Canonical suppression and the contraction constant}

Let the inverse-shell linear scaling on the tail be diagonal by grade:
\[
 (\mathcal Dh)_n=\Lambda^{-\delta_n}h_n,
 \qquad
 \Lambda>1,
 \qquad
 \delta_n\ge\delta_*>0.
\]

\begin{lemma}[Canonical tail contraction]
\label{lem:v5v6-canonical}
On every \(\mathcal E_{\rho,\sigma}\),
\begin{equation}
 \|\mathcal D\|\le\Lambda^{-\delta_*}<1.
 \label{eq:v5v6-canonical-contraction}
\end{equation}
\end{lemma}

\begin{proof}
Insert the component formula into
\eqref{eq:v5v6-Gevrey-space} and use
\(\Lambda^{-\delta_n}\le\Lambda^{-\delta_*}\) term by term.
\end{proof}

\begin{theorem}[Nonlinear contraction budget]
\label{thm:v5v6-contraction-budget}
Suppose
\[
 \mathcal K(g,h)=s(g)+\mathcal Dh+\mathcal N(g,h)
\]
on the weak-radius ball, and
\[
 \|\mathcal N(g,h)-\mathcal N(g,\widetilde h)\|_{\rho_{\rm w},\sigma}
 \le L_{\rm nl}\|h-\widetilde h\|_{\rho_{\rm w},\sigma}.
\]
If
\begin{equation}
 q:=\Lambda^{-\delta_*}+L_{\rm nl}<1
 \label{eq:v5v6-contraction-budget}
\end{equation}
and
\begin{equation}
 \sup_{g\in D}\|s(g)\|_{\rho_{\rm w},\sigma}
 \le(1-q)R,
 \label{eq:v5v6-invariant-budget}
\end{equation}
then the weak ball is invariant and \(\mathcal K\) is a contraction there.
If, in addition, the strong-image estimate
\eqref{eq:v5v6-strong-image} holds, the entire V5 tail-elimination and
truncation-removal theorem applies.
\end{theorem}

\begin{proof}
Lemma~\ref{lem:v5v6-canonical} and the nonlinear Lipschitz estimate give
the contraction constant \(q\).  For \(\|h\|\le R\),
\[
 \|\mathcal K(g,h)\|
 \le\|s(g)\|+q\|h\|
 \le(1-q)R+qR=R.
\]
The remaining conclusions are
Theorem~\ref{thm:v5v5-tail-elimination} and
Corollary~\ref{cor:v5v6-V5-tail}.
\end{proof}

\begin{remark}[Where asymptotic safety enters]
\label{rem:v5v6-anomalous}
At a non-Gaussian fixed point, \(\delta_n\) must be replaced by the exact
spectral decay of the linearized inverse shell map.  Canonical dimension
alone does not prove \eqref{eq:v5v6-contraction-budget}; anomalous mixing
can change both the splitting and the gap.
\end{remark}

\subsection{Test against the inherited smoothing results}

The fourth-series V96 weighted rooted Rellich theorem proves compactness for
nonlocal connected kernels from a derivative, exponential, and arity gap.
V97 then proves, for each fixed renormalized graph, the strong nonlocal
kernel bound created by heat-shell edges and exact Taylor remainders.  It
explicitly assigns the extracted Taylor polynomial to the separate local
coordinate.

\begin{theorem}[Scope test for V96--V97]
\label{thm:v5v6-scope-test}
The inherited V96--V97 results do not imply
\eqref{eq:v5v6-component-majorant} for the local curvature-word tower.
They supply neither:
\begin{enumerate}
 \item an all-grade bound on the extracted local Taylor coefficients; nor
 \item an all-graph majorant of the form
 \eqref{eq:v5v6-absolute-graph-majorant}.
\end{enumerate}
They remain applicable to the rooted nonlocal fibre once the local tower
has independently been controlled.
\end{theorem}

\begin{proof}
By construction, V97 places every Taylor jet of a superficially
nonnegative subgraph in the local coordinate before estimating the
remainder.  The heat-edge derivative gain therefore estimates the
complement of that coordinate and gives no decay in the local grade \(n\).
Furthermore, the V97 theorem fixes a graph and produces a constant
\(C_\Gamma\); its own graph-sum firewall leaves the sum over graphs open.
Neither logical input yields items 1 or 2.  The V96 compact embedding still
applies to the nonlocal remainder because its hypotheses concern the
separate rooted kernel norm.
\end{proof}

\begin{corollary}[Current operator-tower status]
\label{cor:v5v6-status}
The candidate space
\(\mathcal E_{\rho,\sigma}\) makes the operator-tower certificate
quantitative, but the exact SM--gravity shell map has not been proved to
act on it.  The missing estimate is the uniform Gevrey majorant
\eqref{eq:v5v6-component-majorant}, together with the nonlinear contraction
budget \eqref{eq:v5v6-contraction-budget}.
\end{corollary}

\begin{proof}
The positive implication is
Theorems~\ref{thm:v5v6-GCB} and
\ref{thm:v5v6-contraction-budget}.  Their hypotheses are not supplied by
Theorem~\ref{thm:v5v6-scope-test}.
\end{proof}

V6 therefore goes upstream rather than applying another finite
truncation.  The local curvature-word problem is an all-order graph-growth
problem, while the rooted nonlocal compactness problem is already
functionally separated.

V7 will build an abstract forest majorant.  It will determine which
factorial exponent \(\sigma\) follows from labelled connected-graph and
forest counts under a uniform per-vertex shell estimate, and will expose
the exact additional input needed to pass from formal perturbation theory
to a genuine shell map.


\clearpage
\section{V7: forest growth, Gevrey order, and the resummation gate}
\label{sec:v5v7}

V6 reduced local operator-tower compactness to a uniform Gevrey bound for
the exact shell map.  Perturbative graph counting can at best provide such
a bound order by order.  This note counts a broad finite-valence graph and
forest class, identifies the resulting Gevrey order, and proves a
Banach-valued Borel--Leroy sufficiency theorem.  The distinction between a
formal Gevrey series and an actual shell map becomes an explicit analytic
continuation hypothesis.

\subsection{A labelled graph and forest count}

Consider perturbative order \(k\).  Assume temporarily:
\begin{enumerate}
 \item there are at most \(Q\) vertex types;
 \item every vertex has at most \(\Delta\) half-edges;
 \item a subtraction forest is determined by a laminar family of vertex
 subsets, up to at most \(D^k\) additional decorations; and
 \item the action expansion supplies the usual factor \(1/k!\).
\end{enumerate}
The finite-valence hypothesis will be removed from the claim boundary for
gravity below; it is used only for the combinatorial theorem.

\begin{lemma}[Pairing bound]
\label{lem:v5v7-pairings}
The number of labelled half-edge contraction patterns at order \(k\) is at
most
\begin{equation}
 (\Delta k)!
 \le \Delta^{\Delta k}(k!)^\Delta.
 \label{eq:v5v7-pairing-bound}
\end{equation}
After the factor \(1/k!\), their total absolute combinatorial weight is at
most
\[
 \Delta^{\Delta k}(k!)^{\Delta-1}.
\]
\end{lemma}

\begin{proof}
There are at most \(\Delta k\) labelled half-edges, and every pairing or
partial pairing can be encoded by an ordering of those half-edges, giving
the crude upper bound \((\Delta k)!\).  The multinomial coefficient
\[
 \frac{(\Delta k)!}{(k!)^\Delta}
\]
is one term in the expansion of \(\Delta^{\Delta k}\), hence is no larger
than that expansion.  This proves \eqref{eq:v5v7-pairing-bound}; division by
\(k!\) gives the second statement.
\end{proof}

A family \(\mathfrak F\) of nonempty subsets of
\(\{1,\ldots,k\}\) is laminar when any two members are disjoint or one
contains the other.

\begin{lemma}[Laminar forest bound]
\label{lem:v5v7-laminar}
The number of laminar families on \(k\) labelled elements is at most
\begin{equation}
 8^k k!.
 \label{eq:v5v7-laminar-bound}
\end{equation}
\end{lemma}

\begin{proof}
Adjoin the full set and all singletons to a laminar family.  Its inclusion
poset is a rooted hierarchy with the \(k\) labels as leaves.  Refine every
node with more than two children by a fixed deterministic binary rule.
The result is a rooted binary tree with labelled leaves, together with a
set of internal edges whose contraction recovers the original hierarchy.
The number of rooted binary trees with \(k\) labelled leaves is
\[
 (2k-3)!!\le 2^{k-1}(k-1)!.
\]
Such a tree has fewer than \(2k\) edges, so there are fewer than \(2^{2k}\)
edge subsets.  The deterministic refinement makes the hierarchy
recoverable from this pair.  Hence the number is at most
\[
 2^{k-1}(k-1)!\,2^{2k}\le8^k k!.
\]
Removing the sets adjoined at the first step cannot increase the count.
\end{proof}

\begin{theorem}[Absolute forest-count majorant]
\label{thm:v5v7-forest-count}
Under assumptions 1--4, the total absolute graph-and-forest
combinatorial weight at order \(k\) is at most
\begin{equation}
 C_{\rm comb}^k(k!)^\Delta,
 \qquad
 C_{\rm comb}=8QD\Delta^\Delta.
 \label{eq:v5v7-comb-majorant}
\end{equation}
Graph symmetry factors can only decrease this bound.
\end{theorem}

\begin{proof}
Choose the vertex types in at most \(Q^k\) ways, use
Lemma~\ref{lem:v5v7-pairings} after the action factor, and use
Lemma~\ref{lem:v5v7-laminar} with the additional \(D^k\) decorations.
Multiplication gives
\[
 Q^kD^k\Delta^{\Delta k}(k!)^{\Delta-1}
 \,8^k k!
 =
 C_{\rm comb}^k(k!)^\Delta.
\]
\end{proof}

\begin{firewall}[Scope of the hierarchy encoding]
\label{fw:v5v7-forest-scope}
The forest bound assumes that the divergent subgraph is fixed by its vertex
set up to exponentially many decorations.  If the renormalisation
formalism permits superexponentially many edge-decorated subgraphs with the
same vertex set, that extra growth must be added to
\eqref{eq:v5v7-comb-majorant}.
\end{firewall}

\subsection{From graph counting to a formal Banach-valued majorant}

Let
\[
 X_{\rm s}=\mathcal E_{\rho_{\rm s},\sigma_{\rm loc}}
\]
be the strong curvature-word space of V6.  Write the formal shell map at a
scalar bookkeeping coupling \(z\) as
\begin{equation}
 \mathcal K(z;g,h)
 \sim\sum_{k\ge0}z^k\mathcal K^{(k)}(g,h).
 \label{eq:v5v7-formal-shell}
\end{equation}

\begin{theorem}[Per-graph estimate gives a Gevrey perturbation series]
\label{thm:v5v7-formal-Gevrey}
In addition to Theorem~\ref{thm:v5v7-forest-count}, suppose every decorated
renormalized graph at order \(k\) obeys, uniformly on the parameter and
tail balls,
\begin{equation}
 \|\mathcal K_{\Gamma,\mathfrak F}^{(k)}(g,h)\|_{X_{\rm s}}
 \le A B^k.
 \label{eq:v5v7-per-graph}
\end{equation}
Then the absolutely assembled coefficient satisfies
\begin{equation}
 \sup_{g,h}\|\mathcal K^{(k)}(g,h)\|_{X_{\rm s}}
 \le A C^k(k!)^\Delta,
 \qquad
 C=BC_{\rm comb}.
 \label{eq:v5v7-order-Gevrey}
\end{equation}
It also obeys the looser standard Gevrey bound
\begin{equation}
 \sup_{g,h}\|\mathcal K^{(k)}(g,h)\|_{X_{\rm s}}
 \le A C^k\Gamma(1+\Delta k).
 \label{eq:v5v7-Gamma-Gevrey}
\end{equation}
\end{theorem}

\begin{proof}
Sum \eqref{eq:v5v7-per-graph} over the absolute combinatorial bound
\eqref{eq:v5v7-comb-majorant}.  This gives
\eqref{eq:v5v7-order-Gevrey}.  The multinomial coefficient
\[
 \frac{(\Delta k)!}{(k!)^\Delta}
\]
is an integer at least one, so
\((k!)^\Delta\le(\Delta k)!=\Gamma(1+\Delta k)\), proving
\eqref{eq:v5v7-Gamma-Gevrey}.
\end{proof}

\begin{proposition}[The majorant alone gives no convergent coupling series]
\label{prop:v5v7-zero-radius}
For every integer \(p\ge1\) and \(C>0\), the scalar series
\[
 \sum_{k\ge0}C^k\Gamma(1+pk)z^k
\]
has radius of convergence zero.  Consequently the upper bound
\eqref{eq:v5v7-Gamma-Gevrey} does not imply convergence of
\eqref{eq:v5v7-formal-shell} at any nonzero \(z\).
\end{proposition}

\begin{proof}
The ratio of successive scalar coefficients is
\[
 C\frac{\Gamma(1+p(k+1))}{\Gamma(1+pk)}
 =
 C\prod_{j=1}^{p}(pk+j),
\]
which tends to infinity.  The ratio test gives zero radius.  An upper bound
by a zero-radius majorant is insufficient to establish convergence; it is
compatible both with convergent coefficients and with coefficients that
saturate the displayed divergent series.
\end{proof}

\subsection{Banach-valued Borel--Leroy reconstruction}

For an integer \(p\ge1\), define the order-\(p\) Borel--Leroy transform of
the formal coefficients by
\begin{equation}
 \widehat{\mathcal K}_p(\tau;g,h)
 =
 \sum_{k\ge0}
 \frac{\mathcal K^{(k)}(g,h)}
      {\Gamma(1+pk)}\tau^k.
 \label{eq:v5v7-Borel-transform}
\end{equation}
By \eqref{eq:v5v7-Gamma-Gevrey}, this series converges in \(X_{\rm s}\) for
\(|\tau|<C^{-1}\).

\begin{theorem}[Uniform Borel--Leroy sufficiency]
\label{thm:v5v7-Borel-sufficiency}
Let \(p\ge1\).  Suppose
\(\widehat{\mathcal K}_p\) has a common analytic continuation from
\(|\tau|<C^{-1}\) to a sector containing the positive real axis and,
uniformly for \(g\in D\) and
\(\|h\|_{\rho_{\rm w},\sigma_{\rm loc}}\le R\), satisfies
\begin{equation}
 \|\widehat{\mathcal K}_p(\tau;g,h)\|_{X_{\rm s}}
 \le M e^{a|\tau|^{1/p}}
 \qquad(\tau\ge0).
 \label{eq:v5v7-Borel-growth}
\end{equation}
For positive \(z\) with \(a z^{1/p}<1\), the Bochner integral
\begin{equation}
 \mathcal K_{\rm BL}(z;g,h)
 =
 \int_0^\infty e^{-s}
 \widehat{\mathcal K}_p(zs^p;g,h)\,ds
 \label{eq:v5v7-Borel-Laplace}
\end{equation}
converges in \(X_{\rm s}\) and obeys
\begin{equation}
 \sup_{g,h}
 \|\mathcal K_{\rm BL}(z;g,h)\|_{X_{\rm s}}
 \le\frac{M}{1-a z^{1/p}}.
 \label{eq:v5v7-BL-strong}
\end{equation}
Its asymptotic expansion at \(z=0\) in the sector is
\eqref{eq:v5v7-formal-shell}.
\end{theorem}

\begin{proof}
Along the positive ray,
\[
 \|\widehat{\mathcal K}_p(zs^p;g,h)\|_{X_{\rm s}}
 \le M e^{a z^{1/p}s}.
\]
The integrand is dominated by
\(M e^{-(1-a z^{1/p})s}\), which is integrable and gives
\eqref{eq:v5v7-BL-strong}.  For a monomial in
\eqref{eq:v5v7-Borel-transform},
\[
 \int_0^\infty e^{-s}
 \frac{(zs^p)^k}{\Gamma(1+pk)}\,ds
 =
 z^k,
\]
because the numerator integral is \(\Gamma(1+pk)\).
Taylor's formula for the analytic Borel transform, followed by the same
exponential domination on the remainder, gives the asserted Gevrey
asymptotic expansion uniformly on smaller subsectors.
\end{proof}

\begin{corollary}[Borel input to curvature-tail compactness]
\label{cor:v5v7-Borel-compact}
Under Theorem~\ref{thm:v5v7-Borel-sufficiency}, the Borel--Leroy summed map
has a relatively compact image in
\(\mathcal E_{\rho_{\rm w},\sigma_{\rm loc}}\), and
\begin{equation}
 \sup_{g,h}
 \|(I-\Pi_N)\mathcal K_{\rm BL}(z;g,h)\|_{\rho_{\rm w},\sigma_{\rm loc}}
 \le
 \frac{M}{1-a z^{1/p}}
 \left(\frac{\rho_{\rm w}}{\rho_{\rm s}}\right)^{N+1}.
 \label{eq:v5v7-Borel-tail}
\end{equation}
\end{corollary}

\begin{proof}
Use \eqref{eq:v5v7-BL-strong} as the strong-image bound in
Corollary~\ref{cor:v5v6-V5-tail}.
\end{proof}

\begin{firewall}[Positive-ray continuation is new information]
\label{fw:v5v7-Borel-gap}
The local series \eqref{eq:v5v7-Borel-transform} and its disk of convergence
do not imply analytic continuation along the positive ray, the growth bound
\eqref{eq:v5v7-Borel-growth}, or equality of the Borel sum with a
nonperturbatively defined shell integral.  Singularities on the ray and
contour ambiguities are not resolved by factorial coefficient estimates.
\end{firewall}

\subsection{Derivative and Ward requirements}

Tail elimination needs a contraction, not merely a bounded value map.
Let the Borel transform also be differentiable in \(h\).

\begin{proposition}[Borel derivative budget]
\label{prop:v5v7-Borel-derivative}
If, on the same ray,
\[
 \|D_h\widehat{\mathcal K}_p(\tau;g,h)\|
 _{\mathcal E_{\rho_{\rm w},\sigma_{\rm loc}}
   \to\mathcal E_{\rho_{\rm w},\sigma_{\rm loc}}}
 \le M_1e^{a|\tau|^{1/p}},
\]
then differentiation under \eqref{eq:v5v7-Borel-Laplace} is valid and
\begin{equation}
 \|D_h\mathcal K_{\rm BL}(z;g,h)\|
 \le\frac{M_1}{1-a z^{1/p}}.
 \label{eq:v5v7-Borel-derivative-bound}
\end{equation}
In particular, the V5 contraction holds if the right side is strictly
less than one and the invariant-ball source inequality also holds.
\end{proposition}

\begin{proof}
The same integrable majorant used in
Theorem~\ref{thm:v5v7-Borel-sufficiency} applies to the derivative.
The standard dominated differentiation theorem for Bochner integrals gives
the derivative formula and the bound.
\end{proof}

Every coefficient and its Borel continuation must be assembled in a
Ward-complete, diffeomorphism-covariant local quotient.  Summing
gauge-fixed graph norms first and imposing the identities afterward would
not establish that the Borel map descends to the essential theory space.

\subsection{Why the count does not yet apply to Einstein gravity}

The Einstein--Hilbert expansion around a background contains vertices of
unbounded graviton valence.  Therefore no fixed \(\Delta\) from
Theorem~\ref{thm:v5v7-forest-count} applies to the complete expansion.
A replacement is possible only if the valence-\(r\) vertex tensors admit
an analytic generating bound strong enough that the sum over all
half-edge counts contributes at most another Gevrey factor.

\begin{definition}[Gravity forest-resummation certificate]
\label{def:v5v7-GFRC}
A GFRC consists of:
\begin{enumerate}
 \item a background-covariant vertex generating-function bound covering
 all graviton valences;
 \item a forest formula whose decorated hierarchy count obeys a specified
 Gevrey order \(p\);
 \item uniform strong curvature-grade bounds for the complete coefficients
 and their first two parameter derivatives;
 \item positive-ray analytic continuation and exponential-type estimates
 for their Borel--Leroy transforms;
 \item equality of the resulting sum with a nonperturbative,
 Ward-complete shell map on an invariant ball.
\end{enumerate}
\end{definition}

\begin{theorem}[Conditional passage from forests to the operator tower]
\label{thm:v5v7-GFRC}
If a GFRC gives the value and tail-derivative bounds of
Theorem~\ref{thm:v5v7-Borel-sufficiency} and
Proposition~\ref{prop:v5v7-Borel-derivative}, with derivative norm below
one and an invariant-ball estimate, then the V5 operator-tail elimination
holds and finite curvature truncations converge at the rate
\eqref{eq:v5v7-Borel-tail}.
\end{theorem}

\begin{proof}
Theorem~\ref{thm:v5v7-Borel-sufficiency} supplies the strong image.
Proposition~\ref{prop:v5v7-Borel-derivative} supplies the contraction.
Corollary~\ref{cor:v5v7-Borel-compact} supplies tail tightness.  Apply
Theorems~\ref{thm:v5v5-tail-elimination} and
\ref{thm:v5v5-truncation-removal}.
\end{proof}

No GFRC is presently proved.  The inherited V97 fixed-graph estimate does
not give a uniform valence majorant, an all-forest Gevrey bound, or
positive-ray Borel continuation.  The real Euclidean Einstein contour also
cannot supply item 5 through a positive action integral because of the
inherited conformal-factor instability.

V7 has thus identified a genuine upstream fork.  One route proves a
Ward-complete Borel summation theorem with a specified gravitational
contour.  The other constructs the shell map directly from a stable
nonperturbative metric measure and derives the Gevrey estimates afterward.

V8 will analyze the second route at finite cutoff.  It will prove a general
contour-deformation criterion preserving finite-dimensional Gaussian
integrals and reflection conjugation, then test the conformal mode against
it.  This addresses the metric-integration gate needed before a
nonperturbative shell map can be compared with its forest expansion.


\clearpage
\section{V8: finite-cutoff metric contours and the conformal rotation test}
\label{sec:v5v8}

V7 showed that a nonperturbative metric shell map is needed either to
justify the Borel--Leroy continuation or to replace it.  The inherited V98
result excludes an action-derived probability bound on the real Euclidean
Einstein contour because the conformal action is unbounded below.  This
note asks whether a complex contour can repair that problem.  It proves
finite-dimensional contour and reflection criteria, evaluates the
indefinite Gaussian exactly, and then shows that the uniform imaginary
conformal rotation is only quadratically stable: the full nonlinear
Einstein conformal action has no cutoff-uniform coercive bound on that
contour.

\subsection{When a contour deformation preserves an integral}

Let \(\Omega\subset\mathbb C^N\) be open, let \(S\) and \(F\) be
holomorphic, and let \(\Gamma_s\), \(0\le s\le1\), be oriented real
\(N\)-cycles in \(\Omega\).  Write
\[
 I_s(F)=\int_{\Gamma_s}F(z)e^{-S(z)}\,dz_1\wedge\cdots\wedge dz_N.
\]

\begin{theorem}[Admissible contour homotopy]
\label{thm:v5v8-contour-homotopy}
Suppose the portion of the homotopy between
\(\Gamma_0\cap B_R\) and \(\Gamma_1\cap B_R\) forms an
\((N+1)\)-chain whose additional boundary lies on \(\partial B_R\), and
suppose
\begin{equation}
 \lim_{R\to\infty}
 \int_{\text{homotopy boundary on }\partial B_R}
 |F(z)e^{-S(z)}|\,|dz|=0.
 \label{eq:v5v8-boundary-decay}
\end{equation}
If both endpoint integrals are absolutely convergent, then
\begin{equation}
 I_0(F)=I_1(F).
 \label{eq:v5v8-contour-equality}
\end{equation}
\end{theorem}

\begin{proof}
The holomorphic \(N\)-form
\[
 \omega=F(z)e^{-S(z)}\,dz_1\wedge\cdots\wedge dz_N
\]
is closed.  Stokes' theorem on the truncated homotopy chain gives the
difference of the two truncated endpoint integrals as the integral of
\(\omega\) over the boundary at radius \(R\).  This tends to zero by
\eqref{eq:v5v8-boundary-decay}.  Absolute convergence permits passage to
the endpoint limits.
\end{proof}

\begin{firewall}[Rotation of a divergent contour is a new prescription]
\label{fw:v5v8-divergent-start}
If the real endpoint integral diverges, Theorem~\ref{thm:v5v8-contour-homotopy}
does not compare it with a rotated contour.  The rotated integral may
define an analytic continuation, but equality with a pre-existing real
measure has not been proved.
\end{firewall}

\subsection{The indefinite Gaussian}

Let \(A\) be a real symmetric nonsingular matrix on
\(\mathbb R^N=E_+\oplus E_-\), positive definite on \(E_+\) and negative
definite on \(E_-\).  Put \(q=\dim E_-\) and define the oriented contour
\[
 \Gamma_A=E_+\oplus iE_-.
\]

\begin{theorem}[Exact conformally rotated Gaussian]
\label{thm:v5v8-indefinite-Gaussian}
The integral
\begin{equation}
 Z_A=\int_{\Gamma_A}\exp\left(-\frac12z^TAz\right)\,dz
 \label{eq:v5v8-Gaussian-Z}
\end{equation}
is absolutely convergent and equals
\begin{equation}
 Z_A=i^q(2\pi)^{N/2}|\det A|^{-1/2},
 \label{eq:v5v8-Gaussian-value}
\end{equation}
for the product orientation.  Its normalized two-point function in the
original complex coordinate is
\begin{equation}
 \langle zz^T\rangle_A=A^{-1}.
 \label{eq:v5v8-Gaussian-covariance}
\end{equation}
\end{theorem}

\begin{proof}
Choose an orthogonal basis in which
\[
 A=\operatorname{diag}(\lambda_1,\ldots,\lambda_{N-q},
                       -\mu_1,\ldots,-\mu_q)
\]
with all \(\lambda_i,\mu_j>0\).  Parameterize the contour by
\(z=(x,iy)\).  Then
\[
 z^TAz=\sum_i\lambda_ix_i^2+\sum_j\mu_jy_j^2,
 \qquad
 dz=i^q\,dx\,dy.
\]
The product of ordinary one-dimensional Gaussian integrals proves
\eqref{eq:v5v8-Gaussian-value}.  Differentiating the Gaussian generating
function
\[
 Z_A(J)=Z_A\exp\left(\frac12J^TA^{-1}J\right)
\]
twice at \(J=0\) gives \eqref{eq:v5v8-Gaussian-covariance}.
\end{proof}

\begin{corollary}[The rotated negative coordinate is not a real positive
Gaussian field]
\label{cor:v5v8-negative-covariance}
For a nonzero real vector \(v\in E_-\),
\begin{equation}
 \langle(v^Tz)^2\rangle_A=v^TA^{-1}v<0.
 \label{eq:v5v8-negative-variance}
\end{equation}
Thus the original real coordinate \(v^Tz\) cannot have the law of a real
random variable under a positive measure.
\end{corollary}

\begin{proof}
The restriction of \(A^{-1}\) to \(E_-\) is negative definite.  A square
of a real random variable has nonnegative expectation, contradicting
\eqref{eq:v5v8-negative-variance}.
\end{proof}

The parameter \(y=-iz\) on the rotated negative sector does have positive
Gaussian covariance.  It is a different reality assignment.  Recovering a
physical real metric therefore requires a specified involution and
observable algebra, not only the convergent integral
\eqref{eq:v5v8-Gaussian-Z}.

\subsection{Reflection conjugation and its positivity test}

Let \(\vartheta:\mathbb C^N\to\mathbb C^N\) be an antilinear involution and
define
\[
 F^\vartheta(z)=\overline{F(\vartheta z)}.
\]

\begin{proposition}[Reality from an invariant contour]
\label{prop:v5v8-reality}
Suppose \(\vartheta\Gamma=\Gamma\) with the compatible normalized
orientation, and
\begin{equation}
 S(\vartheta z)=\overline{S(z)}
 \qquad(z\in\Gamma).
 \label{eq:v5v8-action-reality}
\end{equation}
If the integrals converge absolutely, then
\begin{equation}
 \overline{\langle F\rangle_\Gamma}
 =\langle F^\vartheta\rangle_\Gamma.
 \label{eq:v5v8-expectation-reality}
\end{equation}
\end{proposition}

\begin{proof}
Complex conjugate the numerator and denominator, then change variables
\(z\mapsto\vartheta z\).  The action condition and contour invariance
return the original normalized integral with \(F^\vartheta\).
\end{proof}

Reality is weaker than reflection positivity.  Let \(P_+\) select the
positive-time coordinates and let the linear part of Euclidean reflection
be \(R\).  For a centered Gaussian with covariance \(C\), the
Osterwalder--Schrader form on linear observables is represented by
\[
 C_{\rm OS}=P_+RC P_+^*.
\]

\begin{proposition}[Necessary linear reflection test]
\label{prop:v5v8-linear-OS}
If a centered Gaussian functional is reflection positive, then
\begin{equation}
 \overline u^{\,T}C_{\rm OS}u\ge0
 \qquad\text{for every positive-time vector }u.
 \label{eq:v5v8-linear-OS-test}
\end{equation}
A reflection-fixed negative Gaussian mode with \(R=1\) and covariance
\(-\mu^{-1}\) violates this condition.
\end{proposition}

\begin{proof}
Apply reflection positivity to the linear cylinder observable
\(F_u(z)=u^TP_+z\).  Its reflected square has expectation
\(\overline u^{\,T}P_+RCP_+^*u\), which proves the necessary condition.
For the stated one-dimensional mode this value is
\(-|u|^2/\mu<0\).
\end{proof}

\begin{firewall}[The conformal coordinate may be unphysical]
\label{fw:v5v8-conformal-physical}
Proposition~\ref{prop:v5v8-linear-OS} tests the raw conformal coordinate.
It does not exclude reflection positivity after a proved gauge/constraint
quotient that removes this coordinate, or after a new involution whose
fixed observable algebra reconstructs the intended Lorentzian metric.
Those operations must themselves be part of the continuum certificate.
\end{firewall}

\subsection{Interacting contour coercivity and Ward identities}

\begin{theorem}[Finite-cutoff integrability]
\label{thm:v5v8-coercive-contour}
Let \(\Gamma\) be parameterized by a proper \(C^1\) map
\(z:\mathbb R^N\to\mathbb C^N\) whose Jacobian grows at most polynomially.
If
\begin{equation}
 \operatorname{Re}S(z(y))
 \ge c\|y\|^\alpha-C
 \qquad(c,\alpha>0),
 \label{eq:v5v8-coercivity}
\end{equation}
then \(e^{-S}dz\) is absolutely integrable on \(\Gamma\), as are all
polynomial observables.
\end{theorem}

\begin{proof}
The absolute value of the integrand is
\(e^{-\operatorname{Re}S}\) times the Jacobian.  By
\eqref{eq:v5v8-coercivity}, it is bounded by a polynomial times
\(e^{-c\|y\|^\alpha+C}\), which is integrable on \(\mathbb R^N\).
Multiplication by another polynomial preserves integrability.
\end{proof}

Let \(V\) be a holomorphic vector field generating an infinitesimal gauge
or diffeomorphism change.

\begin{proposition}[Contour Ward criterion]
\label{prop:v5v8-Ward}
Assume \(V\) is tangent to \(\Gamma\), its flow preserves the contour
orientation, and the boundary term at infinity vanishes.  Then
\begin{equation}
 \int_\Gamma
 \operatorname{div}\!\left(VFe^{-S}\right)\,dz=0,
 \label{eq:v5v8-Ward-divergence}
\end{equation}
which yields the usual finite-dimensional Ward integration-by-parts
identity.  If \(V\) is not tangent, an additional contour-displacement term
must be evaluated.
\end{proposition}

\begin{proof}
Integrate the total derivative along the flow of \(V\) on the truncated
contour.  Tangency makes it an intrinsic divergence, and the assumed
boundary decay removes the boundary integral.  A transverse component
moves the integration cycle and is not represented by an intrinsic
divergence.
\end{proof}

\subsection{The nonlinear Einstein conformal test}

Use the Euclidean sign convention
\[
 S_E(g)=-\kappa\int_M\sqrt g\,R(g),
 \qquad \kappa>0.
\]
On a closed four-manifold, write \(g=e^{2\phi}\bar g\).  The conformal
transformation formula and an integration by parts give
\begin{equation}
 S_E(e^{2\phi}\bar g)
 =
 -\kappa\int_M\sqrt{\bar g}\,e^{2\phi}
 \left(\bar R+6|\bar\nabla\phi|^2\right).
 \label{eq:v5v8-conformal-action}
\end{equation}
For real \(\phi\), the gradient term has the inherited wrong sign.

On the uniform imaginary rotation \(\phi=i\psi\),
\begin{equation}
 \operatorname{Re}S_E(e^{2i\psi}\bar g)
 =
 -\kappa\int_M\sqrt{\bar g}\,
       \cos(2\psi)\bar R
 +6\kappa\int_M\sqrt{\bar g}\,
       \cos(2\psi)|\bar\nabla\psi|^2.
 \label{eq:v5v8-rotated-real-action}
\end{equation}
Near \(\psi=0\), the gradient quadratic form is positive.  Globally,
\(\cos(2\psi)\) changes sign.

\begin{theorem}[No cutoff-uniform coercivity on the straight imaginary
conformal contour]
\label{thm:v5v8-no-uniform-coercivity}
On a flat four-torus, let the finite cutoff contain the Fourier mode
\(\sin(Nx_1)\).  Fix \(0<\varepsilon<\pi/8\) and set
\[
 \psi_N(x)=\frac{\pi}{2}+\varepsilon\sin(Nx_1).
 \label{eq:v5v8-bad-rotated-mode}
\]
Then
\begin{equation}
 \operatorname{Re}S_E(e^{2i\psi_N}\bar g)
 \le-c_0\kappa\varepsilon^2N^2\operatorname{Vol}(M)
 \label{eq:v5v8-negative-N2}
\end{equation}
for a constant \(c_0>0\) independent of \(N\).  Consequently the straight
rotation \(\phi=i\psi\) has no lower bound uniform in the ultraviolet
cutoff of the form required by
\eqref{eq:v5v8-coercivity}.
\end{theorem}

\begin{proof}
Flatness makes \(\bar R=0\).  Since
\[
 \cos(2\psi_N)
 =-\cos\bigl(2\varepsilon\sin(Nx_1)\bigr)
 \le-\cos(2\varepsilon)<0,
\]
equation \eqref{eq:v5v8-rotated-real-action} gives
\[
 \operatorname{Re}S_E
 \le
 -6\kappa\cos(2\varepsilon)
 \int_M|\bar\nabla\psi_N|^2.
\]
The last integral is
\(\frac12\varepsilon^2N^2\operatorname{Vol}(M)\) in the standard torus
normalization.  Take
\(c_0=3\cos(2\varepsilon)\).
\end{proof}

\begin{corollary}[Quadratic stabilization is insufficient]
\label{cor:v5v8-quadratic-insufficient}
The conventional imaginary rotation stabilizes the conformal Hessian near
a background but does not by itself provide a cutoff-uniform interacting
metric measure, a positive-axis Borel prescription, or the MOSC required
by the fourth-series assembly theorem.
\end{corollary}

\begin{proof}
The Gaussian statement is
Theorem~\ref{thm:v5v8-indefinite-Gaussian}.  The nonlinear uniform
coercivity required for
Theorem~\ref{thm:v5v8-coercive-contour} fails by
Theorem~\ref{thm:v5v8-no-uniform-coercivity}.  Reflection positivity also
requires the independent test of
Proposition~\ref{prop:v5v8-linear-OS}; Borel equality and MOSC do not
follow from either quadratic fact.
\end{proof}

\subsection{Metric-contour certificate}

\begin{definition}[Finite-cutoff metric-contour certificate]
\label{def:v5v8-MCC}
An MCC is a cofinal family of finite-dimensional metric cycles
\(\Gamma_\Lambda\) with:
\begin{enumerate}
 \item cutoff-uniform coercivity and moment bounds;
 \item compatibility with the gauge-fixed Ward vector fields, including
 every transverse contour term;
 \item an antilinear reflection involution and positive OS forms on the
 physical cylinder algebra;
 \item compatibility of the cycles and normalized functionals under shell
 projection;
 \item equality of the induced shell map with the Ward-complete map used in
 the V5--V7 operator-tower estimates.
\end{enumerate}
\end{definition}

\begin{theorem}[Conditional contour input to the continuum programme]
\label{thm:v5v8-MCC}
An MCC whose shell maps also satisfy the V7 gravity
forest-resummation certificate supplies the metric-integration,
reflection, and local operator-tail inputs to the conditional assembly
theorem.  The remaining hypercharge, infrared matching, all-graph matter
bound, and OS regularity gates remain separate.
\end{theorem}

\begin{proof}
Items 1 and 4 define cutoff-compatible normalized functionals with uniform
integrability; items 2 and 3 give Ward identities and reflection
positivity; item 5 identifies their shell map.  The V7 certificate then
gives the compact local tail, while the inherited V96 theorem treats the
rooted nonlocal fibre.  None of these implications supplies the separately
listed gates.
\end{proof}

No MCC is presently known.  V8 rules out treating the straight imaginary
conformal rotation as such a certificate.  The next upstream question is
whether the conformal direction can be removed or stabilized before
integration by an exact constraint reduction rather than by a global
linear rotation.

V9 will formulate a finite-dimensional constraint-reduction theorem using
a Schur complement and Faddeev--Popov slice.  It will identify when the
negative conformal block is a gauge direction, when it survives as a
physical scalar determinant, and what positivity estimate the reduced
Hessian must satisfy.


\clearpage
\section{V9: gauge-slice reduction and the surviving conformal block}
\label{sec:v5v9}

V8 ruled out the straight imaginary conformal contour as a complete metric
integration certificate.  A different possibility is that the unstable
conformal coordinate is redundant and disappears after an exact
diffeomorphism constraint reduction.  This note tests that possibility at
finite cutoff.  A linear Faddeev--Popov slice removes precisely the gauge
range, while the action Hessian on the slice is unchanged by the
gauge-fixing term.  On a Ricci-flat background, nonconstant conformal
variations are not diffeomorphism directions and their negative quadratic
form survives every transverse slice.

\subsection{Exact linear gauge slicing}

Let \(X\) be a finite-dimensional real inner-product space of cutoff field
variables, let \(\mathfrak g\) be the finite-dimensional infinitesimal
gauge space, and let
\[
 K:\mathfrak g\longrightarrow X
\]
be injective.  A linear gauge condition is a surjection
\(L:X\to\mathfrak g\) such that
\begin{equation}
 M_{\rm FP}:=LK:\mathfrak g\longrightarrow\mathfrak g
 \quad\text{is invertible}.
 \label{eq:v5v9-FP}
\end{equation}
Set
\begin{equation}
 P=I-K(LK)^{-1}L.
 \label{eq:v5v9-slice-projector}
\end{equation}

\begin{lemma}[Slice decomposition]
\label{lem:v5v9-slice}
The operator \(P\) is a projection with
\[
 \operatorname{Ran}P=\ker L,
 \qquad
 \ker P=\operatorname{Ran}K.
\]
Every \(x\in X\) has the unique decomposition
\begin{equation}
 x=Px+K(LK)^{-1}Lx.
 \label{eq:v5v9-decomposition}
\end{equation}
\end{lemma}

\begin{proof}
Direct multiplication gives
\[
 LP=L-LK(LK)^{-1}L=0,
 \qquad
 PK=K-K(LK)^{-1}LK=0.
\]
Equation \eqref{eq:v5v9-slice-projector} gives \(P^2=P\) and
\eqref{eq:v5v9-decomposition}.  If \(x\in\ker L\), then \(Px=x\); if
\(Px=0\), the decomposition puts \(x\) in \(\operatorname{Ran}K\).
Uniqueness follows because
\(\ker L\cap\operatorname{Ran}K=\{0\}\), a consequence of the
invertibility of \(LK\).
\end{proof}

Let \(H:X\to X\) be a symmetric action Hessian at a stationary background.
Infinitesimal gauge invariance gives
\begin{equation}
 HK=0.
 \label{eq:v5v9-HK-zero}
\end{equation}

\begin{theorem}[Gauge-invariant quadratic form descends exactly]
\label{thm:v5v9-quadratic-descent}
Under \eqref{eq:v5v9-HK-zero},
\begin{equation}
 \langle x,Hx\rangle
 =\langle Px,HPx\rangle
 \qquad(x\in X).
 \label{eq:v5v9-quadratic-descent}
\end{equation}
Consequently a vector \(c\) represents a negative physical slice direction
if and only if
\[
 Pc\ne0
 \quad\text{and}\quad
 \langle c,Hc\rangle<0.
\]
\end{theorem}

\begin{proof}
Write \(x=Px+K\xi\) using
Lemma~\ref{lem:v5v9-slice}.  Symmetry and \(HK=0\) make both the cross term
and the gauge--gauge term vanish.  This proves
\eqref{eq:v5v9-quadratic-descent}.  The last statement follows because
\(Pc=0\) exactly for a gauge vector, whereas a nonzero \(Pc\) has the same
quadratic value as \(c\).
\end{proof}

\subsection{What the Faddeev--Popov factor changes}

Choose Lebesgue measures on \(X\), \(\ker L\), and \(\mathfrak g\).
The linear isomorphism
\[
 (s,\xi)\longmapsto s+K\xi
\]
has a constant Jacobian \(J_{K,L}>0\).

\begin{proposition}[Exact finite-dimensional FP cancellation]
\label{prop:v5v9-FP-cancellation}
For every integrable gauge-invariant function \(f\),
\begin{equation}
 \int_X f(x)\,
 \delta(Lx)\,|\det(LK)|\,dx
 =
 J_{K,L}\int_{\ker L}f(s)\,ds.
 \label{eq:v5v9-FP-integral}
\end{equation}
The determinant removes the gauge-orbit Jacobian; it does not alter the
sign of \(H|_{\ker L}\).
\end{proposition}

\begin{proof}
Use \(x=s+K\xi\).  Gauge invariance gives \(f(s+K\xi)=f(s)\), and
\[
 \delta(Lx)=\delta(LK\xi)
 =|\det(LK)|^{-1}\delta(\xi).
\]
The two determinant factors cancel and the \(\xi\) integral evaluates at
zero, proving \eqref{eq:v5v9-FP-integral}.  The remaining action is its
literal restriction to \(\ker L\), so its Hessian is
\(H|_{\ker L}\).
\end{proof}

A quadratic gauge-fixing term produces
\[
 H_\alpha=H+\alpha^{-1}L^*L,
 \qquad\alpha>0.
\]

\begin{corollary}[Gauge fixing cannot repair a slice-negative mode]
\label{cor:v5v9-gauge-fixing}
For \(s\in\ker L\),
\[
 \langle s,H_\alpha s\rangle=\langle s,Hs\rangle.
\]
Hence a negative direction of the reduced Hessian remains negative for
every \(\alpha>0\).
\end{corollary}

\begin{proof}
The gauge-fixing term is \(\alpha^{-1}\|Ls\|^2=0\) on the slice.
\end{proof}

\subsection{Conformal variations are not generally gauge}

Let \((M,\bar g)\) be a compact Riemannian manifold of dimension \(d\ge3\).
The infinitesimal diffeomorphism action is
\[
 K\xi=\mathcal L_\xi\bar g.
\]
A pure conformal variation is
\[
 c_\phi=2\phi\,\bar g.
\]

\begin{lemma}[Gauge conformal modes are conformal Killing modes]
\label{lem:v5v9-conformal-Killing}
The equality
\[
 \mathcal L_\xi\bar g=2\phi\,\bar g
 \label{eq:v5v9-pure-conformal-gauge}
\]
holds if and only if \(\xi\) is a conformal Killing field and
\begin{equation}
 \phi=\frac1d\operatorname{div}_{\bar g}\xi.
 \label{eq:v5v9-phi-div}
\end{equation}
The space of such \(\phi\) is finite dimensional.
\end{lemma}

\begin{proof}
Taking the trace of \eqref{eq:v5v9-pure-conformal-gauge} gives
\(2\operatorname{div}\xi=2d\phi\), which is
\eqref{eq:v5v9-phi-div}.  Subtracting the trace gives
\[
 \nabla_{(\mu}\xi_{\nu)}
 -\frac1d(\operatorname{div}\xi)\bar g_{\mu\nu}=0,
\]
the conformal Killing equation.  The converse follows by restoring the
trace.  On a compact manifold with \(d\ge3\), the conformal Killing
operator is overdetermined elliptic; its kernel is finite dimensional.
The divergence maps that kernel into a finite-dimensional space of
functions.
\end{proof}

\begin{corollary}[Cofinitely many conformal modes survive the quotient]
\label{cor:v5v9-conformal-survives}
Any cutoff family whose scalar sector exhausts an infinite-dimensional
function space contains, at sufficiently high cutoff, conformal variations
not in \(\operatorname{Ran}K\).  Their slice projections under
\eqref{eq:v5v9-slice-projector} are nonzero.
\end{corollary}

\begin{proof}
By Lemma~\ref{lem:v5v9-conformal-Killing}, the gauge-conformal subspace is
fixed and finite dimensional.  An exhausting scalar cutoff eventually
contains vectors outside it.  Lemma~\ref{lem:v5v9-slice} says precisely that
their projections are nonzero.
\end{proof}

\subsection{Einstein Hessian on a Ricci-flat background}

Take the Euclidean Einstein action and sign convention of V8.  Let
\(\bar g\) be Ricci flat.  For
\(g_\epsilon=e^{2\epsilon\phi}\bar g\), equation
\eqref{eq:v5v8-conformal-action} gives
\begin{equation}
 S_E(g_\epsilon)
 =
 -6\kappa\epsilon^2
 \int_M|\bar\nabla\phi|^2\,dV_{\bar g}
 +O(\epsilon^3).
 \label{eq:v5v9-EH-conformal-Hessian}
\end{equation}

\begin{theorem}[The conformal instability survives FP reduction]
\label{thm:v5v9-surviving-instability}
Let \(L\) be any admissible linear gauge condition satisfying
\eqref{eq:v5v9-FP}.  If \(\phi\) is nonconstant and
\(c_\phi\notin\operatorname{Ran}K\), then
\[
 Pc_\phi\ne0
\]
and the Einstein Hessian is strictly negative on \(Pc_\phi\).
Consequently the real gauge-slice Gaussian is not positive definite.
Along an exhausting cutoff sequence, the number of such negative scalar
directions is unbounded.
\end{theorem}

\begin{proof}
The background is stationary, so diffeomorphism invariance gives
\(HK=0\).  The quadratic coefficient in
\eqref{eq:v5v9-EH-conformal-Hessian} is strictly negative for nonconstant
\(\phi\).  Theorem~\ref{thm:v5v9-quadratic-descent} gives the same value on
\(Pc_\phi\), which is nonzero because \(c_\phi\) is not gauge.
Corollary~\ref{cor:v5v9-conformal-survives} supplies an increasing number
of linearly independent nongauge scalar modes as the cutoff grows; their
Dirichlet quadratic form is negative definite modulo constants.
\end{proof}

\begin{firewall}[Ghost determinants do not supply bosonic coercivity]
\label{fw:v5v9-ghost-coercivity}
The Faddeev--Popov determinant is essential for the gauge quotient, but
Proposition~\ref{prop:v5v9-FP-cancellation} shows that it does not change
the reduced bosonic Hessian.  A cancellation in a formal determinant ratio
is not a positive lower bound for the real metric action.
\end{firewall}

\subsection{Schur complements and auxiliary fields}

Suppose a reduced quadratic action on \(Y\oplus Z\) has block matrix
\[
 \mathbb H=
 \begin{pmatrix}
  A&B\\ B^*&D
 \end{pmatrix},
 \qquad
 A\ge aI>0.
\]

\begin{proposition}[Exact Schur reduction]
\label{prop:v5v9-Schur}
For fixed \(z\in Z\), the minimum over \(y\in Y\) is attained at
\(y=-A^{-1}Bz\), and
\begin{equation}
 \inf_y
 \left\langle
 \binom y z,\mathbb H\binom y z
 \right\rangle
 =
 \langle z,S_Zz\rangle,
 \qquad
 S_Z=D-B^*A^{-1}B.
 \label{eq:v5v9-Schur}
\end{equation}
In particular,
\begin{equation}
 S_Z\le D
 \label{eq:v5v9-Schur-order}
\end{equation}
as quadratic forms.  Mixing with a positive block cannot stabilize an
already negative \(D\).
\end{proposition}

\begin{proof}
Complete the square:
\[
 \left\langle
 \binom y z,\mathbb H\binom y z
 \right\rangle
 =
 \langle y+A^{-1}Bz,\,
 A(y+A^{-1}Bz)\rangle
 +\langle z,(D-B^*A^{-1}B)z\rangle.
\]
The first term has the stated unique minimum.  Since
\(B^*A^{-1}B\ge0\), equation \eqref{eq:v5v9-Schur-order} follows.
\end{proof}

\begin{corollary}[What a stabilizing completion must change]
\label{cor:v5v9-stabilize}
A real positive metric Gaussian can be obtained only if the exact reduced
scalar block itself acquires a positive lower bound after all constraints,
or if the scalar is removed by an additional proved constraint.  Coupling
it to fields with positive quadratic blocks and then integrating those
fields out cannot create that bound.
\end{corollary}

\begin{proof}
Apply Proposition~\ref{prop:v5v9-Schur}.  If the scalar is not removed,
positivity of the full reduced quadratic form requires positivity of every
Schur complement, including \(S_Z\).
\end{proof}

Higher-curvature terms can change \(D\), while lapse or connection
variables can impose additional constraints.  Their effect must be
computed before making a positivity claim; neither is supplied by the
Einstein--Hilbert FP determinant alone.

\subsection{A reduced-Hessian certificate}

Let \(\mathscr S_\Lambda=\ker L_\Lambda\) be the finite-cutoff slice after
all proved algebraic constraints have been solved.  Let
\(H_\Lambda^{\rm red}\) be the exact Hessian on the remaining bosonic
coordinates.

\begin{definition}[Reduced metric Hessian certificate]
\label{def:v5v9-RMHC}
An RMHC consists of:
\begin{enumerate}
 \item uniform invertibility of \(L_\Lambda K_\Lambda\) on the chosen
 nonzero-mode gauge space;
 \item exact identification of every residual scalar coordinate;
 \item a cutoff-independent \(m>0\) such that
 \begin{equation}
  H_\Lambda^{\rm red}\ge mI
  \label{eq:v5v9-uniform-positive-Hessian}
 \end{equation}
 on the physical real slice, or a specified convergent complex contour with
 its own reflection-positive observable algebra;
 \item determinant and Jacobian bounds compatible with shell projection;
 \item nonlinear coercivity and moment estimates extending the Hessian
 bound to the full action.
\end{enumerate}
\end{definition}

\begin{theorem}[RMHC input to the metric-contour certificate]
\label{thm:v5v9-RMHC}
An RMHC with real-slice positivity and nonlinear coercivity supplies the
finite-cutoff integrability and gauge-reduction portions of the V8 MCC.
It does not by itself prove reflection positivity or cofinal shell
compatibility.
\end{theorem}

\begin{proof}
Items 1 and 2 give the exact FP slice by
Lemma~\ref{lem:v5v9-slice} and
Proposition~\ref{prop:v5v9-FP-cancellation}.  Items 3 and 5 give the
coercive real action required by
Theorem~\ref{thm:v5v8-coercive-contour}; item 4 controls the associated
finite-cutoff density.  Reflection positivity and projective compatibility
are additional statements in the MCC definition and do not follow from
ordinary positive definiteness.
\end{proof}

\begin{corollary}[Einstein--Hilbert alone fails the RMHC]
\label{cor:v5v9-EH-fails-RMHC}
On an exhausting cutoff sequence about a compact Ricci-flat background,
the real Euclidean Einstein--Hilbert action fails
\eqref{eq:v5v9-uniform-positive-Hessian} after every admissible
diffeomorphism FP reduction.
\end{corollary}

\begin{proof}
Theorem~\ref{thm:v5v9-surviving-instability} produces nonzero slice vectors
with negative Hessian values at arbitrarily high cutoffs.
\end{proof}

V9 closes one possible escape from the conformal bottleneck: ordinary
diffeomorphism gauge reduction does not remove the nonconstant conformal
sector.  A successful route must add a genuine scalar constraint, modify
the gravitational action so that the reduced scalar block is coercive, or
supply a complex cycle together with a new physical reality and OS
certificate.

V10 will begin with the compulsory companion audit.  It will then test the
action-modification route abstractly: a higher-derivative scalar symbol
can stabilize ultraviolet modes only if all intermediate frequencies and
the associated extra poles satisfy a uniform Schur and reflection test.


\clearpage
\section{V10: companion audit and the higher-derivative reflection test}
\label{sec:v5v10}

This is the second compulsory companion checkpoint of the fifth series.
After recording it, the note tests the action-modification route left open
by V9.  A positive higher-derivative principal symbol can control large
momenta, but the complete scalar symbol must stay positive at every
intermediate momentum and its inverse must pass reflection positivity.
For a raw scalar with a polynomial fourth-order inverse propagator, these
two requirements do not coincide.

\subsection{Compulsory companion audit}

At the time of this audit, the newest live companion files were
\begin{align*}
 &\texttt{Renewal\_Yang\_Mills\_Mass\_Gap\_Research\_Notes\_V28.tex},\\
 &\hspace{1em}\text{SHA--256 }
 \texttt{63DBFA9673192971071F7DD5FFC653BE77F85C69968ACACA96B956280B602AE9},
 \\
 &\texttt{Renewal\_NS\_Stability\_Research\_Notes\_V31.tex},\\
 &\hspace{1em}\text{SHA--256 }
 \texttt{F1121B62238AD5491BB7CF03635EA9934942EDF573ED8FAAA9EE79E1D38A6696}.
\end{align*}

Yang--Mills V28 absorbs arbitrary dualizers inside each fixed local-Casimir
block and sums all heat-weighted output labels.  Its remaining obstruction
is off-diagonal mixing between distinct Casimir tuples: ordinary block
norms do not control an unheated label, and the prescribed next test is a
two-sided heat-weighted Schur bound.  Navier--Stokes V31 is unchanged from
the V5 audit; its exact first-moment calculation continues to show that a
positive marked estimate does not imply the corresponding unmarked
critical estimate.

\begin{theorem}[V10 audit transfer]
\label{thm:v5v10-audit}
For a reduced metric Hessian decomposed into spectral or constraint blocks,
the companion results authorize two rules:
\begin{enumerate}
 \item diagonal block positivity must be supplemented by a two-sided
 weighted bound on every off-diagonal block;
 \item Euclidean coercivity of the resulting operator remains distinct
 from reflection positivity of its inverse covariance.
\end{enumerate}
\end{theorem}

\begin{proof}
The first rule is the abstract Schur mechanism isolated by Yang--Mills V28.
The second is the metric analogue of the NS marked/unmarked distinction:
the positive quadratic action is one weighted statement, while OS
positivity is a positivity statement after inversion and time reflection.
Neither operation may be deleted from the certificate.
\end{proof}

\subsection{A two-sided block coercivity test}

Let
\[
 \mathscr H=\bigoplus_{\alpha\in A}\mathscr H_\alpha,
 \qquad
 D=\bigoplus_{\alpha\in A}D_\alpha,
 \qquad
 D_\alpha\ge m_\alpha I>0.
\]
Let \(E=E^*\) contain all off-diagonal mixing.

\begin{theorem}[Weighted Schur coercivity]
\label{thm:v5v10-weighted-Schur}
Assume the quadratic-form operator
\begin{equation}
 W=D^{-1/2}ED^{-1/2}
 \label{eq:v5v10-weighted-mixing}
\end{equation}
extends boundedly and
\begin{equation}
 \|W\|\le q<1.
 \label{eq:v5v10-Schur-small}
\end{equation}
Then \(H=D+E\) is positive and
\begin{equation}
 (1-q)D\le H\le(1+q)D
 \label{eq:v5v10-form-comparison}
\end{equation}
as quadratic forms.  In particular,
\[
 H\ge(1-q)\left(\inf_\alpha m_\alpha\right)I
\]
when the displayed infimum is positive.
\end{theorem}

\begin{proof}
For \(x\) in the form domain of \(D^{1/2}\), put \(y=D^{1/2}x\).  Then
\[
 \langle x,Hx\rangle
 =\langle y,(I+W)y\rangle.
\]
Self-adjointness and \(\|W\|\le q\) give
\[
 (1-q)\|y\|^2
 \le\langle y,(I+W)y\rangle
 \le(1+q)\|y\|^2,
\]
which is exactly \eqref{eq:v5v10-form-comparison}.
\end{proof}

\begin{firewall}[One-sided block bounds are insufficient]
\label{fw:v5v10-one-sided}
Bounds on each \(E_{\alpha\beta}\) without the two weights in
\eqref{eq:v5v10-weighted-mixing} need not control the sum over
\(\alpha,\beta\).  The weights encode the cost of entering and leaving the
spectral blocks; deleting either one can restore an unbounded label
multiplicity.
\end{firewall}

Theorem~\ref{thm:v5v10-weighted-Schur} is a coercivity theorem.  It says
nothing about the sign of the reflected covariance kernel
\(RH^{-1}\).

\subsection{The scalar polynomial symbol}

Consider a translation-invariant real scalar at finite cutoff with
quadratic symbol
\begin{equation}
 P(s)=m^2-a s+b s^2,
 \qquad
 s=p_0^2+|\mathbf p|^2,
 \qquad
 a,b>0.
 \label{eq:v5v10-fourth-symbol}
\end{equation}
The term \(-as\) models the Euclidean conformal Einstein Hessian, while
\(bs^2\) is a stabilizing fourth-derivative contribution.

\begin{proposition}[Ultraviolet stabilization leaves an intermediate test]
\label{prop:v5v10-intermediate}
If \(m=0\), then
\[
 P(s)<0
 \qquad(0<s<a/b).
 \label{eq:v5v10-negative-window}
\]
If \(m>0\), then \(P(s)>0\) for every \(s\ge0\) if and only if
\begin{equation}
 m^2>\frac{a^2}{4b}.
 \label{eq:v5v10-polynomial-positive}
\end{equation}
Equality gives a zero mode.
\end{proposition}

\begin{proof}
For \(m=0\),
\[
 P(s)=s(bs-a),
\]
which proves \eqref{eq:v5v10-negative-window}.  For \(m>0\), the convex
quadratic has its minimum over \(s\ge0\) at \(s=a/(2b)\), with value
\(m^2-a^2/(4b)\).  The stated alternatives follow.
\end{proof}

Thus a positive principal coefficient \(b\) proves only high-frequency
control.  A uniform RMHC must also control the finite momentum window.

\subsection{Reflection positivity and Stieltjes covariance}

For a Euclidean-invariant scalar two-point function depending on
\(s=p^2\), a necessary consequence of OS positivity and the spectral
condition is a positive Kallen--Lehmann representation
\begin{equation}
 C(s)=\int_{[0,\infty)}
 \frac{d\rho(\mu^2)}{s+\mu^2},
 \qquad
 \rho\ge0.
 \label{eq:v5v10-Stieltjes}
\end{equation}
We call such a function Stieltjes.  The following finite-pole test can also
be proved directly from the linear OS form.

\begin{lemma}[Sign of a finite spectral residue]
\label{lem:v5v10-residue-sign}
Suppose
\[
 C(s)=\sum_{j=1}^J\frac{r_j}{s+\mu_j^2},
 \qquad
 0<\mu_1<\cdots<\mu_J.
 \label{eq:v5v10-partial-fractions}
\]
If the scalar Gaussian is reflection positive, then every \(r_j\ge0\).
\end{lemma}

\begin{proof}
At fixed spatial momentum, inverse Fourier transformation in \(p_0\)
gives the reflected-time kernel
\[
 K(t,u)
 =\sum_{j=1}^J
 \frac{r_j}{2\omega_j}e^{-\omega_j(t+u)},
 \qquad
 \omega_j=(|\mathbf p|^2+\mu_j^2)^{1/2}.
\]
For a positive-time test function \(f\), its OS quadratic form is
\[
 \sum_{j=1}^J\frac{r_j}{2\omega_j}
 \left|\int_0^\infty f(t)e^{-\omega_jt}\,dt\right|^2.
\]
The exponentials with distinct \(\omega_j\) are linearly independent.
Using a finite linear combination of narrow test functions, one can make
the displayed Laplace values vanish at all indices except any prescribed
one.  Positivity of the remaining term forces the corresponding
\(r_j\ge0\).
\end{proof}

\begin{theorem}[Fourth-order raw scalar pole obstruction]
\label{thm:v5v10-fourth-order-no-OS}
Let \(P\) be a real quadratic polynomial with positive leading coefficient,
and suppose \(P(s)>0\) for all \(s\ge0\).  Then the nonconstant covariance
\begin{equation}
 C(s)=\frac1{P(s)}
 \label{eq:v5v10-inverse-polynomial}
\end{equation}
is not Stieltjes and cannot be the reflection-positive covariance of the
raw scalar field.
\end{theorem}

\begin{proof}
There are three possibilities for the two roots of \(P\).
If a root lies on the positive real axis, positivity on \(s\ge0\) fails,
so this case is excluded.  If the roots are nonreal, \(C\) has poles away
from the negative real axis, whereas a Stieltjes function is holomorphic
on \(\mathbb C\setminus(-\infty,0]\).  If the roots are two distinct
negative numbers, write
\[
 P(s)=b(s+\mu_1^2)(s+\mu_2^2),
 \qquad
 0<\mu_1<\mu_2.
\]
Then
\[
 \frac1{P(s)}
 =
 \frac1{b(\mu_2^2-\mu_1^2)}
 \left(\frac1{s+\mu_1^2}-\frac1{s+\mu_2^2}\right),
\]
so one residue is negative, contradicting
Lemma~\ref{lem:v5v10-residue-sign}.  A repeated negative root gives a
double pole.  A positive Stieltjes measure can produce a simple pole from
an atom but not a double pole, again a contradiction.  These cases exhaust
a real quadratic polynomial.
\end{proof}

\begin{corollary}[Coercive fourth-order conformal symbol still fails the raw
OS test]
\label{cor:v5v10-coercive-not-OS}
When \eqref{eq:v5v10-polynomial-positive} holds, the scalar action with
symbol \eqref{eq:v5v10-fourth-symbol} is positive definite, but its raw
covariance fails reflection positivity by
Theorem~\ref{thm:v5v10-fourth-order-no-OS}.
\end{corollary}

\begin{proof}
Proposition~\ref{prop:v5v10-intermediate} gives positivity and
Theorem~\ref{thm:v5v10-fourth-order-no-OS} gives the independent pole
obstruction.
\end{proof}

There is also a degree-independent ultraviolet obstruction.

\begin{proposition}[A positive spectral covariance cannot decay faster than
\(s^{-1}\)]
\label{prop:v5v10-spectral-tail}
Let \(C\) have the nonzero positive representation
\eqref{eq:v5v10-Stieltjes}.  Then
\begin{equation}
 \liminf_{s\to\infty}sC(s)>0,
 \label{eq:v5v10-no-fast-decay}
\end{equation}
where the limit may be infinite.  In particular,
\(C(s)=O(s^{-1-\epsilon})\) for any \(\epsilon>0\) is impossible.
\end{proposition}

\begin{proof}
For every \(\mu^2\),
\[
 \frac{s}{s+\mu^2}\uparrow1
 \quad\text{as }s\to\infty.
\]
Monotone convergence gives
\[
 sC(s)\longrightarrow\rho([0,\infty]).
\]
Because the measure is nonzero, this total mass is strictly positive,
possibly infinite.
\end{proof}

\begin{firewall}[Observable numerators and constraint quotients]
\label{fw:v5v10-observable-numerator}
The pole obstruction concerns the covariance \(1/P\) of the raw residual
scalar.  A physical observable may have a numerator that cancels an
unphysical pole, or an exact constraint may remove the scalar altogether.
Such a cancellation must be proved on the physical OS algebra; it cannot
be inferred from the positive principal symbol.
\end{firewall}

\subsection{A reflection-compatible stabilization certificate}

\begin{definition}[Scalar stabilization and reflection certificate]
\label{def:v5v10-SSRC}
An SSRC for a cofinal cutoff family consists of:
\begin{enumerate}
 \item an exact reduced scalar block after the V9 gauge and constraint
 quotient;
 \item positive diagonal spectral blocks and the weighted off-diagonal
 estimate \eqref{eq:v5v10-Schur-small};
 \item a cutoff-uniform lower bound for the complete symbol at all
 frequencies;
 \item a positive spectral representation for every physical scalar
 covariance, or a direct OS proof on the physical cylinder algebra;
 \item nonlinear coercivity and shell-compatible determinant estimates.
\end{enumerate}
\end{definition}

\begin{theorem}[Conditional action-modification route]
\label{thm:v5v10-conditional-action}
An SSRC supplies the scalar part of the V9 RMHC and the linear
reflection-positivity part of the V8 MCC.  A finite local
fourth-derivative modification with an uncancelled raw scalar covariance
cannot supply item 4.
\end{theorem}

\begin{proof}
Items 2 and 3 give reduced Hessian positivity by
Theorem~\ref{thm:v5v10-weighted-Schur}; item 5 gives the nonlinear RMHC
bounds.  Item 4 is the scalar OS condition.  Its failure for an
uncancelled fourth-order raw covariance is
Theorem~\ref{thm:v5v10-fourth-order-no-OS}.
\end{proof}

V10 therefore narrows the action-modification branch.  Adding a local
higher-derivative term can stabilize the Euclidean conformal Hessian only
after a full intermediate-frequency and mixing test, and an uncancelled
higher-derivative scalar then fails the positive spectral test.  A viable
construction needs an exact constraint removal, an observable-level pole
cancellation, or a nonlocal covariance with positive spectral measure and
separately controlled locality.

V11 will investigate the nonlocal spectral option without assuming it is
acceptable.  It will characterize completely monotone heat-kernel
mixtures, prove their reflection-positive Gaussian representation, and
measure the quasilocality cost relevant to the V6 curvature-word and V96
rooted-kernel spaces.


\clearpage
\section{V11: positive spectral mixtures and their quasilocality cost}
\label{sec:v5v11}

V10 showed that an uncancelled polynomial fourth-order scalar covariance
cannot combine ultraviolet coercivity with reflection positivity.  This
note analyzes the positive spectral alternative.  Stieltjes covariances
are exactly Laplace transforms whose heat weight is completely monotone;
they are reflection positive because they are positive mixtures of massive
free covariances.  Their heat-shell kernels satisfy the weighted smoothing
bounds used by V96--V97 under explicit spectral moment assumptions.
However, a nonzero positive spectral covariance cannot improve the
ultraviolet power beyond \(s^{-1}\), and a mass gap is necessary for
uniform exponential localization.

\subsection{Stieltjes and heat-mixture representations}

Let \(\rho\) be a positive Borel measure on \([0,\infty)\) satisfying
\begin{equation}
 \int_{[0,\infty)}\frac{d\rho(\mu^2)}{1+\mu^2}<\infty.
 \label{eq:v5v11-spectral-integrability}
\end{equation}
Define
\begin{align}
 C(s)&=\int_{[0,\infty)}
       \frac{d\rho(\mu^2)}{s+\mu^2},
 \qquad s>0,
 \label{eq:v5v11-Stieltjes}\\
 w(t)&=\int_{[0,\infty)}
       e^{-t\mu^2}\,d\rho(\mu^2),
 \qquad t>0.
 \label{eq:v5v11-heat-weight}
\end{align}

\begin{theorem}[Positive spectral heat representation]
\label{thm:v5v11-heat-representation}
Under \eqref{eq:v5v11-spectral-integrability},
\begin{equation}
 C(s)=\int_0^\infty e^{-ts}w(t)\,dt.
 \label{eq:v5v11-Laplace-C}
\end{equation}
Moreover, \(w\) is completely monotone:
\begin{equation}
 (-1)^kw^{(k)}(t)
 =
 \int(\mu^2)^ke^{-t\mu^2}\,d\rho(\mu^2)
 \ge0
 \qquad(k\ge0).
 \label{eq:v5v11-complete-monotone}
\end{equation}
Conversely, if a locally finite nonnegative function \(w\) is completely
monotone and the right side of \eqref{eq:v5v11-Laplace-C} is finite, then
Bernstein's theorem gives a positive measure \(\rho\) for which
\eqref{eq:v5v11-heat-weight} and \eqref{eq:v5v11-Stieltjes} hold.
\end{theorem}

\begin{proof}
The identity
\[
 \frac1{s+\mu^2}=\int_0^\infty e^{-t(s+\mu^2)}\,dt
\]
and Tonelli's theorem give \eqref{eq:v5v11-Laplace-C}.  Differentiation
under the nonnegative integral gives
\eqref{eq:v5v11-complete-monotone}; finiteness at positive \(t\) follows
from the exponential damping and
\eqref{eq:v5v11-spectral-integrability}.  Conversely, Bernstein's theorem
represents every completely monotone \(w\) as the Laplace transform of a
positive measure.  A second use of Tonelli then recovers
\eqref{eq:v5v11-Stieltjes}.
\end{proof}

\begin{corollary}[Derivative signs of the covariance]
\label{cor:v5v11-C-monotone}
For every \(k\ge0\),
\begin{equation}
 (-1)^kC^{(k)}(s)
 =k!\int\frac{d\rho(\mu^2)}{(s+\mu^2)^{k+1}}
 \ge0.
 \label{eq:v5v11-C-signs}
\end{equation}
\end{corollary}

\begin{proof}
Differentiate \eqref{eq:v5v11-Stieltjes} under the integral; the
integrability assumption supplies domination on every compact subset of
\(s>0\).
\end{proof}

Complete monotonicity of \(C\) alone is weaker than being Stieltjes.  The
additional statement is that its Laplace weight \(w\) is itself completely
monotone.

\subsection{Reflection positivity of the Gaussian mixture}

Let \(C_\mu=(-\Delta+\mu^2)^{-1}\) be the Euclidean free covariance on
\(\mathbb R^d\), and define the mixed covariance
\begin{equation}
 C_\rho=\int C_\mu\,d\rho(\mu^2)
 \label{eq:v5v11-mixed-covariance}
\end{equation}
on a test-function class where the integral is finite.

\begin{theorem}[OS positivity of a positive spectral mixture]
\label{thm:v5v11-OS-mixture}
The centered Gaussian functional with covariance \(C_\rho\) is reflection
positive on scalar cylinder observables supported at positive Euclidean
time.
\end{theorem}

\begin{proof}
It is enough to prove positivity of the reflected covariance form; Gaussian
reflection positivity then follows by the standard exponential-vector
density argument.  Fourier transform only in the spatial variables.  For
a positive-time test function \(f\),
\begin{align*}
 \langle\theta f,C_\mu f\rangle
 &=
 \int_{\mathbb R^{d-1}}
 \frac{d\mathbf p}{2\omega_\mu(\mathbf p)}
 \left|
  \int_0^\infty
  e^{-\omega_\mu(\mathbf p)t}
  \widehat f(t,\mathbf p)\,dt
 \right|^2,\\
 \omega_\mu(\mathbf p)
 &=\sqrt{|\mathbf p|^2+\mu^2}.
\end{align*}
This is nonnegative for every \(\mu\).  Integrating the nonnegative form
against \(d\rho(\mu^2)\) proves positivity for \(C_\rho\).  For Gaussian
exponentials, the OS Gram matrix is the entrywise exponential of this
positive covariance Gram matrix and is positive; finite linear
combinations are therefore positive, and density gives the cylinder
algebra.
\end{proof}

\begin{firewall}[Shell covariances need not be separately OS positive]
\label{fw:v5v11-shell-OS}
A proper-time band is a difference of two low-pass covariances.  Positivity
of its Fourier multiplier does not by itself imply reflection positivity
of that band.  Theorem~\ref{thm:v5v11-OS-mixture} applies to the complete
spectral covariance.  A Wilsonian decomposition must retain the V99
shell-sewing or an equivalent proof that the assembled functional is OS
positive.
\end{firewall}

\subsection{Weighted heat-shell estimates}

Let
\[
 h_t(x)=(4\pi t)^{-d/2}e^{-|x|^2/(4t)}
\]
be the Euclidean heat kernel.  For a multi-index \(\gamma\), set
\[
 \|f\|_{L^1_\eta}
 =\int_{\mathbb R^d}e^{\eta|x|}|f(x)|\,dx.
\]

\begin{lemma}[Weighted heat derivative]
\label{lem:v5v11-heat-derivative}
For every \(\gamma\) there are constants \(A_\gamma,c_\gamma>0\) such that
\begin{equation}
 \|\partial^\gamma h_t\|_{L^1_\eta}
 \le
 A_\gamma t^{-|\gamma|/2}e^{c_\gamma\eta^2t}
 \qquad(t>0,\ \eta\ge0).
 \label{eq:v5v11-heat-derivative}
\end{equation}
\end{lemma}

\begin{proof}
After the change of variables \(x=\sqrt t\,u\), every derivative has the
form
\[
 \partial^\gamma h_t(\sqrt t\,u)
 =
 t^{-(d+|\gamma|)/2}P_\gamma(u)e^{-|u|^2/4}
\]
for a polynomial \(P_\gamma\).  The inequality
\[
 \eta\sqrt t\,|u|
 \le\frac{|u|^2}{8}+2\eta^2t
\]
absorbs the exponential weight into half of the Gaussian.  The remaining
polynomial-Gaussian integral is finite and gives
\eqref{eq:v5v11-heat-derivative}.
\end{proof}

For \(0<a<b<\infty\), define the spectral proper-time shell
\begin{equation}
 C_{[a,b]}(x)
 =
 \int_{[0,\infty)}d\rho(\mu^2)
 \int_a^b e^{-t\mu^2}h_t(x)\,dt.
 \label{eq:v5v11-proper-shell}
\end{equation}

\begin{theorem}[Spectral shell smoothing bound]
\label{thm:v5v11-shell-bound}
For every integer \(M\ge0\),
\begin{equation}
 \sum_{|\gamma|\le M}
 \|\partial^\gamma C_{[a,b]}\|_{L^1_\eta}
 \le
 \sum_{|\gamma|\le M}A_\gamma
 \int_a^b
 t^{-|\gamma|/2}e^{c_\gamma\eta^2t}w(t)\,dt.
 \label{eq:v5v11-shell-bound}
\end{equation}
In particular, the right side is finite whenever \(w\) is finite on
\([a,b]\).
\end{theorem}

\begin{proof}
Differentiate \eqref{eq:v5v11-proper-shell}, use Tonelli and the triangle
inequality, then apply
Lemma~\ref{lem:v5v11-heat-derivative}.  The spectral integral is exactly
\(w(t)\).
\end{proof}

\begin{corollary}[Input to the inherited graphwise smoother]
\label{cor:v5v11-V97-input}
If adjacent proper-time shells have a fixed ratio and the rescaled
right side of \eqref{eq:v5v11-shell-bound} is uniform through derivative
order \(M\), then the spectral shell covariance satisfies the heat-kernel
hypothesis of the inherited V97 graphwise theorem through order \(M\).
\end{corollary}

\begin{proof}
The V97 hypothesis is precisely a rescaled weighted \(L^1\) bound for the
shell kernel and its derivatives.  Equation
\eqref{eq:v5v11-shell-bound} supplies it under the stated uniformity.
\end{proof}

This result addresses only the nonlocal rooted kernel.  It does not supply
the all-grade local curvature coefficient bound required in V6.

\subsection{Mass gap and spatial localization}

\begin{proposition}[Mass-gap localization]
\label{prop:v5v11-mass-localization}
Assume
\begin{equation}
 \operatorname{supp}\rho\subset[m_0^2,\infty)
 \qquad(m_0>0).
 \label{eq:v5v11-spectral-gap}
\end{equation}
Then \(w(t)\le e^{-m_0^2t}\rho([0,\infty))\) when the total mass is finite.
More generally, every finite spectral moment bound inserted in
\eqref{eq:v5v11-shell-bound} gains the factor \(e^{-m_0^2t}\).
The complete covariance is exponentially localized at every rate
\(\eta<m_0\), away from its coincident-point ultraviolet singularity.
\end{proposition}

\begin{proof}
The heat-weight estimate follows term by term from
\(e^{-t\mu^2}\le e^{-tm_0^2}\).  For the complete resolvent, use the
Schwinger representation
\[
 C_\rho(x)=\int_0^\infty w(t)h_t(x)\,dt.
\]
The exponent
\[
 -m_0^2t-\frac{|x|^2}{4t}
\]
is at most \(-m_0|x|\) by the arithmetic-geometric mean inequality.
Polynomial factors from derivatives can be absorbed at any strictly
smaller exponential rate.  The short-time endpoint contains the standard
coincident singularity and is excluded from the off-diagonal statement.
\end{proof}

\begin{proposition}[No gap, no uniform exponential rate]
\label{prop:v5v11-no-gap}
If every interval \([0,\varepsilon]\) has positive \(\rho\)-measure, then
no bound
\[
 |C_\rho(x)|\le A e^{-\eta|x|}
\]
with fixed \(\eta>0\) can hold at all sufficiently large \(|x|\).
\end{proposition}

\begin{proof}
Choose \(0<m<\eta\) such that
\(\rho([0,m^2])>0\).  The positive mixture dominates its part with masses
at most \(m\).  Each such massive Green function has large-distance
exponential rate no faster than \(e^{-m|x|}\), up to powers of
\(|x|\).  A positive amount of this contribution cannot be bounded by a
fixed multiple of the faster \(e^{-\eta|x|}\) for all large \(|x|\).
\end{proof}

\subsection{The unavoidable ultraviolet power}

\begin{theorem}[Positive spectral mixtures do not give higher-derivative
decay]
\label{thm:v5v11-no-UV-improvement}
If \(\rho\ne0\), then
\begin{equation}
 \lim_{s\to\infty}sC(s)
 =\rho([0,\infty])\in(0,\infty].
 \label{eq:v5v11-sC-limit}
\end{equation}
Consequently a reflection-positive spectral mixture cannot satisfy
\(C(s)=O(s^{-1-\epsilon})\) for any \(\epsilon>0\).
\end{theorem}

\begin{proof}
This is monotone convergence applied to
\[
 sC(s)=\int\frac{s}{s+\mu^2}\,d\rho(\mu^2),
\]
whose integrand increases pointwise to one.
\end{proof}

\begin{corollary}[The spectral option does not regularize by a fourth-order
propagator]
\label{cor:v5v11-not-fourth-order}
A positive spectral scalar can be quasilocal and OS positive, but it cannot
obtain ultraviolet suppression from a covariance behaving as \(s^{-2}\).
Any improvement in graph power counting must instead come from
cancellations, constraints, composite-observable numerators, or
interactions, while the physical two-point spectral tail remains
compatible with \eqref{eq:v5v11-sC-limit}.
\end{corollary}

\begin{proof}
Apply Theorem~\ref{thm:v5v11-no-UV-improvement}.  The listed mechanisms do
not assert a faster positive raw covariance and are therefore the remaining
logical possibilities.
\end{proof}

\subsection{Locality and matching costs}

Except for a single spectral atom, the inverse
\[
 P(s)=C(s)^{-1}
\]
is generally not a finite polynomial.  It defines a nonlocal
pseudodifferential quadratic action.  Reflection positivity has therefore
been gained by enlarging the spectral content, with a locality cost that
must be measured rather than hidden.

\begin{definition}[Positive spectral metric-scalar certificate]
\label{def:v5v11-PSMC}
A PSMC consists of:
\begin{enumerate}
 \item a positive spectral measure satisfying
 \eqref{eq:v5v11-spectral-integrability};
 \item the shell-uniform derivative bounds
 \eqref{eq:v5v11-shell-bound};
 \item either a spectral gap or a separately proved weighted localization
 estimate;
 \item a proof that the scalar is the correct physical observable after the
 V9 constraint quotient;
 \item a quasilocal expansion of \(C^{-1}\) compatible with the V6
 essential curvature-word space and the exact Ward identities;
 \item infrared matching and cofinal OS shell sewing.
\end{enumerate}
\end{definition}

\begin{theorem}[Conditional spectral branch]
\label{thm:v5v11-conditional-spectral}
A PSMC supplies a reflection-positive Gaussian scalar and the spectral
shell input to the inherited graphwise nonlocal smoother.  It does not
supply higher-derivative ultraviolet power improvement, local
operator-tower compactness, or the nonlinear metric measure.
\end{theorem}

\begin{proof}
Theorem~\ref{thm:v5v11-OS-mixture} gives Gaussian OS positivity, and
Theorem~\ref{thm:v5v11-shell-bound} with
Corollary~\ref{cor:v5v11-V97-input} gives shell smoothing.  The negative
claims are the boundaries proved in
Theorem~\ref{thm:v5v11-no-UV-improvement}, the separation in V6, and the
absence of nonlinear control from a covariance statement.
\end{proof}

No PSMC is presently proved for the gravitational conformal sector.  In
particular, its positive spectral covariance cannot match the negative raw
Einstein conformal residue.  It could be relevant only after a physical
constraint/observable map changes which scalar, if any, survives.

V12 will therefore go upstream to the physical quotient.  It will
formulate a finite-dimensional BRST/Hodge reduction theorem for a
Gaussian gauge complex, prove positivity descent to cohomology under an
explicit contracting homotopy, and test whether the conformal direction
can belong to a gauge quartet.  This separates a true physical negative
mode from an indefinite auxiliary representative.


\clearpage
\section{V12: finite-cutoff BRST--Hodge descent and the conformal quartet test}
\label{sec:v5v12}

V9 proved that a nonconstant conformal variation is not generally an
infinitesimal diffeomorphism.  V10--V11 then showed that treating it as an
uncancelled physical scalar creates a coercivity or reflection-positivity
obstruction.  A third possibility remains: the conformal coordinate may be
removed by a constraint cohomology even though it is not gauge exact.  This
note gives the finite-dimensional Hodge criterion for that statement and
proves precisely when an indefinite auxiliary Gaussian disappears from
physical expectations.

\subsection{A finite-dimensional gauge--constraint complex}

Let
\begin{equation}
 C^{-1}\xrightarrow{\,Q_{-1}\,}
 C^0\xrightarrow{\,Q_0\,}C^1,
 \qquad
 Q_0Q_{-1}=0,
 \label{eq:v5v12-complex}
\end{equation}
be a finite-dimensional real Hilbert complex.  Here \(C^0\) contains the
cutoff bosonic fluctuations, \(\operatorname{Ran}Q_{-1}\) contains
infinitesimal gauge directions, and \(Q_0x=0\) is the linearized constraint.
The physical degree-zero cohomology is
\[
 H^0(Q)=\ker Q_0/\operatorname{Ran}Q_{-1}.
\]
Define the degree-zero Hodge Laplacian
\begin{equation}
 \Delta_0=Q_{-1}Q_{-1}^*+Q_0^*Q_0
 \label{eq:v5v12-Hodge-Laplacian}
\end{equation}
and harmonic space
\begin{equation}
 \mathcal H^0=\ker Q_0\cap\ker Q_{-1}^*.
 \label{eq:v5v12-harmonic}
\end{equation}

\begin{theorem}[Finite-dimensional Hodge decomposition]
\label{thm:v5v12-Hodge}
There is an orthogonal direct sum
\begin{equation}
 C^0
 =
 \operatorname{Ran}Q_{-1}
 \oplus\mathcal H^0
 \oplus\operatorname{Ran}Q_0^*.
 \label{eq:v5v12-Hodge-decomposition}
\end{equation}
The natural map
\(\mathcal H^0\to H^0(Q)\) is an isomorphism.
\end{theorem}

\begin{proof}
The complex identity gives
\[
 \operatorname{Ran}Q_{-1}\perp\operatorname{Ran}Q_0^*.
\]
The orthogonal complement of their sum is
\[
 \ker Q_{-1}^*\cap\ker Q_0=\mathcal H^0,
\]
proving \eqref{eq:v5v12-Hodge-decomposition}.  If \(x\in\ker Q_0\),
its coexact component in \(\operatorname{Ran}Q_0^*\) vanishes: writing
that component as \(Q_0^*v\), orthogonality gives
\[
 \|Q_0^*v\|^2
 =\langle v,Q_0Q_0^*v\rangle
 =\langle Q_0^*v,x\rangle=0.
\]
Thus every cocycle is a unique sum of an exact vector and a harmonic
vector.  Harmonic vectors cannot be nonzero exact vectors because the sum
is orthogonal.  This proves the cohomology isomorphism.
\end{proof}

Let \(\Pi_{\mathcal H}\) be the orthogonal projection onto
\(\mathcal H^0\).  Let \(G_0\) be \(\Delta_0^{-1}\) on
\((\mathcal H^0)^\perp\) and zero on \(\mathcal H^0\).

\begin{proposition}[Explicit contracting homotopy]
\label{prop:v5v12-contracting}
On \(C^0\),
\begin{equation}
 I-\Pi_{\mathcal H}
 =
 Q_{-1}Q_{-1}^*G_0+Q_0^*Q_0G_0.
 \label{eq:v5v12-contracting}
\end{equation}
If the smallest nonzero eigenvalue of \(\Delta_0\) is \(\lambda_{\rm H}\),
then
\begin{equation}
 \|G_0\|=\lambda_{\rm H}^{-1}.
 \label{eq:v5v12-G-bound}
\end{equation}
\end{proposition}

\begin{proof}
The two summands in \(\Delta_0\) preserve the three orthogonal spaces in
\eqref{eq:v5v12-Hodge-decomposition}.  On the exact range only
\(Q_{-1}Q_{-1}^*\) acts nontrivially, on the coexact range only
\(Q_0^*Q_0\) acts nontrivially, and both vanish on the harmonic space.
Thus \(\Delta_0G_0=I-\Pi_{\mathcal H}\), which is
\eqref{eq:v5v12-contracting}.  The norm of the inverse on the nonzero
spectral subspace is the reciprocal of its smallest eigenvalue.
\end{proof}

A cofinal Hodge reduction needs a cutoff-independent lower bound on
\(\lambda_{\rm H}\), or a weighted substitute.  Pointwise finite-dimensional
decomposition alone gives no uniform control of the contracting homotopy.

\subsection{The exact quartet criterion}

For \(c\in C^0\), call its exact and coexact components auxiliary and its
harmonic component physical.

\begin{theorem}[Conformal quartet test]
\label{thm:v5v12-quartet-test}
For \(c\in C^0\), the following are equivalent:
\begin{enumerate}
 \item \(c\) has no physical component;
 \item \(\Pi_{\mathcal H}c=0\);
 \item
 \begin{equation}
 c=Q_{-1}u+Q_0^*v
 \label{eq:v5v12-quartet-decomposition}
 \end{equation}
 for some \(u\in C^{-1}\) and \(v\in C^1\);
 \item every linear functional annihilating
 \(\operatorname{Ran}Q_{-1}+\operatorname{Ran}Q_0^*\) also annihilates
 \(c\).
\end{enumerate}
\end{theorem}

\begin{proof}
The equivalence of items 1--3 is exactly the orthogonal decomposition
\eqref{eq:v5v12-Hodge-decomposition}.  If item 3 holds, item 4 is
immediate.  If item 3 fails, \(\Pi_{\mathcal H}c\ne0\); the functional
\[
 \ell_c(x)=\langle\Pi_{\mathcal H}c,x\rangle
\]
annihilates the exact and coexact ranges but satisfies
\(\ell_c(c)=\|\Pi_{\mathcal H}c\|^2>0\), contradicting item 4.
\end{proof}

\begin{corollary}[Diffeomorphism nonexactness is not the full test]
\label{cor:v5v12-nonexact-not-enough}
The V9 conclusion
\(c\notin\operatorname{Ran}Q_{-1}\) does not imply that \(c\) is physical.
It can still have \(\Pi_{\mathcal H}c=0\) by possessing a nonzero coexact
constraint component \(Q_0^*v\).
\end{corollary}

\begin{proof}
Apply item 3 of Theorem~\ref{thm:v5v12-quartet-test}.
\end{proof}

This is the precise upstream correction: gauge slicing and constraint
reduction are different operations.

\subsection{When an indefinite auxiliary integral disappears}

Let
\[
 C^0=\mathcal H^0\oplus\mathcal A,
 \qquad
 \mathcal A=(\mathcal H^0)^\perp.
\]
Suppose a finite-cutoff integration cycle factorizes as
\(\Gamma=\mathcal H^0\times\Gamma_{\mathcal A}\), and the action has the
form
\begin{equation}
 S(h,a)=\frac12\langle h,A_{\mathcal H}h\rangle+S_{\mathcal A}(a),
 \label{eq:v5v12-split-action}
\end{equation}
where \(A_{\mathcal H}>0\) and
\[
 0<|Z_{\mathcal A}|
 :=
 \left|\int_{\Gamma_{\mathcal A}}e^{-S_{\mathcal A}(a)}\,da\right|
 <\infty.
\]

\begin{theorem}[Normalized physical factorization]
\label{thm:v5v12-physical-factorization}
For every integrable observable \(F\) depending only on \(h\),
\begin{equation}
 \langle F\rangle_\Gamma
 =
 \frac{\displaystyle
 \int_{\mathcal H^0}F(h)e^{-\frac12\langle h,A_{\mathcal H}h\rangle}\,dh}
 {\displaystyle
 \int_{\mathcal H^0}e^{-\frac12\langle h,A_{\mathcal H}h\rangle}\,dh}.
 \label{eq:v5v12-factorization}
\end{equation}
The auxiliary phase and every indefinite auxiliary covariance cancel from
normalized physical expectations.
\end{theorem}

\begin{proof}
Fubini's theorem applies by the assumed integrability and factorized
cycle.  Both numerator and denominator contain the same nonzero factor
\(Z_{\mathcal A}\), which cancels.
\end{proof}

\begin{proposition}[Reflection positivity descends under factorization]
\label{prop:v5v12-OS-descent}
Assume an antilinear reflection preserves \(\mathcal H^0\) and
\(\Gamma_{\mathcal A}\), and the harmonic Gaussian is reflection positive
on the physical positive-time algebra.  Then the complete normalized
functional is reflection positive on observables depending only on
\(\mathcal H^0\), regardless of the sign of the auxiliary Hessian.
\end{proposition}

\begin{proof}
For every physical \(F\), the OS expectation
\(\langle F^\theta F\rangle_\Gamma\) reduces by
Theorem~\ref{thm:v5v12-physical-factorization} to the harmonic Gaussian OS
form, which is nonnegative by assumption.
\end{proof}

\begin{firewall}[Factorization is a theorem hypothesis]
\label{fw:v5v12-factorization}
BRST terminology does not prove
\eqref{eq:v5v12-split-action}.  Harmonic--auxiliary mixing, a
nonfactorizing contour, a zero auxiliary partition function, an anomaly,
or an observable that depends on auxiliary representatives invalidates
Theorem~\ref{thm:v5v12-physical-factorization}.
\end{firewall}

\subsection{Block diagonalization from the complex}

Let \(A:C^0\to C^0\) be a symmetric quadratic operator.

\begin{proposition}[Hodge-compatible Hessian]
\label{prop:v5v12-compatible-Hessian}
If
\begin{equation}
 [A,\Pi_{\mathcal H}]=0,
 \label{eq:v5v12-A-commutes}
\end{equation}
then \(A\) has no harmonic--auxiliary cross term and
\[
 A=A_{\mathcal H}\oplus A_{\mathcal A}.
\]
If the action and cycle also factorize, the physical Gaussian covariance is
\(A_{\mathcal H}^{-1}\).
\end{proposition}

\begin{proof}
Commutation with an orthogonal projection makes both its range and kernel
invariant under the symmetric operator.  Hence the off-diagonal blocks
vanish.  The covariance conclusion is the ordinary Gaussian calculation
on the harmonic block.
\end{proof}

If \(A\) does not commute with \(\Pi_{\mathcal H}\), integrating out an
invertible auxiliary block gives a harmonic Schur complement.  The
two-sided weighted estimate of V10 is then required; replacing the
auxiliary block by zero is not a cohomological reduction.

\begin{theorem}[Auxiliary factors cannot hide a physical negative mode]
\label{thm:v5v12-physical-negative}
Under the factorization hypotheses, if
\(A_{\mathcal H}\) has a negative direction seen by the physical
observable algebra, no nonzero auxiliary determinant or contour phase can
make the normalized physical Gaussian positive.
\end{theorem}

\begin{proof}
The normalized physical two-point function is
\(A_{\mathcal H}^{-1}\) by
Theorem~\ref{thm:v5v12-physical-factorization}.  A negative eigenvector of
\(A_{\mathcal H}\) yields a negative square expectation in that direction.
The auxiliary factor has already canceled from the ratio and cannot alter
this value.
\end{proof}

\subsection{Cofinal stability of the physical quotient}

Let the cutoff complexes be indexed by \(\Lambda\), with embeddings or
shell projections \(J_{\Lambda'\Lambda}\).

\begin{definition}[Cofinal Hodge--BRST certificate]
\label{def:v5v12-CHBC}
A CHBC consists of:
\begin{enumerate}
 \item complexes \eqref{eq:v5v12-complex} commuting with the shell maps;
 \item a uniform nonharmonic Hodge gap, or a weighted bound on
 \(G_{0,\Lambda}\);
 \item convergence of the harmonic projections in the physical kernel
 norm;
 \item a Ward-complete factorization or a certified weighted Schur
 reduction of the auxiliary variables;
 \item OS positivity on the harmonic cylinder algebra and compatibility
 of reflection with the complexes;
 \item \(\Pi_{\mathcal H,\Lambda}c_\Lambda=0\) for every claimed
 unphysical conformal coordinate.
\end{enumerate}
\end{definition}

\begin{theorem}[Conditional physical positivity]
\label{thm:v5v12-CHBC}
A CHBC with a positive harmonic covariance and uniform physical moment
bounds removes the auxiliary conformal variables from normalized physical
Schwinger functions and supplies the linear metric-positivity part of the
V8 MCC.
\end{theorem}

\begin{proof}
Items 1--3 stabilize the Hodge decomposition under cutoff passage.
Item 4 permits
Theorem~\ref{thm:v5v12-physical-factorization}, or its V10 Schur
replacement, uniformly.  Item 5 gives OS positivity of the physical
functional by Proposition~\ref{prop:v5v12-OS-descent}.  Item 6 removes the
claimed conformal coordinates from the harmonic algebra.
\end{proof}

No CHBC has been proved for the nonlinear SM--Einstein measure.  The V9
calculation excludes only the assertion that nonconstant conformal
variations are diffeomorphism exact.  It leaves open the concrete,
testable possibility that the Hamiltonian constraint places them in the
coexact sector.

V13 will test that possibility for linearized ADM gravity on a flat torus.
It will solve the Fourier-space momentum and Hamiltonian constraints,
identify the transverse-traceless harmonic quotient, and determine
whether nonzero conformal modes survive in the physical quadratic
Hamiltonian.  The result will be kept separate from the nonlinear
Euclidean measure problem.


\clearpage
\section{V13: linearized ADM reduction and the physical graviton block}
\label{sec:v5v13}

V12 left open whether the Hamiltonian constraint, rather than an ordinary
diffeomorphism slice, removes the conformal mode.  This note answers that
question for linearized vacuum gravity about a flat three-torus.  Every
nonzero spatial Fourier mode reduces to two transverse-traceless
polarizations with a positive oscillator Hamiltonian.  A nonzero conformal
spatial mode violates the linearized Hamiltonian constraint and is
therefore absent from the physical phase space.  The homogeneous volume
mode and the nonlinear, sourced constraint remain separate.

\subsection{Linearized variables and constraints}

Let the spatial background be the flat torus
\((\mathbb T^3,\delta_{ij})\).  Write
\[
 \gamma_{ij}=\delta_{ij}+h_{ij},
 \qquad
 \pi^{ij}=\text{linearized canonical momentum}.
\]
At first order, the vacuum ADM constraints are, up to nonzero conventional
normalizations,
\begin{align}
 \mathcal C(h)
 &=
 \partial_i\partial_jh_{ij}-\Delta h=0,
 \qquad h=\delta^{ij}h_{ij},
 \label{eq:v5v13-Hamiltonian-constraint}\\
 \mathcal C_i(\pi)
 &=
 \partial_j\pi_{ij}=0.
 \label{eq:v5v13-momentum-constraint}
\end{align}
Spatial diffeomorphisms act by
\begin{equation}
 \delta_\xi h_{ij}=\partial_i\xi_j+\partial_j\xi_i.
 \label{eq:v5v13-spatial-gauge}
\end{equation}
The normal deformation generated by
\(\mathcal C(h)\) acts on momentum, at the same linearized level, by
\begin{equation}
 \delta_\eta\pi_{ij}
 =
 \bigl(\partial_i\partial_j-\delta_{ij}\Delta\bigr)\eta
 \label{eq:v5v13-normal-gauge}
\end{equation}
up to an overall nonzero sign.

Fix a nonzero Fourier wavevector \(k\).  Rotate coordinates so that
\(k=(|k|,0,0)\), and use indices \(A,B\in\{2,3\}\).

\begin{lemma}[Fourier form of the constraints]
\label{lem:v5v13-Fourier-constraints}
At the fixed nonzero mode,
\begin{align}
 \mathcal C(h)=0
 &\quad\Longleftrightarrow\quad
 h_{22}+h_{33}=0,
 \label{eq:v5v13-transverse-trace}\\
 \mathcal C_i(\pi)=0
 &\quad\Longleftrightarrow\quad
 \pi_{11}=\pi_{12}=\pi_{13}=0.
 \label{eq:v5v13-momentum-components}
\end{align}
\end{lemma}

\begin{proof}
Fourier transformation replaces \(\partial_i\) by \(ik_i\).  Therefore
\[
 \widehat{\mathcal C}
 =-k_ik_jh_{ij}+|k|^2h
 =|k|^2(h_{22}+h_{33}),
\]
which proves \eqref{eq:v5v13-transverse-trace}.  The momentum constraint is
\(k_j\pi_{ij}=|k|\pi_{i1}=0\), proving
\eqref{eq:v5v13-momentum-components}.
\end{proof}

\subsection{Exact mode-by-mode reduction}

For the chosen wavevector, the spatial gauge transformation gives
\begin{align}
 \delta h_{11}&=2i|k|\xi_1,
 &
 \delta h_{1A}&=i|k|\xi_A,
 &
 \delta h_{AB}&=0.
 \label{eq:v5v13-gauge-components}
\end{align}

\begin{theorem}[Nonzero-mode ADM quotient]
\label{thm:v5v13-ADM-quotient}
For every \(k\ne0\), each constraint and gauge orbit has a unique
representative satisfying
\begin{align}
 h_{1i}&=0,
 &
 h_{22}+h_{33}&=0,
 \label{eq:v5v13-h-TT}\\
 \pi_{1i}&=0,
 &
 \pi_{22}+\pi_{33}&=0.
 \label{eq:v5v13-pi-TT}
\end{align}
The remaining independent canonical pairs are
\begin{align}
 h_+&=\frac1{\sqrt2}(h_{22}-h_{33}),
 &
 h_\times&=\sqrt2\,h_{23},
 \label{eq:v5v13-h-polarizations}\\
 \pi_+&=\frac1{\sqrt2}(\pi_{22}-\pi_{33}),
 &
 \pi_\times&=\sqrt2\,\pi_{23}.
 \label{eq:v5v13-pi-polarizations}
\end{align}
Thus the reduced phase space has exactly two graviton polarizations at
each nonzero \(k\).
\end{theorem}

\begin{proof}
Because \(|k|>0\), equation \eqref{eq:v5v13-gauge-components} uniquely
chooses \(\xi_1,\xi_2,\xi_3\) to set \(h_{11},h_{12},h_{13}\) to zero.
The Hamiltonian constraint then imposes the transverse trace condition in
\eqref{eq:v5v13-h-TT}.  The momentum constraint already gives
\(\pi_{1i}=0\).  In Fourier variables, the normal gauge transformation
\eqref{eq:v5v13-normal-gauge} has
\[
 \delta_\eta\pi_{11}=0,
 \qquad
 \delta_\eta\pi_{AB}=|k|^2\delta_{AB}\eta
\]
up to the common sign.  It therefore fixes the remaining transverse trace
\(\pi_{22}+\pi_{33}\) uniquely.  What remains is a symmetric trace-free
\(2\times2\) tensor in each of \(h\) and \(\pi\), with the two components
listed in \eqref{eq:v5v13-h-polarizations} and
\eqref{eq:v5v13-pi-polarizations}.
\end{proof}

\begin{corollary}[The nonzero conformal mode is constrained out]
\label{cor:v5v13-conformal-out}
For a pure spatial conformal mode
\[
 h_{ij}=2\phi\,\delta_{ij}
\]
at \(k\ne0\), the Hamiltonian constraint gives
\begin{equation}
 4|k|^2\phi=0.
 \label{eq:v5v13-conformal-constraint}
\end{equation}
Hence \(\phi=0\).  The mode is absent from the physical quotient even
though it is not a spatial diffeomorphism direction.
\end{corollary}

\begin{proof}
Substitution gives \(h_{22}+h_{33}=4\phi\).  Apply
Lemma~\ref{lem:v5v13-Fourier-constraints}.
\end{proof}

This realizes the coexact alternative in
Corollary~\ref{cor:v5v12-nonexact-not-enough}: the conformal coordinate is
removed by the Hamiltonian constraint, not by the spatial gauge orbit.

\subsection{Positive physical Hamiltonian}

The quadratic ADM Hamiltonian density about flat space can be normalized
so that its restriction to transverse-traceless variables is
\begin{equation}
 H^{(2)}_{\rm phys}
 =
 \frac12\sum_{k\ne0}\sum_{\lambda\in\{+,\times\}}
 \left(
  |\pi_\lambda(k)|^2
  +|k|^2|h_\lambda(k)|^2
 \right).
 \label{eq:v5v13-positive-Hamiltonian}
\end{equation}
Reality imposes
\(h_\lambda(-k)=\overline{h_\lambda(k)}\) and similarly for momentum.

\begin{theorem}[Positive linearized physical block]
\label{thm:v5v13-positive-block}
After solving the constraints and quotienting the gauge transformations,
the nonzero-mode vacuum ADM Hamiltonian is positive definite and consists
of two massless harmonic oscillators per spatial wavevector.
\end{theorem}

\begin{proof}
The standard quadratic ADM Hamiltonian before reduction is
\begin{align*}
 \int_{\mathbb T^3}\Bigl[
 &\pi_{ij}\pi_{ij}-\frac12\pi^2
 +\frac14(\partial_kh_{ij})(\partial_kh_{ij})
 -\frac14(\partial_kh)(\partial_kh)\\
 &+\frac12(\partial_ih_{ij})(\partial_jh)
 -\frac12(\partial_ih_{ij})(\partial_kh_{kj})
 \Bigr]\,dx,
\end{align*}
up to an overall positive canonical normalization and terms multiplying
the linear constraints.  On
\eqref{eq:v5v13-h-TT}--\eqref{eq:v5v13-pi-TT}, the traces and divergences
vanish.  The surviving terms are positive squares,
\[
 \int\left(\pi_{ij}^{\rm TT}\pi_{ij}^{\rm TT}
 +\frac14\partial_kh_{ij}^{\rm TT}
                  \partial_kh_{ij}^{\rm TT}\right)dx.
\]
A constant symplectic rescaling of \(h^{\rm TT}\) and
\(\pi^{\rm TT}\) puts each mode in the form
\eqref{eq:v5v13-positive-Hamiltonian}.
\end{proof}

\subsection{OS positivity of the free physical modes}

Euclidean time evolution of each physical oscillator has covariance
\begin{equation}
 C_k(t-s)=\frac{e^{-|k||t-s|}}{2|k|}.
 \label{eq:v5v13-oscillator-covariance}
\end{equation}

\begin{theorem}[Free-graviton physical reflection positivity]
\label{thm:v5v13-free-OS}
The centered Gaussian functional generated by
\eqref{eq:v5v13-oscillator-covariance}, summed over the two TT
polarizations and nonzero torus modes, is reflection positive on the TT
positive-time cylinder algebra.
\end{theorem}

\begin{proof}
For a positive-time test family \(f_\lambda(t,k)\), the reflected covariance
form is
\begin{equation}
 \sum_{k\ne0}\sum_{\lambda}
 \frac1{2|k|}
 \left|
  \int_0^\infty e^{-|k|t}f_\lambda(t,k)\,dt
 \right|^2
 \ge0.
 \label{eq:v5v13-TT-OS-form}
\end{equation}
Gaussian reflection positivity follows from the same exponential-vector
argument used in Theorem~\ref{thm:v5v11-OS-mixture}.
\end{proof}

\begin{corollary}[Raw covariant negativity need not be physical negativity]
\label{cor:v5v13-raw-vs-physical}
At the free linearized level, the negative covariant conformal Hessian from
V9 does not appear in the constrained TT Schwinger functional.  The
physical nonzero-mode covariance is positive and OS positive.
\end{corollary}

\begin{proof}
Corollary~\ref{cor:v5v13-conformal-out} removes the nonzero conformal mode,
while Theorems~\ref{thm:v5v13-positive-block} and
\ref{thm:v5v13-free-OS} give the physical assertions.
\end{proof}

\subsection{The homogeneous and sourced sectors}

At \(k=0\), the symbols in
\eqref{eq:v5v13-Hamiltonian-constraint}--\eqref{eq:v5v13-normal-gauge}
vanish.  The preceding reduction therefore does not remove the
homogeneous volume mode.

\begin{firewall}[Zero modes require a separate global reduction]
\label{fw:v5v13-zero-mode}
Theorem~\ref{thm:v5v13-ADM-quotient} uses division by \(|k|\) and
\(|k|^2\).  It makes no claim about torus moduli, the homogeneous volume,
global lapse constraints, or cosmological-constant sectors.
\end{firewall}

With matter, the linear Hamiltonian constraint has the schematic form
\begin{equation}
 \partial_i\partial_jh_{ij}-\Delta h
 =\kappa_{\rm N}\rho_{\rm matter}.
 \label{eq:v5v13-sourced-constraint}
\end{equation}
A conformal component is then determined nonlocally by the source rather
than set to zero.

\begin{proposition}[Sourced conformal determination]
\label{prop:v5v13-sourced-conformal}
For a nonzero Fourier mode in the pure conformal ansatz,
\eqref{eq:v5v13-sourced-constraint} determines
\begin{equation}
 \phi(k)
 =
 \frac{\kappa_{\rm N}}{4|k|^2}\rho_{\rm matter}(k)
 \label{eq:v5v13-Poisson-conformal}
\end{equation}
up to the sign convention absorbed in \(\kappa_{\rm N}\).
It is not an independent Gaussian coordinate.
\end{proposition}

\begin{proof}
The left side was computed in
\eqref{eq:v5v13-conformal-constraint}.  Solve the resulting scalar
equation for \(\phi(k)\).
\end{proof}

The inverse Laplacian in \eqref{eq:v5v13-Poisson-conformal} creates a
nonlocal instantaneous interaction.  Positivity and stability of the
complete matter--gravity Hamiltonian cannot be inferred from the vacuum TT
block alone.

\subsection{Linearized physical-gravity certificate}

\begin{definition}[Linear constrained graviton certificate]
\label{def:v5v13-LCGC}
An LCGC consists of:
\begin{enumerate}
 \item exact solution of all nonzero-mode linear constraints and gauge
 orbits;
 \item the positive TT Hamiltonian
 \eqref{eq:v5v13-positive-Hamiltonian};
 \item the TT reflection form \eqref{eq:v5v13-TT-OS-form};
 \item a separate finite-dimensional treatment of global zero modes;
 \item a source map controlling the inverse-Laplacian constraint solution
 in the chosen critical norms.
\end{enumerate}
\end{definition}

\begin{theorem}[What linear reduction supplies]
\label{thm:v5v13-LCGC}
The vacuum nonzero-mode sector on a flat torus satisfies items 1--3 of the
LCGC.  This supplies the Gaussian physical metric block required by the
V12 CHBC at linear order.  It does not supply items 4--5, a nonlinear
constraint reduction, an interacting Euclidean measure, or the
curvature-operator tail bounds of V5--V7.
\end{theorem}

\begin{proof}
Items 1--3 are
Theorems~\ref{thm:v5v13-ADM-quotient},
\ref{thm:v5v13-positive-block}, and
\ref{thm:v5v13-free-OS}.  The zero-mode exclusion and sourced inverse are
the firewalls and Proposition~\ref{prop:v5v13-sourced-conformal}; none of
the nonlinear or local-tail conclusions follows from a quadratic Fourier
calculation.
\end{proof}

V13 establishes a genuine positive result: the free physical graviton
sector has no nonzero conformal degree of freedom and is OS positive after
canonical constraint reduction.  The metric bottleneck has moved from the
quadratic physical propagator to the nonlinear solution of the
Hamiltonian and momentum constraints, their Jacobian, zero modes, and
compatibility with covariant shell integration.

V14 will formulate a quantitative implicit-function theorem for the
nonlinear constraint map on a fixed gauge slice.  It will give a radius and
Lipschitz bound for solving the auxiliary metric variables as functions of
TT and Standard Model data, and will identify the elliptic inverse norm
that must remain uniform along the shell trajectory.


\clearpage
\section{V14: quantitative nonlinear constraint elimination}
\label{sec:v5v14}

V13 solved the vacuum linearized ADM constraints mode by mode.  To use
that mechanism in an interacting shell construction, the auxiliary metric
variables must be solved as functions of the transverse gravitational and
Standard Model data with constants uniform in the cutoff.  This note gives
a quantitative implicit theorem, its first two derivative bounds, the
finite-dimensional constraint Jacobian, and the exact elliptic inverse
that becomes the next analytic gate.

\subsection{A contraction form of the constraint equation}

Let \(X\) be a Banach space of physical data, \(Z\) a Banach space of
auxiliary metric variables, and \(Y\) a Banach space of constraints.  Let
\[
 \mathcal C:U_X\times U_Z\longrightarrow Y
\]
be \(C^1\), with
\[
 \mathcal C(0,0)=0.
\]
Assume
\begin{equation}
 L:=D_z\mathcal C(0,0):Z\longrightarrow Y
 \label{eq:v5v14-linear-constraint}
\end{equation}
is an isomorphism and
\begin{equation}
 \|L^{-1}\|_{Y\to Z}\le M.
 \label{eq:v5v14-inverse-bound}
\end{equation}
For fixed \(x\), define
\begin{equation}
 \mathcal T_x(z)=z-L^{-1}\mathcal C(x,z).
 \label{eq:v5v14-Newton-map}
\end{equation}
Its fixed points are exactly the solutions of
\(\mathcal C(x,z)=0\).

\begin{theorem}[Quantitative nonlinear constraint solver]
\label{thm:v5v14-constraint-solver}
Let \(D\subset X\) be a parameter set and let
\(\overline B_R^Z\subset U_Z\).  Suppose
\begin{align}
 \sup_{\substack{x\in D\\\|z\|_Z\le R}}
 \|I-L^{-1}D_z\mathcal C(x,z)\|_{Z\to Z}
 &\le q<1,
 \label{eq:v5v14-derivative-small}\\
 \eta
 :=\sup_{x\in D}\|L^{-1}\mathcal C(x,0)\|_Z
 &\le(1-q)R.
 \label{eq:v5v14-residual-radius}
\end{align}
Then, for every \(x\in D\), there is a unique
\(z=Z(x)\in\overline B_R^Z\) satisfying
\(\mathcal C(x,Z(x))=0\), and
\begin{equation}
 \|Z(x)\|_Z
 \le\frac{\|L^{-1}\mathcal C(x,0)\|_Z}{1-q}
 \le\frac{\eta}{1-q}.
 \label{eq:v5v14-solution-bound}
\end{equation}
\end{theorem}

\begin{proof}
The derivative of \(\mathcal T_x\) in \(z\) has norm at most \(q\), so the
mean-value formula makes it a contraction on the ball.  Moreover,
\[
 \|\mathcal T_x(z)\|
 \le\|\mathcal T_x(z)-\mathcal T_x(0)\|
     +\|\mathcal T_x(0)\|
 \le qR+\eta\le R.
\]
Banach's theorem gives the unique fixed point.  Applying the contraction
estimate to \(Z(x)\) and zero yields
\[
 \|Z(x)\|
 \le q\|Z(x)\|+\|L^{-1}\mathcal C(x,0)\|,
\]
which is \eqref{eq:v5v14-solution-bound}.
\end{proof}

\begin{corollary}[Lipschitz dependence on physical data]
\label{cor:v5v14-Lipschitz}
If, in addition,
\begin{equation}
 \|L^{-1}(\mathcal C(x,z)-\mathcal C(\widetilde x,z))\|_Z
 \le K_x\|x-\widetilde x\|_X
 \label{eq:v5v14-x-Lipschitz}
\end{equation}
uniformly on the ball, then
\begin{equation}
 \|Z(x)-Z(\widetilde x)\|_Z
 \le\frac{K_x}{1-q}\|x-\widetilde x\|_X.
 \label{eq:v5v14-Z-Lipschitz}
\end{equation}
\end{corollary}

\begin{proof}
Insert and subtract
\(\mathcal T_x(Z(\widetilde x))\) between the two fixed-point equations.
The \(z\)-difference costs \(q\), while the \(x\)-difference is bounded by
\eqref{eq:v5v14-x-Lipschitz}.  Rearrangement gives the result.
\end{proof}

\subsection{Marked constraint derivatives}

Assume now that \(\mathcal C\) is \(C^2\).  At a solved point write
\[
 A(x)=D_z\mathcal C(x,Z(x)).
\]
Equation \eqref{eq:v5v14-derivative-small} gives
\begin{equation}
 \|A(x)^{-1}\|_{Y\to Z}
 \le\frac{M}{1-q}.
 \label{eq:v5v14-A-inverse}
\end{equation}

\begin{theorem}[Two marked derivatives of the constraint solution]
\label{thm:v5v14-two-marks}
Suppose
\[
 \|D_x\mathcal C\|_{X\to Y}\le B_1
\]
and every second partial derivative of \(\mathcal C\), with the
corresponding \(X,Z\) product norms, is bounded by \(B_2\).  Then \(Z\) is
\(C^2\) and
\begin{equation}
 DZ(x)=-A(x)^{-1}D_x\mathcal C(x,Z(x)),
 \qquad
 \|DZ(x)\|\le d_1:=\frac{MB_1}{1-q}.
 \label{eq:v5v14-first-mark}
\end{equation}
For \(u,v\in X\),
\begin{align}
 D^2Z(x)[u,v]
 =-A(x)^{-1}\Bigl(
 &D_{xx}^2\mathcal C[u,v]
 +D_{xz}^2\mathcal C[u,DZ\,v]\notag\\
 &+D_{zx}^2\mathcal C[DZ\,u,v]
 +D_{zz}^2\mathcal C[DZ\,u,DZ\,v]
 \Bigr),
 \label{eq:v5v14-second-mark}
\end{align}
and
\begin{equation}
 \|D^2Z(x)\|
 \le
 \frac{MB_2}{1-q}(1+2d_1+d_1^2).
 \label{eq:v5v14-second-mark-bound}
\end{equation}
\end{theorem}

\begin{proof}
The implicit-function theorem applies because \(A(x)\) is invertible.
Differentiate
\[
 \mathcal C(x,Z(x))=0
\]
once to obtain
\(D_x\mathcal C+A\,DZ=0\), which gives
\eqref{eq:v5v14-first-mark}.  Differentiate that identity in a second
direction, keeping both ordered mixed derivatives, to obtain
\eqref{eq:v5v14-second-mark}.  Use
\eqref{eq:v5v14-A-inverse}, the \(B_2\) bounds, and
\(\|DZ\|\le d_1\) to get
\eqref{eq:v5v14-second-mark-bound}.
\end{proof}

These are the constraint analogues of the marked shell bounds in V2.
They are needed because stress tensors and response observables
differentiate the solved auxiliary metric with respect to matter and TT
sources.

\subsection{The finite-dimensional constraint Jacobian}

At a finite cutoff, suppose \(Z=Y=\mathbb R^m\) and the hypotheses of
Theorem~\ref{thm:v5v14-constraint-solver} hold.  Let
\(F(x,z)\) be continuous and integrable.

\begin{theorem}[Exact delta-function reduction]
\label{thm:v5v14-delta-reduction}
If \(z=Z(x)\) is the only constraint root in the integration domain, then
\begin{equation}
 \int_{\mathbb R^m}
 F(x,z)\,
 \delta(\mathcal C(x,z))\,
 |\det D_z\mathcal C(x,z)|\,dz
 =
 F(x,Z(x)).
 \label{eq:v5v14-delta-reduction}
\end{equation}
\end{theorem}

\begin{proof}
The inverse-function theorem and
\eqref{eq:v5v14-A-inverse} make \(z\mapsto\mathcal C(x,z)\) a local
diffeomorphism at the root.  The global uniqueness in the stated
integration domain allows the change of variables
\(y=\mathcal C(x,z)\) on the root contribution.  Its Jacobian is
\(|\det D_z\mathcal C|\), which cancels the displayed determinant, and
\(\delta(y)\) evaluates the remaining function at \(y=0\).
\end{proof}

\begin{firewall}[Uniqueness on one ball is not a global gauge theorem]
\label{fw:v5v14-global-roots}
Theorem~\ref{thm:v5v14-delta-reduction} applies to the certified root
domain.  Additional roots outside that domain, Gribov copies, determinant
sign changes, and zero modes require separate treatment.  Replacing the
absolute determinant by an oriented or ghost determinant also requires a
proved orientation convention.
\end{firewall}

\subsection{Elliptic realization and scale dependence}

Let \(Z=H^{s+2}_0(\mathbb T^3)\) and
\(Y=H^s_0(\mathbb T^3)\), where the subscript removes the spatial mean.
For the flat conformal Hamiltonian constraint,
\[
 L=-4\Delta.
\]

\begin{proposition}[Flat nonzero-mode inverse]
\label{prop:v5v14-flat-inverse}
If \(k_{\min}>0\) is the smallest nonzero torus frequency, then
\begin{equation}
 \|L^{-1}f\|_{H^{s+2}}
 \le M_{\rm flat}\|f\|_{H^s},
 \qquad
 M_{\rm flat}\le C(1+k_{\min}^{-2}).
 \label{eq:v5v14-flat-elliptic}
\end{equation}
The constant is independent of the ultraviolet cutoff.
\end{proposition}

\begin{proof}
In Fourier variables,
\[
 \widehat{L^{-1}f}(k)=\frac{\widehat f(k)}{4|k|^2}
 \qquad(k\ne0).
\]
The ratio
\[
 \frac{(1+|k|^2)^{s+2}}{16|k|^4(1+|k|^2)^s}
\]
is uniformly bounded for \(|k|\ge k_{\min}\) by a constant of order
\((1+k_{\min}^{-2})^2\).  Taking square roots gives
\eqref{eq:v5v14-flat-elliptic}.
\end{proof}

\begin{corollary}[Infrared, not ultraviolet, loss of the flat inverse]
\label{cor:v5v14-IR-loss}
Adding higher Fourier modes does not enlarge \(M_{\rm flat}\), but sending
the spatial volume to infinity makes \(k_{\min}\downarrow0\) and destroys
the bound unless the source has stronger cancellation or a weighted
infrared norm is used.
\end{corollary}

\begin{proof}
This is immediate from
\eqref{eq:v5v14-flat-elliptic}.
\end{proof}

On a curved or interacting background, \(L\) is replaced by the
linearized lapse--conformal--longitudinal constraint operator.  Uniform
invertibility can fail through a Killing field, a conformal zero mode, a
change of ellipticity, or a small eigenvalue.  These are spectral facts,
not removable choices of norm.

\subsection{Reduced action and its Schur Hessian}

Suppose the auxiliary equation is a stationarity condition
\[
 \mathcal C(x,z)=D_zS(x,z)
\]
for a \(C^2\) action, and define
\[
 S_{\rm red}(x)=S(x,Z(x)).
\]

\begin{proposition}[Exact reduced Hessian]
\label{prop:v5v14-reduced-Hessian}
At a solved point,
\begin{equation}
 D^2S_{\rm red}
 =
 S_{xx}-S_{xz}S_{zz}^{-1}S_{zx}.
 \label{eq:v5v14-reduced-Schur}
\end{equation}
\end{proposition}

\begin{proof}
The first derivative is
\[
 DS_{\rm red}=S_x+S_zDZ=S_x
\]
because \(S_z=0\).  Differentiating again gives
\[
 D^2S_{\rm red}=S_{xx}+S_{xz}DZ.
\]
The first-mark identity for
\(\mathcal C=S_z\) gives
\(DZ=-S_{zz}^{-1}S_{zx}\), proving
\eqref{eq:v5v14-reduced-Schur}.
\end{proof}

Thus solving the constraint does not permit deletion of its response.
The V10 two-sided weighted Schur estimate must be applied to
\eqref{eq:v5v14-reduced-Schur} when positivity is claimed.

\subsection{The nonlinear constraint certificate}

\begin{definition}[Uniform nonlinear constraint certificate]
\label{def:v5v14-UNCC}
A UNCC along a shell trajectory consists of:
\begin{enumerate}
 \item a gauge-fixed constraint map and a declared separation
 \(X\oplus Z\);
 \item a mean-zero or otherwise globally fixed zero-mode convention;
 \item a cutoff-uniform inverse bound \eqref{eq:v5v14-inverse-bound};
 \item a common solution ball satisfying
 \eqref{eq:v5v14-derivative-small}--\eqref{eq:v5v14-residual-radius}
 through two marks;
 \item the finite-cutoff Jacobian, orientation, and all-root analysis;
 \item the reduced Schur and OS estimates for the physical action;
 \item compatibility of \(Z(x)\) with the Ward identities and shell
 projections.
\end{enumerate}
\end{definition}

\begin{theorem}[Conditional nonlinear physical reduction]
\label{thm:v5v14-UNCC}
A UNCC supplies a unique twice-marked auxiliary metric as a function of
the physical TT and Standard Model data, with the explicit bounds
\eqref{eq:v5v14-solution-bound},
\eqref{eq:v5v14-first-mark}, and
\eqref{eq:v5v14-second-mark-bound}.  Together with the V12 cohomology
factorization and V10 OS certificate for the reduced action, it removes
the solved conformal and longitudinal variables from the physical
finite-cutoff functional.
\end{theorem}

\begin{proof}
Existence, uniqueness, and derivative bounds are
Theorems~\ref{thm:v5v14-constraint-solver} and
\ref{thm:v5v14-two-marks}.  The Jacobian reduction is
Theorem~\ref{thm:v5v14-delta-reduction}.  The remaining two stated inputs
identify the physical algebra and prove positivity of its reduced
functional.
\end{proof}

No UNCC has been proved for the nonlinear coupled SM--Einstein constraint
map.  V14 nevertheless identifies a favorable scale fact: on a fixed flat
torus the basic inverse-Laplacian constant is ultraviolet-uniform.  The
new risks are the nonlinear derivative budget, background-dependent small
eigenvalues, global roots, and the thermodynamic infrared limit.

V15 will begin with the compulsory companion audit.  It will then derive a
source-space estimate for the quadratic and bilinear terms in the
constraint map, using Sobolev multiplication on the three-dimensional
slice, to determine a concrete small-data radius for
\eqref{eq:v5v14-derivative-small}.


\clearpage
\section{V15: companion audit, Sobolev constraint radius, and the quantum
regularity mismatch}
\label{sec:v5v15}

This is the third compulsory companion checkpoint of the fifth series.
The mathematical work then turns the abstract V14 derivative budget into
a concrete small-data radius on a fixed three-torus.  Sobolev multiplication
does give cutoff-uniform classical bounds.  The same calculation exposes a
new upstream obstruction: equal-time quantum fields have distributional
regularity far below the positive Sobolev algebra used by the classical
constraint theorem.  A continuum constraint map therefore requires
renormalized products, not only a smaller classical ball.

\subsection{Compulsory companion audit}

At the time of this audit, the newest live companion files were
\begin{align*}
 &\texttt{Renewal\_Yang\_Mills\_Mass\_Gap\_Research\_Notes\_V30.tex},\\
 &\hspace{1em}\text{SHA--256 }
 \texttt{69940730FFC793D014225869A6532FEEA7E2EF30660B94762150305BE4EDE522},
 \\
 &\texttt{Renewal\_NS\_Stability\_Research\_Notes\_V31.tex},\\
 &\hspace{1em}\text{SHA--256 }
 \texttt{F1121B62238AD5491BB7CF03635EA9934942EDF573ED8FAAA9EE79E1D38A6696}.
\end{align*}

Yang--Mills V30 closes arbitrary local incidence at three certified
coordinates by retaining complete moment partitions, reducible
complements, the unique payer, and arbitrary label-preserving output
dualizers.  Its remaining obstruction is a dimension-growing active row
outside those coordinates.  Navier--Stokes V31 still proves that lifting a
weak source to a critical first moment returns the missing critical
source weight.

\begin{theorem}[V15 audit transfer]
\label{thm:v5v15-audit}
For the nonlinear gravitational constraint map, the companion audit gives
the following acquisition rules:
\begin{enumerate}
 \item assemble all fixed-arity local products before estimating them in a
 Sobolev algebra;
 \item count each derivative loss and elliptic inverse once;
 \item retain any unbounded exterior curvature or source sector rather
 than absorbing it into a dimension-free constant; and
 \item test the actual quantum source regularity against the norm that pays
 the nonlinear product.
\end{enumerate}
\end{theorem}

\begin{proof}
Items 1 and 3 are the local-incidence and exterior-row separation proved in
Yang--Mills V30.  Items 2 and 4 are the constraint analogue of the exact
source rung in Navier--Stokes V31.  Each is a restriction on the estimates
used below, not a claim that the companion theorem itself proves a gravity
bound.
\end{proof}

\subsection{Three-dimensional Sobolev products}

Let \(s>3/2\) and work on a fixed compact three-manifold of bounded
geometry, initially \(\mathbb T^3\).

\begin{theorem}[Sobolev algebra and derivative products]
\label{thm:v5v15-Sobolev-algebra}
There is \(C_s<\infty\) such that
\begin{equation}
 \|uv\|_{H^s}
 \le C_s\|u\|_{H^s}\|v\|_{H^s}.
 \label{eq:v5v15-algebra}
\end{equation}
Consequently, for \(u,v\in H^{s+2}\),
\begin{align}
 \|u\,\nabla^2v\|_{H^s}
 &\le C_s'\|u\|_{H^{s+2}}\|v\|_{H^{s+2}},
 \label{eq:v5v15-uD2v}\\
 \|\nabla u\,\nabla v\|_{H^s}
 &\le C_s'\|u\|_{H^{s+2}}\|v\|_{H^{s+2}}.
 \label{eq:v5v15-DuDv}
\end{align}
\end{theorem}

\begin{proof}
For \(s>3/2\), Sobolev embedding gives
\(H^s\hookrightarrow L^\infty\), and the fractional Leibniz estimate gives
\[
 \|uv\|_{H^s}
 \le C\bigl(\|u\|_{L^\infty}\|v\|_{H^s}
           +\|v\|_{L^\infty}\|u\|_{H^s}\bigr),
\]
which proves \eqref{eq:v5v15-algebra}.  Since
\(H^{s+2}\hookrightarrow H^s\) and
\(\nabla^2v\in H^s\), equation
\eqref{eq:v5v15-uD2v} follows.  For
\eqref{eq:v5v15-DuDv}, both first derivatives lie in
\(H^{s+1}\hookrightarrow H^s\), and another application of
\eqref{eq:v5v15-algebra} completes the proof.
\end{proof}

Let \(P_\Lambda\) be the Fourier projection to \(|k|\le\Lambda\).

\begin{corollary}[Cutoff-uniform projected products]
\label{cor:v5v15-projected-products}
The constants in
\eqref{eq:v5v15-algebra}--\eqref{eq:v5v15-DuDv} remain unchanged when
each output is replaced by \(P_\Lambda\) and each input is cutoff
supported.
\end{corollary}

\begin{proof}
The Fourier projection has operator norm one on every \(H^r\).  Apply it
after the unprojected estimates.
\end{proof}

Thus the local product estimate itself has no ultraviolet dimension
factor.  This is the constraint analogue of retaining the complete local
incidence before taking the norm.

\subsection{A bilinear majorant for the constraint}

Let
\[
 X=H^{s+2}_{\rm TT}\oplus\mathcal X_{\rm SM},
 \qquad
 Z=H^{s+2}_{0,\rm aux},
 \qquad
 Y=H^s_0,
\]
where the zero-mode convention is fixed and the Standard Model norm is
chosen so that its classical stress source lies in \(H^s\).  Suppose the
constraint has the assembled form
\begin{equation}
 \mathcal C(x,z)
 =
 Lz+\mathcal N_{zz}(z,z)
 +\mathcal N_{xz}(x,z)
 +\mathcal S(x),
 \label{eq:v5v15-constraint-majorant}
\end{equation}
where the bilinear maps obey
\begin{align}
 \|\mathcal N_{zz}(z_1,z_2)\|_Y
 &\le C_{zz}\|z_1\|_Z\|z_2\|_Z,
 \label{eq:v5v15-Nzz}\\
 \|\mathcal N_{xz}(x,z)\|_Y
 &\le C_{xz}\|x\|_X\|z\|_Z,
 \label{eq:v5v15-Nxz}\\
 \|\mathcal S(x)\|_Y
 &\le C_{\rm s}\bigl(\|x\|_X+\|x\|_X^2\bigr).
 \label{eq:v5v15-source}
\end{align}
On a small metric ball, the inverse metric and volume density are expanded
by convergent Neumann and analytic composition series; Theorem
\ref{thm:v5v15-Sobolev-algebra} bounds every fixed-arity term of this form.

\begin{theorem}[Explicit classical constraint radius]
\label{thm:v5v15-explicit-radius}
Assume \(\|L^{-1}\|\le M\), and let \(R_{\max}\) be a radius on which
\eqref{eq:v5v15-constraint-majorant}--\eqref{eq:v5v15-source} hold.  Put
\begin{align}
 C_0&=2C_{\rm s},
 \label{eq:v5v15-C0}\\
 K_0&=M\bigl(4MC_{zz}C_0+C_{xz}\bigr),
 \label{eq:v5v15-K0}\\
 r_*&=\min\left\{
  1,\,
  \frac{R_{\max}}{2MC_0},\,
  \frac1{2K_0}
 \right\}.
 \label{eq:v5v15-rstar}
\end{align}
For every \(\|x\|_X\le r_*\), set
\begin{equation}
 R(x)=2MC_0\|x\|_X.
 \label{eq:v5v15-Rx}
\end{equation}
Then the V14 hypotheses hold on the \(Z\)-ball of radius \(R(x)\) with
\(q\le1/2\).  There is a unique solution of
\(\mathcal C(x,z)=0\) in that ball, and
\begin{equation}
 \|Z(x)\|_Z\le R(x).
 \label{eq:v5v15-Z-bound}
\end{equation}
All constants are independent of a Fourier ultraviolet cutoff when the
bounds \eqref{eq:v5v15-Nzz}--\eqref{eq:v5v15-source} are.
\end{theorem}

\begin{proof}
For \(\|x\|\le1\), equation \eqref{eq:v5v15-source} gives
\[
 \|\mathcal S(x)\|_Y\le C_0\|x\|_X.
\]
Hence the V14 residual obeys
\[
 \eta(x)=\|L^{-1}\mathcal C(x,0)\|_Z
 \le MC_0\|x\|_X=\frac12R(x).
\]
The \(z\)-derivative of the nonlinear part satisfies on the radius
\(R(x)\)
\begin{align*}
 \|L^{-1}D_z(\mathcal C-Lz)\|
 &\le
 M\bigl(2C_{zz}R(x)+C_{xz}\|x\|\bigr)\\
 &=
 M\bigl(4MC_{zz}C_0+C_{xz}\bigr)\|x\|
 =K_0\|x\|\le\frac12.
\end{align*}
Thus \(q\le1/2\) and
\[
 \eta(x)\le\frac12R(x)\le(1-q)R(x).
\]
The definition of \(r_*\) also ensures \(R(x)\le R_{\max}\).
Apply Theorem~\ref{thm:v5v14-constraint-solver}.  Cutoff independence
follows when every input constant is cutoff independent.
\end{proof}

\begin{corollary}[Classical source examples]
\label{cor:v5v15-classical-sources}
For \(s>3/2\), the bosonic Yang--Mills and Higgs stress terms are in
\(H^s\) when gauge potentials and Higgs fields lie in \(H^{s+1}\), on a
small \(H^{s+1}\) ball.  Their fixed-arity source bounds can therefore be
included in \(C_{\rm s}\).
\end{corollary}

\begin{proof}
A curvature has the schematic form
\[
 F=\nabla A+A^2.
\]
For \(A\in H^{s+1}\), both terms lie in \(H^s\) by
Theorem~\ref{thm:v5v15-Sobolev-algebra}.  The stress is quadratic in \(F\)
and remains in \(H^s\).  The Higgs covariant derivative has the form
\(\nabla\Phi+A\Phi\), and the potential is a fixed polynomial; the same
algebra estimate applies.
\end{proof}

Fermionic quantum stress is not covered by this classical Banach
argument.  It must be defined through renormalized determinant responses
and the still-open finite marginal packet identified in the fourth-series
ledger.

\subsection{The equal-time quantum regularity test}

The preceding theorem assumes positive Sobolev regularity.  Consider a
free massless scalar on \(\mathbb T^3\), with zero mode removed.  Its
equal-time Fourier covariance is
\begin{equation}
 \mathbb E|\widehat\Phi(k)|^2=\frac1{2|k|}.
 \label{eq:v5v15-free-covariance}
\end{equation}

\begin{proposition}[Gaussian Sobolev threshold]
\label{prop:v5v15-Gaussian-threshold}
For the Gaussian field \eqref{eq:v5v15-free-covariance},
\begin{equation}
 \mathbb E\|\Phi\|_{H^r}^2
 =
 \frac12\sum_{k\ne0}
 \frac{(1+|k|^2)^r}{|k|}
 \label{eq:v5v15-Hr-expectation}
\end{equation}
is finite if and only if
\begin{equation}
 r<-1.
 \label{eq:v5v15-threshold}
\end{equation}
Moreover, the field belongs to \(H^r\) almost surely for \(r<-1\) and
does not belong to \(H^r\) almost surely for \(r\ge-1\).
\end{proposition}

\begin{proof}
The number of lattice modes in a shell of radius \(K\) and unit width is
of order \(K^2\).  The shell contribution to
\eqref{eq:v5v15-Hr-expectation} is therefore comparable to
\(K^{2r+1}\).  The sum converges exactly when \(2r+1<-1\), which is
\eqref{eq:v5v15-threshold}.  In the convergent case, finite expectation
implies finite \(H^r\) norm almost surely.  In the divergent case, group
independent Gaussian modes into dyadic shells.  Their weighted squared
norms have expectations bounded below by the divergent shell series and
uniformly controlled relative variances after shell aggregation.
Independence and the second Borel--Cantelli lemma then force infinitely
many nonnegligible shell contributions, so the \(H^r\) norm diverges
almost surely.
\end{proof}

\begin{corollary}[The classical constraint ball has Gaussian probability
zero]
\label{cor:v5v15-zero-probability}
For \(s>3/2\), the free equal-time scalar Gaussian assigns probability zero
to \(H^{s+1}\), and hence to every small \(H^{s+1}\) ball used in
Corollary~\ref{cor:v5v15-classical-sources}.  The same ultraviolet
regularity mismatch occurs for free gauge potentials, up to polarization
and gauge projection.
\end{corollary}

\begin{proof}
Here \(s+1>5/2>-1\), so
Proposition~\ref{prop:v5v15-Gaussian-threshold} applies.  Removing finitely
many polarizations or imposing a transverse projection does not change the
large-\(|k|\) covariance power or the threshold.
\end{proof}

\begin{firewall}[Cutoff-uniform constants do not give a typical-field
domain]
\label{fw:v5v15-typical-domain}
The radius \(r_*\) in
Theorem~\ref{thm:v5v15-explicit-radius} is independent of the ultraviolet
cutoff, but typical quantum configurations leave its positive-Sobolev
domain as the cutoff is removed.  This is a support obstruction, not a
loss in the numerical value of \(r_*\).
\end{firewall}

\subsection{Renormalized constraint certificate}

A distributional quantum constraint must replace products such as
\(F^2\), \((\nabla\Phi)^2\), and nonlinear curvature terms by
renormalized composite fields.  The counterterms must agree with the
essential local tower of V5--V7.

\begin{definition}[Renormalized quantum constraint certificate]
\label{def:v5v15-RQCC}
An RQCC consists of:
\begin{enumerate}
 \item a negative-regularity support space carrying the cutoff field
 measures;
 \item renormalized stress and curvature composites converging in a source
 space \(Y_{\rm ren}\);
 \item a constraint inverse mapping \(Y_{\rm ren}\) to an auxiliary space
 with cutoff-uniform bounds;
 \item locally Lipschitz or controlled rough dependence of the nonlinear
 constraint solution on enhanced field data;
 \item counterterms compatible with the V5 essential quotient and V7
 resummation;
 \item the V14 Jacobian, all-root, Ward, and shell-compatibility
 properties in the renormalized setting.
\end{enumerate}
\end{definition}

\begin{theorem}[Classical versus quantum constraint status]
\label{thm:v5v15-status}
On a fixed torus, Theorem~\ref{thm:v5v15-explicit-radius} gives a
cutoff-uniform classical constraint solution for sufficiently regular
small data.  It does not give an RQCC or a quantum SM--Einstein constraint
solution, because the free-field support already fails the required
Sobolev regularity.
\end{theorem}

\begin{proof}
The positive classical statement is
Theorem~\ref{thm:v5v15-explicit-radius}.  The support failure is
Corollary~\ref{cor:v5v15-zero-probability}.  None of the renormalized
objects in Definition~\ref{def:v5v15-RQCC} is constructed by a classical
Sobolev product estimate.
\end{proof}

V15 therefore prevents a circular use of smooth constraint data inside a
quantum continuum proof.  The classical nonlinear reduction is
ultraviolet-uniform on its domain, but that domain is negligible for the
limiting field measure.  The next attack must construct and estimate the
renormalized source before applying the elliptic inverse.

V16 will begin with the Gaussian model.  It will compute the regularity of
a Wick square and its inverse-Laplacian image on \(\mathbb T^3\), identify
the exact improvement supplied by the constraint Green operator, and test
whether the solved conformal response is a function or still a
distribution.


\clearpage
\section{V16: the equal-time Wick-source obstruction}
\label{sec:v5v16}

V15 showed that typical quantum fields do not lie in the classical
Sobolev algebra used by the nonlinear constraint solver.  The simplest
possible repair would be to Wick order the quadratic source and then use
the two-derivative gain of the constraint Green operator.  This note proves
that the repair fails already for a free scalar at equal time.  The
nonzero Fourier modes of the Wick square have variance growing linearly
with the ultraviolet cutoff, and an inverse Laplacian on the output mode
does not affect the internal high--high to low loop.  Smooth time smearing
does cure the divergence, so the missing theorem is an instantaneous
boundary-value or complete-source cancellation theorem.

\subsection{The cutoff free field}

On \(\mathbb T^3\), remove the zero mode and let
\(\Phi_\Lambda\) be the centered real Gaussian field with Fourier
covariance
\begin{equation}
 \mathbb E\bigl[
  \widehat\Phi_\Lambda(p)
  \widehat\Phi_\Lambda(q)
 \bigr]
 =
 \frac{\mathbf 1_{\{0<|p|\le\Lambda\}}}{2|p|}
 \,\mathbf 1_{\{p+q=0\}}.
 \label{eq:v5v16-field-covariance}
\end{equation}
Define the centered square
\[
 Q_\Lambda(x)
 =
 \Phi_\Lambda(x)^2-\mathbb E\Phi_\Lambda(x)^2.
\]
For \(k\ne0\), centering does not change its Fourier coefficient:
\begin{equation}
 \widehat Q_\Lambda(k)
 =
 \sum_{\substack{p\ne0,k\\
       |p|\le\Lambda,\ |k-p|\le\Lambda}}
 \widehat\Phi_\Lambda(p)
 \widehat\Phi_\Lambda(k-p).
 \label{eq:v5v16-Wick-mode}
\end{equation}

\begin{lemma}[Exact Wick variance]
\label{lem:v5v16-Wick-variance}
For fixed \(k\ne0\),
\begin{equation}
 \mathbb E|\widehat Q_\Lambda(k)|^2
 =
 \frac12
 \sum_{\substack{p\ne0,k\\
       |p|\le\Lambda,\ |k-p|\le\Lambda}}
 \frac1{|p|\,|k-p|},
 \label{eq:v5v16-exact-variance}
\end{equation}
up to the harmless convention for choosing one representative of the real
Fourier pairs.
\end{lemma}

\begin{proof}
Expand the square of \eqref{eq:v5v16-Wick-mode}.  Wick's theorem has two
cross contractions, each imposing the same momentum sum and contributing
\((2|p|)^{-1}(2|k-p|)^{-1}\).  Their sum gives the factor \(1/2\).
The contraction within one copy is absent for \(k\ne0\).
\end{proof}

\subsection{Linear ultraviolet divergence at every fixed output mode}

\begin{theorem}[Fixed-mode Wick-square divergence]
\label{thm:v5v16-fixed-mode-divergence}
For every fixed \(k\ne0\), there are constants
\(0<c_k<C_k<\infty\) and \(\Lambda_k\) such that
\begin{equation}
 c_k\Lambda
 \le
 \mathbb E|\widehat Q_\Lambda(k)|^2
 \le
 C_k\Lambda
 \qquad(\Lambda\ge\Lambda_k).
 \label{eq:v5v16-linear-divergence}
\end{equation}
Consequently \(Q_\Lambda\) does not converge in
\(L^2(\Omega;\mathcal D'(\mathbb T^3))\).
\end{theorem}

\begin{proof}
For the lower bound, restrict the sum in
\eqref{eq:v5v16-exact-variance} to
\[
 2|k|\le|p|\le\Lambda/2.
\]
Then
\[
 \frac12|p|\le|k-p|\le\frac32|p|,
\]
so every summand is comparable to \(|p|^{-2}\).  A lattice shell of radius
\(r\) and unit width contains order \(r^2\) points.  Each such shell
therefore contributes an amount bounded below by a positive constant, and
there are order \(\Lambda\) shells.

For the upper bound, split into
\(|p|\le2|k|\) and \(|p|>2|k|\).  The first set has finitely many terms
depending on \(k\).  On the second,
\(|k-p|\ge|p|/2\), so the summand is at most \(2|p|^{-2}\).
Summing the order-one contribution of each shell up to \(\Lambda+|k|\)
gives the upper bound.

If \(Q_\Lambda\) converged in
\(L^2(\Omega;\mathcal D')\), its pairing with the smooth Fourier test
\(e^{-ik\cdot x}\) would be Cauchy in \(L^2(\Omega)\), hence bounded there.
Equation \eqref{eq:v5v16-linear-divergence} contradicts boundedness.
\end{proof}

\begin{corollary}[Identity Wick subtraction cannot remove the divergence]
\label{cor:v5v16-identity-no-cure}
Changing the subtracted scalar
\(\mathbb E\Phi_\Lambda(x)^2\) affects only the zero Fourier mode and
cannot change \eqref{eq:v5v16-linear-divergence} at \(k\ne0\).
\end{corollary}

\begin{proof}
A spatially constant counterterm has Fourier support only at \(k=0\).
\end{proof}

This is stronger than the statement that the unrenormalized expectation
diverges.  The fluctuation of every fixed nonzero mode diverges.

\subsection{The constraint inverse does not cure an internal loop}

Let \(G\) be a translation-invariant output operator with Fourier symbol
\(g(k)\), and set
\[
 Z_\Lambda=GQ_\Lambda.
\]

\begin{theorem}[Output smoothing leaves the fixed-mode divergence]
\label{thm:v5v16-output-no-cure}
If \(g(k_0)\ne0\) for some fixed \(k_0\ne0\), then
\begin{equation}
 \mathbb E|\widehat Z_\Lambda(k_0)|^2
 =
 |g(k_0)|^2
 \mathbb E|\widehat Q_\Lambda(k_0)|^2
 \asymp_{k_0}\Lambda.
 \label{eq:v5v16-output-divergence}
\end{equation}
In particular, for the mean-zero constraint inverse
\(G=(-\Delta)^{-1}\),
\begin{equation}
 \mathbb E|\widehat Z_\Lambda(k_0)|^2
 \asymp_{k_0}|k_0|^{-4}\Lambda.
 \label{eq:v5v16-Poisson-divergence}
\end{equation}
No amount of regularity gain acting only on the output frequency removes
the high--high to fixed-low divergence.
\end{theorem}

\begin{proof}
Fourier multiplication gives
\(\widehat Z_\Lambda(k)=g(k)\widehat Q_\Lambda(k)\).
Apply Theorem~\ref{thm:v5v16-fixed-mode-divergence}.  For the inverse
Laplacian, \(g(k)=|k|^{-2}\).
\end{proof}

\begin{firewall}[Elliptic gain and composite renormalization are different]
\label{fw:v5v16-elliptic-vs-loop}
The V14 estimate
\(L^{-1}:H^s\to H^{s+2}\) begins after a source has been defined in
\(H^s\).  Theorem~\ref{thm:v5v16-output-no-cure} shows that the inverse
cannot define a source whose fixed low modes already have divergent
variance.  Quoting two derivatives of gain before renormalizing the
internal contraction is circular.
\end{firewall}

\subsection{Smooth time smearing supplies the missing internal decay}

The obstruction concerns an instantaneous field.  Let \(\Phi(t,x)\) be
the free massless field and smear its normal-ordered square with
\(f\in C^\infty(\mathbb T^3)\) and a Schwartz time function
\(a\in\mathcal S(\mathbb R)\).  The two-particle creation amplitude from
the vacuum has the form
\begin{equation}
 \mathcal A(p,q)
 =
 \frac{\widehat f(p+q)\,
       \widehat a(|p|+|q|)}
      {2\sqrt{|p|\,|q|}},
 \label{eq:v5v16-time-smeared-amplitude}
\end{equation}
up to normalization.

\begin{theorem}[Spacetime smearing makes the free Wick square finite]
\label{thm:v5v16-time-smearing}
For smooth \(f\) and Schwartz \(a\),
\begin{equation}
 \sum_{p,q\ne0}|\mathcal A(p,q)|^2<\infty.
 \label{eq:v5v16-time-smearing-finite}
\end{equation}
Hence the spacetime-smeared Wick square applied to the vacuum has finite
norm.
\end{theorem}

\begin{proof}
For every \(N\), Schwartz decay gives
\[
 |\widehat a(|p|+|q|)|
 \le C_N(1+|p|+|q|)^{-N}.
\]
The smooth spatial function also has
\[
 |\widehat f(p+q)|
 \le C_N(1+|p+q|)^{-N}.
\]
Choose \(N\) larger than the lattice dimension.  The time factor controls
the dangerous region \(p\approx-q\), where the spatial factor alone does
not decay, and the spatial factor controls large total momentum.  After
discarding the favorable denominators \(|p|^{-1}|q|^{-1}\) outside a
finite set, the resulting double lattice sum is dominated by a sufficiently
high negative power of \(1+|p|+|q|\), and therefore converges.
\end{proof}

\begin{corollary}[The instantaneous limit is a boundary problem]
\label{cor:v5v16-instantaneous-boundary}
The equal-time constraint source cannot be obtained merely by deleting
the time smearing in
Theorem~\ref{thm:v5v16-time-smearing}.  Such a limit must use a separate
energy bound, cancellation among the complete constraint terms, or a
renormalized time-zero theorem.
\end{corollary}

\begin{proof}
Deleting the time smearing replaces
\(\widehat a(|p|+|q|)\) by a nondecaying factor.  The fixed-output sector
\(p+q=k_0\) then reduces to the divergent sum of
Theorem~\ref{thm:v5v16-fixed-mode-divergence}.
\end{proof}

\subsection{Why the complete stress tensor must remain assembled}

The scalar square is a diagnostic packet, not the full SM stress tensor.
Derivative terms, improvement terms, gauge constraints, fermionic
determinant responses, and local counterterms may cancel contact
singularities only when their coefficients and Ward identities are kept
together.

\begin{proposition}[Termwise absolute estimates cannot prove a cancellation]
\label{prop:v5v16-termwise}
Suppose a complete source has
\[
 \mathcal S_\Lambda
 =\sum_{\alpha=1}^J c_\alpha Q_{\alpha,\Lambda}.
\]
If one component has variance of order \(\Lambda\), an estimate using
\[
 \|\mathcal S_\Lambda\|_{L^2(\Omega)}
 \le\sum_\alpha|c_\alpha|
       \|Q_{\alpha,\Lambda}\|_{L^2(\Omega)}
\]
cannot establish a uniform bound by cancellation.  A successful proof
must evaluate the covariance matrix of the assembled vector and show that
its divergent part annihilates the coefficient vector \(c\).
\end{proposition}

\begin{proof}
The triangle inequality replaces every cross covariance by its absolute
value and therefore discards its sign.  In contrast,
\[
 \mathbb E|\mathcal S_\Lambda|^2
 =c^*\Sigma_\Lambda c
\]
contains the signed off-diagonal covariances.  Cancellation of a divergent
matrix component is exactly the condition that its leading quadratic form
vanish on \(c\).
\end{proof}

This is the quantum-constraint version of the companion-programme rule to
assemble the complete local source before taking block norms.

\subsection{An instantaneous composite-source certificate}

\begin{definition}[Time-zero composite constraint certificate]
\label{def:v5v16-TZCC}
A TZCC consists of:
\begin{enumerate}
 \item a spacetime-renormalized, Ward-complete stress and constraint
 insertion;
 \item a proof that its positive-time boundary value exists in a source
 space \(Y_{\rm ren}\), or a proof that the constraint can be imposed
 without taking an equal-time composite;
 \item cancellation or subtraction of every fixed-output internal
 divergence, including the covariance divergence
 \eqref{eq:v5v16-linear-divergence};
 \item cutoff-uniform bounds after the complete source is assembled;
 \item compatibility of the source counterterms with the V5 essential
 local quotient and V7 resummation;
 \item application of the V14 inverse only after items 1--4.
\end{enumerate}
\end{definition}

\begin{theorem}[Gaussian model boundary]
\label{thm:v5v16-status}
The centered equal-time square of a free four-dimensional scalar restricted
to a spatial slice does not satisfy a TZCC, and applying the spatial
constraint inverse does not repair it.  Smooth spacetime smearing is
well-defined, so the failure is localized to the time-zero restriction.
\end{theorem}

\begin{proof}
The fixed-mode failure is
Theorem~\ref{thm:v5v16-fixed-mode-divergence}; persistence after the
constraint inverse is
Theorem~\ref{thm:v5v16-output-no-cure}.  The positive spacetime statement
is Theorem~\ref{thm:v5v16-time-smearing}.
\end{proof}

V16 rules out a tempting shortcut in the nonlinear constraint programme:
Wick centering followed by an inverse Laplacian is not enough.  The next
step must retain the full constraint source and exploit its Ward
structure, or formulate the physical theory without an instantaneous
composite insertion.

V17 will derive the linearized Ward identities for the complete improved
scalar stress tensor and compute its two-point transverse projector.  The
goal is to determine which divergent contact pieces are killed by
constraint contraction and which require independent local gravitational
counterterms.


\clearpage
\section{V17: stress-tensor Ward decomposition and the surviving spin-two
source}
\label{sec:v5v17}

V16 used a scalar square as a diagnostic and warned that only the complete
constraint source can exhibit Ward cancellation.  This note performs the
linearized tensor decomposition for the improved free-scalar stress
tensor.  Conservation removes longitudinal tensor structures, and the
conformal improvement removes the separated-point spin-zero trace
structure.  The transverse-traceless spin-two structure survives, is
independent of the improvement coefficient, and has a nonzero two-particle
spectral density.  Its ultraviolet contact terms require local
gravitational counterterms; they are not killed by the constraint
projection.

\subsection{The improved scalar stress tensor}

On flat Euclidean \(\mathbb R^4\), let \(\phi\) be a massless free scalar.
For a real improvement parameter \(\xi\), define
\begin{equation}
 T_{\mu\nu}^{(\xi)}
 =
 \partial_\mu\phi\,\partial_\nu\phi
 -\frac12\delta_{\mu\nu}(\partial\phi)^2
 +\xi\bigl(\delta_{\mu\nu}\Box-\partial_\mu\partial_\nu\bigr)\phi^2.
 \label{eq:v5v17-improved-stress}
\end{equation}

\begin{lemma}[On-shell Ward identities]
\label{lem:v5v17-Ward}
Modulo contact terms and the free equation \(\Box\phi=0\),
\begin{equation}
 \partial^\mu T_{\mu\nu}^{(\xi)}=0.
 \label{eq:v5v17-conservation}
\end{equation}
For \(\xi=1/6\), the separated-point trace also vanishes:
\begin{equation}
 \delta^{\mu\nu}T_{\mu\nu}^{(1/6)}=0.
 \label{eq:v5v17-trace}
\end{equation}
\end{lemma}

\begin{proof}
The divergence of the first two terms is
\[
 (\Box\phi)\partial_\nu\phi.
\]
The divergence of the improvement vanishes identically because commuting
flat derivatives gives
\[
 \partial_\nu\Box\phi^2-\Box\partial_\nu\phi^2=0.
\]
This proves \eqref{eq:v5v17-conservation} on the free equation.  In four
dimensions, the trace is
\[
 -(\partial\phi)^2+3\xi\Box\phi^2.
\]
Using
\(\Box\phi^2=2(\partial\phi)^2\) on shell gives
\((-1+6\xi)(\partial\phi)^2\), which vanishes at \(\xi=1/6\).
Coincident insertions can add local contact terms and are excluded from
the separated-point identity.
\end{proof}

\subsection{Transverse tensor projectors}

For \(p\ne0\), set
\begin{align}
 \pi_{\mu\nu}(p)
 &=\delta_{\mu\nu}-\frac{p_\mu p_\nu}{p^2},
 \label{eq:v5v17-pi}\\
 \Pi^{(2)}_{\mu\nu,\rho\sigma}
 &=
 \frac12\bigl(
  \pi_{\mu\rho}\pi_{\nu\sigma}
 +\pi_{\mu\sigma}\pi_{\nu\rho}
 \bigr)
 -\frac13\pi_{\mu\nu}\pi_{\rho\sigma},
 \label{eq:v5v17-Pi2}\\
 \Pi^{(0)}_{\mu\nu,\rho\sigma}
 &=
 \frac13\pi_{\mu\nu}\pi_{\rho\sigma}.
 \label{eq:v5v17-Pi0}
\end{align}

\begin{lemma}[Projector algebra]
\label{lem:v5v17-projectors}
The operators \(\Pi^{(2)}\) and \(\Pi^{(0)}\) are orthogonal
self-adjoint projections on the transverse symmetric tensors.  Their
ranges are respectively the transverse-traceless tensors and the
one-dimensional transverse trace sector.
\end{lemma}

\begin{proof}
The matrix \(\pi\) is the orthogonal projection onto \(p^\perp\), a
three-dimensional space.  A symmetric tensor on that space splits
orthogonally into its trace-free part and its scalar multiple of \(\pi\).
Equations \eqref{eq:v5v17-Pi2}--\eqref{eq:v5v17-Pi0} are exactly those two
orthogonal projections.
\end{proof}

Let the renormalized connected stress two-point function be
\[
 \mathcal W_{\mu\nu,\rho\sigma}(p)
 =
 \langle T_{\mu\nu}(p)T_{\rho\sigma}(-p)\rangle_{\rm conn}.
\]

\begin{theorem}[Ward-complete two-point decomposition]
\label{thm:v5v17-decomposition}
At nonzero momentum and away from local polynomial contact terms,
Euclidean invariance, index symmetry, and
\eqref{eq:v5v17-conservation} imply
\begin{equation}
 \mathcal W(p)
 =
 A(p^2)\Pi^{(2)}(p)
 +B(p^2)\Pi^{(0)}(p).
 \label{eq:v5v17-W-decomposition}
\end{equation}
For the conformally improved massless scalar, \(B=0\) at separated points.
\end{theorem}

\begin{proof}
Conservation places every tensor index pair in the transverse symmetric
space.  The stabilizer of \(p\) acts as the orthogonal group on
\(p^\perp\).  Its symmetric rank-two representation splits into the
irreducible trace-free part and the scalar trace part.  An invariant
two-point operator acts by a scalar on each irreducible component, giving
\eqref{eq:v5v17-W-decomposition}.  The trace identity
\eqref{eq:v5v17-trace} annihilates the scalar component at separated
points, so \(B=0\).
\end{proof}

Local contact terms can be transverse polynomial tensors.  At fourth
momentum degree they include \(p^4\Pi^{(2)}\) and
\(p^4\Pi^{(0)}\), which are polynomial despite the notation with
\(\pi\).  They correspond to second variations of curvature-squared local
actions.

\subsection{Improvement cannot change the spin-two component}

In momentum space, the difference of two improvements is
\begin{equation}
 T_{\mu\nu}^{(\xi)}(p)-T_{\mu\nu}^{(\xi')}(p)
 =
 -(\xi-\xi')p^2\pi_{\mu\nu}(p)\,
 \widehat{\phi^2}(p).
 \label{eq:v5v17-improvement-difference}
\end{equation}

\begin{proposition}[Spin-two improvement invariance]
\label{prop:v5v17-spin2-invariant}
For every \(\xi,\xi'\),
\begin{equation}
 \Pi^{(2)}
 \bigl(T^{(\xi)}-T^{(\xi')}\bigr)=0.
 \label{eq:v5v17-improvement-killed}
\end{equation}
Consequently the nonlocal coefficient \(A(p^2)\) in
\eqref{eq:v5v17-W-decomposition} is independent of \(\xi\).
\end{proposition}

\begin{proof}
The tensor in \eqref{eq:v5v17-improvement-difference} is proportional to
\(\pi_{\mu\nu}\), while
\(\Pi^{(2)}\) is trace free on the transverse space.  Hence the projection
vanishes.  Applying the projection to both stress insertions shows that
their spin-two two-point functions agree.
\end{proof}

Thus improvement may move the transverse scalar and contact owners but
cannot tune away the physical spin-two source.

\subsection{Nonzero spin-two spectral weight}

Let \(e_{\mu\nu}\) be a nonzero transverse-traceless polarization:
\[
 e_{\mu\nu}p^\mu=0,
 \qquad
 e_\mu{}^\mu=0.
\]
The vacuum-to-two-particle vertex of the stress tensor satisfies, up to a
nonzero normalization,
\begin{equation}
 \mathcal V_e(q,p-q)
 =
 e^{\mu\nu}q_\mu(p-q)_\nu
 =
 -e^{\mu\nu}q_\mu q_\nu.
 \label{eq:v5v17-TT-vertex}
\end{equation}
The trace and improvement terms vanish after contraction with \(e\).

\begin{theorem}[The scalar has a genuine spin-two stress channel]
\label{thm:v5v17-nonzero-spin2}
For every nonzero external momentum for which the two-particle channel is
kinematically available after continuation, the spin-two spectral density
of the free-scalar stress tensor is nonnegative and not identically zero.
Equivalently, \(A\) in
\eqref{eq:v5v17-W-decomposition} is a genuine nonlocal coefficient, not a
pure contact term.
\end{theorem}

\begin{proof}
The two-particle contribution to the spectral density is the phase-space
integral of
\[
 |\mathcal V_e(q,p-q)|^2
\]
against a positive measure.  It is nonnegative.  The quadratic polynomial
\(e^{\mu\nu}q_\mu q_\nu\) is not identically zero when \(e\ne0\), so it is
nonzero on an open subset of the two-particle phase space.  The positive
integral is therefore nonzero.  A pure contact term has no nontrivial
spectral density, proving the last assertion.
\end{proof}

\begin{corollary}[Constraint projection cannot remove all quantum stress]
\label{cor:v5v17-constraint-not-zero}
A constraint or gauge projection that removes only longitudinal and
transverse-scalar sectors leaves the nonzero spin-two stress channel of
Theorem~\ref{thm:v5v17-nonzero-spin2}.
\end{corollary}

\begin{proof}
The surviving channel lies in
\(\operatorname{Ran}\Pi^{(2)}\), which is orthogonal to the removed
sectors by Lemma~\ref{lem:v5v17-projectors}.
\end{proof}

\subsection{Ultraviolet degree and local owners}

The cutoff one-loop spin-two integral has the schematic form
\begin{equation}
 \int_{|q|\le\Lambda}
 \frac{|e^{\mu\nu}q_\mu q_\nu|^2}
      {q^2(p-q)^2}\,d^4q.
 \label{eq:v5v17-spin2-loop}
\end{equation}
At large \(q\), its numerator and denominator both have degree four.
Power divergences and the logarithmic fourth-order momentum term are local
after the Ward-complete Taylor subtraction.  The remaining nonlocal
coefficient has the scale form \(p^4\log(p^2/\mu^2)\), up to a constant
and scheme-dependent local polynomial.

\begin{proposition}[Local versus nonlocal stress response]
\label{prop:v5v17-local-nonlocal}
A Ward-complete renormalization of
\eqref{eq:v5v17-spin2-loop} requires local metric counterterms through
curvature-squared order.  Those counterterms can remove the polynomial
cutoff dependence but cannot remove the nonzero spin-two spectral density.
\end{proposition}

\begin{proof}
Taylor subtraction through superficial degree four produces a local
polynomial in the external momentum.  Diffeomorphism Ward identities
assemble that polynomial into variations of the cosmological,
Einstein, and curvature-squared owners.  Subtracting a polynomial changes
only contact terms.  The spectral density from
Theorem~\ref{thm:v5v17-nonzero-spin2} is supported away from coincidence
and is unchanged by such a subtraction.
\end{proof}

This is compatible with the inherited V90 matter heat-counterterm card.
V17 does not repeat its determinant construction; it identifies why those
local owners remain necessary even after the physical constraint quotient.

\subsection{Why the equal-time limit remains open}

The Ward identity is a spacetime distributional identity.  At the
two-point level it reads
\[
 p^\mu\mathcal W_{\mu\nu,\rho\sigma}(p)
 =\text{local contact polynomial}.
\]
It does not provide decay in the loop momentum \(q\) at fixed external
spatial momentum.

\begin{proposition}[Spacetime transversality is not a time-zero bound]
\label{prop:v5v17-Ward-not-timezero}
The decomposition
\eqref{eq:v5v17-W-decomposition} does not imply existence of the
equal-time random distribution \(T_{\mu\nu}(0,\cdot)\).  Such a boundary
value requires an independent estimate in the energy variable \(p_0\).
\end{proposition}

\begin{proof}
Conservation is an algebraic relation among tensor components at each full
four-momentum \(p=(p_0,\mathbf p)\).  Restriction to time zero integrates
the Fourier transform over all \(p_0\) without a decaying time test.
Neither multiplication by the bounded projectors at fixed nonzero \(p\)
nor the identity \(p^\mu\mathcal W_{\mu\nu,\rho\sigma}=0\) supplies
integrability in \(p_0\).  The V16 two-particle calculation explicitly
shows how removing time smearing restores an internal divergence.
\end{proof}

\begin{firewall}[Ward cancellation must precede restriction]
\label{fw:v5v17-Ward-order}
A proposed TZCC must first renormalize and assemble the complete spacetime
stress tensor, including its local contacts, then prove a boundary value
for the constrained combination.  Restricting individual products to a
time slice and imposing the Ward identity afterward is not equivalent.
\end{firewall}

\subsection{Stress-to-constraint certificate}

\begin{definition}[Ward-complete stress boundary certificate]
\label{def:v5v17-WSBC}
A WSBC consists of:
\begin{enumerate}
 \item a spacetime stress tensor with exact diffeomorphism Ward identities
 and all local contact owners;
 \item the transverse decomposition
 \eqref{eq:v5v17-W-decomposition} through the required marked insertions;
 \item uniform bounds for the nonlocal spin-two and scalar spectral
 coefficients;
 \item a time-zero boundary theorem for the exact combination entering the
 Hamiltonian and momentum constraints;
 \item compatibility of the local contacts with the V5 curvature quotient
 and V7 all-order shell map;
 \item insertion of the resulting source into the V14 constraint solver.
\end{enumerate}
\end{definition}

\begin{theorem}[Current free-scalar stress status]
\label{thm:v5v17-status}
For the free scalar, the separated-point Ward decomposition and the
nonzero spin-two channel are explicit.  They prove neither item 4 of the
WSBC nor the nonlinear SM--Einstein constraint measure.  Local
curvature-squared counterterms remain necessary, and the physical
spin-two matter source remains nontrivial after constraint reduction.
\end{theorem}

\begin{proof}
The decomposition is
Theorem~\ref{thm:v5v17-decomposition}, the spin-two statement is
Theorem~\ref{thm:v5v17-nonzero-spin2}, and the local allocation is
Proposition~\ref{prop:v5v17-local-nonlocal}.  The missing time-zero
implication is
Proposition~\ref{prop:v5v17-Ward-not-timezero}.
\end{proof}

V17 shows exactly what the Ward structure saves and what it does not.
Longitudinal and conformal trace sectors can be removed at separated
points, but the TT stress response and its local gravitational
counterterms survive.  The instantaneous constraint source remains the
bottleneck.

V18 will replace the sharp time slice by a positive-time slab of width
\(\varepsilon\).  It will compute the dependence of the free two-particle
source norm on \(\varepsilon\), determine whether the V14 constraint
inverse gives a uniform boundary limit, and identify the precise
boundary counterterm or energy estimate still required.


\clearpage
\section{V18: slab-to-boundary scaling of the quadratic source}
\label{sec:v5v18}

V16 proved that an equal-time scalar square does not exist as an
\(L^2\) random distribution, while smooth spacetime smearing is finite.
This note measures the transition between the two statements.  An
\(L^1\)-normalized time mollifier of width \(\varepsilon\), which genuinely
approximates a time slice, gives a fixed-spatial-mode two-particle norm of
order \(\varepsilon^{-1}\).  The spatial constraint inverse does not alter
that rate.  An \(L^2\)-normalized slab remains bounded but tends to zero as
a time distribution, so it is not a boundary value.  The result identifies
the exact positive contact matrix that a complete constraint source must
annihilate.

\subsection{Approximate time slices}

Let \(a\in\mathcal S(\mathbb R)\) be real with
\begin{equation}
 \int_{\mathbb R}a(t)\,dt=1.
 \label{eq:v5v18-mollifier-normalization}
\end{equation}
Define the approximate identity
\begin{equation}
 a_\varepsilon(t)=\varepsilon^{-1}a(t/\varepsilon),
 \qquad
 \widehat a_\varepsilon(\omega)=\widehat a(\varepsilon\omega).
 \label{eq:v5v18-approximate-identity}
\end{equation}
Fix a nonzero spatial Fourier mode \(k\).  For the normal-ordered free
scalar square, the squared vacuum norm of the \(a_\varepsilon\)-smeared
\(k\)-mode is, up to a fixed positive normalization,
\begin{equation}
 V_\varepsilon(k)
 =
 \sum_{\substack{p\ne0,k}}
 \frac{
 |\widehat a(\varepsilon(|p|+|k-p|))|^2}
 {|p|\,|k-p|}.
 \label{eq:v5v18-slab-variance}
\end{equation}

\begin{theorem}[Sharp slab-to-boundary divergence]
\label{thm:v5v18-sharp-slab}
For every fixed \(k\ne0\), there are
\(0<c_{a,k}<C_{a,k}<\infty\) and \(\varepsilon_{a,k}>0\) such that
\begin{equation}
 c_{a,k}\varepsilon^{-1}
 \le V_\varepsilon(k)
 \le C_{a,k}\varepsilon^{-1}
 \qquad(0<\varepsilon<\varepsilon_{a,k}).
 \label{eq:v5v18-epsilon-inverse}
\end{equation}
\end{theorem}

\begin{proof}
Continuity and
\(\widehat a(0)=1\) give a number \(\delta>0\) for which
\[
 |\widehat a(\omega)|\ge\frac12
 \qquad(|\omega|\le\delta).
\]
Restrict the sum to
\[
 2|k|\le|p|\le\frac{\delta}{4\varepsilon}.
\]
For sufficiently small \(\varepsilon\),
\[
 |p|+|k-p|\le\frac52|p|
 \le\frac{5\delta}{8\varepsilon}<\frac{\delta}{\varepsilon},
\]
after decreasing the upper constant if necessary.  Also
\(|k-p|\asymp|p|\).  Each unit radial shell then contributes a fixed
positive amount, and there are order \(\varepsilon^{-1}\) shells.  This
proves the lower bound.

For the upper bound, the finitely many modes with \(|p|\le2|k|\) are
harmless.  On the remainder,
\[
 |k-p|\asymp|p|,
 \qquad
 |p|+|k-p|\ge c|p|.
\]
A unit shell of radius \(r\) contains \(O(r^2)\) lattice points and each
denominator is of order \(r^2\).  Thus its contribution is at most
\[
 C_N(1+\varepsilon r)^{-N}
\]
for every \(N\), by Schwartz decay of \(\widehat a\).  Summing over
\(r\ge1\) and comparing with an integral gives
\[
 \sum_{r\ge1}(1+\varepsilon r)^{-N}
 \le C_N'\varepsilon^{-1}
\]
when \(N>1\).
\end{proof}

\begin{corollary}[No \(L^2\) time-zero boundary vector]
\label{cor:v5v18-no-boundary-vector}
The vectors obtained by applying the smeared square to the vacuum have no
bounded \(L^2\) limit as \(a_\varepsilon\to\delta_0\), even after fixing a
smooth nonzero spatial Fourier test.
\end{corollary}

\begin{proof}
Their squared norms are proportional to
\(V_\varepsilon(k)\), which diverges by
Theorem~\ref{thm:v5v18-sharp-slab}.
\end{proof}

\subsection{The constraint inverse retains the boundary cost}

Let a spatial output operator \(G\) have multiplier \(g(k)\).

\begin{theorem}[Slab divergence after constraint solving]
\label{thm:v5v18-constraint-slab}
For every fixed \(k_0\ne0\) with \(g(k_0)\ne0\), the squared norm of the
smeared solved response \(G{:}\phi^2{:}\) is
\begin{equation}
 |g(k_0)|^2V_\varepsilon(k_0)
 \asymp_{a,k_0}\varepsilon^{-1}.
 \label{eq:v5v18-G-slab}
\end{equation}
For \(G=(-\Delta)^{-1}\), the prefactor is \(|k_0|^{-4}\).
\end{theorem}

\begin{proof}
The time smearing acts on the internal energy
\(|p|+|k-p|\), whereas \(G\) acts only on the fixed external spatial
momentum.  Fourier multiplication gives the displayed prefactor, and
Theorem~\ref{thm:v5v18-sharp-slab} gives the rate.
\end{proof}

This proves that the ultraviolet-uniform spatial elliptic estimate of V14
does not imply a uniform time-boundary estimate.

\subsection{The normalization tradeoff}

Define instead
\begin{equation}
 b_\varepsilon(t)=\varepsilon^{-1/2}a(t/\varepsilon).
 \label{eq:v5v18-L2-slab}
\end{equation}
Then
\[
 \widehat b_\varepsilon(\omega)
 =\varepsilon^{1/2}\widehat a(\varepsilon\omega).
\]

\begin{proposition}[Bounded \(L^2\) slab is not an approximate identity]
\label{prop:v5v18-L2-tradeoff}
The fixed-\(k\) two-particle norm for \(b_\varepsilon\) is bounded above
and below by positive constants as \(\varepsilon\downarrow0\).  However,
\begin{equation}
 \int b_\varepsilon(t)\,dt=\varepsilon^{1/2}\longrightarrow0,
 \label{eq:v5v18-L2-mass}
\end{equation}
so \(b_\varepsilon\to0\) as a distribution rather than to
\(\delta_0\).
\end{proposition}

\begin{proof}
The squared amplitude gains a factor \(\varepsilon\) relative to
\eqref{eq:v5v18-slab-variance}; multiply
\eqref{eq:v5v18-epsilon-inverse} by \(\varepsilon\).
Equation \eqref{eq:v5v18-L2-mass} follows by a change of variables.  For a test function \(\psi\), the change of variables \(t=\varepsilon u\)
gives
\[
 \left|\int b_\varepsilon(t)\psi(t)\,dt\right|
 =\varepsilon^{1/2}
  \left|\int a(u)\psi(\varepsilon u)\,du\right|
 \le\varepsilon^{1/2}\|a\|_{L^1}\|\psi\|_{L^\infty}
 \longrightarrow0.
\]
\end{proof}

\begin{firewall}[A bounded vanishing slab is not a constraint boundary]
\label{fw:v5v18-normalization}
Replacing \(a_\varepsilon\) by \(b_\varepsilon\) removes the divergent
norm by shrinking the total time mass.  It does not define the
instantaneous Hamiltonian constraint.  Restoring unit time mass restores
the factor \(\varepsilon^{-1}\).
\end{firewall}

\subsection{The leading boundary contact matrix}

Let a complete finite family of renormalized source components be
\(\mathcal S_{\alpha,\varepsilon}(k)\), \(1\le\alpha\le J\).  Suppose their
smeared covariance has the expansion
\begin{equation}
 \mathbb E\left[
 \mathcal S_{\alpha,\varepsilon}(k)
 \overline{\mathcal S_{\beta,\varepsilon}(k)}
 \right]
 =
 \varepsilon^{-r}M_{\alpha\beta}(k)
 +R_{\alpha\beta,\varepsilon}(k),
 \label{eq:v5v18-contact-matrix}
\end{equation}
where \(r>0\), \(M(k)=M(k)^*\), and \(R_\varepsilon\) is uniformly bounded.

\begin{lemma}[Positivity of the leading contact]
\label{lem:v5v18-contact-positive}
The matrix \(M(k)\) is positive semidefinite.
\end{lemma}

\begin{proof}
For every coefficient vector \(c\), covariance positivity gives
\[
 0\le
 \mathbb E\left|\sum_\alpha
 c_\alpha\mathcal S_{\alpha,\varepsilon}(k)\right|^2
 =
 \varepsilon^{-r}c^*M(k)c+O_c(1).
\]
Multiply by \(\varepsilon^r\) and let
\(\varepsilon\downarrow0\).  The limit is \(c^*M(k)c\ge0\).
\end{proof}

\begin{theorem}[Exact leading-contact cancellation criterion]
\label{thm:v5v18-contact-null}
For a constraint combination
\[
 \mathcal C_\varepsilon(k)
 =\sum_\alpha c_\alpha(k)
   \mathcal S_{\alpha,\varepsilon}(k),
\]
its variance is uniformly bounded as
\(\varepsilon\downarrow0\) if and only if
\begin{equation}
 M(k)^{1/2}c(k)=0
 \label{eq:v5v18-contact-null}
\end{equation}
and the remainders are uniformly bounded on that vector.  Equivalently,
\(c(k)^*M(k)c(k)=0\).
\end{theorem}

\begin{proof}
Substitution into \eqref{eq:v5v18-contact-matrix} gives
\[
 \mathbb E|\mathcal C_\varepsilon(k)|^2
 =
 \varepsilon^{-r}c^*M(k)c+c^*R_\varepsilon(k)c.
\]
The second term is bounded.  Since \(M\ge0\) by
Lemma~\ref{lem:v5v18-contact-positive}, the leading coefficient vanishes
if and only if
\(\|M^{1/2}c\|^2=0\).  This is necessary and sufficient under the stated
remainder bound.
\end{proof}

\begin{corollary}[Termwise renormalization is insufficient]
\label{cor:v5v18-matrix-needed}
A Ward identity can produce a time-zero source only if the exact
constraint coefficient vector lies in the nullspace of the complete
leading contact matrix.  Bounds on individual diagonal entries cannot
prove this cancellation.
\end{corollary}

\begin{proof}
The nullspace condition contains the signed off-diagonal covariances.
Replacing the covariance form by the sum of diagonal norms deletes those
terms and cannot establish
\eqref{eq:v5v18-contact-null}.
\end{proof}

\subsection{Local counterterms and boundary contacts}

A one-point identity counterterm changes a mean but not the positive
matrix \(M(k)\) at nonzero \(k\).  A renormalization of a stress
two-point function can change local contact distributions and hence the
matrix in \eqref{eq:v5v18-contact-matrix}; that choice must be fixed by
the Ward identities, reflection positivity, and the gravitational local
counterterm scheme.

\begin{firewall}[Subtracting positive covariance is not automatically
admissible]
\label{fw:v5v18-positive-contact}
Removing a positive divergent contact from a covariance by hand need not
leave a positive Schwinger functional.  An admissible subtraction must
arise from a common local action or operator renormalization and pass the
OS form.  The V5 essential quotient and V10 reflection test therefore
remain active.
\end{firewall}

\subsection{Boundary constraint certificate}

\begin{definition}[Slab-to-boundary constraint certificate]
\label{def:v5v18-SBCC}
An SBCC consists of:
\begin{enumerate}
 \item a unit-mass positive-time approximate identity;
 \item a Ward-complete covariance expansion
 \eqref{eq:v5v18-contact-matrix} through the required marked sources;
 \item the exact nullspace condition
 \eqref{eq:v5v18-contact-null} for every divergent boundary exponent and
 external mode;
 \item a uniformly bounded remainder in the renormalized constraint source
 space;
 \item compatibility of all contact choices with OS positivity and the
 local gravitational counterterms;
 \item application of the V14 spatial constraint inverse after the
 boundary limit.
\end{enumerate}
\end{definition}

\begin{theorem}[Free-square slab status]
\label{thm:v5v18-status}
The single free-scalar square has a leading boundary exponent \(r=1\) and
a strictly positive one-component contact coefficient at every fixed
nonzero spatial mode.  It therefore fails the SBCC.  A complete
Ward-related source can pass only through a genuine matrix nullspace
cancellation or a separately constructed boundary operator.
\end{theorem}

\begin{proof}
The exponent and positivity are
Theorem~\ref{thm:v5v18-sharp-slab}.  For a one-component matrix,
strict positivity gives a trivial nullspace, so
Theorem~\ref{thm:v5v18-contact-null} excludes a nonzero constraint
coefficient.
\end{proof}

V18 turns the time-zero problem into a finite algebraic test after the
asymptotic covariance is known.  The constraint Green operator cannot pay
the boundary contact; the exact source combination must annihilate it
before spatial inversion.

V19 will apply this test to conserved sources at the spectral level.  It
will prove that a nonzero positive spin-two spectral measure cannot lie in
a contact nullspace, while longitudinal constraint combinations may.
This will separate the physical TT matter backreaction from auxiliary
Hamiltonian components and decide which source can admit a time-zero
boundary.



\clearpage
\section{V19: spectral mass and the time-zero constraint boundary}
\label{sec:v5v19}

V18 formulated the time-zero problem through a leading contact matrix.
This note gives the corresponding exact spectral criterion.  For a
positive-time source, the norm of an approximate boundary value increases
to the total mass of its OS spectral measure.  Longitudinal Ward
combinations can lie in the spectral nullspace, but the free-scalar
spin-two stress channel has a positive scale-invariant density with
infinite total mass.  A spatial elliptic inverse cannot change that
energy-tail conclusion.

\subsection{A one-sided spectral boundary theorem}

Fix a nonzero spatial momentum \(\mathbf k\).  Let
\(\Sigma_{\mathbf k}\) be a positive operator-valued measure on
\([0,\infty)\) acting on a finite source-component space \(E\).  Suppose
the reflected two-point kernel has the representation
\begin{equation}
 K_{\mathbf k}(t,u)
 =
 \int_{[0,\infty)}
 e^{-\omega(t+u)}\,d\Sigma_{\mathbf k}(\omega),
 \qquad t,u>0.
 \label{eq:v5v19-OS-spectral-kernel}
\end{equation}
For \(c\in E\), write
\[
 d\nu_c(\omega)
 =\langle c,d\Sigma_{\mathbf k}(\omega)c\rangle.
\]
Use the positive-time approximate identity
\begin{equation}
 a_\varepsilon(t)
 =\varepsilon^{-1}e^{-t/\varepsilon}\mathbf 1_{\{t\ge0\}}.
 \label{eq:v5v19-exponential-slab}
\end{equation}
Its Laplace transform is
\[
 A_\varepsilon(\omega)
 =\int_0^\infty a_\varepsilon(t)e^{-\omega t}\,dt
 =\frac1{1+\varepsilon\omega}.
\]

\begin{theorem}[Exact spectral boundary criterion]
\label{thm:v5v19-spectral-boundary}
The OS squared norm of the smeared source component \(c\) is
\begin{equation}
 \mathcal N_\varepsilon(c)
 =
 \int_{[0,\infty)}
 \frac{d\nu_c(\omega)}{(1+\varepsilon\omega)^2}.
 \label{eq:v5v19-boundary-norm}
\end{equation}
Moreover,
\begin{equation}
 \lim_{\varepsilon\downarrow0}\mathcal N_\varepsilon(c)
 =
 \nu_c([0,\infty])\in[0,\infty].
 \label{eq:v5v19-total-mass}
\end{equation}
Thus the source has a finite Hilbert-space time-zero boundary vector in
direction \(c\) if and only if its scalar spectral measure has finite total
mass.
\end{theorem}

\begin{proof}
Insert the smeared source in both variables of
\eqref{eq:v5v19-OS-spectral-kernel}.  Tonelli's theorem gives
\[
 \mathcal N_\varepsilon(c)
 =
 \int |A_\varepsilon(\omega)|^2\,d\nu_c(\omega),
\]
which is \eqref{eq:v5v19-boundary-norm}.  The integrand increases
pointwise to one as \(\varepsilon\downarrow0\).  Monotone convergence gives
\eqref{eq:v5v19-total-mass}.  Finite limiting norm is exactly the
Cauchy boundary-vector condition for this monotone spectral
regularization.
\end{proof}

\begin{corollary}[Spatial multipliers do not improve energy mass]
\label{cor:v5v19-spatial-no-help}
Let \(G(\mathbf k)\) be a fixed bounded matrix at the chosen nonzero
spatial momentum.  The boundary criterion for \(Gc\) is the total mass of
\[
 \langle c,G(\mathbf k)^*
 d\Sigma_{\mathbf k}(\omega)G(\mathbf k)c\rangle.
\]
A nonzero scalar spatial multiplier changes this mass by a fixed factor
and cannot turn an infinite energy tail into a finite one.
\end{corollary}

\begin{proof}
The multiplier is independent of \(\omega\) and can be moved into the two
test vectors in \eqref{eq:v5v19-boundary-norm}.
\end{proof}

This is the spectral form of the V16--V18 conclusion for the constraint
Green operator.

\subsection{Spectral nullspaces and Ward directions}

Define the common spectral nullspace
\begin{equation}
 \mathcal Z_{\mathbf k}
 =
 \left\{c\in E:
 \Sigma_{\mathbf k}(B)c=0
 \text{ for every Borel }B\subset[0,\infty)\right\}.
 \label{eq:v5v19-spectral-nullspace}
\end{equation}

\begin{lemma}[Positive-measure null criterion]
\label{lem:v5v19-null}
For \(c\in E\), the following are equivalent:
\begin{enumerate}
 \item \(c\in\mathcal Z_{\mathbf k}\);
 \item \(\nu_c([0,\infty])=0\);
 \item \(\mathcal N_\varepsilon(c)=0\) for one, hence every,
 \(\varepsilon>0\).
\end{enumerate}
\end{lemma}

\begin{proof}
Positive operator-valued measure implies
\[
 \nu_c(B)=\|\Sigma_{\mathbf k}(B)^{1/2}c\|^2.
\]
Thus vanishing of the total scalar measure is equivalent to vanishing on
every Borel set and to item 1.  The weight in
\eqref{eq:v5v19-boundary-norm} is strictly positive, so its integral
vanishes exactly when the positive measure \(\nu_c\) vanishes.
\end{proof}

For a conserved stress tensor, every separated-point spectral tensor is
transverse to the full four-momentum \(p\).  Consequently a tensor test of
the form
\begin{equation}
 c_{\mu\nu}=p_\mu v_\nu+p_\nu v_\mu
 \label{eq:v5v19-longitudinal-test}
\end{equation}
lies in the separated spectral nullspace.

\begin{proposition}[Longitudinal Ward nullspace]
\label{prop:v5v19-longitudinal-null}
After the Ward contact terms have been assigned to their local owners,
the nonlocal spectral measure of a conserved stress tensor annihilates
every test \eqref{eq:v5v19-longitudinal-test}.
\end{proposition}

\begin{proof}
Each stress index pair is transverse:
\(p^\mu d\Sigma_{\mu\nu,\rho\sigma}=0\).
Contracting either term of
\eqref{eq:v5v19-longitudinal-test} with the first index pair therefore
gives zero.  Positivity then places the test in the nullspace of the
nonlocal measure.  Local contact polynomials were explicitly excluded and
must be treated by their counterterms.
\end{proof}

\begin{firewall}[A component is not automatically longitudinal]
\label{fw:v5v19-component}
The energy component \(T_{00}\), or the matter combination appearing in a
Hamiltonian constraint, need not have the form
\eqref{eq:v5v19-longitudinal-test}.  Conservation relates components but
does not imply that each component lies in the spectral nullspace.
The exact constraint tensor must be projected before the proposition is
used.
\end{firewall}

\subsection{The spin-two measure has infinite boundary mass}

For the four-dimensional massless free scalar, V17 proved that the
spin-two stress spectral density is positive and nonzero.  Scale
invariance fixes its dependence on the invariant mass \(s\).

\begin{lemma}[Scale of the free-scalar spin-two density]
\label{lem:v5v19-spin2-scale}
Up to a positive constant \(c_T\), the separated spin-two spectral density
is
\begin{equation}
 d\rho_2(s)=c_Ts^2\,ds
 \qquad(s>0).
 \label{eq:v5v19-spin2-density}
\end{equation}
\end{lemma}

\begin{proof}
The stress tensor has scaling dimension four.  Its two-point function has
momentum degree four after removal of local contacts, so its spectral
density has degree two in \(s=p^2\).  The massless theory has no other
scale.  Positivity and the nonzero two-particle matrix element from
Theorem~\ref{thm:v5v17-nonzero-spin2} give \(c_T>0\).
\end{proof}

At fixed \(\mathbf k\), the energy variable satisfies
\(s=\omega^2-|\mathbf k|^2\) in the Lorentzian spectral parametrization.
The corresponding positive measure inherits polynomial growth.

\begin{theorem}[No spin-two stress boundary vector]
\label{thm:v5v19-no-TT-boundary}
For every nonzero TT polarization that couples to the free-scalar stress
at fixed \(\mathbf k\), the associated energy spectral measure has
infinite total mass.  Hence the stress applied at a sharp time does not
define a Hilbert-space vector in that direction.
\end{theorem}

\begin{proof}
Equation \eqref{eq:v5v19-spin2-density} is positive for arbitrarily large
\(s\) and grows as \(s^2\).  Changing variables from \(s\) to
\(\omega\) multiplies by \(2\omega\); the remaining fixed-momentum
polarization and free phase-space factors are polynomial and do not
supply ultraviolet decay.  Therefore the energy spectral measure has
infinite total mass.  Apply
Theorem~\ref{thm:v5v19-spectral-boundary}.
\end{proof}

\begin{corollary}[The physical TT source cannot be a contact cancellation]
\label{cor:v5v19-TT-not-null}
The free-scalar spin-two stress direction is not in
\(\mathcal Z_{\mathbf k}\).  No Ward cancellation with a purely
longitudinal source can remove its positive spectral mass.
\end{corollary}

\begin{proof}
Its scalar measure is nonzero, indeed infinite, by
Theorem~\ref{thm:v5v19-no-TT-boundary}; use
Lemma~\ref{lem:v5v19-null}.  Longitudinal and spin-two sectors are
orthogonal under the projectors of V17.
\end{proof}

\subsection{Boundary forms instead of boundary vectors}

Failure of a time-zero vector does not prevent a spacetime operator-valued
distribution or an integrated evolution equation.  It means that the
constraint cannot be imposed by applying the unsmeared local stress to the
vacuum as an ordinary Hilbert vector.

\begin{proposition}[A weighted energy domain can define a boundary form]
\label{prop:v5v19-boundary-form}
Let \(H\ge0\) be the reconstructed Hamiltonian and let a source have
spectral measure \(\nu_c\).  If
\begin{equation}
 \int(1+\omega)^{-2r}\,d\nu_c(\omega)<\infty
 \label{eq:v5v19-negative-energy-domain}
\end{equation}
for some \(r>0\), then the time-zero source defines a continuous linear
functional on the domain of \((1+H)^r\), although it need not define a
Hilbert-space vector.
\end{proposition}

\begin{proof}
Insert \((1+H)^{-r}\) on the source vector.  Its squared spectral norm is
exactly the integral in
\eqref{eq:v5v19-negative-energy-domain}.  Cauchy--Schwarz then bounds the
source pairing by the graph norm of \((1+H)^r\).
\end{proof}

This is a possible route for canonical constraints, but the nonlinear
constraint solver V14 must then be rebuilt in a scale of operator forms
rather than random functions.

\subsection{Spectral boundary certificate}

\begin{definition}[Spectral constraint-boundary certificate]
\label{def:v5v19-SCBC}
An SCBC consists of:
\begin{enumerate}
 \item the positive operator-valued energy spectral measure of the complete
 Ward-renormalized constraint source;
 \item its exact longitudinal, scalar, and spin-two projections;
 \item finite total spectral mass for every claimed boundary-vector
 component, or a declared weighted form domain satisfying
 \eqref{eq:v5v19-negative-energy-domain};
 \item local contact owners compatible with the V5 curvature quotient;
 \item a version of the V14 constraint solver on the resulting boundary
 vector or form space;
 \item OS and shell compatibility of that solution.
\end{enumerate}
\end{definition}

\begin{theorem}[Free matter boundary status]
\label{thm:v5v19-status}
For free scalar matter, longitudinal separated-point stress combinations
satisfy the spectral null test after local contacts are removed, while the
physical spin-two stress has infinite time-zero spectral mass.  The latter
can enter a constrained continuum theory only as a spacetime distribution
or on a proved negative energy form domain, not as the \(H^s\) random
source assumed in V14--V15.
\end{theorem}

\begin{proof}
The longitudinal statement is
Proposition~\ref{prop:v5v19-longitudinal-null}; the spin-two statement is
Theorem~\ref{thm:v5v19-no-TT-boundary}.  The alternative form
interpretation is Proposition~\ref{prop:v5v19-boundary-form}.
\end{proof}

V19 completes the Gaussian boundary diagnosis.  The free physical
graviton covariance is positive after constraints, but the quantum matter
backreaction is not an instantaneous random Sobolev source.  A viable
canonical construction must use energy forms or spacetime-smeared
constraints, while a covariant construction must keep the full
renormalized stress insertion and its local curvature owners.

V20 will begin with the compulsory companion audit.  It will then formulate
a scale of negative Hamiltonian-energy spaces and prove a contraction
theorem for constraint solutions valued in forms, testing whether the
quadratic free stress lies within the required loss budget.


\clearpage
\section{V20: companion audit and constraint solving in negative energy
spaces}
\label{sec:v5v20}

This is the fourth compulsory companion checkpoint of the fifth series.
V19 showed that the physical stress source may fail to be a time-zero
Hilbert vector while remaining a continuous form on a positive-energy
domain.  This note constructs that Hilbert scale, proves the exact spectral
membership threshold for scale-invariant composites, and gives a
quantitative constraint contraction theorem in the form space.  The free
scalar square belongs after more than one half energy derivative of loss;
the spin-two stress requires more than five halves.  Nonlinear closure
still needs bounded products at that same loss.

\subsection{Compulsory companion audit}

At the time of this audit, the newest live companion files were
\begin{align*}
 &\texttt{Renewal\_Yang\_Mills\_Mass\_Gap\_Research\_Notes\_V33.tex},\\
 &\hspace{1em}\text{SHA--256 }
 \texttt{EFAD39097A75E6186ADA84F5AE52CE7E8207429AC7C460328199434AF3852446},
 \\
 &\texttt{Renewal\_NS\_Stability\_Research\_Notes\_V31.tex},\\
 &\hspace{1em}\text{SHA--256 }
 \texttt{F1121B62238AD5491BB7CF03635EA9934942EDF573ED8FAAA9EE79E1D38A6696}.
\end{align*}

Yang--Mills V33 closes its one-strong \(q\ge2\) direct fibres by keeping
positive moment words, exact duality, and all local factors assembled
before taking absolute values.  It leaves the \(q=0,1\) response deficits
explicit rather than absorbing them in an ordinary norm.  Navier--Stokes
V31 continues to show that the critical response requires the exact
critical source rung.

\begin{theorem}[V20 audit transfer]
\label{thm:v5v20-audit}
For a constraint source represented as an energy form:
\begin{enumerate}
 \item preserve the positive spectral measure until after the complete
 constraint combination is assembled;
 \item record the exact negative-energy loss of each source component;
 \item require the nonlinear maps to close at that same loss; and
 \item do not use a bounded weaker form norm as evidence for a time-zero
 vector or a stronger response.
\end{enumerate}
\end{theorem}

\begin{proof}
Items 1 and 2 are the form-space versions of the positive-word and explicit
deficit ledgers in Yang--Mills V33.  Items 3 and 4 transfer the exact source
rung of Navier--Stokes V31.  These rules determine the norms in the
theorems below.
\end{proof}

\subsection{The Hamiltonian Hilbert scale}

Let \(\mathscr H\) be a Hilbert space with a nonnegative self-adjoint
Hamiltonian \(H\).  For \(r\ge0\), define
\begin{align}
 \mathscr H_r
 &=
 D((1+H)^r),
 &
 \|u\|_r
 &=
 \|(1+H)^ru\|_{\mathscr H},
 \label{eq:v5v20-positive-scale}\\
 \mathscr H_{-r}
 &=
 \text{completion of }\mathscr H
 \text{ in }
 \|v\|_{-r}
 =
 \|(1+H)^{-r}v\|_{\mathscr H}.
 \label{eq:v5v20-negative-scale}
\end{align}

\begin{proposition}[Duality of the energy scale]
\label{prop:v5v20-duality}
With the pivot inner product of \(\mathscr H\),
\begin{equation}
 (\mathscr H_r)^*\cong\mathscr H_{-r}
 \label{eq:v5v20-duality}
\end{equation}
isometrically.  A vector or distribution \(S\) defines a continuous form
on \(\mathscr H_r\) precisely when
\((1+H)^{-r}S\in\mathscr H\), and
\begin{equation}
 |\langle S,u\rangle|
 \le\|S\|_{-r}\|u\|_r.
 \label{eq:v5v20-form-bound}
\end{equation}
\end{proposition}

\begin{proof}
For pivot vectors,
\[
 \langle S,u\rangle
 =
 \left\langle
  (1+H)^{-r}S,\,
  (1+H)^ru
 \right\rangle_{\mathscr H}.
\]
Cauchy--Schwarz gives \eqref{eq:v5v20-form-bound}.  Conversely, the Riesz
representation theorem applied after the isometry
\(u\mapsto(1+H)^ru\) identifies every continuous functional with a unique
element of the completion \(\mathscr H_{-r}\).
\end{proof}

Let \(S\Omega\) be a source applied to the vacuum, interpreted first with a
positive-time regulator.  If its scalar energy measure is \(d\nu(\omega)\),
then
\begin{equation}
 \|S\Omega\|_{-r}^2
 =
 \int_{[0,\infty)}(1+\omega)^{-2r}\,d\nu(\omega).
 \label{eq:v5v20-spectral-negative-norm}
\end{equation}

\begin{corollary}[V19 boundary vector and form alternatives]
\label{cor:v5v20-V19}
The V19 time-zero vector criterion is the case \(r=0\) of
\eqref{eq:v5v20-spectral-negative-norm}.  If it fails but the integral is
finite for some \(r>0\), the source has a time-zero value only as an
element of \(\mathscr H_{-r}\).
\end{corollary}

\begin{proof}
At \(r=0\), the integral is the total spectral mass from
Theorem~\ref{thm:v5v19-spectral-boundary}.  The general case is
Proposition~\ref{prop:v5v20-duality}.
\end{proof}

\subsection{Exact scale-invariant thresholds}

Suppose a scalar composite has positive invariant-mass density
\begin{equation}
 d\rho_\alpha(s)=c_\alpha s^\alpha\,ds,
 \qquad c_\alpha>0,
 \qquad s>0.
 \label{eq:v5v20-scale-density}
\end{equation}
At fixed spatial momentum \(\mathbf k\), put
\(\omega=(s+|\mathbf k|^2)^{1/2}\).  The OS time kernel contains
\[
 \frac{e^{-\omega(t+u)}}{2\omega}\,d\rho_\alpha(s).
\]

\begin{theorem}[Negative-energy membership threshold]
\label{thm:v5v20-threshold}
For the density \eqref{eq:v5v20-scale-density}, the fixed-spatial-mode
source lies in \(\mathscr H_{-r}\) if and only if
\begin{equation}
 r>\alpha+\frac12.
 \label{eq:v5v20-threshold}
\end{equation}
\end{theorem}

\begin{proof}
Since \(ds=2\omega\,d\omega\), the factor \(2\omega\) cancels the
\(1/(2\omega)\) in the time kernel.  Thus the scalar energy measure is
\begin{equation}
 d\nu_{\alpha,\mathbf k}(\omega)
 =
 c_\alpha(\omega^2-|\mathbf k|^2)^\alpha
 \mathbf 1_{\{\omega\ge|\mathbf k|\}}\,d\omega.
 \label{eq:v5v20-energy-density}
\end{equation}
At large \(\omega\), its density is comparable to
\(\omega^{2\alpha}\).  Therefore
\eqref{eq:v5v20-spectral-negative-norm} converges exactly when
\[
 2\alpha-2r<-1,
\]
which is \eqref{eq:v5v20-threshold}.  The lower endpoint is harmless for
the fixed nonzero spatial mode; the zero-mode infrared endpoint requires a
separate test.
\end{proof}

\begin{corollary}[Free square and stress losses]
\label{cor:v5v20-free-losses}
For a free massless scalar in four dimensions:
\begin{enumerate}
 \item the separated scalar square has \(\alpha=0\) and belongs to
 \(\mathscr H_{-r}\) exactly for \(r>1/2\);
 \item the spin-two stress channel has \(\alpha=2\) and belongs to
 \(\mathscr H_{-r}\) exactly for \(r>5/2\).
\end{enumerate}
Neither is a time-zero Hilbert vector.
\end{corollary}

\begin{proof}
The scalar square has scaling dimension two, so its two-point invariant
density has degree zero.  Lemma~\ref{lem:v5v19-spin2-scale} gives degree
two for the stress.  Apply
Theorem~\ref{thm:v5v20-threshold}.  The exclusion of \(r=0\) also follows
from V18--V19.
\end{proof}

The exponent \(5/2\) is the exact free spin-two energy loss.  A constraint
theorem using less loss cannot contain that source, while one using more
loss must show that every nonlinear operation remains defined there.

\subsection{Spatial constraint operators on the energy scale}

Let \(\mathscr Y_{-r}\) and \(\mathscr Z_{-r}\) be source and auxiliary
form spaces built over the same Hamiltonian scale.  A spatial constraint
inverse \(G\) is admissible at loss \(r\) when
\begin{equation}
 \|Gf\|_{\mathscr Z_{-r}}
 \le M_r\|f\|_{\mathscr Y_{-r}}.
 \label{eq:v5v20-G-weighted}
\end{equation}

\begin{proposition}[Commuting spatial inverse]
\label{prop:v5v20-commuting-G}
If \(G\) acts on spatial labels, commutes with \(H\), and is bounded from
the pivot source space to the pivot auxiliary space with norm \(M_0\),
then \eqref{eq:v5v20-G-weighted} holds with \(M_r=M_0\) for every
\(r\ge0\).
\end{proposition}

\begin{proof}
Functional calculus and commutation give
\[
 (1+H)^{-r}Gf=G(1+H)^{-r}f.
\]
Apply the pivot bound for \(G\).
\end{proof}

\begin{firewall}[Interacting time evolution need not commute]
\label{fw:v5v20-no-commutation}
For the exact interacting Hamiltonian, a spatial inverse with
field-dependent coefficients need not commute with \(H\).  The required
quantity is the weighted conjugation norm in
\eqref{eq:v5v20-G-weighted}, not its free pivot norm.
\end{firewall}

\subsection{A constraint contraction in form space}

Let \(X\) be a Banach space of physical input forms, and let
\(Z_{-r},Y_{-r}\) be Banach spaces modeled on the energy loss \(r\).
Assume an invertible linear constraint
\[
 L:Z_{-r}\to Y_{-r},
 \qquad
 \|L^{-1}\|\le M_r,
\]
and an assembled constraint
\begin{equation}
 \mathcal C(x,z)
 =
 Lz+B(z,z)+D(x,z)+S(x).
 \label{eq:v5v20-form-constraint}
\end{equation}
Suppose
\begin{align}
 \|B(z_1,z_2)\|_{Y_{-r}}
 &\le C_B\|z_1\|_{Z_{-r}}\|z_2\|_{Z_{-r}},
 \label{eq:v5v20-form-B}\\
 \|D(x,z)\|_{Y_{-r}}
 &\le C_D\|x\|_X\|z\|_{Z_{-r}},
 \label{eq:v5v20-form-D}\\
 \|S(x)\|_{Y_{-r}}
 &\le C_S(\|x\|_X+\|x\|_X^2).
 \label{eq:v5v20-form-S}
\end{align}

\begin{theorem}[Negative-energy constraint contraction]
\label{thm:v5v20-form-contraction}
Define
\begin{align}
 C_0&=2C_S,\\
 K_r&=M_r(4M_rC_BC_0+C_D),\\
 \delta_r&=\min\left\{
  1,\,
  \frac{R_{\max}}{2M_rC_0},\,
  \frac1{2K_r}
 \right\}.
 \label{eq:v5v20-form-radius}
\end{align}
For \(\|x\|_X\le\delta_r\), the equation
\(\mathcal C(x,z)=0\) has a unique solution in the ball
\begin{equation}
 \|z\|_{Z_{-r}}\le2M_rC_0\|x\|_X.
 \label{eq:v5v20-form-solution}
\end{equation}
The contraction constant is at most \(1/2\).
\end{theorem}

\begin{proof}
Use the fixed-point map
\[
 \mathcal T_xz=-L^{-1}\bigl(B(z,z)+D(x,z)+S(x)\bigr).
\]
On the displayed ball, the difference in \(z\) is bounded by
\[
 M_r(2C_BR+C_D\|x\|)\|z-\widetilde z\|
 \le K_r\|x\|\|z-\widetilde z\|
 \le\frac12\|z-\widetilde z\|.
\]
For \(\|x\|\le1\),
\(\|S(x)\|\le C_0\|x\|\).  Exactly as in
Theorem~\ref{thm:v5v15-explicit-radius}, the source occupies half of the
chosen radius and the nonlinear difference occupies at most the other
half.  Banach's theorem gives existence and uniqueness.
\end{proof}

\begin{firewall}[The bilinear form bound is the hard hypothesis]
\label{fw:v5v20-bilinear-hard}
Negative Hamiltonian-energy spaces are not automatically algebras.
Corollary~\ref{cor:v5v20-free-losses} proves membership of one renormalized
source, but it does not prove
\eqref{eq:v5v20-form-B} for products of two auxiliary form-valued metrics.
If the product loses additional energy regularity, the fixed Banach-space
iteration in Theorem~\ref{thm:v5v20-form-contraction} does not close.
\end{firewall}

\subsection{A precise loss-of-regularity boundary}

Let a scale \(Z_{-r}\) be nested so that increasing \(r\) weakens the norm.
Suppose instead of \eqref{eq:v5v20-form-B} one only has
\begin{equation}
 B:Z_{-r}\times Z_{-r}\longrightarrow Y_{-(r+\ell)}
 \qquad(\ell>0).
 \label{eq:v5v20-lossy-B}
\end{equation}

\begin{proposition}[Fixed-space iteration cannot absorb an uncompensated
loss]
\label{prop:v5v20-loss-no-close}
The estimate \eqref{eq:v5v20-lossy-B} alone does not define the
fixed-point map of
Theorem~\ref{thm:v5v20-form-contraction} on \(Z_{-r}\).  Closure requires
either:
\begin{enumerate}
 \item an inverse \(L^{-1}\) that gains at least \(\ell\) energy
 derivatives;
 \item a shell or time smoother supplying that gain before multiplication;
 or
 \item a tame iteration on a scale with an explicit shrinking-radius
 budget.
\end{enumerate}
\end{proposition}

\begin{proof}
The nonlinear image lies only in \(Y_{-(r+\ell)}\).  If
\(L^{-1}\) is bounded merely from \(Y_{-r}\) to \(Z_{-r}\), its application
to that image is not defined with the target norm used by the contraction.
Each listed alternative supplies the missing map back to \(Z_{-r}\);
without one, the iteration has no self-map property.
\end{proof}

The spatial inverse of V14 gains spatial derivatives but, by itself, does
not gain Hamiltonian energy derivatives.  These two regularity ledgers
must remain distinct.

\subsection{The form-valued constraint certificate}

\begin{definition}[Hamiltonian-form constraint certificate]
\label{def:v5v20-HFCC}
An HFCC consists of:
\begin{enumerate}
 \item an OS-reconstructed nonnegative Hamiltonian and its scale
 \(\mathscr H_{\pm r}\);
 \item Ward-renormalized constraint sources with uniform spectral bounds
 in \(\mathscr H_{-r}\);
 \item a weighted spatial inverse satisfying
 \eqref{eq:v5v20-G-weighted};
 \item the closed bilinear bounds
 \eqref{eq:v5v20-form-B}--\eqref{eq:v5v20-form-D}, or a proved
 loss-compensating scale iteration;
 \item the V14 marked, Jacobian, all-root, and shell compatibility
 statements in the form topology;
 \item physical OS positivity after the constraint reduction.
\end{enumerate}
\end{definition}

\begin{theorem}[Conditional form-valued reduction]
\label{thm:v5v20-HFCC}
An HFCC supplies a unique small form-valued auxiliary metric by
Theorem~\ref{thm:v5v20-form-contraction}.  The free scalar square is
eligible at every loss \(r>1/2\), and its spin-two stress is eligible at
every loss \(r>5/2\), provided the remaining HFCC bounds hold.  No such
HFCC is presently proved for the interacting SM--Einstein system.
\end{theorem}

\begin{proof}
The thresholds are Corollary~\ref{cor:v5v20-free-losses}.  The contraction
is Theorem~\ref{thm:v5v20-form-contraction}.  The eligibility statement is
conditional because the bilinear, interacting Hamiltonian, and shell
bounds are separate items of the definition.
\end{proof}

V20 converts the time-zero divergence from an absolute obstruction into a
quantified form loss.  For free matter the loss is finite, but the
nonlinear constraint map must close at that loss or supply a compensating
smoother.  This is now the immediate metric-source bottleneck.

V21 will construct a two-index spatial/energy Hilbert scale and prove a
loss-compensating shell-smoothing contraction.  The theorem will show
exactly how many spatial and Hamiltonian derivatives a shell must gain to
multiply two form-valued auxiliary metrics without an infinite descent of
regularity.


\clearpage
\section{V21: two-index smoothing and the nonlinear loss budget}
\label{sec:v5v21}

V20 placed quantum constraint sources in negative Hamiltonian-energy
spaces and showed that a fixed-space contraction fails if nonlinear
products lose further regularity.  This note introduces a joint
spatial/energy scale and proves a loss-compensating shell theorem.
Spatial heat time \(\tau\) gains \(a\) derivatives at cost
\(\tau^{-a/2}\), while positive Hamiltonian time \(t\) gains \(b\) energy
derivatives at cost \(t^{-b}\).  At ultraviolet scale \(K\), a canonical
suppression \(K^{-\delta}\) beats the loss only when
\(\delta>a+b\).  The equality case is genuinely critical.

\subsection{A joint spatial and Hamiltonian scale}

Let \(\mathscr H\) be a Hilbert space.  Let \(A\ge1\) and \(1+H\ge1\) be
strongly commuting positive self-adjoint operators, representing spatial
ellipticity and reconstructed energy.  For \(s\in\mathbb R\) and
\(r\ge0\), define
\begin{equation}
 \|u\|_{s,-r}
 =
 \|A^{s/2}(1+H)^{-r}u\|_{\mathscr H},
 \label{eq:v5v21-two-index-norm}
\end{equation}
and let \(\mathscr H^{s,-r}\) be the corresponding completion of the
common spectral core.

\begin{proposition}[Monotone embeddings]
\label{prop:v5v21-embeddings}
If \(s_1\ge s_0\) and \(r_1\le r_0\), then
\[
 \mathscr H^{s_1,-r_1}
 \hookrightarrow
 \mathscr H^{s_0,-r_0}
\]
continuously with norm at most one.
\end{proposition}

\begin{proof}
By joint spectral calculus, the quotient multiplier is
\[
 \lambda_A^{(s_0-s_1)/2}
 (1+\lambda_H)^{-(r_0-r_1)},
\]
whose absolute value is at most one because
\(\lambda_A,1+\lambda_H\ge1\).
\end{proof}

The indices record different mechanisms.  A spatial inverse Laplacian may
increase \(s\) without changing \(r\); a positive-time semigroup may
decrease the energy loss \(r\) without changing spatial regularity.

\subsection{Sharp smoothing estimates}

Define
\[
 \mathcal S_{\tau,t}=e^{-\tau A}e^{-tH},
 \qquad
 \tau,t>0.
\]

\begin{theorem}[Two-index semigroup smoothing]
\label{thm:v5v21-semigroup}
For \(a,b\ge0\),
\begin{equation}
 \|\mathcal S_{\tau,t}u\|_{s+a,-r}
 \le
 C_{a,b}\tau^{-a/2}t^{-b}
 \|u\|_{s,-(r+b)}.
 \label{eq:v5v21-semigroup-bound}
\end{equation}
The powers of \(\tau\) and \(t\) are sharp whenever the joint spectrum
contains arbitrarily large comparable spatial and energy eigenvalues.
\end{theorem}

\begin{proof}
Joint spectral calculus reduces the operator norm to
\[
 \sup_{\lambda_A\ge1,\lambda_H\ge0}
 \lambda_A^{a/2}e^{-\tau\lambda_A}
 (1+\lambda_H)^be^{-t\lambda_H}.
\]
The elementary maxima satisfy
\[
 \sup_{\lambda\ge1}\lambda^{a/2}e^{-\tau\lambda}
 \le C_a\tau^{-a/2},
 \qquad
 \sup_{\lambda\ge0}(1+\lambda)^be^{-t\lambda}
 \le C_b t^{-b}
\]
for \(0<\tau,t\le1\); larger times only improve the bound after adjusting
the constant.  This proves \eqref{eq:v5v21-semigroup-bound}.

For sharpness, choose spectral vectors with
\(\lambda_A\asymp\tau^{-1}\) and
\(\lambda_H\asymp t^{-1}\).  The two multipliers are then bounded below
by constant multiples of
\(\tau^{-a/2}\) and \(t^{-b}\), respectively.
\end{proof}

\begin{corollary}[Dyadic ultraviolet cost]
\label{cor:v5v21-dyadic-cost}
At a shell of spatial scale \(K\), take
\[
 \tau=c_\tau K^{-2},
 \qquad
 t=c_tK^{-1}.
\]
Then
\begin{equation}
 \|\mathcal S_{\tau,t}u\|_{s+a,-r}
 \le C K^{a+b}\|u\|_{s,-(r+b)}.
 \label{eq:v5v21-K-cost}
\end{equation}
\end{corollary}

\begin{proof}
Insert the two scale relations into
\eqref{eq:v5v21-semigroup-bound}.
\end{proof}

Positive Hamiltonian time is the same regulator whose removal was studied
in V18--V20.  Equation \eqref{eq:v5v21-K-cost} records its cost rather
than treating it as a free regularization.

\subsection{A loss-compensating nonlinear map}

Let \(X_K\) be a physical input space at shell scale \(K\).  Let
\[
 Z=\mathscr H^{s+a,-r},
 \qquad
 Y=\mathscr H^{s,-(r+b)}.
\]
Assume bilinear and mixed maps satisfy
\begin{align}
 \|B(z_1,z_2)\|_Y
 &\le C_B\|z_1\|_Z\|z_2\|_Z,
 \label{eq:v5v21-loss-B}\\
 \|D(x,z)\|_Y
 &\le C_D\|x\|_{X_K}\|z\|_Z,
 \label{eq:v5v21-loss-D}\\
 \|S(x)\|_Y
 &\le C_S(\|x\|_{X_K}+\|x\|_{X_K}^2).
 \label{eq:v5v21-loss-S}
\end{align}
Let the inverse shell response be
\begin{equation}
 \mathcal G_K
 =
 K^{-\delta}\mathcal S_{c_\tau K^{-2},c_tK^{-1}}
 \mathcal L^{-1},
 \label{eq:v5v21-shell-response}
\end{equation}
where
\(\mathcal L^{-1}:Y\to Y\) has norm at most \(M\).  The regulated
constraint equation is
\begin{equation}
 z=-\mathcal G_K\bigl(B(z,z)+D(x,z)+S(x)\bigr).
 \label{eq:v5v21-regulated-constraint}
\end{equation}

\begin{lemma}[Net response norm]
\label{lem:v5v21-net-response}
There is \(C_G\), independent of \(K\ge1\), such that
\begin{equation}
 \|\mathcal G_K\|_{Y\to Z}
 \le C_GK^{-g},
 \qquad
 g:=\delta-a-b.
 \label{eq:v5v21-net-gap}
\end{equation}
\end{lemma}

\begin{proof}
Apply Corollary~\ref{cor:v5v21-dyadic-cost} after
\(\mathcal L^{-1}\) and multiply by \(K^{-\delta}\).
\end{proof}

\begin{theorem}[Loss-compensating shell contraction]
\label{thm:v5v21-loss-contraction}
Assume \(g\ge0\).  Put
\[
 M_K=C_GK^{-g}
\]
and define
\begin{align}
 C_0&=2C_S,\\
 \mathcal K_K&=M_K(4M_KC_BC_0+C_D),\\
 \delta_K&=\min\left\{
  1,\,
  \frac{R_{\max}}{2M_KC_0},\,
  \frac1{2\mathcal K_K}
 \right\}.
 \label{eq:v5v21-shell-radius}
\end{align}
For every \(\|x\|_{X_K}\le\delta_K\),
equation \eqref{eq:v5v21-regulated-constraint} has a unique solution in
\[
 \|z\|_Z\le2M_KC_0\|x\|_{X_K},
 \label{eq:v5v21-shell-solution}
\]
with contraction constant at most \(1/2\).  If \(g>0\), then
\(M_K\to0\), and for fixed sufficiently small input radius the contraction
improves along ultraviolet shells.
\end{theorem}

\begin{proof}
The proof is the V20 contraction with \(M_r\) replaced by the net response
norm \(M_K\).  On the displayed ball,
\[
 \|\mathcal G_K(B(z,z)-B(\widetilde z,\widetilde z)
 +D(x,z-\widetilde z))\|_Z
 \le
 M_K(2C_BR+C_D\|x\|)\|z-\widetilde z\|_Z.
\]
The definitions in \eqref{eq:v5v21-shell-radius} make this at most one
half of the difference.  The source occupies at most half of the invariant
radius.  Banach's theorem gives the result.  When \(g>0\),
\eqref{eq:v5v21-net-gap} tends to zero.
\end{proof}

\begin{corollary}[The critical and supercritical alternatives]
\label{cor:v5v21-gap-trichotomy}
The shell response has three regimes:
\begin{enumerate}
 \item if \(\delta>a+b\), the ultraviolet response gains the strict factor
 \(K^{-(\delta-a-b)}\);
 \item if \(\delta=a+b\), smoothing only repays the product loss and no
 ultraviolet smallness follows;
 \item if \(\delta<a+b\), the response norm grows like
 \(K^{a+b-\delta}\), so a scale-independent input ball cannot be obtained
 from these estimates.
\end{enumerate}
\end{corollary}

\begin{proof}
This is the sign trichotomy for \(g\) in
\eqref{eq:v5v21-net-gap}.  Sharpness of the semigroup powers in
Theorem~\ref{thm:v5v21-semigroup} prevents improvement without additional
structure.
\end{proof}

\subsection{A compactness refinement}

Let \(s_{\rm s}>s_{\rm w}\) and \(r_{\rm s}<r_{\rm w}\).  On a compact
spatial manifold with discrete joint spectrum and finite joint spectral
multiplicities, the inclusion
\[
 \mathscr H^{s_{\rm s},-r_{\rm s}}
 \hookrightarrow
 \mathscr H^{s_{\rm w},-r_{\rm w}}
\]
is compact when the quotient spectral weight tends to zero along every
sequence escaping to infinity.

\begin{proposition}[Two-index compact embedding]
\label{prop:v5v21-compact}
Suppose
\begin{equation}
 \lambda_A^{(s_{\rm w}-s_{\rm s})/2}
 (1+\lambda_H)^{-(r_{\rm w}-r_{\rm s})}
 \longrightarrow0
 \label{eq:v5v21-compact-ratio}
\end{equation}
along the joint spectrum at infinity.  Then the displayed inclusion is
compact.
\end{proposition}

\begin{proof}
Joint spectral projections below a fixed compact spectral rectangle have
finite rank.  On its complement, the norm of the inclusion is the
supremum of the quotient multiplier in
\eqref{eq:v5v21-compact-ratio}, which tends to zero.  The inclusion is
therefore an operator-norm limit of finite-rank maps.
\end{proof}

A uniform strong bound for the regulated constraint image would then give
the tail tightness needed by V5, now simultaneously in spatial and energy
labels.

\subsection{Why positive-time smoothing does not yet solve the constraint}

The map \(\mathcal S_{\tau,t}\) uses \(t>0\).  At shell scale \(K\),
\(t\asymp K^{-1}\downarrow0\).  V18 proves that a unit-mass boundary
source may diverge with precisely such a shrinking time width.

\begin{firewall}[Regulated contraction is not a sharp-time constraint]
\label{fw:v5v21-regulated-only}
Theorem~\ref{thm:v5v21-loss-contraction} solves the regulated equation
\eqref{eq:v5v21-regulated-constraint}.  Passing to the instantaneous
constraint requires a uniform weak-space limit of the complete source and
solution as \(t\downarrow0\), including every contact nullspace condition
of V18 and spectral loss of V20.
\end{firewall}

\begin{proposition}[Critical loss cannot be hidden in the semigroup]
\label{prop:v5v21-critical-sharp}
At \(g=0\), there are joint spectral packets saturating
\[
 \|\mathcal G_K\|_{Y\to Z}\asymp1.
\]
Thus a contraction at the critical balance must come from a small coupling,
a Ward cancellation, a positive block structure, or a nonlinear
coefficient bound; shell smoothing alone gives none.
\end{proposition}

\begin{proof}
Choose the saturating spectral vectors from the sharpness part of
Theorem~\ref{thm:v5v21-semigroup}, with spatial eigenvalue of order
\(K^2\) and energy of order \(K\).  Multiplication by
\(K^{-\delta}\) cancels the factor \(K^{a+b}\) when
\(\delta=a+b\).
\end{proof}

\subsection{The two-index constraint certificate}

\begin{definition}[Loss-compensated shell constraint certificate]
\label{def:v5v21-LSCC}
An LSCC consists of:
\begin{enumerate}
 \item a strongly commuting or otherwise controlled spatial/energy scale;
 \item the complete product estimates
 \eqref{eq:v5v21-loss-B}--\eqref{eq:v5v21-loss-S};
 \item a net gap \(g\ge0\), with an additional smallness mechanism when
 \(g=0\);
 \item a compact strong-to-weak embedding satisfying
 \eqref{eq:v5v21-compact-ratio};
 \item a Ward-complete sharp-time limit satisfying the V18--V20 boundary
 criteria;
 \item compatibility with the V14 Jacobian, marked response, and physical
 OS quotient.
\end{enumerate}
\end{definition}

\begin{theorem}[Conditional two-index closure]
\label{thm:v5v21-LSCC}
An LSCC gives unique regulated form-valued constraint solutions at every
sufficiently high shell, a relatively compact weak family, and a
sharp-time physical constraint solution along every convergent subsequence
for which item 5 identifies a unique limit.
\end{theorem}

\begin{proof}
Existence and uniqueness are
Theorem~\ref{thm:v5v21-loss-contraction}.  Relative compactness is
Proposition~\ref{prop:v5v21-compact}.  The boundary and physical
identification follow from items 5--6; uniqueness makes all subsequential
limits agree whenever the limiting equation retains the contraction
property.
\end{proof}

No LSCC is presently proved for the coupled quantum constraint.  V21 gives
the exact acquisition target: determine the product losses \(a,b\) and the
canonical shell suppression \(\delta\) for the Ward-complete source.  If
the balance is critical, a structural cancellation must be retained before
norms are taken.

V22 will compute the two-index loss in a free quadratic model with a
positive-time shell.  It will compare the scalar-square threshold
\(r>1/2\) and stress threshold \(r>5/2\) with the energy gain of one
Duhamel integration, determining whether time integration rather than a
sharp constraint closes the source.


\clearpage
\section{V22: the Duhamel energy-gain ladder}
\label{sec:v5v22}

V20 found the negative-energy losses of free quadratic sources, and V21
showed how positive-time smoothing can repay a specified loss.  This note
computes the exact gain from genuine time integration.  One Duhamel
integration supplies one power of energy.  It turns the free scalar square
into a Hilbert vector, but the spin-two stress requires three integrations;
a second-order covariant Green response supplies only two high-energy
powers and remains at a half-derivative form loss.  This separates
instantaneous constraints, first-order evolution, and covariant
second-order response.

\subsection{Repeated positive-time integration}

Let \(H\ge0\) be self-adjoint.  For an integer \(m\ge1\), define
\begin{equation}
 \mathcal D_T^{(m)}
 =
 \int_0^T
 \frac{t^{m-1}}{(m-1)!}e^{-tH}\,dt,
 \qquad T>0.
 \label{eq:v5v22-Dm}
\end{equation}
By spectral calculus its multiplier is
\begin{equation}
 d_{m,T}(\omega)
 =
 \omega^{-m}\left[
 1-e^{-T\omega}
 \sum_{j=0}^{m-1}\frac{(T\omega)^j}{j!}
 \right],
 \label{eq:v5v22-Dm-multiplier}
\end{equation}
with continuous value \(T^m/m!\) at \(\omega=0\).

\begin{lemma}[High-energy Duhamel power]
\label{lem:v5v22-D-power}
For every \(m\ge1\) and fixed \(T>0\), there are constants
\(0<c_{m,T}<C_{m,T}<\infty\) such that
\begin{equation}
 c_{m,T}(1+\omega)^{-m}
 \le d_{m,T}(\omega)
 \le C_{m,T}(1+\omega)^{-m}
 \qquad(\omega\ge0).
 \label{eq:v5v22-D-comparison}
\end{equation}
\end{lemma}

\begin{proof}
The bracket in \eqref{eq:v5v22-Dm-multiplier} is the normalized lower
incomplete gamma function and is positive.  Near zero, Taylor expansion
gives
\[
 d_{m,T}(\omega)=\frac{T^m}{m!}+O(\omega).
\]
At infinity, the exponential polynomial tends to zero, so
\[
 d_{m,T}(\omega)=\omega^{-m}(1+o(1)).
\]
Continuity and positivity on the remaining compact interval give the two
global comparison constants.
\end{proof}

\begin{theorem}[Exact energy gain]
\label{thm:v5v22-energy-gain}
For every \(r\in\mathbb R\),
\begin{equation}
 \mathcal D_T^{(m)}:
 \mathscr H_{-r}\longrightarrow\mathscr H_{m-r}
 \label{eq:v5v22-energy-map}
\end{equation}
is bounded, using the natural extension of the V20 scale to all real
indices.  Equivalently, repeated time integration gains exactly \(m\)
Hamiltonian-energy derivatives.
\end{theorem}

\begin{proof}
The norm ratio in spectral variables is
\[
 (1+\omega)^{m-r}
 d_{m,T}(\omega)
 (1+\omega)^r
 =(1+\omega)^m d_{m,T}(\omega),
\]
which is bounded by Lemma~\ref{lem:v5v22-D-power}.  Sharpness of the lower
comparison in \eqref{eq:v5v22-D-comparison} excludes any larger uniform
gain on unbounded spectral support.
\end{proof}

\subsection{Scale-invariant source responses}

Use the invariant density of V20,
\[
 d\rho_\alpha(s)=c_\alpha s^\alpha\,ds.
\]
At fixed nonzero spatial momentum, its energy measure behaves as
\[
 d\nu_\alpha(\omega)\asymp\omega^{2\alpha}\,d\omega.
\]

\begin{theorem}[Integrated-source threshold]
\label{thm:v5v22-integrated-threshold}
The response
\(\mathcal D_T^{(m)}S_\alpha\Omega\) is a Hilbert-space vector if and only if
\begin{equation}
 m>\alpha+\frac12.
 \label{eq:v5v22-vector-threshold}
\end{equation}
More generally, it lies in \(\mathscr H_{-r}\) precisely when
\begin{equation}
 r+m>\alpha+\frac12.
 \label{eq:v5v22-form-threshold}
\end{equation}
\end{theorem}

\begin{proof}
The squared response norm has high-energy integrand comparable to
\[
 \omega^{-2m}\omega^{2\alpha}\,d\omega.
\]
It is integrable exactly when
\(2\alpha-2m<-1\), which is
\eqref{eq:v5v22-vector-threshold}.  Inserting the additional
\((1+\omega)^{-2r}\) form weight gives
\eqref{eq:v5v22-form-threshold}.
\end{proof}

\begin{corollary}[Free quadratic response ladder]
\label{cor:v5v22-free-ladder}
For a four-dimensional free scalar:
\begin{enumerate}
 \item the scalar square has \(\alpha=0\); one time integration makes its
 response a Hilbert vector;
 \item the TT stress has \(\alpha=2\); after one integration it lies in
 \(\mathscr H_{-r}\) for every \(r>3/2\), after two for every \(r>1/2\),
 and after three it is a Hilbert vector.
\end{enumerate}
\end{corollary}

\begin{proof}
Insert \(\alpha=0\) and \(\alpha=2\) into
Theorem~\ref{thm:v5v22-integrated-threshold}.
\end{proof}

This is the energy analogue of the weak, intermediate, and critical
response ladder identified in the NS companion programme.  Each rung is
paid by an actual time integration.

\subsection{First-order and second-order responses}

A first-order reconstructed evolution equation
\[
 \dot u+Hu=S
\]
has the mild response
\[
 u(T)=\int_0^Te^{-(T-t)H}S(t)\,dt.
\]
For a stationary spectral source its high-energy multiplier has one
inverse power and therefore corresponds to \(m=1\).

A covariant elliptic or Euclidean wave response has symbol
\[
 (p_0^2+\Omega(\mathbf p)^2)^{-1},
\]
which decays as \(|p_0|^{-2}\) at fixed spatial momentum and corresponds
to two energy powers away from its low-energy structure.

\begin{proposition}[A second-order propagator does not make TT stress a
vector]
\label{prop:v5v22-second-order}
For the free spin-two stress density \(\alpha=2\), multiplying the
high-energy spectral amplitude by two inverse energy powers leaves a
squared norm comparable to
\[
 \int^\infty \omega^{4-4}\,d\omega,
\]
which diverges.  The response lies in
\(\mathscr H_{-r}\) for every \(r>1/2\), but not in \(\mathscr H\).
\end{proposition}

\begin{proof}
This is the case \(m=2\), \(\alpha=2\) of
Theorem~\ref{thm:v5v22-integrated-threshold}.
\end{proof}

\begin{firewall}[A propagator is not three independent integrations]
\label{fw:v5v22-propagator-count}
The second-order gravitational Green operator supplies its actual
two-power energy decay.  Iterating a formal inverse or using two external
metric legs in a covariance graph does not grant a third power to one
metric response.  A vector-level TT response needs a genuine additional
integration or an exact cancellation of the leading spectral density.
\end{firewall}

\subsection{Constraint, evolution, and observable levels}

The Hamiltonian constraint is instantaneous and has \(m=0\).  The
canonical evolution of a solved datum supplies its dynamical Duhamel
factor, while a spacetime Schwinger function retains all time variables
and need not define a sharp-time source vector.

\begin{theorem}[Three distinct regularity claims]
\label{thm:v5v22-three-claims}
For a source \(S\), the following statements are logically distinct:
\begin{enumerate}
 \item \(S(0)\Omega\in\mathscr H\), requiring finite total spectral mass;
 \item the integrated response
 \(\mathcal D_T^{(m)}S\Omega\in\mathscr H\), requiring
 \eqref{eq:v5v22-vector-threshold};
 \item the spacetime-smeared distribution \(S(f)\Omega\in\mathscr H\),
 where rapid decay of the time test can control every polynomial spectral
 tail.
\end{enumerate}
Neither item 2 nor item 3 proves item 1.
\end{theorem}

\begin{proof}
The three norms contain respectively the spectral weights
\(1\), \(|d_{m,T}(\omega)|^2\), and
\(|\widehat f(\omega)|^2\).  Lemma~\ref{lem:v5v22-D-power} gives a fixed
polynomial weight for item 2, while a smooth compactly supported or
Schwartz time test has superpolynomial Fourier decay for item 3.  Removing
either weight returns the total-mass condition of item 1.
\end{proof}

\subsection{A Duhamel constraint alternative}

One may avoid a sharp-time auxiliary metric by solving a spacetime
evolution for the constrained variable.  Abstractly, let
\begin{equation}
 z(T)
 =
 \mathcal D_T^{(m)}
 \mathcal G_{\rm sp}
 \bigl(B(z,z)+D(x,z)+S(x)\bigr),
 \label{eq:v5v22-evolutionary-constraint}
\end{equation}
where \(\mathcal G_{\rm sp}\) is the spatial inverse.

\begin{theorem}[Conditional integrated constraint contraction]
\label{thm:v5v22-integrated-contraction}
Suppose the source and nonlinear maps take values in
\(\mathscr H_{-r}\), the spatial inverse is bounded there, and
\(m\) is large enough that
\(\mathcal D_T^{(m)}\mathcal G_{\rm sp}\) maps the nonlinear source space
back to the chosen solution space.  If its norm and the bilinear constants
satisfy the V20 radius inequalities, then
\eqref{eq:v5v22-evolutionary-constraint} has a unique small solution.
\end{theorem}

\begin{proof}
Theorem~\ref{thm:v5v22-energy-gain} supplies the energy recovery.  Compose
it with the spatial inverse and apply
Theorem~\ref{thm:v5v20-form-contraction} using the norm of the composite
linear response.
\end{proof}

\begin{firewall}[Changing the equation changes the construction]
\label{fw:v5v22-change-equation}
Replacing the instantaneous ADM constraint by
\eqref{eq:v5v22-evolutionary-constraint} is legitimate only if it is
derived from the same constrained Hamiltonian or covariant field equations
and has equivalent physical solutions.  The analytic gain alone does not
prove equivalence.
\end{firewall}

\subsection{Energy-gain certificate}

\begin{definition}[Duhamel energy-gain certificate]
\label{def:v5v22-DEGC}
A DEGC consists of:
\begin{enumerate}
 \item the exact spectral degree \(\alpha\) of each complete source sector;
 \item the actual number \(m\) of energy powers supplied by the physical
 response;
 \item the residual loss
 \[
  r_{\rm min}=\max\left\{0,\alpha+\frac12-m\right\};
 \]
 \item nonlinear product estimates closing at every residual loss;
 \item derivation of the integrated equation from the constrained theory;
 \item sharp-time or OS compatibility of the resulting observables.
\end{enumerate}
\end{definition}

\begin{theorem}[Free-source DEGC ledger]
\label{thm:v5v22-ledger}
The free scalar square has zero residual loss after one integration.  The
free TT stress has residual loss greater than \(3/2\) after one
integration and greater than \(1/2\) after a second-order response; it
becomes a vector only after three energy powers.  These facts do not yet
establish a nonlinear DEGC for SM--Einstein theory.
\end{theorem}

\begin{proof}
This is Corollary~\ref{cor:v5v22-free-ladder} and
Proposition~\ref{prop:v5v22-second-order}.  The nonlinear and equivalence
items of Definition~\ref{def:v5v22-DEGC} remain unproved.
\end{proof}

V22 shows that evolution can soften the boundary obstruction but does not
automatically eliminate it for stress-tensor backreaction.  The next
question is whether conservation supplies an extra energy denominator in
the exact constrained combination without dividing by a zero-frequency
mode.

V23 will prove a frequency-local Ward division theorem.  It will split the
energy axis into a nonresonant region, where conservation can trade one
time component for a spatial component at cost \(|p_0|^{-1}\), and a
low-energy collar, which must be controlled without division.


\clearpage
\section{V23: Ward division for the Hamiltonian constraint source}
\label{sec:v5v23}

V22 counted the full spin-two stress as a source and found a loss exceeding
five halves at sharp time.  The ADM Hamiltonian constraint does not use an
arbitrary stress component: it uses the energy density, and conservation
relates that component to spatial stress.  This note performs that division
on the positive spectral support.  At fixed nonzero spatial momentum,
\(T_{0i}\) gains one inverse energy and \(T_{00}\) gains two.  The latter
therefore has at most the spectral growth of a scalar square and belongs to
every \(\mathscr H_{-r}\) with \(r>1/2\).  The threshold-energy collar and
the zero spatial mode remain separate.

\subsection{Ward identities on spectral matrix elements}

Let \(T_{\mu\nu}\) be a symmetric conserved stress tensor reconstructed on
a positive-energy Hilbert space.  Fix a spatial momentum
\(\mathbf k\ne0\).  For a generalized joint energy-momentum state with
energy \(\omega>0\), write
\[
 M_{\mu\nu}(\omega,\mathbf k)
 =
 \langle\omega,\mathbf k|T_{\mu\nu}(0)|\Omega\rangle.
\]
After local contact terms have been assigned to their counterterm owners,
conservation gives
\begin{equation}
 \omega M_{0\nu}(\omega,\mathbf k)
 =
 k_iM_{i\nu}(\omega,\mathbf k)
 \label{eq:v5v23-spectral-Ward}
\end{equation}
up to the harmless common sign fixed by Fourier convention.

\begin{lemma}[One and two Ward divisions]
\label{lem:v5v23-Ward-division}
For \(\omega>0\),
\begin{align}
 M_{0j}
 &=
 \frac{k_i}{\omega}M_{ij},
 \label{eq:v5v23-one-division}\\
 M_{00}
 &=
 \frac{k_ik_j}{\omega^2}M_{ij}.
 \label{eq:v5v23-two-divisions}
\end{align}
\end{lemma}

\begin{proof}
Equation \eqref{eq:v5v23-one-division} is
\eqref{eq:v5v23-spectral-Ward} with \(\nu=j\).  With \(\nu=0\),
symmetry gives
\[
 \omega M_{00}=k_iM_{i0}=k_iM_{0i}.
\]
Insert \eqref{eq:v5v23-one-division} to obtain
\eqref{eq:v5v23-two-divisions}.
\end{proof}

\begin{firewall}[Contact terms are not divided]
\label{fw:v5v23-contact}
The spectral Ward identity concerns separated positive-energy matrix
elements.  Polynomial contact terms in the stress two-point function do
not have this representation and may not be divided by \(\omega\).
They remain in the cosmological, Einstein, and curvature-squared local
owners identified in V17.
\end{firewall}

\subsection{Positive spectral inequalities}

Let \(d\nu_{ij,\ell m}\) be the positive matrix-valued energy measure of
the spatial stress components at fixed \(\mathbf k\).  Define the summed
spatial measure by
\[
 d\nu_{\rm sp}
 =
 \sum_{i,j}d\nu_{ij,ij}.
\]

\begin{theorem}[Ward spectral domination]
\label{thm:v5v23-spectral-domination}
The energy and momentum component measures satisfy, as positive scalar
measures,
\begin{align}
 d\nu_{0j}
 &\le
 \frac{|\mathbf k|^2}{\omega^2}\,d\nu_{\rm sp},
 \label{eq:v5v23-momentum-measure}\\
 d\nu_{00}
 &\le
 \frac{|\mathbf k|^4}{\omega^4}\,d\nu_{\rm sp}.
 \label{eq:v5v23-energy-measure}
\end{align}
\end{theorem}

\begin{proof}
At each generalized spectral state,
Cauchy--Schwarz and
\eqref{eq:v5v23-one-division} give
\[
 |M_{0j}|^2
 \le\frac{|\mathbf k|^2}{\omega^2}
      \sum_i|M_{ij}|^2.
\]
Summing over \(j\) and integrating the positive spectral resolution gives
\eqref{eq:v5v23-momentum-measure}.  Similarly,
\[
 |M_{00}|^2
 =
 \frac{|k_ik_jM_{ij}|^2}{\omega^4}
 \le
 \frac{|\mathbf k|^4}{\omega^4}
 \sum_{i,j}|M_{ij}|^2,
\]
which gives \eqref{eq:v5v23-energy-measure}.
\end{proof}

The inequalities keep the complete positive spectral form until after the
Ward combination is assembled; no cancellation is inferred from separate
component norms.

\subsection{High-energy loss of the constraint components}

Assume the spatial stress measure obeys the scale-invariant upper bound
\begin{equation}
 d\nu_{\rm sp}(\omega,\mathbf k)
 \le C_{\mathbf k}(1+\omega)^4\,d\omega
 \label{eq:v5v23-spatial-growth}
\end{equation}
above a fixed multiple of \(|\mathbf k|\).  This is the free
four-dimensional stress degree from V19--V20.

\begin{theorem}[Constraint-component energy losses]
\label{thm:v5v23-component-losses}
Under \eqref{eq:v5v23-spatial-growth}, the high-energy portions satisfy:
\begin{enumerate}
 \item \(T_{ij}\Omega\in\mathscr H_{-r}\) for \(r>5/2\);
 \item \(T_{0j}\Omega\in\mathscr H_{-r}\) for \(r>3/2\);
 \item \(T_{00}\Omega\in\mathscr H_{-r}\) for \(r>1/2\).
\end{enumerate}
\end{theorem}

\begin{proof}
The spatial density grows at most as \(\omega^4\), so the weighted integral
\[
 \int^\infty(1+\omega)^{-2r}\omega^4\,d\omega
\]
converges for \(r>5/2\).  Equations
\eqref{eq:v5v23-momentum-measure} and
\eqref{eq:v5v23-energy-measure} reduce the powers respectively to
\(\omega^2\) and \(\omega^0\) at fixed \(\mathbf k\).
The corresponding integrals converge for \(r>3/2\) and \(r>1/2\).
\end{proof}

\begin{corollary}[Poisson-solved conformal response]
\label{cor:v5v23-Poisson-response}
At fixed \(\mathbf k\ne0\), the linear Hamiltonian-constraint response
\[
 \phi(\mathbf k)=
 c\,|\mathbf k|^{-2}T_{00}(\mathbf k)
\]
belongs to \(\mathscr H_{-r}\) for every \(r>1/2\).  At high energy,
the two factors of \(\mathbf k\) in
\eqref{eq:v5v23-two-divisions} cancel the spatial inverse, leaving an
effective \(\omega^{-2}\) contraction of spatial stress.
\end{corollary}

\begin{proof}
The fixed spatial multiplier does not change the energy exponent.  More
explicitly,
\[
 |\mathbf k|^{-2}M_{00}
 =
 \frac{\widehat k_i\widehat k_j}{\omega^2}M_{ij},
\]
where \(\widehat{\mathbf k}=\mathbf k/|\mathbf k|\).
Apply Theorem~\ref{thm:v5v23-component-losses}.
\end{proof}

Thus the Hamiltonian constraint source is substantially better than the
unprojected TT stress source.  It is still a form rather than a sharp-time
Hilbert vector.

\subsection{Nonresonant and threshold collars}

Choose \(\chi_{\rm hi},\chi_{\rm lo}\in C^\infty([0,\infty))\) with
\[
 \chi_{\rm hi}+\chi_{\rm lo}=1,
 \qquad
 \chi_{\rm hi}(\omega)=0
 \text{ for }\omega\le2|\mathbf k|,
 \qquad
 \chi_{\rm hi}(\omega)=1
 \text{ for }\omega\ge3|\mathbf k|.
\]

\begin{proposition}[Threshold collar is finite at fixed nonzero momentum]
\label{prop:v5v23-collar}
If the renormalized positive spectral measure is locally finite, then the
\(\chi_{\rm lo}\) portion has finite mass at fixed
\(\mathbf k\ne0\).  All ultraviolet loss is carried by the
\(\chi_{\rm hi}\) portion, where
\eqref{eq:v5v23-one-division}--\eqref{eq:v5v23-two-divisions} give uniform
inverse-energy gains.
\end{proposition}

\begin{proof}
The low factor is supported in the compact energy interval
\([0,3|\mathbf k|]\), so local finiteness gives finite mass.  On the high
support, \(\omega\ge2|\mathbf k|>0\) and both Ward divisions are uniformly
valid.  The ultraviolet tail lies entirely there.
\end{proof}

\begin{firewall}[No uniform conclusion at the spatial zero mode]
\label{fw:v5v23-zero-k}
The collar shrinks and the Poisson inverse diverges as
\(|\mathbf k|\downarrow0\).  The spatial zero mode is the global
Hamiltonian, volume, and total-energy constraint and must be treated
before the thermodynamic limit.  Fixed-nonzero-mode estimates do not
control it.
\end{firewall}

\subsection{One dynamical integration after Ward division}

The effective high-energy degree of \(T_{00}\) is \(\alpha=0\) in the
notation of V22.

\begin{corollary}[One integration makes the constraint response a vector]
\label{cor:v5v23-one-integration}
At fixed \(\mathbf k\ne0\), if the bounds above hold and the low collar has
finite mass, one genuine Duhamel energy integration maps the
Poisson-solved \(T_{00}\) response into \(\mathscr H\).
\end{corollary}

\begin{proof}
Theorem~\ref{thm:v5v23-component-losses} gives the scalar-square threshold
\(r>1/2\), equivalently \(\alpha=0\).  Apply
Theorem~\ref{thm:v5v22-integrated-threshold} with \(m=1\).
The low compact-energy collar is unaffected.
\end{proof}

\begin{firewall}[The instantaneous constraint is still form-valued]
\label{fw:v5v23-instantaneous}
Corollary~\ref{cor:v5v23-one-integration} concerns a subsequent physical
evolution or derived integrated equation.  It does not turn the original
sharp-time ADM constraint into a vector equation.  At that stage the
source remains in \(\mathscr H_{-r}\), \(r>1/2\).
\end{firewall}

\subsection{Interacting and marked Ward requirements}

In the interacting SM theory, conservation includes gauge-covariant
equations, improvement terms, anomalies, and local contacts.  The
spectral division is available only after the exact Ward identity has
passed to the reconstructed physical Hilbert space.

\begin{definition}[Ward-divided constraint certificate]
\label{def:v5v23-WDCC}
A WDCC consists of:
\begin{enumerate}
 \item an anomaly-free, renormalized conserved stress tensor on the
 physical OS Hilbert space;
 \item positive operator-valued spectral measures through the required
 marked insertions;
 \item the exact matrix-element identities
 \eqref{eq:v5v23-one-division}--\eqref{eq:v5v23-two-divisions};
 \item polynomial high-energy bounds for the spatial stress measure;
 \item locally finite threshold collars and a separate zero-mode theorem;
 \item closed nonlinear form products at the reduced \(r>1/2\) loss;
 \item compatibility with the V14 constraint Jacobian and shell maps.
\end{enumerate}
\end{definition}

\begin{theorem}[Conditional improvement of the HFCC]
\label{thm:v5v23-HFCC-improvement}
A WDCC reduces the energy loss of the Hamiltonian-constraint matter source
from the generic stress threshold \(r>5/2\) to \(r>1/2\) at every fixed
nonzero spatial mode.  If the V20 bilinear form bounds close at that loss,
the small form-valued constraint solution exists.
\end{theorem}

\begin{proof}
The loss reduction is
Theorem~\ref{thm:v5v23-component-losses} and
Corollary~\ref{cor:v5v23-Poisson-response}.  Apply
Theorem~\ref{thm:v5v20-form-contraction} under the stated bilinear bounds.
\end{proof}

No WDCC has been proved for the interacting continuum.  V23 nevertheless
changes the correct target: the constraint sector need not pay the full
TT stress loss if conservation is used before component norms.  The open
analytic question is whether nonlinear constraint products close in the
form space with loss just above one half.

V24 will test that closure in a Gaussian Fock model.  It will derive the
energy convolution rule for products of two form-valued sources and
determine whether the \(r>1/2\) class is stable, or whether a new energy
loss appears at each nonlinear iteration.


\clearpage
\section{V24: Fock-energy convolution and failure of a fixed form loss}
\label{sec:v5v24}

V23 reduced the Hamiltonian-density source to an energy form with loss just
above one half.  The V20 contraction would close if quadratic products
mapped that form space to itself.  This note tests the claim in a Gaussian
Fock model.  Positive spectral densities convolve under the highest
particle-number part of a normal-ordered product.  Their boundary losses
add: two sources at the one-half threshold produce a source at the
one-derivative threshold, and the \(n\)-fold product requires loss
\(n/2\).  Therefore no fixed negative-energy space contains the full
quadratic iteration.  A particle-number graded loss restores an algebra
and becomes the next candidate norm.

\subsection{Convolution of positive energy measures}

Let two source vectors have absolutely continuous positive energy measures
\begin{equation}
 d\nu_f(\omega)=f(\omega)\,d\omega,
 \qquad
 d\nu_g(\omega)=g(\omega)\,d\omega.
 \label{eq:v5v24-two-measures}
\end{equation}
For independent particle sectors, the tensor-product Hamiltonian is
\[
 H_{12}=H_1\otimes I+I\otimes H_2.
\]

\begin{lemma}[Energy measure of a tensor product]
\label{lem:v5v24-convolution}
The energy density of the tensor-product vector is
\begin{equation}
 (f*g)(\Omega)
 =
 \int_0^\Omega f(\omega)g(\Omega-\omega)\,d\omega.
 \label{eq:v5v24-convolution}
\end{equation}
\end{lemma}

\begin{proof}
The joint spectral measure is the product
\(d\nu_f(\omega_1)d\nu_g(\omega_2)\).  Pushing it forward under
\((\omega_1,\omega_2)\mapsto\omega_1+\omega_2\) gives
\eqref{eq:v5v24-convolution}.
\end{proof}

\begin{theorem}[Exact power convolution]
\label{thm:v5v24-power-convolution}
If
\[
 f(\omega)=c_f\omega^p,
 \qquad
 g(\omega)=c_g\omega^q,
 \qquad
 p,q>-1,
\]
then
\begin{equation}
 (f*g)(\Omega)
 =
 c_fc_g\,\mathrm B(p+1,q+1)\,
 \Omega^{p+q+1},
 \label{eq:v5v24-beta-convolution}
\end{equation}
where \(\mathrm B\) is the beta function.
\end{theorem}

\begin{proof}
Set \(\omega=\Omega u\) in
\eqref{eq:v5v24-convolution}:
\[
 (f*g)(\Omega)
 =
 c_fc_g\Omega^{p+q+1}
 \int_0^1u^p(1-u)^q\,du.
\]
The remaining integral is \(\mathrm B(p+1,q+1)\).
\end{proof}

\subsection{Addition of negative-energy losses}

A density behaving as \(\omega^p\,d\omega\) at infinity belongs to
\(\mathscr H_{-r}\) precisely when
\[
 \int^\infty(1+\omega)^{-2r}\omega^p\,d\omega<\infty,
\]
or
\begin{equation}
 r>\frac{p+1}{2}.
 \label{eq:v5v24-single-threshold}
\end{equation}

\begin{corollary}[Losses add under tensor products]
\label{cor:v5v24-loss-adds}
For the power densities of
Theorem~\ref{thm:v5v24-power-convolution}, the tensor-product source lies
in \(\mathscr H_{-r}\) exactly when
\begin{equation}
 r>\frac{p+q+2}{2}
 =
 \frac{p+1}{2}+\frac{q+1}{2}.
 \label{eq:v5v24-added-loss}
\end{equation}
\end{corollary}

\begin{proof}
The convolution exponent is \(p+q+1\).  Insert it into
\eqref{eq:v5v24-single-threshold}.
\end{proof}

\begin{theorem}[The one-half class is not a quadratic algebra]
\label{thm:v5v24-half-not-algebra}
Let each of two independent source sectors have a density bounded above and
below by positive constants at high energy.  Each source lies in
\(\mathscr H_{-r}\) for every \(r>1/2\).  Their tensor-product source lies
in \(\mathscr H_{-r}\) only for \(r>1\).  Hence no estimate
\[
 \|uv\|_{-r}\le C_r\|u\|_{-r}\|v\|_{-r}
 \label{eq:v5v24-false-algebra}
\]
can hold universally at a fixed \(r\in(1/2,1]\).
\end{theorem}

\begin{proof}
Take \(p=q=0\) in
Corollary~\ref{cor:v5v24-loss-adds}.  The tensor product has density
proportional to \(\Omega\) and threshold \(r>1\).  If
\eqref{eq:v5v24-false-algebra} held, the product of two admissible
one-half-loss vectors would remain in the same space, contradicting that
threshold.
\end{proof}

\subsection{Normal ordering does not remove the highest chaos}

In a bosonic Gaussian Fock space, the product of two Wick polynomials
contains contractions of lower particle number and a no-contraction term
whose vacuum vector is the symmetrized tensor product of the two highest
chaos vectors.

\begin{proposition}[Highest-particle positive packet]
\label{prop:v5v24-highest-chaos}
Let \(U\) and \(V\) be Wick polynomials whose highest vacuum sectors are
nonzero vectors \(u_n\) and \(v_m\).  The \((n+m)\)-particle component of
\({:}UV{:}\Omega\) is
\begin{equation}
 \operatorname{Sym}(u_n\otimes v_m).
 \label{eq:v5v24-highest-chaos}
\end{equation}
It is orthogonal to every contraction term.  For independent species, its
energy measure is exactly the convolution of the two positive measures.
\end{proposition}

\begin{proof}
Wick multiplication partitions pairings into contractions and the
no-contraction term.  Every contraction lowers particle number by two,
while the no-contraction term is the symmetrized tensor product.
Different particle sectors are orthogonal in Fock space.  For independent
species the symmetrization introduces no destructive identification, and
Lemma~\ref{lem:v5v24-convolution} applies.
\end{proof}

\begin{corollary}[Lower-chaos counterterms cannot cancel the loss]
\label{cor:v5v24-no-lower-cancel}
A local counterterm or Wick contraction supported in lower particle number
cannot cancel the positive norm of
\eqref{eq:v5v24-highest-chaos}.  Any uniform nonlinear estimate must
control that sector directly.
\end{corollary}

\begin{proof}
Orthogonal particle sectors have squared norms that add.  Cancellation
between them is impossible.
\end{proof}

\subsection{Iteration of a quadratic constraint response}

Suppose the Ward-divided first response \(z_1\) has high-energy density of
degree zero, as in V23.  The highest chaos of the quadratic term
\(z_1^2\) has degree one.  Iterating a quadratic fixed-point equation
generates products of arbitrarily many copies of the initial source.

\begin{theorem}[Unbounded loss along the perturbative constraint tree]
\label{thm:v5v24-tree-loss}
The \(n\)-fold no-contraction product of degree-zero source packets has
energy density proportional to
\begin{equation}
 \Omega^{n-1}\,d\Omega
 \label{eq:v5v24-nfold-density}
\end{equation}
and belongs to \(\mathscr H_{-r}\) exactly when
\begin{equation}
 r>\frac n2.
 \label{eq:v5v24-nfold-loss}
\end{equation}
Therefore an all-order quadratic constraint expansion is not contained in
any fixed \(\mathscr H_{-r}\) by energy weights alone.
\end{theorem}

\begin{proof}
Induct on \(n\).  The density is constant at \(n=1\).
Convolution of \(\Omega^{n-1}\) with a constant gives a positive multiple
of \(\Omega^n\) by
Theorem~\ref{thm:v5v24-power-convolution}.  This proves
\eqref{eq:v5v24-nfold-density}.  Equation
\eqref{eq:v5v24-single-threshold} gives
\eqref{eq:v5v24-nfold-loss}.  For every fixed \(r\), choose \(n\ge2r\).
\end{proof}

\begin{firewall}[Small coupling does not change membership]
\label{fw:v5v24-small-coupling}
Multiplying the \(n\)-particle packet by a small nonzero coupling changes
its norm only when that norm is finite.  It cannot place a vector with
infinite \(\mathscr H_{-r}\) norm into the space.  Coupling smallness must
be combined with a graded norm or genuine energy smoothing.
\end{firewall}

This proves that the fixed-space bilinear hypothesis in V20 fails in the
unstructured Gaussian tensor-product model.

\subsection{A particle-number graded rescue space}

Let \(\mathscr F=\bigoplus_{n\ge0}\mathscr H^{(n)}\) be a Fock-type direct
sum with
\[
 H^{(n)}=\sum_{j=1}^nH_j.
\]
For \(b\ge0\), \(\eta>0\), and \(\sigma\ge0\), define
\begin{equation}
 \|\Psi\|_{\eta,\sigma,b}
 =
 \sum_{n\ge0}
 \frac{\eta^n}{(n!)^\sigma}
 \|(1+H^{(n)})^{-bn}\Psi_n\|.
 \label{eq:v5v24-graded-Fock}
\end{equation}

\begin{lemma}[Additive energy weight]
\label{lem:v5v24-additive-weight}
For \(E_1,E_2\ge0\) and integers \(n,m\ge0\),
\begin{equation}
 (1+E_1+E_2)^{-b(n+m)}
 \le
 (1+E_1)^{-bn}(1+E_2)^{-bm}.
 \label{eq:v5v24-energy-weight}
\end{equation}
\end{lemma}

\begin{proof}
Both \(1+E_1\) and \(1+E_2\) are at most
\(1+E_1+E_2\).  Raise the first inequality to \(bn\), the second to
\(bm\), multiply, and invert.
\end{proof}

Let \(\diamond\) denote a graded product whose
\((n+m)\)-particle component is a contraction of
\(\Psi_n\otimes\Phi_m\) with operator norm at most \(C_\diamond\), and
sum over \(n+m=N\).

\begin{theorem}[Graded energy-loss algebra]
\label{thm:v5v24-graded-algebra}
Under the stated product bound,
\begin{equation}
 \|\Psi\diamond\Phi\|_{\eta,\sigma,b}
 \le
 C_\diamond
 \|\Psi\|_{\eta,\sigma,b}
 \|\Phi\|_{\eta,\sigma,b}.
 \label{eq:v5v24-graded-algebra}
\end{equation}
\end{theorem}

\begin{proof}
Lemma~\ref{lem:v5v24-additive-weight} bounds the energy norm of every
tensor product by the product of the two weighted energy norms.
Furthermore,
\[
 \frac{\eta^{n+m}}{((n+m)!)^\sigma}
 \le
 \frac{\eta^n}{(n!)^\sigma}
 \frac{\eta^m}{(m!)^\sigma}.
\]
Sum the nonnegative bounds over \(n,m\), exactly as in the V6 Gevrey word
algebra.
\end{proof}

\begin{corollary}[Minimal loss slope for degree-zero packets]
\label{cor:v5v24-minimal-slope}
For the \(n\)-fold degree-zero source family of
Theorem~\ref{thm:v5v24-tree-loss}, any \(b>1/2\) places every
\(n\)-particle packet in its assigned component of
\eqref{eq:v5v24-graded-Fock}.  The boundary slope \(b=1/2\) remains
critical.
\end{corollary}

\begin{proof}
The assigned loss is \(bn\), while the exact membership condition is
strictly greater than \(n/2\).
\end{proof}

\subsection{What the graded rescue still requires}

The norm \eqref{eq:v5v24-graded-Fock} weakens with particle number.  To use
it in a continuum construction one must prove that physical observables,
Ward operations, shell maps, and OS reconstruction act continuously on
the resulting common domain.

\begin{definition}[Graded form constraint certificate]
\label{def:v5v24-GFCC}
A GFCC consists of:
\begin{enumerate}
 \item a Fock or chaos decomposition compatible with the interacting shell
 map;
 \item a loss slope \(b>1/2\) for the Ward-divided constraint sector;
 \item factorial weights controlling graph and partition growth;
 \item bounded complete-source products as in
 \eqref{eq:v5v24-graded-algebra};
 \item shell contraction and compact radius gaps compatible with V6--V7
 and V21;
 \item a common physical observable domain, OS positivity, and removal of
 the auxiliary constraint variables.
\end{enumerate}
\end{definition}

\begin{theorem}[Gaussian nonlinear status]
\label{thm:v5v24-status}
The naive fixed form space with \(r\) just above \(1/2\) is not stable under
quadratic Gaussian constraint iteration.  The graded space
\eqref{eq:v5v24-graded-Fock} is an algebra under the declared bounded
graded contractions and contains the highest-chaos source packets for
every \(b>1/2\).  No interacting GFCC is presently proved.
\end{theorem}

\begin{proof}
Failure of the fixed space is
Theorem~\ref{thm:v5v24-tree-loss}.  The positive graded statements are
Theorem~\ref{thm:v5v24-graded-algebra} and
Corollary~\ref{cor:v5v24-minimal-slope}.  The remaining GFCC items are
additional hypotheses.
\end{proof}

V24 identifies the correct nonlinear response space at the Gaussian
boundary: energy loss must grow with particle number, just as curvature
grade and graph order require Gevrey weights.  This links the metric
constraint problem back to the V5--V7 operator-tower construction.

V25 will begin with the compulsory companion audit.  It will then combine
the curvature-grade and particle-number weights into one tri-graded shell
space and prove a compact algebra embedding with explicit radius gaps.

\clearpage
\section{V25: A signed tri-graded compact algebra for the coupled shell map}
\label{sec:v5v25}

\subsection{Purpose and compulsory companion audit}

V24 showed that the Ward-divided gravitational constraint source does not
remain in a fixed energy form space under unrestricted quadratic iteration:
the required negative energy order grows linearly with particle number.  The
present note combines that particle-number loss with the curvature or local
operator grade used in V5--V7 and with the spatial regularity scale used in
V21.  A second issue is forced by the companion programmes.  A norm which
takes absolute values diagram by diagram can erase an exact Ward or cumulant
cancellation before it is used.  We therefore keep a declared low-label
packet assembled as a signed vector and take its norm only after assembly.

At this fifth-version checkpoint the newest live companion files were read
before choosing the construction:
\begin{align*}
 &\texttt{Renewal\_Yang\_Mills\_Mass\_Gap\_Research\_Notes\_V35.tex},\\
 &\qquad
 \operatorname{SHA256}=
 \texttt{37472D4EF58415B9F6CD310614691F53F76B6336175CDC4411B2D8D081D5716C},\\
 &\texttt{Renewal\_NS\_Stability\_Research\_Notes\_V31.tex},\\
 &\qquad
 \operatorname{SHA256}=
 \texttt{F1121B62238AD5491BB7CF03635EA9934942EDF573ED8FAAA9EE79E1D38A6696}.
\end{align*}
The YM notes close several one-strong sectors and reduce the remaining
single-coordinate terminal to a low-output packet.  In its positive
projection the available gain is one power short, so any further gain must
come from a complete signed local moment--cumulant symbol evaluated before
the triangle inequality.  The NS notes retain an exact marked first moment
through heat formation and likewise isolate a signed balanced-sector
nullspace.  These are structural warnings, not estimates for the
SM--Einstein map: the present note imports only the rule that exact signed
assemblies must precede absolute summation.

\begin{firewall}[No transfer of companion conclusions]
\label{fw:v5v25-companion}
Neither companion audit proves a Standard Model Ward identity, an Einstein
constraint estimate, or a mass gap.  Every such assertion below is either
proved from stated hypotheses or explicitly retained as an open certificate.
\end{firewall}

\subsection{The tri-graded space}

Let (p\in\mathbb N_0) be curvature or local-operator grade and let
(m\in\mathbb N_0) be particle or highest-chaos number.  For every pair
((p,m)), let (X_{p,m}) be a Hilbert space.  Let (J_{p,m}\ge1) measure
spatial derivatives and let (H_{p,m}\ge0) be the total positive energy.
We assume that their spectral measures commute.  This holds, for example,
when both are functions of a common joint spectral resolution.  For
(s\ge0) and (b\ge0), set
\begin{equation}
 \|u\|_{p,m;s,b}
 :=\bigl\|J_{p,m}^{s}(1+H_{p,m})^{-bm}u\bigr\|_{X_{p,m}}.
 \label{eq:v5v25-block-norm}
\end{equation}
The factor (m) in the energy loss is the conclusion forced by V24.

Fix exponents (\alpha,\beta\ge0) and radii (\rho,\eta>0).  For a
double sequence (u=(u_{p,m})), define
\begin{equation}
 \|u\|_{\rho,\eta;s,b}^{\alpha,\beta}
 :=\sum_{p,m\ge0}
 \frac{\rho^p}{(p!)^\alpha}
 \frac{\eta^m}{(m!)^\beta}
 \|u_{p,m}\|_{p,m;s,b}.
 \label{eq:v5v25-trigraded-norm}
\end{equation}
Its completion is denoted
(\mathbb B_{\rho,\eta;s,b}^{\alpha,\beta}).  It is tri-graded in the
sense that (p), (m), and the spectral regularity measured by (J) are
independent coordinates; the energy loss (bm) is tied to the particle
coordinate rather than treated as a fourth free grade.

\begin{lemma}[Monotonicity of the scale]
\label{lem:v5v25-monotonicity}
Suppose
\begin{equation}
 \rho_+\ge\rho_-,\qquad \eta_+\ge\eta_-,\qquad
 s_+\ge s_-,\qquad b_+\le b_- .
 \label{eq:v5v25-parameter-order}
\end{equation}
Then the identity on finite sequences extends with norm at most one from
(\mathbb B_{\rho_+,\eta_+;s_+,b_+}^{\alpha,\beta}) to
(\mathbb B_{\rho_-,\eta_-;s_-,b_-}^{\alpha,\beta}).
\end{lemma}

\begin{proof}
The radius quotients are at most one.  By joint spectral calculus,
\[
 J^{s_-}(1+H)^{-b_-m}
 =J^{-(s_+-s_-)}(1+H)^{-(b_--b_+)m}
  J^{s_+}(1+H)^{-b_+m}.
\]
Both prefactors have operator norm at most one.  Apply this blockwise and
sum with the two radius quotients.
\end{proof}

The direction of the (b)-inequality deserves emphasis: a larger loss
slope gives a weaker topology.  Thus a useful compact radius gap has
\(1/2<b_+<b_-), so that the stronger space still contains the Gaussian
highest-chaos packets identified in V24.

\subsection{A product theorem with all three grades visible}

Let (\star_{p,m;q,n}:X_{p,m}\times X_{q,n}\to X_{p+q,m+n}) be bilinear,
and put
\begin{equation}
 (u\star v)_{r,k}
 =\sum_{p+q=r\atop m+n=k}u_{p,m}\star_{p,m;q,n}v_{q,n}.
 \label{eq:v5v25-Cauchy-product}
\end{equation}
We isolate the analytic input in a block estimate.

\begin{definition}[Uniform tri-graded product certificate]
\label{def:v5v25-UTPC}
A UTPC at ((s,b)) is a constant (C_\star<\infty) such that
\begin{equation}
 \|x\star_{p,m;q,n}y\|_{p+q,m+n;s,b}
 \le C_\star
 \|x\|_{p,m;s,b}\|y\|_{q,n;s,b}
 \label{eq:v5v25-block-product}
\end{equation}
for every four grades and all vectors in the indicated block domains.
\end{definition}

For pointwise spatial multiplication, the usual sufficient condition is
(s>d/2), together with compatibility of (J) with the Sobolev scale.
For the energy factor, Lemma~\ref{lem:v5v24-additive-weight} supplies the
needed estimate for no-contraction tensor products.  Contraction terms and
Ward projections must still be verified as part of the UTPC.

\begin{theorem}[Tri-graded Banach algebra]
\label{thm:v5v25-algebra}
If the UTPC holds, then
\begin{equation}
 \|u\star v\|_{\rho,\eta;s,b}^{\alpha,\beta}
 \le C_\star
 \|u\|_{\rho,\eta;s,b}^{\alpha,\beta}
 \|v\|_{\rho,\eta;s,b}^{\alpha,\beta}.
 \label{eq:v5v25-algebra}
\end{equation}
Consequently (\mathbb B_{\rho,\eta;s,b}^{\alpha,\beta}) is a Banach
algebra after adjoining the product continuously.
\end{theorem}

\begin{proof}
For (r=p+q) and (k=m+n),
\[
 \frac{\rho^{p+q}}{((p+q)!)^\alpha}
 \frac{\eta^{m+n}}{((m+n)!)^\beta}
 \le
 \frac{\rho^p}{(p!)^\alpha}\frac{\eta^m}{(m!)^\beta}
 \frac{\rho^q}{(q!)^\alpha}\frac{\eta^n}{(n!)^\beta},
\]
because ((p+q)!\ge p!q!) and ((m+n)!\ge m!n!).
Apply \eqref{eq:v5v25-block-product}, use the triangle inequality only
after each declared block product has been formed, and sum over
((p,m,q,n)).  Tonelli's theorem applies to the resulting nonnegative
series and gives the product of the two norms.  Completeness follows from
the weighted \(\ell^1\)-sum of complete block spaces.
\end{proof}

\begin{remark}[Where combinatorial constants go]
If a chosen normalization produces a factor
(C_0^{p+q+m+n}) in \eqref{eq:v5v25-block-product}, replace
((\rho,\eta)) on the output by ((\rho/C_0,\eta/C_0)).  If the factor is
factorial rather than exponential, the exponents \(\alpha\) and \(\beta\)
must dominate it.  This is an estimate to be checked for the exact graph
normalization; it cannot be declared away by small coupling.
\end{remark}

\subsection{Why the protected symbol must be assembled first}

The word ``signed'' refers to the order of operations, not to an indefinite
physical inner product.  Let a protected block be presented by elementary
diagrammatic contributions (a_1,\ldots,a_N\) in a normed space (Y).
Its physical value is (A=\sum_j a_j).

\begin{lemma}[Triangle inequality loses exact cancellation]
\label{lem:v5v25-cancellation-loss}
There is no constant (C) such that
\begin{equation}
 \sum_{j=1}^N\|a_j\|_Y\le C\Bigl\|\sum_{j=1}^Na_j\Bigr\|_Y
 \label{eq:v5v25-false-reverse}
\end{equation}
for all finite signed families, even when (N=2).  Therefore an estimate
made after replacing every contribution by its norm cannot prove a Ward or
cumulant cancellation of the assembled symbol.
\end{lemma}

\begin{proof}
Take a nonzero (a\in Y), (a_1=a), and (a_2=-a).  The left side of
\eqref{eq:v5v25-false-reverse} is (2\|a\|_Y>0), while the right side is
zero for every finite (C).
\end{proof}

For every ((p,m)), choose a finite-rank orthogonal projection
\(P_{p,m}\) containing the predeclared low-label Ward and cumulant packet,
and put (Q_{p,m}=1-P_{p,m}).  Let (\|\cdot\|_{A_{p,m}}) be an absolute
coefficient norm on the tail.  Define the hybrid block norm
\begin{equation}
 \|u\|^{\mathrm{hyb}}_{p,m;s,b}
 :=
 \bigl\|P_{p,m}J^s(1+H)^{-bm}u\bigr\|_{X_{p,m}}
 +
 \bigl\|Q_{p,m}J^s(1+H)^{-bm}u\bigr\|_{A_{p,m}}.
 \label{eq:v5v25-hybrid-block}
\end{equation}
All diagrammatic or counterterm contributions to (P_{p,m}u) are summed
as vectors before the first norm in \eqref{eq:v5v25-hybrid-block} is
taken.  Replace the block norm in \eqref{eq:v5v25-trigraded-norm} by
\eqref{eq:v5v25-hybrid-block} to obtain
(\mathbb B^{\mathrm{hyb}}_{\rho,\eta;s,b}).

\begin{theorem}[Hybrid algebra criterion]
\label{thm:v5v25-hybrid-algebra}
Suppose, for every pair of input blocks, the fully assembled product obeys
\begin{align}
 \|P(x\star y)\|_X
 &\le C_P(\|Px\|_X+\|Qx\|_A)
              (\|Py\|_X+\|Qy\|_A),
 \label{eq:v5v25-P-bound}\\
 \|Q(x\star y)\|_A
 &\le C_Q(\|Px\|_X+\|Qx\|_A)
              (\|Py\|_X+\|Qy\|_A),
 \label{eq:v5v25-Q-bound}
\end{align}
with the same spectral weights as in
\eqref{eq:v5v25-hybrid-block}, uniformly in the grades.  Then the hybrid
space is a Banach algebra with constant (C_P+C_Q).
\end{theorem}

\begin{proof}
Add \eqref{eq:v5v25-P-bound} and \eqref{eq:v5v25-Q-bound}.  This gives the
block estimate of Definition~\ref{def:v5v25-UTPC} for the hybrid norm.
The proof of Theorem~\ref{thm:v5v25-algebra} then applies verbatim.  The
ordering is essential: the left side of \eqref{eq:v5v25-P-bound} is the
norm of the assembled protected projection, not the sum of norms of its
diagrammatic summands.
\end{proof}

\begin{firewall}[Growing protected windows]
\label{fw:v5v25-growing-window}
Finite rank at each shell is not a uniform compactness estimate if the
rank grows with the shell.  A window of size (O(\sqrt{\log d})), as in
the audited YM terminal, requires a separate uniform operator bound and a
tail estimate.  The theorem above covers the fixed SM protected sector
only after its ranks and constants are controlled uniformly.
\end{firewall}

\subsection{Compact embedding with explicit radius gaps}

We next state the compactness result needed for removal of a finite shell
truncation.  Assume that (J_{p,m}) has compact resolvent for each fixed
((p,m)), and that the protected range lies in its form domain.  Let
(R_{p,m}(L)) be the joint spectral projection of (J_{p,m}) onto
([1,L]).  It has finite rank.

Choose strict gaps
\begin{equation}
 \rho_+>\rho_->0,qquad \eta_+>\eta_->0,qquad
 s_+>s_-,\qquad \frac12<b_+<b_- .
 \label{eq:v5v25-strict-gaps}
\end{equation}
Let (\iota_{+-}) denote the identity from the plus space to the minus
space, with the same factorial exponents.

\begin{theorem}[Tri-graded compact embedding]
\label{thm:v5v25-compact}
Under the compact-resolvent assumption, (\iota_{+-}) is compact.  If
(\Pi_{P,M,L}) retains (p\le P), (m\le M), and the range of
(R_{p,m}(L)), then
\begin{equation}
 \|(1-\Pi_{P,M,L})\iota_{+-}\|
 \le
 \left(\frac{\rho_-}{\rho_+}\right)^{P+1}
 +\left(\frac{\eta_-}{\eta_+}\right)^{M+1}
 +L^{-(s_+-s_-)}.
 \label{eq:v5v25-compact-tail}
\end{equation}
The same conclusion holds for the hybrid spaces if the protected
projections have finite rank in every retained block and the constants in
the block norm equivalences are uniform on retained rectangles.
\end{theorem}

\begin{proof}
Write (w_{p,m}^\pm=\rho_\pm^p\eta_\pm^m/
((p!)^\alpha(m!)^\beta)).  For (p>P),
(w_{p,m}^-/w_{p,m}^+\le(\rho_-/\rho_+)^{P+1}); summing the plus block
norms gives the first term.  On (p\le P) but (m>M), the analogous
ratio is at most ((\eta_-/\eta_+)^{M+1}).  On the remaining rectangle,
joint spectral calculus gives, on (1-R_{p,m}(L)),
\[
 \|J^{s_-}(1+H)^{-b_-m}u\|
 \le L^{-(s_+-s_-)}
 \|J^{s_+}(1+H)^{-b_+m}u\|,
\]
where the extra energy multiplier has norm at most one.  This proves
\eqref{eq:v5v25-compact-tail}; overlaps only improve the estimate.

The range of (\Pi_{P,M,L}) is a finite direct sum of finite-dimensional
spaces, hence is finite dimensional.  The right side of
\eqref{eq:v5v25-compact-tail} tends to zero successively as
(P,M,L\to\infty).  Thus (\iota_{+-}) is an operator-norm limit of
finite-rank maps and is compact.  In the hybrid case the protected part is
already finite dimensional on each retained block, and uniform equivalence
transfers the same estimate.
\end{proof}

\begin{firewall}[Infinite-volume translations]
\label{fw:v5v25-translations}
On (\mathbb R^d), the embedding (H^{s_+}\hookrightarrow H^{s_-}) is
not globally compact: translations of one fixed nonzero function give a
counterexample.  The compact-resolvent hypothesis must therefore arise
from a finite-volume shell construction, localization with a proved
tightness estimate, or quotienting the translation orbit.  It is not an
automatic consequence of a Sobolev gap.
\end{firewall}

\subsection{A compact contraction consequence}

The preceding theorem isolates the functional-analytic part of the shell
limit.  It does not supply the dynamical estimates, but it shows exactly
which ones suffice.

\begin{theorem}[Strong-image contraction theorem]
\label{thm:v5v25-strong-image}
Let \(B_-\) be a closed ball in the minus hybrid space.  Suppose a shell
map (\mathcal T:B_-\to B_-) satisfies
\begin{equation}
 \|\mathcal T u-\mathcal T v\|_-
 \le q\|u-v\|_-,\qquad 0\le q<1,
 \label{eq:v5v25-contraction}
\end{equation}
and
\begin{equation}
 \sup_{u\in B_-}\|\mathcal T u\|_+\le M_+<\infty.
 \label{eq:v5v25-strong-image-bound}
\end{equation}
Then (\mathcal T) has a unique fixed point (u_*\in B_-), its image is
relatively compact in the minus space, and
\begin{equation}
 \sup_{u\in B_-}
 \|(1-\Pi_{P,M,L})\mathcal T u\|_-
 \le M_+\varepsilon_{P,M,L},
 \label{eq:v5v25-uniform-truncation}
\end{equation}
where (\varepsilon_{P,M,L}) is the right side of
\eqref{eq:v5v25-compact-tail}.
\end{theorem}

\begin{proof}
The Banach fixed-point theorem gives existence and uniqueness in the
complete closed ball.  Equation \eqref{eq:v5v25-strong-image-bound} puts
the image in a bounded subset of the plus space, whose image under the
compact embedding is relatively compact.  Apply
\eqref{eq:v5v25-compact-tail} to each (\mathcal T u) and take the
supremum to obtain \eqref{eq:v5v25-uniform-truncation}.
\end{proof}

\begin{corollary}[Convergence of consistently truncated fixed points]
\label{cor:v5v25-truncated-fixed-points}
Assume in addition that
(\mathcal T_N=\Pi_N\mathcal T\Pi_N) maps (B_-\cap\operatorname{ran}\Pi_N)
to itself with the same contraction constant (q), and that
\begin{equation}
 \delta_N:=\sup_{u\in B_-}
 \|\mathcal T_Nu-\mathcal Tu\|_-\longrightarrow0.
 \label{eq:v5v25-consistency}
\end{equation}
If (u_N) and (u_*) are the respective fixed points, then
\begin{equation}
 \|u_N-u_*\|_-\le\frac{\delta_N}{1-q}.
 \label{eq:v5v25-fixed-point-error}
\end{equation}
\end{corollary}

\begin{proof}
Using both fixed-point equations,
\[
 \|u_N-u_*\|_-
 \le\|\mathcal T_Nu_N-\mathcal Tu_N\|_-
    +\|\mathcal Tu_N-\mathcal Tu_*\|_-
 \le\delta_N+q\|u_N-u_*\|_-.
\]
Rearrange.
\end{proof}

\subsection{Application ledger and next obstruction}

The new space resolves one precise obstruction.  V24's (n/2) energy
loss is compatible with every particle sector once (b_+>1/2), while
Theorem~\ref{thm:v5v25-algebra} permits products without returning to a
fixed-loss space.  Theorem~\ref{thm:v5v25-compact} simultaneously controls
large curvature grade, large particle number, and large spatial frequency.
The hybrid norm prevents a declared finite Ward packet from being split by
the triangle inequality.

It does not yet establish the following dynamical inputs:
\begin{enumerate}
 \item the exact SM--Einstein shell product bounds
 \eqref{eq:v5v25-P-bound}--\eqref{eq:v5v25-Q-bound}, including ghosts,
 gauge fixing, and the improvement terms in the stress tensor;
 \item a uniform strong-image bound \eqref{eq:v5v25-strong-image-bound}
 with (1/2<b_+<b_-) and radii independent of the ultraviolet cutoff;
 \item an assembled protected Ward-symbol identity before absolute graph
 summation;
 \item infinite-volume tightness, OS positivity, and recovery of the
 physical constraint surface.
\end{enumerate}

\begin{theorem}[V25 status]
\label{thm:v5v25-status}
Subject to the explicit UTPC and compact-resolvent hypotheses, the
tri-graded hybrid space is a Banach algebra, its strict-radius embedding is
compact with the quantitative error
\eqref{eq:v5v25-compact-tail}, and a shell contraction with a bounded strong
image has cutoff-independent fixed points satisfying
\eqref{eq:v5v25-fixed-point-error}.  No exact interacting SM--Einstein shell
map has yet been shown to satisfy those hypotheses.
\end{theorem}

\begin{proof}
Combine Theorems~\ref{thm:v5v25-hybrid-algebra},
\ref{thm:v5v25-compact}, and \ref{thm:v5v25-strong-image} with
Corollary~\ref{cor:v5v25-truncated-fixed-points}.  The last sentence records
the four unproved inputs in the preceding ledger.
\end{proof}

V26 will attack the first and third inputs at the abstract operator level.
It will derive a head--protected--tail Schur estimate for one nonlinear
shell step, keeping the complete Ward packet assembled, and will identify
the exact smallness condition needed for a contraction on the hybrid
tri-graded ball.

\clearpage
\section{V26: Head--protected--tail reduction of one nonlinear shell step}
\label{sec:v5v26}

\subsection{The obstruction left by V25}

The compact algebra of V25 is useful only after the exact shell map is
shown to preserve its signed protected packet and to contract after the
finite relevant parameters have been fixed.  This note separates these
tasks.  It gives a quantitative block criterion which allows large-looking
individual cross estimates provided their complete Schur feedback is small.
It also formulates the Ward identity on the assembled shell map, where the
cancellation is meaningful.

Fix either the plus or minus hybrid tri-graded space of V25 and write it as
a topological direct sum
\begin{equation}
 \mathbb B=\mathbb H\oplus\mathbb P\oplus\mathbb T .
 \label{eq:v5v26-HPT}
\end{equation}
Here \(\mathbb H\) is the finite-dimensional head of relevant and marginal
renormalization data, \(\mathbb P\) is the assembled protected Ward packet,
and \(\mathbb T\) is the irrelevant tail.  The norm on \(\mathbb P\) is
taken after all graphs contributing to the same Ward symbol have been
summed, as required by Lemma~\ref{lem:v5v25-cancellation-loss}.

\subsection{A nonnegative Lipschitz matrix}

Let \(\mathcal T=(\mathcal T_H,\mathcal T_P,\mathcal T_T)\) be a nonlinear
shell step on a convex set \(D\subset\mathbb B\).  Suppose there are
nonnegative numbers \(L_{ij}\), \(i,j\in\{H,P,T\}\), such that
\begin{equation}
 \|\mathcal T_i(x)-\mathcal T_i(y)\|_i
 \le\sum_{j\in\{H,P,T\}}L_{ij}\|x_j-y_j\|_j
 \quad (x,y\in D).
 \label{eq:v5v26-block-Lipschitz}
\end{equation}
Write
\begin{equation}
 L=
 \begin{pmatrix}
  A&u\\
  v^{\mathsf T}&\kappa
 \end{pmatrix},
 \qquad
 A=
 \begin{pmatrix}
 L_{HH}&L_{HP}\\
 L_{PH}&L_{PP}
 \end{pmatrix},
 \quad
 u=\binom{L_{HT}}{L_{PT}},
 \quad
 v=\binom{L_{TH}}{L_{TP}},
 \quad
 \kappa=L_{TT}.
 \label{eq:v5v26-Schur-data}
\end{equation}
Thus \(A\) is the coupled head--protected block and \(\kappa\) is the
direct tail contraction.

\begin{lemma}[Weighted max norm from a positive supersolution]
\label{lem:v5v26-weighted-max}
If a vector \(w=(w_H,w_P,w_T)^{\mathsf T}\) has strictly positive entries
and
\begin{equation}
 Lw\le q w
 \quad\text{componentwise for some }q<1,
 \label{eq:v5v26-supersolution}
\end{equation}
then \(\mathcal T\) is \(q\)-Lipschitz for
\begin{equation}
 \|x\|_w=\max_{i\in\{H,P,T\}}\frac{\|x_i\|_i}{w_i}.
 \label{eq:v5v26-weighted-max-norm}
\end{equation}
\end{lemma}

\begin{proof}
Put \(d=\|x-y\|_w\).  Then
\(\|x_j-y_j\|_j\le w_jd\).  By
\eqref{eq:v5v26-block-Lipschitz},
\[
 \|\mathcal T_i(x)-\mathcal T_i(y)\|_i
 \le (Lw)_i d\le q w_i d.
\]
Divide by \(w_i\) and take the maximum.
\end{proof}

\begin{theorem}[Head--protected--tail Schur criterion]
\label{thm:v5v26-Schur}
Assume \(A\ge0\), \(u,v\ge0\), \(\kappa\ge0\), and
\begin{equation}
 \rho(A)<1,\qquad
 \kappa+v^{\mathsf T}(I-A)^{-1}u<1.
 \label{eq:v5v26-Schur-condition}
\end{equation}
Then there are \(w_H,w_P,w_T>0\) and \(q<1\) satisfying
\eqref{eq:v5v26-supersolution}.  Consequently every map satisfying
\eqref{eq:v5v26-block-Lipschitz} is a strict contraction in an equivalent
weighted block norm.
\end{theorem}

\begin{proof}
Since \(A\ge0\) and \(\rho(A)<1\), the Neumann series
\[
 (I-A)^{-1}=\sum_{n\ge0}A^n
\]
converges and is nonnegative.  Choose a strictly positive vector \(e\) such
that \(Ae<e\).  One construction is
\(e=(I-A)^{-1}{\bf 1}\), for then \(e-Ae={\bf 1}>0\).
Set \(w_T=1\) and
\[
 w_{HP}=(I-A)^{-1}u+\varepsilon e.
\]
This vector is strictly positive for every \(\varepsilon>0\), and
\[
 Aw_{HP}+u
 =(I-A)^{-1}u+\varepsilon Ae
 <w_{HP}.
\]
The last component is
\[
 v^{\mathsf T}w_{HP}+\kappa
 =v^{\mathsf T}(I-A)^{-1}u+\kappa
   +\varepsilon v^{\mathsf T}e.
\]
By the strict second inequality in \eqref{eq:v5v26-Schur-condition},
choose \(\varepsilon>0\) small enough that this is less than \(1=w_T\).
Thus \(Lw<w\) componentwise.  The maximum of the three ratios
\((Lw)_i/w_i\) is a number \(q<1\).  Apply
Lemma~\ref{lem:v5v26-weighted-max}.
\end{proof}

\begin{remark}[Meaning of the feedback term]
The scalar
\begin{equation}
 v^{\mathsf T}(I-A)^{-1}u
 =\sum_{n\ge0}v^{\mathsf T}A^nu
 \label{eq:v5v26-feedback-series}
\end{equation}
is the total tail-to-head/protected-to-tail feedback, including arbitrarily
many returns inside the finite block.  Bounding the three rows separately
by their unweighted sums would discard this routing information and can
fail even when \eqref{eq:v5v26-Schur-condition} holds.
\end{remark}

\begin{corollary}[Explicit two-by-two head condition]
\label{cor:v5v26-two-by-two}
For
\[
 A=\begin{pmatrix}a&b\\c&d\end{pmatrix}\ge0,
\]
the condition \(\rho(A)<1\) is equivalent to
\begin{equation}
 a<1,\qquad d<1,\qquad (1-a)(1-d)>bc.
 \label{eq:v5v26-2x2}
\end{equation}
Under these inequalities,
\begin{equation}
 (I-A)^{-1}
 =\frac1{(1-a)(1-d)-bc}
 \begin{pmatrix}1-d&b\\c&1-a\end{pmatrix}.
 \label{eq:v5v26-2x2-inverse}
\end{equation}
Hence \eqref{eq:v5v26-Schur-condition} is a completely explicit inequality
in the nine block constants.
\end{corollary}

\begin{proof}
For a nonnegative \(2\times2\) matrix, \(I-A\) is a nonsingular
\(M\)-matrix exactly when its two principal minors are positive.  Those
minors give \eqref{eq:v5v26-2x2}; equivalently, the larger eigenvalue of
\(A\) is below one.  Direct inversion gives
\eqref{eq:v5v26-2x2-inverse}.
\end{proof}

\subsection{Invariance and assembly of the Ward surface}

Let \(\mathcal W:\mathbb B\to\mathbb Z\) be the finite collection of
renormalized gauge, BRST, and diffeomorphism Ward defects assigned to the
protected sector.  The physical slice at this stage is
\begin{equation}
 \mathfrak S=\ker\mathcal W.
 \label{eq:v5v26-Ward-slice}
\end{equation}
The phrase finite collection refers to the identities imposed at a
fixed local grade and shell; compatibility across all grades is encoded by
the factorial sum.

\begin{definition}[Assembled shell Ward certificate]
\label{def:v5v26-ASWC}
An ASWC is an identity
\begin{equation}
 \mathcal W\mathcal T(x)=\mathcal R(x)\mathcal W x
 \label{eq:v5v26-intertwining}
\end{equation}
on \(D\), where \(\mathcal R(x)\) is a bounded operator on \(\mathbb Z\)
and the left side is evaluated after summing the complete shell insertion,
counterterm, ghost, and improvement contributions belonging to the same
Ward symbol.
\end{definition}

\begin{proposition}[Exact Ward invariance]
\label{prop:v5v26-Ward-invariance}
If the ASWC holds, then
\begin{equation}
 \mathcal T(D\cap\mathfrak S)\subseteq\mathfrak S.
 \label{eq:v5v26-invariance}
\end{equation}
If \(\mathcal T\) is a contraction on a closed invariant subset
\(B\subset D\cap\mathfrak S\), its unique fixed point belongs to
\(\mathfrak S\).
\end{proposition}

\begin{proof}
For \(x\in\ker\mathcal W\), equation
\eqref{eq:v5v26-intertwining} gives
\(\mathcal W\mathcal T(x)=0\).  Thus the slice is invariant.  Banach's
fixed-point theorem constructs the fixed point as the limit of iterates
starting from any point of \(B\); since \(\mathcal T\) maps \(B\) to
itself, that point lies in \(B\subset\mathfrak S\).
\end{proof}

\begin{lemma}[Diagramwise identities are sufficient but not necessary]
\label{lem:v5v26-assembly}
Suppose
\(\mathcal T=\sum_{\gamma\in\Gamma}\mathcal T_\gamma\) converges in the
hybrid space.  If each graph obeys
\(\mathcal W\mathcal T_\gamma=\mathcal R_\gamma\mathcal W\), then the ASWC
holds with \(\mathcal R=\sum_\gamma\mathcal R_\gamma\), provided the latter
series converges in operator norm.  Conversely, the ASWC may hold although
no individual graph has this property.
\end{lemma}

\begin{proof}
The sufficient statement follows by continuity of \(\mathcal W\) and
termwise summation.  For the converse, take two contributions with Ward
defects \(z\) and \(-z\).  Their sum has zero defect although neither
contribution does.  This is the operator form of
Lemma~\ref{lem:v5v25-cancellation-loss}.
\end{proof}

\begin{firewall}[The certificate is not yet an identity]
\label{fw:v5v26-ASWC-open}
Definition~\ref{def:v5v26-ASWC} states what the regularized coupled
SM--Einstein shell step must satisfy.  Its validity requires a regulator
and counterterm prescription obeying the joint Slavnov--Taylor and
diffeomorphism identities.  No such all-order construction was proved in
the preceding notes, so \eqref{eq:v5v26-intertwining} remains a certificate
rather than a conclusion for the physical model.
\end{firewall}

\subsection{Eliminating the renormalized head}

Relevant and marginal couplings are generally not contracted until
renormalization conditions determine them.  We make that elimination
quantitative.  Write \(y=(p,t)\in\mathbb P\oplus\mathbb T\), and let
\(\mathcal F(h,y)=0\) be the finite system of head renormalization
conditions.

\begin{theorem}[Quantitative head solver]
\label{thm:v5v26-head-solver}
Let \(h_0\in\mathbb H\), let \(U\) be a closed ball about \(h_0\), and
suppose
\begin{equation}
 \mathcal F(h,y)=0
 \quad\Longleftrightarrow\quad
 h=\mathcal G(h,y)
 \label{eq:v5v26-head-fixed}
\end{equation}
on \(U\times Y\).  Assume
\begin{align}
 \|\mathcal G(h,y)-\mathcal G(\widetilde h,y)\|_H
 &\le a\|h-\widetilde h\|_H,\qquad a<1,
 \label{eq:v5v26-head-h}\\
 \|\mathcal G(h,y)-\mathcal G(h,\widetilde y)\|_H
 &\le c\|y-\widetilde y\|_Y,
 \label{eq:v5v26-head-y}
\end{align}
and that \(\mathcal G(\,\cdot\,,y)\) maps \(U\) to itself for every
\(y\in Y\).  Then there is a unique solution \(h=\chi(y)\in U\), and
\begin{equation}
 \|\chi(y)-\chi(\widetilde y)\|_H
 \le\frac{c}{1-a}\|y-\widetilde y\|_Y.
 \label{eq:v5v26-head-Lipschitz}
\end{equation}
\end{theorem}

\begin{proof}
For fixed \(y\), equation \eqref{eq:v5v26-head-h} and Banach's theorem give
a unique fixed point.  For two data values,
\begin{align*}
 \|\chi(y)-\chi(\widetilde y)\|_H
 &\le
 \|\mathcal G(\chi(y),y)-\mathcal G(\chi(\widetilde y),y)\|_H\\
 &\quad+
 \|\mathcal G(\chi(\widetilde y),y)
       -\mathcal G(\chi(\widetilde y),\widetilde y)\|_H\\
 &\le a\|\chi(y)-\chi(\widetilde y)\|_H+c\|y-\widetilde y\|_Y.
\end{align*}
Move the first term to the left.
\end{proof}

Substituting \(h=\chi(p,t)\) into the protected and tail rows changes their
Lipschitz constants by at most the head-column constants multiplied by
\(c/(1-a)\).  Thus the Schur test is to be applied after head tuning, unless
the head row is already contractive.  This prevents a relevant coupling
with eigenvalue one from being hidden inside a claimed full-space
contraction.

\subsection{Ball invariance and the exact smallness condition}

\begin{theorem}[Certified nonlinear shell contraction]
\label{thm:v5v26-certified-shell}
Assume the Schur conditions \eqref{eq:v5v26-Schur-condition}, and let
\(w,q\) be constructed in Theorem~\ref{thm:v5v26-Schur}.  Let
\[
 B_w(R)=\{x:\|x\|_w\le R\}\subset D.
\]
If
\begin{equation}
 \|\mathcal T(0)\|_w\le(1-q)R,
 \label{eq:v5v26-ball-smallness}
\end{equation}
then \(\mathcal T\) maps \(B_w(R)\) to itself and has a unique fixed point
there.  If the ASWC holds and \(0\in\mathfrak S\), the same conclusion
holds on \(B_w(R)\cap\mathfrak S\).
\end{theorem}

\begin{proof}
For \(x\in B_w(R)\), Lemma~\ref{lem:v5v26-weighted-max} gives
\[
 \|\mathcal T(x)\|_w
 \le\|\mathcal T(x)-\mathcal T(0)\|_w+\|\mathcal T(0)\|_w
 \le qR+(1-q)R=R.
\]
The contraction theorem gives the unique fixed point.  Under the ASWC,
Proposition~\ref{prop:v5v26-Ward-invariance} makes the Ward slice invariant,
so the restricted conclusion follows.
\end{proof}

\begin{corollary}[Perturbative shell-size test]
\label{cor:v5v26-perturbative-test}
Suppose, uniformly in the ultraviolet cutoff,
\[
 A=O(g),\quad u=O(g),\quad v=O(g),\quad
 \kappa\le\kappa_0+O(g),\quad \kappa_0<1,
\]
and \(\|\mathcal T(0)\|_w=O(g)\) as \(g\to0\).  Then for sufficiently small
\(g\), the Schur and ball conditions hold for some fixed \(R>0\).
\end{corollary}

\begin{proof}
For small \(g\), \(\rho(A)=O(g)<1\) and
\[
 v^{\mathsf T}(I-A)^{-1}u=O(g^2).
\]
Hence the left side of the second Schur inequality tends to
\(\kappa_0<1\).  The resulting \(q\) remains below one after decreasing
the smallness range.  Finally \(O(g)\le(1-q)R\) for sufficiently small
\(g\).
\end{proof}

\begin{firewall}[Uniformity is the substantive estimate]
\label{fw:v5v26-uniformity}
Corollary~\ref{cor:v5v26-perturbative-test} is useless if the constants
hidden in \(O(g)\) grow with the ultraviolet cutoff, the volume, or the
local grade.  Uniform tri-graded bounds and the strong-image estimate of
V25 must be proved before the corollary yields a continuum limit.
\end{firewall}

\subsection{V26 status and the next calculation}

\begin{theorem}[V26 reduction]
\label{thm:v5v26-status}
For any regularized coupled shell map with certified block Lipschitz
constants, head renormalization solver, ASWC, and cutoff-uniform strong
image, the inequalities
\eqref{eq:v5v26-Schur-condition} and
\eqref{eq:v5v26-ball-smallness} imply a unique Ward-compatible fixed point.
Together with V25, consistently truncated fixed points then converge in the
weaker hybrid tri-graded topology.
\end{theorem}

\begin{proof}
Use Theorem~\ref{thm:v5v26-head-solver} if the relevant head must first be
eliminated.  Apply Theorem~\ref{thm:v5v26-certified-shell} to the resulting
map on the Ward slice.  The strong-image hypothesis and
Theorem~\ref{thm:v5v25-strong-image} give relative compactness, while
Corollary~\ref{cor:v5v25-truncated-fixed-points} gives convergence of
consistent truncations.
\end{proof}

The theorem is a reduction, not yet the SM--Einstein construction.  The
first unresolved calculation is the ASWC for a concrete regularized
coupled action.  V27 will construct the classical joint gauge--diffeomorphism
Ward complex for the Yang--Mills, Higgs, fermion, and Einstein sectors and
prove closure of its linearized differential.  That calculation will
identify which regularization and anomaly statements are needed before an
all-order shell ASWC can be claimed.

\clearpage
\section{V27: The joint gauge--diffeomorphism shell Ward complex}
\label{sec:v5v27}

\subsection{Scope relative to the inherited work}

The fifth-series summary records an anomaly-normalized three-generation
fermion determinant line, and V12 proves a finite-dimensional Hodge
reduction for a given complex.  Repeating either construction would not
advance the present bottleneck.  This note supplies the missing link: the
complex generated by the coupled classical SM--Einstein action and the
condition under which one shell integration is a chain map for that
complex.

Let \(M\) be a compact oriented spin four-manifold at finite volume and
let \(P\to M\) carry the faithful Standard Model gauge group fixed in the
archived construction.  Let \(\mathcal C_\Lambda\) be a finite-cutoff
configuration manifold containing the metric, gauge connections, Higgs
field, fermions, ghosts, and any auxiliary fields retained at cutoff
\(\Lambda\).  Finite cutoff is used so all derivatives and changes of
variables below are ordinary finite-dimensional operations.  A cofinal
limit requires the uniform estimates of V25--V26.

\subsection{The semidirect symmetry algebra and BRST closure}

The infinitesimal symmetry algebra is the semidirect product
\begin{equation}
 \mathfrak a
 =
 \Gamma(TM)\ltimes\Gamma(\operatorname{ad}P).
 \label{eq:v5v27-symmetry-algebra}
\end{equation}
With the convention that vector fields act on adjoint sections by the
chosen natural lift, its bracket is
\begin{equation}
 [(\xi,\lambda),(\eta,\mu)]
 =
 \bigl([\xi,\eta],
 \mathcal L_\xi\mu-\mathcal L_\eta\lambda+[\lambda,\mu]\bigr).
 \label{eq:v5v27-semidirect-bracket}
\end{equation}
The construction is independent of a local trivialization because every
term is the infinitesimal form of the bundle automorphism action.

Let \(R_\epsilon\) be the fundamental vector field on
\(\mathcal C_\Lambda\) generated by \(\epsilon\in\mathfrak a\).  We use
the left-action convention
\begin{equation}
 [R_\epsilon,R_\zeta]=R_{[\epsilon,\zeta]}.
 \label{eq:v5v27-representation}
\end{equation}
For example, its metric component is \(\mathcal L_\xi g\); its connection
component is the sum of the lifted diffeomorphism variation and the
covariant gauge variation; and its matter components are the induced
associated-bundle actions.  This formulation includes the spin lift
without replacing it by an ill-defined componentwise Lie derivative.

Introduce an odd ghost \(C\in\Pi\mathfrak a\).  The Chevalley--Eilenberg
differential \(s\) is fixed by
\begin{equation}
 s\phi=R_C\phi,\qquad
 sC=-\frac12[C,C],
 \label{eq:v5v27-BRST}
\end{equation}
and the graded Leibniz rule.  Equation \eqref{eq:v5v27-BRST} also fixes all
sign conventions in what follows.

\begin{theorem}[Off-shell classical BRST closure]
\label{thm:v5v27-BRST-closure}
On polynomial functions of the cutoff fields and ghosts,
\begin{equation}
 s^2=0.
 \label{eq:v5v27-nilpotent}
\end{equation}
This conclusion uses no field equation.
\end{theorem}

\begin{proof}
On the ghost, the graded Jacobi identity gives
\[
 s^2C
 =\frac14\bigl([[C,C],C]-[C,[C,C]]\bigr)=0
\]
with the two terms interpreted in the graded extension of
\(\mathfrak a\).  On a field coordinate, applying the graded Leibniz rule
to \(R_C\phi\) gives two contributions: the ghost variation
\(-\frac12R_{[C,C]}\phi\) and the antisymmetrized double action
\(\frac12[R_C,R_C]\phi\).  They cancel by
\eqref{eq:v5v27-representation}.  Nilpotence on all polynomial functions
then follows from the derivation rule.
\end{proof}

\begin{firewall}[Closure versus quantization]
\label{fw:v5v27-classical-only}
Theorem~\ref{thm:v5v27-BRST-closure} is an algebraic statement about the
classical symmetry action.  It neither proves invariance of a regulator or
fermion measure nor removes a global anomaly of the determinant line.
\end{firewall}

\subsection{The linearized Noether complex}

Let \(S_\Lambda:\mathcal C_\Lambda\to\mathbb R\) be a twice
differentiable, invariant cutoff action:
\begin{equation}
 DS_\Lambda(\phi)R_\epsilon(\phi)=0
 \quad\text{for all }\phi,\epsilon.
 \label{eq:v5v27-invariance}
\end{equation}
Let \(\bar\phi\) be a stationary background,
\(DS_\Lambda(\bar\phi)=0\).  Put
\[
 E_\Lambda=T_{\bar\phi}\mathcal C_\Lambda,\qquad
 R_{\bar\phi}\epsilon=R_\epsilon(\bar\phi),\qquad
 K_{\bar\phi}=D^2S_\Lambda(\bar\phi).
\]
After choosing the cutoff inner products, \(K_{\bar\phi}\) is identified
with a symmetric map \(E_\Lambda\to E_\Lambda^*\), and
\(R_{\bar\phi}^*:E_\Lambda^*\to\mathfrak a_\Lambda^*\) is the adjoint
Noether map.

\begin{theorem}[Coupled gauge--gravity Noether complex]
\label{thm:v5v27-Noether-complex}
At every stationary background of an invariant cutoff action,
\begin{equation}
 \mathfrak a_\Lambda
 \xrightarrow{\ R_{\bar\phi}\ }
 E_\Lambda
 \xrightarrow{\ K_{\bar\phi}\ }
 E_\Lambda^*
 \xrightarrow{\ R_{\bar\phi}^*\ }
 \mathfrak a_\Lambda^*
 \label{eq:v5v27-Noether-complex}
\end{equation}
is a complex:
\begin{equation}
 K_{\bar\phi}R_{\bar\phi}=0,\qquad
 R_{\bar\phi}^*K_{\bar\phi}=0.
 \label{eq:v5v27-Noether-zero}
\end{equation}
It includes the Yang--Mills, hypercharge, Higgs, chiral-fermion, and metric
blocks and all their classical cross terms.
\end{theorem}

\begin{proof}
Differentiate \eqref{eq:v5v27-invariance} at \(\bar\phi\) in the direction
\(v\in E_\Lambda\):
\[
 K_{\bar\phi}(v,R_{\bar\phi}\epsilon)
 +DS_\Lambda(\bar\phi)
   \bigl(DR_\epsilon(\bar\phi)v\bigr)=0.
\]
The second term vanishes because the background is stationary.  Since this
holds for all \(v\), symmetry of the Hessian gives
\(K_{\bar\phi}R_{\bar\phi}\epsilon=0\).  Taking the adjoint yields
\(R_{\bar\phi}^*K_{\bar\phi}=0\).
\end{proof}

\begin{remark}[Why stationarity is needed]
Away from a critical background the differentiated identity contains the
term
\(DS_\Lambda(\phi)DR_\epsilon(\phi)\).  Thus the Hessian alone does not
annihilate gauge directions off shell.  The BV complex restores off-shell
closure by adjoining antifields; equation
\eqref{eq:v5v27-Noether-complex} is the on-background deformation complex.
\end{remark}

\begin{corollary}[Hodge input for V12]
\label{cor:v5v27-Hodge-input}
After gauge fixing, any choice of cutoff adjoints turns
\eqref{eq:v5v27-Noether-complex} into the finite Hilbert-complex input of
V12.  A uniform Hodge gap is additional analytic information and does not
follow from the two zero compositions.
\end{corollary}

\begin{proof}
The zero compositions are precisely the algebraic complex identities
needed to form the associated Laplacians.  Spectral lower bounds depend on
the chosen background, cutoff, gauge fixing, and boundary conditions, so
they are not consequences of algebraic closure.
\end{proof}

\subsection{One exact shell integral}

We now state a finite-cutoff theorem that produces the assembled Ward
identity rather than assuming it graph by graph.  Split the variables as
\(\phi=(x,y)\), with \(x\) retained and \(y\) integrated over one shell.
Assume the infinitesimal action splits as
\begin{equation}
 R_\epsilon(x,y)=
 \bigl(R_\epsilon^{\mathrm{lo}}(x),
       R_\epsilon^{\mathrm{hi}}(x,y)\bigr).
 \label{eq:v5v27-equivariant-split}
\end{equation}
Let \(d\mu_x(y)\) be a possibly \(x\)-dependent shell measure and
\(\Gamma\) an invariant integration cycle.  Define
\begin{equation}
 e^{-S_{\mathrm{eff}}(x)}
 =
 Z(x):=\int_\Gamma e^{-S(x,y)}\,d\mu_x(y),
 \qquad 0<Z(x)<\infty.
 \label{eq:v5v27-shell-integral}
\end{equation}

\begin{theorem}[Equivariant shell integration]
\label{thm:v5v27-equivariant-shell}
Suppose:
\begin{enumerate}
 \item \(S\) is invariant under the simultaneous low and high action;
 \item the family \(d\mu_x\) and the cycle \(\Gamma\) are equivariant, so
 the change of variables generated by
 \(R_\epsilon^{\mathrm{hi}}\) has unit Jacobian and no boundary term;
 \item differentiation under the integral is justified.
\end{enumerate}
Then
\begin{equation}
 R_\epsilon^{\mathrm{lo}}S_{\mathrm{eff}}(x)=0.
 \label{eq:v5v27-effective-invariance}
\end{equation}
Equivalently, the complete shell integration preserves the classical Ward
surface.
\end{theorem}

\begin{proof}
Let \(g_t\) be the one-parameter symmetry generated by \(\epsilon\).
Equivariance of the split, measure, and cycle permits the high-variable
change of variables
\(y\mapsto g_t^{\mathrm{hi}}\cdot y\).  Simultaneous invariance of \(S\)
then gives
\[
 Z(g_t^{\mathrm{lo}}\cdot x)=Z(x).
\]
Differentiate at \(t=0\).  Since \(Z>0\),
\[
 R_\epsilon^{\mathrm{lo}}S_{\mathrm{eff}}
 =-Z^{-1}R_\epsilon^{\mathrm{lo}}Z=0.
\]
All shell diagrams, ghosts, and Jacobian contributions have already been
combined inside \(Z\), so no premature triangle inequality occurs.
\end{proof}

\begin{corollary}[Shell chain map at a stationary point]
\label{cor:v5v27-shell-chain}
If \(\bar x\) is stationary for \(S_{\mathrm{eff}}\), its Hessian and
retained symmetry generator form the Noether complex
\eqref{eq:v5v27-Noether-complex}.  Thus an exact equivariant shell step is
a chain map between the corresponding symmetry deformation problems.
\end{corollary}

\begin{proof}
Apply Theorem~\ref{thm:v5v27-Noether-complex} to
\eqref{eq:v5v27-effective-invariance}.
\end{proof}

\subsection{The regulator defect and the anomaly class}

When the high-variable change of variables has logarithmic Jacobian
\(\mathcal A_\epsilon(x,y)\), the same calculation gives the precise
defect
\begin{equation}
 R_\epsilon^{\mathrm{lo}}S_{\mathrm{eff}}(x)
 =
 -\frac{\displaystyle
 \int_\Gamma \mathcal A_\epsilon(x,y)e^{-S(x,y)}\,d\mu_x(y)}
 {\displaystyle
 \int_\Gamma e^{-S(x,y)}\,d\mu_x(y)}
 \label{eq:v5v27-anomaly-expectation}
\end{equation}
up to the sign chosen in the definition of the logarithmic Jacobian.  We
fix the sign by \eqref{eq:v5v27-anomaly-expectation}.

\begin{proposition}[Wess--Zumino consistency]
\label{prop:v5v27-WZ}
Let \(\mathfrak A\) be the ghost-number-one Ward defect of a regularized
effective action.  If the classical BRST operator remains nilpotent on the
regularized field algebra, then
\begin{equation}
 s\mathfrak A=0.
 \label{eq:v5v27-WZ}
\end{equation}
If \(\mathfrak A=sB\) for a local ghost-number-zero functional \(B\), the
counterterm \(-B\), with the sign matched to the Ward convention, removes
the defect.  A nonzero cohomology class \([\mathfrak A]\in H^1(s)\) cannot
be removed by a local counterterm.
\end{proposition}

\begin{proof}
Write the broken Ward equation as \(s\Gamma=\mathfrak A\).  Applying \(s\)
and using \(s^2=0\) gives \eqref{eq:v5v27-WZ}.  Replacing
\(\Gamma\) by \(\Gamma-B\) changes the defect by \(-sB\).  Hence an exact
defect is removable.  Conversely, every local counterterm changes the
defect by an exact term, so it cannot alter its cohomology class.
\end{proof}

\begin{assessment}[Inherited anomaly input]
\label{ass:v5v27-inherited-anomaly}
The archived determinant-line construction already supplies the
three-generation Standard Model anomaly normalization and zero-mode
bookkeeping.  V27 does not recompute the
\(SU(3)^2U(1)\), \(SU(2)^2U(1)\), \(U(1)^3\), mixed
gauge--gravitational, or global weak-sector tests.  What remains to be
proved is that the chosen joint metric and gauge shell regulator realizes
that anomaly-free line as an equivariant measure family satisfying
Theorem~\ref{thm:v5v27-equivariant-shell}.
\end{assessment}

\begin{firewall}[Anomaly cancellation is not ultraviolet closure]
\label{fw:v5v27-anomaly-not-UV}
Vanishing of the Ward cohomology obstruction does not bound the shell
integral, close the gravitational counterterm tower, cure the hypercharge
trajectory, or prove reflection positivity.  It removes only the
cohomological obstruction to restoring the stated symmetry identity.
\end{firewall}

\subsection{Passage of the Ward identity to a limit}

\begin{lemma}[Closedness of the Ward surface]
\label{lem:v5v27-Ward-closed}
Let \(x_n\to x\) in the weak hybrid tri-graded space.  Suppose
\(\mathcal W_nx_n=0\), the operators \(\mathcal W_n\) converge strongly to
a bounded operator \(\mathcal W\), and
\begin{equation}
 \|(\mathcal W_n-\mathcal W)x_n\|_{\mathbb Z}\longrightarrow0.
 \label{eq:v5v27-Ward-uniform}
\end{equation}
Then \(\mathcal Wx=0\).
\end{lemma}

\begin{proof}
Boundedness of \(\mathcal W\) gives
\[
 \|\mathcal Wx\|
 \le\|\mathcal W(x-x_n)\|
   +\|(\mathcal W-\mathcal W_n)x_n\|
   +\|\mathcal W_nx_n\|.
\]
All three terms tend to zero.
\end{proof}

\begin{theorem}[Conditional cofinal Ward compatibility]
\label{thm:v5v27-cofinal-Ward}
Assume the hypotheses of the V25 strong-image contraction theorem and V26
Schur criterion uniformly along a cofinal cutoff sequence.  If every
truncated fixed point obeys its assembled shell Ward identity and
\eqref{eq:v5v27-Ward-uniform} holds, then the limiting fixed point lies on
the limiting Ward surface.
\end{theorem}

\begin{proof}
V25 gives convergence of the consistent truncated fixed points in the weak
hybrid topology.  Apply Lemma~\ref{lem:v5v27-Ward-closed}.
\end{proof}

\subsection{V27 status and next upstream step}

\begin{theorem}[V27 conclusion]
\label{thm:v5v27-status}
The classical cutoff SM--Einstein symmetry action has a nilpotent
semidirect BRST differential.  At a stationary background its Hessian
belongs to the coupled Noether complex
\eqref{eq:v5v27-Noether-complex}.  A shell integral preserves that complex
when its split, measure, cycle, and differentiation are equivariant; any
failure is the assembled Jacobian defect
\eqref{eq:v5v27-anomaly-expectation}, whose removable part is exactly its
BRST-exact part.  These statements do not yet verify equivariance or
uniform bounds for a concrete joint regulator.
\end{theorem}

\begin{proof}
Combine Theorems~\ref{thm:v5v27-BRST-closure},
\ref{thm:v5v27-Noether-complex}, and
\ref{thm:v5v27-equivariant-shell} with
Proposition~\ref{prop:v5v27-WZ}.
\end{proof}

The new bottleneck is therefore a repair problem rather than another
formal Ward declaration.  V28 will derive a homological equation for the
finite-cutoff regulator defect.  It will prove that a uniformly bounded
contracting homotopy produces symmetry-restoring counterterms for the
exact component, while a harmonic ghost-number-one component survives as
the genuine anomaly obstruction.

\clearpage
\section{V28: Homological repair of a finite-cutoff Ward defect}
\label{sec:v5v28}

\subsection{From a defect to a solvable equation}

V27 reduces failure of an equivariant shell step to an assembled
ghost-number-one defect.  The equation for a restoring counterterm is not
merely that the defect be small.  It must be exact in the local Ward
complex, and the inverse producing its primitive must be uniformly bounded
in the tri-graded topology.  This note proves that statement and separates
three obstructions: a harmonic anomaly, loss of locality, and collapse of
the homotopy bound.

At a fixed cutoff, consider a finite Hilbert cochain complex
\begin{equation}
 C^0_\Lambda
 \xrightarrow{\ s_0\ }
 C^1_\Lambda
 \xrightarrow{\ s_1\ }
 C^2_\Lambda,
 \qquad s_1s_0=0.
 \label{eq:v5v28-complex}
\end{equation}
Here \(C^0_\Lambda\) is the allowed counterterm space,
\(C^1_\Lambda\) is the space of assembled Ward defects, and
\(C^2_\Lambda\) contains the consistency defects.  In the local theory
these spaces must be read as local functionals modulo total derivatives;
the finite-dimensional model is the cutoff realization of that quotient.

Define
\begin{equation}
 \Delta_1=s_0s_0^*+s_1^*s_1,\qquad
 \mathcal H^1_\Lambda=\ker s_1\cap\ker s_0^*.
 \label{eq:v5v28-Hodge-data}
\end{equation}
Let \(\Pi_1\) be the orthogonal projection onto
\(\mathcal H^1_\Lambda\), and let \(G_1\) be
\(\Delta_1^{-1}\) on \((\mathcal H^1_\Lambda)^\perp\), extended by zero on
\(\mathcal H^1_\Lambda\).

\begin{theorem}[Closed-defect decomposition]
\label{thm:v5v28-closed-decomposition}
Every consistency-closed defect \(A\in\ker s_1\) has the orthogonal
decomposition
\begin{equation}
 A=s_0B+A_{\mathrm{harm}},
 \qquad
 B=s_0^*G_1A,\qquad
 A_{\mathrm{harm}}=\Pi_1A.
 \label{eq:v5v28-defect-decomposition}
\end{equation}
Consequently \(A\) is removable by an allowed counterterm if
\(\Pi_1A=0\).  If its harmonic component is nonzero, no counterterm in
\(C^0_\Lambda\) removes it.
\end{theorem}

\begin{proof}
Finite-dimensional Hodge decomposition in degree one gives
\[
 C^1_\Lambda
 =\operatorname{Ran}s_0
  \oplus\mathcal H^1_\Lambda
  \oplus\operatorname{Ran}s_1^*.
\]
If the coexact component of \(A\in\ker s_1\) is \(s_1^*z\), then
\[
 \|s_1^*z\|^2
 =\langle s_1^*z,A\rangle
 =\langle z,s_1A\rangle=0,
\]
so that component vanishes.  The Green formula
\[
 I-\Pi_1=s_0s_0^*G_1+s_1^*s_1G_1
\]
therefore reduces on \(A\) to
\((I-\Pi_1)A=s_0s_0^*G_1A\), proving
\eqref{eq:v5v28-defect-decomposition}.  Every exact counterterm shift
changes \(A\) by an element of \(\operatorname{Ran}s_0\), which is
orthogonal to \(\mathcal H^1_\Lambda\); hence a nonzero harmonic component
cannot be removed.
\end{proof}

\begin{corollary}[Quantitative repair bound]
\label{cor:v5v28-repair-bound}
If the smallest positive eigenvalue of \(\Delta_1\) is
\(\lambda_{1,\Lambda}>0\), then
\begin{equation}
 \|B\|_{C^0_\Lambda}
 \le
 \frac{\|s_0^*\|}{\lambda_{1,\Lambda}}
 \|A\|_{C^1_\Lambda}.
 \label{eq:v5v28-repair-bound}
\end{equation}
\end{corollary}

\begin{proof}
The spectral theorem gives
\(\|G_1\|=\lambda_{1,\Lambda}^{-1}\).  Apply the definition
\(B=s_0^*G_1A\).
\end{proof}

\begin{firewall}[Small defects can require large counterterms]
\label{fw:v5v28-small-large}
Even when \(\|A_\Lambda\|\to0\), the canonical primitive can diverge if
\(\lambda_{1,\Lambda}\to0\) faster.  Pointwise anomaly cancellation is
therefore insufficient for a cofinal construction; the homotopy must be
uniformly bounded in the topology used by the shell map.
\end{firewall}

\subsection{The locality obstruction}

The Green operator \(G_1\) above is an inverse on the full cutoff cochain
space.  It need not preserve local functionals.  Let
\(C^\bullet_{\Lambda,\mathrm{loc}}\subset C^\bullet_\Lambda\) be a
subcomplex of allowed local counterterms and local defects.

\begin{definition}[Local repair certificate]
\label{def:v5v28-LRC}
A local repair certificate consists of maps
\begin{equation}
 h_1:C^1_{\Lambda,\mathrm{loc}}\to
 C^0_{\Lambda,\mathrm{loc}},\qquad
 \Pi_{\mathrm{an}}:C^1_{\Lambda,\mathrm{loc}}\to
 \mathcal H^1_{\Lambda,\mathrm{loc}}
 \label{eq:v5v28-local-homotopy}
\end{equation}
such that, on \(\ker s_1\cap C^1_{\Lambda,\mathrm{loc}}\),
\begin{equation}
 s_0h_1=I-\Pi_{\mathrm{an}}.
 \label{eq:v5v28-local-contraction}
\end{equation}
The range of \(\Pi_{\mathrm{an}}\) represents the local
ghost-number-one cohomology.
\end{definition}

\begin{proposition}[Global Hodge inversion does not prove local repair]
\label{prop:v5v28-global-not-local}
The condition \(\Pi_1A=0\) in the full Hilbert complex does not imply the
existence of a local primitive unless \(s_0^*G_1A\) is local or an
independent local repair certificate is supplied.
\end{proposition}

\begin{proof}
Theorem~\ref{thm:v5v28-closed-decomposition} constructs a primitive in
\(C^0_\Lambda\).  Exactness in a complex can change after restricting its
degree-zero space: if
\(\operatorname{Ran}(s_0|_{C^0_{\rm loc}})\) is a proper subset of
\(\operatorname{Ran}s_0\), a vector can belong to the latter but not the
former.  Since Green kernels are generally nonlocal, membership in the
smaller range needs a separate proof.
\end{proof}

\begin{theorem}[Local symmetry restoration]
\label{thm:v5v28-local-restoration}
Let \(A\in C^1_{\Lambda,\mathrm{loc}}\) satisfy \(s_1A=0\).  Under a local
repair certificate,
\begin{equation}
 A=s_0(h_1A)+\Pi_{\mathrm{an}}A.
 \label{eq:v5v28-local-decomposition}
\end{equation}
Thus the local counterterm \(-h_1A\), with sign adjusted to the convention
for the broken Ward equation, removes exactly the component
\((I-\Pi_{\mathrm{an}})A\).  It removes all of \(A\) if and only if
\(\Pi_{\mathrm{an}}A=0\).
\end{theorem}

\begin{proof}
Equation \eqref{eq:v5v28-local-decomposition} is
\eqref{eq:v5v28-local-contraction} applied to \(A\).  A local counterterm
changes the Ward defect by an \(s_0\)-exact local functional, so the exact
component is canceled.  The projected component represents the quotient
\(\ker s_1/\operatorname{Ran}s_0\) and is unchanged by such a shift.
\end{proof}

\subsection{Tri-graded bounds for the repair operator}

Give each local cochain space the curvature-grade, particle-number, and
spatial spectral weights of V25.  The cochain degree is separate from those
three analytic grades.  Denote the strong and weak norms by
\(\|\cdot\|_+\) and \(\|\cdot\|_-\).

\begin{definition}[Uniform graded homotopy certificate]
\label{def:v5v28-UGHC}
A UGHC is a local repair certificate for every cutoff such that
\begin{equation}
 \|h_{1,\Lambda}A\|_-
 \le C_h\|A\|_+
 \label{eq:v5v28-uniform-homotopy}
\end{equation}
with \(C_h\) independent of cutoff, volume, curvature grade, and particle
number.
\end{definition}

\begin{theorem}[Gevrey stability of local repair]
\label{thm:v5v28-Gevrey-repair}
Assume \(h_1\) acts blockwise without increasing curvature grade or
particle number and obeys, uniformly in \((p,m)\),
\begin{equation}
 \|h_1A_{p,m}\|_{p,m;s_-,b_-}
 \le C_h\|A_{p,m}\|_{p,m;s_+,b_+}.
 \label{eq:v5v28-block-h}
\end{equation}
Then it extends to a bounded map from the strong to the weak hybrid
tri-graded space with the same constant:
\begin{equation}
 \|h_1A\|_{\rho_-,\eta_-;s_-,b_-}
 \le C_h\|A\|_{\rho_+,\eta_+;s_+,b_+}.
 \label{eq:v5v28-trigraded-h}
\end{equation}
This includes both factorial weights unchanged.
\end{theorem}

\begin{proof}
Multiply \eqref{eq:v5v28-block-h} by the weak block weight.  The quotient
of weak and strong weights is at most one by
Lemma~\ref{lem:v5v25-monotonicity}.  Sum over \((p,m)\).  On the protected
sector the signed packet is assembled before applying \(h_1\) and before
taking its Hilbert norm; on the tail use the declared absolute norm.
\end{proof}

\begin{remark}[Controlled grade loss]
If \(h_1\) raises curvature grade by at most \(r\), particle number by at
most \(k\), or loses \(a\) spatial derivatives, the same proof works after
paying the corresponding strict radius and regularity gaps.  The constants
then contain finite factors depending on \(r,k,a\).  An unbounded grade
shift is not covered by the theorem.
\end{remark}

\subsection{Order-by-order restoration}

Let \(g\) be a formal coupling parameter and suppose the broken Slavnov
identity has been restored through order \(g^{N-1}\).  At order \(g^N\),
write its local defect as \(A_N\).

\begin{theorem}[Filtered homological induction]
\label{thm:v5v28-induction}
Assume at every order:
\begin{enumerate}
 \item lower-order restoration implies the consistency condition
 \(s_1A_N=0\);
 \item a single local repair certificate applies;
 \item \(\Pi_{\mathrm{an}}A_N=0\).
\end{enumerate}
Then local counterterms can restore the identity through every finite
order.  One may choose the order-\(N\) counterterm
\begin{equation}
 B_N=h_1A_N.
 \label{eq:v5v28-order-counterterm}
\end{equation}
\end{theorem}

\begin{proof}
Induct on \(N\).  The hypotheses make \(A_N\) closed and annihilate its
local cohomology component.  Equation
\eqref{eq:v5v28-local-contraction} gives
\(s_0B_N=A_N\), so adding the counterterm with the cancelling sign removes
the order-\(N\) defect without changing lower orders.  This proves the
inductive step.
\end{proof}

\begin{corollary}[Norm control of the formal repair series]
\label{cor:v5v28-series-bound}
If the UGHC holds and
\begin{equation}
 \sum_{N\ge1}\frac{\rho^N}{(N!)^\alpha}\|A_N\|_+<\infty,
 \label{eq:v5v28-defect-series}
\end{equation}
then the chosen primitives obey
\begin{equation}
 \sum_{N\ge1}\frac{\rho^N}{(N!)^\alpha}\|B_N\|_-
 \le C_h
 \sum_{N\ge1}\frac{\rho^N}{(N!)^\alpha}\|A_N\|_+.
 \label{eq:v5v28-counterterm-series}
\end{equation}
\end{corollary}

\begin{proof}
Apply \eqref{eq:v5v28-uniform-homotopy} at each order and sum the
nonnegative inequalities.
\end{proof}

\begin{firewall}[Formal restoration versus an exact shell map]
\label{fw:v5v28-formal}
Theorem~\ref{thm:v5v28-induction} proves restoration to arbitrary finite
formal order.  An exact Ward-compatible shell map additionally requires
convergence or a summation theorem for the counterterm series and
compatibility of that summation with the shell integral.  The bound
\eqref{eq:v5v28-counterterm-series} supplies this only when the defect
series itself converges in the stated topology.
\end{firewall}

\subsection{Preserving the V26 contraction margin}

Let \(\mathcal T\) be the uncorrected shell map and let
\(\mathcal C\) be the counterterm correction constructed from the defect.
Set \(\mathcal T^{\mathrm{rep}}=\mathcal T+\mathcal C\).  Work in the
weighted block norm supplied by V26.

\begin{proposition}[Stability under a bounded repair]
\label{prop:v5v28-contraction-stability}
Suppose
\begin{align}
 \|\mathcal T(x)-\mathcal T(y)\|_w
 &\le q\|x-y\|_w,\qquad q<1,
 \label{eq:v5v28-old-q}\\
 \|\mathcal C(x)-\mathcal C(y)\|_w
 &\le\varepsilon\|x-y\|_w.
 \label{eq:v5v28-repair-epsilon}
\end{align}
If \(q+\varepsilon<1\), then
\(\mathcal T^{\mathrm{rep}}\) is a contraction.  It preserves a radius
\(R\) ball whenever
\begin{equation}
 \|\mathcal T(0)+\mathcal C(0)\|_w
 \le(1-q-\varepsilon)R.
 \label{eq:v5v28-repaired-ball}
\end{equation}
\end{proposition}

\begin{proof}
The triangle inequality gives Lipschitz constant at most
\(q+\varepsilon\).  For \(\|x\|_w\le R\),
\[
 \|\mathcal T^{\mathrm{rep}}(x)\|_w
 \le(q+\varepsilon)R+
 \|\mathcal T(0)+\mathcal C(0)\|_w\le R.
\]
Banach's theorem then applies.
\end{proof}

\begin{corollary}[Defect-to-contraction budget]
\label{cor:v5v28-defect-budget}
If the map from a defect \(A(x)\) to its inserted counterterm has norm
\(C_{\mathrm{ins}}\), the UGHC holds with \(C_h\), and
\[
 \|A(x)-A(y)\|_+\le L_A\|x-y\|_w,
\]
then one may take
\begin{equation}
 \varepsilon\le C_{\mathrm{ins}}C_hL_A.
 \label{eq:v5v28-epsilon-budget}
\end{equation}
Thus a sufficient repair margin is
\begin{equation}
 q+C_{\mathrm{ins}}C_hL_A<1.
 \label{eq:v5v28-final-margin}
\end{equation}
\end{corollary}

\begin{proof}
Compose the defect Lipschitz estimate, the homotopy bound
\eqref{eq:v5v28-uniform-homotopy}, and the insertion bound, then apply
Proposition~\ref{prop:v5v28-contraction-stability}.
\end{proof}

\subsection{Cofinal obstruction trichotomy}

\begin{theorem}[Repair trichotomy]
\label{thm:v5v28-trichotomy}
For a cofinal family of consistency-closed local Ward defects, the
homological repair route can fail in exactly one or more of the following
ways:
\begin{enumerate}
 \item \(\Pi_{\mathrm{an}}A_\Lambda\ne0\), a cohomological anomaly;
 \item no primitive lies in the allowed local counterterm complex, a
 locality obstruction;
 \item local primitives exist but every suitable homotopy constant or
 insertion norm is unbounded, an analytic obstruction.
\end{enumerate}
If none occurs and the repaired contraction margin
\eqref{eq:v5v28-final-margin} is uniform, the repaired shell maps preserve
the Ward surface and remain strict contractions.
\end{theorem}

\begin{proof}
Theorem~\ref{thm:v5v28-local-restoration} separates the harmonic and exact
local components, giving the first two alternatives.  The UGHC and
Corollary~\ref{cor:v5v28-defect-budget} give the analytic alternative and
the positive conclusion when its constants are uniform.
\end{proof}

\begin{assessment}[Application to the inherited determinant line]
\label{ass:v5v28-determinant}
The archived anomaly normalization addresses the first branch for the
three-generation chiral determinant.  It does not by itself construct the
local homotopy for a moving metric-dependent cutoff or bound
\(C_h\) and \(C_{\mathrm{ins}}\) uniformly.  Those are the live second and
third branches.
\end{assessment}

\subsection{V28 status and next concrete regulator test}

\begin{theorem}[V28 conclusion]
\label{thm:v5v28-status}
A consistency-closed finite-cutoff Ward defect decomposes uniquely into an
exact component and a harmonic anomaly.  A local uniformly bounded
contracting homotopy repairs the exact component in the V25 tri-graded
space, and the repaired V26 shell map remains contractive under the
explicit margin \eqref{eq:v5v28-final-margin}.  Existence and uniformity of
that local homotopy for the coupled spectral regulator remain unproved.
\end{theorem}

\begin{proof}
Use Theorems~\ref{thm:v5v28-closed-decomposition},
\ref{thm:v5v28-local-restoration}, and
\ref{thm:v5v28-Gevrey-repair}, followed by
Corollary~\ref{cor:v5v28-defect-budget}.
\end{proof}

V29 will test a concrete covariant spectral shell.  It will prove
equivariance of functional calculus under gauge and diffeomorphism
conjugation, derive the divided-difference formula for a moving spectral
cutoff, and isolate the terms that obstruct an equivariant fixed
low--high splitting.

\clearpage
\section{V29: Covariant spectral shells and the moving-splitting defect}
\label{sec:v5v29}

\subsection{A field-dependent spectral regulator}

V28 leaves open whether a concrete regulator has a local, uniformly
repairable Ward defect.  The natural candidate is functional calculus of a
covariant positive operator.  This note proves its exact equivariance and
then isolates a different obstruction: because the operator depends on the
retained fields, its low and high spectral subspaces move.  Ignoring that
motion creates a false fixed splitting and a spurious Ward defect.

At finite cutoff, let \(\mathcal E_\phi\) be a finite-dimensional Hermitian
field fibre over \(\phi\in\mathcal C_\Lambda\), and let
\begin{equation}
 \Delta_\phi=\Delta_\phi^*\ge0
 \label{eq:v5v29-Delta}
\end{equation}
be a gauge- and diffeomorphism-covariant Laplace-type matrix.  Bundle
automorphisms act by unitary maps
\(U_g:\mathcal E_\phi\to\mathcal E_{g\phi}\) and satisfy
\begin{equation}
 \Delta_{g\phi}=U_g\Delta_\phi U_g^{-1}.
 \label{eq:v5v29-covariance}
\end{equation}
For metric transformations, unitarity is between the metric-dependent
\(L^2\) fibres; a fixed reference-fibre representation requires the
half-density or microlocal normalization already retained in the archived
metric construction.

\begin{theorem}[Equivariance of spectral functional calculus]
\label{thm:v5v29-functional-equivariance}
For every bounded Borel function \(f:[0,\infty)\to\mathbb C\),
\begin{equation}
 f(\Delta_{g\phi})
 =
 U_g f(\Delta_\phi)U_g^{-1}.
 \label{eq:v5v29-f-equivariance}
\end{equation}
In particular, spectral projectors, heat kernels, and heat-shell
differences are equivariant.
\end{theorem}

\begin{proof}
If
\(\Delta_\phi=\sum_\lambda\lambda P_\lambda\), then
\eqref{eq:v5v29-covariance} gives
\[
 \Delta_{g\phi}
 =\sum_\lambda\lambda\,U_gP_\lambda U_g^{-1}.
\]
The spectral theorem therefore gives
\[
 f(\Delta_{g\phi})
 =\sum_\lambda f(\lambda)U_gP_\lambda U_g^{-1}
 =U_gf(\Delta_\phi)U_g^{-1}.
\]
\end{proof}

\begin{corollary}[Invariant spectral traces]
\label{cor:v5v29-traces}
Whenever the indicated finite-cutoff trace is defined,
\begin{equation}
 \operatorname{Tr}f(\Delta_{g\phi})
 =\operatorname{Tr}f(\Delta_\phi).
 \label{eq:v5v29-trace-invariance}
\end{equation}
The same holds for a determinant of an invertible spectral function.
\end{corollary}

\begin{proof}
Trace and determinant are invariant under conjugation.
\end{proof}

\begin{firewall}[Spectral covariance is not measure covariance]
\label{fw:v5v29-measure}
Equation \eqref{eq:v5v29-f-equivariance} controls the operator and its
bosonic spectral fibres.  For a chiral fermion measure, the phases of
bases in those fibres form a determinant line; conjugation of
\(\Delta_\phi\) alone does not prove that its measure Jacobian is trivial.
The inherited anomaly-normalized determinant line remains essential.
\end{firewall}

\subsection{Smooth heat shells}

For \(0<a<b\) and scale \(L>0\), define the positive heat-shell covariance
\begin{equation}
 C_{a,b;L}(\phi)
 =
 e^{-a\Delta_\phi/L^2}-e^{-b\Delta_\phi/L^2}.
 \label{eq:v5v29-heat-shell}
\end{equation}
It is nonnegative because its scalar multiplier is nonnegative on
\([0,\infty)\).

\begin{proposition}[Exact bosonic shell equivariance]
\label{prop:v5v29-bosonic-shell}
The centered finite-dimensional Gaussian measure with covariance
\(C_{a,b;L}(\phi)\), restricted to the support of that covariance, is
carried by \(U_g\) to the Gaussian measure with covariance
\(C_{a,b;L}(g\phi)\).  Hence a bosonic shell integral with a covariant
interaction and equivariant cycle satisfies the measure hypothesis of
Theorem~\ref{thm:v5v27-equivariant-shell}.
\end{proposition}

\begin{proof}
The characteristic function of the first Gaussian is
\[
 \exp\!\left(-\frac12
 \langle \zeta,C_{a,b;L}(\phi)\zeta\rangle\right).
\]
Under the unitary change of variables, this becomes the characteristic
function with covariance
\(U_gC_{a,b;L}(\phi)U_g^{-1}\), which equals
\(C_{a,b;L}(g\phi)\) by
Theorem~\ref{thm:v5v29-functional-equivariance}.  Equality of
characteristic functions proves equality of the Gaussian measures.
\end{proof}

\begin{lemma}[Duhamel derivative of a heat shell]
\label{lem:v5v29-Duhamel}
Let \(t>0\) and let \(B=D\Delta_\phi[\dot\phi]\).  Then
\begin{equation}
 D(e^{-t\Delta_\phi})[\dot\phi]
 =
 -\int_0^t e^{-(t-r)\Delta_\phi}B e^{-r\Delta_\phi}\,dr.
 \label{eq:v5v29-Duhamel}
\end{equation}
Consequently, in operator norm,
\begin{equation}
 \bigl\|D(e^{-t\Delta_\phi})[\dot\phi]\bigr\|
 \le t\|B\|.
 \label{eq:v5v29-Duhamel-bound}
\end{equation}
\end{lemma}

\begin{proof}
Differentiate
\(F(r)=e^{-(t-r)(\Delta_\phi+\varepsilon B)}
e^{-r\Delta_\phi}\) with respect to \(r\), integrate from \(0\) to \(t\),
divide by \(\varepsilon\), and let \(\varepsilon\to0\).  This gives
\eqref{eq:v5v29-Duhamel}.  Since \(\Delta_\phi\ge0\), both heat factors
have norm at most one, and integration gives
\eqref{eq:v5v29-Duhamel-bound}.
\end{proof}

\begin{corollary}[Moving heat-shell bound]
\label{cor:v5v29-heat-shell-derivative}
With \(C_{a,b;L}\) as in \eqref{eq:v5v29-heat-shell},
\begin{equation}
 \|DC_{a,b;L}(\phi)[\dot\phi]\|
 \le\frac{a+b}{L^2}\,
 \|D\Delta_\phi[\dot\phi]\|.
 \label{eq:v5v29-heat-shell-derivative}
\end{equation}
\end{corollary}

\begin{proof}
Apply Lemma~\ref{lem:v5v29-Duhamel} with
\(t=a/L^2\) and \(t=b/L^2\), then use the triangle inequality.
\end{proof}

The factor \(L^{-2}\) is useful only after
\(D\Delta_\phi[\dot\phi]\) is measured between the correct Sobolev
spaces.  Since a Laplace-type operator has order two, the apparent gain
can be exactly spent by its metric variation.  The V21 two-index estimate
must therefore be used rather than reading
\eqref{eq:v5v29-heat-shell-derivative} as an automatic ultraviolet gain.

\subsection{Sharp projectors and their divided differences}

Let
\begin{equation}
 P_\ell(\phi)=\mathbf 1_{[0,\ell)}(\Delta_\phi),\qquad
 Q_\ell(\phi)=I-P_\ell(\phi).
 \label{eq:v5v29-sharp-projector}
\end{equation}
Assume a spectral collar
\begin{equation}
 \operatorname{dist}\bigl(\ell,\operatorname{spec}\Delta_\phi\bigr)
 \ge\gamma>0
 \label{eq:v5v29-gap}
\end{equation}
on the parameter neighborhood under consideration.

\begin{theorem}[Derivative of the moving sharp split]
\label{thm:v5v29-projector-derivative}
Let
\(\Delta_\phi=\sum_\lambda\lambda P_\lambda\) and
\(B=D\Delta_\phi[\dot\phi]\).  Under
\eqref{eq:v5v29-gap},
\begin{equation}
 DP_\ell(\phi)[\dot\phi]
 =
 \sum_{\lambda<\ell\atop\mu>\ell}
 \frac{P_\lambda BP_\mu+P_\mu BP_\lambda}{\lambda-\mu}.
 \label{eq:v5v29-divided-projector}
\end{equation}
In Hilbert--Schmidt norm,
\begin{equation}
 \|DP_\ell(\phi)[\dot\phi]\|_{\mathrm{HS}}
 \le\gamma^{-1}\|B\|_{\mathrm{HS}}.
 \label{eq:v5v29-projector-bound}
\end{equation}
\end{theorem}

\begin{proof}
Choose a contour \(\Gamma\) enclosing the eigenvalues below \(\ell\) and
no eigenvalue above it.  The Riesz formula is
\[
 P_\ell=\frac1{2\pi i}\int_\Gamma(z-\Delta_\phi)^{-1}\,dz.
\]
Differentiation gives a resolvent on each side of \(B\).  Inserting the
spectral resolution and evaluating the contour integral gives zero when
both eigenvalues lie on the same side of \(\ell\), and
\((\lambda-\mu)^{-1}\) in the two cross cases.  This proves
\eqref{eq:v5v29-divided-projector}.  The matrix blocks
\(P_\lambda BP_\mu\) are orthogonal in Hilbert--Schmidt inner product.
Every cross denominator has magnitude at least \(\gamma\), so squaring and
summing yields \eqref{eq:v5v29-projector-bound}.
\end{proof}

\begin{corollary}[The derivative is off diagonal]
\label{cor:v5v29-off-diagonal}
For every differentiable family of orthogonal projections,
\begin{equation}
 P_\ell(DP_\ell)P_\ell=0,\qquad
 Q_\ell(DP_\ell)Q_\ell=0.
 \label{eq:v5v29-off-diagonal}
\end{equation}
Thus motion of the cutoff subspace is entirely low--high mixing.
\end{corollary}

\begin{proof}
Differentiate \(P_\ell^2=P_\ell\):
\[
 (DP_\ell)P_\ell+P_\ell(DP_\ell)=DP_\ell.
\]
Multiply on both sides by \(P_\ell\) to obtain the first identity.
The second follows by applying the first to \(Q_\ell=I-P_\ell\).
\end{proof}

\subsection{The connection term in moving coordinates}

For a field vector \(v(\phi)\in\mathcal E_\phi\), its low coordinate is
\(v_{\mathrm{lo}}=P_\ell(\phi)v(\phi)\).  Differentiation gives
\begin{equation}
 Dv_{\mathrm{lo}}[\dot\phi]
 =
 P_\ell Dv[\dot\phi]
 +(DP_\ell[\dot\phi])v.
 \label{eq:v5v29-moving-coordinate}
\end{equation}
The second summand is absent from a fixed-split calculation.

\begin{proposition}[Size of the moving-split correction]
\label{prop:v5v29-moving-correction}
Under the spectral collar,
\begin{equation}
 \|(DP_\ell[\dot\phi])v\|
 \le
 \gamma^{-1}\|D\Delta_\phi[\dot\phi]\|_{\mathrm{HS}}\|v\|.
 \label{eq:v5v29-moving-correction}
\end{equation}
Any Ward or Lipschitz estimate for a sharp covariant split must include
this low--high term.
\end{proposition}

\begin{proof}
The operator norm is bounded by the Hilbert--Schmidt norm.  Combine this
fact with \eqref{eq:v5v29-projector-bound}.
\end{proof}

\begin{theorem}[Covariance of the moving spectral subbundle]
\label{thm:v5v29-subbundle-covariance}
Equation \eqref{eq:v5v29-covariance} implies
\begin{equation}
 P_\ell(g\phi)U_g=U_gP_\ell(\phi),\qquad
 Q_\ell(g\phi)U_g=U_gQ_\ell(\phi).
 \label{eq:v5v29-projector-covariance}
\end{equation}
Hence the low and high ranges form equivariant vector bundles on every
parameter region with constant spectral rank.  A shell integral over the
high bundle can satisfy V27's equivariance condition, but only when its
bundle density, connection term \eqref{eq:v5v29-moving-coordinate}, and
determinant-line orientation are included.
\end{theorem}

\begin{proof}
Apply Theorem~\ref{thm:v5v29-functional-equivariance} to the indicator
functions defining \(P_\ell\) and \(Q_\ell\).  Constant rank and the
spectral collar give smooth subbundles by
Theorem~\ref{thm:v5v29-projector-derivative}.  The final statement is
exactly the change-of-variables hypothesis of
Theorem~\ref{thm:v5v27-equivariant-shell} expressed on those bundles.
\end{proof}

\begin{firewall}[Fixed projectors generally break covariance]
\label{fw:v5v29-fixed-projector}
Freezing \(P_\ell(\phi)\) at a reference field removes
\(DP_\ell\) but generally destroys
\eqref{eq:v5v29-projector-covariance}.  Keeping a fixed projector is
legitimate only after estimating its commutator with the symmetry action
as a V28 repair defect.  It cannot be called covariant by definition.
\end{firewall}

\subsection{Why a uniform sharp collar can fail}

\begin{proposition}[Rank-change obstruction]
\label{prop:v5v29-rank-change}
Let \(t\mapsto\Delta_t\) be a continuous family of self-adjoint matrices.
If one eigenvalue crosses \(\ell\) between \(t_0\) and \(t_1\), then
\(t\mapsto P_\ell(\Delta_t)\) is discontinuous in operator norm at a
crossing and no positive uniform collar \(\gamma\) can hold on the entire
path.
\end{proposition}

\begin{proof}
The rank of \(P_\ell(\Delta_t)\) changes at the crossing.  Two orthogonal
projections of different finite rank have operator-norm distance one:
the larger range contains a unit vector orthogonal to the smaller range.
Therefore the projector cannot be norm continuous.  At the crossing,
\(\ell\) belongs to the spectrum, so the distance in
\eqref{eq:v5v29-gap} is zero.
\end{proof}

\begin{corollary}[Smooth shells avoid the rank jump]
\label{cor:v5v29-smooth-avoid}
The heat-shell covariance \eqref{eq:v5v29-heat-shell} varies
differentiably whenever \(\Delta_\phi\) does and obeys
\eqref{eq:v5v29-heat-shell-derivative}, even when eigenvalues pass through
a nominal shell scale.
\end{corollary}

\begin{proof}
The exponential is an entire function, and
Lemma~\ref{lem:v5v29-Duhamel} requires no spectral collar.
\end{proof}

Smooth shells replace a direct low--high sum by an overlapping covariance
decomposition.  Constructive independence is recovered by introducing
independent Gaussian increments with those covariances, not by treating
the smooth multipliers as complementary orthogonal projectors.

\subsection{Insertion into the V28 repair budget}

Let \(L_\Delta\) be a tri-graded Lipschitz bound for the field derivative
of \(\Delta_\phi\).  For a sharp split, the moving-coordinate contribution
to the Ward defect has the schematic certified bound
\begin{equation}
 L_A^{\mathrm{sharp}}
 \le C_{\mathrm{geom}}\gamma^{-1}L_\Delta.
 \label{eq:v5v29-sharp-defect}
\end{equation}
For the heat shell,
\begin{equation}
 L_A^{\mathrm{heat}}
 \le C_{\mathrm{geom}}\frac{a+b}{L^2}L_\Delta,
 \label{eq:v5v29-heat-defect}
\end{equation}
where \(C_{\mathrm{geom}}\) includes the interaction insertion and measure
derivatives.  These equations define the constants to be established;
they do not assume that \(L_\Delta\) is cutoff uniform.

\begin{corollary}[Regulator contribution to the contraction margin]
\label{cor:v5v29-margin}
Under the V28 insertion and homotopy bounds, sufficient margins are
\begin{align}
 q+C_{\mathrm{ins}}C_hC_{\mathrm{geom}}
 \gamma^{-1}L_\Delta&<1
 &&\text{for a sharp shell},\label{eq:v5v29-sharp-margin}\\
 q+C_{\mathrm{ins}}C_hC_{\mathrm{geom}}
 \frac{a+b}{L^2}L_\Delta&<1
 &&\text{for a heat shell}.
 \label{eq:v5v29-heat-margin}
\end{align}
\end{corollary}

\begin{proof}
Insert \eqref{eq:v5v29-sharp-defect} or
\eqref{eq:v5v29-heat-defect} into
\eqref{eq:v5v28-final-margin}.
\end{proof}

\begin{firewall}[The conformal problem remains]
\label{fw:v5v29-conformal}
A covariant spectral cutoff organizes the shell and its Ward identities.
It does not make the real Euclidean Einstein action bounded below, prove
metric reflection positivity, or supply the missing ultraviolet fixed
point.  Those independent gates remain in force.
\end{firewall}

\subsection{V29 status}

\begin{theorem}[V29 conclusion]
\label{thm:v5v29-status}
Functional calculus of a covariant positive cutoff operator is exactly
gauge- and diffeomorphism-equivariant.  Smooth heat shells inherit that
equivariance and vary by the Duhamel formula
\eqref{eq:v5v29-Duhamel}.  Sharp spectral subspaces are equivariant only as
moving bundles; their derivative is the off-diagonal divided difference
\eqref{eq:v5v29-divided-projector} and costs the inverse spectral collar.
Consequently a field-independent sharp split is not a valid substitute
without a separately repaired commutator estimate.
\end{theorem}

\begin{proof}
Combine Theorems~\ref{thm:v5v29-functional-equivariance},
\ref{thm:v5v29-projector-derivative}, and
\ref{thm:v5v29-subbundle-covariance} with
Proposition~\ref{prop:v5v29-rank-change}.
\end{proof}

V30 begins with the compulsory companion audit.  It will then use the
smooth covariant heat-shell choice to derive a Duhamel bound for the
variation of one shell Gaussian expectation, including the covariance
derivative and the normalized-measure cancellation.  The aim is a
cutoff-uniform candidate for \(L_A\), rather than another formal
equivariance statement.

\clearpage
\section{V30: Centered Gaussian shell variation and the representation-summed tail}
\label{sec:v5v30}

\subsection{Compulsory companion audit and direction change}

The fifth-version audit used the newest live companion files:
\begin{align*}
 &\texttt{Renewal\_Yang\_Mills\_Mass\_Gap\_Research\_Notes\_V38.tex},
 \qquad
 \operatorname{SHA256}=
 \texttt{74C5155726682E8BAA66494FD43530FF8829B62877FD906753600F3991855527},
 \\
 &\texttt{Renewal\_NS\_Stability\_Research\_Notes\_V31.tex},
 \qquad
 \operatorname{SHA256}=
 \texttt{F1121B62238AD5491BB7CF03635EA9934942EDF573ED8FAAA9EE79E1D38A6696}.
\end{align*}
The YM V38 file is byte-for-byte identical to the simultaneously present
V37 file and its internal record ends at V37; it therefore contributes no
additional mathematical step beyond that record.  No companion file is
modified here.

The YM result constructs a legal saturated five-row support on which the
exact paid connected coefficient is larger than the complete classical
tree target by a factor comparable to the representation dimension.  Thus
a coefficientwise absolute tail norm is not an admissible premise for the
full shell theorem merely because the finite protected sector is kept
signed.  The remaining candidate must retain a representation sum,
Plancherel weight, Schatten norm, or a quantum bridge before taking the
high-label supremum.  The NS V31 result proves an exact heat-formation
first-moment ladder and shows that a further gain must use the complete
signed nonlinear source before its norm is taken.

These companion results do not prove an SM estimate.  They change the
acquisition order in two ways:
\begin{enumerate}
 \item the covariance derivative below is kept as a centered Gaussian
 insertion until after normalization;
 \item the V25 tail norm is now required to be representation-summed rather
 than a classical coefficientwise tree majorant.
\end{enumerate}

\begin{firewall}[Status of the V25 algebra]
\label{fw:v5v30-v25-status}
The abstract V25 hybrid algebra theorem remains correct under its stated
block product hypotheses.  The audit shows that one proposed way of
verifying those hypotheses, a uniform coefficientwise classical quintic
tree estimate, is false in the companion Yang--Mills model.  It must not be
used as an SM--Einstein input.
\end{firewall}

\subsection{Variation of a normalized Gaussian}

Let \(E\) be a finite-dimensional real Hilbert space and let
\(C_\theta>0\) be a differentiable family of covariances.  Write
\(\mu_\theta=N(0,C_\theta)\), let a dot denote
\(\partial_\theta\), and set
\begin{equation}
 Z_{\dot C}(y)
 =
 \frac12\left(
 \langle y,C^{-1}\dot C C^{-1}y\rangle
 -\operatorname{Tr}(C^{-1}\dot C)
 \right).
 \label{eq:v5v30-centered-insertion}
\end{equation}

\begin{lemma}[Centered covariance score]
\label{lem:v5v30-centered-score}
The logarithmic derivative of the normalized Gaussian density is
\(Z_{\dot C}\), and
\begin{equation}
 \mathbb E_{\mu_\theta}Z_{\dot C}=0.
 \label{eq:v5v30-score-zero}
\end{equation}
\end{lemma}

\begin{proof}
The density is
\[
 (2\pi)^{-\dim E/2}(\det C)^{-1/2}
 \exp\!\left(-\frac12\langle y,C^{-1}y\rangle\right).
\]
Use
\[
 \partial_\theta\log\det C=\operatorname{Tr}(C^{-1}\dot C),
 \qquad
 \partial_\theta C^{-1}=-C^{-1}\dot C C^{-1}
\]
to obtain \eqref{eq:v5v30-centered-insertion}.  Since
\(\mathbb E(y\otimes y)=C\),
\[
 \mathbb E\langle y,C^{-1}\dot C C^{-1}y\rangle
 =\operatorname{Tr}(C^{-1}\dot C),
\]
which proves \eqref{eq:v5v30-score-zero}.
\end{proof}

\begin{theorem}[Exact normalized Gaussian variation]
\label{thm:v5v30-Gaussian-variation}
If \(F_\theta\) and its derivative are integrable against \(\mu_\theta\),
then
\begin{equation}
 \frac{d}{d\theta}\mathbb E_{\mu_\theta}F_\theta
 =
 \mathbb E_{\mu_\theta}\dot F_\theta
 +\mathbb E_{\mu_\theta}(F_\theta Z_{\dot C}).
 \label{eq:v5v30-score-formula}
\end{equation}
If \(F\) is twice differentiable with integrable second derivative, then
the covariance term has the inverse-free form
\begin{equation}
 \mathbb E_{\mu_\theta}(F Z_{\dot C})
 =
 \frac12\mathbb E_{\mu_\theta}
 \operatorname{Tr}(\dot C D^2F).
 \label{eq:v5v30-IBP-formula}
\end{equation}
\end{theorem}

\begin{proof}
Differentiate the normalized density using
Lemma~\ref{lem:v5v30-centered-score}; this gives
\eqref{eq:v5v30-score-formula}.  For
\eqref{eq:v5v30-IBP-formula}, choose an orthonormal basis and apply
Gaussian integration by parts twice:
\[
 \mathbb E(y_iy_jF)
 =
 C_{ij}\mathbb EF+
 \sum_{k,\ell}C_{ik}C_{j\ell}
 \mathbb E(\partial_k\partial_\ell F).
\]
Contract with
\(\frac12(C^{-1}\dot C C^{-1})_{ij}\).  The first term is exactly
\(\frac12\operatorname{Tr}(C^{-1}\dot C)\mathbb EF\) and cancels the trace
part of \(Z_{\dot C}\).  The second is
\(\frac12\mathbb E\operatorname{Tr}(\dot C D^2F)\).
\end{proof}

\begin{corollary}[Dimension-free Hilbert--Schmidt estimate]
\label{cor:v5v30-HS-bound}
Under the hypotheses of \eqref{eq:v5v30-IBP-formula},
\begin{equation}
 \left|
 \frac{d}{d\theta}\mathbb E_{\mu_\theta}F_\theta
 \right|
 \le
 \mathbb E|\dot F_\theta|
 +\frac12\|\dot C\|_{\mathcal S_2}
  \mathbb E\|D^2F_\theta\|_{\mathcal S_2}.
 \label{eq:v5v30-HS-bound}
\end{equation}
\end{corollary}

\begin{proof}
Use \eqref{eq:v5v30-IBP-formula} and the Hilbert--Schmidt duality
\(|\operatorname{Tr}(AB)|\le\|A\|_{\mathcal S_2}
\|B\|_{\mathcal S_2}\).
\end{proof}

\begin{assessment}[The cancellation that must precede the norm]
\label{ass:v5v30-centering}
If the two terms in \eqref{eq:v5v30-centered-insertion} are bounded
separately, each carries a trace proportional to the number of shell
modes.  Their exact equality for \(F=1\) would be lost.  Formula
\eqref{eq:v5v30-IBP-formula} performs the normalized-measure cancellation
first and leaves a representation-summed Hilbert--Schmidt pairing.  This is
the Gaussian analogue of the audit lesson.
\end{assessment}

\subsection{Interacting normalized expectations}

Let \(V_\theta:E\to\mathbb R\) be an interaction with
\(0<\mathbb E_{\mu_\theta}e^{-V_\theta}<\infty\), and define
\begin{equation}
 \langle F\rangle_{\theta,V}
 =
 \frac{\mathbb E_{\mu_\theta}(Fe^{-V_\theta})}
 {\mathbb E_{\mu_\theta}e^{-V_\theta}}.
 \label{eq:v5v30-interacting-expectation}
\end{equation}

\begin{theorem}[Connected shell-variation identity]
\label{thm:v5v30-connected-variation}
For differentiable \(F_\theta,V_\theta,C_\theta\),
\begin{equation}
 \frac{d}{d\theta}\langle F\rangle_{\theta,V}
 =
 \langle\dot F\rangle_{\theta,V}
 -\operatorname{Cov}_{\theta,V}(F,\dot V)
 +\operatorname{Cov}_{\theta,V}(F,Z_{\dot C}),
 \label{eq:v5v30-connected-score}
\end{equation}
where
\(\operatorname{Cov}(F,G)=\langle FG\rangle-\langle F\rangle\langle
G\rangle\).  If the required field derivatives are integrable, the last
term equals
\begin{align}
 \operatorname{Cov}_{\theta,V}(F,Z_{\dot C})
 &=
 \frac12\left\langle
 \operatorname{Tr}\dot C
 \bigl(D^2F-DF\otimes DV-DV\otimes DF\bigr)
 \right\rangle_{\theta,V}
 \nonumber\\
 &\quad+
 \frac12\operatorname{Cov}_{\theta,V}
 \left(
 F,\operatorname{Tr}\dot C
 \bigl(DV\otimes DV-D^2V\bigr)
 \right).
 \label{eq:v5v30-connected-IBP}
\end{align}
\end{theorem}

\begin{proof}
Differentiate the numerator and denominator in
\eqref{eq:v5v30-interacting-expectation} using
\eqref{eq:v5v30-score-formula}, then apply the quotient rule.  The
interaction derivative gives
\(-\operatorname{Cov}(F,\dot V)\), and the score gives the last covariance
in \eqref{eq:v5v30-connected-score}.

For the inverse-free formula, apply
\eqref{eq:v5v30-IBP-formula} first to \(Fe^{-V}\) and then to \(e^{-V}\).
The product rules give
\begin{align*}
 e^VD^2(Fe^{-V})
 &=D^2F-DF\otimes DV-DV\otimes DF\\
 &\quad+F(DV\otimes DV-D^2V),\\
 e^VD^2(e^{-V})
 &=DV\otimes DV-D^2V.
\end{align*}
Substitution and subtraction of the denominator derivative give
\eqref{eq:v5v30-connected-IBP}.
\end{proof}

\begin{corollary}[A connected Schatten bound]
\label{cor:v5v30-connected-bound}
Let
\[
 H_V=DV\otimes DV-D^2V.
\]
For an \(\mathcal S_2\)-valued random operator, write
\[
 \operatorname{Var}_{\mathcal S_2}(H_V)
 =
 \left\langle
 \|H_V-\langle H_V\rangle_{\theta,V}\|_{\mathcal S_2}^2
 \right\rangle_{\theta,V}.
\]
Then the covariance part of \eqref{eq:v5v30-connected-score} obeys
\begin{align}
 |\operatorname{Cov}_{\theta,V}(F,Z_{\dot C})|
 &\le\frac12\|\dot C\|_{\mathcal S_2}
 \Bigl\langle
 \|D^2F\|_{\mathcal S_2}
 +2\|DF\|\,\|DV\|
 \Bigr\rangle_{\theta,V}
 \nonumber\\
 &\quad+
 \frac12\|\dot C\|_{\mathcal S_2}
 \sqrt{\operatorname{Var}_{\theta,V}(F)}
 \sqrt{\operatorname{Var}_{\mathcal S_2}(H_V)}.
 \label{eq:v5v30-connected-bound}
\end{align}
\end{corollary}

\begin{proof}
Apply Hilbert--Schmidt duality to the first line of
\eqref{eq:v5v30-connected-IBP}.  For the second line, first bound the trace
by
\(\|\dot C\|_{\mathcal S_2}\|H_V\|_{\mathcal S_2}\), then apply
Cauchy--Schwarz after centering both \(F\) and the
\(\mathcal S_2\)-valued operator \(H_V\).
\end{proof}

\begin{firewall}[No automatic perturbative gain]
\label{fw:v5v30-no-free-gain}
The right side of \eqref{eq:v5v30-connected-bound} still requires moment
bounds for \(DV,D^2V\) and the observable derivatives.  Heat smoothing
does not create an extra critical scale factor after the order of
\(D\Delta\) and the interaction insertions are paid.  This is consistent
with the sharp NS formation ladder in the audit.
\end{firewall}

\subsection{Application to the covariant heat shell}

Take \(C_\theta=C_{a,b;L}(\phi_\theta)\) from
\eqref{eq:v5v29-heat-shell}.  Corollary
\ref{cor:v5v29-heat-shell-derivative} gives
\begin{equation}
 \|\dot C_\theta\|_{\mathcal S_2}
 \le\frac{a+b}{L^2}
 \|D\Delta_{\phi_\theta}[\dot\phi_\theta]\|_{\mathcal S_2}.
 \label{eq:v5v30-heat-HS}
\end{equation}

\begin{theorem}[One-shell variation certificate]
\label{thm:v5v30-one-shell}
Suppose, uniformly over the admissible retained-field ball,
\begin{align}
 \frac1{L^2}
 \|D\Delta_\phi[\dot\phi]\|_{\mathcal S_2}
 &\le K_\Delta\|\dot\phi\|_+,
 \label{eq:v5v30-KDelta}\\
 \left\langle
 \|D^2F\|_{\mathcal S_2}
 +2\|DF\|\|DV\|
 \right\rangle_V
 &\le K_{F,V},
 \label{eq:v5v30-KFV}\\
 \sqrt{\operatorname{Var}_V(F)}
 \sqrt{\operatorname{Var}_{\mathcal S_2}(H_V)}
 &\le K_{\mathrm{conn}}.
 \label{eq:v5v30-Kconn}
\end{align}
Then the covariance part of the retained-field derivative satisfies
\begin{equation}
 |\delta_C\langle F\rangle_V|
 \le
 \frac{a+b}{2}K_\Delta
 (K_{F,V}+K_{\mathrm{conn}})
 \|\dot\phi\|_+.
 \label{eq:v5v30-one-shell-bound}
\end{equation}
\end{theorem}

\begin{proof}
Insert \eqref{eq:v5v30-heat-HS} and
\eqref{eq:v5v30-KDelta}--\eqref{eq:v5v30-Kconn} into
\eqref{eq:v5v30-connected-bound}.
\end{proof}

The theorem identifies a candidate contribution to the V28 constant
\(L_A\).  Its hypotheses are deliberately representation-summed.  In four
dimensions the raw Hilbert--Schmidt norm in
\eqref{eq:v5v30-KDelta} can still grow with the shell volume, so the
covariant Taylor subtraction and rooted localization of V6 must be used
before calling \(K_\Delta\) uniform.

\subsection{Revised tail certificate}

\begin{definition}[Representation-summed tail certificate]
\label{def:v5v30-RSTC}
An RSTC for one shell consists of:
\begin{enumerate}
 \item a decomposition into complete Ward and cumulant packets before
 taking norms;
 \item a Plancherel-weighted Hilbert or Schatten norm on each packet,
 retaining every representation sum used by orthogonality;
 \item a rooted spatial summation after the packet norm;
 \item tri-graded factorial weights after the representation and rooted
 sums;
 \item product, covariance-variation, and counterterm-insertion estimates
 in that order;
 \item constants uniform in cutoff, shell volume, representation label,
 curvature grade, and particle number.
\end{enumerate}
\end{definition}

\begin{proposition}[Schatten multiplication routes]
\label{prop:v5v30-Schatten-products}
For finite-cutoff operators,
\begin{align}
 \|AB\|_{\mathcal S_1}
 &\le\|A\|_{\mathcal S_2}\|B\|_{\mathcal S_2},
 \label{eq:v5v30-S2S2}\\
 \|AB\|_{\mathcal S_2}
 &\le\|A\|_{\mathrm{op}}\|B\|_{\mathcal S_2}.
 \label{eq:v5v30-opS2}
\end{align}
Thus a representation-summed algebra can close either through an
\(\mathcal S_2\times\mathcal S_2\to\mathcal S_1\) two-tier norm or through
an operator-norm multiplier acting on an \(\mathcal S_2\) tail.  It need
not imply a coefficientwise bound.
\end{proposition}

\begin{proof}
These are the Schatten Hölder inequalities with
\(1=1/2+1/2\) and \(1/2=0+1/2\), respectively.
\end{proof}

\begin{assessment}[What the audit rules out and leaves open]
\label{ass:v5v30-revised-tail}
The audited YM example rules out the specified classical coefficientwise
quintic marked-tree target.  It does not rule out either route in
Proposition~\ref{prop:v5v30-Schatten-products}.  Conversely, the proposition
alone does not prove the required rooted graph sum; the exact RSTC remains
an SM--Einstein shell estimate to be established.
\end{assessment}

\subsection{V30 conclusion}

\begin{theorem}[V30 result]
\label{thm:v5v30-status}
Variation of a normalized Gaussian shell has the exact centered form
\eqref{eq:v5v30-score-formula} and the inverse-free,
representation-summed form \eqref{eq:v5v30-IBP-formula}.  For an interacting
normalized expectation, all covariance dependence is the connected
identity \eqref{eq:v5v30-connected-IBP}; under
\eqref{eq:v5v30-KDelta}--\eqref{eq:v5v30-Kconn} it obeys the uniform
candidate bound \eqref{eq:v5v30-one-shell-bound}.  The companion audit
requires replacing any coefficientwise classical tail premise by an RSTC.
No RSTC or uniform moment bounds have yet been proved for the complete
coupled shell.
\end{theorem}

\begin{proof}
Use Theorems~\ref{thm:v5v30-Gaussian-variation},
\ref{thm:v5v30-connected-variation}, and
\ref{thm:v5v30-one-shell}, together with the audit counterexample recorded
at the start of this note.
\end{proof}

V31 will attack the first uniform estimate in
\eqref{eq:v5v30-KDelta}.  It will subtract the local heat coefficients
before taking the rooted Schatten norm and determine the exact
four-dimensional shell scaling of the remaining covariance derivative.

\clearpage
\section{V31: Failure of global Schatten uniformity and the rooted heat-shell repair}
\label{sec:v5v31}

\subsection{Why the V30 Hilbert--Schmidt target is too strong}

V30 proposed a cutoff-uniform Hilbert--Schmidt estimate for the derivative
of a covariant heat shell.  The simplest covariant family disproves that
target.  This is not a defect of heat regularization; it is the ordinary
extensive number of modes in a four-dimensional shell.  The correction is
to separate the internal representation norm from the rooted spatial
summation, as in the V6 polymer space.

Let \(\mathbb T^4\) be a flat torus of fixed positive volume, let
\(\Delta=-\partial^2\), and consider the scaling family
\begin{equation}
 \Delta_\theta=(1+\theta)\Delta.
 \label{eq:v5v31-scaling-family}
\end{equation}
For \(0<a<b\), define
\[
 C_L(\theta)
 =
 e^{-a\Delta_\theta/L^2}-e^{-b\Delta_\theta/L^2}.
\]
At \(\theta=0\),
\begin{equation}
 \dot C_L=h(\Delta/L^2),\qquad
 h(x)=-ax e^{-ax}+bx e^{-bx}.
 \label{eq:v5v31-shell-multiplier}
\end{equation}

\begin{theorem}[Global Schatten no-go in four dimensions]
\label{thm:v5v31-Schatten-no-go}
For every finite \(1\le p<\infty\), there are constants
\(c_p,C_p>0\) such that, for all sufficiently large \(L\),
\begin{equation}
 c_pL^{4/p}
 \le\|\dot C_L\|_{\mathcal S_p}
 \le C_pL^{4/p}.
 \label{eq:v5v31-Schatten-scaling}
\end{equation}
In particular,
\(\|\dot C_L\|_{\mathcal S_2}\asymp L^2\), so the global
Hilbert--Schmidt constant in \eqref{eq:v5v30-KDelta} cannot be cutoff
uniform.
\end{theorem}

\begin{proof}
The function \(h\) is continuous, rapidly decreasing, and not identically
zero because \(a\ne b\).  Choose a compact interval
\(I\subset(0,\infty)\) and \(c>0\) such that \(|h(x)|\ge c\) on \(I\).
The torus Weyl law gives
\[
 \#\{\lambda\in\operatorname{spec}\Delta:\lambda/L^2\in I\}
 \asymp L^4.
\]
Each corresponding singular value of \(h(\Delta/L^2)\) is at least \(c\),
which proves the lower bound.

For the upper bound, rapid decay gives
\(|h(x)|^p\le C_N(1+x)^{-N}\) for any \(N\).  Choose \(N>2\).  The Weyl
counting bound \(N_\Delta(R)\le C(1+R)^2\), followed by a dyadic spectral
sum, gives
\[
 \sum_\lambda|h(\lambda/L^2)|^p\le C_pL^4.
\]
Taking the \(p\)-th root proves the upper bound.
\end{proof}

\begin{corollary}[Operator norm is the nonextensive endpoint]
\label{cor:v5v31-op-endpoint}
\begin{equation}
 \|\dot C_L\|_{\mathrm{op}}
 \le\sup_{x\ge0}|h(x)|
 \label{eq:v5v31-op-bound}
\end{equation}
uniformly in \(L\).
\end{corollary}

\begin{proof}
Apply the spectral theorem to \eqref{eq:v5v31-shell-multiplier}.
\end{proof}

\begin{firewall}[Local heat coefficients do not change this operator count]
\label{fw:v5v31-local-not-Schatten}
The local heat coefficients summarized in V1 subtract divergent local
terms from traces and effective actions.  They do not alter the singular
values of the covariance derivative \(\dot C_L\).  Therefore they cannot
turn \eqref{eq:v5v31-Schatten-scaling} into a uniform global Schatten
bound unless the regulator itself is changed.
\end{firewall}

\subsection{Finite-rank subtractions cannot remove the bulk}

The preceding conclusion remains true after any correction that touches
only a negligible number of shell modes.

\begin{proposition}[Rank lower bound]
\label{prop:v5v31-rank-lower}
Let \(N_L\ge cL^4\) singular values of \(\dot C_L\) be at least \(c_0>0\),
as in the proof of Theorem~\ref{thm:v5v31-Schatten-no-go}.  If
\(\operatorname{rank}R_L=r_L<N_L\), then
\begin{equation}
 \|\dot C_L-R_L\|_{\mathcal S_p}
 \ge c_0(N_L-r_L)^{1/p}.
 \label{eq:v5v31-rank-lower}
\end{equation}
\end{proposition}

\begin{proof}
The singular-value approximation inequality says that subtraction of a
rank-\(r_L\) operator cannot remove more than the first \(r_L\) singular
directions:
\[
 s_k(\dot C_L-R_L)\ge s_{k+r_L}(\dot C_L).
\]
For \(1\le k\le N_L-r_L\), the right side is at least \(c_0\).  Sum the
\(p\)-th powers.
\end{proof}

\begin{corollary}[No finite-head repair]
\label{cor:v5v31-no-finite-head}
If \(r_L=o(L^4)\), then
\(\|\dot C_L-R_L\|_{\mathcal S_p}\gtrsim L^{4/p}\).
Thus the finite head and protected blocks of V25 cannot absorb the
extensive covariance derivative.
\end{corollary}

\begin{proof}
Insert \(N_L\asymp L^4\) and \(r_L=o(L^4)\) into
\eqref{eq:v5v31-rank-lower}.
\end{proof}

\subsection{The rooted spatial norm}

Return first to \(\mathbb R^4\).  Let \(h\) be a Schwartz function of
\(|\xi|^2\), and define the Fourier multiplier
\begin{equation}
 T_L=h(-\partial^2/L^2).
 \label{eq:v5v31-TL}
\end{equation}
With a fixed Fourier convention, its kernel has the scaling form
\begin{equation}
 K_L(x,y)=L^4\kappa(L(x-y)),
 \qquad
 \kappa=\mathcal F^{-1}\bigl(h(|\cdot|^2)\bigr).
 \label{eq:v5v31-kernel-scaling}
\end{equation}
Define the rooted moment norm
\begin{equation}
 \|K\|_{\mathrm{root},r}
 =
 \sup_x\int_{\mathbb R^4}|x-y|^r|K(x,y)|\,dy,
 \qquad r\ge0.
 \label{eq:v5v31-root-norm}
\end{equation}

\begin{theorem}[Exact rooted heat-shell scaling]
\label{thm:v5v31-root-scaling}
For every \(r\ge0\),
\begin{equation}
 \|K_L\|_{\mathrm{root},r}
 =
 L^{-r}\int_{\mathbb R^4}|z|^r|\kappa(z)|\,dz.
 \label{eq:v5v31-root-scaling}
\end{equation}
In particular, the unmarked rooted norm is uniform in \(L\), and every
spatial Taylor remainder of order \(r\) gains \(L^{-r}\).
\end{theorem}

\begin{proof}
Insert \eqref{eq:v5v31-kernel-scaling} into
\eqref{eq:v5v31-root-norm} and use \(z=L(x-y)\):
\[
 \int |x-y|^rL^4|\kappa(L(x-y))|\,dy
 =
 L^{-r}\int|z|^r|\kappa(z)|\,dz.
\]
The expression is independent of \(x\).  All moments are finite because
\(\kappa\) is Schwartz.
\end{proof}

\begin{corollary}[Periodic rooted bound]
\label{cor:v5v31-periodic-root}
On a flat torus, periodize \(K_L\).  Then
\begin{equation}
 \sup_x\int_{\mathbb T^4}|K_L^{\mathrm{per}}(x,y)|\,dy
 \le\|\kappa\|_{L^1(\mathbb R^4)}
 \label{eq:v5v31-periodic-root}
\end{equation}
uniformly in \(L\).
\end{corollary}

\begin{proof}
Write the periodized kernel as the sum over lattice translates, apply the
triangle inequality, and unfold the fundamental domain to
\(\mathbb R^4\).
\end{proof}

For vector bundles, replace the scalar absolute value inside
\eqref{eq:v5v31-root-norm} by the Hilbert--Schmidt norm on the
finite-dimensional internal fibre.  This retains the representation sum
before the spatial root integral and matches the ordering required by the
V30 audit.

\subsection{Taylor subtraction and local owners}

Let \(G(x,y)\) be the kernel paired with \(K_L(x,y)\) in the connected
Gaussian variation.  For an integer \(r\ge1\), let
\(\operatorname{Tay}_{x}^{r-1}G(x,y)\) be its covariant Taylor polynomial
in \(y\) about \(x\), with parallel transport used to compare fibres.
Write
\begin{equation}
 \mathcal R_x^rG(x,y)
 =
 G(x,y)-\operatorname{Tay}_{x}^{r-1}G(x,y).
 \label{eq:v5v31-Taylor-remainder}
\end{equation}

\begin{theorem}[Rooted Taylor gain]
\label{thm:v5v31-Taylor-gain}
Assume
\begin{equation}
 \|\mathcal R_x^rG(x,y)\|_{\mathcal S_2}
 \le M_r|x-y|^r
 \label{eq:v5v31-remainder-bound}
\end{equation}
on the localization chart.  Then
\begin{equation}
 \sup_x\left|
 \int\operatorname{tr}
 \bigl(K_L(x,y)\mathcal R_x^rG(y,x)\bigr)\,dy
 \right|
 \le
 M_rL^{-r}
 \int|z|^r\|\kappa(z)\|_{\mathcal S_2}\,dz.
 \label{eq:v5v31-Taylor-gain}
\end{equation}
\end{theorem}

\begin{proof}
Use Hilbert--Schmidt duality in the internal fibre, then
\eqref{eq:v5v31-remainder-bound}:
\[
 |\operatorname{tr}(K_L\mathcal R_x^rG)|
 \le\|K_L\|_{\mathcal S_2}\|\mathcal R_x^rG\|_{\mathcal S_2}
 \le M_r|x-y|^r\|K_L(x,y)\|_{\mathcal S_2}.
\]
Integrate in \(y\), take the supremum in \(x\), and apply
Theorem~\ref{thm:v5v31-root-scaling}.
\end{proof}

\begin{assessment}[Ownership of the Taylor polynomial]
\label{ass:v5v31-local-owner}
The Taylor polynomial removed in
\eqref{eq:v5v31-Taylor-remainder} is local at the root.  Its contractions
are assigned to the cosmological, Einstein--Hilbert, gauge, Higgs,
nonminimal, curvature-tower, and boundary owners already identified in the
archived heat analysis.  V31 does not recompute those coefficients.  Its
new statement is the explicit \(L^{-r}\) bound for the rooted nonlocal
remainder.
\end{assessment}

\begin{firewall}[Subtraction must match the exact contraction]
\label{fw:v5v31-exact-subtraction}
A formal Taylor polynomial does not cancel a shell contribution unless the
counterterm insertion uses the same covariant parallel transport,
normalization, boundary convention, and complete Ward packet.  Mismatched
subtractions leave a local V28 defect.
\end{firewall}

\subsection{A corrected covariance-variation certificate}

For a bundle kernel \(A(x,y)\), define
\begin{equation}
 \|A\|_{\mathfrak R_r}
 =
 \sup_x\int d(x,y)^r
 \|A(x,y)\|_{\mathcal S_2(\mathrm{int})}\,dy,
 \label{eq:v5v31-bundle-root}
\end{equation}
with a partition of unity and bounded-overlap chart atlas on a curved
finite-volume background.  Uniform bounded geometry controls the
equivalence of chart norms.

\begin{definition}[Rooted covariance derivative certificate]
\label{def:v5v31-RCDC}
An RCDC of order \(r\) is a bound
\begin{equation}
 \|D C_L(\phi)[\dot\phi]\|_{\mathfrak R_r}
 \le C_rL^{-r}\|\dot\phi\|_+
 \label{eq:v5v31-RCDC}
\end{equation}
after the local Taylor owners of order below \(r\) have been extracted,
with \(C_r\) uniform in cutoff, volume chart, local grade, and internal
representation packet.
\end{definition}

\begin{theorem}[Flat heat shell satisfies the rooted model certificate]
\label{thm:v5v31-flat-RCDC}
For the scaling family \eqref{eq:v5v31-scaling-family}, the derivative
kernel \(\dot C_L\) satisfies
\begin{equation}
 \|\dot C_L\|_{\mathfrak R_r}
 \le C_{a,b,r}L^{-r}
 \label{eq:v5v31-flat-RCDC}
\end{equation}
for every \(r\ge0\), both on \(\mathbb R^4\) and, with the periodic
distance convention and a harmless chart constant, on the flat torus.
\end{theorem}

\begin{proof}
The multiplier \(h\) in \eqref{eq:v5v31-shell-multiplier} is Schwartz as a
function of momentum.  Apply Theorem~\ref{thm:v5v31-root-scaling} and its
periodic corollary.
\end{proof}

\begin{firewall}[Curved interacting RCDC remains open]
\label{fw:v5v31-curved-open}
The flat theorem does not prove the RCDC for a rough dynamical metric.
That requires uniform bounded-geometry or parametrix estimates for the
metric-dependent heat kernel, control of \(D\Delta_\phi\) in the V21
two-index scale, and Ward-complete local subtractions.  Concentrating
metrics can destroy the constants.
\end{firewall}

\subsection{Consequences for the V30 estimate}

The factorization in \eqref{eq:v5v30-HS-bound} is valid at each finite
cutoff, but its two global Hilbert--Schmidt factors need not be uniformly
bounded separately.  The exact trace pairing must instead be estimated
after spatial rooting and internal representation summation.

\begin{proposition}[Rooted pairing bound]
\label{prop:v5v31-rooted-pairing}
Let \(A(x,y)\) and \(G(y,x)\) be internal Hilbert--Schmidt kernels.  Then
\begin{equation}
 \sup_x\left|
 \int\operatorname{tr}(A(x,y)G(y,x))\,dy
 \right|
 \le
 \left(\sup_x\int\|A(x,y)\|_{\mathcal S_2}\,dy\right)
 \sup_{x,y}\|G(y,x)\|_{\mathcal S_2}.
 \label{eq:v5v31-rooted-pairing}
\end{equation}
With an order-\(r\) Taylor remainder, the right side gains \(L^{-r}\)
under the RCDC.
\end{proposition}

\begin{proof}
Apply Hilbert--Schmidt duality pointwise under the integral.  For the last
claim use Theorem~\ref{thm:v5v31-Taylor-gain}.
\end{proof}

\begin{corollary}[Corrected candidate for \(L_A\)]
\label{cor:v5v31-corrected-LA}
If the connected interacting Hessian kernels in
\eqref{eq:v5v30-connected-IBP} have a uniform rooted dual bound
\(K_{\mathrm{conn},r}\), and the RCDC holds, then the nonlocal covariance
contribution to the Ward-defect Lipschitz constant obeys
\begin{equation}
 L_A^{\mathrm{nonloc}}
 \le\frac12 C_rK_{\mathrm{conn},r}L^{-r}.
 \label{eq:v5v31-corrected-LA}
\end{equation}
The extracted Taylor terms belong to the local V28 repair problem.
\end{corollary}

\begin{proof}
Apply Proposition~\ref{prop:v5v31-rooted-pairing} to each term of the
connected identity \eqref{eq:v5v30-connected-IBP}, after its common Taylor
forest has been extracted.
\end{proof}

\subsection{V31 status}

\begin{theorem}[V31 conclusion]
\label{thm:v5v31-status}
No finite Schatten norm of a four-dimensional heat-shell covariance
derivative is globally cutoff uniform, even for a flat scaling family; its
exact growth is \(L^{4/p}\).  The corresponding rooted kernel norm is
uniform, and an order-\(r\) covariant Taylor remainder gains exactly
\(L^{-r}\).  Hence the viable shell certificate pairs a rooted spatial
norm with an internal representation-summed Schatten norm and assigns the
Taylor polynomial to the established local counterterm owners.  The RCDC
for rough interacting metrics remains unproved.
\end{theorem}

\begin{proof}
Combine Theorem~\ref{thm:v5v31-Schatten-no-go},
Theorem~\ref{thm:v5v31-root-scaling}, and
Theorem~\ref{thm:v5v31-Taylor-gain}.
\end{proof}

V32 will attack the curved RCDC by a resolvent-parametrix identity.  It
will express the difference between two metric-dependent heat shells
through one insertion of \(\Delta_g-\Delta_{\tilde g}\), keep two-sided
heat localization, and identify the bounded-geometry quantities that must
be controlled in the strong tri-graded space.

\clearpage
\section{V32: Curved heat-shell variation in the rooted kernel norm}
\label{sec:v5v32}

\subsection{Uniform geometric hypotheses}

V31 proves the rooted covariance derivative certificate on a flat
background.  This note proves its curved analogue under explicit heat
kernel and coefficient hypotheses.  The hypotheses are strong enough to
show where metric concentration enters and weak enough to retain the
principal second-order variation.

Let \(M\) be a closed \(d\)-dimensional manifold, or a member of a family
with uniform bounded geometry.  Identify metric-dependent \(L^2\) fibres
by the half-density and bundle isometries retained in the archived
microlocal normalization.  Let \(\Delta_\phi\ge0\) be a self-adjoint
Laplace-type operator with heat kernel \(K_t^\phi(x,y)\).  Assume, for
\(i+j\le1\),
\begin{equation}
 \|\nabla_x^i\nabla_y^jK_t^\phi(x,y)\|_{\mathcal S_2(\mathrm{int})}
 \le
 C_{ij}t^{-(d+i+j)/2}
 \exp\!\left(-\frac{d(x,y)^2}{c t}\right)
 \label{eq:v5v32-Gaussian}
\end{equation}
for \(0<t\le t_0\), uniformly over the admissible field ball.  The
finite-dimensional internal Hilbert--Schmidt norm is taken before the
spatial integration.

For a kernel \(K\), retain the rooted moment norm
\[
 \|K\|_{\mathfrak R_r}
 =
 \sup_x\int_Md(x,y)^r
 \|K(x,y)\|_{\mathcal S_2(\mathrm{int})}\,dy .
\]

\begin{lemma}[Heat-kernel rooted moments]
\label{lem:v5v32-heat-moments}
Under \eqref{eq:v5v32-Gaussian}, for \(k=0,1\) and every \(r\ge0\),
\begin{equation}
 \|\nabla^kK_t^\phi\|_{\mathfrak R_r}
 \le C_{r,k}t^{(r-k)/2}.
 \label{eq:v5v32-heat-moments}
\end{equation}
\end{lemma}

\begin{proof}
Uniform volume comparison under bounded geometry reduces the integral to
\[
 Ct^{-(d+k)/2}\int_0^\infty
 \rho^{d-1+r}e^{-\rho^2/(ct)}\,d\rho.
\]
The change of variable \(\rho=\sqrt t\,u\) gives
\(Ct^{(r-k)/2}\) times a finite Gaussian moment.  The same estimate holds
for a derivative in either kernel variable.
\end{proof}

\subsection{Rooted convolution calculus}

If \(A\) and \(B\) are bundle kernels, write
\[
 (A*B)(x,y)=\int_MA(x,z)B(z,y)\,dz.
\]

\begin{lemma}[Weighted rooted convolution]
\label{lem:v5v32-root-convolution}
For \(r\ge0\), set \(c_r=1\) for \(0\le r\le1\) and
\(c_r=2^{r-1}\) for \(r\ge1\).  Then
\begin{equation}
 \|A*B\|_{\mathfrak R_r}
 \le c_r\bigl(
 \|A\|_{\mathfrak R_r}\|B\|_{\mathfrak R_0}
 +\|A\|_{\mathfrak R_0}\|B\|_{\mathfrak R_r}
 \bigr).
 \label{eq:v5v32-root-convolution}
\end{equation}
\end{lemma}

\begin{proof}
The triangle inequality for the metric and the scalar inequality
\[
 (u+v)^r\le c_r(u^r+v^r)
\]
give
\(d(x,y)^r\le c_r(d(x,z)^r+d(z,y)^r)\).  Internal
Hilbert--Schmidt/operator Hölder bounds the norm of the kernel product.
Tonelli's theorem and successive integration in \(y\) and \(z\) give
\eqref{eq:v5v32-root-convolution}.  If both internal factors are only
Hilbert--Schmidt, their product is trace class and its trace norm obeys the
same estimate; the chosen RSTC must record which of these two tiers is
used.
\end{proof}

\subsection{Divergence-form perturbations}

Let \(\phi,\widetilde\phi\) be two admissible retained fields.  After the
common Hilbert-fibre identification, assume
\begin{equation}
 B:=\Delta_\phi-\Delta_{\widetilde\phi}
 =
 \nabla^*(A_2\nabla)+A_1\nabla+A_0
 \label{eq:v5v32-divergence-B}
\end{equation}
in local covariant notation.  Let
\begin{equation}
 \varepsilon_k=\|A_k\|_{L^\infty},\qquad k=0,1,2,
 \label{eq:v5v32-epsilons}
\end{equation}
including the bounded changes of connection and volume density in these
coefficients.

\begin{lemma}[Two-operator Duhamel identity]
\label{lem:v5v32-two-Duhamel}
For nonnegative self-adjoint \(\Delta_\phi,\Delta_{\widetilde\phi}\),
\begin{equation}
 e^{-t\Delta_\phi}-e^{-t\Delta_{\widetilde\phi}}
 =
 -\int_0^t
 e^{-(t-s)\Delta_\phi}
 (\Delta_\phi-\Delta_{\widetilde\phi})
 e^{-s\Delta_{\widetilde\phi}}\,ds.
 \label{eq:v5v32-two-Duhamel}
\end{equation}
\end{lemma}

\begin{proof}
Differentiate
\(F(s)=e^{-(t-s)\Delta_\phi}e^{-s\Delta_{\widetilde\phi}}\)
and integrate \(F'\) from \(0\) to \(t\).  Rearrangement gives the stated
sign.
\end{proof}

\begin{theorem}[Curved rooted heat-difference bound]
\label{thm:v5v32-curved-difference}
Under \eqref{eq:v5v32-Gaussian}--\eqref{eq:v5v32-divergence-B}, for every
\(r\ge0\) and \(0<t\le t_0\),
\begin{equation}
 \|e^{-t\Delta_\phi}-e^{-t\Delta_{\widetilde\phi}}\|_{\mathfrak R_r}
 \le
 C_r\left(
 \varepsilon_2t^{r/2}
 +\varepsilon_1t^{(r+1)/2}
 +\varepsilon_0t^{(r+2)/2}
 \right).
 \label{eq:v5v32-curved-difference}
\end{equation}
\end{theorem}

\begin{proof}
Insert \eqref{eq:v5v32-divergence-B} into
\eqref{eq:v5v32-two-Duhamel}.  For the principal term, integrate the
divergence derivative by parts in the intermediate variable.  Its kernel
is a sum of terms bounded by
\[
 \varepsilon_2
 \int_0^t
 |\nabla K_{t-s}^{\phi}|*|\nabla K_s^{\widetilde\phi}|\,ds.
\]
There is no boundary term on the closed manifold.  Apply
Lemmas~\ref{lem:v5v32-heat-moments} and
\ref{lem:v5v32-root-convolution}.  The two contributions to the
\(r\)-moment are bounded by constants times
\begin{align*}
 \varepsilon_2\int_0^t
 (t-s)^{(r-1)/2}s^{-1/2}\,ds
 &=
 C_r\varepsilon_2t^{r/2},\\
 \varepsilon_2\int_0^t
 (t-s)^{-1/2}s^{(r-1)/2}\,ds
 &=
 C_r\varepsilon_2t^{r/2}.
\end{align*}
Both integrals are finite beta integrals for every \(r\ge0\).

For the first-order term, put its single derivative on either heat kernel.
The analogous beta integrals scale as \(t^{(r+1)/2}\).  With no derivative,
the zeroth-order term scales as \(t^{(r+2)/2}\).  Sum the three estimates.
\end{proof}

\begin{assessment}[Why divergence form is decisive]
\label{ass:v5v32-divergence}
Estimating \(A_2\nabla^2e^{-s\Delta}\) on only one side produces the
nonintegrable endpoint majorant \(s^{-1}\).  Integration by parts splits
the principal derivatives between the two heat kernels and replaces it by
\((t-s)^{-1/2}s^{-1/2}\), whose integral is finite.  This is the exact
two-sided smoothing required by the V21 index ledger.
\end{assessment}

\begin{firewall}[Boundary terms]
\label{fw:v5v32-boundary}
On artificial core boundaries, the integration by parts in
Theorem~\ref{thm:v5v32-curved-difference} produces boundary kernels.
They must be assigned to the half-integer boundary heat owners summarized
in V1 and estimated in boundary rooted norms.  The closed-manifold proof
does not silently remove them.
\end{firewall}

\subsection{The curved heat-shell RCDC}

Define
\[
 C_{a,b;L}(\phi)
 =
 e^{-a\Delta_\phi/L^2}-e^{-b\Delta_\phi/L^2}.
\]

\begin{corollary}[Lipschitz curved RCDC]
\label{cor:v5v32-curved-RCDC}
For \(L\) large enough that \(b/L^2\le t_0\),
\begin{align}
 \|C_{a,b;L}(\phi)-C_{a,b;L}(\widetilde\phi)\|_{\mathfrak R_r}
 \le C_{a,b,r}\bigl(
 &\varepsilon_2L^{-r}
 +\varepsilon_1L^{-(r+1)}
 \nonumber\\
 &+\varepsilon_0L^{-(r+2)}
 \bigr).
 \label{eq:v5v32-curved-RCDC}
\end{align}
\end{corollary}

\begin{proof}
Apply Theorem~\ref{thm:v5v32-curved-difference} at
\(t=a/L^2\) and \(t=b/L^2\), then add the bounds.
\end{proof}

\begin{corollary}[Differential RCDC]
\label{cor:v5v32-differential-RCDC}
Suppose the coefficients of the infinitesimal variation
\(D\Delta_\phi[\dot\phi]\) obey
\begin{equation}
 \varepsilon_k\le M_k\|\dot\phi\|_+,\qquad k=0,1,2,
 \label{eq:v5v32-coefficient-Lipschitz}
\end{equation}
uniformly on the admissible ball.  Then
\begin{equation}
 \|DC_{a,b;L}(\phi)[\dot\phi]\|_{\mathfrak R_r}
 \le
 C_{a,b,r}
 \left(M_2L^{-r}+M_1L^{-r-1}+M_0L^{-r-2}\right)
 \|\dot\phi\|_+.
 \label{eq:v5v32-differential-RCDC}
\end{equation}
\end{corollary}

\begin{proof}
Set \(\widetilde\phi=\phi+\tau\dot\phi\) in
\eqref{eq:v5v32-curved-RCDC}, divide by \(|\tau|\), and pass to the limit.
\end{proof}

\begin{theorem}[Uniform RCDC on a bounded-geometry ball]
\label{thm:v5v32-RCDC}
If the Gaussian heat bounds \eqref{eq:v5v32-Gaussian} and coefficient
bounds \eqref{eq:v5v32-coefficient-Lipschitz} hold uniformly, then the
covariant heat shell satisfies the V31 RCDC of every fixed order \(r\).
Its leading constant is controlled by the principal metric-variation
coefficient \(M_2\).
\end{theorem}

\begin{proof}
Discarding the two additional powers of \(L^{-1}\) in
\eqref{eq:v5v32-differential-RCDC} gives
\[
 \|DC_{a,b;L}[\dot\phi]\|_{\mathfrak R_r}
 \le C_r(M_2+M_1+M_0)L^{-r}\|\dot\phi\|_+
\]
for \(L\ge1\), which is Definition~\ref{def:v5v31-RCDC}.
\end{proof}

\subsection{Gauge and diffeomorphism covariance of the estimate}

Every object in the proof can be transported by the unitary
\(U_g\) of V29.  The rooted distance, volume, and bundle norms are evaluated
in the transformed metric fibre.  Uniform bounded geometry makes their
chart constants common.

\begin{proposition}[Equivariance of the RCDC]
\label{prop:v5v32-equivariant-RCDC}
If the admissible field ball is symmetry invariant and its bounded-geometry
constants are uniform, the constants in
\eqref{eq:v5v32-differential-RCDC} can be chosen on symmetry orbits.
The RCDC therefore creates no new Ward defect.
\end{proposition}

\begin{proof}
The functional calculus is equivariant by
Theorem~\ref{thm:v5v29-functional-equivariance}.  Unitary conjugation
preserves the internal Schatten norms.  Pullback by the transformed metric
preserves its distance and volume integral.  Hence the rooted norm and the
Duhamel identity transform covariantly.  Uniformity of the declared
geometric constants completes the proof.
\end{proof}

\subsection{Tri-graded cost of geometric coefficients}

The coefficients \(A_2,A_1,A_0\) in
\eqref{eq:v5v32-divergence-B} are nonlinear functions of the inverse
metric, connection, curvature, gauge field, and Higgs background.  Their
uniform bounds must be derived from the strong tri-graded ball.

\begin{definition}[Bounded-geometry shell certificate]
\label{def:v5v32-BGSC}
A BGSC consists of:
\begin{enumerate}
 \item a uniform ellipticity interval for every \(\Delta_\phi\);
 \item injectivity-radius and curvature-derivative bounds sufficient for
 \eqref{eq:v5v32-Gaussian};
 \item the coefficient Lipschitz estimates
 \eqref{eq:v5v32-coefficient-Lipschitz} in the strong hybrid norm;
 \item boundary heat-kernel analogues when artificial boundaries occur;
 \item constants compatible with curvature-grade factorial weights and
 the representation-summed packet norm.
\end{enumerate}
\end{definition}

\begin{firewall}[Critical rough metrics]
\label{fw:v5v32-rough-metrics}
An \(H^s\) or curvature-action bound at a critical exponent need not give a
positive injectivity radius, pointwise ellipticity chart, or bounded
curvature derivatives.  Consequently the BGSC is not automatic on the
rough critical profile space demanded by the full continuum programme.
\end{firewall}

\begin{corollary}[Net shell gap]
\label{cor:v5v32-net-gap}
Suppose a connected insertion paired with the covariance derivative costs
\(a\) spatial derivatives and \(b\) energy indices in the V21 scale.  An
order-\(r\) rooted Taylor remainder gives a decaying shell factor whenever
\begin{equation}
 r>a+b.
 \label{eq:v5v32-net-gap}
\end{equation}
The net gain is \(L^{-(r-a-b)}\).
\end{corollary}

\begin{proof}
Equation \eqref{eq:v5v32-differential-RCDC} supplies \(L^{-r}\) in the
principal term.  V21's two-index smoothing ledger charges
\(L^{a+b}\).  Multiply the factors.
\end{proof}

The local Taylor polynomial of order below \(r\) increases the curvature
or local-operator grade.  It is precisely here that the infinite
gravitational operator tower and the V25 factorial radius are needed:
choosing large \(r\) is not a free finite-counterterm operation.

\subsection{V32 conclusion}

\begin{theorem}[V32 result]
\label{thm:v5v32-status}
On a symmetry-invariant bounded-geometry ball, a metric-dependent
Laplace-type heat shell satisfies the curved rooted covariance derivative
bound \eqref{eq:v5v32-differential-RCDC}.  The principal second-order
variation is integrable because divergence form splits one derivative onto
each heat kernel.  An order-\(r\) Taylor remainder overcomes a V21 cost
\(a+b\) exactly when \(r>a+b\).  Extending the BGSC to the required rough
critical metric space remains open.
\end{theorem}

\begin{proof}
Use Theorem~\ref{thm:v5v32-curved-difference},
Corollary~\ref{cor:v5v32-differential-RCDC}, and
Corollary~\ref{cor:v5v32-net-gap}.
\end{proof}

V33 will go upstream on the remaining rough-metric gate.  It will test
whether the critical curvature control available in four dimensions can
imply the BGSC, construct concentrating metric families when it cannot,
and identify the weakest geometric quantity that the admissible shell
space must add.

\clearpage
\section{V33: Critical curvature does not control the heat radius}
\label{sec:v5v33}

\subsection{The geometric implication to be tested}

V32 proves the curved rooted covariance derivative certificate on a
uniform bounded-geometry ball.  The continuum programme, however, is
naturally controlled by four-dimensional critical curvature quantities
such as
\begin{equation}
 \int_M|\operatorname{Rm}_g|^2\,dV_g.
 \label{eq:v5v33-critical-curvature}
\end{equation}
This note tests whether such a bound supplies the geometric hypotheses of
V32.  It does not: collapse is invisible to
\eqref{eq:v5v33-critical-curvature}, even at fixed total volume.

\subsection{A flat fixed-volume collapsing family}

Let
\[
 \mathbb T^4=\mathbb R^4/(2\pi\mathbb Z)^4
\]
and, for \(0<\varepsilon\le1\), define the constant flat metric
\begin{equation}
 g_\varepsilon
 =
 \varepsilon^2(dx^1)^2
 +\varepsilon^{-2/3}\bigl((dx^2)^2+(dx^3)^2+(dx^4)^2\bigr).
 \label{eq:v5v33-flat-collapse}
\end{equation}

\begin{theorem}[Critical curvature and volume do not prevent collapse]
\label{thm:v5v33-flat-collapse}
The family \eqref{eq:v5v33-flat-collapse} satisfies
\begin{align}
 \operatorname{Vol}(\mathbb T^4,g_\varepsilon)
 &=\operatorname{Vol}(\mathbb T^4,g_1),
 \label{eq:v5v33-fixed-volume}\\
 \int_{\mathbb T^4}|\operatorname{Rm}_{g_\varepsilon}|^2\,dV_{g_\varepsilon}
 &=0,
 \label{eq:v5v33-zero-curvature}\\
 \operatorname{inj}(\mathbb T^4,g_\varepsilon)
 &=\pi\varepsilon\longrightarrow0.
 \label{eq:v5v33-inj-collapse}
\end{align}
It is not uniformly elliptic relative to \(g_1\).
\end{theorem}

\begin{proof}
The determinant of the diagonal metric matrix is
\[
 \varepsilon^2(\varepsilon^{-2/3})^3=1,
\]
which proves \eqref{eq:v5v33-fixed-volume}.  Its coefficients are constant,
so its Levi--Civita connection is flat and
\eqref{eq:v5v33-zero-curvature} follows.  The shortest nontrivial closed
geodesic is the \(x^1\)-circle of length \(2\pi\varepsilon\), hence the
injectivity radius is half that length.  Finally, the smallest eigenvalue
of \(g_\varepsilon\) relative to \(g_1\) is \(\varepsilon^2\), and the
largest is \(\varepsilon^{-2/3}\); no common positive ellipticity interval
exists.
\end{proof}

\begin{corollary}[No BGSC from the critical action alone]
\label{cor:v5v33-no-BGSC}
There is no function of an upper bound for
\(\int|\operatorname{Rm}|^2\) and two-sided total-volume bounds alone that
gives a positive lower bound for injectivity radius or a uniform
ellipticity constant.  Therefore those data alone do not imply the V32
bounded-geometry shell certificate.
\end{corollary}

\begin{proof}
The family in Theorem~\ref{thm:v5v33-flat-collapse} has identical volume
and zero critical curvature energy while its injectivity radius and
ellipticity constants degenerate.
\end{proof}

\begin{firewall}[Curvature counterterms do not see flat collapse]
\label{fw:v5v33-flat-invisible}
Adding higher local polynomials in the Riemann tensor does not penalize the
family \eqref{eq:v5v33-flat-collapse}: every such curvature polynomial
vanishes.  A noncollapse condition, moduli treatment, or dynamical
mechanism outside the local curvature tower is required.
\end{firewall}

\subsection{The volume radius}

For \(0<\nu<\omega_4\), define the lower volume radius at \(x\) by
\begin{equation}
 r_{\mathrm{vol}}(x;g,\nu)
 =
 \sup\left\{r>0:
 \operatorname{Vol}_gB_g(x,s)\ge\nu s^4
 \text{ for every }0<s\le r\right\}.
 \label{eq:v5v33-volume-radius}
\end{equation}
Let \(r_{\mathrm{vol}}(g,\nu)=\inf_xr_{\mathrm{vol}}(x;g,\nu)\).

\begin{proposition}[The flat family loses volume radius]
\label{prop:v5v33-volume-radius-collapse}
For every fixed admissible \(\nu\), there is \(C_\nu\) such that
\begin{equation}
 r_{\mathrm{vol}}(g_\varepsilon,\nu)
 \le C_\nu\varepsilon.
 \label{eq:v5v33-volume-radius-collapse}
\end{equation}
\end{proposition}

\begin{proof}
At radii \(s\) larger than the short circumference
\(2\pi\varepsilon\) but smaller than the other three circumferences, a
geodesic ball wraps around the short circle.  Its volume is bounded above
by
\[
 C\varepsilon\,s^3.
\]
Choose \(s=C_\nu\varepsilon\) with \(C_\nu\) large enough that
\(C\varepsilon s^3<\nu s^4\).  This violates the defining lower-volume
inequality at that scale, proving
\eqref{eq:v5v33-volume-radius-collapse}.
\end{proof}

Thus a uniform positive volume radius detects the collapse missed by
critical curvature.  It is necessary for the present uniform
four-dimensional heat-shell route, but by itself does not control
derivatives of the metric coefficients.

\subsection{Concentration counting and its limitation}

Let a metric satisfy the critical energy bound
\begin{equation}
 \int_M|\operatorname{Rm}_g|^2\,dV_g\le E.
 \label{eq:v5v33-energy-E}
\end{equation}
Fix a scale \(r>0\) and threshold \(\epsilon_*>0\).  Call \(x\)
\(r\)-bad if
\begin{equation}
 \int_{B_g(x,2r)}|\operatorname{Rm}_g|^2\,dV_g\ge\epsilon_*.
 \label{eq:v5v33-bad}
\end{equation}

\begin{lemma}[Energy-counted bad balls]
\label{lem:v5v33-bad-balls}
Any disjoint family of \(2r\)-balls satisfying
\eqref{eq:v5v33-bad} has cardinality at most
\begin{equation}
 N_{\mathrm{bad}}\le E/\epsilon_*.
 \label{eq:v5v33-bad-count}
\end{equation}
A maximal disjoint family has enlarged balls that cover every \(r\)-bad
point.
\end{lemma}

\begin{proof}
Disjointness and \eqref{eq:v5v33-energy-E} give
\[
 N_{\mathrm{bad}}\epsilon_*
 \le\sum_{j=1}^{N_{\mathrm{bad}}}
 \int_{B_g(x_j,2r)}|\operatorname{Rm}|^2\,dV_g
 \le E.
\]
For the covering statement, greedily choose a maximal disjoint family.
If an \(r\)-bad point were outside the corresponding fixed enlargement,
its \(2r\)-ball would be disjoint from every chosen smaller core after
adjusting the harmless enlargement factor, contradicting maximality.
\end{proof}

\begin{firewall}[Energy counting is not epsilon regularity]
\label{fw:v5v33-no-epsilon-regularity}
Lemma~\ref{lem:v5v33-bad-balls} counts balls on which curvature energy is
large.  It does not prove that the complement has harmonic-radius or heat
kernel bounds.  Such an epsilon-regularity implication requires an
additional geometric equation, gauge, Ricci control, or analytic
regularity theorem.  The flat collapse example has no curvature-bad balls
at all and still degenerates.
\end{firewall}

\subsection{An analytic heat radius}

The exact input used by V32 can be stated without imposing unnecessarily
strong classical geometry.  Fix constants \(c,C_{ij}\) and define
\(r_{\mathrm{heat}}(x;\phi)\) as the supremum of \(r>0\) such that, on the
ball \(B_\phi(x,2r)\), the heat kernel and its first derivatives obey the
localized form of \eqref{eq:v5v32-Gaussian} for \(0<t\le r^2\), with those
fixed constants.  Set
\begin{equation}
 r_{\mathrm{heat}}(\phi)=\inf_xr_{\mathrm{heat}}(x;\phi).
 \label{eq:v5v33-heat-radius}
\end{equation}

\begin{definition}[Noncollapsed heat-radius certificate]
\label{def:v5v33-NHRC}
An NHRC on an admissible shell set consists of constants
\(r_0,\nu,M_0,M_1,M_2>0\) such that
\begin{equation}
 r_{\mathrm{heat}}(\phi)\ge r_0,\qquad
 r_{\mathrm{vol}}(\phi,\nu)\ge r_0,
 \label{eq:v5v33-radii-lower}
\end{equation}
and the coefficient bounds
\eqref{eq:v5v32-coefficient-Lipschitz} hold with \(M_k\), uniformly over
that set.
\end{definition}

\begin{theorem}[NHRC implies the analytic part of the BGSC]
\label{thm:v5v33-NHRC-BGSC}
An NHRC gives the heat-kernel, noncollapse, and coefficient inputs used in
Theorem~\ref{thm:v5v32-RCDC} for every shell time
\begin{equation}
 b/L^2\le r_0^2.
 \label{eq:v5v33-shell-time}
\end{equation}
Consequently the curved RCDC holds on those ultraviolet shells with
uniform constants.
\end{theorem}

\begin{proof}
The definition of \(r_{\mathrm{heat}}\) supplies
\eqref{eq:v5v32-Gaussian} with common constants up to time \(r_0^2\).
The NHRC supplies the coefficient bounds directly.  Condition
\eqref{eq:v5v33-shell-time} places both heat times \(a/L^2,b/L^2\) inside
that interval.  Apply Theorem~\ref{thm:v5v32-RCDC}.  The volume-radius
condition prevents the dimensional collapse exhibited in V33 and makes
the four-dimensional chart normalization uniform.
\end{proof}

\begin{assessment}[Weakest input for the proved route]
\label{ass:v5v33-weakest}
The NHRC is not asserted to be a necessary condition for every possible
continuum construction.  It is the direct analytic condition sufficient
for the V32 proof.  Replacing it by an intrinsic action bound is a separate
regularity theorem, and Corollary~\ref{cor:v5v33-no-BGSC} shows that the
critical curvature action plus total volume cannot be that theorem.
\end{assessment}

\subsection{Regular region and terminal cores}

If a global NHRC is unavailable, split the manifold into a regular region
\(\mathcal R_r\) where the heat radius and volume radius exceed \(r\), and
a bad region \(\mathcal B_r\).  A partition of unity
\(\{\chi_\alpha\}\) of overlap at most \(N_0\) localizes a kernel by
\(K=\sum_{\alpha,\beta}\chi_\alpha K\chi_\beta\).

\begin{lemma}[Bounded-overlap rooted sewing]
\label{lem:v5v33-rooted-sewing}
Suppose each root belongs to at most \(N_0\) left localization charts and,
for each such chart, at most \(N_1\) right charts contribute to a
finite-range kernel.  If every localized block has rooted norm at most
\(C\), then
\begin{equation}
 \|K\|_{\mathfrak R_0}\le N_0N_1C.
 \label{eq:v5v33-rooted-sewing}
\end{equation}
\end{lemma}

\begin{proof}
For a fixed root \(x\), at most \(N_0\) left cutoffs are nonzero.  Each
has at most \(N_1\) right partners.  Apply the triangle inequality after
the internal packet norm and sum their individual rooted bounds.
\end{proof}

\begin{proposition}[Conditional regular-core decomposition]
\label{prop:v5v33-regular-core}
Assume the regular charts satisfy an NHRC with common constants, the bad
region is covered by finitely many terminal cores, every core and
regular--core bridge has an RCDC bound in its assigned boundary/rooted
norm, and the cover has uniform overlap.  Then the global shell has an
RCDC with a constant equal to the bounded-overlap sum of the regular,
core, and bridge constants.
\end{proposition}

\begin{proof}
Decompose the kernel into regular--regular, core--core, and the two bridge
types.  Apply Theorem~\ref{thm:v5v33-NHRC-BGSC} to the regular blocks and
the assumed estimates to the others.  Sew them using
Lemma~\ref{lem:v5v33-rooted-sewing}.
\end{proof}

\begin{firewall}[The terminal-core estimates are not yet supplied]
\label{fw:v5v33-core-open}
The archived determinant and heat-owner machinery organizes terminal
cores and their boundaries, but V33 has not proved the RCDC for a
collapsing or concentrating metric core, nor a uniform bridge estimate.
Proposition~\ref{prop:v5v33-regular-core} is therefore a precise assembly
theorem, not closure of the rough-metric gate.
\end{firewall}

\subsection{Consequences for the admissible metric space}

\begin{theorem}[Necessary enlargement of the metric ledger]
\label{thm:v5v33-metric-ledger}
Any use of the V32 heat-kernel proof on a metric class controlled only by
total volume and four-dimensional critical curvature energy is invalid.
A sufficient admissible ledger must add either:
\begin{enumerate}
 \item a uniform NHRC;
 \item a theorem deriving an NHRC from stronger dynamical/geometric
 hypotheses; or
 \item a regular--core decomposition with independently proved core and
 bridge RCDC estimates.
\end{enumerate}
\end{theorem}

\begin{proof}
Corollary~\ref{cor:v5v33-no-BGSC} rules out the claimed implication from
the two critical quantities.  The three alternatives are sufficient by
Theorem~\ref{thm:v5v33-NHRC-BGSC} or
Proposition~\ref{prop:v5v33-regular-core}.
\end{proof}

\subsection{V33 status}

\begin{theorem}[V33 conclusion]
\label{thm:v5v33-status}
Fixed total volume and bounded \(L^2\) curvature do not control
injectivity radius, ellipticity, volume radius, or the V32 heat constants.
The flat family \eqref{eq:v5v33-flat-collapse} proves this without any
curvature concentration.  Curvature energy can count large-energy cores
but cannot detect flat collapse.  An NHRC closes the regular-region RCDC;
the collapsing and concentrating core channel remains open.
\end{theorem}

\begin{proof}
Use Theorem~\ref{thm:v5v33-flat-collapse},
Lemma~\ref{lem:v5v33-bad-balls}, and
Theorem~\ref{thm:v5v33-NHRC-BGSC}.
\end{proof}

V34 will go upstream on the flat-collapse channel.  It will decompose the
heat shell on a product with one short circle into fibre zero modes and
massive Kaluza--Klein modes, prove the exact crossover bounds, and test
whether the collapsing modulus belongs in the finite head rather than the
irrelevant metric tail.

\clearpage
\section{V34: Kaluza--Klein crossover of the flat collapsing shell}
\label{sec:v5v34}

\subsection{Exact spectrum of the V33 family}

The flat collapse in V33 is analytically tractable and therefore tests the
placement of collapsing moduli in the shell coordinates.  On the torus
with metric \eqref{eq:v5v33-flat-collapse}, the scalar Laplacian has
eigenfunctions
\[
 e_n(x)=e^{in\cdot x},\qquad n=(n_1,n')\in\mathbb Z\times\mathbb Z^3,
\]
and eigenvalues
\begin{equation}
 \lambda_{\varepsilon,n}
 =
 \frac{n_1^2}{\varepsilon^2}
 +\varepsilon^{2/3}|n'|^2.
 \label{eq:v5v34-eigenvalues}
\end{equation}
At shell scale \(L\), introduce
\begin{equation}
 u=L\varepsilon,\qquad
 R=L\varepsilon^{-1/3}.
 \label{eq:v5v34-uR}
\end{equation}
Then
\begin{equation}
 \frac{\lambda_{\varepsilon,n}}{L^2}
 =
 \frac{n_1^2}{u^2}+\frac{|n'|^2}{R^2}.
 \label{eq:v5v34-scaled-spectrum}
\end{equation}
The dimensionless number \(u\) is the number of short-circle wavelengths
resolved by the shell, up to fixed constants.

\subsection{Heat-shell trace crossover}

Let
\begin{equation}
 f(x)=e^{-ax}-e^{-bx},\qquad 0<a<b.
 \label{eq:v5v34-f}
\end{equation}
This is the positive multiplier of the V29 heat-shell covariance.  Define
\begin{equation}
 \mathcal N_{\varepsilon,L}
 =
 \operatorname{Tr} f(\Delta_\varepsilon/L^2)
 =
 \sum_{n_1\in\mathbb Z}\sum_{n'\in\mathbb Z^3}
 f\left(\frac{n_1^2}{u^2}+\frac{|n'|^2}{R^2}\right).
 \label{eq:v5v34-effective-count}
\end{equation}

\begin{lemma}[Lattice Riemann sums]
\label{lem:v5v34-Riemann}
If \(F\) is a Schwartz function on \(\mathbb R^k\), then
\begin{equation}
 q^{-k}\sum_{n\in\mathbb Z^k}F(n/q)
 \longrightarrow\int_{\mathbb R^k}F(\xi)\,d\xi
 \quad(q\to\infty).
 \label{eq:v5v34-Riemann}
\end{equation}
The convergence is dominated for the Gaussian functions used below.
\end{lemma}

\begin{proof}
On every fixed cube, ordinary Riemann-sum convergence applies.  Outside a
large cube, the Schwartz bound gives a summable majorant uniformly for
\(q\ge1\), by comparison with the integral over unit lattice cubes.
First choose the cube so that both tails are small, then pass to the limit
inside it.
\end{proof}

\begin{theorem}[Three- and four-dimensional shell regimes]
\label{thm:v5v34-crossover}
Let \(L\to\infty\) and allow \(\varepsilon=\varepsilon_L\).
\begin{enumerate}
 \item If \(u=L\varepsilon\to0\), then
 \begin{equation}
  \frac{\varepsilon}{L^3}\mathcal N_{\varepsilon,L}
  \longrightarrow
  I_3(a,b):=
  \int_{\mathbb R^3}f(|\xi|^2)\,d\xi>0.
  \label{eq:v5v34-3d-limit}
 \end{equation}
 \item If \(u=L\varepsilon\to\infty\), then
 \begin{equation}
  L^{-4}\mathcal N_{\varepsilon,L}
  \longrightarrow
  I_4(a,b):=
  \int_{\mathbb R^4}f(|\xi|^2)\,d\xi>0.
  \label{eq:v5v34-4d-limit}
 \end{equation}
\end{enumerate}
\end{theorem}

\begin{proof}
In the first regime, the \(n_1=0\) contribution is
\[
 \sum_{n'\in\mathbb Z^3}f(|n'|^2/R^2).
\]
Since \(R=L\varepsilon^{-1/3}\to\infty\), Lemma
\ref{lem:v5v34-Riemann} makes this asymptotic to \(R^3I_3\).  Because
\(R^3=L^3/\varepsilon\), it gives \eqref{eq:v5v34-3d-limit}.

For \(n_1\ne0\), use
\[
 0\le f(x)\le e^{-ax}
\]
and factor the Gaussian sums:
\[
 \sum_{n_1\ne0,n'}f\left(\frac{n_1^2}{u^2}
                         +\frac{|n'|^2}{R^2}\right)
 \le
 \left(\sum_{n_1\ne0}e^{-an_1^2/u^2}\right)
 \left(\sum_{n'}e^{-a|n'|^2/R^2}\right).
\]
The first factor is \(O(e^{-a/u^2})\) and the second is \(O(R^3)\).
After division by \(R^3=L^3/\varepsilon\), the result tends to zero.
This proves the first statement.

In the second regime, apply the anisotropic four-dimensional Riemann sum
to
\[
 F(\xi_1,\xi')=f(\xi_1^2+|\xi'|^2)
\]
with mesh sizes \(u^{-1}\) and \(R^{-1}\).  Both \(u,R\to\infty\), and
\[
 uR^3=(L\varepsilon)(L^3/\varepsilon)=L^4.
\]
Lemma~\ref{lem:v5v34-Riemann}, with the same Gaussian domination, gives
\eqref{eq:v5v34-4d-limit}.  Positivity of \(I_3,I_4\) follows from
\(f(x)>0\) for \(x>0\).
\end{proof}

\begin{corollary}[The crossover occurs at \(L\varepsilon\asymp1\)]
\label{cor:v5v34-crossover}
The leading mode count is of order
\begin{equation}
 \mathcal N_{\varepsilon,L}
 \asymp
 \begin{cases}
  L^3/\varepsilon,&L\varepsilon\ll1,\\
  L^4,&L\varepsilon\gg1.
 \end{cases}
 \label{eq:v5v34-count-scaling}
\end{equation}
The two expressions agree in order when \(L\varepsilon\asymp1\).
\end{corollary}

\begin{proof}
This is Theorem~\ref{thm:v5v34-crossover} and the identity
\(L^3/\varepsilon=L^4/(L\varepsilon)\).
\end{proof}

\subsection{Exact zero-mode and massive-mode decomposition}

Let \(P_0\) project onto functions independent of \(x^1\), and
\(Q_0=I-P_0\).  Both commute with \(\Delta_\varepsilon\).  On the zero-mode
sector,
\begin{equation}
 P_0\Delta_\varepsilon P_0
 =
 \varepsilon^{2/3}\Delta_{\mathbb T^3}.
 \label{eq:v5v34-zero-operator}
\end{equation}

\begin{proposition}[Massive Kaluza--Klein suppression]
\label{prop:v5v34-KK-suppression}
For the heat shell \(C_{\varepsilon,L}=f(\Delta_\varepsilon/L^2)\),
\begin{equation}
 \|Q_0C_{\varepsilon,L}Q_0\|_{\mathrm{op}}
 \le e^{-a/(L\varepsilon)^2}.
 \label{eq:v5v34-KK-suppression}
\end{equation}
Meanwhile
\begin{equation}
 P_0C_{\varepsilon,L}P_0
 =
 P_0f(\varepsilon^{2/3}\Delta_{\mathbb T^3}/L^2)P_0
 \label{eq:v5v34-zero-shell}
\end{equation}
exactly.
\end{proposition}

\begin{proof}
On \(Q_0\), \(n_1\ne0\), so
\(\lambda_{\varepsilon,n}/L^2\ge(L\varepsilon)^{-2}\).  Since
\(0\le f(x)\le e^{-ax}\), the spectral theorem gives
\eqref{eq:v5v34-KK-suppression}.  Equation
\eqref{eq:v5v34-zero-shell} follows from
\eqref{eq:v5v34-zero-operator} and functional calculus.
\end{proof}

\begin{corollary}[Operator-level dimensional reduction]
\label{cor:v5v34-dimensional-reduction}
If \(L\varepsilon\to0\), then
\begin{equation}
 \|C_{\varepsilon,L}
 -P_0f(\varepsilon^{2/3}\Delta_{\mathbb T^3}/L^2)P_0\|_{\mathrm{op}}
 \longrightarrow0
 \label{eq:v5v34-dimensional-reduction}
\end{equation}
faster than every power of \(L\varepsilon\).
\end{corollary}

\begin{proof}
The difference is the \(Q_0\) block because \(P_0\) commutes with the
operator.  Apply \eqref{eq:v5v34-KK-suppression} and the elementary fact
that \(e^{-a/u^2}=O(u^N)\) as \(u\downarrow0\) for every \(N\).
\end{proof}

\subsection{Variation of the collapse modulus}

Let \(\tau=\log\varepsilon\).  Differentiating
\eqref{eq:v5v34-eigenvalues} gives
\begin{equation}
 \partial_\tau\lambda_{\varepsilon,n}
 =
 -2\frac{n_1^2}{\varepsilon^2}
 +\frac23\varepsilon^{2/3}|n'|^2.
 \label{eq:v5v34-modulus-variation}
\end{equation}
The relative variation is order one on either part of a scale-matched
shell; it does not gain an ultraviolet power.

\begin{proposition}[The modulus is not an irrelevant shell perturbation]
\label{prop:v5v34-modulus-not-irrelevant}
There is no \(\delta>0\) and constant \(C\), uniform in
\(\varepsilon,L\), for which
\begin{equation}
 \|\partial_\tau C_{\varepsilon,L}\|_{\mathrm{op}}
 \le CL^{-\delta}
 \label{eq:v5v34-false-irrelevant}
\end{equation}
whenever the shell contains a nonzero mode.  Along sequences with
\(L\varepsilon\to0\), the surviving zero-mode sector has an order-one
response to \(\tau\) at transverse shell frequency.
\end{proposition}

\begin{proof}
On the zero-mode sector, choose lattice vectors \(n'_L\) such that
\[
 x_L=\varepsilon^{2/3}|n'_L|^2/L^2
\]
converges to a point \(x_*>0\) where \(f'(x_*)x_*\ne0\); such a point
exists because \(f\) is nonconstant, and such lattice vectors exist since
\(R\to\infty\).  Functional calculus and
\eqref{eq:v5v34-modulus-variation} give the corresponding eigenvalue of
\(\partial_\tau C_{\varepsilon,L}\) as
\[
 \frac23x_Lf'(x_L)\longrightarrow\frac23x_*f'(x_*)\ne0.
\]
This contradicts \eqref{eq:v5v34-false-irrelevant} for every
\(\delta>0\).
\end{proof}

\begin{theorem}[Collapse modulus belongs to the head or a separate core]
\label{thm:v5v34-modulus-head}
In any shell coordinate system whose irrelevant metric tail is required to
gain a positive power of \(L^{-1}\), the modulus
\(\tau=\log\varepsilon\) of the V33 flat family cannot be placed in that
tail uniformly through the crossover.  It must be:
\begin{enumerate}
 \item retained as a finite-dimensional head/moduli coordinate with a
 prescribed trajectory; or
 \item assigned to a separate collapsing-core channel with its own
 lower-dimensional shell map.
\end{enumerate}
\end{theorem}

\begin{proof}
Proposition~\ref{prop:v5v34-modulus-not-irrelevant} excludes the defining
positive shell gain of the irrelevant tail.  The exact commuting
decomposition in
Proposition~\ref{prop:v5v34-KK-suppression} leaves only the stated two
placements: retain the modulus while integrating its spectral fibres, or
separate the zero-mode effective theory from the massive modes.
\end{proof}

\begin{firewall}[Dimensional reduction is not decoupling of interactions]
\label{fw:v5v34-interacting-decoupling}
Corollary~\ref{cor:v5v34-dimensional-reduction} is an operator-norm
statement for the Gaussian heat shell.  An interacting lower-dimensional
effective action additionally requires uniform integration of the massive
Kaluza--Klein fields, matching of induced local operators, control of
zero-mode couplings, and the Ward and OS certificates.  None follows from
the Gaussian estimate alone.
\end{firewall}

\subsection{Cofinal ambiguity of the spectral dimension}

\begin{theorem}[Different cofinal modulus paths give different shell
asymptotics]
\label{thm:v5v34-cofinal-paths}
There exist fixed-volume flat metric sequences with identical zero
curvature energy for which the ultraviolet heat-shell count has either the
three-dimensional asymptotic \eqref{eq:v5v34-3d-limit} or the
four-dimensional asymptotic \eqref{eq:v5v34-4d-limit}.
\end{theorem}

\begin{proof}
Take, for example,
\(\varepsilon_L=L^{-2}\), so \(L\varepsilon_L\to0\), and
\(\widetilde\varepsilon_L=L^{-1/2}\), so
\(L\widetilde\varepsilon_L\to\infty\).  Both families are instances of
\eqref{eq:v5v33-flat-collapse}, hence have fixed volume and zero curvature
energy.  Apply Theorem~\ref{thm:v5v34-crossover}.
\end{proof}

\begin{corollary}[Critical metric data do not select a universality class]
\label{cor:v5v34-universality}
The critical curvature and total-volume ledger does not determine the
ultraviolet spectral dimension of a cofinal metric sequence.  A compact
inverse-limit construction using that ledger must add a modulus/core
selection rule.
\end{corollary}

\begin{proof}
The two sequences in
Theorem~\ref{thm:v5v34-cofinal-paths} have identical values of those
ledger quantities and different normalized shell asymptotics.
\end{proof}

\subsection{Insertion into the V26 head solver}

Let \(\tau_j\) denote all noncollapsing and collapsing moduli retained at
shell \(j\).  They enlarge the finite head \(\mathbb H\) in
\eqref{eq:v5v26-HPT}.  The renormalization/matching equations must
determine a trajectory
\begin{equation}
 \mathcal F_j(\tau_j,\tau_{j+1},p_{j+1},t_{j+1})=0.
 \label{eq:v5v34-modulus-matching}
\end{equation}
The V26 head solver applies only if the corresponding finite-dimensional
map has a uniform inverse or contraction constant.  Neither curvature
energy nor the Gaussian KK calculation supplies that dynamical estimate.

\begin{definition}[Modulus ancestry certificate]
\label{def:v5v34-MAC}
A MAC consists of:
\begin{enumerate}
 \item a finite or compactly parametrized moduli head containing every
 collapse direction not detected by the local curvature tower;
 \item compatible equations \eqref{eq:v5v34-modulus-matching} with a
 uniform V26 head-solver bound;
 \item a declared choice of four-dimensional or dimensionally reduced
 cofinal regime for every short modulus;
 \item uniform integration and matching of massive modes;
 \item preservation of the joint Ward complex and metric OS certificate.
\end{enumerate}
\end{definition}

\begin{firewall}[The moduli head need not be finite globally]
\label{fw:v5v34-moduli-dimension}
The flat torus has a finite-dimensional moduli space, but general
collapsing geometries can have stratified or unbounded degeneration data.
V34 proves the need to extract the displayed modulus; it does not prove
that all collapse channels fit a fixed finite-dimensional head.
\end{firewall}

\subsection{V34 status}

\begin{theorem}[V34 conclusion]
\label{thm:v5v34-status}
The V33 flat collapse has an exact Kaluza--Klein crossover governed by
\(u=L\varepsilon\).  For \(u\to0\), massive circle modes are exponentially
suppressed and the shell is three-dimensional with count
\(L^3/\varepsilon\); for \(u\to\infty\), the four-dimensional
\(L^4\) count returns.  The logarithmic collapse modulus has an order-one
shell response and therefore cannot be a uniformly irrelevant coordinate.
A MAC, or a separate collapsing-core theory, is required.
\end{theorem}

\begin{proof}
Combine Theorem~\ref{thm:v5v34-crossover},
Proposition~\ref{prop:v5v34-KK-suppression}, and
Theorem~\ref{thm:v5v34-modulus-head}.
\end{proof}

V35 will begin with the compulsory companion audit.  It will then decide
whether the new YM representation-summed programme and the NS signed
formation programme alter the modulus/core route before constructing a
two-regime shell map across \(L\varepsilon\asymp1\).

\clearpage
\section{V35: Exact two-regime shell atlas across a collapsing circle}
\label{sec:v5v35}

\subsection{Compulsory companion audit}

The newest companion notes at this fifth-version checkpoint are
\begin{align*}
 &\texttt{Renewal\_Yang\_Mills\_Mass\_Gap\_Research\_Notes\_V41.tex},
 \qquad
 \operatorname{SHA256}=
 \texttt{F2FB159D2F6725E054C41A69E9883C27176FF537F2D0A4577916A4A0A7D5908A},
 \\
 &\texttt{Renewal\_NS\_Stability\_Research\_Notes\_V31.tex},
 \qquad
 \operatorname{SHA256}=
 \texttt{F1121B62238AD5491BB7CF03635EA9934942EDF573ED8FAAA9EE79E1D38A6696}.
\end{align*}
The NS file is byte-identical to the file audited at V30.  Its conclusion
is unchanged: heat formation pays exactly the critical first moment, and
any improvement must retain the signed nonlinear correlation before
absolute values.

YM V41 supplies new evidence for the representation-summed route selected
in V30.  On the sparse transition packet that defeated a coefficientwise
classical estimate, the exact Peter--Weyl normalization turns the output
weights into a pinching subprobability.  Three assembled fixed connected
prefix differences each pay one Wilson clock; together they dominate the
hot Cartan payer and leave an inverse-\(A_j\) surplus over the marked-tree
target.  This does not prove the SM RSTC, but it shows why a complete
connected output sum can close when a selected local scalar does not.

Accordingly, V35 does not approximate the collapsing shell mode by mode.
It first factors the complete Gaussian measure and defines the reduced
interaction by its logarithm, so every Kaluza--Klein contribution is
assembled as a connected cumulant before a norm is applied.

\begin{firewall}[No companion transfer]
\label{fw:v5v35-companion}
The YM BFA estimate and the NS formation ledger are not estimates for a
collapsing SM--Einstein metric.  They justify only the ordering of
operations used below: exact representation or mode summation, connected
assembly, and then the rooted tri-graded norm.
\end{firewall}

\subsection{Exact Gaussian factorization}

Retain the flat family of V33--V34.  Let
\(P_0\) be the short-circle zero-mode projection and \(Q_0=I-P_0\).
For the heat-shell covariance
\[
 C_{\varepsilon,L}=f(\Delta_\varepsilon/L^2),
 \qquad f(x)=e^{-ax}-e^{-bx},
\]
write
\begin{equation}
 C_0=P_0C_{\varepsilon,L}P_0,\qquad
 C_{\mathrm{KK}}=Q_0C_{\varepsilon,L}Q_0.
 \label{eq:v5v35-covariance-blocks}
\end{equation}
Both cross blocks vanish because \(P_0\) commutes with
\(\Delta_\varepsilon\).

\begin{theorem}[Exact zero--massive Gaussian product]
\label{thm:v5v35-Gaussian-product}
On the support of \(C_{\varepsilon,L}\),
\begin{equation}
 \mu_{C_{\varepsilon,L}}
 =
 \mu_{C_0}\otimes\mu_{C_{\mathrm{KK}}}
 \label{eq:v5v35-Gaussian-product}
\end{equation}
under the orthogonal identification
\(\mathcal E=P_0\mathcal E\oplus Q_0\mathcal E\).
\end{theorem}

\begin{proof}
The characteristic function of the right side at
\(\zeta=\zeta_0+\zeta_{\mathrm{KK}}\) is
\[
 \exp\!\left(-\frac12\langle\zeta_0,C_0\zeta_0\rangle
             -\frac12\langle\zeta_{\mathrm{KK}},
              C_{\mathrm{KK}}\zeta_{\mathrm{KK}}\rangle\right).
\]
The block decomposition \eqref{eq:v5v35-covariance-blocks} makes the
exponent equal to
\(-\frac12\langle\zeta,C_{\varepsilon,L}\zeta\rangle\), the characteristic
function of the left side.  Equality of finite-dimensional Gaussian
measures follows.
\end{proof}

Let \(p\in P_0\mathcal E\), \(q\in Q_0\mathcal E\), and let
\(V(p,q)\) be a real interaction for which the following integrals are
finite.  Define the exact massive-mode effective interaction by
\begin{equation}
 e^{-V_{\mathrm{red}}(p)}
 =
 \int e^{-V(p,q)}\,d\mu_{C_{\mathrm{KK}}}(q).
 \label{eq:v5v35-reduced-interaction}
\end{equation}

\begin{theorem}[Exact interacting reduction]
\label{thm:v5v35-exact-reduction}
For every integrable observable \(F(p,q)\), define
\begin{equation}
 F_{\mathrm{red}}(p)
 =
 e^{V_{\mathrm{red}}(p)}
 \int F(p,q)e^{-V(p,q)}
 \,d\mu_{C_{\mathrm{KK}}}(q).
 \label{eq:v5v35-reduced-observable}
\end{equation}
Then
\begin{equation}
 \frac{\displaystyle
 \int F(p,q)e^{-V(p,q)}
 \,d\mu_{C_0}(p)d\mu_{C_{\mathrm{KK}}}(q)}
 {\displaystyle
 \int e^{-V(p,q)}
 \,d\mu_{C_0}(p)d\mu_{C_{\mathrm{KK}}}(q)}
 =
 \frac{\displaystyle
 \int F_{\mathrm{red}}(p)e^{-V_{\mathrm{red}}(p)}
 \,d\mu_{C_0}(p)}
 {\displaystyle
 \int e^{-V_{\mathrm{red}}(p)}\,d\mu_{C_0}(p)}.
 \label{eq:v5v35-exact-reduction}
\end{equation}
\end{theorem}

\begin{proof}
By \eqref{eq:v5v35-reduced-observable},
\[
 F_{\mathrm{red}}(p)e^{-V_{\mathrm{red}}(p)}
 =
 \int F(p,q)e^{-V(p,q)}\,d\mu_{C_{\mathrm{KK}}}(q).
\]
The analogous identity with \(F=1\) is
\eqref{eq:v5v35-reduced-interaction}.  Substitute both identities into
the right side of \eqref{eq:v5v35-exact-reduction} and apply Fubini's
theorem.
\end{proof}

\begin{assessment}[An atlas, not a blended action]
\label{ass:v5v35-atlas}
Equation \eqref{eq:v5v35-exact-reduction} gives two coordinate
descriptions of the same shell integral: a four-dimensional chart in
\((p,q)\) and a reduced chart in \(p\) with \(V_{\mathrm{red}}\).  Taking a
convex combination of two separately truncated nonlinear shell maps would
not equal this integral and is not used.
\end{assessment}

\subsection{Connected massive-mode assembly}

Write
\[
 V(p,q)=V_0(p)+W(p,q)
\]
and, for fixed \(p\), define
\[
 M_p(z)=\mathbb E_{\mathrm{KK}}e^{zW(p,q)}.
\]

\begin{theorem}[Cumulant representation of the reduced interaction]
\label{thm:v5v35-cumulant}
Suppose \(M_p(z)\) is analytic and nonzero on a disc
\(|z|<1+\delta\), uniformly for \(p\) in the retained-field ball.  Let
\(\kappa_n^{\mathrm{KK}}(W(p,\cdot))\) be the derivatives at zero of
\(\log M_p(z)\).  Then
\begin{equation}
 V_{\mathrm{red}}(p)
 =
 V_0(p)
 +\sum_{n\ge1}\frac{(-1)^{n+1}}{n!}
 \kappa_n^{\mathrm{KK}}(W(p,\cdot)),
 \label{eq:v5v35-cumulant-expansion}
\end{equation}
and the series converges at \(z=-1\).
\end{theorem}

\begin{proof}
By the nonvanishing hypothesis there is an analytic branch of
\(\log M_p(z)\) on the disc, fixed by \(\log M_p(0)=0\).  Its Taylor
series is
\[
 \log M_p(z)
 =\sum_{n\ge1}\frac{z^n}{n!}
 \kappa_n^{\mathrm{KK}}(W).
\]
Since
\[
 V_{\mathrm{red}}=V_0-\log M_p(-1),
\]
substitution of \(z=-1\) gives
\eqref{eq:v5v35-cumulant-expansion}.  The radius exceeds one, so the
series converges there.
\end{proof}

\begin{firewall}[Momentwise bounds can erase the gain]
\label{fw:v5v35-momentwise}
The terms in \(\kappa_n\) are signed combinations of moment partitions.
They must be assembled before the RSTC norm.  Bounding
\(\mathbb E W^k\) separately can replace a small connected difference by
order-one moments, exactly the loss excluded by the V35 YM audit.
\end{firewall}

\subsection{Operator suppression versus measure suppression}

V34 proves the operator bound
\[
 \|C_{\mathrm{KK}}\|_{\mathrm{op}}
 \le e^{-a/u^2},\qquad u=L\varepsilon.
\]
The trace of the covariance, rather than its operator norm, controls the
mean Hilbert-space size of a Gaussian massive field.

\begin{lemma}[Massive heat-shell trace bound]
\label{lem:v5v35-massive-trace}
For \(0<u\le1\), \(R=L\varepsilon^{-1/3}\ge1\), and the scalar flat model,
\begin{equation}
 \operatorname{Tr}C_{\mathrm{KK}}
 \le C_aR^3e^{-a/u^2}
 =
 C_a\frac{L^4}{u}e^{-a/u^2}.
 \label{eq:v5v35-massive-trace}
\end{equation}
\end{lemma}

\begin{proof}
Since \(0\le f(x)\le e^{-ax}\),
\begin{align*}
 \operatorname{Tr}C_{\mathrm{KK}}
 &\le
 \sum_{n_1\ne0,n'}
 \exp\!\left(-a\frac{n_1^2}{u^2}
             -a\frac{|n'|^2}{R^2}\right)\\
 &=
 \left(\sum_{n_1\ne0}e^{-an_1^2/u^2}\right)
 \left(\sum_{n'\in\mathbb Z^3}e^{-a|n'|^2/R^2}\right).
\end{align*}
For \(u\le1\), the first factor is at most
\(C_ae^{-a/u^2}\); for \(R\ge1\), comparison with a Gaussian integral
bounds the second by \(C_aR^3\).  Finally
\(R^3=L^3/\varepsilon=L^4/u\).
\end{proof}

\begin{lemma}[Gaussian Lipschitz decoupling]
\label{lem:v5v35-Lipschitz-decoupling}
Let \(q\) be a centered Gaussian with trace-class covariance \(C\).  If
\(G\) is \(M\)-Lipschitz on its Hilbert space, then
\begin{equation}
 |\mathbb EG(q)-G(0)|
 \le M(\operatorname{Tr}C)^{1/2}.
 \label{eq:v5v35-Lipschitz-decoupling}
\end{equation}
\end{lemma}

\begin{proof}
By Lipschitz continuity and Cauchy--Schwarz,
\[
 |\mathbb EG(q)-G(0)|
 \le M\mathbb E\|q\|
 \le M(\mathbb E\|q\|^2)^{1/2}.
\]
For a centered Gaussian,
\(\mathbb E\|q\|^2=\operatorname{Tr}C\).
\end{proof}

\begin{corollary}[Measure-decoupling parameter]
\label{cor:v5v35-measure-parameter}
For uniformly Lipschitz massive-mode observables, a sufficient scalar
decoupling condition is
\begin{equation}
 \mathfrak d_{\varepsilon,L}
 :=
 \frac{L^4}{u}e^{-a/u^2}
 \longrightarrow0.
 \label{eq:v5v35-decoupling-parameter}
\end{equation}
The weaker condition \(u\to0\) gives operator decoupling but does not by
itself imply \eqref{eq:v5v35-decoupling-parameter}.
\end{corollary}

\begin{proof}
Combine Lemmas~\ref{lem:v5v35-massive-trace} and
\ref{lem:v5v35-Lipschitz-decoupling}.  To see that \(u\to0\) alone is
insufficient, choose for example
\[
 u_L=\sqrt{\frac{a}{2\log L}}.
\]
Then \(e^{-a/u_L^2}=L^{-2}\), so
\(\mathfrak d_{\varepsilon,L}=L^2/u_L\to\infty\).
\end{proof}

\begin{firewall}[Infinite-dimensional interactions]
\label{fw:v5v35-Lipschitz-limit}
Uniform Hilbert-space Lipschitz bounds are rarely available for local
polynomial interactions in the continuum.  The actual decoupling estimate
must use rooted connected cumulants and local counterterm subtraction.
Corollary~\ref{cor:v5v35-measure-parameter} is a sharp warning about mode
multiplicity, not a complete interacting theorem.
\end{firewall}

\subsection{Two exact charts and their overlap}

Let \(\mathscr Z_{\varepsilon,L}(V,F)\) denote the normalized full
expectation on the left of \eqref{eq:v5v35-exact-reduction}.  Define:
\begin{enumerate}
 \item the four-dimensional chart
 \(\mathcal U_4\), in which \((C_0,C_{\mathrm{KK}},V,F)\) are retained;
 \item the reduced chart \(\mathcal U_3\), in which
 \((C_0,V_{\mathrm{red}},F_{\mathrm{red}})\) are retained.
\end{enumerate}
The chart transition \(\mathcal S_{\mathrm{KK}}\) is the exact integration
\eqref{eq:v5v35-reduced-interaction}--%
\eqref{eq:v5v35-reduced-observable}.

\begin{theorem}[Exact overlap compatibility]
\label{thm:v5v35-overlap}
Whenever both chart integrals and the Fubini interchange are valid,
\begin{equation}
 \mathscr Z_{\varepsilon,L}(V,F)
 =
 \mathscr Z_{\varepsilon,L}
 \bigl(\mathcal S_{\mathrm{KK}}(V,F)\bigr).
 \label{eq:v5v35-overlap}
\end{equation}
Successive integrations of disjoint orthogonal Kaluza--Klein packets are
associative.  Therefore the two descriptions define one shell functional
on their overlap, not two competing approximations.
\end{theorem}

\begin{proof}
The first assertion is
Theorem~\ref{thm:v5v35-exact-reduction}.  For three orthogonal Gaussian
blocks, both orders of successive integration equal the same absolutely
integrable triple integral by Fubini's theorem.  Induction gives the
statement for every finite family of disjoint packets.
\end{proof}

\begin{definition}[Two-regime shell atlas certificate]
\label{def:v5v35-TRAC}
A TRAC consists of:
\begin{enumerate}
 \item a four-dimensional RSTC and V26 contraction on a chart containing
 \(u\ge u_-\);
 \item a reduced connected-cumulant RSTC and V26 contraction on a chart
 containing \(u\le u_+\), where \(0<u_-<u_+<\infty\);
 \item exact transition \(\mathcal S_{\mathrm{KK}}\) on
 \(u_-\le u\le u_+\);
 \item uniform bounds for the transition and its inverse on their images;
 \item compatible head equations for the modulus \(\tau=\log\varepsilon\);
 \item common Ward, anomaly, and OS certificates in both charts.
\end{enumerate}
\end{definition}

\begin{theorem}[Fixed points agree across an exact chart transition]
\label{thm:v5v35-fixed-point-charts}
Let \(T_4\) and \(T_3\) be shell maps on the two charts, and suppose
\begin{equation}
 \mathcal S_{\mathrm{KK}}T_4
 =
 T_3\mathcal S_{\mathrm{KK}}
 \label{eq:v5v35-intertwining}
\end{equation}
on the overlap.  If either map has a unique fixed point there and
\(\mathcal S_{\mathrm{KK}}\) is injective on the relevant physical
observables, the two chart fixed points represent the same shell
functional.
\end{theorem}

\begin{proof}
If \(x_4=T_4x_4\), then
\[
 T_3\mathcal S_{\mathrm{KK}}x_4
 =\mathcal S_{\mathrm{KK}}T_4x_4
 =\mathcal S_{\mathrm{KK}}x_4.
\]
Thus its image is a \(T_3\) fixed point and equals the unique one.
Injectivity on physical observables identifies their shell functionals.
The reverse implication is identical when the inverse chart is defined.
\end{proof}

\subsection{Ward compatibility of the zero-mode split}

On the exact product background, \(P_0\) commutes with product-preserving
gauge transformations and diffeomorphisms.  It does not commute with an
arbitrary diffeomorphism that changes the circle fibration.

\begin{proposition}[Restricted equivariance]
\label{prop:v5v35-restricted-equivariance}
For the subgroup preserving the short-circle fibration and the product
metric ansatz, the factorization
\eqref{eq:v5v35-Gaussian-product} is equivariant and the reduced shell
inherits the V27 Ward identity, assuming the determinant measure and
cycle are equivariant.  Full diffeomorphism covariance requires treating
\(P_0\) as a moving spectral projection and including the V29 connection
term.
\end{proposition}

\begin{proof}
For the preserving subgroup,
\([U_g,P_0]=0\), so both covariance blocks transform by conjugation.
The change-of-variables proof of
Theorem~\ref{thm:v5v27-equivariant-shell} applies first to the massive
integral and then to the zero-mode integral.  For a general
diffeomorphism, \(U_gP_0U_g^{-1}\ne P_0\) in the fixed product
trivialization; Theorem~\ref{thm:v5v29-projector-derivative} supplies the
missing moving-projection term.
\end{proof}

\begin{firewall}[Reduction structure is additional head data]
\label{fw:v5v35-fibration}
A short length alone does not canonically select a circle fibration on a
general metric.  A collapsing-core chart must carry its fibration or
nilpotent reduction structure as head/core data and prove compatibility
of transitions.  Otherwise the apparent zero-mode theory is
coordinate-dependent.
\end{firewall}

\subsection{Cofinal regime control}

For a shell sequence \(L_j\to\infty\), let
\(u_j=L_j\varepsilon_j\).  Without a modulus ancestry equation, \(u_j\)
may cross the overlap infinitely often.

\begin{lemma}[Finite transition under monotone logarithmic drift]
\label{lem:v5v35-finite-transition}
If
\begin{equation}
 \log u_{j+1}-\log u_j\ge\delta>0
 \label{eq:v5v35-monotone-drift}
\end{equation}
whenever \(u_j\in[u_-/2,2u_+]\), then the sequence visits the overlap
\([u_-,u_+]\) at most
\begin{equation}
 1+\left\lceil
 \frac{\log(u_+/u_-)}{\delta}
 \right\rceil
 \label{eq:v5v35-transition-count}
\end{equation}
times.
\end{lemma}

\begin{proof}
Each visit after the first increases \(\log u\) by at least \(\delta\).
The total logarithmic width of the overlap is
\(\log(u_+/u_-)\), yielding the stated count.
\end{proof}

\begin{assessment}[Fixed versus running modulus]
\label{ass:v5v35-running-modulus}
For fixed \(\varepsilon>0\) and geometrically increasing ultraviolet
\(L_j\), condition \eqref{eq:v5v35-monotone-drift} is automatic and the
flow crosses once from the reduced chart to the four-dimensional chart.
For a dynamical \(\varepsilon_j\), the MAC must prove such control or pay
every chart transition in the head ledger.
\end{assessment}

\subsection{V35 status}

\begin{theorem}[V35 conclusion]
\label{thm:v5v35-status}
The collapsing flat heat shell has an exact two-chart description.  Its
zero and massive Gaussian sectors factor, massive integration produces
the connected interaction \eqref{eq:v5v35-cumulant-expansion}, and the
four- and three-dimensional charts agree exactly on their common domain.
Operator decoupling requires only \(u\to0\), whereas elementary
measure decoupling must also defeat the multiplicity parameter
\(\mathfrak d_{\varepsilon,L}\).  A full TRAC still requires
representation-summed interacting bounds, modulus ancestry, moving-split
diffeomorphism covariance, and OS compatibility.
\end{theorem}

\begin{proof}
Combine Theorems~\ref{thm:v5v35-Gaussian-product},
\ref{thm:v5v35-exact-reduction},
\ref{thm:v5v35-cumulant}, and
\ref{thm:v5v35-overlap}, together with
Lemma~\ref{lem:v5v35-massive-trace}.
\end{proof}

V36 will quantify the exact chart transition in the hybrid tri-graded
space.  It will derive derivative and Lipschitz formulae for
\(V_{\mathrm{red}}=-\log\mathbb E_{\mathrm{KK}}e^{-V}\), keep their
connected covariance form, and insert the resulting constants into the
V26 Schur margin.

\clearpage
\section{V36: Derivatives and injectivity of the Kaluza--Klein chart}
\label{sec:v5v36}

\subsection{The exact logarithmic derivative}

Let \(P\) denote retained zero-mode variables and \(Q\) the massive
Kaluza--Klein variables.  For each \(p\in P\), let
\[
 Z(p)=\int_Qe^{-V(p,q)}\,d\mu_{C(p)}(q),\qquad
 V_{\mathrm{red}}(p)=-\log Z(p).
\]
Both the interaction and the massive covariance may depend on \(p\).
Assume throughout this section that differentiation under the integral and
the Gaussian integration by parts are justified.

For a tangent \(h\in T_pP\), write
\[
 \dot C_h=DC(p)[h].
\]
All \(q\)-derivatives below act in the massive field fibre.

\begin{theorem}[Derivative of the exact reduced interaction]
\label{thm:v5v36-first-derivative}
The reduced interaction obeys
\begin{align}
 DV_{\mathrm{red}}(p)[h]
 &=
 \left\langle D_pV(p,q)[h]\right\rangle_{p,V}
 -\left\langle Z_{\dot C_h}(q)\right\rangle_{p,V}
 \label{eq:v5v36-score-derivative}\\
 &=
 \left\langle D_pV(p,q)[h]\right\rangle_{p,V}
 -\frac12\left\langle
 \operatorname{Tr}\!\left[
 \dot C_h\bigl(D_qV\otimes D_qV-D_q^2V\bigr)
 \right]\right\rangle_{p,V}.
 \label{eq:v5v36-inverse-free-derivative}
\end{align}
Here \(\langle\cdot\rangle_{p,V}\) is the normalized massive Gaussian
expectation with density \(e^{-V(p,q)}\), and
\(Z_{\dot C_h}\) is the centered score
\eqref{eq:v5v30-centered-insertion}.
\end{theorem}

\begin{proof}
Differentiate \(Z(p)\) using
Theorem~\ref{thm:v5v30-Gaussian-variation}:
\[
 DZ(p)[h]
 =
 \int e^{-V}
 \left(-D_pV[h]+Z_{\dot C_h}\right)d\mu_C.
\]
Since \(DV_{\mathrm{red}}=-Z^{-1}DZ\), this gives
\eqref{eq:v5v36-score-derivative}.  Apply
\eqref{eq:v5v30-IBP-formula} to \(F=e^{-V}\).  The identity
\[
 e^VD_q^2(e^{-V})=D_qV\otimes D_qV-D_q^2V
\]
gives \eqref{eq:v5v36-inverse-free-derivative}.
\end{proof}

\begin{corollary}[Fixed-covariance Hessian]
\label{cor:v5v36-Hessian}
If \(C\) is independent of \(p\), then
\begin{equation}
 D^2V_{\mathrm{red}}(p)[h,k]
 =
 \left\langle D_{pp}^2V[h,k]\right\rangle_{p,V}
 -\operatorname{Cov}_{p,V}
 \bigl(D_pV[h],D_pV[k]\bigr).
 \label{eq:v5v36-Hessian}
\end{equation}
\end{corollary}

\begin{proof}
For a fixed reference measure and any differentiable \(A(p,q)\),
\[
 D_p\langle A\rangle[h]
 =
 \langle D_pA[h]\rangle
 -\operatorname{Cov}(A,D_pV[h]).
\]
Apply this identity to
\(A=D_pV[k]\), using symmetry of the second derivative.
\end{proof}

\begin{corollary}[Concavity correction from integrated fields]
\label{cor:v5v36-concavity}
For fixed covariance and an interaction affine in a retained coordinate,
the contribution of massive fluctuations to the Hessian of
\(V_{\mathrm{red}}\) is nonpositive:
\begin{equation}
 D^2V_{\mathrm{red}}[h,h]
 =
 -\operatorname{Var}_{p,V}(D_pV[h])\le0.
 \label{eq:v5v36-concavity}
\end{equation}
\end{corollary}

\begin{proof}
Set \(D_{pp}^2V=0\) in \eqref{eq:v5v36-Hessian}.
\end{proof}

\begin{firewall}[Integrating out fields need not preserve convexity]
\label{fw:v5v36-convexity}
Even if the direct retained-field Hessian is positive, the connected
covariance in \eqref{eq:v5v36-Hessian} subtracts from it.  A lower Hessian
bound for the reduced metric or Higgs action therefore requires a
quantitative variance estimate; it is not inherited by inspection.
\end{firewall}

\subsection{Lipschitz bounds at fixed covariance}

\begin{lemma}[Log-Laplace contraction in the supremum norm]
\label{lem:v5v36-log-Laplace}
For two real interactions \(V,W\) on the same massive Gaussian space,
\begin{equation}
 \left|
 -\log\mathbb Ee^{-V}
 +\log\mathbb Ee^{-W}
 \right|
 \le\|V-W\|_\infty .
 \label{eq:v5v36-log-Laplace}
\end{equation}
\end{lemma}

\begin{proof}
Let \(\delta=\|V-W\|_\infty\).  Then
\[
 e^{-\delta}e^{-W}\le e^{-V}\le e^\delta e^{-W}.
\]
Integrate, take logarithms, and rearrange.
\end{proof}

\begin{lemma}[Change of a Gibbs expectation]
\label{lem:v5v36-Gibbs-change}
Let
\[
 \langle A\rangle_t
 =
 \frac{\mathbb E(Ae^{-V_t})}{\mathbb Ee^{-V_t}},
 \qquad
 V_t=W+t(V-W).
\]
For a fixed bounded \(A\),
\begin{equation}
 |\langle A\rangle_V-\langle A\rangle_W|
 \le
 2\|A\|_\infty\|V-W\|_\infty.
 \label{eq:v5v36-Gibbs-change}
\end{equation}
\end{lemma}

\begin{proof}
Differentiate along the interpolation:
\[
 \frac d{dt}\langle A\rangle_t
 =
 -\operatorname{Cov}_t(A,V-W).
\]
For bounded random variables,
\[
 |\operatorname{Cov}(A,B)|
 =|\langle A(B-\langle B\rangle)\rangle|
 \le2\|A\|_\infty\|B\|_\infty.
\]
Integrate from \(0\) to \(1\).
\end{proof}

\begin{theorem}[First-derivative stability of reduction]
\label{thm:v5v36-C1-stability}
Assume fixed covariance and
\begin{align}
 \|D_pV-D_pW\|_\infty&\le\delta_1,
 \label{eq:v5v36-delta1}\\
 \|D_pW\|_\infty&\le M_1,\qquad
 \|V-W\|_\infty\le\delta_0.
 \label{eq:v5v36-delta0}
\end{align}
Then
\begin{equation}
 \|DV_{\mathrm{red}}-DW_{\mathrm{red}}\|
 \le\delta_1+2M_1\delta_0.
 \label{eq:v5v36-C1-stability}
\end{equation}
\end{theorem}

\begin{proof}
At fixed covariance,
\(DV_{\mathrm{red}}=\langle D_pV\rangle_V\).  Add and subtract
\(\langle D_pW\rangle_V\).  The first difference is at most
\(\delta_1\); apply Lemma~\ref{lem:v5v36-Gibbs-change} to the second with
\(A=D_pW\).
\end{proof}

\begin{corollary}[Hessian upper bound]
\label{cor:v5v36-Hessian-bound}
If
\[
 \|D_pV\|_\infty\le M_1,\qquad
 \|D_{pp}^2V\|_\infty\le M_2,
\]
then
\begin{equation}
 \|D^2V_{\mathrm{red}}\|\le M_2+M_1^2.
 \label{eq:v5v36-Hessian-bound}
\end{equation}
\end{corollary}

\begin{proof}
For unit \(h,k\), the first term in
\eqref{eq:v5v36-Hessian} is bounded by \(M_2\).  Cauchy--Schwarz gives
\[
 |\operatorname{Cov}(D_pV[h],D_pV[k])|
 \le
 \sqrt{\mathbb E|D_pV[h]|^2}
 \sqrt{\mathbb E|D_pV[k]|^2}
 \le M_1^2.
\]
\end{proof}

\subsection{The rooted moving-covariance contribution}

Put
\begin{equation}
 H_{q,V}=D_qV\otimes D_qV-D_q^2V.
 \label{eq:v5v36-Hq}
\end{equation}
Split its kernel into the local covariant Taylor polynomial and a remainder
\(\mathcal R^rH_{q,V}\), as in V31.

\begin{theorem}[Rooted derivative bound for the chart transition]
\label{thm:v5v36-rooted-chart}
Assume:
\begin{enumerate}
 \item the curved RCDC
 \[
  \|DC(p)[h]\|_{\mathfrak R_r}
  \le C_rL^{-r}\|h\|_+;
 \]
 \item the connected massive Hessian remainder obeys the rooted dual bound
 \[
  \sup_{x,y}
  \|\mathcal R^rH_{q,V}(y,x)\|_{\mathcal S_2(\mathrm{int})}
  \le K_{H,r};
 \]
 \item \(\|D_pV[h]\|_{L^1(\langle\cdot\rangle_{p,V})}
 \le M_p\|h\|_+\).
\end{enumerate}
Then the nonlocal part of the reduced derivative satisfies
\begin{equation}
 |DV_{\mathrm{red}}^{\mathrm{nonloc}}(p)[h]|
 \le
 \left(M_p+\frac12C_rK_{H,r}L^{-r}\right)\|h\|_+.
 \label{eq:v5v36-rooted-chart}
\end{equation}
The Taylor polynomial contributes only to the declared local owner
coordinates.
\end{theorem}

\begin{proof}
Use \eqref{eq:v5v36-inverse-free-derivative}.  The first expectation is
bounded by item 3.  Apply the rooted pairing
Proposition~\ref{prop:v5v31-rooted-pairing} to
\(DC[h]\) and \(\mathcal R^rH_{q,V}\); items 1 and 2 give the second term.
The extracted polynomial is local at the root and is not included in the
nonlocal norm.
\end{proof}

\begin{corollary}[Chart contribution to the V26 matrix]
\label{cor:v5v36-chart-Schur}
If insertion of the local Taylor owners followed by the V28 homotopy has
Lipschitz constant \(L_{\mathrm{loc}}\), then a certified retained-field
column bound for the reduced chart is
\begin{equation}
 L_{\mathrm{chart}}
 \le
 M_p+\frac12C_rK_{H,r}L^{-r}+L_{\mathrm{loc}}.
 \label{eq:v5v36-chart-Schur}
\end{equation}
This number must be placed in the appropriate head, protected, or tail
entry of the V26 nonnegative Lipschitz matrix before applying its Schur
test.
\end{corollary}

\begin{proof}
Add the nonlocal bound of
Theorem~\ref{thm:v5v36-rooted-chart} and the declared local insertion
bound.  The block receiving the insertion determines the matrix entry.
\end{proof}

\subsection{The unaugmented reduction is not an invertible chart}

The map \(V\mapsto V_{\mathrm{red}}\) forgets conditional massive-mode
observables.  This matters on the overlap \(L\varepsilon\asymp1\), where
those modes have not decoupled.

\begin{proposition}[Noninjectivity of massive integration]
\label{prop:v5v36-noninjective}
Let \(q\) be a one-dimensional centered Gaussian of variance
\(\sigma^2>0\).  The two interactions
\begin{equation}
 V(q)=0,\qquad
 W_t(q)=tq+\frac12t^2\sigma^2,\qquad t\ne0,
 \label{eq:v5v36-noninjective-example}
\end{equation}
have the same reduced interaction, but different normalized massive-field
expectations:
\begin{equation}
 V_{\mathrm{red}}=W_{t,\mathrm{red}}=0,\qquad
 \langle q\rangle_V=0,\quad
 \langle q\rangle_{W_t}=-t\sigma^2.
 \label{eq:v5v36-noninjective-result}
\end{equation}
\end{proposition}

\begin{proof}
The Gaussian moment-generating formula gives
\[
 \mathbb Ee^{-tq}=\exp\!\left(\frac12t^2\sigma^2\right).
\]
Hence
\(\mathbb Ee^{-W_t}=1=\mathbb Ee^{-V}\), proving equality of the reduced
interactions.  Exponential tilting shifts the Gaussian mean by
\(-t\sigma^2\), giving the second assertion.
\end{proof}

\begin{corollary}[A reduced fixed point need not determine a full one]
\label{cor:v5v36-reduced-not-full}
A contraction theorem proved only for \(V_{\mathrm{red}}\) cannot establish
uniqueness of the full shell functional on observables containing massive
fields.  Kernel directions of the reduction map remain uncontrolled.
\end{corollary}

\begin{proof}
Proposition~\ref{prop:v5v36-noninjective} gives distinct full functionals
with identical reduced data.  A norm or fixed-point equation on that data
cannot distinguish them.
\end{proof}

\subsection{The augmented conditional chart}

Define the normalized conditional source functional
\begin{equation}
 \mathcal K_p(J)
 =
 \log
 \frac{\displaystyle
 \int_Qe^{\langle J,q\rangle-V(p,q)}\,d\mu_{C(p)}(q)}
 {\displaystyle
 \int_Qe^{-V(p,q)}\,d\mu_{C(p)}(q)}.
 \label{eq:v5v36-conditional-generator}
\end{equation}
Then \(\mathcal K_p(0)=0\), and its derivatives at zero are the connected
conditional Kaluza--Klein correlations.

\begin{theorem}[Exact augmented chart]
\label{thm:v5v36-augmented-chart}
Suppose the conditional exponential moment in
\eqref{eq:v5v36-conditional-generator} is finite on a neighborhood of
\(J=0\).  The data
\begin{equation}
 \bigl(V_{\mathrm{red}}(p),\mathcal K_p(J)\bigr)
 \label{eq:v5v36-augmented-data}
\end{equation}
determine:
\begin{enumerate}
 \item the zero-mode weight \(e^{-V_{\mathrm{red}}(p)}\);
 \item every connected conditional massive correlation;
 \item the conditional massive probability law whenever its Laplace
 transform determines the measure.
\end{enumerate}
If the conditional law is absolutely continuous with respect to
\(\mu_{C(p)}\), these data reconstruct \(V(p,q)\) up to a null set.
\end{theorem}

\begin{proof}
The first assertion is the definition of \(V_{\mathrm{red}}\).  Derivatives
of the logarithm in \eqref{eq:v5v36-conditional-generator} are cumulants,
proving the second.  A finite Laplace transform on a neighborhood of zero
uniquely determines a finite-dimensional probability measure.  Denote its
density relative to \(\mu_C\) by
\[
 \frac{d\nu_p}{d\mu_C}(q)
 =
 e^{V_{\mathrm{red}}(p)-V(p,q)}.
\]
Thus
\[
 V(p,q)
 =
 V_{\mathrm{red}}(p)
 -\log\frac{d\nu_p}{d\mu_C}(q)
\]
almost everywhere, proving reconstruction.
\end{proof}

\begin{assessment}[Relation to the grand tensor]
\label{ass:v5v36-grand-tensor}
The Taylor coefficients of \(\mathcal K_p(J)\) form the complete connected
massive-mode tensor tower.  The curvature-grade and particle-number
weights of V25 are therefore not auxiliary decoration: they are the
coordinates needed to make the reduced overlap chart injective without
retaining the original massive integration variables.
\end{assessment}

\begin{firewall}[Moment determinacy]
\label{fw:v5v36-moment-determinacy}
A formal list of moments need not determine a measure.  The analytic
Laplace-transform hypothesis in
Theorem~\ref{thm:v5v36-augmented-chart} is substantive and must be uniform
in the cofinal limit.  Without it, the augmented formal tensor can still
lose nonperturbative information.
\end{firewall}

\subsection{Contraction transported through a genuine chart}

\begin{theorem}[Bi-Lipschitz conjugation of a contraction]
\label{thm:v5v36-conjugated-contraction}
Let \(\mathcal S:X\to Y\) be injective on an invariant domain and suppose
\begin{equation}
 \|\mathcal Sx-\mathcal Sy\|_Y\le L_S\|x-y\|_X,\qquad
 \|\mathcal S^{-1}u-\mathcal S^{-1}v\|_X
 \le L_{S^{-1}}\|u-v\|_Y.
 \label{eq:v5v36-bi-Lipschitz}
\end{equation}
If \(T_X\) is \(q_X\)-Lipschitz and
\[
 T_Y=\mathcal S T_X\mathcal S^{-1},
\]
then
\begin{equation}
 \operatorname{Lip}(T_Y)
 \le L_SL_{S^{-1}}q_X.
 \label{eq:v5v36-transported-q}
\end{equation}
It is a contraction whenever the right side is below one.
\end{theorem}

\begin{proof}
For \(u,v\in Y\),
\begin{align*}
 \|T_Yu-T_Yv\|_Y
 &\le L_S
 \|T_X\mathcal S^{-1}u-T_X\mathcal S^{-1}v\|_X\\
 &\le L_Sq_XL_{S^{-1}}\|u-v\|_Y.
\end{align*}
\end{proof}

\begin{corollary}[Corrected TRAC requirement]
\label{cor:v5v36-corrected-TRAC}
On a crossover chart intended to retain all physical observables, item 4
of Definition~\ref{def:v5v35-TRAC} must be verified for the augmented data
\eqref{eq:v5v36-augmented-data}, with
\begin{equation}
 L_SL_{S^{-1}}q_X<1.
 \label{eq:v5v36-TRAC-margin}
\end{equation}
If only zero-mode observables are retained, the reduction may instead be
treated as a quotient, but it cannot be used to infer uniqueness in the
discarded directions.
\end{corollary}

\begin{proof}
The unaugmented map is noninjective by
Proposition~\ref{prop:v5v36-noninjective}.  The augmented map is injective
under Theorem~\ref{thm:v5v36-augmented-chart}.  Apply
Theorem~\ref{thm:v5v36-conjugated-contraction}; the quotient alternative is
exactly Corollary~\ref{cor:v5v36-reduced-not-full}.
\end{proof}

\subsection{V36 status}

\begin{theorem}[V36 conclusion]
\label{thm:v5v36-status}
The exact reduced interaction has derivative
\eqref{eq:v5v36-inverse-free-derivative}; its covariance dependence is a
centered, representation-summed massive Hessian insertion and obeys the
rooted chart bound \eqref{eq:v5v36-rooted-chart}.  The reduced interaction
alone is not an invertible overlap coordinate.  Adding the analytic
conditional generator \(\mathcal K_p\) restores the complete connected
massive tensor and yields an injective chart when its Laplace transform
determines the conditional law.  Uniform bi-Lipschitz estimates for this
augmented chart remain unproved.
\end{theorem}

\begin{proof}
Use Theorems~\ref{thm:v5v36-first-derivative},
\ref{thm:v5v36-rooted-chart}, and
\ref{thm:v5v36-augmented-chart}, together with
Proposition~\ref{prop:v5v36-noninjective}.
\end{proof}

V37 will estimate the conditional generator directly.  It will derive a
Gaussian interpolation formula for its connected tensors and determine
whether the massive trace parameter of V35 can be replaced by rooted
connected bounds that remain uniform at the crossover.

\clearpage
\section{V37: Connected interpolation for the conditional massive tensor}
\label{sec:v5v37}

\subsection{Interaction interpolation}

V36 shows that an injective crossover chart must retain the conditional
source functional
\[
 \mathcal K_p(J)
 =
 \log
 \frac{\mathbb E_{\mathrm{KK}}
 e^{\langle J,q\rangle-V(p,q)}}
 {\mathbb E_{\mathrm{KK}}e^{-V(p,q)}}.
\]
This note derives exact identities for its connected tensors and shows why
local rooted observables can decouple under a weaker condition than the
global trace criterion of V35.

Fix \(p\) and suppress it from the notation.  Write
\[
 V=V_0+W(q)
\]
where \(V_0\) is independent of \(q\).  For \(0\le s\le1\), define
\begin{equation}
 \mathcal K_s(J)
 =
 \log
 \frac{\mathbb E_Ce^{\langle J,q\rangle-sW(q)}}
 {\mathbb E_Ce^{-sW(q)}}.
 \label{eq:v5v37-Ks}
\end{equation}
Let \(\langle\cdot\rangle_{s,J}\) be the corresponding normalized tilted
expectation and let
\(\kappa_s(X_1,\ldots,X_n)\) denote its joint cumulant at \(J=0\).

\begin{theorem}[Connected interaction interpolation]
\label{thm:v5v37-interpolation}
Whenever the exponential moments are differentiable in \(s\) and analytic
in \(J\) near zero,
\begin{equation}
 \partial_s\mathcal K_s(J)
 =
 -\langle W\rangle_{s,J}+\langle W\rangle_{s,0}.
 \label{eq:v5v37-Ks-derivative}
\end{equation}
For every \(n\ge1\) and test vectors \(f_1,\ldots,f_n\),
\begin{equation}
 \partial_s
 D_J^n\mathcal K_s(0)[f_1,\ldots,f_n]
 =
 -\kappa_s
 \bigl(q(f_1),\ldots,q(f_n),W\bigr).
 \label{eq:v5v37-tensor-interpolation}
\end{equation}
Consequently,
\begin{align}
 \kappa_{V}^{(n)}(f_1,\ldots,f_n)
 &=
 \kappa_{0}^{(n)}(f_1,\ldots,f_n)
 \nonumber\\
 &\quad-
 \int_0^1
 \kappa_s
 \bigl(q(f_1),\ldots,q(f_n),W\bigr)\,ds.
 \label{eq:v5v37-integrated-interpolation}
\end{align}
\end{theorem}

\begin{proof}
Differentiate the two logarithms in \eqref{eq:v5v37-Ks}; their logarithmic
derivatives are respectively
\(-\langle W\rangle_{s,J}\) and
\(-\langle W\rangle_{s,0}\), proving
\eqref{eq:v5v37-Ks-derivative}.  Derivatives in \(J\) of a tilted
expectation are joint cumulants with the inserted source variables.  This
follows directly by differentiating the logarithmic generating function
for \(W,q(f_1),\ldots,q(f_n)\).  The second term in
\eqref{eq:v5v37-Ks-derivative} is independent of \(J\), giving
\eqref{eq:v5v37-tensor-interpolation}.  Integrate from \(s=0\) to \(s=1\).
\end{proof}

\begin{corollary}[Gaussian endpoint]
\label{cor:v5v37-Gaussian-endpoint}
At \(s=0\),
\begin{equation}
 \mathcal K_0(J)=\frac12\langle J,CJ\rangle,
 \label{eq:v5v37-K0}
\end{equation}
so
\begin{equation}
 \kappa_0^{(1)}=0,\qquad
 \kappa_0^{(2)}=C,\qquad
 \kappa_0^{(n)}=0\quad(n\ge3).
 \label{eq:v5v37-Gaussian-cumulants}
\end{equation}
\end{corollary}

\begin{proof}
This is the moment-generating formula for a centered Gaussian.
\end{proof}

\subsection{Every Gaussian external leg carries a covariance}

\begin{theorem}[Gaussian connected-leg identity]
\label{thm:v5v37-Gaussian-leg}
Let \(q\sim N(0,C)\) on a finite-dimensional Hilbert space, and let
\(W\) be \(n\) times differentiable with the required Gaussian
integrability.  Then
\begin{equation}
 \kappa_0\bigl(q(f_1),\ldots,q(f_n),W(q)\bigr)
 =
 \mathbb E_C
 D^nW(q)[Cf_1,\ldots,Cf_n].
 \label{eq:v5v37-Gaussian-leg}
\end{equation}
\end{theorem}

\begin{proof}
Gaussian completion of the square gives, for a source
\(J=\sum_it_if_i\),
\[
 \frac{\mathbb E_C\bigl(W(q)e^{\langle J,q\rangle}\bigr)}
 {\mathbb E_Ce^{\langle J,q\rangle}}
 =
 \mathbb E_C W(q+CJ).
\]
The derivative in \(t_1,\ldots,t_n\) of the left side at zero is exactly
the joint cumulant of the \(n\) linear fields with \(W\).  Differentiating
the translated right side gives
\eqref{eq:v5v37-Gaussian-leg}.
\end{proof}

\begin{assessment}[Connected-prefix mechanism]
\label{ass:v5v37-prefix}
Equation \eqref{eq:v5v37-Gaussian-leg} assigns one full covariance to every
external massive leg before any norm is taken.  This is the continuum
Gaussian form of the fixed connected-prefix payment observed in the V35 YM
audit.  Replacing the joint cumulant by a product of raw moment bounds can
lose these factors.
\end{assessment}

\subsection{Rooted norm of the massive covariance}

On the flat product of V34, the heat kernel factorizes into its short-circle
and transverse parts.  Let \(k_t^{S^1_\varepsilon}\) be the normalized
circle heat kernel and let \(P_0\) be its constant-mode projection.  The
massive circle kernel is
\begin{equation}
 k_{t,\mathrm{KK}}^\varepsilon
 =
 k_t^{S^1_\varepsilon}-P_0
 =
 \frac1{2\pi\varepsilon}
 \sum_{n_1\ne0}
 e^{-tn_1^2/\varepsilon^2}
 e^{in_1(x^1-y^1)}.
 \label{eq:v5v37-circle-KK}
\end{equation}
Coordinates in the exponential are understood in their dimensionless
\(2\pi\)-periodic normalization.

\begin{lemma}[Massive rooted suppression]
\label{lem:v5v37-rooted-massive}
For \(t=a/L^2\) and \(0<u=L\varepsilon\le1\),
\begin{equation}
 \sup_x\int
 \left|
 k_{t,\mathrm{KK}}^\varepsilon(x,y)
 \right|\,dV_{S^1_\varepsilon}(y)
 \le C_ae^{-a/u^2}.
 \label{eq:v5v37-circle-root}
\end{equation}
For the four-dimensional product heat kernel,
\begin{equation}
 \|Q_0e^{-a\Delta_\varepsilon/L^2}Q_0\|_{\mathfrak R_0}
 \le C_ae^{-a/u^2}.
 \label{eq:v5v37-product-root}
\end{equation}
\end{lemma}

\begin{proof}
Integrate the absolute Fourier series in
\eqref{eq:v5v37-circle-KK}.  The normalized volume cancels the prefactor,
and
\[
 \int|k_{t,\mathrm{KK}}^\varepsilon(x,y)|\,dy
 \le\sum_{n_1\ne0}e^{-an_1^2/u^2}
 \le C_ae^{-a/u^2}
\]
for \(u\le1\).  The transverse heat kernel is positive and has integral
one.  Tensor-product factorization therefore gives
\eqref{eq:v5v37-product-root}.
\end{proof}

\begin{corollary}[Massive heat-shell rooted suppression]
\label{cor:v5v37-shell-root}
For \(C_{\mathrm{KK}}=Q_0f(\Delta_\varepsilon/L^2)Q_0\),
\begin{equation}
 \|C_{\mathrm{KK}}\|_{\mathfrak R_0}
 \le C_{a,b}e^{-a/u^2}.
 \label{eq:v5v37-shell-root}
\end{equation}
For every \(r\ge0\),
\begin{equation}
 \|C_{\mathrm{KK}}\|_{\mathfrak R_r}
 \le
 C_{a,b,r}L^{-r}e^{-a/u^2}
 \qquad(0<u\le1).
 \label{eq:v5v37-shell-root-r}
\end{equation}
\end{corollary}

\begin{proof}
Apply Lemma~\ref{lem:v5v37-rooted-massive} at heat times
\(a/L^2\) and \(b/L^2\) and add the bounds.  For the \(r\)-moment, use
\[
 d_{\mathrm{prod}}(x,y)^r
 \le c_r\bigl(d_{S^1}(x^1,y^1)^r+d_{\mathbb T^3}(x',y')^r\bigr).
\]
The short-circle distance is at most \(C\varepsilon=Cu/L\le C/L\), and
the transverse heat-kernel moment is \(O(L^{-r})\) in the physical
transverse metric.  The same Fourier-series bound supplies the common
exponential factor.
\end{proof}

\begin{theorem}[Local and global massive decoupling differ]
\label{thm:v5v37-local-global}
As \(u\to0\):
\begin{enumerate}
 \item the massive covariance tends to zero exponentially in every fixed
 rooted moment norm by \eqref{eq:v5v37-shell-root-r};
 \item its trace need not tend to zero unless the stronger multiplicity
 condition \eqref{eq:v5v35-decoupling-parameter} holds.
\end{enumerate}
Therefore local rooted observables can have a decoupling limit on cofinal
paths where the global massive Hilbert norm does not.
\end{theorem}

\begin{proof}
The first assertion is
Corollary~\ref{cor:v5v37-shell-root}.  The trace bound and counterexample
path in Corollary~\ref{cor:v5v35-measure-parameter} prove the second.
\end{proof}

\begin{firewall}[Observable algebra must be declared]
\label{fw:v5v37-observable-algebra}
Theorem~\ref{thm:v5v37-local-global} does not permit a local limit to be
reported as convergence of the full global field measure.  The continuum
statement must specify whether its observable algebra is local/rooted or
contains extensive massive-field observables.
\end{firewall}

\subsection{First connected correction in the rooted norm}

Let the kernel norm of an \(n\)-linear functional be the V25 rooted,
internally representation-summed norm, denoted
\(\|\cdot\|_{\mathfrak T_n}\).  Assume composition with a covariance in
one argument obeys
\begin{equation}
 \|A\circ_i C_{\mathrm{KK}}\|_{\mathfrak T_n}
 \le
 \|A\|_{\mathfrak T_n}
 \|C_{\mathrm{KK}}\|_{\mathfrak R_0}.
 \label{eq:v5v37-leg-composition}
\end{equation}
This is the rooted Schur convolution estimate with internal Schatten
Hölder.

\begin{corollary}[Gaussian \(n\)-leg suppression]
\label{cor:v5v37-nleg}
If
\[
 \sup_q\|D^nW(q)\|_{\mathfrak T_n}\le M_n,
\]
then
\begin{equation}
 \left\|
 \kappa_0(q^{\otimes n},W)
 \right\|_{\mathfrak T_n}
 \le
 M_n
 \left(C_{a,b}e^{-a/u^2}\right)^n.
 \label{eq:v5v37-nleg}
\end{equation}
\end{corollary}

\begin{proof}
Use Theorem~\ref{thm:v5v37-Gaussian-leg}.  Each of its \(n\) arguments is
composed with \(C_{\mathrm{KK}}\).  Apply
\eqref{eq:v5v37-leg-composition} successively and then
\eqref{eq:v5v37-shell-root}.
\end{proof}

This is an exact first interpolation derivative at \(s=0\).  Controlling
the integral over \(0<s\le1\) in
\eqref{eq:v5v37-integrated-interpolation} requires a connected cluster
bound for the interacting tilted measures.

\subsection{A sufficient all-\(s\) tensor certificate}

\begin{definition}[Conditional connected tensor certificate]
\label{def:v5v37-CCTC}
A CCTC with constants \(A,B,\sigma\ge0\) and decay
\(\chi(u,L)\) is the uniform estimate
\begin{equation}
 \left\|
 \kappa_s(q^{\otimes n},W)
 \right\|_{\mathfrak T_n}
 \le
 A B^n(n!)^\sigma\chi(u,L)^n
 \label{eq:v5v37-CCTC}
\end{equation}
for all \(n\ge1\), \(0\le s\le1\), and all admissible retained fields.
\end{definition}

\begin{theorem}[CCTC controls the augmented chart]
\label{thm:v5v37-CCTC-chart}
Assume a CCTC.  Let
\(\kappa_V^{(n)}\) be the Taylor coefficients of
\(\mathcal K_1(J)\).  Then
\begin{equation}
 \|\kappa_V^{(n)}-\kappa_0^{(n)}\|_{\mathfrak T_n}
 \le AB^n(n!)^\sigma\chi(u,L)^n.
 \label{eq:v5v37-CCTC-coefficients}
\end{equation}
Moreover,
\begin{equation}
 \sum_{n\ge1}
 \frac{\eta^n}{(n!)^\beta}
 \|\kappa_V^{(n)}-\kappa_0^{(n)}\|_{\mathfrak T_n}
 <\infty
 \label{eq:v5v37-CCTC-sum}
\end{equation}
if either \(\beta>\sigma\), or
\(\beta=\sigma\) and \(\eta B\chi(u,L)<1\).
\end{theorem}

\begin{proof}
Integrate \eqref{eq:v5v37-CCTC} in
\eqref{eq:v5v37-integrated-interpolation}; the interval has length one,
giving \eqref{eq:v5v37-CCTC-coefficients}.  The series in
\eqref{eq:v5v37-CCTC-sum} is bounded by
\[
 A\sum_{n\ge1}
 \frac{(\eta B\chi)^n}{(n!)^{\beta-\sigma}}.
\]
It converges for every finite \(\eta B\chi\) when
\(\beta>\sigma\), and is geometric under the stated condition when
\(\beta=\sigma\).
\end{proof}

\begin{corollary}[Rooted local chart decoupling]
\label{cor:v5v37-rooted-chart-decoupling}
If the CCTC holds with
\begin{equation}
 \chi(u,L)\le C e^{-a/u^2},
 \label{eq:v5v37-chi}
\end{equation}
then the interacting correction to every fixed conditional tensor tends
to zero as \(u\to0\), and the complete factorially weighted correction
tends to zero whenever its series is dominated uniformly in a
neighborhood of that limit.
\end{corollary}

\begin{proof}
Use \eqref{eq:v5v37-CCTC-coefficients} for fixed \(n\).
For the full series, apply dominated convergence to
\eqref{eq:v5v37-CCTC-sum}.
\end{proof}

\begin{firewall}[The CCTC is the live estimate]
\label{fw:v5v37-CCTC-open}
The Gaussian endpoint estimate
\eqref{eq:v5v37-nleg} does not prove the all-\(s\) CCTC.  Interacting
connected graphs can contain arbitrarily many internal vertices.  A
Ward-complete rooted cluster expansion must show that a spanning family of
massive covariance lines retains enough factors
\(e^{-a/u^2}\) after all representation and graph sums.
\end{firewall}

\subsection{Insertion into the crossover contraction}

Let
\[
 \|\mathcal K\|_{\eta,\beta}
 =
 \sum_{n\ge1}\frac{\eta^n}{(n!)^\beta}
 \|D_J^n\mathcal K(0)\|_{\mathfrak T_n}.
\]
Suppose the full shell map has V26 contraction constant \(q\), and the
augmented-chart correction has Lipschitz constant \(L_{\mathcal K}\) in
this norm.

\begin{proposition}[Conditional-tensor Schur budget]
\label{prop:v5v37-Schur-budget}
If the conditional tensor is added as a fourth block to the V26
head--protected--tail decomposition and its self-row contracts by
\(q_{\mathcal K}<1\), then the nonnegative matrix criterion of
Theorem~\ref{thm:v5v26-Schur} applies after enlarging \(A\) or the tail
block.  A sufficient scalar perturbative condition in a pre-existing
weighted norm is
\begin{equation}
 q+L_{\mathcal K}<1.
 \label{eq:v5v37-simple-budget}
\end{equation}
\end{proposition}

\begin{proof}
The block theorem is unchanged for a larger finite block matrix: a
positive supersolution gives a weighted max-norm contraction exactly as
in Lemma~\ref{lem:v5v26-weighted-max}.  If the augmented contribution is
estimated as one perturbation of the old map, its Lipschitz constant adds
to \(q\), proving \eqref{eq:v5v37-simple-budget}.
\end{proof}

\subsection{V37 status}

\begin{theorem}[V37 conclusion]
\label{thm:v5v37-status}
The conditional massive tensor has the exact connected interpolation
\eqref{eq:v5v37-integrated-interpolation}.  At the Gaussian endpoint,
every external massive leg carries one covariance, and the massive
covariance is exponentially small in rooted norms as \(u\to0\), without
the global trace multiplicity.  An all-\(s\) CCTC would control the complete
augmented chart in the V25 factorial norm, but that interacting connected
estimate remains open.
\end{theorem}

\begin{proof}
Combine Theorem~\ref{thm:v5v37-interpolation},
Theorem~\ref{thm:v5v37-Gaussian-leg},
Corollary~\ref{cor:v5v37-shell-root}, and
Theorem~\ref{thm:v5v37-CCTC-chart}.
\end{proof}

V38 will connect the CCTC to the inherited
Koteck\'y--Preiss polymer theorem.  It will state a massive-line incidence
condition, prove the resulting tree extraction, and determine the exact
smallness parameter needed to sum all conditional connected graphs.

\clearpage
\section{V38: Massive-line tree extraction for the conditional tensor}
\label{sec:v5v38}

\subsection{Purpose and relation to the inherited polymer theorem}

The fifth-series summary already records a Banach-valued
Koteck\'y--Preiss theorem for Ward-complete small-field activities.  V38
does not repeat that general theorem.  It identifies a sufficient
massive-line incidence hypothesis under which its connected cluster sum
implies the CCTC left open in V37, and it computes the resulting smallness
parameter.

Let
\begin{equation}
 \chi_u=C_Ce^{-a/u^2},\qquad 0<u\le1,
 \label{eq:v5v38-chi}
\end{equation}
where \(C_C\) is the rooted covariance constant in
\eqref{eq:v5v37-shell-root}.  Let \(\zeta\ge0\) be the normed load of one
renormalized, Ward-complete massive interaction activity after all
self-contractions and local Taylor owners have been assigned.  No
coefficientwise representation supremum is taken in either constant.

\subsection{Source shifting and external covariance lines}

For a centered massive Gaussian \(q\sim N(0,C)\), completion of the square
gives
\begin{equation}
 \mathbb E_Ce^{\langle J,q\rangle-sW(q)}
 =
 e^{\frac12\langle J,CJ\rangle}
 \mathbb E_Ce^{-sW(q+CJ)}.
 \label{eq:v5v38-source-shift}
\end{equation}

\begin{lemma}[Every interaction-source derivative carries \(C\)]
\label{lem:v5v38-source-C}
For every \(n\ge1\), each derivative in \(J\) of the interaction factor on
the right of \eqref{eq:v5v38-source-shift} differentiates \(W(q+CJ)\) in
a direction \(Cf_i\).  Thus every external source leg attached to an
interaction activity contains one massive covariance before any graph or
representation norm is taken.
\end{lemma}

\begin{proof}
The chain rule gives
\[
 D_JW(q+CJ)[f]=DW(q+CJ)[Cf].
\]
Iterating the chain rule proves the assertion.  Derivatives of the
separate Gaussian factor
\(e^{\langle J,CJ\rangle/2}\) generate only the base covariance
\eqref{eq:v5v37-Gaussian-cumulants}; they are not interaction attachments.
\end{proof}

\begin{corollary}[Rooted external-leg bound]
\label{cor:v5v38-external-leg}
If attachment of one differentiated interaction leg has rooted operator
constant \(B\), then \(n\) labelled external legs contribute at most
\begin{equation}
 (B\chi_u)^n
 \label{eq:v5v38-external-factor}
\end{equation}
before the choice of their incident interaction vertices is counted.
\end{corollary}

\begin{proof}
Apply the rooted convolution estimate
\eqref{eq:v5v37-leg-composition} once for every covariance \(C\) supplied
by Lemma~\ref{lem:v5v38-source-C}, then use
\(\|C\|_{\mathfrak R_0}\le\chi_u\).
\end{proof}

\subsection{The massive incidence tree certificate}

\begin{definition}[Massive incidence tree certificate]
\label{def:v5v38-MITC}
An MITC with constants \((\zeta,B,\chi_u)\) consists of the following
post-assembly bounds for every connected conditional cluster:
\begin{enumerate}
 \item a cluster with \(N\ge1\) internal Ward-complete activities admits a
 tree-graph majorant summed over the \(N^{N-2}\) labelled trees, with the
 convention \(1^{-1}=1\);
 \item every internal spanning-tree edge is a massive covariance
 contraction of rooted norm at most \(\chi_u\);
 \item the product of the \(N\) vertex loads, including symmetry factor,
 is at most \(\zeta^N/N!\);
 \item each labelled external source leg chooses at most \(N\) incident
 vertices and, after that choice, costs at most \(B\chi_u\);
 \item loop, representation, local-cumulant, and rooted spatial sums left
 after tree extraction have total norm at most one, or are included in
 \(\zeta\) and \(B\).
\end{enumerate}
\end{definition}

The MITC is a specialization of the inherited Banach-valued polymer
bound.  Its substantive content is item 2: no internal connectedness edge
may be replaced by an uncentered complete moment before the tree is
extracted.

\begin{lemma}[Cayley majorant]
\label{lem:v5v38-Cayley-majorant}
For \(N\ge1\),
\begin{equation}
 \frac{N^{N-2}}{N!}\le\frac{e^N}{N^2},
 \label{eq:v5v38-Cayley-majorant}
\end{equation}
with the right side interpreted harmlessly at \(N=1\).
\end{lemma}

\begin{proof}
The elementary Stirling lower bound
\(N!\ge(N/e)^N\) gives
\[
 \frac{N^{N-2}}{N!}
 \le
 \frac{N^{N-2}}{(N/e)^N}
 =\frac{e^N}{N^2}.
\]
\end{proof}

\begin{lemma}[Power-series attachment bound]
\label{lem:v5v38-attachment-series}
For \(0\le q<1\) and every integer \(n\ge0\),
\begin{equation}
 \sum_{N\ge1}N^nq^{N-1}
 \le\frac{n!}{(1-q)^{n+1}}.
 \label{eq:v5v38-attachment-series}
\end{equation}
\end{lemma}

\begin{proof}
For \(N\ge1\),
\[
 N^n\le N(N+1)\cdots(N+n-1)
 =n!\binom{N+n-1}{n}.
\]
The binomial generating function gives
\[
 \sum_{N\ge1}\binom{N+n-1}{n}q^{N-1}
 =(1-q)^{-n-1}.
\]
\end{proof}

\subsection{All-order connected tensor bound}

\begin{theorem}[MITC implies a CCTC]
\label{thm:v5v38-MITC-CCTC}
Assume an MITC and set
\begin{equation}
 q_u=e\zeta\chi_u.
 \label{eq:v5v38-qu}
\end{equation}
If \(q_u<1\), the sum of all connected interaction clusters with \(n\ge1\)
labelled external massive source legs obeys
\begin{equation}
 \|\mathcal C_n\|_{\mathfrak T_n}
 \le
 \frac{e\zeta}{1-q_u}\,
 n!\left(\frac{B\chi_u}{1-q_u}\right)^n.
 \label{eq:v5v38-Cn}
\end{equation}
Consequently the CCTC holds with
\begin{equation}
 A=\frac{e\zeta}{1-q_u},\qquad
 \sigma=1,\qquad
 B_{\mathrm{CCTC}}=\frac{B}{1-q_u},\qquad
 \chi(u,L)=\chi_u.
 \label{eq:v5v38-CCTC-constants}
\end{equation}
\end{theorem}

\begin{proof}
For fixed \(N\), the MITC gives the connected-tree majorant
\[
 \frac{\zeta^N}{N!}
 N^{N-2}\chi_u^{N-1}
 N^n(B\chi_u)^n.
\]
Sum over \(N\ge1\).  By
Lemma~\ref{lem:v5v38-Cayley-majorant},
\begin{align*}
 \|\mathcal C_n\|
 &\le
 e\zeta(B\chi_u)^n
 \sum_{N\ge1}N^{n-2}q_u^{N-1}\\
 &\le
 e\zeta(B\chi_u)^n
 \sum_{N\ge1}N^nq_u^{N-1}.
\end{align*}
Apply Lemma~\ref{lem:v5v38-attachment-series} and rearrange to obtain
\eqref{eq:v5v38-Cn}.  Comparison with
Definition~\ref{def:v5v37-CCTC} gives
\eqref{eq:v5v38-CCTC-constants}.
\end{proof}

\begin{remark}[The constant is sufficient, not sharp]
Cayley counting, the Stirling lower bound, and the replacement
\(N^{n-2}\le N^n\) deliberately favor a transparent sufficient condition.
The important structural fact is the product
\(\zeta\chi_u\): one activity load is paid for each vertex and one massive
rooted covariance for each internal tree edge.
\end{remark}

\begin{corollary}[Factorially weighted augmented chart]
\label{cor:v5v38-weighted-chart}
Under the MITC and \(q_u<1\), the interaction part of the conditional
generator belongs to the V25 particle-number norm
\begin{equation}
 \sum_{n\ge1}
 \frac{\eta^n}{(n!)^\beta}
 \|\mathcal C_n\|_{\mathfrak T_n}
 \label{eq:v5v38-weighted-chart}
\end{equation}
whenever either \(\beta>1\), or
\begin{equation}
 \beta=1,\qquad
 \frac{\eta B\chi_u}{1-q_u}<1.
 \label{eq:v5v38-beta-one}
\end{equation}
\end{corollary}

\begin{proof}
Insert \eqref{eq:v5v38-Cn}.  The remaining scalar series is
\[
 \frac{e\zeta}{1-q_u}
 \sum_{n\ge1}
 \frac1{(n!)^{\beta-1}}
 \left(\frac{\eta B\chi_u}{1-q_u}\right)^n.
\]
It converges under the stated alternatives.
\end{proof}

\subsection{Uniformity along the interaction interpolation}

The CCTC in V37 is required for \(0\le s\le1\), not only at the final
interaction.  Let \(\zeta_s\) be the activity load for \(sW\).

\begin{proposition}[Uniform interpolation criterion]
\label{prop:v5v38-uniform-s}
If the MITC holds for every \(s\in[0,1]\) with common \(B,\chi_u\) and
\begin{equation}
 \sup_{0\le s\le1}\zeta_s\le\zeta,
 \qquad e\zeta\chi_u<1,
 \label{eq:v5v38-uniform-s}
\end{equation}
then the bound \eqref{eq:v5v38-Cn} holds uniformly along the interpolation.
If one distinguished derivative activity represents
\(\partial_s(sW)=W\) and obeys the same vertex bound, then the joint
cumulants in \eqref{eq:v5v37-tensor-interpolation} satisfy the CCTC.
\end{proposition}

\begin{proof}
The proof of Theorem~\ref{thm:v5v38-MITC-CCTC} uses only the common upper
bound for the vertex load and the common covariance and attachment
constants.  Marking one vertex as the \(s\)-derivative does not change the
tree count; its norm is absorbed in the same \(\zeta\) bound.  Therefore
the result is uniform in \(s\) and applies to the differentiated connected
clusters.
\end{proof}

\begin{corollary}[Completion of the V37 implication]
\label{cor:v5v38-complete-V37}
Under Proposition~\ref{prop:v5v38-uniform-s}, the conditional tensors obey
\eqref{eq:v5v37-CCTC-coefficients} with the constants
\eqref{eq:v5v38-CCTC-constants}.  Hence the augmented chart is
factorially summable under
Corollary~\ref{cor:v5v38-weighted-chart}.
\end{corollary}

\begin{proof}
Apply Proposition~\ref{prop:v5v38-uniform-s} in
Theorem~\ref{thm:v5v37-CCTC-chart}.
\end{proof}

\subsection{Where the massive incidence can fail}

\begin{theorem}[MITC failure alternatives]
\label{thm:v5v38-failure-alternatives}
For a connected conditional graph expansion, failure of item 2 of the
MITC requires at least one of:
\begin{enumerate}
 \item a connectedness link is supplied by an uncentered moment rather
 than a massive covariance;
 \item a local cumulant and an external or internal role coincide before
 the Ward-complete operator is assembled;
 \item a zero-mode propagator enters what was declared to be the massive
 conditional cluster;
 \item the representation or rooted spatial sum is taken before tree
 extraction and destroys the covariance incidence;
 \item a non-Gaussian reference measure has no covariance-line
 interpolation with an analogous small connector.
\end{enumerate}
\end{theorem}

\begin{proof}
If none occurs, every inter-activity connectedness link in the Gaussian
massive integration is a contraction by \(C_{\mathrm{KK}}\).  Any connected
\(N\)-vertex graph contains a spanning tree of \(N-1\) such links.  The
external source shift supplies its own covariance by
Lemma~\ref{lem:v5v38-source-C}, and the sums are retained until after this
selection.  These are exactly items 2 and 4 of the MITC.
\end{proof}

\begin{assessment}[Companion alignment]
\label{ass:v5v38-companion-alignment}
The alternatives mirror the boundary in YM V41: a dangerous residual must
lose a connected prefix, use an uncentered moment, grow a representation
inside the prefix, or merge roles before scalar factorization.  V38 does
not import the YM estimate; it derives the analogous Gaussian incidence
list for the conditional massive field.
\end{assessment}

\subsection{Collapsed and crossover regimes}

\begin{corollary}[Automatic collapsed-regime smallness]
\label{cor:v5v38-collapsed-smallness}
If \(\zeta\) is uniformly bounded along a cofinal family and
\(u\to0\), then
\begin{equation}
 q_u=e\zeta C_Ce^{-a/u^2}\longrightarrow0.
 \label{eq:v5v38-qu-zero}
\end{equation}
Thus the MITC tree series is eventually summable independently of the
global trace multiplicity.
\end{corollary}

\begin{proof}
Exponential decay of \(e^{-a/u^2}\) dominates the fixed activity bound.
\end{proof}

\begin{corollary}[Crossover smallness condition]
\label{cor:v5v38-crossover-smallness}
On a fixed overlap \(u\in[u_-,u_+]\subset(0,\infty)\), a sufficient
uniform tree condition is
\begin{equation}
 e\zeta C_Ce^{-a/u_+^2}<1.
 \label{eq:v5v38-crossover-smallness}
\end{equation}
\end{corollary}

\begin{proof}
The function \(e^{-a/u^2}\) increases with \(u\).  Therefore its largest
value on the overlap occurs at \(u_+\).  Insert it into
\eqref{eq:v5v38-qu}.
\end{proof}

\begin{firewall}[Large-field and metric sectors]
\label{fw:v5v38-large-field}
The activity load \(\zeta\) is not known to be uniformly bounded for the
complete SM--Einstein shell.  The real conformal instability and large
Higgs/gauge fields lie outside a small-field KP estimate unless a separate
large-field domination theorem is supplied.  Exponential massive-line
suppression does not control an unbounded vertex load.
\end{firewall}

\subsection{Ward and local-owner compatibility}

\begin{proposition}[Assembly order preserves the protected packet]
\label{prop:v5v38-Ward-assembly}
Suppose each activity is the complete Ward packet assigned to one polymer,
all local self-contractions are placed in their V28 counterterm owners, and
the tree-graph inequality is applied only after those assemblies.  Then
the MITC majorant does not split a protected Ward cancellation.
\end{proposition}

\begin{proof}
The tree extraction acts on the incompatibility graph of already assembled
Banach-valued activities.  It never replaces an activity by the absolute
values of its diagrammatic summands.  Thus the protected projection in the
V25 hybrid norm is evaluated before the nonnegative cluster majorant.
\end{proof}

\begin{firewall}[Verification still required]
\label{fw:v5v38-MITC-open}
Proposition~\ref{prop:v5v38-Ward-assembly} explains how an MITC can coexist
with the signed protected sector.  It does not prove that the actual
gauge-fixed chiral SM--Einstein activities satisfy the MITC, the inherited
KP hypothesis, or the required regulator Ward identities.
\end{firewall}

\subsection{V38 status}

\begin{theorem}[V38 conclusion]
\label{thm:v5v38-status}
An MITC converts the rooted massive covariance suppression into an
all-order conditional connected tensor bound.  The explicit cluster
parameter is
\(q_u=e\zeta C_Ce^{-a/u^2}\), and the \(n\)-tensor bound is
\eqref{eq:v5v38-Cn}.  It closes the factorial augmented chart whenever
\eqref{eq:v5v38-beta-one} holds.  Verification of the MITC and a uniform
activity load for the complete coupled shell remain open.
\end{theorem}

\begin{proof}
Use Theorem~\ref{thm:v5v38-MITC-CCTC},
Proposition~\ref{prop:v5v38-uniform-s}, and
Corollary~\ref{cor:v5v38-weighted-chart}.
\end{proof}

V39 will combine the MITC bound with the two-chart transition.  It will
test whether a uniformly bi-Lipschitz inverse is possible as
\(u\to0\), and, if not, replace it by a precise local-observable quotient
with a controlled reconstruction only on the finite crossover overlap.

\clearpage
\section{V39: The collapsed boundary is a quotient, not a bi-Lipschitz chart}
\label{sec:v5v39}

\subsection{Inverse conditioning of the augmented data}

V36 restores injectivity of the reduced chart by adjoining the conditional
generator \(\mathcal K_p\).  Injectivity at each finite cutoff does not
imply a uniformly bounded inverse as the massive covariance tends to zero.
The same one-dimensional tilted Gaussian used in V36 gives the sharp
obstruction.

Let \(q\sim N(0,\sigma^2)\) and
\[
 V_t(q)=tq+\frac12t^2\sigma^2.
\]
Its reduced interaction is zero and its conditional law is
\(N(-t\sigma^2,\sigma^2)\).  Hence
\begin{equation}
 \mathcal K_t(J)
 =
 -t\sigma^2J+\frac12\sigma^2J^2.
 \label{eq:v5v39-linear-generator}
\end{equation}

\begin{theorem}[Inverse chart blow-up]
\label{thm:v5v39-inverse-blowup}
Measure the affine interaction difference \(V_t-V_0\) by the absolute
value \(|t|\) of its linear coefficient, and measure the corresponding
conditional-generator difference in any tensor norm whose first-tensor
weight is \(\eta>0\).  Every inverse Lipschitz constant satisfies
\begin{equation}
 L_{\mathcal S^{-1}}
 \ge\frac1{\eta\sigma^2}.
 \label{eq:v5v39-inverse-blowup}
\end{equation}
For a massive mode with
\(\sigma^2\le C e^{-a/u^2}\), this lower bound grows at least like a
constant times \(e^{a/u^2}\) as \(u\to0\).
\end{theorem}

\begin{proof}
By \eqref{eq:v5v39-linear-generator}, the difference of the augmented
conditional data has only the first connected tensor
\(-t\sigma^2\); the common Gaussian covariance cancels.  Its weighted norm
is \(\eta|t|\sigma^2\).  An inverse Lipschitz estimate would require
\[
 |t|
 \le L_{\mathcal S^{-1}}\eta|t|\sigma^2.
\]
Cancel nonzero \(|t|\) to obtain
\eqref{eq:v5v39-inverse-blowup}.  Insert the massive covariance bound.
\end{proof}

\begin{corollary}[No uniform bi-Lipschitz collapsed atlas]
\label{cor:v5v39-no-bilipschitz}
The augmented four- to three-dimensional chart cannot satisfy the uniform
bi-Lipschitz condition \eqref{eq:v5v36-bi-Lipschitz} on a family with
\(u\to0\), when the full interaction norm assigns a nonvanishing weight to
the disappearing massive linear coefficient.
\end{corollary}

\begin{proof}
The lower bound in
Theorem~\ref{thm:v5v39-inverse-blowup} diverges.
\end{proof}

\begin{assessment}[Geometric meaning]
\label{ass:v5v39-blowdown}
At the collapsed boundary, many distinct massive interactions have the
same limiting zero-mode functional because the massive fluctuations cease
to probe them.  The boundary is a blowdown of those fibre directions, not
an ordinary Banach-manifold chart overlap.
\end{assessment}

\subsection{The exact zero-mode quotient}

Let \(X\) be the space of full shell data and let
\[
 S_0:X\to Y,\qquad S_0(V)=V_{\mathrm{red}},
\]
be exact massive integration.  Define
\begin{equation}
 x\sim_0x'
 \quad\Longleftrightarrow\quad
 S_0x=S_0x'.
 \label{eq:v5v39-equivalence}
\end{equation}
The quotient \(X/{\sim_0}\) is identified with the image of \(S_0\).

\begin{theorem}[Exactness on the zero-mode observable algebra]
\label{thm:v5v39-zero-observables}
For every observable \(F(p)\) independent of the massive field,
its normalized expectation depends on full shell data only through
\(V_{\mathrm{red}}\).  Therefore equivalent data in
\eqref{eq:v5v39-equivalence} define the same state on the zero-mode
observable algebra.
\end{theorem}

\begin{proof}
For \(F(p,q)=F(p)\), equation
\eqref{eq:v5v35-reduced-observable} gives
\(F_{\mathrm{red}}(p)=F(p)\).  The exact reduction identity
\eqref{eq:v5v35-exact-reduction} then depends only on
\(V_{\mathrm{red}}\).
\end{proof}

\begin{definition}[Quotient-compatible shell map]
\label{def:v5v39-quotient-map}
A full shell map \(T:X\to X\) is quotient compatible if
\begin{equation}
 S_0T=\overline T S_0
 \label{eq:v5v39-quotient-intertwining}
\end{equation}
for a map \(\overline T:Y\to Y\).
\end{definition}

\begin{theorem}[Contraction on the physical quotient]
\label{thm:v5v39-quotient-contraction}
Suppose \(S_0\) is onto \(Y\), \(T\) is quotient compatible, and
\(\overline T\) is a strict contraction on a closed ball in \(Y\).
Then there is a unique fixed equivalence class of shell data on that ball.
All zero-mode observables have unique fixed-point expectations.
\end{theorem}

\begin{proof}
Banach's theorem gives a unique \(y_*\) satisfying
\(\overline T y_*=y_*\).  The fibre \(S_0^{-1}(y_*)\) is the unique fixed
class in \(X/{\sim_0}\).  Indeed, the induced quotient map is identified
with \(\overline T\) by \eqref{eq:v5v39-quotient-intertwining}.  Apply
Theorem~\ref{thm:v5v39-zero-observables} to obtain uniqueness of the
stated expectations.
\end{proof}

\begin{firewall}[A fixed class is not a fixed representative]
\label{fw:v5v39-class-not-point}
Theorem~\ref{thm:v5v39-quotient-contraction} does not produce
\(x_*\in X\) with \(Tx_*=x_*\).  The map can drift or have multiple fixed
points inside the fibre \(S_0^{-1}(y_*)\).  Full four-dimensional
observables remain undetermined without a fibre estimate.
\end{firewall}

\subsection{Lifting the quotient by a fibre contraction}

Assume locally that full data have coordinates \((y,z)\in Y\oplus F\),
where \(y\) is the reduced functional and \(z\) is the conditional massive
tensor.  A quotient-compatible map is triangular:
\begin{equation}
 T(y,z)=\bigl(\overline T(y),G(y,z)\bigr).
 \label{eq:v5v39-triangular-map}
\end{equation}

\begin{theorem}[Base--fibre contraction]
\label{thm:v5v39-base-fibre}
Suppose
\begin{align}
 \|\overline T(y)-\overline T(\widetilde y)\|_Y
 &\le q_B\|y-\widetilde y\|_Y,\qquad q_B<1,
 \label{eq:v5v39-base-q}\\
 \|G(y,z)-G(y,\widetilde z)\|_F
 &\le q_F\|z-\widetilde z\|_F,\qquad q_F<1,
 \label{eq:v5v39-fibre-q}\\
 \|G(y,z)-G(\widetilde y,z)\|_F
 &\le L_{FB}\|y-\widetilde y\|_Y.
 \label{eq:v5v39-fibre-base}
\end{align}
Then \(T\) is a strict contraction in an equivalent weighted max norm and
has a unique full fixed point on every invariant complete ball.
\end{theorem}

\begin{proof}
The nonnegative Lipschitz matrix is
\[
 L=
 \begin{pmatrix}
 q_B&0\\
 L_{FB}&q_F
 \end{pmatrix}.
\]
Its spectral radius is \(\max(q_B,q_F)<1\).  Explicitly, set \(w_B=1\)
and choose
\[
 w_F>\frac{L_{FB}}{1-q_F}.
\]
Then
\[
 L\binom{w_B}{w_F}
 =
 \binom{q_B}{L_{FB}+q_Fw_F}
 <
 \binom{w_B}{w_F}.
\]
Lemma~\ref{lem:v5v26-weighted-max} gives a contraction in the associated
weighted norm.  Banach's theorem gives the fixed point.
\end{proof}

\begin{corollary}[Collapsed massive fibre under the MITC]
\label{cor:v5v39-MITC-fibre}
If the massive conditional-tensor row obeys the V38 MITC and its nonlinear
Lipschitz constant satisfies
\begin{equation}
 q_F(u)\le
 \frac{K\zeta\chi_u}{1-e\zeta\chi_u},
 \label{eq:v5v39-qF}
\end{equation}
then \(q_F(u)\to0\) as \(u\to0\) for bounded \(\zeta\).  The fibre
contraction is eventually stronger even though the inverse chart constant
diverges.
\end{corollary}

\begin{proof}
By \eqref{eq:v5v38-chi}, \(\chi_u\to0\).  The numerator in
\eqref{eq:v5v39-qF} tends to zero and the denominator tends to one.
\end{proof}

There is no contradiction between
Theorem~\ref{thm:v5v39-inverse-blowup} and
Corollary~\ref{cor:v5v39-MITC-fibre}: reconstruction amplifies vanishing
coordinates, while the forward dynamics contracts those same coordinates.

\subsection{Smeared local observables}

The zero-mode algebra is smaller than the local limiting algebra suggested
by rooted covariance decay.  Let \(f_1,\ldots,f_k\) be fixed normalized
test sections of bounded support and put
\[
 Q_f=(q(f_1),\ldots,q(f_k)).
\]

\begin{lemma}[Gaussian local-cylinder decay]
\label{lem:v5v39-cylinder}
Suppose the massive covariance obeys both rooted Schur bounds
\[
 \sup_x\int\|C(x,y)\|\,dy\le\chi_u,\qquad
 \sup_y\int\|C(x,y)\|\,dx\le\chi_u.
\]
Then the covariance operator has \(L^2\)-norm at most \(\chi_u\), and
\begin{equation}
 \mathbb E|Q_f|^2
 \le\chi_u\sum_{i=1}^k\|f_i\|_2^2.
 \label{eq:v5v39-cylinder-variance}
\end{equation}
\end{lemma}

\begin{proof}
The Schur test gives \(\|C\|_{L^2\to L^2}\le\chi_u\).  Therefore
\[
 \mathbb E|Q_f|^2
 =\sum_i\langle f_i,Cf_i\rangle
 \le\chi_u\sum_i\|f_i\|_2^2.
\]
\end{proof}

\begin{corollary}[Gaussian cylinder quotient]
\label{cor:v5v39-cylinder-quotient}
If \(F:\mathbb R^k\to\mathbb R\) is \(M\)-Lipschitz, then
\begin{equation}
 |\mathbb EF(Q_f)-F(0)|
 \le
 M\chi_u^{1/2}
 \left(\sum_i\|f_i\|_2^2\right)^{1/2}.
 \label{eq:v5v39-cylinder-decay}
\end{equation}
\end{corollary}

\begin{proof}
Use the Lipschitz bound, Cauchy--Schwarz, and
\eqref{eq:v5v39-cylinder-variance}.
\end{proof}

\begin{assessment}[Renormalized composites]
\label{ass:v5v39-composites}
Point fields and local polynomial composites are not Lipschitz cylinder
functions; their coincident contractions can diverge even when the rooted
operator norm is small.  They enter the quotient only after the local
Taylor/heat coefficients are assigned to renormalized composite owners and
their connected remainders satisfy the MITC.
\end{assessment}

\subsection{An active crossover band}

A stable inverse is possible on massive directions whose covariance is
bounded below.  Fix \(0<\delta<D<\infty\) and let
\begin{equation}
 P_{\mathrm{act}}
 =
 Q_0\mathbf1_{[\delta,D]}(\Delta_\varepsilon/L^2).
 \label{eq:v5v39-active-band}
\end{equation}
Since \(f(x)>0\) for \(x>0\),
\begin{equation}
 m_{\delta,D}:=\min_{\delta\le x\le D}f(x)>0.
 \label{eq:v5v39-band-floor}
\end{equation}

\begin{proposition}[Stable affine reconstruction on the active band]
\label{prop:v5v39-active-reconstruction}
For an affine active-band interaction
\[
 V_t(q)=\langle t,q\rangle+\frac12\langle t,Ct\rangle,
\]
the conditional mean is \(m=-Ct\).  On
\(\operatorname{Ran}P_{\mathrm{act}}\),
\begin{equation}
 \|t-\widetilde t\|
 \le m_{\delta,D}^{-1}\|m-\widetilde m\|.
 \label{eq:v5v39-active-inverse}
\end{equation}
\end{proposition}

\begin{proof}
The covariance eigenvalues on the active band are
\(f(\lambda/L^2)\ge m_{\delta,D}\).  Thus
\(\|C^{-1}\|\le m_{\delta,D}^{-1}\), and
\(t-\widetilde t=-C^{-1}(m-\widetilde m)\).
\end{proof}

\begin{definition}[Stratified crossover certificate]
\label{def:v5v39-SCC}
An SCC consists of:
\begin{enumerate}
 \item the reduced base \(V_{\mathrm{red}}\);
 \item augmented conditional tensors on an active band with a uniform
 covariance floor and a proved inverse estimate;
 \item a silent massive tail treated as a quotient for the declared local
 observable algebra;
 \item a V38 fibre contraction for that tail;
 \item local composite counterterms and Ward identities compatible with
 the three blocks;
 \item convergence of the active-band thresholds without spectral
 crossings, or the V29 moving-projector connection estimates.
\end{enumerate}
\end{definition}

\begin{firewall}[Affine stability is not nonlinear reconstruction]
\label{fw:v5v39-nonlinear-inverse}
Proposition~\ref{prop:v5v39-active-reconstruction} proves only the affine
model.  For nonlinear conditional densities, a uniform inverse requires
lower and upper density bounds, analytic control of the logarithm, and
moment determinacy in the augmented tensor norm.  These remain part of the
SCC.
\end{firewall}

\subsection{V39 status}

\begin{theorem}[V39 conclusion]
\label{thm:v5v39-status}
The augmented conditional chart is injective at finite cutoff but cannot
have a uniformly bounded inverse at the collapsed boundary; the inverse
constant grows at least as the reciprocal massive covariance.  The exact
limiting structure is a reduced quotient together with a separately
contracting conditional fibre.  A triangular base--fibre estimate yields a
unique full fixed point when both rows contract.  Stable reconstruction can
be confined to an active crossover band, while the silent tail requires a
declared local-observable quotient and renormalized composite control.
\end{theorem}

\begin{proof}
Combine Theorem~\ref{thm:v5v39-inverse-blowup},
Theorem~\ref{thm:v5v39-quotient-contraction},
Theorem~\ref{thm:v5v39-base-fibre}, and
Proposition~\ref{prop:v5v39-active-reconstruction}.
\end{proof}

V40 begins with the compulsory companion audit.  It will then formulate a
stratified inverse-limit theorem for shell sequences that move between the
four-dimensional chart, the active crossover band, and the collapsed
quotient, with explicit compatibility and compactness conditions.

\clearpage
\section{V40: A stratified inverse-limit theorem for dimensional crossover}
\label{sec:v5v40}

\subsection{Compulsory companion audit}

The newest companion filenames at this checkpoint are
\begin{align*}
 &\texttt{Renewal\_Yang\_Mills\_Mass\_Gap\_Research\_Notes\_V46.tex},
 \qquad
 \operatorname{SHA256}=
 \texttt{17188B235E52F6199045644679680C8F00F07F315C2365CF88260D014C9159B5},
 \\
 &\texttt{Renewal\_NS\_Stability\_Research\_Notes\_V31.tex},
 \qquad
 \operatorname{SHA256}=
 \texttt{F1121B62238AD5491BB7CF03635EA9934942EDF573ED8FAAA9EE79E1D38A6696}.
\end{align*}
The YM V46 file is byte-for-byte identical to the simultaneously present
V45 file.  Its mathematical end record is V46.  The NS file is unchanged
from the V35 audit.

YM V46 proves that, at every fixed cumulant order, two selected centered
legs retain their character-gap debits even when the Casimir payer and
reduced resolvent act in the same unresolved multiplicity block.  An
equivariant weighted partial trace creates no multiplicity-count factor.
Its residual is a genuine same-coordinate joining, where the complete
local carrier must be estimated before factorization.  This supports the
representation-summed transition norm of V30--V39 and adds a restriction:
chart transition, protected Ward packet, and local payer must remain one
assembled equivariant block when their roles coincide.

NS V31 continues to require the signed formation correlation before
positive-part or annular norms.  Neither companion result proves any
SM--Einstein inverse-limit hypothesis below.

\subsection{Finite-cutoff stratified state spaces}

At shell \(j\), let
\[
 X_j^{(4)},\qquad X_j^{(A)},\qquad X_j^{(0)}
\]
denote respectively the four-dimensional chart, the active crossover
chart of Definition~\ref{def:v5v39-SCC}, and the collapsed quotient chart.
Choose compact admissible cores
\[
 K_j^{(r)}\subset X_j^{(r)},\qquad r\in\{4,A,0\},
\]
in the weak hybrid topology.  In applications their compactness is to
come from a uniform strong V25 bound followed by the strict-radius compact
embedding.

Let
\begin{equation}
 \widehat K_j=
 K_j^{(4)}\sqcup K_j^{(A)}\sqcup K_j^{(0)}
 \label{eq:v5v40-disjoint-union}
\end{equation}
and impose the equivalence relation generated by:
\begin{enumerate}
 \item the exact augmented Kaluza--Klein chart on the
 \(4\)--\(A\) overlap;
 \item the exact zero-mode reduction quotient on the \(A\)--\(0\)
 boundary;
 \item equality inside each individual chart.
\end{enumerate}
Write
\begin{equation}
 K_j=\widehat K_j/{\sim_j}.
 \label{eq:v5v40-stratified-K}
\end{equation}

\begin{theorem}[Compact Hausdorff stratified cutoff space]
\label{thm:v5v40-compact-quotient}
Assume every \(K_j^{(r)}\) is compact Hausdorff and the graph of
\(\sim_j\) is closed in
\(\widehat K_j\times\widehat K_j\).  Then \(K_j\) is compact Hausdorff.
\end{theorem}

\begin{proof}
The finite disjoint union \(\widehat K_j\) is compact Hausdorff.  Its
continuous quotient image is compact.  To prove Hausdorffness, let
\([x]\ne[y]\).  Their equivalence classes are disjoint compact subsets
because a closed equivalence relation on a compact space has closed
classes.  The closed relation prevents nets from distinct classes from
having a common quotient limit.  Equivalently, the saturation of a closed
set is the projection of a closed subset of the compact product and is
closed.  Normality of \(\widehat K_j\) separates the two compact classes
by open sets; taking complements of the saturations of their closed
complements yields disjoint saturated neighborhoods.  Their quotient
images separate \([x]\) and \([y]\).  Thus the quotient is Hausdorff.
\end{proof}

\begin{firewall}[Closedness is substantive at the collapsed boundary]
\label{fw:v5v40-closed-relation}
Pointwise injectivity of the augmented chart does not prove that
\(\sim_j\) is closed uniformly as \(u\to0\).  V39's inverse blow-up makes
this a real issue.  The quotient relation must be defined by convergence
of the declared local observable functionals, or closedness must be proved
directly.
\end{firewall}

\subsection{Descent of the shell map}

Let
\(\widehat{\mathcal R}_j:\widehat K_{j+1}\to\widehat K_j\) be the shell
map written in the appropriate source and target charts.

\begin{definition}[Stratified compatibility]
\label{def:v5v40-stratified-compatibility}
The shell map is stratified compatible if
\begin{equation}
 x\sim_{j+1}y
 \quad\Longrightarrow\quad
 \widehat{\mathcal R}_j x
 \sim_j
 \widehat{\mathcal R}_j y.
 \label{eq:v5v40-equivalence-preserved}
\end{equation}
\end{definition}

\begin{theorem}[Continuous descended shell map]
\label{thm:v5v40-descended-map}
If \(\widehat{\mathcal R}_j\) is continuous and stratified compatible,
there is a unique continuous map
\begin{equation}
 \mathcal R_j:K_{j+1}\to K_j
 \label{eq:v5v40-descended-map}
\end{equation}
satisfying
\[
 \mathcal R_j\pi_{j+1}
 =
 \pi_j\widehat{\mathcal R}_j,
\]
where \(\pi_j\) is the quotient projection.
\end{theorem}

\begin{proof}
Equation \eqref{eq:v5v40-equivalence-preserved} makes
\(\pi_j\widehat{\mathcal R}_j\) constant on every fibre of
\(\pi_{j+1}\), so it defines a unique set map \(\mathcal R_j\).
The quotient topology is final: a map from \(K_{j+1}\) is continuous
exactly when its composition with \(\pi_{j+1}\) is continuous.  The
displayed composition is continuous, proving the claim.
\end{proof}

\begin{corollary}[No inverse is needed at collapse]
\label{cor:v5v40-no-inverse}
Construction of \(\mathcal R_j\) across the \(A\)--\(0\) boundary uses
only quotient compatibility.  It does not require the divergent inverse
chart constant excluded by Theorem~\ref{thm:v5v39-inverse-blowup}.
\end{corollary}

\begin{proof}
The proof of Theorem~\ref{thm:v5v40-descended-map} uses only the forward
quotient projections and \eqref{eq:v5v40-equivalence-preserved}.
\end{proof}

\subsection{Assembled transition blocks}

Let \(\mathscr M_T\) be an unresolved multiplicity block shared by an
active-band chart transition, a protected Ward carrier, and a local
counterterm payer.  The companion audit requires their product to be
formed before the partial trace.

\begin{definition}[Coincident transition certificate]
\label{def:v5v40-CTC}
A CTC is a bound
\begin{equation}
 \sum_T w_T
 \left\|
 \operatorname{Tr}_{\mathscr M_T}
 \bigl[
 X_T\,\mathcal U_T\mathcal W_T\mathcal P_T
 \bigr]
 \right\|_1
 \le
 C_{\mathrm{tr}}
 \sum_T\|X_T\|_1,
 \label{eq:v5v40-CTC}
\end{equation}
where \(\mathcal U_T\) is the exact equivariant chart transition,
\(\mathcal W_T\) the assembled Ward/cumulant carrier,
\(\mathcal P_T\) the payer or resolvent, and all output weights are retained
inside the block.  The constant must be independent of
\(\dim\mathscr M_T\).
\end{definition}

\begin{proposition}[No-count sufficient condition]
\label{prop:v5v40-no-count}
Suppose the three carriers in \eqref{eq:v5v40-CTC} are equivariant on
complete isotypic blocks and
\begin{equation}
 \sup_T
 \|w_T\mathcal U_T\mathcal W_T\mathcal P_T\|
 \le C_{\mathrm{tr}}.
 \label{eq:v5v40-block-sup}
\end{equation}
Then the CTC holds with no multiplicity factor.
\end{proposition}

\begin{proof}
Trace-norm duality gives
\[
 \left\|\operatorname{Tr}_{\mathscr M_T}(X_TK_T)\right\|_1
 \le\|X_T\|_1\|K_T\|
\]
for the complete block operator
\(K_T=w_T\mathcal U_T\mathcal W_T\mathcal P_T\).  Sum over \(T\) and use
\eqref{eq:v5v40-block-sup}.  At no point is the multiplicity identity
expanded in a basis.
\end{proof}

\begin{firewall}[Same-coordinate transition]
\label{fw:v5v40-same-coordinate}
Separate bounds for \(\mathcal U_T,\mathcal W_T,\mathcal P_T\) do not imply
\eqref{eq:v5v40-block-sup} when their selected-leg or connectedness roles
coincide.  The complete local block must be estimated, just as the YM V46
audit leaves genuine same-coordinate joining outside its factorized
theorem.
\end{firewall}

\subsection{Existence of a stratified cofinal ancestry}

Fix a matched infrared point \(x_0\in K_0\).  For \(N\ge1\), let
\(\mathcal A_N\) be the set of finite compatible ancestries
\begin{equation}
 (x_0,x_1,\ldots,x_N),\qquad
 \mathcal R_jx_{j+1}=x_j
 \quad(0\le j<N).
 \label{eq:v5v40-finite-ancestry}
\end{equation}

\begin{theorem}[Stratified inverse-limit existence]
\label{thm:v5v40-inverse-existence}
Assume:
\begin{enumerate}
 \item every \(K_j\) is nonempty compact Hausdorff;
 \item every \(\mathcal R_j\) is continuous;
 \item \(\mathcal A_N\ne\varnothing\) for every \(N\).
\end{enumerate}
Then there is an all-scale ancestry
\begin{equation}
 (x_j)_{j\ge0}\in\prod_{j\ge0}K_j,\qquad
 \mathcal R_jx_{j+1}=x_j
 \quad\text{for every }j.
 \label{eq:v5v40-infinite-ancestry}
\end{equation}
\end{theorem}

\begin{proof}
The product \(\prod_{j\ge0}K_j\) is compact.  For each \(N\), let
\(F_N\) be the closed subset on which the first \(N\) compatibility
equations hold and the zeroth coordinate is the prescribed \(x_0\).
Continuity makes \(F_N\) closed.  The finite-ancestry hypothesis, followed
by arbitrary choices in the remaining nonempty factors, makes \(F_N\)
nonempty.  The sets are nested.  Compactness gives
\(\bigcap_{N\ge1}F_N\ne\varnothing\), and every point of this intersection
satisfies \eqref{eq:v5v40-infinite-ancestry}.
\end{proof}

\begin{assessment}[Relation to the inherited inverse-limit theorem]
\label{ass:v5v40-nonduplication}
The compact finite-intersection argument is inherited in the earlier
programme.  The new content is its application after constructing the
Hausdorff stratified quotient
\eqref{eq:v5v40-stratified-K} and proving that shell maps descend through
the noninvertible collapsed boundary.
\end{assessment}

\subsection{Uniqueness under cumulative contraction}

Let \(d_j\) be compatible metrics on the stratified compact sets.  Suppose
two ancestries obey a uniform diameter bound \(d_j(x_j,y_j)\le D\) and
\begin{equation}
 d_j(\mathcal R_jx,\mathcal R_jy)
 \le\theta_jd_{j+1}(x,y)
 \label{eq:v5v40-theta}
\end{equation}
on the relevant fibres.

\begin{theorem}[Cumulative contraction criterion]
\label{thm:v5v40-cumulative-contraction}
If, for every fixed \(j\),
\begin{equation}
 \prod_{k=j}^{N}\theta_k\longrightarrow0
 \quad(N\to\infty),
 \label{eq:v5v40-product-zero}
\end{equation}
then the compatible ancestry over \(x_0\) is unique in the contracted
fibre coordinates.
\end{theorem}

\begin{proof}
For two compatible ancestries and \(N>j\), iterate
\eqref{eq:v5v40-theta}:
\[
 d_j(x_j,y_j)
 \le
 \left(\prod_{k=j}^{N-1}\theta_k\right)
 d_N(x_N,y_N)
 \le
 D\prod_{k=j}^{N-1}\theta_k.
\]
Let \(N\to\infty\) and use
\eqref{eq:v5v40-product-zero}.
\end{proof}

If a shell contraction \(q_j\) is expressed through a chart transition
with forward and inverse constants \(L_j^+,L_j^-\), one may take
\begin{equation}
 \theta_j\le q_jL_j^+L_j^-.
 \label{eq:v5v40-chart-theta}
\end{equation}
At the collapsed quotient no inverse factor is used; the reduced base and
conditional fibre instead use the triangular contraction of V39.

\begin{corollary}[Finite crossover cost]
\label{cor:v5v40-finite-crossover}
If all but finitely many steps have \(\theta_j\le q<1\) and every
exceptional transition has finite distortion, then
\eqref{eq:v5v40-product-zero} holds.
\end{corollary}

\begin{proof}
The finite product of exceptional constants is fixed, while the remaining
tail contributes a power of \(q\) tending to zero.
\end{proof}

\begin{corollary}[Infinite transition budget]
\label{cor:v5v40-infinite-budget}
If \(\theta_j>0\) and
\begin{equation}
 \sum_{j\ge j_0}-\log\theta_j=+\infty
 \label{eq:v5v40-log-budget}
\end{equation}
for every \(j_0\), then
\eqref{eq:v5v40-product-zero} holds.
\end{corollary}

\begin{proof}
Take logarithms of the product in
\eqref{eq:v5v40-product-zero}.
\end{proof}

\begin{firewall}[Existence does not imply uniqueness]
\label{fw:v5v40-existence-uniqueness}
Compact finite-depth ancestry gives at least one inverse-limit point but no
contraction.  Conversely, contraction estimates do not create a point if
some finite-depth ancestry set is empty.  Both hypotheses are required.
\end{firewall}

\subsection{Closed physical certificates}

Let \(K_j^{\mathrm{phys}}\subset K_j\) be the subset satisfying the
finite-cutoff Ward identities, anomaly normalization, reflection-positive
cylinder inequalities, and the declared quotient-observable conditions.

\begin{theorem}[Physical inverse-limit inheritance]
\label{thm:v5v40-physical-limit}
Assume:
\begin{enumerate}
 \item every \(K_j^{\mathrm{phys}}\) is nonempty and closed;
 \item \(\mathcal R_j(K_{j+1}^{\mathrm{phys}})
 \subseteq K_j^{\mathrm{phys}}\);
 \item physical finite-depth ancestries exist for every depth;
 \item cylinder Schwinger functions converge along the weak hybrid
 topology.
\end{enumerate}
Then an inverse-limit ancestry exists entirely in the physical sets.
Its limiting cylinder functional satisfies the closed Ward identities and
reflection-positive matrix inequalities.
\end{theorem}

\begin{proof}
Apply Theorem~\ref{thm:v5v40-inverse-existence} to the compact closed
subsets \(K_j^{\mathrm{phys}}\).  Ward identities pass to limits by
Lemma~\ref{lem:v5v27-Ward-closed}.  Every reflection-positive cylinder
condition is a finite matrix inequality.  Entrywise convergence preserves
positive semidefiniteness, as recorded in the inherited OS closure
theorem.
\end{proof}

\begin{firewall}[The physical sets are not known nonempty]
\label{fw:v5v40-physical-nonempty}
Theorem~\ref{thm:v5v40-physical-limit} is conditional.  In particular,
the conformal metric contour, hypercharge ultraviolet trajectory,
large-field domination, full RSTC/MITC, and same-coordinate transition
carrier have not been proved.  They can make the physical finite-depth
sets empty or destroy their uniform compactness.
\end{firewall}

\subsection{A stratified continuum certificate}

\begin{definition}[Stratified continuum ancestry certificate]
\label{def:v5v40-SCAC}
An SCAC consists of:
\begin{enumerate}
 \item V25 strong compact cores in all three strata;
 \item closed chart/quotient equivalence relations;
 \item continuous stratified-compatible shell maps;
 \item exact CTC bounds for coincident transition/Ward/payer blocks;
 \item finite-depth physical ancestries at every depth;
 \item cumulative contraction in the irrelevant and conditional-fibre
 coordinates;
 \item a MAC controlling the stratum itinerary of every collapse modulus;
 \item common anomaly, OS, tightness, clustering, and reconstruction
 certificates on the declared local observable algebra.
\end{enumerate}
\end{definition}

\begin{theorem}[What an SCAC would provide]
\label{thm:v5v40-SCAC}
An SCAC produces at least one all-scale stratified shell ancestry.  Its
contracted fibre coordinates are unique, its Ward and reflection-positive
cylinder conditions survive the cofinal limit, and its chart description
is independent of representatives in the collapsed quotient.
\end{theorem}

\begin{proof}
Use Theorems~\ref{thm:v5v40-compact-quotient},
\ref{thm:v5v40-descended-map},
\ref{thm:v5v40-inverse-existence},
\ref{thm:v5v40-cumulative-contraction}, and
\ref{thm:v5v40-physical-limit}.
\end{proof}

\subsection{V40 status}

\begin{theorem}[V40 conclusion]
\label{thm:v5v40-status}
The four-dimensional, active-band, and collapsed descriptions form a
valid compact inverse system when their exact chart and quotient relation
is closed and every shell map preserves it.  Finite-depth compatible
ancestors then give an all-scale stratified ancestry; cumulative
contraction gives uniqueness of the contracted coordinates without a
bi-Lipschitz inverse at collapse.  The theorem reduces the continuum step
to the SCAC but does not verify its physical nonemptiness or analytic
estimates for the coupled model.
\end{theorem}

\begin{proof}
This is the combined content of
Theorems~\ref{thm:v5v40-compact-quotient},
\ref{thm:v5v40-descended-map},
\ref{thm:v5v40-inverse-existence}, and
\ref{thm:v5v40-cumulative-contraction}.
\end{proof}

V41 will return to the first unresolved physical base direction:
hypercharge.  It will formulate the exact alternatives by which a
non-Gaussian ultraviolet fixed point or a finite compositeness boundary
could replace the Landau-pole trajectory, and prove the matching
conditions required for either mechanism to enter the SCAC without
claiming that such a fixed point exists.

\clearpage
\section{V41: Hypercharge exit, ultraviolet replacement, and exact matching}
\label{sec:v5v41}

\subsection{What is new relative to V3}

V3 proves that the Gaussian-dominance hypercharge face points outward and
states a degree and stable-manifold certificate for a hypothetical joint
fixed point.  V41 does not repeat those results.  It proves a quantitative
exit alternative for the positive inverse coupling and constructs the
inverse-system splice required when the Standard Model chart is replaced
at a finite scale by a genuine ultraviolet theory.  It also proves why a
bare compositeness endpoint is not an all-scale continuum.

Use ultraviolet time \(t=\log(\mu/\mu_0)\) and write
\begin{equation}
 z(t)=g_Y(t)^{-2}>0.
 \label{eq:v5v41-z}
\end{equation}
Let \(\mathcal W\) be the Gaussian-dominance region in the complete
marginal/tower state space on which the inherited estimate gives
\begin{equation}
 \dot z(t)\le-\kappa<0.
 \label{eq:v5v41-negative-z}
\end{equation}

\begin{theorem}[Quantitative hypercharge exit]
\label{thm:v5v41-exit}
Let a trajectory start at \(z(0)=z_0>0\) and remain in
\(\mathcal W\cap\{z>0\}\) on \(0\le t<T\).  Then
\begin{equation}
 z(t)\le z_0-\kappa t.
 \label{eq:v5v41-linear-exit}
\end{equation}
Consequently it cannot remain in that region beyond
\begin{equation}
 T_{\max}=z_0/\kappa.
 \label{eq:v5v41-Tmax}
\end{equation}
\end{theorem}

\begin{proof}
Integrate \eqref{eq:v5v41-negative-z} from \(0\) to \(t\):
\[
 z(t)-z_0=\int_0^t\dot z(s)\,ds\le-\kappa t.
\]
At \(t>z_0/\kappa\), the right side is negative, contradicting \(z(t)>0\).
\end{proof}

\begin{corollary}[All-scale alternative]
\label{cor:v5v41-alternative}
Every all-scale trajectory matching the observed positive hypercharge
coupling must do at least one of the following before \(T_{\max}\):
\begin{enumerate}
 \item leave the Gaussian-dominance region because the exact joint beta
 field changes;
 \item enter the stable set of a non-Gaussian ultraviolet invariant set;
 \item pass through a matched interface into a different ultraviolet
 theory or chart;
 \item cease to define a positive hypercharge kinetic state.
\end{enumerate}
\end{corollary}

\begin{proof}
If none of the first three events occurs, the trajectory stays in the
hypotheses of Theorem~\ref{thm:v5v41-exit}.  It then reaches the boundary
\(z=0\) by \(T_{\max}\), which is the fourth alternative.
\end{proof}

\begin{firewall}[A coordinate change is not a new mechanism]
\label{fw:v5v41-coordinate-change}
A regular invertible reparametrization of couplings transports the same
trajectory and cannot remove its finite-time exit from the physical
positive kinetic domain.  An ultraviolet chart counts as alternative 3
only when it includes additional states or a nontrivial completion and has
an exact matching map on observables.
\end{firewall}

\subsection{A one-dimensional non-Gaussian entrance test}

The actual coupled beta field is multidimensional, as in V3.  The following
one-dimensional theorem isolates the minimum hypercharge sign pattern
along an invariant curve.

\begin{theorem}[Monotone entrance to a positive fixed point]
\label{thm:v5v41-1d-fixed}
Let \(\beta\in C^1([g_0,g_*])\) satisfy
\begin{equation}
 0<g_0<g_*,\qquad
 \beta(g)>0\ (g_0\le g<g_*),\qquad
 \beta(g_*)=0,\qquad
 \beta'(g_*)=-\lambda<0.
 \label{eq:v5v41-1d-hyp}
\end{equation}
Then the solution of
\(\dot g=\beta(g)\), \(g(0)=g_0\), exists for every \(t\ge0\), is
increasing, and converges to \(g_*\).  For every
\(0<\lambda'<\lambda\), it obeys
\begin{equation}
 0<g_*-g(t)\le C_{\lambda'}e^{-\lambda't}
 \label{eq:v5v41-1d-decay}
\end{equation}
once it enters a sufficiently small neighborhood of \(g_*\).
\end{theorem}

\begin{proof}
Local uniqueness and the sign of \(\beta\) make the solution increase
while it lies below \(g_*\).  It cannot cross \(g_*\), because the constant
solution at \(g_*\) and uniqueness would then be violated at the first
crossing.  Thus it is bounded and has a limit \(g_\infty\le g_*\).  If
\(g_\infty<g_*\), continuity gives a positive lower bound for \(\beta\) on
a neighborhood of \(g_\infty\), contradicting convergence.  Hence
\(g_\infty=g_*\), and boundedness prevents finite-time escape.

For \(h=g_*-g\), Taylor's theorem gives
\[
 \dot h=-\beta(g_*-h)=-\lambda h+o(h).
\]
For sufficiently small \(h>0\),
\(\dot h\le-\lambda'h\).  Gronwall's inequality gives
\eqref{eq:v5v41-1d-decay}.
\end{proof}

\begin{firewall}[Invariant-curve hypothesis]
\label{fw:v5v41-invariant-curve}
Theorem~\ref{thm:v5v41-1d-fixed} applies only after the other Standard
Model, gravitational, and operator-tower couplings have been solved as an
invariant curve.  A zero of a truncated scalar beta function does not prove
that such a curve exists or intersects the observed infrared surface.
\end{firewall}

\subsection{A genuine ultraviolet replacement system}

Let \(J\) be a matching shell.  For \(j<J\), let
\[
 K_j^{\mathrm L},\qquad
 \mathcal R_j^{\mathrm L}:K_{j+1}^{\mathrm L}\to K_j^{\mathrm L}
\]
be the low-energy SM--Einstein cutoff system.  For \(j\ge J\), let
\[
 K_j^{\mathrm U},\qquad
 \mathcal R_j^{\mathrm U}:K_{j+1}^{\mathrm U}\to K_j^{\mathrm U}
\]
be a proposed ultraviolet completion.  Let
\begin{equation}
 \mathcal M_J:K_J^{\mathrm U}\to K_J^{\mathrm L}
 \label{eq:v5v41-matching-map}
\end{equation}
be a continuous matching map.

\begin{definition}[Ultraviolet interface certificate]
\label{def:v5v41-UIC}
A UIC consists of:
\begin{enumerate}
 \item nonempty compact physical sets \(K_j^{\mathrm L},K_j^{\mathrm U}\);
 \item continuous physical shell maps on both sides;
 \item the matching map \eqref{eq:v5v41-matching-map};
 \item equality of the normalized Schwinger functionals on the common
 matched local observable algebra;
 \item compatibility of determinant lines, zero-mode selection rules,
 gauge/diffeomorphism Ward identities, and local counterterm owners;
 \item preservation of every reflected cylinder matrix used in the OS
 certificate;
 \item finite-depth ultraviolet ancestors matching the prescribed
 low-energy point at every depth.
\end{enumerate}
\end{definition}

Define the spliced bonding map at the interface by
\begin{equation}
 \mathcal R_{J-1}^{\mathrm{splice}}
 =
 \mathcal R_{J-1}^{\mathrm L}\mathcal M_J.
 \label{eq:v5v41-splice-map}
\end{equation}
Below and above it use the corresponding low and ultraviolet maps.

\begin{theorem}[Ultraviolet splice theorem]
\label{thm:v5v41-splice}
A UIC produces an all-scale inverse-limit ancestry matching the prescribed
low-energy state.  If the low system, interface, and ultraviolet system
have cumulative contraction product tending to zero in the contracted
coordinates, those coordinates are unique.
\end{theorem}

\begin{proof}
The spliced state spaces are compact and nonempty, and all bonding maps are
continuous; continuity of the interface map makes
\eqref{eq:v5v41-splice-map} continuous.  UIC item 7 supplies compatible
finite-depth ancestors for this single projective system.  Apply
Theorem~\ref{thm:v5v40-inverse-existence}.  For uniqueness, the interface
contributes one finite Lipschitz factor to the product in
Theorem~\ref{thm:v5v40-cumulative-contraction}; the assumed cumulative
product still tends to zero.
\end{proof}

\begin{corollary}[Closed physical properties survive the splice]
\label{cor:v5v41-physical-splice}
Under a UIC, the cofinal cylinder functional inherits the common closed
Ward and OS conditions, provided the uniform convergence and tightness
hypotheses of V40 hold on both sides.
\end{corollary}

\begin{proof}
UIC items 4--6 make the physical subsets invariant across the interface.
Apply Theorem~\ref{thm:v5v40-physical-limit} to the spliced system.
\end{proof}

\subsection{Why a bare compositeness scale is not a continuum}

\begin{theorem}[Finite terminal boundary no-go]
\label{thm:v5v41-terminal-no-go}
Suppose the low-energy shell system is defined only through a finite
terminal shell \(J\), with no ultraviolet state spaces or matching map
beyond it.  Then it cannot satisfy the all-depth ancestry hypothesis of
Theorem~\ref{thm:v5v40-inverse-existence} and does not define an
all-scale continuum by that construction.
\end{theorem}

\begin{proof}
For every \(N>J\), the finite ancestry set
\(\mathcal A_N\) requires a state in \(K_N\) and bonding maps from shell
\(N\) down to shell zero.  These objects do not exist in the terminal
system.  Hence the hypothesis
\(\mathcal A_N\ne\varnothing\) for every \(N\) fails.
\end{proof}

\begin{corollary}[Meaning of a compositeness condition]
\label{cor:v5v41-compositeness}
A condition such as \(z_Y(J)=0\) at a finite scale can be part of a
continuum construction only if it is an interface boundary condition for
a genuine ultraviolet system satisfying a UIC.  By itself it is a finite
cutoff model.
\end{corollary}

\begin{proof}
Apply Theorem~\ref{thm:v5v41-terminal-no-go}.  A UIC supplies the missing
states and maps by Theorem~\ref{thm:v5v41-splice}.
\end{proof}

\subsection{Gauge-kinetic matching at an embedding}

One concrete ultraviolet replacement embeds hypercharge into a
nonabelian gauge algebra.  Let
\(\iota:\mathfrak u(1)_Y\to\mathfrak g_{\mathrm U}\) be an embedding, and
normalize invariant quadratic forms so that
\begin{equation}
 \langle\iota(Y),\iota(Y)\rangle_{\mathrm U}
 =
 k_Y\langle Y,Y\rangle_Y,\qquad k_Y>0.
 \label{eq:v5v41-embedding-index}
\end{equation}

\begin{proposition}[Quadratic kinetic matching]
\label{prop:v5v41-kinetic-matching}
If the ultraviolet gauge action has inverse coupling \(z_{\mathrm U}\),
its restriction to the hypercharge direction contributes inverse coupling
\(k_Yz_{\mathrm U}\).  After integrating threshold fields, the most general
local quadratic matching has the form
\begin{equation}
 z_Y(J^-)=k_Yz_{\mathrm U}(J^+)+\Delta z_Y(J),
 \label{eq:v5v41-kinetic-matching}
\end{equation}
where \(\Delta z_Y\) is the renormalized threshold owner.
\end{proposition}

\begin{proof}
Restrict
\[
 \frac{z_{\mathrm U}}4
 \int\langle F_{\mathrm U},F_{\mathrm U}\rangle_{\mathrm U}
\]
to \(F_{\mathrm U}=\iota(F_Y)\).  Equation
\eqref{eq:v5v41-embedding-index} gives coefficient \(k_Yz_{\mathrm U}\).
Every additional gauge-invariant local quadratic contribution generated
at the interface has the same \(F_Y^2\) form and is, by definition, the
threshold coefficient \(\Delta z_Y\).
\end{proof}

\begin{firewall}[Quadratic matching is not a UIC]
\label{fw:v5v41-quadratic-not-UIC}
Equation \eqref{eq:v5v41-kinetic-matching} matches one two-point local
coefficient.  A UIC also requires all common Schwinger functions,
anomaly/determinant-line data, massive threshold tensors, Ward identities,
OS matrices, and compact shell estimates.  Gauge-coupling unification
alone is insufficient.
\end{firewall}

\subsection{A fixed point as a special ultraviolet system}

Let \(K_*^{\mathrm U}\) be a compact neighborhood of a non-Gaussian
ultraviolet fixed state \(x_*\), including the complete gravitational
operator tower.  Suppose its shell map has a stable graph and a physical
matching point, as in the V2--V3 certificates.

\begin{definition}[Hypercharge fixed-point interface certificate]
\label{def:v5v41-HFIC}
An HFIC is a UIC in which:
\begin{enumerate}
 \item the ultraviolet maps approach a fixed shell map with fixed point
 \(x_*\);
 \item the positive hypercharge kinetic coordinate of \(x_*\) is finite or
 is represented by a regular equivalent physical coordinate;
 \item the matched infrared point lies on the ultraviolet stable graph;
 \item the exact beta/shell remainder and the infinite local tower obey
 the V25--V26 contraction and compactness estimates.
\end{enumerate}
\end{definition}

\begin{theorem}[Conditional fixed-point completion]
\label{thm:v5v41-fixed-completion}
An HFIC replaces the Gaussian hypercharge exit by an all-scale matched
trajectory and supplies the base ancestry needed by the V2 irrelevant
graph and the V40 stratified inverse system.
\end{theorem}

\begin{proof}
The HFIC stable graph supplies ultraviolet ancestors at every depth and
its interface matching supplies item 7 of the UIC.  Apply
Theorem~\ref{thm:v5v41-splice}.  The resulting marginal/tower ancestry is
the prescribed base for Theorem~\ref{thm:v5v2-certificate}, while its
stratified compact realization is covered by V40.
\end{proof}

\begin{firewall}[No fixed point has been constructed]
\label{fw:v5v41-no-fixed-point}
V41 proves what a non-Gaussian fixed point or ultraviolet replacement must
do.  It does not establish an exact SM--gravity beta zero, an ultraviolet
gauge embedding, a UIC, or an HFIC.  The hypercharge gate remains open.
\end{firewall}

\subsection{Interface distortion and the Schur budget}

Let \(L_{\mathcal M}\) be the matching-map Lipschitz constant in the
hybrid physical norm.  If a contracted ultraviolet step has constant
\(q_{\mathrm U}\), transporting it through an injective active interface
with inverse constant \(L_{\mathcal M^{-1}}\) costs
\begin{equation}
 q_{\mathrm{match}}
 \le L_{\mathcal M}L_{\mathcal M^{-1}}q_{\mathrm U}.
 \label{eq:v5v41-match-q}
\end{equation}
At a quotient interface, use the base--fibre theorem V39 instead.

\begin{proposition}[Finite interface cost]
\label{prop:v5v41-finite-interface}
If only finitely many matching shells occur and all interface constants
are finite, their product does not obstruct cumulative contraction of a
uniformly contracting ultraviolet tail.
\end{proposition}

\begin{proof}
This is Corollary~\ref{cor:v5v40-finite-crossover} with each interface
distortion inserted into \(\theta_j\).
\end{proof}

The V40 audit adds that \(L_{\mathcal M}\) must be estimated after complete
equivariant multiplicity pinching.  If matching, a Ward carrier, and a
threshold payer coincide, the required input is the assembled CTC
\eqref{eq:v5v40-CTC}, not the product of three coefficientwise bounds.

\subsection{V41 status}

\begin{theorem}[V41 conclusion]
\label{thm:v5v41-status}
The positive hypercharge trajectory exits the inherited
Gaussian-dominance region within the explicit time
\(z_0/\kappa\).  A non-Gaussian fixed point can replace that exit only
through a matched stable trajectory, while a finite compositeness boundary
defines a continuum only when it is the interface to a genuine ultraviolet
inverse system.  A UIC is sufficient to splice such a system into the V40
ancestry.  No UIC or HFIC is presently proved.
\end{theorem}

\begin{proof}
Combine Theorem~\ref{thm:v5v41-exit},
Theorem~\ref{thm:v5v41-splice}, and
Theorem~\ref{thm:v5v41-terminal-no-go}.
\end{proof}

V42 will examine whether a finite truncation of the gravitational
operator tower can certify an HFIC.  It will prove a projection-error
theorem for beta zeros and show exactly which tail inverse and residual
bounds are required before a fixed point of a truncated flow can be lifted
to the full tri-graded tower.

\clearpage
\section{V42: Lifting a truncated beta zero to the full operator tower}
\label{sec:v5v42}

\subsection{The projection problem}

An HFIC cannot be certified by finding a zero of finitely many beta
functions while discarding the generated gravitational tower.  Let \(X\)
be the full hybrid tri-graded Banach space of marginal and local operator
coordinates, and let
\begin{equation}
 F:X\to X
 \label{eq:v5v42-beta-map}
\end{equation}
be the exact beta vector field or fixed-shell defect.  An exact ultraviolet
fixed state satisfies \(F(x)=0\).

Let \(P_N:X\to X_N\) be a bounded finite-rank projection onto all declared
head and local grades through \(N\), and put \(Q_N=I-P_N\).  A numerical or
symbolic truncation supplies \(x_N\in X_N\) satisfying
\begin{equation}
 P_NF(x_N)=0.
 \label{eq:v5v42-truncated-zero}
\end{equation}
Its full residual is
\begin{equation}
 r_N=F(x_N)=Q_NF(x_N).
 \label{eq:v5v42-residual}
\end{equation}

\begin{firewall}[A projected zero is not an approximate zero without a
tail norm]
\label{fw:v5v42-projected-zero}
Equation \eqref{eq:v5v42-truncated-zero} gives no bound on
\(\|r_N\|_X\).  The omitted beta functions can be order one or undefined
in the chosen topology.  A truncation becomes evidence for a full zero
only after \eqref{eq:v5v42-residual} is estimated.
\end{firewall}

\subsection{A Newton--Kantorovich lifting theorem}

\begin{theorem}[Quantitative full-space zero lift]
\label{thm:v5v42-Newton-lift}
Let \(F\) be continuously differentiable on the closed ball
\(\overline B(x_N,R)\subset X\).  Set
\[
 L_N=DF(x_N).
\]
Assume:
\begin{align}
 L_N^{-1}\text{ exists},\qquad
 \|L_N^{-1}\|&\le M,
 \label{eq:v5v42-inverse-M}\\
 \|DF(x)-DF(y)\|&\le K\|x-y\|
 \quad(x,y\in\overline B(x_N,R)),
 \label{eq:v5v42-DF-Lipschitz}\\
 \eta:=\|F(x_N)\|,\qquad
 2M\eta&\le R,
 \label{eq:v5v42-radius}\\
 2M^2K\eta&<1.
 \label{eq:v5v42-Newton-smallness}
\end{align}
Then \(F\) has a unique zero in
\(\overline B(x_N,2M\eta)\), and
\begin{equation}
 \|x_*-x_N\|\le2M\eta.
 \label{eq:v5v42-zero-error}
\end{equation}
\end{theorem}

\begin{proof}
Write \(x=x_N+h\) and
\[
 N(h)=F(x_N+h)-F(x_N)-L_Nh.
\]
The derivative Lipschitz bound and the fundamental theorem of calculus give
\begin{align}
 \|N(h)\|&\le\frac K2\|h\|^2,
 \label{eq:v5v42-N-quadratic}\\
 \|N(h)-N(k)\|&\le
 K\max\{\|h\|,\|k\|\}\|h-k\|.
 \label{eq:v5v42-N-Lipschitz}
\end{align}
Define
\[
 \mathcal T(h)
 =
 -L_N^{-1}\bigl(F(x_N)+N(h)\bigr)
\]
on the ball of radius \(R_*=2M\eta\), which lies inside the declared ball
by \eqref{eq:v5v42-radius}.  For \(\|h\|\le R_*\),
\[
 \|\mathcal T(h)\|
 \le M\eta+\frac{MK}{2}R_*^2
 =
 \frac{R_*}{2}+M^2K\eta R_*<R_*.
\]
For \(h,k\) in the ball,
\[
 \|\mathcal T(h)-\mathcal T(k)\|
 \le MKR_*\|h-k\|
 =2M^2K\eta\|h-k\|.
\]
This is a strict contraction by
\eqref{eq:v5v42-Newton-smallness}.  Its unique fixed point gives
\(F(x_N+h)=0\) and obeys \eqref{eq:v5v42-zero-error}.  Every zero in the
ball is a fixed point of the same map, proving local uniqueness.
\end{proof}

\begin{assessment}[Three independent numerical obligations]
\label{ass:v5v42-three-obligations}
The lift needs a full residual \(\eta\), a full derivative inverse bound
\(M\), and a nonlinear derivative bound \(K\).  Agreement of successive
truncated coordinates estimates none of these by itself.
\end{assessment}

\subsection{A block certificate for the full derivative inverse}

Relative to \(X=X_N\oplus Q_NX\), write
\begin{equation}
 L_N=
 \begin{pmatrix}
 A&B\\
 C&D
 \end{pmatrix}.
 \label{eq:v5v42-block-L}
\end{equation}
The finite truncation computes \(A\), but the other three blocks contain
the tower feedback.

\begin{theorem}[Schur inverse certificate]
\label{thm:v5v42-Schur-inverse}
Assume \(A\) and \(D\) are invertible and put
\begin{equation}
 a=\|A^{-1}\|,\quad b=\|B\|,\quad
 c=\|C\|,\quad d=\|D^{-1}\|,\quad
 \vartheta=abdc.
 \label{eq:v5v42-block-constants}
\end{equation}
If \(\vartheta<1\), then \(L_N\) is invertible.  With the block max norm,
\begin{equation}
 \|L_N^{-1}\|
 \le
 M_{\mathrm{Schur}}
 :=
 \max\left\{
 s(1+bd),\
 d+dc\,s(1+bd)
 \right\},
 \qquad
 s=\frac a{1-\vartheta}.
 \label{eq:v5v42-Schur-M}
\end{equation}
\end{theorem}

\begin{proof}
The head Schur complement is
\[
 S=A-BD^{-1}C
 =A\bigl(I-A^{-1}BD^{-1}C\bigr).
\]
The norm of the second term is at most \(\vartheta<1\), so its inverse is
the Neumann series and
\(\|S^{-1}\|\le a/(1-\vartheta)=s\).  The block inverse is
\[
 L_N^{-1}
 =
 \begin{pmatrix}
 S^{-1}&-S^{-1}BD^{-1}\\
 -D^{-1}CS^{-1}&
 D^{-1}+D^{-1}CS^{-1}BD^{-1}
 \end{pmatrix}.
\]
The sum of the norms in its first row is at most \(s(1+bd)\); the second
row sum is at most \(dc\,s+d+dc\,sbd=d+dc\,s(1+bd)\).  The induced max
norm is bounded by their maximum.
\end{proof}

\begin{corollary}[Tail inverse from canonical scaling]
\label{cor:v5v42-tail-inverse}
If
\begin{equation}
 D=D_0(I-E),\qquad
 \|D_0^{-1}\|\le d_0,\qquad
 \|E\|\le e_0<1,
 \label{eq:v5v42-tail-D}
\end{equation}
then
\begin{equation}
 \|D^{-1}\|\le\frac{d_0}{1-e_0}.
 \label{eq:v5v42-tail-d}
\end{equation}
\end{corollary}

\begin{proof}
The inverse is
\((I-E)^{-1}D_0^{-1}\).  Sum the Neumann series for \(I-E\).
\end{proof}

\begin{firewall}[The finite Hessian is insufficient]
\label{fw:v5v42-finite-Hessian}
Invertibility of the computed finite matrix \(A\) does not imply
invertibility of \(L_N\).  The tail can have a kernel, and the feedback
\(BD^{-1}C\) can make the Schur complement singular.  Bounds for
\(b,c,d\) and \(\vartheta\) are mandatory.
\end{firewall}

\subsection{Factorial tail control of the residual}

Use the strong and weak curvature radii
\(\rho_+>\rho_->0\) of V25, with the same particle and spatial parameters
for simplicity.  Suppose the residual has no components through grade
\(N\):
\[
 P_NF(x_N)=0.
\]

\begin{lemma}[Tri-graded residual tail]
\label{lem:v5v42-residual-tail}
If
\begin{equation}
 \|F(x_N)\|_{\rho_+,\eta;s,b}^{\alpha,\beta}
 \le R_F,
 \label{eq:v5v42-strong-residual}
\end{equation}
then
\begin{equation}
 \eta_N:=
 \|F(x_N)\|_{\rho_-,\eta;s,b}^{\alpha,\beta}
 \le
 R_F
 \left(\frac{\rho_-}{\rho_+}\right)^{N+1}.
 \label{eq:v5v42-residual-decay}
\end{equation}
\end{lemma}

\begin{proof}
Every nonzero component has curvature grade \(p\ge N+1\).  In each such
block, the quotient of weak and strong weights is
\((\rho_-/\rho_+)^p\), at most its value at \(N+1\).  Factor out that
number and sum the strong norm.
\end{proof}

\begin{theorem}[Convergent truncation lift]
\label{thm:v5v42-convergent-lift}
Suppose a sequence of truncated zeros \(x_N\) satisfies:
\begin{enumerate}
 \item the strong residual bound
 \eqref{eq:v5v42-strong-residual} with common \(R_F\);
 \item the Schur inverse bound \(\|DF(x_N)^{-1}\|\le M\) uniformly;
 \item the derivative Lipschitz bound
 \eqref{eq:v5v42-DF-Lipschitz} with common \(K\) on balls of a fixed
 radius;
 \item the radius and Newton conditions hold for all sufficiently large
 \(N\).
\end{enumerate}
Then every sufficiently large truncation lifts to an exact full zero
\(x_N^*\), with
\begin{equation}
 \|x_N^*-x_N\|_-
 \le
 2MR_F
 \left(\frac{\rho_-}{\rho_+}\right)^{N+1}.
 \label{eq:v5v42-lift-rate}
\end{equation}
\end{theorem}

\begin{proof}
Lemma~\ref{lem:v5v42-residual-tail} bounds the weak residual by the right
side of \eqref{eq:v5v42-residual-decay}.  Insert that \(\eta_N\) into
Theorem~\ref{thm:v5v42-Newton-lift}.  Its conclusion is
\eqref{eq:v5v42-lift-rate}.
\end{proof}

\begin{firewall}[Strong boundedness must be proved]
\label{fw:v5v42-strong-residual}
The decay in \eqref{eq:v5v42-residual-decay} comes from a strict radius
gap and a uniform strong norm.  Merely setting all computed high-grade
coordinates to zero does not establish
\eqref{eq:v5v42-strong-residual}; the exact beta field can generate a
factorially large tail.
\end{firewall}

\subsection{Compatibility with the Ward surface}

Let \(\mathfrak S\subset X\) be a closed linearized Ward slice.  Suppose
\[
 F(\mathfrak S)\subset\mathfrak S
\]
in the chosen local chart and \(x_N\in\mathfrak S\).

\begin{theorem}[Ward-compatible zero lift]
\label{thm:v5v42-Ward-lift}
If the hypotheses of Theorem~\ref{thm:v5v42-Newton-lift} hold after
restricting \(F\) and \(L_N\) to \(\mathfrak S\), then the lifted zero
belongs to \(\mathfrak S\).
\end{theorem}

\begin{proof}
The Newton map in the proof of
Theorem~\ref{thm:v5v42-Newton-lift} maps the tangent space
\(\mathfrak S-x_N\) to itself because the residual, nonlinear remainder,
and restricted inverse all lie there.  Start its iteration at zero.  The
limit \(h_*\) lies in the closed subspace, so
\(x_N+h_*\in\mathfrak S\).
\end{proof}

\begin{firewall}[An anomalous slice is not invariant]
\label{fw:v5v42-anomalous-slice}
Theorem~\ref{thm:v5v42-Ward-lift} requires the exact assembled Ward
identity.  If the V28 harmonic anomaly is nonzero or the regulator defect
is unrepaired, \(F(\mathfrak S)\subset\mathfrak S\) fails and Newton
iteration can leave the physical slice.
\end{firewall}

\subsection{Hypercharge positivity after lifting}

Let \(\ell_Y\in X^*\) extract the inverse hypercharge coupling:
\[
 z_Y(x)=\ell_Y(x).
\]

\begin{proposition}[Positive kinetic margin]
\label{prop:v5v42-positive-z}
If a truncated zero has
\begin{equation}
 z_Y(x_N)\ge z_*>0
 \label{eq:v5v42-z-margin}
\end{equation}
and
\begin{equation}
 2\|\ell_Y\|M\eta<z_*,
 \label{eq:v5v42-z-smallness}
\end{equation}
then its lifted zero satisfies \(z_Y(x_N^*)>0\).
\end{proposition}

\begin{proof}
By \eqref{eq:v5v42-zero-error},
\[
 z_Y(x_N^*)
 \ge z_Y(x_N)-\|\ell_Y\|\|x_N^*-x_N\|
 \ge z_*-2\|\ell_Y\|M\eta>0.
\]
\end{proof}

\begin{firewall}[A zero near \(z_Y=0\)]
\label{fw:v5v42-z-boundary}
If the truncated fixed point approaches the boundary \(z_Y=0\) at the
same rate as its uncontrolled tail error, the lift theorem cannot certify
a positive hypercharge kinetic term.  Positivity needs a margin larger
than the full correction.
\end{firewall}

\subsection{Stability data after the lift}

An HFIC also needs an ultraviolet stable graph.  Let
\(A_N=DF(x_N)\) and \(A_*=DF(x_N^*)\).  The derivative Lipschitz estimate
gives
\begin{equation}
 \|A_*-A_N\|
 \le K\|x_N^*-x_N\|
 \le2KM\eta.
 \label{eq:v5v42-derivative-error}
\end{equation}

\begin{definition}[Dichotomy robustness certificate]
\label{def:v5v42-DRC}
A DRC consists of a stable/unstable spectral splitting for \(A_N\), a
resolvent contour \(\Gamma\) separating the two spectral sets, and a
constant \(R_\Gamma\) such that
\begin{equation}
 \sup_{\lambda\in\Gamma}
 \|(\lambda-A_N)^{-1}\|\le R_\Gamma,\qquad
 2KM\eta\,R_\Gamma<1.
 \label{eq:v5v42-DRC}
\end{equation}
\end{definition}

\begin{proposition}[Persistence of the spectral splitting]
\label{prop:v5v42-splitting}
Under a DRC, \(\Gamma\) lies in the resolvent set of \(A_*\), and the Riesz
projections
\begin{equation}
 P_N^\Gamma=\frac1{2\pi i}\int_\Gamma
 (\lambda-A_N)^{-1}\,d\lambda,\qquad
 P_*^\Gamma=\frac1{2\pi i}\int_\Gamma
 (\lambda-A_*)^{-1}\,d\lambda
 \label{eq:v5v42-Riesz}
\end{equation}
are well defined.  Their difference is bounded by
\begin{equation}
 \|P_*^\Gamma-P_N^\Gamma\|
 \le
 \frac{|\Gamma|}{2\pi}
 \frac{R_\Gamma^2\,2KM\eta}
 {1-R_\Gamma\,2KM\eta}.
 \label{eq:v5v42-projection-error}
\end{equation}
\end{proposition}

\begin{proof}
For \(\lambda\in\Gamma\),
\[
 \lambda-A_*
 =
 \bigl[I-(A_*-A_N)(\lambda-A_N)^{-1}\bigr]
 (\lambda-A_N).
\]
The bracket is invertible by the Neumann series and the DRC.  Thus
\[
 \|(\lambda-A_*)^{-1}\|
 \le\frac{R_\Gamma}{1-R_\Gamma\|A_*-A_N\|}.
\]
Use the resolvent identity
\[
 (\lambda-A_*)^{-1}-(\lambda-A_N)^{-1}
 =
 (\lambda-A_*)^{-1}(A_*-A_N)(\lambda-A_N)^{-1},
\]
insert \eqref{eq:v5v42-derivative-error}, and integrate around
\(\Gamma\).
\end{proof}

\begin{assessment}[What remains for a stable manifold]
\label{ass:v5v42-stable-manifold}
Proposition~\ref{prop:v5v42-splitting} preserves the spectral separation.
An actual Lyapunov--Perron stable graph still needs semigroup dichotomy
bounds and a nonlinear Lipschitz constant satisfying the V3 contraction
inequality.  Spectral plots from a finite matrix do not provide those
infinite-dimensional estimates.
\end{assessment}

\subsection{A full fixed-point lift certificate}

\begin{definition}[Full tower fixed-point lift certificate]
\label{def:v5v42-FTFLC}
An FTFLC at truncation \(N\) consists of:
\begin{enumerate}
 \item the exact projected equation \eqref{eq:v5v42-truncated-zero};
 \item a strong tri-graded bound on the full residual;
 \item finite-head and infinite-tail inverse bounds satisfying
 Theorem~\ref{thm:v5v42-Schur-inverse};
 \item a full derivative Lipschitz bound and Newton smallness;
 \item an assembled invariant Ward slice and vanishing anomaly component;
 \item a positive hypercharge margin;
 \item a DRC and the V3 semigroup/nonlinear stable-graph estimates;
 \item compatibility of the lifted zero with the V41 interface and
 infrared matching surface.
\end{enumerate}
\end{definition}

\begin{theorem}[FTFLC implication]
\label{thm:v5v42-FTFLC}
An FTFLC produces a genuine positive-hypercharge fixed state of the full
tri-graded tower near \(x_N\), preserves its declared Ward surface and
hyperbolic splitting, and supplies the local fixed-point data required by
an HFIC.
\end{theorem}

\begin{proof}
Use Theorem~\ref{thm:v5v42-Schur-inverse} to obtain \(M\), then
Theorem~\ref{thm:v5v42-Newton-lift} for the exact zero.
Theorem~\ref{thm:v5v42-Ward-lift} and
Proposition~\ref{prop:v5v42-positive-z} give the physical slice and
positivity.  Proposition~\ref{prop:v5v42-splitting}, followed by the
declared V3 stable-graph estimates, gives the ultraviolet local stable
set.  The last item supplies the interface matching.
\end{proof}

\begin{firewall}[Current fixed-point evidence]
\label{fw:v5v42-current-evidence}
No FTFLC has been supplied for a computed Standard Model--Einstein
truncation.  V42 therefore does not promote any finite beta-function zero
to a continuum fixed point.
\end{firewall}

\subsection{V42 status}

\begin{theorem}[V42 conclusion]
\label{thm:v5v42-status}
A finite truncation zero lifts to the full operator tower only when its
omitted beta residual is small in a strict-radius norm, the full derivative
has a certified Schur inverse including tail feedback, and the
Newton--Kantorovich product \(2M^2K\eta\) is below one.  Ward invariance,
positive hypercharge margin, and stable-spectrum robustness require
separate quantitative inequalities.  These conditions define an FTFLC;
none is presently verified for the physical coupled flow.
\end{theorem}

\begin{proof}
This is the combined content of
Theorems~\ref{thm:v5v42-Newton-lift},
\ref{thm:v5v42-Schur-inverse},
\ref{thm:v5v42-convergent-lift}, and
\ref{thm:v5v42-FTFLC}.
\end{proof}

V43 will go upstream to the tail inverse \(D^{-1}\).  It will use canonical
operator dimensions to construct a diagonal tail parametrix and determine
whether gravitational anomalous mixing is small enough, in the factorial
radius norm, for the Neumann condition
\eqref{eq:v5v42-tail-D}.

\clearpage
\section{V43: Canonical tail inversion and the grade-lowering obstruction}
\label{sec:v5v43}

\subsection{The tail derivative}

V42 reduces a full fixed-point lift to a tail inverse estimate.  Write the
tail block of the exact derivative as
\begin{equation}
 D=D_0+K,
 \label{eq:v5v43-D}
\end{equation}
where \(D_0\) is the canonical-dimension diagonal and \(K\) is anomalous
mixing at the candidate fixed point.  This note derives the exact weighted
matrix criterion for inversion in the V25 factorial curvature norm and
shows why grade-lowering mixing is the dangerous direction.

Suppress particle and spatial indices temporarily; they remain inside the
block norm \(X_p\).  On the tail \(p>N\), define
\begin{equation}
 \|x\|_{\rho,\alpha}
 =
 \sum_{p>N}w_p\|x_p\|_{X_p},
 \qquad
 w_p=\frac{\rho^p}{(p!)^\alpha}.
 \label{eq:v5v43-tail-norm}
\end{equation}
Assume
\begin{equation}
 (D_0x)_p=\lambda_px_p,\qquad
 |\lambda_p|\ge c_0+c_1p,\qquad c_1>0.
 \label{eq:v5v43-canonical-gap}
\end{equation}

\begin{lemma}[Canonical inverse]
\label{lem:v5v43-D0-inverse}
The diagonal tail operator is invertible and
\begin{equation}
 \|D_0^{-1}\|_{\rho,\alpha}
 \le
 d_{0,N}:=
 \frac1{c_0+c_1(N+1)}
 \label{eq:v5v43-D0-bound}
\end{equation}
when the denominator is positive.
\end{lemma}

\begin{proof}
Blockwise,
\(\|(D_0^{-1}y)_p\|\le|\lambda_p|^{-1}\|y_p\|\).
For \(p>N\), the reciprocal is at most the right side of
\eqref{eq:v5v43-D0-bound}.  Multiply by \(w_p\) and sum.
\end{proof}

\subsection{The exact weighted column criterion}

Write \(K_{pq}:X_q\to X_p\) for the mixing blocks.  Define
\begin{equation}
 \epsilon_N(\rho,\alpha)
 =
 \sup_{q>N}
 \sum_{p>N}
 \frac{w_p}{w_q}
 \frac{\|K_{pq}\|}{|\lambda_p|}.
 \label{eq:v5v43-epsilon}
\end{equation}

\begin{theorem}[Weighted tail inverse certificate]
\label{thm:v5v43-WTIC}
If
\begin{equation}
 \epsilon_N(\rho,\alpha)<1,
 \label{eq:v5v43-epsilon-small}
\end{equation}
then \(D=D_0+K\) is invertible on the weighted tail and
\begin{equation}
 \|D^{-1}\|_{\rho,\alpha}
 \le
 \frac{d_{0,N}}{1-\epsilon_N(\rho,\alpha)}.
 \label{eq:v5v43-D-inverse}
\end{equation}
\end{theorem}

\begin{proof}
Let \(E=D_0^{-1}K\).  For \(x\) of finite support,
\begin{align*}
 \|Ex\|_{\rho,\alpha}
 &\le
 \sum_{p>N}w_p
 \sum_{q>N}\frac{\|K_{pq}\|}{|\lambda_p|}\|x_q\|\\
 &=
 \sum_{q>N}w_q\|x_q\|
 \sum_{p>N}\frac{w_p}{w_q}
 \frac{\|K_{pq}\|}{|\lambda_p|}\\
 &\le\epsilon_N(\rho,\alpha)\|x\|_{\rho,\alpha}.
\end{align*}
Density extends \(E\) boundedly.  Since
\[
 D=D_0(I+E),
\]
the Neumann series gives
\[
 D^{-1}=(I+E)^{-1}D_0^{-1},\qquad
 \|(I+E)^{-1}\|\le(1-\epsilon_N)^{-1}.
\]
Use Lemma~\ref{lem:v5v43-D0-inverse}.
\end{proof}

\begin{assessment}[What must be computed]
\label{ass:v5v43-column}
The relevant object is the weighted column sum
\eqref{eq:v5v43-epsilon}, not the largest displayed anomalous-dimension
matrix element.  It retains all output grades generated by one input
operator and includes the inverse canonical eigenvalue at the output.
\end{assessment}

\subsection{Factorial weight ratios}

For \(k\ge1\),
\begin{align}
 \frac{w_{q+k}}{w_q}
 &=
 \rho^k
 \left(\frac{q!}{(q+k)!}\right)^\alpha
 \le
 \rho^k(q+1)^{-\alpha k},
 \label{eq:v5v43-raising-ratio}\\
 \frac{w_{q-k}}{w_q}
 &=
 \rho^{-k}
 \left(\frac{q!}{(q-k)!}\right)^\alpha
 \le
 \rho^{-k}q^{\alpha k}.
 \label{eq:v5v43-lowering-ratio}
\end{align}
Thus the algebra norm damps grade-raising output and amplifies
grade-lowering output.

\begin{theorem}[Finite-band polynomial mixing test]
\label{thm:v5v43-banded}
Suppose \(K_{pq}=0\) unless \(p=q+k\) for
\(-r\le k\le r\), and
\begin{equation}
 \|K_{q+k,q}\|
 \le gA_k(1+q)^{m_k}
 \label{eq:v5v43-banded-bound}
\end{equation}
whenever both indices lie in the tail.  Assume
\begin{align}
 m_k&\le1+\alpha k &&(0\le k\le r),
 \label{eq:v5v43-raising-condition}\\
 m_{-k}&\le1-\alpha k &&(1\le k\le r).
 \label{eq:v5v43-lowering-condition}
\end{align}
Then
\begin{equation}
 \epsilon_N(\rho,\alpha)
 \le
 g\,C(c_0,c_1,N,r)
 \left(
 \sum_{k=0}^rA_k\rho^k+
 \sum_{k=1}^rA_{-k}\rho^{-k}
 \right).
 \label{eq:v5v43-banded-epsilon}
\end{equation}
In particular, the WTIC holds when the right side is below one.
\end{theorem}

\begin{proof}
For \(k\ge0\), use
\eqref{eq:v5v43-raising-ratio},
\eqref{eq:v5v43-banded-bound}, and
\(|\lambda_{q+k}|^{-1}\le C(1+q)^{-1}\).  The resulting power of
\(1+q\) is
\[
 m_k-1-\alpha k\le0.
\]
For output \(p=q-k\), use
\eqref{eq:v5v43-lowering-ratio} and
\(|\lambda_{q-k}|^{-1}\le C_N(1+q)^{-1}\), where the tail boundary only
changes the fixed constant.  The exponent is
\[
 m_{-k}-1+\alpha k\le0.
\]
There are at most \(2r+1\) nonzero outputs per column.  Sum the bounds and
take the supremum in \(q\).
\end{proof}

\begin{corollary}[Nonlowering anomalous mixing]
\label{cor:v5v43-nonlowering}
If \(K_{pq}=0\) for \(p<q\), has finite upward bandwidth, and its
\(k\)-th raised diagonal grows no faster than
\((1+q)^{1+\alpha k}\), then sufficiently small overall load \(g\)
gives a tail inverse.
\end{corollary}

\begin{proof}
Only \eqref{eq:v5v43-raising-condition} remains in
Theorem~\ref{thm:v5v43-banded}.
\end{proof}

\subsection{A sharp lowering counterexample}

Let \(X_p=\mathbb R\), \(\lambda_p=p\), and define the one-step lowering
operator
\begin{equation}
 (Kx)_{p}=x_{p+1}.
 \label{eq:v5v43-lowering-shift}
\end{equation}
Thus \(K_{q-1,q}=1\).

\begin{proposition}[One-step lowering threshold]
\label{prop:v5v43-lowering-threshold}
For the operator \(E=D_0^{-1}K\),
\begin{equation}
 \|E\|_{\rho,\alpha}
 =
 \sup_{q>N+1}
 \frac1{q-1}\frac{w_{q-1}}{w_q}
 =
 \frac1\rho
 \sup_{q>N+1}\frac{q^\alpha}{q-1}.
 \label{eq:v5v43-lowering-norm}
\end{equation}
Hence:
\begin{enumerate}
 \item if \(\alpha>1\), \(E\) is unbounded;
 \item if \(\alpha=1\), \(\|E\|\) is finite and tends to \(\rho^{-1}\)
 as the tail starts farther out;
 \item if \(0\le\alpha<1\), the high-grade part tends to zero.
\end{enumerate}
\end{proposition}

\begin{proof}
Each input column \(q\) has the single output \(p=q-1\), so the weighted
column formula in \eqref{eq:v5v43-epsilon} is exactly
\eqref{eq:v5v43-lowering-norm}.  The three conclusions follow from the
asymptotic \(q^\alpha/(q-1)\sim q^{\alpha-1}\).
\end{proof}

\begin{corollary}[Two-step lowering at Gevrey exponent one]
\label{cor:v5v43-two-step}
At \(\alpha=1\), the bounded two-step shift
\((Kx)_p=x_{p+2}\) makes \(D_0^{-1}K\) unbounded.
\end{corollary}

\begin{proof}
The single column ratio is
\[
 \frac1p\frac{w_p}{w_{p+2}}
 =
 \frac1p\rho^{-2}(p+1)(p+2)
 \sim\rho^{-2}p,
\]
which diverges.
\end{proof}

\begin{firewall}[Canonical growth does not dominate arbitrary lowering]
\label{fw:v5v43-canonical-not-enough}
Linear canonical growth pays only one power of grade.  Factorial norms can
charge \(p^{\alpha k}\) for a \(k\)-step lowering transition.  Therefore
canonical dominance is not a proof without a mixing-direction estimate.
\end{firewall}

\subsection{Strict radius gaps do not directly give a Neumann inverse}

Let \(\rho_+>\rho_->0\).  For fixed-band lowering, the quotient
\[
 \frac{w_p^-}{w_q^+}
 =
 \frac{\rho_-^p}{\rho_+^q}
 \left(\frac{q!}{p!}\right)^\alpha
\]
contains the exponentially decaying factor
\((\rho_-/\rho_+)^q\), which dominates every fixed polynomial in \(q\).
Thus a lowering operator can be bounded from the strong radius to the weak
radius even when it is unbounded on either single-radius space.

\begin{proposition}[Two-radius boundedness for finite lowering bandwidth]
\label{prop:v5v43-two-radius}
Assume the finite-band polynomial bound
\eqref{eq:v5v43-banded-bound}.  Then
\begin{equation}
 D_0^{-1}K:
 \mathbb B_{\rho_+,\eta;s,b}^{\alpha,\beta}
 \longrightarrow
 \mathbb B_{\rho_-,\eta;s,b}^{\alpha,\beta}
 \label{eq:v5v43-two-radius-map}
\end{equation}
is bounded for every fixed bandwidth \(r\).
\end{proposition}

\begin{proof}
For \(p=q+k\) with fixed \(k\), the weak-output/strong-input weight ratio
is the same polynomial factorial quotient as
\eqref{eq:v5v43-raising-ratio} or
\eqref{eq:v5v43-lowering-ratio}, multiplied by
\[
 \left(\frac{\rho_-}{\rho_+}\right)^q
\]
and a fixed \(k\)-dependent radius factor.  Exponential decay in \(q\)
dominates the polynomial growth in
\eqref{eq:v5v43-banded-bound} and the factorial ratio.  The finite number
of diagonals can therefore be summed uniformly.
\end{proof}

\begin{firewall}[A two-radius map is not an endomorphism]
\label{fw:v5v43-two-radius-not-inverse}
Proposition~\ref{prop:v5v43-two-radius} does not justify the Neumann series
in Theorem~\ref{thm:v5v43-WTIC}.  Repeated powers would consume a radius
gap at every iteration.  A scale-of-spaces Newton or Nash--Moser theorem,
or a triangular normal form that removes the lowering terms, is required.
The single-space V42 Newton theorem cannot be cited.
\end{firewall}

\subsection{Separating head feedback from tail lowering}

Some grade-lowering terms land in the finite head \(X_N\); these belong to
the block \(B\) or \(C\) of the V42 Schur matrix rather than to \(D\).
Only lowering transitions with both input and output above \(N\) enter
\eqref{eq:v5v43-epsilon}.

\begin{proposition}[Finite landing is a Schur block]
\label{prop:v5v43-finite-landing}
If \(K_{pq}=0\) whenever \(p,q>N\) and \(p<q\), then the infinite tail
block is nonlowering.  All lowering output lies in the finite head and is
accounted for by the finite--tail cross block in
\eqref{eq:v5v42-block-L}.
\end{proposition}

\begin{proof}
The premise says exactly that the restriction
\(Q_NKQ_N\) has no matrix blocks below its diagonal.  Every remaining
lowering block has \(p\le N<q\), so it equals a component of
\(P_NKQ_N\), one of the Schur cross maps.
\end{proof}

\begin{firewall}[The cross block still needs a bound]
\label{fw:v5v43-cross-block}
Moving a lowering term into the finite head does not make it harmless.
The weighted norm of \(P_NKQ_N\) and the feedback product
\(\|A^{-1}\|\|B\|\|D^{-1}\|\|C\|\) must satisfy the V42 Schur condition.
\end{firewall}

\subsection{The full multi-index certificate}

Restore particle number \(m\), spatial index \(s\), and energy loss
\(bm\).  Let a combined block index be \(I=(p,m,\lambda)\), with weight
\(W_I\) from \eqref{eq:v5v25-trigraded-norm}.  If
\(\Lambda_I\) is the canonical diagonal, define
\begin{equation}
 \epsilon_{\mathrm{full}}
 =
 \sup_J
 \sum_I
 \frac{W_I}{W_J}
 \frac{\|K_{IJ}\|}{|\Lambda_I|}.
 \label{eq:v5v43-full-epsilon}
\end{equation}

\begin{theorem}[Full tri-graded tail criterion]
\label{thm:v5v43-full-tail}
If the canonical inverse is bounded on every block and
\(\epsilon_{\mathrm{full}}<1\), the full tail derivative is invertible in
the hybrid tri-graded norm, with inverse bound
\begin{equation}
 \|D^{-1}\|
 \le
 \frac{\|D_0^{-1}\|}{1-\epsilon_{\mathrm{full}}}.
 \label{eq:v5v43-full-inverse}
\end{equation}
Protected blocks must be assembled before the matrix block norm in
\eqref{eq:v5v43-full-epsilon} is taken.
\end{theorem}

\begin{proof}
Repeat the weighted column-sum proof of
Theorem~\ref{thm:v5v43-WTIC} with the combined indices.  The hybrid
protected component is already one normed block, so no signed cancellation
is split.
\end{proof}

\begin{definition}[Canonical tail inverse certificate]
\label{def:v5v43-CTIC}
A CTIC consists of:
\begin{enumerate}
 \item a canonical block diagonal with a proved spectral gap;
 \item a complete assembled anomalous-mixing matrix;
 \item the column bound \eqref{eq:v5v43-full-epsilon} below one;
 \item explicit treatment of every grade-lowering block, either inside
 that bound, in a finite Schur cross block, or by a proved scale-of-spaces
 inverse theorem;
 \item compatibility with the Ward slice and the CTC multiplicity
 ordering.
\end{enumerate}
\end{definition}

\subsection{V43 status}

\begin{theorem}[V43 conclusion]
\label{thm:v5v43-status}
Canonical dimension yields a diagonal tail inverse, but the exact
criterion for the anomalously mixed inverse is the weighted column sum
\eqref{eq:v5v43-full-epsilon}.  Factorial weights suppress
grade-raising mixing and can amplify grade-lowering mixing; even a bounded
two-step lowering shift is unbounded at Gevrey exponent one after only a
linear canonical inverse.  A CTIC therefore requires a structural
mixing-direction theorem or a scale-of-spaces inverse, neither of which is
yet proved for the full gravitational tower.
\end{theorem}

\begin{proof}
Use Theorem~\ref{thm:v5v43-WTIC},
Theorem~\ref{thm:v5v43-banded}, and
Proposition~\ref{prop:v5v43-lowering-threshold}.
\end{proof}

V44 will attack the structural alternative.  It will grade local
operators by canonical dimension and field number, examine which
linearized shell contractions can lower those grades after local Taylor
subtraction, and formulate the triangularity statement actually needed
for a CTIC.

\clearpage
\section{V44: Wick-degree triangularity, exact lowering, and operator provenance}
\label{sec:v5v44}

V43 reduced the high-grade inverse problem to a weighted column estimate.
This note tests the most direct structural proposal: perhaps Wick ordering
makes the linearized shell map nonlowering.  The statement is exactly true
at the Gaussian base point, where integration of a covariance-matched Wick
polynomial preserves its degree.  It is false at an interacting background.
Already a one-dimensional linear or quadratic Gaussian tilt creates
degree-lowering blocks with polynomial combinatorial growth.  Thus normal
ordering by itself cannot supply the canonical tail inverse certificate.

Two possible repairs are separated.  The first is a quantitative
selected-slot debit: every lost field slot must carry a reciprocal falling
factorial before the block norm is taken.  The second is an operator
provenance lift which retains the grade of the source operator.  A lift is
useful only if its evaluation map is bounded, its physical subspace is a
commuting retract, and its diagonal gap grows in the lifted grade.  These
requirements are made explicit below; none is silently inferred from the
bookkeeping label.

\subsection{Exact Gaussian preservation of Wick degree}

It suffices to work first in a finite-dimensional real vector space
\(E\).  Let \(C_<\) and \(C_>\) be nonnegative covariance forms and put
\(C=C_<+C_>\).  For \(t\in E^*\), define the Wick exponential
\begin{equation}
 {:}\!e^{\langle t,x\rangle}\!{:}_C
 =
 \exp\!\left(\langle t,x\rangle-\frac12 C(t,t)\right).
 \label{eq:v5v44-Wick-exponential}
\end{equation}
Its homogeneous Taylor coefficient of order \(n\) is the Wick tensor
\({:}x^{\otimes n}{:}_C\).

\begin{theorem}[Covariance-matched shell identity]
\label{thm:v5v44-Wick-preservation}
If \(\eta\) is a centred Gaussian vector with covariance \(C_>\), then
\begin{equation}
 \mathbb E_{C_>}
 {:}\!e^{\langle t,x+\eta\rangle}\!{:}_{C}
 =
 {:}\!e^{\langle t,x\rangle}\!{:}_{C_<}.
 \label{eq:v5v44-Wick-shell}
\end{equation}
Consequently, coefficient by coefficient,
\begin{equation}
 \mathbb E_{C_>}
 {:}(x+\eta)^{\otimes n}{:}_{C}
 =
 {:}x^{\otimes n}{:}_{C_<}.
 \label{eq:v5v44-Wick-tensor-preservation}
\end{equation}
Gaussian shell integration therefore preserves Wick degree exactly on
cylinder polynomials.
\end{theorem}

\begin{proof}
The Gaussian moment-generating identity gives
\[
 \mathbb E_{C_>}e^{\langle t,\eta\rangle}
 =e^{C_>(t,t)/2}.
\]
Multiplication by the two remaining factors in
\eqref{eq:v5v44-Wick-exponential} leaves
\[
 \exp\!\left(
 \langle t,x\rangle
 -\frac12(C_<+C_>)(t,t)+\frac12C_>(t,t)
 \right),
\]
which is the right side of \eqref{eq:v5v44-Wick-shell}.  Both sides are
entire in \(t\), so equality of their Taylor coefficients proves
\eqref{eq:v5v44-Wick-tensor-preservation}.
\end{proof}

Consider the normalized shell map
\begin{equation}
 \mathcal R(V)(x)
 =
 -\log\frac{
 \mathbb E_{C_>}e^{-V(x+\eta)}}{
 \mathbb E_{C_>}e^{-V(\eta)}}.
 \label{eq:v5v44-shell-map}
\end{equation}
The denominator fixes \(\mathcal R(V)(0)=0\).

\begin{corollary}[Triangularity at the free base point]
\label{cor:v5v44-free-triangularity}
The derivative at zero is
\begin{equation}
 D\mathcal R(0)h(x)
 =
 \mathbb E_{C_>}h(x+\eta)-\mathbb E_{C_>}h(\eta).
 \label{eq:v5v44-free-derivative}
\end{equation}
On a homogeneous covariance-matched Wick tensor of positive degree
\(n\), it equals the corresponding \(C_<\)-Wick tensor, up to its value
at zero.  In particular, its nonconstant part has degree \(n\) and no
lower positive-degree component.
\end{corollary}

\begin{proof}
Differentiate \eqref{eq:v5v44-shell-map} at \(V=0\).  Then apply
Theorem~\ref{thm:v5v44-Wick-preservation}.  Subtraction at zero affects
only the constant component.
\end{proof}

\subsection{An interacting background lowers Wick degree}

For a general background \(V\), differentiation gives the conditional
expectation
\begin{equation}
 D\mathcal R(V)h(x)
 =
 \frac{\mathbb E[h(x+\eta)e^{-V(x+\eta)}]}
      {\mathbb E[e^{-V(x+\eta)}]}
 -
 \frac{\mathbb E[h(\eta)e^{-V(\eta)}]}
      {\mathbb E[e^{-V(\eta)}]}.
 \label{eq:v5v44-background-derivative}
\end{equation}
The background changes the conditional mean, covariance, and higher
cumulants.  Matching Wick order to the untitled covariance no longer
cancels those changes.

\begin{proposition}[Exact lowering under a linear tilt]
\label{prop:v5v44-linear-tilt}
Let \(E=\mathbb R\), let \(C_< =c_<\ge0\), \(C_>=c_> >0\), and let
\(H_n(\,cdot\,;c)\) denote the Wick polynomial generated by
\[
 e^{tx-ct^2/2}
 =\sum_{n\ge0}H_n(x;c)\frac{t^n}{n!}.
\]
For \(V_g(y)=gy\) and \(h_n(y)=H_n(y;c_<+c_>)\), the first quotient in
\eqref{eq:v5v44-background-derivative} is
\begin{equation}
 \sum_{p=0}^{n}
 \binom np(-gc_>)^{n-p}H_p(x;c_<).
 \label{eq:v5v44-linear-lowering}
\end{equation}
After the normalization at \(x=0\), all nonconstant terms in this sum
remain.  For \(g\ne0\), a degree-\(n\) input therefore has nonzero output
in every parity-allowed or non-parity-restricted lower degree; in
particular its \((n-1)\)-degree coefficient is \(-ngc_>\).
\end{proposition}

\begin{proof}
Under the density proportional to \(e^{-g\eta}\), the Gaussian variable
\(\eta\) has mean \(-gc_>\) and covariance \(c_>\).  Equivalently,
\begin{align*}
 &\frac{\mathbb E\!left[
 e^{t(x+\eta)-(c_<+c_>)t^2/2}e^{-g(x+\eta)}
 \right]}
 {\mathbb E[e^{-g(x+\eta)}]}
 \\
 &\qquad=
 e^{tx-c_<t^2/2}e^{-gc_>t}.
\end{align*}
Expand the first factor in \(C_<\)-Wick polynomials and the second in
powers of \(t\).  The coefficient of \(t^n/n!\) is
\eqref{eq:v5v44-linear-lowering}.  The second quotient in
\eqref{eq:v5v44-background-derivative} subtracts precisely the value of
that polynomial at \(x=0\).
\end{proof}

Odd backgrounds may be excluded by a symmetry.  The obstruction persists
for an even background.

\begin{proposition}[Exact two-step lowering under a covariance tilt]
\label{prop:v5v44-quadratic-tilt}
In the same scalar model, take
\(V_g(y)=g y^2/2\) with \(1+gc_> >0\), and set
\begin{equation}
 c_g=\frac{c_>}{1+gc_>},
 \qquad
 \delta c=c_g-c_>.
 \label{eq:v5v44-tilted-covariance}
\end{equation}
Apart from the normalization at \(x=0\), the derivative sends
\(H_n(\,cdot\,;c_<+c_>)\) to
\begin{equation}
 \sum_{k=0}^{\lfloor n/2\rfloor}
 \frac{n!}{(n-2k)!\,k!\,2^k}
 (\delta c)^k H_{n-2k}(x;c_<).
 \label{eq:v5v44-quadratic-lowering}
\end{equation}
Thus the first lowered diagonal has coefficient
\(n(n-1)\delta c/2\), which grows quadratically in the input degree.
\end{proposition}

\begin{proof}
The quadratic tilt changes the shell covariance from \(c_>\) to \(c_g\)
and does not change its mean.  The generating function of the first
quotient is therefore
\[
 e^{tx-c_<t^2/2}e^{\delta c\,t^2/2}.
\]
Multiplying the two Taylor series and extracting the coefficient of
\(t^n/n!\) gives \eqref{eq:v5v44-quadratic-lowering}.  Evaluation at zero
again changes only the constant output.
\end{proof}

\begin{firewall}[Wick order is not interacting triangularity]
\label{fw:v5v44-Wick-not-triangular}
Theorem~\ref{thm:v5v44-Wick-preservation} concerns integration against
the same Gaussian covariance used to define the Wick tensors.
Propositions~\ref{prop:v5v44-linear-tilt} and
\ref{prop:v5v44-quadratic-tilt} show that the derivative at a nonzero
background contains growing lower diagonals.  Neither normal ordering nor
an even symmetry proves the nonlowering hypothesis of
Corollary~\ref{cor:v5v43-nonlowering}.
\end{firewall}

\subsection{The exact debit needed for physical lowering}

Retain the V43 weights
\begin{equation}
 w_p=\frac{\rho^p}{(p!)^\alpha},
 \qquad \alpha\ge0,
 \label{eq:v5v44-weights}
\end{equation}
and the canonical diagonal \(D_0\) with eigenvalues \(\lambda_p\).
The following criterion permits arbitrary lowering distance but records
the exact combinatorial price required by the norm.

\begin{theorem}[Selected-slot debit criterion]
\label{thm:v5v44-slot-debit}
Suppose that for \(q\ge p>N\),
\begin{equation}
 \|K_{pq}\|
 \le
 g\,|\lambda_p|\,b_{q-p}
 \left(\frac{p!}{q!}\right)^\alpha,
 \label{eq:v5v44-debit-bound}
\end{equation}
where \(b_k\ge0\), and suppose the nonlowering part of the weighted
column sum is at most \(\epsilon_+\).  Then
\begin{equation}
 \epsilon_N(\rho,\alpha)
 \le
 \epsilon_+
 +g\sum_{k\ge0}b_k\rho^{-k}.
 \label{eq:v5v44-debit-column}
\end{equation}
If the right side is below one, the tail derivative is invertible by
Theorem~\ref{thm:v5v43-WTIC}.
\end{theorem}

\begin{proof}
For a lowering block write \(k=q-p\).  Multiplying
\eqref{eq:v5v44-debit-bound} by the weighted inverse factor gives
\begin{align*}
 \frac{w_p}{w_q}\frac{\|K_{pq}\|}{{|\lambda_p|}}
 &\le
 \rho^{p-q}
 \left(\frac{q!}{p!}\right)^\alpha
 g b_{q-p}
 \left(\frac{p!}{q!}\right)^\alpha
 \\
 &=g b_k\rho^{-k}.
\end{align*}
Sum over the possible outputs in each input column and add the assumed
nonlowering contribution.  The result is
\eqref{eq:v5v44-debit-column}; V43 then supplies the Neumann inverse.
\end{proof}

\begin{corollary}[Geometric debit]
\label{cor:v5v44-geometric-debit}
If \(b_k=A^k\), \(\rho>A\), and
\begin{equation}
 \epsilon_++\frac{g}{1-A/\rho}<1,
 \label{eq:v5v44-geometric-condition}
\end{equation}
then the CTIC tail inverse exists.
\end{corollary}

\begin{proof}
Sum the geometric series in
\eqref{eq:v5v44-debit-column}.
\end{proof}

The factor in \eqref{eq:v5v44-debit-bound} is a reciprocal falling
factorial:
\[
 \left(\frac{p!}{q!}\right)^\alpha
 =
 \left((p+1)(p+2)\cdots q\right)^{-\alpha}.
\]
It must be present in the assembled block itself.  The scalar examples
show that ordinary Wick contractions generally produce falling
factorials in the opposite direction.

\begin{corollary}[The tilted Gaussian violates the debit]
\label{cor:v5v44-tilt-violates-debit}
For \(\alpha>0\), fixed nonzero \(g\), and a canonical eigenvalue growing
at most polynomially, neither the one-step coefficient
\(-ngc_>\) in Proposition~\ref{prop:v5v44-linear-tilt} nor the two-step
coefficient \(n(n-1)\delta c/2\) in
Proposition~\ref{prop:v5v44-quadratic-tilt} satisfies a uniform bound of
the form \eqref{eq:v5v44-debit-bound} on the corresponding diagonal.
\end{corollary}

\begin{proof}
For \(q=n,p=n-1\), the debit contributes \(n^{-\alpha}\), while the
actual coefficient grows as \(n\).  For \(p=n-2\), it contributes
\(((n-1)n)^{-\alpha}\), while the actual coefficient grows as \(n^2\).
Multiplication by a fixed polynomial canonical factor cannot reverse the
specified reciprocal-factorial requirement uniformly unless an
additional compensating estimate is proved.  In the linear canonical
case of V43 the failure is immediate.
\end{proof}

\begin{assessment}[Where a debit could arise]
\label{ass:v5v44-debit-source}
The YM coincident-transition calculation inspected at V40 shows a
plausible source of a debit: selected centred legs can be paid inside a
complete cumulant and an equivariant multiplicity partial trace before
norms are taken.  To use Theorem~\ref{thm:v5v44-slot-debit}, the full
SM--Einstein operator block must retain enough such payments to establish
\eqref{eq:v5v44-debit-bound}.  A diagram count performed after absolute
values is not sufficient.
\end{assessment}

\subsection{A rigorous operator-provenance lift}

An alternative is to retain the source grade after contractions.  To
avoid confusing this label with the inverse-limit ancestries of V40--V41,
call it the \emph{operator provenance} \(a\).  A typical lifted block has
indices \((a,p)\), where \(a\) is the source grade and \(p\le a\) is the
physical degree left after contractions.

The lift has mathematical content only if it returns the physical
operator space through controlled maps.  The following elementary
retraction theorem states the exact requirement.

\begin{theorem}[Inverse transfer through a commuting retract]
\label{thm:v5v44-retract-inverse}
Let \(X\) and \(\widetilde X\) be Banach spaces.  Suppose
\begin{equation}
 I:X\longrightarrow\widetilde X,
 \qquad
 E:\widetilde X\longrightarrow X,
 \qquad EI=1_X
 \label{eq:v5v44-retract}
\end{equation}
are bounded.  Put \(P=IE\).  Let \(\widetilde D\) be a boundedly
invertible operator on \(\widetilde X\) such that
\begin{equation}
 P\widetilde D=\widetilde D P,
 \label{eq:v5v44-commuting-projection}
\end{equation}
and define
\begin{equation}
 D=E\widetilde D I.
 \label{eq:v5v44-physical-D}
\end{equation}
Then \(D\) is boundedly invertible and
\begin{equation}
 D^{-1}=E\widetilde D^{-1}I,
 \qquad
 \|D^{-1}\|
 \le\|E\|\,\|\widetilde D^{-1}\|\,\|I\|.
 \label{eq:v5v44-transferred-inverse}
\end{equation}
\end{theorem}

\begin{proof}
The relation \(P\widetilde D=\widetilde DP\) also implies
\(P\widetilde D^{-1}=\widetilde D^{-1}P\).  Using \(EI=1_X\), hence
\(PI=I\) and \(EP=E\), compute
\begin{align*}
 D(E\widetilde D^{-1}I)
 &=E\widetilde D I E\widetilde D^{-1}I
 =E\widetilde D P\widetilde D^{-1}I
 =EPI=1_X,\\
 (E\widetilde D^{-1}I)D
 &=E\widetilde D^{-1}I E\widetilde D I
 =E\widetilde D^{-1}P\widetilde D I
 =EPI=1_X.
\end{align*}
The norm estimate follows by submultiplicativity.
\end{proof}

\begin{corollary}[Lifted CTIC]
\label{cor:v5v44-lifted-CTIC}
Suppose a provenance-lifted derivative
\(\widetilde D=\widetilde D_0+\widetilde K\) satisfies:
\begin{enumerate}
 \item a block diagonal estimate
 \( |\widetilde\lambda_{a,p}|\ge c_0+c_1a\);
 \item the full V43 weighted column bound for the lifted indices is below
 one;
 \item the physical evaluation and section satisfy
 \eqref{eq:v5v44-retract}--\eqref{eq:v5v44-commuting-projection}.
\end{enumerate}
Then the physical derivative has a bounded inverse with the estimate
\eqref{eq:v5v44-transferred-inverse} and the lifted V43 Neumann bound.
\end{corollary}

\begin{proof}
The first two assumptions give a bounded inverse for
\(\widetilde D\) by Theorem~\ref{thm:v5v43-full-tail}.  Apply
Theorem~\ref{thm:v5v44-retract-inverse}.
\end{proof}

\begin{firewall}[A provenance label supplies no spectral gap]
\label{fw:v5v44-label-no-gap}
If the rescaling eigenvalue of a lifted block \((a,p)\) still grows only
with the surviving physical grade \(p\), then retaining \(a\) does not
prove item 1 of Corollary~\ref{cor:v5v44-lifted-CTIC}.  The V43 lowering
counterexample embeds unchanged in a fixed-\(p\) slice with
\(a\to\infty\).  A provenance construction must derive its diagonal gap
from the actual shell dynamics or absorb the low physical descendants
into a controlled finite Schur block.
\end{firewall}

\begin{firewall}[Formal evaluation may be unbounded]
\label{fw:v5v44-evaluation-unbounded}
The physical evaluation usually sums all provenance sectors that produce
the same local operator.  Such a sum need not be bounded in the lifted
factorial norm.  The maps \(E\) and \(I\), their norms, and the commuting
identity \eqref{eq:v5v44-commuting-projection} are analytic parts of the
certificate, not bookkeeping conventions.
\end{firewall}

\begin{definition}[Provenance tail certificate]
\label{def:v5v44-PTC}
A provenance tail certificate consists of:
\begin{enumerate}
 \item a lifted local-operator space with source grade \(a\), physical
 degree \(p\), particle number, spatial regularity, and energy loss;
 \item an exact lifted shell derivative whose evaluation is the physical
 derivative;
 \item a proved lifted canonical gap and a full weighted column sum below
 one;
 \item bounded section and evaluation maps forming a commuting retract;
 \item Ward, BRST, multiplicity, and local-subtraction compatibility
 before passage to block norms;
 \item either a selected-slot debit for physical lowering or a finite
 Schur allocation of every low descendant.
\end{enumerate}
\end{definition}

\subsection{V44 status}

\begin{theorem}[V44 conclusion]
\label{thm:v5v44-status}
Covariance-matched Gaussian shell integration preserves Wick degree at
the free base point, but the linearized map at a nonzero background need
not be triangular.  Linear and quadratic Gaussian tilts give exact
lowering coefficients growing respectively as \(n\) and \(n^2\).
Therefore the CTIC requires either the quantitative reciprocal-factorial
debit \eqref{eq:v5v44-debit-bound}, a bounded provenance lift satisfying
the full PTC, or a scale-of-spaces inverse theorem.  None of these three
structures has yet been proved for the full SM--Einstein tower.
\end{theorem}

\begin{proof}
Use Theorem~\ref{thm:v5v44-Wick-preservation},
Propositions~\ref{prop:v5v44-linear-tilt} and
\ref{prop:v5v44-quadratic-tilt},
Theorem~\ref{thm:v5v44-slot-debit}, and
Theorem~\ref{thm:v5v44-retract-inverse}.
\end{proof}

V45 will perform the required companion audit.  It will then test whether
the YM selected-leg mechanism supplies the reciprocal-factorial debit in
a representation-summed operator column, or only a fixed finite number
of boundary factors.  The distinction decides whether the physical CTIC
can close without a provenance lift.

\clearpage
\section{V45: companion audit, corrected covariance tilt, and a moving Gaussian frame}
\label{sec:v5v45}

V44 showed that a fixed Wick basis is not generally triangular at a
nonzero background.  Before using that statement, this note corrects the
quadratic-tilt formula: the conditional shell Gaussian acquires an
\(x\)-dependent mean as well as a new covariance.  The correction leaves
the qualitative conclusion intact but changes the high-degree size of the
lowering coefficients.  It also reveals a better coordinate choice.  In a
quadratic background, an exactly transported Wick frame restores degree
preservation and turns the apparent lowering into a basis-change matrix.

This is the compulsory fifth-version cross-program audit.  The newest YM
note is V51 and the NS note remains V31.  The YM result supports a strict
``numerator first, one payer last'' rule and proves that representation
multiplicity need not be lost.  It does not provide a debit whose order
grows with the number of lowered field slots.  The NS endpoint is unchanged
and adds no new estimate to the present tail-inverse problem.

\subsection{Corrigendum to the V44 quadratic tilt}

Retain the notation of Proposition~\ref{prop:v5v44-quadratic-tilt}.  Put
\begin{equation}
 c=c_<+c_>,
 \qquad
 a_g=\frac1{1+gc_>},
 \qquad
 c_g=\frac{c_>}{1+gc_>}.
 \label{eq:v5v45-quadratic-parameters}
\end{equation}
The conditional law of \(\eta\), at fixed \(x\), under the density
proportional to \(e^{-g(x+\eta)^2/2}\) has mean
\(-gc_gx\) and covariance \(c_g\).  Hence \(x+\eta\) has mean \(a_gx\),
not \(x\).

\begin{theorem}[Correct quadratic-tilt formula]
\label{thm:v5v45-corrected-quadratic}
Assume \(1+gc_> >0\), and define
\begin{equation}
 \Delta_g=c_g-c+c_<a_g^2.
 \label{eq:v5v45-Delta}
\end{equation}
For \(h_n(y)=H_n(y;c)\), the first quotient in
\eqref{eq:v5v44-background-derivative} equals
\begin{equation}
 \sum_{k=0}^{\lfloor n/2\rfloor}
 \frac{n!}{(n-2k)!\,k!\,2^k}
 \Delta_g^k a_g^{,n-2k}
 H_{n-2k}(x;c_<).
 \label{eq:v5v45-corrected-expansion}
\end{equation}
The normalized derivative subtracts the value at \(x=0\).  In
particular, its first nonconstant lowered coefficient is
\begin{equation}
 \frac{n(n-1)}2\Delta_g a_g^{,n-2}.
 \label{eq:v5v45-first-lowered}
\end{equation}
\end{theorem}

\begin{proof}
Completing the square gives
\[
 -\frac{\eta^2}{2c_>}-\frac g2(x+\eta)^2
 =
 -\frac{(\eta+gc_gx)^2}{2c_g}
 -\frac{g a_gx^2}{2}.
\]
The last factor cancels between numerator and denominator.  Therefore
\begin{align*}
 &\frac{\mathbb E\!left[
 e^{t(x+\eta)-ct^2/2}e^{-g(x+\eta)^2/2}
 \right]}
 {\mathbb E[e^{-g(x+\eta)^2/2}]}
 \\
 &\qquad=
 \exp\!\left(a_gtx-\frac12(c-c_g)t^2\right)\\
 &\qquad=
 \exp\!\left(a_gtx-\frac12c_<a_g^2t^2\right)
 \exp\!\left(\frac12\Delta_gt^2\right).
\end{align*}
Expand the first factor in \(H_p(x;c_<)(a_gt)^p/p!\), expand the
second in powers of \(t^2\), and take the coefficient of \(t^n/n!\).
This proves \eqref{eq:v5v45-corrected-expansion}.  The second quotient in
\eqref{eq:v5v44-background-derivative} is evaluation at zero.
\end{proof}

\begin{assessment}[Scope of the correction]
\label{ass:v5v45-correction-scope}
Equation~\eqref{eq:v5v45-corrected-expansion} supersedes
\eqref{eq:v5v44-quadratic-lowering}.  The combinatorial prefactor of the
two-step block is quadratic in \(n\), but it is multiplied by
\(a_g^{n-2}\).  For \(g>0\), this factor decays exponentially; for
\(-c_>^{-1}<g<0\), it grows exponentially.  Thus V44's assertion of
unqualified quadratic growth and the quadratic part of
Corollary~\ref{cor:v5v44-tilt-violates-debit} are withdrawn.  The exact
degree-lowering statement remains true whenever \(\Delta_g\ne0\).
The linear-tilt obstruction remains valid for the V43 linearly growing
canonical diagonal; more generally it obstructs a degree-\(d\) canonical
upper bound whenever \(\alpha>d-1\).
\end{assessment}

\subsection{Mandatory V45 companion audit}

The files inspected were
\begin{align}
 &\texttt{Renewal\_Yang\_Mills\_Mass\_Gap\_Research\_Notes\_V51.tex},
 \notag\\
 &\hspace{1em}\text{SHA--256 }
 \texttt{1D3BEF0BBD29D2E7A8858A84FDFA549EB95BED49F741687F13BD6047567CC497},
 \label{eq:v5v45-YM-hash}\\
 &\texttt{Renewal\_NS\_Stability\_Research\_Notes\_V31.tex},
 \notag\\
 &\hspace{1em}\text{SHA--256 }
 \texttt{F1121B62238AD5491BB7CF03635EA9934942EDF573ED8FAAA9EE79E1D38A6696}.
 \label{eq:v5v45-NS-hash}
\end{align}

The relevant new YM statement is its projection-complete ternary prefix
numerator.  It assembles the entire local third cumulant before output
projection, includes the scalar and nontrivial channels, retains capped
pair influences, and lifts to raw BFA coefficients with exact cancellation
of output dimension and no output-multiplicity factor.  No reduced
resolvent is inserted in the prefix; one mass--resolvent payer is used only
after the final joining coordinate.  Its remaining hot-center branch is
explicitly unresolved because singleton spectator means were bounded by
one too early.

The V46--V51 YM chain also proves two-selectable-leg bounds for each fixed
cumulant order and preserves them when a Casimir payer occupies the same
equivariant multiplicity block.  The number of selected legs in that
theorem is two; it does not increase with cumulant order or with an
arbitrary operator-degree loss.

The NS V31 hash is unchanged from the V40 audit.  Its endpoint remains the
stopping-time-correlated formation estimate for
\(\mathscr Q_{\rm form}\), with a required gain of one critical record
factor.  It neither supplies a new factorial operator-column estimate nor
changes the present SM direction.

\begin{definition}[One-payer numerator rule]
\label{def:v5v45-OPNR}
For one linearized SM shell column, the OPNR requires:
\begin{enumerate}
 \item assemble all connected local cumulants, Ward partners, scalar and
 nontrivial isotypic outputs, and multiplicity partial traces into the
 physical block \(K_{pq}\);
 \item estimate that complete numerator before applying a canonical
 inverse;
 \item apply the single output factor \(|\lambda_p|^{-1}\) only in the
 weighted column sum \eqref{eq:v5v43-full-epsilon}.
\end{enumerate}
An internal inverse is permitted only if it occurs in the exact shell
formula and is not the final Newton or tail inverse counted again.
\end{definition}

\begin{proposition}[Audit transfer]
\label{prop:v5v45-audit-transfer}
The YM V51 numerator theorem supports items 1--2 of the OPNR for its
ternary Wilson sector and proves that output isotypic multiplicity can be
summed without a counting loss.  It does not imply the all-grade lowering
estimate \eqref{eq:v5v44-debit-bound}.
\end{proposition}

\begin{proof}
The first assertion is the content of the projection-complete numerator
and raw pinching estimates just summarized.  For the second, the available
selected-leg theorem supplies two Wilson debits at every fixed cumulant
order.  Bound \eqref{eq:v5v44-debit-bound}, for a loss of
\(k=q-p\) slots, requires a reciprocal falling factorial of order
\(\alpha k\).  No statement in V51 makes the number of selected debits
grow with \(k\).  Therefore the fixed-order result does not imply the
required all-column estimate.
\end{proof}

\begin{firewall}[A companion result transfers only in its proved scope]
\label{fw:v5v45-audit-scope}
The YM estimate concerns compact-group Wilson cumulants in its declared
coordinate and representation norm.  Its exact assembly order may be
reused.  Its numerical debit cannot be assigned to SM matter, Higgs,
ghost, or metric blocks until the corresponding complete cumulant and Ward
identities are proved in those blocks.
\end{firewall}

\subsection{A necessary estimate on every lowering diagonal}

V44 gave a sufficient selected-slot debit.  The converse one-column test
shows how much decay is unavoidable.

\begin{theorem}[Necessary lowering-diagonal decay]
\label{thm:v5v45-necessary-diagonal}
Let \(D_0\) be block diagonal and let
\(E=D_0^{-1}K\) be bounded on the V43 weighted \(\ell^1\) tail with
\(\|E\|\le M\).  Then every block satisfies
\begin{equation}
 \|K_{pq}\|
 \le
 M|\lambda_p|\frac{w_q}{w_p}.
 \label{eq:v5v45-block-necessary}
\end{equation}
In particular, for \(p=q-k>N\),
\begin{equation}
 \|K_{q-k,q}\|
 \le
 M|\lambda_{q-k}|\rho^k
 \left(\frac{(q-k)!}{q!}\right)^\alpha.
 \label{eq:v5v45-lowering-necessary}
\end{equation}
If \(|\lambda_p|\le C(1+p)^d\), then for each fixed \(k\),
\begin{equation}
 \|K_{q-k,q}\|=O\!\left(q^{d-\alpha k}\right).
 \label{eq:v5v45-power-necessary}
\end{equation}
\end{theorem}

\begin{proof}
Choose an input supported only in block \(q\), with a vector approaching
the norm of \(D_{0,p}^{-1}K_{pq}\).  Boundedness of \(E\) gives
\[
 w_p\|D_{0,p}^{-1}K_{pq}\|
 \le M w_q.
\]
Use \(\|D_{0,p}^{-1}\|=|\lambda_p|^{-1}\) in the scalar-block model, or
the corresponding lower singular-value bound in a finite block, to obtain
\eqref{eq:v5v45-block-necessary}.  Insert
\(w_j=\rho^j/(j!)^\alpha\) to obtain
\eqref{eq:v5v45-lowering-necessary}.  Finally,
\[
 \frac{(q-k)!}{q!}
 =\frac1{q(q-1)\cdots(q-k+1)}=O(q^{-k})
\]
for fixed \(k\), which proves \eqref{eq:v5v45-power-necessary}.
\end{proof}

\begin{corollary}[A fixed debit depth cannot close unbounded lowering]
\label{cor:v5v45-fixed-debit-insufficient}
Suppose the available numerator estimate pays at most a fixed power
\(q^{-\sigma}\), independently of the lowering distance \(k\), while the
canonical eigenvalue has degree \(d\).  On the V43 norm with
\(\alpha>0\), that estimate can verify
\eqref{eq:v5v45-power-necessary} only for
\begin{equation}
 k\le\frac{d+\sigma-r}{\alpha}
 \label{eq:v5v45-fixed-depth}
\end{equation}
if the unpaid numerator grows like \(q^r\).  It cannot verify all
lowering diagonals unless the lowering bandwidth is finite or the number
of debits grows with \(k\).
\end{corollary}

\begin{proof}
After the fixed debit, the available upper bound is of order
\(q^{r-\sigma}\).  The necessary target from
\eqref{eq:v5v45-power-necessary} is \(q^{d-\alpha k}\).  The former is no
larger at high \(q\) only when
\(r-\sigma\le d-\alpha k\), which is
\eqref{eq:v5v45-fixed-depth}.  The right side is a finite bound on \(k\).
\end{proof}

\begin{firewall}[An upper bound failing the test is not a divergence proof]
\label{fw:v5v45-bound-not-divergence}
Corollary~\ref{cor:v5v45-fixed-debit-insufficient} says that a fixed-debit
estimate cannot prove the CTIC.  It does not say that the exact blocks
must saturate that estimate.  Additional cancellation or an adapted frame
may make the true columns smaller.
\end{firewall}

\subsection{Finite primitive valence does not give finite bandwidth}

The corrected quadratic formula also closes a logical loophole in the V43
finite-band test.

\begin{proposition}[Quadratic resummation has unbounded lowering distance]
\label{prop:v5v45-quadratic-unbounded-band}
If \(\Delta_g\ne0\), then for every \(k\ge1\) and every \(n\ge2k\), the
quadratic-background derivative contains a block from degree \(n\) to
degree \(n-2k\), with coefficient
\begin{equation}
 \frac{n!}{(n-2k)!\,k!\,2^k}
 \Delta_g^k a_g^{n-2k}.
 \label{eq:v5v45-all-even-diagonals}
\end{equation}
Thus a background of primitive valence two can generate infinitely many
lowering diagonals after exact resummation.
\end{proposition}

\begin{proof}
This is the \(k\)-th term of
\eqref{eq:v5v45-corrected-expansion}.  For every fixed \(k\), it is
nonzero when \(n\ge2k\), \(\Delta_g\ne0\), and \(a_g\ne0\).
\end{proof}

\subsection{The moving Gaussian frame}

The same calculation which produces all lower diagonals also sums them
exactly.  This shows that, for a quadratic background, lowering is a
connection between two Wick frames rather than an intrinsic loss of
degree.

\begin{theorem}[Exact transported Wick basis in one dimension]
\label{thm:v5v45-moving-Wick-scalar}
Under the assumptions of
Theorem~\ref{thm:v5v45-corrected-quadratic}, define
\begin{equation}
 \widehat c_g=\frac{c-c_g}{a_g^2}.
 \label{eq:v5v45-transported-covariance}
\end{equation}
Then
\begin{equation}
 \frac{\mathbb E[H_n(x+\eta;c)e^{-g(x+\eta)^2/2}]}
      {\mathbb E[e^{-g(x+\eta)^2/2}]}
 =
 a_g^n H_n(x;\widehat c_g).
 \label{eq:v5v45-moving-Wick-identity}
\end{equation}
Thus the quadratic-background derivative is diagonal in degree when the
output Wick covariance is transported from \(c\) to \(\widehat c_g\).
For \(g>0\), \(|a_g|<1\), so the diagonal factor gives exponential
high-degree contraction.
\end{theorem}

\begin{proof}
The generating function computed in
Theorem~\ref{thm:v5v45-corrected-quadratic} is
\[
 \exp\!\left(a_gtx-\frac12(c-c_g)t^2\right)
 =
 \exp\!\left(
 (a_gt)x-\frac12\widehat c_g(a_gt)^2
 \right).
\]
The coefficient of \(t^n/n!\) on the right is
\(a_g^nH_n(x;\widehat c_g)\).  If \(g>0\), then
\(0<a_g<1\).
\end{proof}

The finite-dimensional form is equally direct.  Let \(C_>\) be positive
definite, let \(G=G^*\) with \(C_>^{-1}+G>0\), and put
\begin{equation}
 C_G=(C_>^{-1}+G)^{-1},
 \qquad
 A_G=I-C_GG.
 \label{eq:v5v45-matrix-parameters}
\end{equation}
Assume \(A_G\) is invertible and let \(C=C_<+C_>\).

\begin{theorem}[Finite-dimensional transported Gaussian frame]
\label{thm:v5v45-moving-Wick-matrix}
Define the output covariance
\begin{equation}
 \widehat C_G
 =A_G^{-1}(C-C_G)A_G^{-*}.
 \label{eq:v5v45-matrix-output-covariance}
\end{equation}
For every \(t\in E^*\),
\begin{align}
 &\frac{\mathbb E_{C_>}\!left[
 {:}e^{\langle t,x+\eta\rangle}{:}_{C}
 e^{-\langle x+\eta,G(x+\eta)\rangle/2}
 \right]}
 {\mathbb E_{C_>}e^{-\langle x+\eta,G(x+\eta)\rangle/2}}
 \notag\\
 &\qquad=
 {:}e^{\langle A_G^*t,x\rangle}{:}_{\widehat C_G}.
 \label{eq:v5v45-matrix-moving-frame}
\end{align}
Consequently the order-\(n\) Wick tensor maps to the order-\(n\) Wick
tensor under the linear argument transformation \(A_G^*\), with no
lower-degree term in the transported frame.
\end{theorem}

\begin{proof}
Completion of the square shows that the tilted law of \(\eta\) has
covariance \(C_G\) and mean \(-C_GGx\).  Thus \(x+\eta\) has mean
\(A_Gx\), and the left side of
\eqref{eq:v5v45-matrix-moving-frame} is
\[
 \exp\!\left(
 \langle t,A_Gx\rangle
 -\frac12(C-C_G)(t,t)
 \right).
\]
By \eqref{eq:v5v45-matrix-output-covariance},
\[
 \widehat C_G(A_G^*t,A_G^*t)=(C-C_G)(t,t).
\]
The last exponential is therefore the right side of
\eqref{eq:v5v45-matrix-moving-frame}.  Equality of homogeneous Taylor
coefficients proves the tensor statement.
\end{proof}

\begin{firewall}[The transported frame is not yet an SM construction]
\label{fw:v5v45-frame-not-SM}
Theorem~\ref{thm:v5v45-moving-Wick-matrix} is finite dimensional and
Gaussian.  In the full theory, the Hessian may be differential,
gauge-degenerate before slicing, indefinite in conformal directions, and
dependent on the low field.  The transported covariance must be local or
quasilocal in the V25 norm, respect BRST and Ward identities, and have a
bounded connection between adjacent shells.  None of these properties is
automatic.
\end{firewall}

\subsection{Gaussian-frame reduction of a general background}

Let a background split as
\begin{equation}
 V(y)=\frac12\langle y,Gy\rangle+W(y).
 \label{eq:v5v45-background-split}
\end{equation}
The preceding change of Gaussian measure is exact.

\begin{proposition}[Exact reduction to the non-Gaussian remainder]
\label{prop:v5v45-remainder-reduction}
Under the hypotheses of
Theorem~\ref{thm:v5v45-moving-Wick-matrix}, write
\(y=A_Gx+\zeta\), where \(\zeta\) is centred Gaussian with covariance
\(C_G\).  Then the first quotient in
\eqref{eq:v5v44-background-derivative} is exactly
\begin{equation}
 \frac{
 \mathbb E_{C_G}[h(A_Gx+\zeta)e^{-W(A_Gx+\zeta)}]}
 {
 \mathbb E_{C_G}[e^{-W(A_Gx+\zeta)}]}.
 \label{eq:v5v45-remainder-conditional}
\end{equation}
All degree lowering beyond the transported Gaussian frame is therefore
carried by connected cumulants containing at least one insertion of
\(W\).
\end{proposition}

\begin{proof}
Complete the square in the Gaussian density as in
Theorem~\ref{thm:v5v45-moving-Wick-matrix}.  The factor depending only on
\(x\) cancels between numerator and denominator, and the affine change
\(x+\eta=A_Gx+\zeta\) gives
\eqref{eq:v5v45-remainder-conditional}.  Expanding the quotient by
connected cumulants, its zeroth term is the transported Gaussian map and
every remaining term contains \(W\).
\end{proof}

\begin{definition}[Moving-frame tail certificate]
\label{def:v5v45-MFTC}
An MFTC consists of:
\begin{enumerate}
 \item a positive sliced quadratic shell Hessian and its exact transported
 Wick frame;
 \item uniform bounds for the affine contraction \(A_G\), transported
 covariance, and adjacent-shell frame connection in the tri-graded norm;
 \item a complete connected-cumulant numerator estimate for the remainder
 \(W\), assembled according to the OPNR;
 \item an all-lowering column bound for that remainder, with selected
 debits increasing with lowering distance or an equivalent summable
 capped-influence estimate;
 \item Ward, BRST, OS, and multiplicity compatibility of the frame change.
\end{enumerate}
\end{definition}

\subsection{V45 status}

\begin{theorem}[V45 conclusion]
\label{thm:v5v45-status}
The corrected quadratic calculation has infinitely many lowering
diagonals in a fixed Wick basis, but all of them are summed by an exact
transported Gaussian frame in which degree is preserved.  The YM V51
audit justifies complete numerator assembly and one final payer, including
scalar output and multiplicity pinching, but its two-selected-leg estimate
does not imply the lowering-distance-dependent decay forced by
Theorem~\ref{thm:v5v45-necessary-diagonal}.  The next viable tail route is
therefore an MFTC: absorb the exact Hessian into a moving frame and apply
connected selected-leg estimates only to the non-Gaussian remainder.
This certificate is not yet proved for the SM--Einstein shell map.
\end{theorem}

\begin{proof}
Use Theorems~\ref{thm:v5v45-corrected-quadratic},
\ref{thm:v5v45-necessary-diagonal}, and
\ref{thm:v5v45-moving-Wick-matrix}, together with
Proposition~\ref{prop:v5v45-audit-transfer} and
Proposition~\ref{prop:v5v45-remainder-reduction}.
\end{proof}

V46 will derive the connection between two transported Wick frames.  It
will determine whether a small covariance change is bounded on one V25
factorial radius, only across a radius scale, or only after the connection
is kept inside the complete cumulant numerator.  This is the first missing
analytic item of the MFTC.

\clearpage
\section{V46: Wick-frame connection bounds and a split Gevrey grading}
\label{sec:v5v46}

V45 diagonalized a quadratic-background shell derivative by transporting
the Wick covariance.  The missing question was whether changing from one
transported frame to the next is bounded in the V25 factorial norm.  This
note computes that connection exactly.  With Wick tensors normalized by
\(n!\), the connection is a constant-coefficient heat operator.  It is
bounded on a fixed analytic grade norm, but it is unbounded on every fixed
positive-factorial-exponent norm.  Across a strict radius gap it is bounded
for exponent below one half, and at exponent one half under an explicit
smallness condition.  It is unbounded across every radius gap when the
exponent is above one half.

This threshold forces a refinement of the V25 grade.  Field number, on
which covariance connections act, must be distinguished from the
provenance or graph grade which may require a larger Gevrey exponent.  A
split norm is constructed below and remains a Banach algebra when both
grades add under products.  The construction does not yet prove that the
full renormalized SM--Einstein operator quotient admits this split.

\subsection{Normalized Wick tensors and the exact connection}

For \(c\ge0\), write
\begin{equation}
 e_n^{(c)}(x)=\frac1{n!}H_n(x;c),
 \label{eq:v5v46-normalized-Wick}
\end{equation}
where the Hermite--Wick polynomials are those used in V44--V45.

\begin{lemma}[Connection formula]
\label{lem:v5v46-connection-formula}
For any two covariance parameters \(c_1,c_2\),
\begin{equation}
 e_q^{(c_1)}
 =
 \sum_{k=0}^{\lfloor q/2\rfloor}
 \frac1{k!}left(\frac{c_2-c_1}{2}\right)^k
 e_{q-2k}^{(c_2)}.
 \label{eq:v5v46-Wick-connection}
\end{equation}
\end{lemma}

\begin{proof}
Factor the generating function as
\[
 e^{tx-c_1t^2/2}
 =e^{tx-c_2t^2/2}e^{(c_2-c_1)t^2/2}.
\]
Expand both factors and compare the coefficient of \(t^q\).
\end{proof}

For \(\delta=c_2-c_1\), the induced coefficient connection is
\begin{equation}
 (T_\delta u)_p
 =
 \sum_{k\ge0}
 \frac1{k!}\left(\frac\delta2\right)^k u_{p+2k}.
 \label{eq:v5v46-Tdelta}
\end{equation}
On polynomials this is the transpose coefficient action of
\(e^{(\delta/2)\partial_x^2}\).

\begin{lemma}[Connection group law]
\label{lem:v5v46-group-law}
On finite sequences,
\begin{equation}
 T_{\delta_1}T_{\delta_2}=T_{\delta_1+\delta_2},
 \qquad T_0=I,
 \qquad T_\delta^{-1}=T_{-\delta}.
 \label{eq:v5v46-group-law}
\end{equation}
\end{lemma}

\begin{proof}
Constant-coefficient derivatives commute, so their exponentials add:
\[
 e^{(\delta_1/2)\partial_x^2}
 e^{(\delta_2/2)\partial_x^2}
 =e^{((\delta_1+\delta_2)/2)\partial_x^2}.
\]
Equivalently, substitute \eqref{eq:v5v46-Tdelta} twice and use
\[
 \sum_{j=0}^k\frac{\delta_1^j\delta_2^{k-j}}{j!(k-j)!}
 =\frac{(\delta_1+\delta_2)^k}{k!}.
\]
\end{proof}

\subsection{Sharp obstruction on a single factorial radius}

For a scalar coefficient sequence, set
\begin{equation}
 \|u\|_{\rho,\alpha}
 =\sum_{q\ge0}\frac{\rho^q}{(q!)^\alpha}|u_q|.
 \label{eq:v5v46-scalar-norm}
\end{equation}

\begin{theorem}[Fixed-radius connection dichotomy]
\label{thm:v5v46-fixed-radius}
Let \(\rho>0\) and \(\delta\ne0\).
\begin{enumerate}
 \item If \(\alpha=0\), then
 \begin{equation}
  \|T_\delta\|_{\rho,0\to\rho,0}
  \le
  \exp\!\left(\frac{|\delta|}{2\rho^2}\right).
  \label{eq:v5v46-analytic-connection}
 \end{equation}
 \item If \(\alpha>0\), then \(T_\delta\) is unbounded on
 \((\rho,\alpha)\).
\end{enumerate}
\end{theorem}

\begin{proof}
For an input supported in degree \(q\), the weighted contribution of the
output \(p=q-2k\) is
\begin{equation}
 R_{p,k}(\rho,\alpha)
 =
 \rho^{-2k}
 \left(\frac{(p+2k)!}{p!}\right)^\alpha
 \frac{(|\delta|/2)^k}{k!}.
 \label{eq:v5v46-column-ratio}
\end{equation}
When \(\alpha=0\), summing over \(k\) gives at most the exponential in
\eqref{eq:v5v46-analytic-connection}.  When \(\alpha>0\), retain only
\(k=1\).  Then
\[
 R_{p,1}
 =\frac{|\delta|}{2\rho^2}
 ((p+1)(p+2))^\alpha\longrightarrow\infty.
\]
The one-column test proves unboundedness.
\end{proof}

\begin{firewall}[Small covariance change does not cure a fixed radius]
\label{fw:v5v46-small-delta}
For every nonzero \(\delta\), however small, the polynomial growth in the
last line of the proof is unbounded as \(p\to\infty\).  Small shell width
alone therefore does not make adjacent moving frames bounded on a fixed
positive-exponent V25 grade norm.
\end{firewall}

\subsection{A strict radius gap and the one-half threshold}

Let \(0<\rho_-<\rho_+\), put
\begin{equation}
 r=\frac{\rho_-}{\rho_+},
 \qquad
 A=\frac{|\delta|}{2\rho_+^2}.
 \label{eq:v5v46-rA}
\end{equation}
The weighted contribution of the block \(q=p+2k\), now from the strong
radius to the weak radius, is
\begin{equation}
 R_{p,k}^{-,+}
 =
 r^p A^k
 \frac1{k!}
 \left(\frac{(p+2k)!}{p!}\right)^\alpha.
 \label{eq:v5v46-two-radius-ratio}
\end{equation}

\begin{theorem}[Two-radius Wick-connection theorem]
\label{thm:v5v46-two-radius}
The following statements hold.
\begin{enumerate}
 \item If \(0\le\alpha<1/2\), then
 \begin{equation}
  T_\delta:(\rho_+,\alpha)\longrightarrow(\rho_-,\alpha)
  \label{eq:v5v46-subcritical-map}
 \end{equation}
 is bounded for every finite \(\delta\).
 \item If \(\alpha=1/2\), the same map is bounded under the sufficient
 conditions
 \begin{equation}
  re^{2A}<1,
  \qquad 4eA<1.
  \label{eq:v5v46-critical-smallness}
 \end{equation}
 \item If \(\alpha>1/2\) and \(\delta\ne0\), the map is unbounded for
 every pair \(\rho_->0<\rho_+\).
\end{enumerate}
\end{theorem}

\begin{proof}
It is enough to bound the double nonnegative sum of
\eqref{eq:v5v46-two-radius-ratio}, since that sum dominates every input
column.  Split the pairs \((p,k)\) into \(p\ge2k\) and \(p<2k\).

If \(p\ge2k\), then
\[
 \frac{(p+2k)!}{p!}\le(p+2k)^{2k}\le(2p)^{2k}.
\]
For each \(p\), summing all \(k\le p/2\) is bounded by
\begin{equation}
 r^p\exp\!\left(A(2p)^{2\alpha}\right).
 \label{eq:v5v46-first-region}
\end{equation}
The sum over \(p\) converges when \(2\alpha<1\).  At
\(\alpha=1/2\), it converges if \(re^{2A}<1\).

If \(p<2k\), then \(p+2k<4k\), so summing first over \(p\) gives the
bound
\begin{equation}
 \frac1{1-r}
 \sum_{k\ge1}
 \frac{A^k(4k)^{2\alpha k}}{k!}.
 \label{eq:v5v46-second-region}
\end{equation}
Stirling's lower bound \(k!\ge(k/e)^k\) shows convergence when
\(2\alpha-1<0\).  At \(\alpha=1/2\), the series is dominated by a
geometric series with ratio \(4eA\).  This proves the first two claims.

For \(\alpha>1/2\), fix any output \(p\) in the tail and let
\(k\to\infty\).  Stirling's formula in
\eqref{eq:v5v46-two-radius-ratio} gives a factor of the form
\[
 k^{(2\alpha-1)k}C^k
\]
with \(C>0\).  It diverges, so even one block in each of those columns is
unbounded.
\end{proof}

\begin{assessment}[Meaning of the threshold]
\label{ass:v5v46-threshold}
The factorial \(1/k!\) in the normalized heat connection pays one
factorial for every two removed fields.  The V25 weight demands
approximately \(((2k)!)^\alpha\) when a high input lands near a fixed
tail output.  Stirling balance occurs at \(2\alpha=1\).  This is why the
threshold is one half rather than one.
\end{assessment}

\subsection{Cofinal frame transport does not spend one radius gap per shell}

The V43 two-radius firewall concerned repeated Neumann powers of a general
operator.  Wick-frame connections have the additional group law
\eqref{eq:v5v46-group-law}, so consecutive covariance changes telescope
before a norm is taken.

\begin{theorem}[Telescoping frame connection]
\label{thm:v5v46-telescoping}
Let \(c_j\) be a sequence of scalar covariance parameters and set
\(\delta_j=c_{j+1}-c_j\).  For every \(N\),
\begin{equation}
 T_{\delta_{N-1}}\cdots T_{\delta_0}
 =T_{c_N-c_0}.
 \label{eq:v5v46-finite-telescope}
\end{equation}
If \(c_N\to c_\infty\), then on finite sequences the left side converges
coefficientwise to \(T_{c_\infty-c_0}\).  Moreover:
\begin{enumerate}
 \item at \(\alpha=0\), it extends on one fixed radius with norm at most
 \begin{equation}
  \exp\!\left(\frac{|c_\infty-c_0|}{2\rho^2}\right);
  \label{eq:v5v46-telescope-analytic}
 \end{equation}
 \item for \(0<\alpha<1/2\), it extends across one chosen strict radius
 gap by Theorem~\ref{thm:v5v46-two-radius}.
\end{enumerate}
Thus the exact Gaussian frame transport consumes at most one global radius
gap, not one gap per shell.
\end{theorem}

\begin{proof}
Iterate Lemma~\ref{lem:v5v46-group-law}; the parameter sum telescopes to
\(c_N-c_0\).  Coefficientwise convergence on finite sequences is
immediate.  Apply Theorem~\ref{thm:v5v46-fixed-radius} or
Theorem~\ref{thm:v5v46-two-radius} to the single limiting connection.
The same majorants dominate partial connections once \(c_N-c_0\) stays in
a compact interval, giving the continuous extension.
\end{proof}

\begin{firewall}[Field-dependent connections need not commute]
\label{fw:v5v46-noncommuting}
The telescope uses constant covariance differences.  If the transported
Hessian depends on the low field or if gauge projection changes with the
shell, the connection contains variable-coefficient derivatives and
commutators.  Lemma~\ref{lem:v5v46-group-law} then fails in this form.  A
path-ordered connection and a curvature estimate are required.
\end{firewall}

\subsection{Separating field degree from Gevrey provenance}

The V6--V7 factorial exponent was introduced to absorb graph and forest
growth, while the connection above acts on field number.  Using one grade
for both purposes imposes the incompatible requirements
\(\alpha\) large enough for graph growth and \(\alpha\le1/2\) for frame
transport.  Define separate indices:
\begin{itemize}
 \item \(a\): counterterm, forest, or operator-provenance grade;
 \item \(f\): normalized Wick field degree;
 \item \(m\): particle or highest-chaos number.
\end{itemize}
For blocks \(X_{a,f,m}\), define
\begin{equation}
 \|u\|_{R,\rho,\eta;s,b}^{\sigma,\alpha,\beta}
 =
 \sum_{a,f,m\ge0}
 \frac{R^a}{(a!)^\sigma}
 \frac{\rho^f}{(f!)^\alpha}
 \frac{\eta^m}{(m!)^\beta}
 \|u_{a,f,m}\|_{a,f,m;s,b}.
 \label{eq:v5v46-split-norm}
\end{equation}

\begin{theorem}[Split tri-factorial Banach algebra]
\label{thm:v5v46-split-algebra}
Suppose the physical product maps
\begin{equation}
 X_{a,f,m}\times X_{a',f',m'}
 \longrightarrow
 X_{a+a',f+f',m+m'}
 \label{eq:v5v46-additive-product}
\end{equation}
with a block norm bounded by a uniform constant \(C_\star\), including the
spatial and energy weights.  Then for all
\(\sigma,\alpha,\beta\ge0\),
\begin{equation}
 \|u\star v\|_{R,\rho,\eta;s,b}^{\sigma,\alpha,\beta}
 \le C_\star
 \|u\|_{R,\rho,\eta;s,b}^{\sigma,\alpha,\beta}
 \|v\|_{R,\rho,\eta;s,b}^{\sigma,\alpha,\beta}.
 \label{eq:v5v46-split-algebra-bound}
\end{equation}
\end{theorem}

\begin{proof}
Use
\[
 (a+a')!\ge a!a'!,\qquad
 (f+f')!\ge f!f'!,\qquad
 (m+m')!\ge m!m'!.
\]
The output weight is at most the product of the two input weights.  Apply
the block estimate and sum the resulting nonnegative Cauchy product by
Tonelli's theorem, exactly as in V25.
\end{proof}

\begin{theorem}[Frame connection on the split norm]
\label{thm:v5v46-split-connection}
Assume \(T_\delta\) acts only on the field index \(f\), is the identity on
\((a,m)\), and is contractive on the remaining block norm.  Then:
\begin{enumerate}
 \item for \(\alpha=0\), it is bounded on the same split radii with the
 factor \eqref{eq:v5v46-analytic-connection};
 \item for \(0<\alpha<1/2\), it is bounded from \(\rho_+\) to
 \(\rho_-\), leaving \(R,\eta,\sigma,\beta,s,b\) unchanged;
 \item the provenance exponent \(\sigma\) may be chosen independently to
 dominate forest growth.
\end{enumerate}
\end{theorem}

\begin{proof}
For every fixed \((a,m)\), apply
Theorem~\ref{thm:v5v46-fixed-radius} or
Theorem~\ref{thm:v5v46-two-radius} to the \(f\)-sequence.  Its constant is
independent of \((a,m)\), so summation with the untouched weights proves
the first two claims.  The third follows because \(T_\delta\) does not
change \(a\).
\end{proof}

\begin{firewall}[The split must descend to the renormalized quotient]
\label{fw:v5v46-split-quotient}
Curvature commutators, equations of motion, integration by parts, BRST
relations, and local Taylor subtraction can exchange displayed field
number with derivative and curvature factors.  Therefore
\eqref{eq:v5v46-additive-product} and invariance of \(a\) under
\(T_\delta\) must be proved after forming the renormalized operator
quotient.  A formal assignment of two labels is not an MFTC.
\end{firewall}

\begin{definition}[Split moving-frame certificate]
\label{def:v5v46-SMFC}
An SMFC consists of:
\begin{enumerate}
 \item a renormalized operator quotient carrying the three independent
 grades \((a,f,m)\) and the norm \eqref{eq:v5v46-split-norm};
 \item normalized Wick field tensors and an exact adjacent-shell
 connection;
 \item \(\alpha<1/2\), or the critical smallness
 \eqref{eq:v5v46-critical-smallness}, for the field grade;
 \item a provenance exponent \(\sigma\) large enough for the proved
 forest majorant;
 \item either the telescoping identity
 \eqref{eq:v5v46-finite-telescope} or a summable curvature bound for the
 path-ordered connection;
 \item complete OPNR estimates for the non-Gaussian remainder and
 compatibility with the CTIC diagonal.
\end{enumerate}
\end{definition}

\subsection{V46 status}

\begin{theorem}[V46 conclusion]
\label{thm:v5v46-status}
The covariance connection between normalized Wick frames is bounded on a
fixed analytic field-degree norm, is bounded across a radius gap below
factorial exponent one half, and is unbounded across every radius gap
above one half.  Constant Gaussian frame changes telescope exactly across
shells.  Separating field degree from graph provenance removes the direct
conflict with a larger forest Gevrey exponent and preserves the Banach
algebra property.  What remains unproved is that the complete
renormalized SM--Einstein operator quotient carries this split and that
the field-dependent frame connection has summable curvature.
\end{theorem}

\begin{proof}
Use Theorems~\ref{thm:v5v46-fixed-radius},
\ref{thm:v5v46-two-radius},
\ref{thm:v5v46-telescoping},
\ref{thm:v5v46-split-algebra}, and
\ref{thm:v5v46-split-connection}.
\end{proof}

V47 will attack the variable-frame obstruction.  It will write the
connection generated by a low-field-dependent covariance, compute its
curvature commutator, and identify a sufficient flatness or summable
curvature condition under which the cofinal frame transport remains
bounded without spending an independent radius gap at every shell.

\clearpage
\section{V47: a normalization-independent product--connection obstruction}
\label{sec:v5v47}

V46 proved that the normalized Wick-frame connection is well behaved only
for a field-degree factorial exponent at or below one half.  Its split
Banach-algebra theorem was conditional on a uniform block product.  This
note checks that condition in the same Wick basis.  The normalization which
improves the covariance connection transfers factorials into the highest
degree part of ordinary multiplication.  Varying the normalization cannot
remove the conflict: in the full one-parameter family of Wick
normalizations, same-space bounded multiplication forces a parameter range
in which every nontrivial covariance connection is unbounded, even across
a strict radius gap.

This is a genuine structural obstruction, rather than a failed estimate.
It shows that the moving-frame programme cannot close in a single V25-type
factorial Banach algebra.  A scale of spaces with an explicit loss, a
nuclear test-function algebra, or a different exact product is required.

\subsection{The full normalization family}

For \(\theta\in\mathbb R\), define
\begin{equation}
 e_n^{(c,\theta)}(x)
 =\frac{H_n(x;c)}{(n!)^\theta}.
 \label{eq:v5v47-theta-basis}
\end{equation}
The V46 normalization is \(\theta=1\), while the unnormalized Wick basis
is \(\theta=0\).  Coefficients are measured with
\begin{equation}
 \|u\|_{\rho,\alpha}
 =\sum_{n\ge0}\frac{\rho^n}{(n!)^\alpha}|u_n|.
 \label{eq:v5v47-weight}
\end{equation}

\begin{lemma}[General covariance-connection coefficient]
\label{lem:v5v47-general-connection}
Let \(\delta=c_2-c_1\).  Then
\begin{equation}
 e_q^{(c_1,\theta)}
 =
 \sum_{k=0}^{\lfloor q/2\rfloor}
 \left(\frac{q!}{(q-2k)!}\right)^{1-\theta}
 \frac{(\delta/2)^k}{k!}
 e_{q-2k}^{(c_2,\theta)}.
 \label{eq:v5v47-general-connection}
\end{equation}
\end{lemma}

\begin{proof}
The unnormalized connection is
\[
 H_q(x;c_1)
 =\sum_{k=0}^{\lfloor q/2\rfloor}
 \frac{q!}{(q-2k)!k!}
 \left(\frac\delta2\right)^k
 H_{q-2k}(x;c_2).
\]
Divide by \((q!)^\theta\) and substitute
\(H_{q-2k}=((q-2k)!)^\theta e_{q-2k}^{(c_2,\theta)}\).
\end{proof}

Put
\begin{equation}
 \gamma=\alpha+1-\theta.
 \label{eq:v5v47-gamma}
\end{equation}
For a block \(q=p+2k\), the weighted strong-to-weak connection ratio is
\begin{equation}
 r^p A^k\frac1{k!}
 \left(\frac{(p+2k)!}{p!}\right)^\gamma,
 \qquad
 r=\frac{\rho_-}{\rho_+},
 \quad
 A=\frac{|\delta|}{2\rho_+^2}.
 \label{eq:v5v47-general-ratio}
\end{equation}

\begin{theorem}[General connection threshold]
\label{thm:v5v47-general-threshold}
Assume \(\delta\ne0\) and \(0<\rho_-<\rho_+\).
\begin{enumerate}
 \item If \(0\le\gamma<1/2\), the covariance connection is bounded from
 \((\rho_+,\alpha)\) to \((\rho_-,\alpha)\).
 \item If \(\gamma=1/2\), it is bounded under the sufficient conditions
 \eqref{eq:v5v46-critical-smallness} with the \(A,r\) in
 \eqref{eq:v5v47-general-ratio}.
 \item If \(\gamma>1/2\), it is unbounded across every strict radius gap.
\end{enumerate}
\end{theorem}

\begin{proof}
Equation~\eqref{eq:v5v47-general-ratio} is exactly the ratio
\eqref{eq:v5v46-two-radius-ratio} with \(\alpha\) replaced by
\(\gamma\).  The two-region proof of
Theorem~\ref{thm:v5v46-two-radius} applies verbatim.  For the last claim,
fix \(p\) and let \(k\to\infty\); Stirling's formula gives the divergent
factor \(k^{(2\gamma-1)k}C^k\).
\end{proof}

\subsection{The highest-degree product packet}

The Wick product formula begins with
\begin{equation}
 H_p(x;c)H_q(x;c)
 =H_{p+q}(x;c)+\text{terms of degree at most }p+q-2.
 \label{eq:v5v47-Wick-product-top}
\end{equation}
Therefore the coefficient of the top output in the \(\theta\)-basis is
\begin{equation}
 e_p^{(c,\theta)}e_q^{(c,\theta)}
 =\binom{p+q}{p}^{\theta}
 e_{p+q}^{(c,\theta)}+\text{lower degrees}.
 \label{eq:v5v47-top-coefficient}
\end{equation}

\begin{lemma}[Necessary same-space product condition]
\label{lem:v5v47-product-necessary}
If multiplication extends to a bounded bilinear map
\begin{equation}
 (\rho,\alpha)\times(\rho,\alpha)
 \longrightarrow(\rho,\alpha)
 \label{eq:v5v47-product-map}
\end{equation}
in a direct grade norm, then
\begin{equation}
 \sup_{p,q}\binom{p+q}{p}^{\theta-\alpha}<\infty.
 \label{eq:v5v47-binomial-condition}
\end{equation}
Consequently, a necessary condition is
\begin{equation}
 \theta\le\alpha.
 \label{eq:v5v47-theta-alpha}
\end{equation}
\end{lemma}

\begin{proof}
Take inputs supported only in grades \(p\) and \(q\).  The ratio of the
weighted norm of the highest output in
\eqref{eq:v5v47-top-coefficient} to the two input norms is
\begin{align*}
 &\frac{\rho^{p+q}}{((p+q)!)^\alpha}
 \binom{p+q}{p}^{\theta}
 \frac{(p!)^\alpha(q!)^\alpha}{\rho^{p+q}}
 =\binom{p+q}{p}^{\theta-\alpha}.
\end{align*}
Lower degrees lie in different direct-sum blocks and cannot cancel this
positive norm contribution.  If \(\theta>\alpha\), take \(p=q\to\infty\);
the central binomial coefficient diverges exponentially.
\end{proof}

\begin{theorem}[Single-factorial-space no-go theorem]
\label{thm:v5v47-no-go}
Within the basis family \eqref{eq:v5v47-theta-basis} and weights
\eqref{eq:v5v47-weight}, no choice of \((\theta,\alpha)\) can make both
of the following true:
\begin{enumerate}
 \item ordinary multiplication is a bounded bilinear endomorphism;
 \item a nonzero covariance connection is bounded even from one strict
 radius to a smaller radius.
\end{enumerate}
\end{theorem}

\begin{proof}
Item 1 requires \(\theta\le\alpha\) by
Lemma~\ref{lem:v5v47-product-necessary}.  Then
\[
 \gamma=\alpha+1-\theta\ge1.
\]
In particular \(\gamma>1/2\), so item 2 is impossible by
Theorem~\ref{thm:v5v47-general-threshold}.
\end{proof}

\begin{corollary}[The V46 conditional algebra does not hold for the
ordinary normalized Wick product]
\label{cor:v5v47-v46-scope}
For \(\theta=1\), the connection-favourable range in V46 is
\(\alpha<1/2\), whereas boundedness of the highest ordinary product packet
requires \(\alpha\ge1\).  Hence the uniform product hypothesis
\eqref{eq:v5v46-additive-product} fails in that natural realization of
the field grade.
\end{corollary}

\begin{proof}
Insert \(\theta=1\) into
Theorem~\ref{thm:v5v47-general-threshold} and
Lemma~\ref{lem:v5v47-product-necessary}.
\end{proof}

\begin{firewall}[Changing basis only moves the factorial]
\label{fw:v5v47-basis-no-cure}
Choosing \(\theta\) larger improves the covariance connection because it
adds denominator factorials to contractions.  The same choice enlarges
the binomial coefficient in the top no-contraction product.  Theorem
\ref{thm:v5v47-no-go} shows that this trade is invariant throughout the
entire power-factorial normalization family.
\end{firewall}

\subsection{What a radius gap can and cannot do for products}

Although no same-space solution exists, a strict radius gap can absorb the
highest product binomial.

\begin{proposition}[Highest-product radius estimate]
\label{prop:v5v47-product-radius}
Let two inputs have radius \(\rho_+\) and measure the output at
\(\rho_-<\rho_+\).  The highest packet in
\eqref{eq:v5v47-top-coefficient} is uniformly bounded if either
\(\theta\le\alpha\), or
\begin{equation}
 \frac{\rho_-}{\rho_+},2^{\theta-\alpha}\le1.
 \label{eq:v5v47-product-radius-condition}
\end{equation}
\end{proposition}

\begin{proof}
The weighted ratio is
\[
 \left(\frac{\rho_-}{\rho_+}\right)^{p+q}
 \binom{p+q}{p}^{\theta-\alpha}.
\]
If \(\theta\le\alpha\), it is at most one.  Otherwise use
\(\binom{p+q}{p}\le2^{p+q}\) and
\eqref{eq:v5v47-product-radius-condition}.
\end{proof}

\begin{firewall}[The highest packet is only a necessary product test]
\label{fw:v5v47-lower-products}
Proposition~\ref{prop:v5v47-product-radius} controls the no-contraction
output only.  The complete Wick product contains all lower contractions,
and the interacting SM product contains local subtractions and Ward
partners.  A full product theorem must assemble and bound those terms as
well.
\end{firewall}

The combination of Theorem~\ref{thm:v5v47-general-threshold} and
Proposition~\ref{prop:v5v47-product-radius} explains why a scale of spaces
is unavoidable.  A connection-favourable normalization has
\(\theta>\alpha+1/2\), while multiplication tends to move toward a weaker
radius or a larger admissible factorial exponent.  Repeated alternation
must be controlled by a smoothing or tame theorem; an ordinary Banach
fixed point would spend its regularity budget indefinitely.

\subsection{A typed scale formulation}

Let \(\mathcal X_{\rho,\alpha}^{\theta}\) denote the completion in the
\(\theta\)-basis and weight \((\rho,\alpha)\).  A proposed construction
must type each operation, rather than calling all of them bounded on one
space:
\begin{align}
 T_\delta&:
 \mathcal X_{\rho_+,\alpha}^{\theta}
 \longrightarrow
 \mathcal X_{\rho_-,\alpha}^{\theta},
 &&\alpha+1-\theta\le\frac12,
 \label{eq:v5v47-typed-connection}\\
 \mathcal M&:
 \mathcal X_{\rho_+,alpha}^{\theta}
 \times\mathcal X_{\rho_+,alpha}^{\theta}
 \longrightarrow
 \mathcal X_{\rho_-,\alpha}^{\theta},
 &&\text{with at least \eqref{eq:v5v47-product-radius-condition}.}
 \label{eq:v5v47-typed-product}
\end{align}

\begin{definition}[Field-scale compatibility certificate]
\label{def:v5v47-FSCC}
An FSCC consists of:
\begin{enumerate}
 \item a nested field-function scale containing every transported Wick
 frame used by the shell system;
 \item complete, not merely highest-packet, estimates for multiplication,
 covariance connection, local subtraction, and conditional cumulants with
 their exact source and target indices;
 \item smoothing projections with quantitative approximation estimates;
 \item a Newton or Nash--Moser theorem whose regularity losses match those
 operation types;
 \item a cofinal radius or scale parameter bounded away from its degenerate
 endpoint;
 \item compatibility with the independent provenance and particle grades.
\end{enumerate}
\end{definition}

\begin{proposition}[Why a formal projective intersection is insufficient]
\label{prop:v5v47-projective-not-enough}
Taking the intersection of all radii does not by itself solve the no-go
theorem.  To obtain a continuous nonlinear shell map, for every target
seminorm one must exhibit source seminorms for which all operations in the
same formula are bounded, and these choices must satisfy the hypotheses of
the chosen inverse theorem.
\end{proposition}

\begin{proof}
Continuity in a projective-limit topology means that each target seminorm
is bounded by finitely many source seminorms.  Theorems
\ref{thm:v5v47-general-threshold} and
\ref{thm:v5v47-no-go} show that connection and multiplication require
different source-target assignments.  Merely belonging to every
individual space supplies no uniform assignment and no convergence of a
Newton iteration.
\end{proof}

\subsection{V47 status}

\begin{theorem}[V47 conclusion]
\label{thm:v5v47-status}
For every power-factorial normalization of Wick tensors, bounded ordinary
multiplication on one factorial grade space and bounded transport through
a nonzero covariance change are incompatible.  The conflict is already
visible in the highest no-contraction product and one exact covariance
connection column, so lower diagrams cannot repair it by cancellation
between grades.  The moving-frame route must therefore be formulated on a
typed scale and paired with a genuine loss-tolerant inverse theorem.  No
FSCC has yet been proved for the SM--Einstein operator system.
\end{theorem}

\begin{proof}
The statement is Theorem~\ref{thm:v5v47-no-go}, with the scale consequence
given by Proposition~\ref{prop:v5v47-projective-not-enough}.
\end{proof}

V48 will construct the simplest such scale for entire functions of
Gaussian type.  It will test both multiplication and the heat connection
in a supremum norm with an explicit type parameter, then determine whether
finite interaction valence and a cofinal type budget can replace the
failed single-space Banach algebra.

\clearpage
\section{V48: Gaussian-type entire functions and the shell type budget}
\label{sec:v5v48}

V47 proved that a single factorial coefficient norm cannot support both
ordinary multiplication and covariance-frame transport.  This note replaces
the factorial exponent by a continuous Gaussian type.  In that scale,
multiplication adds types, field dilation multiplies type by the square of
the dilation, and the heat connection changes type by an explicit fractional
linear rule.  These three exact formulas produce a closed type budget for a
finite-degree shell operation.

The same calculation also finds the boundary of this repair.  A
term-by-term cumulant expansion has unbounded multiplication order, and no
fixed positive Gaussian type survives all orders under a fixed nonzero
field dilation.  Moreover, the exponential of a polynomial interaction of
degree greater than two generally has growth above Gaussian type in complex
directions.  Therefore the Gaussian-type scale is a rigorous local model
for a finite-valence numerator, but it is not yet a space for the full
Boltzmann shell map.

\subsection{The Gaussian-type scale}

For \(\tau>0\), let \(\mathcal G_\tau\) be the Banach space of entire
functions \(f:\mathbb C\to\mathbb C\) such that
\begin{equation}
 \|f\|_{\mathcal G_\tau}
 =
 \sup_{z\in\mathbb C}
 |f(z)|e^{-\tau|z|^2}<\infty.
 \label{eq:v5v48-Gtau}
\end{equation}
If \(0<\tau_1\le\tau_2\), then
\(\mathcal G_{\tau_1}\subset\mathcal G_{\tau_2}\) contractively.  Larger
type is a weaker space.

\begin{theorem}[Exact product typing]
\label{thm:v5v48-product}
For \(\tau_1,\tau_2>0\), pointwise multiplication is bounded as
\begin{equation}
 \mathcal G_{\tau_1}\times\mathcal G_{\tau_2}
 \longrightarrow\mathcal G_{\tau_1+\tau_2},
 \qquad
 \|fg\|_{\mathcal G_{\tau_1+\tau_2}}
 \le
 \|f\|_{\mathcal G_{\tau_1}}
 \|g\|_{\mathcal G_{\tau_2}}.
 \label{eq:v5v48-product}
\end{equation}
Consequently,
\begin{equation}
 \|f^r\|_{\mathcal G_{r\tau}}
 \le\|f\|_{\mathcal G_\tau}^r.
 \label{eq:v5v48-power}
\end{equation}
\end{theorem}

\begin{proof}
Multiply the two pointwise estimates
\[
 |f(z)|\le\|f\|_{\mathcal G_{\tau_1}}e^{\tau_1|z|^2},
 \qquad
 |g(z)|\le\|g\|_{\mathcal G_{\tau_2}}e^{\tau_2|z|^2}.
\]
This proves \eqref{eq:v5v48-product}; iteration proves
\eqref{eq:v5v48-power}.
\end{proof}

\begin{theorem}[Exact dilation typing]
\label{thm:v5v48-dilation}
For \(a\in\mathbb C\), let \(S_af(z)=f(az)\).  Then
\begin{equation}
 S_a:\mathcal G_\tau\longrightarrow\mathcal G_{|a|^2\tau},
 \qquad
 \|S_af\|_{\mathcal G_{|a|^2\tau}}
 \le\|f\|_{\mathcal G_\tau}.
 \label{eq:v5v48-dilation}
\end{equation}
\end{theorem}

\begin{proof}
Use
\[
 |f(az)|\le\|f\|_{\mathcal G_\tau}
 e^{\tau|az|^2}.
\]
\end{proof}

\subsection{Heat connection on Gaussian type}

For real \(\delta\), define on polynomials
\begin{equation}
 H_\delta
 =\exp\!\left(\frac\delta2\partial_z^2\right).
 \label{eq:v5v48-heat}
\end{equation}
This is the function-side form of the Wick-frame connection in V46.

\begin{theorem}[Gaussian-type heat estimate]
\label{thm:v5v48-heat-bound}
Let \(\tau>0\) and suppose
\begin{equation}
 2\tau|\delta|<1.
 \label{eq:v5v48-heat-smallness}
\end{equation}
Then \(H_\delta\) extends to a bounded map
\begin{equation}
 H_\delta:\mathcal G_\tau
 \longrightarrow
 \mathcal G_{\Phi_\delta(\tau)},
 \qquad
 \Phi_\delta(\tau)
 =
 \frac{\tau}{1-2\tau|\delta|},
 \label{eq:v5v48-heat-typing}
\end{equation}
with
\begin{equation}
 \|H_\delta f\|_{\mathcal G_{\Phi_\delta(\tau)}}
 \le
 (1-2\tau|\delta|)^{-1/2}
 \|f\|_{\mathcal G_\tau}.
 \label{eq:v5v48-heat-norm}
\end{equation}
\end{theorem}

\begin{proof}
Let \(Z\) be a standard real Gaussian.  Choose
\(\sigma\in\mathbb C\) with \(\sigma^2=\delta\) and
\(|\sigma|^2=|\delta|\).  On polynomials,
\[
 H_\delta f(z)=\mathbb E f(z+\sigma Z),
\]
by expansion and the Gaussian moments.  For \(f\in\mathcal G_\tau\),
\[
 |H_\delta f(z)|
 \le
 \|f\|_{\mathcal G_\tau}
 \mathbb E e^{\tau|z+\sigma Z|^2}.
\]
The inequality
\[
 |z+\sigma Z|^2
 =
 |z|^2+2\operatorname{Re}(\overline z\sigma)Z
 +|\sigma|^2Z^2
\]
and the exact Gaussian integral give
\begin{align*}
 \mathbb E e^{\tau|z+\sigma Z|^2}
 &=
 (1-2\tau|\sigma|^2)^{-1/2}
 \exp\!\left(
 \tau|z|^2+
 \frac{2\tau^2\operatorname{Re}(\overline z\sigma)^2}
 {1-2\tau|\sigma|^2}
 \right)\\
 &\le
 (1-2\tau|\delta|)^{-1/2}
 \exp\!\left(
 \frac{\tau}{1-2\tau|\delta|}|z|^2
 \right).
\end{align*}
This proves the asserted bound first for polynomials.  Polynomial
approximants obtained from the Taylor series converge locally uniformly;
the displayed majorant gives the extension to \(\mathcal G_\tau\).
\end{proof}

\begin{corollary}[Heat connection is a type loss]
\label{cor:v5v48-heat-loss}
For \(\delta\ne0\),
\begin{equation}
 \Phi_\delta(\tau)>\tau.
 \label{eq:v5v48-strict-heat-loss}
\end{equation}
Thus the absolute Gaussian-type estimate treats both signs of the
covariance change as a loss of type.
\end{corollary}

\begin{proof}
The denominator in \(\Phi_\delta(\tau)\) lies strictly between zero and
one.
\end{proof}

\begin{assessment}[Sign information is discarded]
\label{ass:v5v48-sign}
For special functions and a favourable sign, a backward or forward heat
operator can reduce Gaussian growth.  Theorem~\ref{thm:v5v48-heat-bound}
uses absolute values and a complex Gaussian shift, so its uniform bound
depends on \(|\delta|\).  Any improvement using the sign must be proved
inside the complete physical numerator before this absolute estimate.
\end{assessment}

\subsection{A closed finite-valence type budget}

Consider an operation which first dilates every input by \(a\), multiplies
\(r\) inputs, and then applies the frame connection:
\begin{equation}
 \mathcal Q_{r,a,\delta}(f_1,\ldots,f_r)
 =
 H_\delta\prod_{j=1}^r S_af_j.
 \label{eq:v5v48-Q}
\end{equation}

\begin{theorem}[Finite-valence invariant-type criterion]
\label{thm:v5v48-type-budget}
Let every \(f_j\in\mathcal G_\tau\).  If
\begin{equation}
 2r|a|^2\tau|\delta|<1,
 \qquad
 r|a|^2(1+2\tau|\delta|)\le1,
 \label{eq:v5v48-budget-condition}
\end{equation}
then
\begin{equation}
 \mathcal Q_{r,a,\delta}:
 (\mathcal G_\tau)^r\longrightarrow\mathcal G_\tau
 \label{eq:v5v48-Q-endomorphism}
\end{equation}
and
\begin{equation}
 \|\mathcal Q_{r,a,\delta}(f_1,\ldots,f_r)\|_{\mathcal G_\tau}
 \le
 (1-2r|a|^2\tau|\delta|)^{-1/2}
 \prod_{j=1}^r\|f_j\|_{\mathcal G_\tau}.
 \label{eq:v5v48-Q-bound}
\end{equation}
\end{theorem}

\begin{proof}
Theorems~\ref{thm:v5v48-product} and
\ref{thm:v5v48-dilation} place the product before heat transport in
\(\mathcal G_{\tau_0}\), where
\[
 \tau_0=r|a|^2\tau.
\]
The heat theorem changes this to
\[
 \Phi_\delta(\tau_0)
 =
 \frac{r|a|^2\tau}
 {1-2r|a|^2\tau|\delta|}.
\]
The second inequality in \eqref{eq:v5v48-budget-condition} is equivalent
to \(\Phi_\delta(\tau_0)\le\tau\).  The first permits application of the
heat theorem.  The norm estimate follows from its prefactor and the
contractive inclusion
\(\mathcal G_{\Phi_\delta(\tau_0)}\subset\mathcal G_\tau\).
\end{proof}

\begin{corollary}[Zero-connection budget]
\label{cor:v5v48-zero-connection}
When \(\delta=0\), an \(r\)-linear shell numerator preserves
\(\mathcal G_\tau\) whenever
\begin{equation}
 r|a|^2\le1.
 \label{eq:v5v48-zero-budget}
\end{equation}
\end{corollary}

\begin{proof}
Set \(\delta=0\) in
Theorem~\ref{thm:v5v48-type-budget}.
\end{proof}

\begin{definition}[Gaussian type budget certificate]
\label{def:v5v48-GTBC}
A GTBC for a finite-valence shell numerator consists of:
\begin{enumerate}
 \item a Gaussian-type field chart with a proved type \(\tau\);
 \item a shell field dilation \(a\) and covariance connection \(\delta\);
 \item a maximal numerator valence \(r_{\max}<\infty\);
 \item condition \eqref{eq:v5v48-budget-condition} for every
 \(1\le r\le r_{\max}\);
 \item compatible estimates for the independent provenance, particle,
 Ward, and multiplicity blocks.
\end{enumerate}
\end{definition}

\begin{firewall}[The dilation must come from the physical shell map]
\label{fw:v5v48-dilation-physical}
The factor \(a\) cannot be chosen to make
\eqref{eq:v5v48-budget-condition} true.  It is fixed by canonical field
rescaling, wave-function renormalization, and the moving-frame affine map.
All relevant singular values must satisfy the multidimensional analogue
of the type budget.
\end{firewall}

\subsection{Why unbounded cumulant order defeats a fixed type}

The shell map
\(-\log\mathbb E e^{-V}\) has connected terms of every order in \(V\),
even when each primitive interaction vertex has finite field valence.
The following scalar test shows why a termwise Gaussian-type estimate
cannot close those orders in one type.

\begin{proposition}[No fixed type for unbounded product order]
\label{prop:v5v48-unbounded-order}
Fix \(\tau>0\) and \(a\ne0\).  There is no finite type
\(\tau_*\) such that
\begin{equation}
 (S_af)^r\in\mathcal G_{\tau_*}
 \quad\text{with a uniform type assignment for every }r\ge1
 \label{eq:v5v48-false-all-order}
\end{equation}
for all \(f\in\mathcal G_\tau\).
\end{proposition}

\begin{proof}
Take \(f(z)=e^{\tau z^2}\), which has
\(\|f\|_{\mathcal G_\tau}\le1\).  Along the real axis,
\[
 |(S_af)(x)|^r
 =
 e^{r\tau\operatorname{Re}(a^2)x^2}
\]
when \(a^2\) is positive real.  For general nonzero \(a\), rotate the test
Gaussian so that its maximal-growth ray aligns with \(a\).  Its minimal
admissible Gaussian type is \(r|a|^2\tau\), which tends to infinity.
Thus no fixed \(\tau_*\) contains every product order.
\end{proof}

\begin{corollary}[Termwise connected expansion does not prove a GTBC]
\label{cor:v5v48-cumulant-not-GTBC}
An all-order connected-cumulant expansion estimated by taking absolute
values of each \(r\)-fold product cannot prove that the full shell map is
an endomorphism of a fixed \(\mathcal G_\tau\) for any fixed nonzero field
dilation.
\end{corollary}

\begin{proof}
The connected coefficients and their \(1/r!\) factors can improve
amplitude convergence, but they do not reduce the Gaussian type
\(r|a|^2\tau\) of a nonzero \(r\)-fold product.  Apply
Proposition~\ref{prop:v5v48-unbounded-order}.
\end{proof}

\begin{firewall}[Amplitude smallness and growth type are different]
\label{fw:v5v48-amplitude-type}
A factor \(g^r/r!\) can make a scalar series converge at each fixed
field value.  It does not place every partial sum in one Gaussian-type
ball when the exponent of \(|z|^2\) grows with \(r\).  Type control
requires exact resummation, stability on the real contour, or a stronger
functional space.
\end{firewall}

\subsection{Polynomial Boltzmann weights exceed Gaussian entire type}

\begin{proposition}[Quartic complex-growth obstruction]
\label{prop:v5v48-quartic-growth}
Let \(g>0\) and \(V(z)=gz^4\).  The entire function \(e^{-V(z)}\) does not
belong to \(\mathcal G_\tau\) for any finite \(\tau\).
\end{proposition}

\begin{proof}
On the ray \(z=re^{i\pi/4}\), one has \(z^4=-r^4\), so
\[
 |e^{-gz^4}|=e^{gr^4}.
\]
For every finite \(\tau\),
\(e^{gr^4-\tau r^2}\to\infty\).  Hence the norm
\eqref{eq:v5v48-Gtau} is infinite.
\end{proof}

\begin{corollary}[Stable real interaction is not globally stable complex
growth]
\label{cor:v5v48-real-complex}
Positivity of a quartic potential on the real integration contour does
not put its Boltzmann factor in a global Gaussian-type entire space.
\end{corollary}

\begin{proof}
The potential in Proposition~\ref{prop:v5v48-quartic-growth} is positive
on the real axis and has the opposite sign on the displayed complex ray.
\end{proof}

\begin{assessment}[Return to a measure-first chart]
\label{ass:v5v48-measure-first}
The obstruction points upstream to V8, V26, and V36.  The exact conditional
measure treats a stable real interaction before complex absolute values
are taken; its log-Laplace transform is contractive in a bounded-potential
norm, and its augmented source functional retains the conditional law.
The next chart should therefore combine the moving Gaussian quadratic
part with a real-contour relative-density norm for the non-Gaussian
remainder, rather than exponentiating the full interaction in
\(\mathcal G_\tau\).
\end{assessment}

\subsection{V48 status}

\begin{theorem}[V48 conclusion]
\label{thm:v5v48-status}
Gaussian-type entire functions give exact typed estimates for
multiplication, dilation, and Wick-frame heat transport.  A finite-valence
numerator closes when the explicit type budget
\eqref{eq:v5v48-budget-condition} holds.  The full exponentiated shell
interaction does not follow: connected order is unbounded, and even a
stable real quartic Boltzmann factor has growth beyond every Gaussian type
in complex directions.  Hence a GTBC can control finite complete
numerators but cannot replace the measure-first conditional chart needed
for the full SM--Einstein shell map.
\end{theorem}

\begin{proof}
Use Theorems~\ref{thm:v5v48-product},
\ref{thm:v5v48-dilation},
\ref{thm:v5v48-heat-bound}, and
\ref{thm:v5v48-type-budget}, followed by
Proposition~\ref{prop:v5v48-unbounded-order} and
Proposition~\ref{prop:v5v48-quartic-growth}.
\end{proof}

V49 will return to the real conditional measure.  It will build a
relative-density chart around the transported Gaussian, prove exact
conditional-expectation and log-Laplace contractions, and determine which
weighted oscillation or entropy norm can include stable Higgs-type
polynomials without reintroducing the global complex-growth obstruction.

\clearpage
\section{V49: a reference-invariant relative-density shell chart}
\label{sec:v5v49}

V47--V48 proved that potential coefficients and global entire-function
norms face incompatible product and covariance-transport requirements.
The obstruction arises after a probability measure has been converted
into a particular potential or Wick-coordinate representation.  This note
returns upstream and keeps the normalized density itself.  At finite
cutoff, integrating a shell is then conditional expectation, hence a
positive \(L^1\) contraction.  Changing the Gaussian reference multiplies
the density by a Radon--Nikodym derivative but is exactly isometric in
weighted \(L^1\).  Thus the moving-covariance connection which was
unbounded in V46--V47 disappears at the measure level.

The chart includes stable Higgs-type polynomial interactions and their
unbounded observables by using physical Lyapunov weights.  It does not
create strict contraction: marginalization preserves every density which
already depends only on the retained variables.  A continuum construction
still needs a local weighted mixing estimate, compact cofinal ancestry, or
a separate beta-function inverse on the retained coordinates.

\subsection{Normalized densities and shell marginalization}

Let \((P,\mu_P)\) and \((Q,\mu_Q)\) be probability spaces and put
\(\mu=\mu_P\otimes\mu_Q\).  A finite-cutoff state is represented by a
nonnegative density
\begin{equation}
 r\in L^1(\mu),
 \qquad
 \int_{P\times Q}r\,d\mu=1.
 \label{eq:v5v49-density}
\end{equation}
Its retained density after integrating the shell \(Q\) is
\begin{equation}
 (\mathsf M r)(p)
 =
 \int_Q r(p,q)\,d\mu_Q(q).
 \label{eq:v5v49-marginal}
\end{equation}

\begin{theorem}[Exact \(L^1\) shell contraction]
\label{thm:v5v49-L1-contraction}
For every signed \(h\in L^1(\mu)\),
\begin{equation}
 \|\mathsf M h\|_{L^1(\mu_P)}
 \le
 \|h\|_{L^1(\mu)}.
 \label{eq:v5v49-L1-contraction}
\end{equation}
Moreover, \(\mathsf M\) preserves positivity and total mass.  Hence it
maps normalized densities to normalized densities and is nonexpansive in
total variation.
\end{theorem}

\begin{proof}
By the triangle inequality inside the conditional integral and Tonelli's
theorem,
\begin{align*}
 \int_P|(\mathsf Mh)(p)|\,d\mu_P(p)
 &\le
 \int_P\int_Q|h(p,q)|\,d\mu_Q(q)\,d\mu_P(p)\\
 &=\|h\|_{L^1(\mu)}.
\end{align*}
Positivity is immediate from \eqref{eq:v5v49-marginal}, and Fubini gives
\(\int\mathsf Mr\,d\mu_P=\int r\,d\mu\).
\end{proof}

\begin{corollary}[Data processing for finite shell kernels]
\label{cor:v5v49-data-processing}
Let \(K\) be any Markov kernel from a finite-cutoff state space \(X\) to a
coarser state space \(Y\).  Then its pushforward obeys
\begin{equation}
 \|K_*\nu-K_*\xi\|_{\mathrm{TV}}
 \le
 \|\nu-\xi\|_{\mathrm{TV}}.
 \label{eq:v5v49-data-processing}
\end{equation}
\end{corollary}

\begin{proof}
For every measurable \(f\) with \(\|f\|_\infty\le1\), the function
\(Kf(x)=\int f(y)K(x,dy)\) also has supremum norm at most one.  Use the
dual characterization of total variation:
\[
 2\|K_*\nu-K_*\xi\|_{\mathrm{TV}}
 =
 \sup_{\|f\|_\infty\le1}
 \left|\int Kf\,d(\nu-\xi)\right|.
\]
\end{proof}

\subsection{Changing the Gaussian frame is an \(L^1\) isometry}

Let \(\mu\) and \(\widetilde\mu\) be equivalent probability references on
the same measurable space, and set
\begin{equation}
 h=\frac{d\mu}{d\widetilde\mu}.
 \label{eq:v5v49-reference-h}
\end{equation}
If \(r=d\nu/d\mu\), then
\(\widetilde r=d\nu/d\widetilde\mu=rh\).

\begin{theorem}[Reference-invariant density distance]
\label{thm:v5v49-reference-isometry}
For two finite signed measures \(\nu,\xi\ll\mu,\widetilde\mu\), with
densities \(r,s\) relative to \(\mu\) and
\(\widetilde r,\widetilde s\) relative to \(\widetilde\mu\),
\begin{equation}
 \|\widetilde r-\widetilde s\|_{L^1(\widetilde\mu)}
 =
 \|r-s\|_{L^1(\mu)}.
 \label{eq:v5v49-reference-isometry}
\end{equation}
More generally, for every nonnegative physical weight \(\Omega\),
\begin{equation}
 \int\Omega|\widetilde r-\widetilde s|\,d\widetilde\mu
 =
 \int\Omega|r-s|\,d\mu.
 \label{eq:v5v49-weighted-reference-isometry}
\end{equation}
\end{theorem}

\begin{proof}
Use
\(\widetilde r-\widetilde s=(r-s)h\) and
\(h\,d\widetilde\mu=d\mu\).  The same calculation with the factor
\(\Omega\) proves the weighted identity.
\end{proof}

\begin{corollary}[The Gaussian Wick connection is a coordinate effect]
\label{cor:v5v49-Wick-coordinate}
At finite cutoff, two nondegenerate Gaussian reference covariances are
equivalent.  Re-expressing the same interacting state relative to either
Gaussian is an isometry in the density distance
\eqref{eq:v5v49-reference-isometry}, even when the corresponding
Wick-coefficient connection is unbounded in the V47 factorial norm.
\end{corollary}

\begin{proof}
Finite-dimensional nondegenerate Gaussian measures have strictly positive
Lebesgue densities and are mutually absolutely continuous.  Apply
Theorem~\ref{thm:v5v49-reference-isometry}.  V47 concerns a different
coordinate norm on polynomial coefficients, so there is no contradiction.
\end{proof}

\begin{firewall}[Finite-cutoff equivalence is not continuum equivalence]
\label{fw:v5v49-Gaussian-singularity}
In infinite dimensions, Gaussian measures with different covariances can
be mutually singular unless the Feldman--Hájek condition holds.  The
finite-cutoff isometry therefore gives a uniform continuum chart only if
the relative densities or their local restrictions have cutoff-independent
weighted bounds.  Formal covariance transport does not prove that.
\end{firewall}

\subsection{The exact augmented density chart}

Let \(m=\mathsf Mr\).  On the set where \(m(p)>0\), define
\begin{equation}
 k_r(q\mid p)=\frac{r(p,q)}{m(p)}.
 \label{eq:v5v49-conditional-density}
\end{equation}
Set \(k_r=1\) on a null set where \(m=0\), using any fixed normalized
reference density there.

\begin{theorem}[Marginal--conditional reconstruction]
\label{thm:v5v49-density-reconstruction}
The pair \((m,k_r)\) satisfies
\begin{equation}
 r(p,q)=m(p)k_r(q\mid p)
 \quad\text{for }\mu\text{-almost every }(p,q),
 \label{eq:v5v49-density-factorization}
\end{equation}
and
\begin{equation}
 \int_Q k_r(q\mid p)\,d\mu_Q(q)=1
 \quad\text{for }m\,d\mu_P\text{-almost every }p.
 \label{eq:v5v49-conditional-normalization}
\end{equation}
Thus the marginal together with the conditional density is an injective
augmented shell chart.
\end{theorem}

\begin{proof}
Both statements follow directly from
\eqref{eq:v5v49-marginal} and
\eqref{eq:v5v49-conditional-density}.  Their product reconstructs \(r\)
where \(m>0\), while the set \(m=0\) carries zero joint mass.
\end{proof}

\begin{theorem}[Stability of the augmented factorization]
\label{thm:v5v49-factorization-stability}
Let \(r=mk\) and \(s=n\ell\) be two normalized joint densities with
conditional densities \(k,\ell\).  Then
\begin{align}
 \|m-n\|_{L^1(\mu_P)}
 &\le\|r-s\|_{L^1(\mu)},\label{eq:v5v49-marginal-stability}\\
 \|r-s\|_{L^1(\mu)}
 &\le
 \|m-n\|_{L^1(\mu_P)}
 +
 \int_P n(p)
 \|k(\cdot\mid p)-\ell(\cdot\mid p)\|_{L^1(\mu_Q)}
 \,d\mu_P(p).
 \label{eq:v5v49-joint-stability}
\end{align}
\end{theorem}

\begin{proof}
The first line is
Theorem~\ref{thm:v5v49-L1-contraction} applied to \(r-s\).  For the
second, write
\[
 mk-n\ell=(m-n)k+n(k-\ell),
\]
take absolute values, integrate over \(q\), use
\(\int k\,d\mu_Q=1\), and then integrate over \(p\).
\end{proof}

\begin{assessment}[Relation to the V36 source functional]
\label{ass:v5v49-V36}
When the conditional Laplace transform exists, the V36 functional
\(\mathcal K_p(J)\) is the logarithmic moment transform of
\(k_r(\cdot\mid p)\).  The density chart retains the same conditional law
without first expanding it into a factorial tensor tower.  Taylor
coefficients may be extracted only after a weighted moment estimate is
available.
\end{assessment}

\subsection{Stable polynomial interactions belong to the density chart}

Let \(\gamma\) be a nondegenerate Gaussian probability measure on
\(\mathbb R^N\).  For a measurable interaction \(U\), define
\begin{equation}
 Z_U=\int e^{-U(y)}\,d\gamma(y),
 \qquad
 r_U(y)=Z_U^{-1}e^{-U(y)}.
 \label{eq:v5v49-Gibbs-density}
\end{equation}

\begin{theorem}[Coercive polynomial density and exponential moments]
\label{thm:v5v49-coercive-density}
Suppose for some integer \(d\ge1\),
\begin{equation}
 U(y)\ge a|y|^{2d}-b,
 \qquad a>0.
 \label{eq:v5v49-coercive-U}
\end{equation}
Then:
\begin{enumerate}
 \item \(0<Z_U<\infty\) and \(r_U\in L^1(\gamma)\cap L^\infty(\gamma)\);
 \item for every \(0\le\kappa<a\),
 \begin{equation}
  \int e^{\kappa|y|^{2d}}r_U(y)\,d\gamma(y)<\infty;
  \label{eq:v5v49-exponential-moment}
 \end{equation}
 \item every polynomial moment and every finite-dimensional linear
 Laplace transform
 \(\int e^{\langle J,y\rangle}r_U\,d\gamma\) is finite.
\end{enumerate}
\end{theorem}

\begin{proof}
The lower bound gives
\[
 0<e^{-U(y)}\le e^b e^{-a|y|^{2d}}\le e^b,
\]
so \(0<Z_U\le e^b\) and
\(\|r_U\|_\infty\le e^b/Z_U\).  For \(\kappa<a\),
\[
 \int e^{\kappa|y|^{2d}}r_U\,d\gamma
 \le
 \frac{e^b}{Z_U}
 \int e^{-(a-\kappa)|y|^{2d}}\,d\gamma
 <\infty.
\]
Every polynomial is bounded by a constant times
\(e^{\kappa|y|^{2d}}\).  Young's inequality also gives, for fixed \(J\),
\(\langle J,y\rangle\le C_J+\kappa|y|^{2d}\).  This proves the last
claim.
\end{proof}

\begin{corollary}[Higgs-type quartics avoid the V48 complex obstruction]
\label{cor:v5v49-Higgs-density}
A real positive quartic interaction defines a bounded normalized density
relative to a finite-cutoff Gaussian and has all polynomial moments,
although its complexified Boltzmann factor lies in no V48 Gaussian-type
entire space.
\end{corollary}

\begin{proof}
Apply Theorem~\ref{thm:v5v49-coercive-density} with \(d=2\), and compare
with Proposition~\ref{prop:v5v48-quartic-growth}.
\end{proof}

\begin{firewall}[The normalization constant must be uniform locally]
\label{fw:v5v49-normalization}
The finite-dimensional bound
\(\|r_U\|_\infty\le e^b/Z_U\) may grow exponentially with the cutoff
volume.  A continuum local chart requires rooted or conditional estimates
whose constants do not contain the full-volume partition function.  The
normalized \(L^1\) identity alone does not provide such locality.
\end{firewall}

\subsection{Physical Lyapunov weights}

Let \(\Omega:P\times Q\to[1,\infty)\) and
\(\Omega_P:P\to[1,\infty)\) be measurable weights.  Define
\begin{equation}
 \|h\|_{1,\Omega}
 =
 \int_{P\times Q}\Omega(p,q)|h(p,q)|\,d\mu(p,q).
 \label{eq:v5v49-weighted-L1}
\end{equation}

\begin{theorem}[Weighted shell contraction]
\label{thm:v5v49-weighted-contraction}
If
\begin{equation}
 \Omega_P(p)\le\Omega(p,q)
 \quad\text{for }\mu\text{-almost every }(p,q),
 \label{eq:v5v49-weight-domination}
\end{equation}
then
\begin{equation}
 \|\mathsf Mh\|_{1,\Omega_P}
 \le\|h\|_{1,\Omega}.
 \label{eq:v5v49-weighted-contraction}
\end{equation}
\end{theorem}

\begin{proof}
Use the proof of Theorem~\ref{thm:v5v49-L1-contraction} and insert the
pointwise inequality:
\begin{align*}
 \int\Omega_P(p)|\mathsf Mh(p)|\,d\mu_P
 &\le
 \int\Omega_P(p)|h(p,q)|\,d\mu\\
 &\le
 \int\Omega(p,q)|h(p,q)|\,d\mu.
\end{align*}
\end{proof}

\begin{corollary}[Dual observable control]
\label{cor:v5v49-observable-dual}
If a physical observable satisfies \(|A|\le C_A\Omega\), then for two
density states \(r,s\),
\begin{equation}
 \left|\int A(r-s)\,d\mu\right|
 \le C_A\|r-s\|_{1,\Omega}.
 \label{eq:v5v49-observable-dual}
\end{equation}
\end{corollary}

\begin{proof}
Integrate \(|A||r-s|\le C_A\Omega|r-s|\).
\end{proof}

\begin{proposition}[Exponential Lyapunov weight for a coercive interaction]
\label{prop:v5v49-coercive-weight}
Under \eqref{eq:v5v49-coercive-U}, every weight
\begin{equation}
 \Omega_\kappa(y)=e^{\kappa|y|^{2d}},
 \qquad 0\le\kappa<a,
 \label{eq:v5v49-Omega-kappa}
\end{equation}
is integrable under the Gibbs state.  Its dual class contains all
polynomial local observables and all fixed linear source exponentials.
\end{proposition}

\begin{proof}
This is items 2--3 of
Theorem~\ref{thm:v5v49-coercive-density}.
\end{proof}

\subsection{Why the density contraction is not strict}

\begin{proposition}[Retained-variable isometric sector]
\label{prop:v5v49-no-strictness}
Let \(r(p,q)=r_0(p)\) and \(s(p,q)=s_0(p)\) be normalized densities
independent of \(q\).  Then
\begin{equation}
 \|\mathsf Mr-\mathsf Ms\|_{L^1(\mu_P)}
 =
 \|r-s\|_{L^1(\mu)}.
 \label{eq:v5v49-isometric-sector}
\end{equation}
Consequently, shell marginalization has contraction coefficient one on
the full density space whenever the retained density sector is
nontrivial.
\end{proposition}

\begin{proof}
The marginal map sends \(r\) to \(r_0\) and \(s\) to \(s_0\).  Because
\(\mu_Q\) has total mass one, both sides of
\eqref{eq:v5v49-isometric-sector} equal
\(\int_P|r_0-s_0|\,d\mu_P\).
\end{proof}

\begin{firewall}[Nonexpansive is not a fixed-point theorem]
\label{fw:v5v49-nonexpansive}
Theorem~\ref{thm:v5v49-L1-contraction} removes the Wick-coordinate
blow-up but supplies no strict factor below one.  It cannot replace the
V26 Schur contraction, the V42 derivative inverse, or the V40 cumulative
diameter condition on retained variables.
\end{firewall}

For a genuinely mixing Markov kernel one may seek a weighted total
variation contraction.  A simple sufficient finite-state form is a
minorization.

\begin{theorem}[Doeblin contraction]
\label{thm:v5v49-Doeblin}
Suppose a Markov kernel \(K\) satisfies
\begin{equation}
 K(x,\cdot)\ge\varepsilon\,\pi(\cdot)
 \quad\text{for every }x
 \label{eq:v5v49-minorization}
\end{equation}
for a probability measure \(\pi\) and \(0<\varepsilon\le1\).  Then
\begin{equation}
 \|K_*\nu-K_*\xi\|_{\mathrm{TV}}
 \le(1-\varepsilon)\|\nu-\xi\|_{\mathrm{TV}}.
 \label{eq:v5v49-Doeblin}
\end{equation}
\end{theorem}

\begin{proof}
Write
\[
 K(x,\cdot)=\varepsilon\pi(\cdot)
 +(1-\varepsilon)R(x,\cdot)
\]
with \(R\) a Markov kernel.  The common \(\varepsilon\pi\) part cancels
between the two pushforwards, and
Corollary~\ref{cor:v5v49-data-processing} applied to \(R\) gives the
bound.
\end{proof}

\begin{assessment}[Uniform minorization is too strong for field space]
\label{ass:v5v49-minorization}
An infinite-dimensional field kernel with unbounded amplitudes rarely has
a global Doeblin constant.  The relevant replacement is a local
minorization on a Lyapunov sublevel set plus drift back to that set, or a
rooted Dobrushin influence estimate.  These constants must be uniform in
cutoff and compatible with gauge constraints.
\end{assessment}

\subsection{Relative-density shell certificate}

\begin{definition}[Relative-density shell certificate]
\label{def:v5v49-RDSC}
An RDSC consists of:
\begin{enumerate}
 \item normalized finite-cutoff physical densities relative to a positive
 sliced reference measure;
 \item a coercive, cutoff-uniform local Lyapunov weight controlling the
 Higgs, matter, gauge, and physical metric amplitudes;
 \item exact Markov pushforward and reference-change identities
 \eqref{eq:v5v49-L1-contraction} and
 \eqref{eq:v5v49-weighted-reference-isometry};
 \item the injective marginal--conditional chart of
 Theorem~\ref{thm:v5v49-density-reconstruction};
 \item a rooted weighted mixing or compactness theorem strong enough to
 control the retained isometric sector;
 \item Ward, BRST, OS reflection, and local composite-observable closure in
 the dual weighted class;
 \item compatibility with the stratified crossover and UV interface.
\end{enumerate}
\end{definition}

\begin{firewall}[Einstein coercivity remains open]
\label{fw:v5v49-Einstein-coercivity}
Theorem~\ref{thm:v5v49-coercive-density} applies once a positive physical
finite-cutoff reference and a coercive reduced interaction are available.
V8--V10 showed that the Euclidean conformal metric direction does not
automatically provide them.  Passing to the density chart does not solve
the conformal-factor, contour, or reflection-positivity problem.
\end{firewall}

\subsection{V49 status}

\begin{theorem}[V49 conclusion]
\label{thm:v5v49-status}
At finite cutoff, normalized relative densities give an exact
reference-invariant shell chart: marginalization is a positive weighted
\(L^1\) contraction, Gaussian frame changes are isometries, and the
marginal together with its conditional density reconstructs the full
state.  Coercive stable polynomial interactions, including positive
quartics, belong to this chart with exponential Lyapunov moments.  The
chart removes the product--connection obstruction of V47 at the measure
level, but its shell map is only nonexpansive on retained densities.
Uniform local mixing and physical Einstein coercivity remain unproved.
\end{theorem}

\begin{proof}
Use Theorems~\ref{thm:v5v49-L1-contraction},
\ref{thm:v5v49-reference-isometry},
\ref{thm:v5v49-density-reconstruction},
\ref{thm:v5v49-coercive-density}, and
\ref{thm:v5v49-weighted-contraction}, together with
Proposition~\ref{prop:v5v49-no-strictness}.
\end{proof}

V50 will perform the next compulsory YM/NS audit.  It will then formulate
a rooted Dobrushin influence matrix for the conditional-density chart and
test whether complete cumulant assembly can yield a cutoff-uniform
contraction on local observables while leaving the retained marginal to
the finite-dimensional Schur or UV fixed-point system.

\clearpage
\section{V50: companion audit and a rooted conditional-influence matrix}
\label{sec:v5v50}

V49 replaced the unstable potential-coordinate shell connection by a
reference-invariant relative-density chart.  Marginalization is
nonexpansive, but it is not strict on retained-variable densities.  This
note performs the compulsory companion audit and constructs the next
quantitative object: a rooted matrix of conditional influences.  If its
weighted row norm is below one, conditional fibres contract and boundary
perturbations decay exponentially with distance.  The retained marginal
must still be controlled by the V26/V39 base system.

The YM audit reinforces the required assembly order.  Its newest work
closes every nondiscrete quintic predecessor partition by complete relative
cumulants, raw pinching, and one final payer.  The sole remaining quintic
source is reduced to a coefficientwise paid selected-junction gate.  Thus
an SM influence coefficient must be assigned only after its complete
conditional cumulant, Ward partners, output compression, and local
normalization have been assembled.

\subsection{Mandatory V50 companion audit}

The newest named YM file is
\begin{equation}
 \texttt{Renewal\_Yang\_Mills\_Mass\_Gap\_Research\_Notes\_V57.tex},
 \qquad
 \mathrm{SHA256}=
 \texttt{183A509AC967230360EB9D0CE000A73684FC1338E5CCCB849F0714E55B6F0F21}.
 \label{eq:v5v50-YM-hash}
\end{equation}
The simultaneously present V56 file has the same byte length and the same
digest, so V57 was used as the newest-name audit target.  No companion file
was modified.

The NS file remains
\begin{equation}
 \texttt{Renewal\_NS\_Stability\_Research\_Notes\_V31.tex},
 \qquad
 \mathrm{SHA256}=
 \texttt{F1121B62238AD5491BB7CF03635EA9934942EDF573ED8FAAA9EE79E1D38A6696}.
 \label{eq:v5v50-NS-hash}
\end{equation}
Its digest and mathematical endpoint are unchanged from V45.

The relevant YM changes since the V45 audit are:
\begin{enumerate}
 \item spectator retention and corrected clock accounting close every
 \(3+2\) first-connectivity source;
 \item payer stripping plus a scalar spectator lift close \(4+1\);
 \item complete relative cumulants of block products close the
 \(3+1+1\), \(2+2+1\), and \(2+1+1+1\) partitions without iterated
 intermediate cumulants or resolvents;
 \item the fully discrete five-block source is reduced to the exact
 coefficientwise gate \(\operatorname{PGSEL}_5\);
 \item coordinate aggregation of its four selected tree marks is proved
 and creates no support count;
 \item the output-dominant and branching-dominant parts of
 \(\operatorname{PGSEL}_5\) remain open.
\end{enumerate}
The YM mass gap, all-order atlas, cutoff removal, OS reconstruction, and
continuum measure remain explicitly unproved.

\begin{assessment}[Audit transfer to the density chart]
\label{ass:v5v50-audit-transfer}
The closed YM partitions show that a local influence entry should be a
complete block-product cumulant rather than a sum of sequential pair
influences.  Raw equivariant pinching can remove output multiplicity, and a
single final inverse may be applied after assembly.  The unresolved
discrete fifth-cumulant gate prevents extrapolation of those fixed-order
facts to an all-order SM influence row.
\end{assessment}

\begin{firewall}[No modification of companion inventories]
\label{fw:v5v50-companion-inventory}
V56 and V57 being byte-identical is an inventory fact about the separate
YM programme.  This SM audit neither deletes nor renames either file,
because the shared folder contains independently managed programmes.
\end{firewall}

\subsection{Variation of a conditional Gibbs density}

Let \(\mu\) be a probability reference and let
\begin{equation}
 d\nu_t=Z_t^{-1}e^{-V_t}\,d\mu,
 \qquad 0\le t\le1,
 \label{eq:v5v50-Gibbs-path}
\end{equation}
where \(V_t\) is differentiable and bounded for the present lemma.  Write
\(\langle\cdot\rangle_t\) for expectation under \(\nu_t\).

\begin{lemma}[Exact density derivative]
\label{lem:v5v50-density-derivative}
The normalized density \(r_t=d\nu_t/d\mu\) satisfies
\begin{equation}
 \partial_t r_t
 =
 -\bigl(\dot V_t-\langle\dot V_t\rangle_t\bigr)r_t.
 \label{eq:v5v50-density-derivative}
\end{equation}
Consequently, for every bounded observable \(A\),
\begin{equation}
 \frac d{dt}\langle A\rangle_t
 =
 -\operatorname{Cov}_t(A,\dot V_t).
 \label{eq:v5v50-expectation-derivative}
\end{equation}
\end{lemma}

\begin{proof}
Differentiate \(r_t=Z_t^{-1}e^{-V_t}\) and use
\[
 \partial_t\log Z_t=-\langle\dot V_t\rangle_t.
\]
Integration of \(A\partial_t r_t\) gives the covariance formula.
\end{proof}

For a bounded real function \(F\), set
\begin{equation}
 \operatorname{osc}(F)
 =
 \operatorname*{ess\,sup}F-\operatorname*{ess\,inf}F.
 \label{eq:v5v50-oscillation}
\end{equation}

\begin{lemma}[Oscillation covariance bound]
\label{lem:v5v50-osc-covariance}
For bounded real random variables \(A,B\),
\begin{equation}
 |\operatorname{Cov}(A,B)|
 \le\frac14\operatorname{osc}(A)\operatorname{osc}(B).
 \label{eq:v5v50-osc-covariance}
\end{equation}
\end{lemma}

\begin{proof}
Popoviciu's inequality gives
\(\operatorname{Var}(A)\le\operatorname{osc}(A)^2/4\) and the analogous
bound for \(B\).  Apply Cauchy--Schwarz to the covariance.
\end{proof}

\begin{theorem}[Conditional total-variation influence bound]
\label{thm:v5v50-TV-influence}
For the path \eqref{eq:v5v50-Gibbs-path},
\begin{equation}
 \|\nu_1-\nu_0\|_{\mathrm{TV}}
 \le
 \frac14\int_0^1\operatorname{osc}(\dot V_t)\,dt.
 \label{eq:v5v50-TV-influence}
\end{equation}
\end{theorem}

\begin{proof}
Use the dual identity
\[
 \|\nu_1-\nu_0\|_{\mathrm{TV}}
 =
 \frac12\sup_{\|A\|_\infty\le1}
 |\langle A\rangle_1-\langle A\rangle_0|.
\]
For such \(A\), \(\operatorname{osc}(A)\le2\).  Integrate
\eqref{eq:v5v50-expectation-derivative} and apply
Lemma~\ref{lem:v5v50-osc-covariance}.
\end{proof}

\begin{firewall}[The oscillation estimate is only the first card]
\label{fw:v5v50-oscillation-card}
Unbounded Higgs and metric interactions may have infinite global
oscillation.  Even for bounded Wilson variables, separate absolute
oscillations can discard the complete cumulant cancellations used in the
YM audit.  Equation~\eqref{eq:v5v50-TV-influence} is a valid crude card;
the physical coefficient should be improved by complete conditional
assembly before it enters the influence matrix.
\end{firewall}

\subsection{A rooted Dobrushin matrix}

Let \(\Lambda\) be a finite set of blocks.  For each \(i\in\Lambda\), let
\(\Gamma_i(\cdot\mid x)\) be the conditional update law of block \(i\)
given the other blocks.  Equip each block with a lower semicontinuous
metric \(d_i\le1\), and let \(W_i\) denote the corresponding
Kantorovich--Rubinstein distance.  Suppose nonnegative coefficients
\(c_{ij}\) satisfy
\begin{equation}
 W_i\bigl(\Gamma_i(\cdot\mid x),
          \Gamma_i(\cdot\mid y)\bigr)
 \le
 \sum_{j\ne i}c_{ij}d_j(x_j,y_j).
 \label{eq:v5v50-local-influence}
\end{equation}
Write \(C=(c_{ij})\), with \(c_{ii}=0\).

\begin{theorem}[Weighted conditional contraction]
\label{thm:v5v50-weighted-Dobrushin}
Assume there are weights \(v_i>0\) and \(0\le\theta<1\) such that
\begin{equation}
 \sum_j c_{ij}v_j\le\theta v_i
 \qquad(i\in\Lambda).
 \label{eq:v5v50-weighted-row}
\end{equation}
If two coupled boundary configurations have discrepancy bounds
\[
 \mathbb E d_j(X_j,Y_j)\le Dv_j,
\]
then optimal couplings of the conditional updates can be chosen so that
\begin{equation}
 \mathbb E d_i(X_i',Y_i')
 \le\theta Dv_i.
 \label{eq:v5v50-one-update-contraction}
\end{equation}
\end{theorem}

\begin{proof}
By the definition of \(W_i\), choose an optimal or arbitrarily near-optimal
coupling of the two conditional laws.  Average
\eqref{eq:v5v50-local-influence} over the boundary coupling:
\[
 \mathbb E d_i(X_i',Y_i')
 \le
 \sum_jc_{ij}\mathbb E d_j(X_j,Y_j)
 \le
 D\sum_jc_{ij}v_j
 \le\theta Dv_i.
\]
Pass to the limit if only near-optimal couplings exist.
\end{proof}

For overlapping block updates, choose colour classes
\(\Lambda=\Lambda_1\cup\cdots\cup\Lambda_s\) such that blocks in one class
can be conditionally coupled in parallel.  Let \(B_a\) be the nonnegative
matrix which replaces row \(i\in\Lambda_a\) by row \(i\) of \(C\) and
leaves every other row equal to the corresponding identity row.

\begin{theorem}[Coloured-sweep contraction]
\label{thm:v5v50-coloured-sweep}
Let
\begin{equation}
 B=B_sB_{s-1}\cdots B_1.
 \label{eq:v5v50-sweep-matrix}
\end{equation}
If, for some \(v_i>0\),
\begin{equation}
 Bv\le\vartheta v
 \qquad\text{with }0\le\vartheta<1,
 \label{eq:v5v50-sweep-contraction}
\end{equation}
then one full coloured conditional sweep contracts the weighted maximal
expected discrepancy by the factor \(\vartheta\).
\end{theorem}

\begin{proof}
At colour \(a\), unchanged blocks retain their previous discrepancy and
updated blocks obey Theorem~\ref{thm:v5v50-weighted-Dobrushin} with the
current discrepancy vector.  Hence the vector after the update is bounded
componentwise by \(B_a\) times the preceding vector.  Iterate through the
colours and use \eqref{eq:v5v50-sweep-contraction}.
\end{proof}

\begin{firewall}[A single-site random scan is volume slow]
\label{fw:v5v50-random-scan}
Updating only one of \(|\Lambda|\) sites per step generally gives a
contraction factor tending to one with the volume.  The continuum shell
certificate needs a parallel finite-colour sweep, a block update of fixed
physical size, or a direct rooted comparison estimate whose constant is
volume independent.
\end{firewall}

\subsection{Rooted spatial decay}

Let \(d(i,j)\) be a graph distance and define
\begin{equation}
 \theta_m
 =
 \sup_i\sum_j e^{m d(i,j)}c_{ij}.
 \label{eq:v5v50-theta-m}
\end{equation}

\begin{theorem}[Exponential Dobrushin Green kernel]
\label{thm:v5v50-Green-decay}
If \(\theta_m<1\) for some \(m>0\), then the Neumann Green matrix
\begin{equation}
 G=(I-C)^{-1}=\sum_{n\ge0}C^n
 \label{eq:v5v50-Green}
\end{equation}
exists on weighted \(\ell^\infty\), and its entries obey
\begin{equation}
 0\le G_{ij}
 \le
 \frac{e^{-m d(i,j)}}{1-\theta_m}.
 \label{eq:v5v50-Green-decay}
\end{equation}
\end{theorem}

\begin{proof}
For every path
\(i=i_0,i_1,\ldots,i_n=j\), the triangle inequality gives
\[
 d(i,j)\le\sum_{\ell=0}^{n-1}d(i_\ell,i_{\ell+1}).
\]
Therefore
\begin{align*}
 e^{m d(i,j)}(C^n)_{ij}
 &\le
 \sum_{i_1,\ldots,i_{n-1}}
 \prod_{\ell=0}^{n-1}
 e^{m d(i_\ell,i_{\ell+1})}
 c_{i_\ell i_{\ell+1}}\\
 &\le\theta_m^n.
\end{align*}
Thus
\((C^n)_{ij}\le e^{-md(i,j)}\theta_m^n\).  Sum the geometric series.
\end{proof}

\begin{corollary}[Rooted boundary stability]
\label{cor:v5v50-rooted-stability}
If a nonnegative local discrepancy vector satisfies
\begin{equation}
 u\le Cu+b,
 \label{eq:v5v50-comparison-inequality}
\end{equation}
then
\begin{equation}
 u_i
 \le
 \frac1{1-\theta_m}
 \sum_j e^{-md(i,j)}b_j.
 \label{eq:v5v50-rooted-stability}
\end{equation}
\end{corollary}

\begin{proof}
Iterate \eqref{eq:v5v50-comparison-inequality}.  The remainder
\(C^nu\) tends to zero in the weighted norm, leaving \(u\le Gb\).
Apply Theorem~\ref{thm:v5v50-Green-decay}.
\end{proof}

\begin{assessment}[Why this matches the rooted continuum norm]
\label{ass:v5v50-rooted-match}
The exponential Green estimate localizes the response to a marked root
before volume summation.  It is the probability-kernel analogue of the
rooted covariance and complete-transition estimates used in V31--V40.
The row sum must be taken only after each protected local block has been
assembled.
\end{assessment}

\subsection{Complete-block influence assignment}

Suppose changing boundary block \(j\) produces an interpolating conditional
action \(V_{i,t}\) at output block \(i\).  The crude assignment is
\begin{equation}
 c_{ij}^{\mathrm{osc}}
 =
 \min\left\{
 1,\,
 \frac14\sup_{\text{\(j\)-boundary pairs}}
 \int_0^1\operatorname{osc}_i(\dot V_{i,t})\,dt
 \right\}.
 \label{eq:v5v50-crude-cij}
\end{equation}
For unbounded blocks this is replaced by a weighted Wasserstein estimate
in \eqref{eq:v5v50-local-influence}.

\begin{definition}[Complete conditional influence entry]
\label{def:v5v50-CCIE}
A CCIE \(c_{ij}\) is an estimate of
\eqref{eq:v5v50-local-influence} obtained only after:
\begin{enumerate}
 \item the full conditional density ratio for the boundary change is
 normalized;
 \item all moment--cumulant partitions joining \(j\) to \(i\) are
 assembled;
 \item scalar, physical, Ward, BRST, and multiplicity partners are placed
 in their complete output blocks;
 \item local spectator multipliers and Lyapunov weights are retained;
 \item any final inverse or payer is applied once.
\end{enumerate}
\end{definition}

\begin{proposition}[Fixed-order YM transfer and its limit]
\label{prop:v5v50-YM-transfer}
The audited YM V52--V56 estimates verify the CCIE assembly order for all
nondiscrete quintic predecessor partitions in their Wilson norm.  V57
reduces the remaining discrete fifth-cumulant entry to
\(\operatorname{PGSEL}_5\), but does not prove it on the
output-dominant or branching-dominant branches.  Therefore those notes do
not yet supply a complete fifth-order YM row, and they do not imply the
all-order SM bound \(\theta_m<1\).
\end{proposition}

\begin{proof}
This is the exact scope recorded in the audit at the beginning of V50.
The Dobrushin norm sums every incident input \(j\) at a root \(i\);
an unresolved positive block cannot be omitted from that row sum.
\end{proof}

\subsection{Base--fibre closure}

Let \(d_{\mathrm b}\) measure the retained marginal or finite coupling
coordinates, and let \(d_{\mathrm f}\) be the rooted conditional-density
distance controlled above.  Suppose one full shell step obeys
\begin{align}
 d_{\mathrm b}'&\le a\,d_{\mathrm b},
 \label{eq:v5v50-base}\\
 d_{\mathrm f}'&\le g\,d_{\mathrm b}+\vartheta\,d_{\mathrm f}.
 \label{eq:v5v50-fibre}
\end{align}

\begin{theorem}[Triangular density-chart contraction]
\label{thm:v5v50-triangular}
If \(a<1\) and \(\vartheta<1\), then there is \(\lambda>0\) and
\(\Theta<1\) such that
\begin{equation}
 d_{\mathrm b}'+\lambda d_{\mathrm f}'
 \le
 \Theta(d_{\mathrm b}+\lambda d_{\mathrm f}).
 \label{eq:v5v50-triangular-contraction}
\end{equation}
One may take any
\begin{equation}
 0<\lambda<\frac{1-a}{g}
 \label{eq:v5v50-lambda}
\end{equation}
when \(g>0\), followed by
\(\Theta=\max\{a+\lambda g,\vartheta\}<1\).
\end{theorem}

\begin{proof}
Multiply \eqref{eq:v5v50-fibre} by \(\lambda\) and add
\eqref{eq:v5v50-base}:
\[
 d_{\mathrm b}'+\lambda d_{\mathrm f}'
 \le
 (a+\lambda g)d_{\mathrm b}
 +\lambda\vartheta d_{\mathrm f}.
\]
Condition \eqref{eq:v5v50-lambda} makes the first coefficient below one;
the second normalized coefficient is \(\vartheta<1\).
\end{proof}

\begin{firewall}[The retained factor is still missing]
\label{fw:v5v50-base-missing}
The conditional influence theorem can provide
\(\vartheta<1\) for the fibre.  V49 proves only
\(a\le1\) for the retained marginal.  A strict \(a<1\) must come from the
finite-dimensional SM beta system, a UV fixed point/interface, or a
quotient that removes the isometric retained directions.
\end{firewall}

\subsection{Rooted density influence certificate}

\begin{definition}[Rooted density influence certificate]
\label{def:v5v50-RDIC}
An RDIC consists of:
\begin{enumerate}
 \item the V49 relative-density shell chart with cutoff-uniform local
 Lyapunov weights;
 \item complete conditional influence entries in the sense of the CCIE;
 \item a finite-colour update or direct block shell kernel;
 \item a weighted exponential row bound
 \(\theta_m<1\) uniform in cutoff and volume;
 \item a strict retained-base estimate \(a<1\) and finite coupling
 \(g\) in \eqref{eq:v5v50-base}--\eqref{eq:v5v50-fibre};
 \item Ward, BRST, OS, stratified crossover, and UV-interface
 compatibility.
\end{enumerate}
\end{definition}

\subsection{V50 status}

\begin{theorem}[V50 conclusion]
\label{thm:v5v50-status}
A cutoff-uniform complete conditional-influence matrix with
\(\theta_m<1\) yields strict fibre contraction and exponentially decaying
rooted boundary response.  Combined with a strict retained-base estimate,
it contracts the full augmented relative-density chart.  The YM V57 audit
supports the complete-block assembly order through all nondiscrete
quintic partitions, but leaves the discrete fifth-cumulant gate partly
open and supplies no all-order SM row bound.  Neither an RDIC nor the
retained-base factor \(a<1\) has yet been proved for the full
SM--Einstein shell system.
\end{theorem}

\begin{proof}
Use Theorems~\ref{thm:v5v50-weighted-Dobrushin},
\ref{thm:v5v50-coloured-sweep},
\ref{thm:v5v50-Green-decay}, and
\ref{thm:v5v50-triangular}, together with
Proposition~\ref{prop:v5v50-YM-transfer}.
\end{proof}

V51 will test the influence row in the first unbounded but exactly
solvable matter sector: a coercive scalar site with quadratic nearest-
neighbour coupling.  It will derive a weighted Wasserstein coefficient
from strong convexity, compare its lattice row sum with one, and identify
which Higgs and metric signs preserve or destroy the contraction.

\clearpage
\section{V51: strong-convexity influence and the continuum scaling test}
\label{sec:v5v51}

V50 reduced strict fibre contraction to a cutoff-uniform conditional
influence matrix.  This note computes that matrix in the first unbounded
matter model where all constants are explicit: a scalar lattice field with
strongly convex on-site potential and quadratic nearest-neighbour coupling.
Strong convexity gives an exact Wasserstein influence card.  At fixed
lattice spacing, a positive on-site curvature makes the Dobrushin row sum
strictly smaller than one; zero curvature is critical, and negative
curvature makes this test fail even when each conditional law remains
integrable.

The continuum scaling exposes a further bottleneck.  For the standard
gradient discretization, the one-update contraction gap is of order
\(a^2\) and vanishes as the lattice spacing \(a\to0\).  Diffusive
acceleration by order \(a^{-2}\) updates restores a finite limiting factor,
and a directionally weighted Green function retains physical exponential
decay.  Thus a continuum RDIC cannot ask a microscopic update for a
cutoff-independent gap; its shell kernel must already contain the
appropriate diffusive block evolution.

\subsection{A one-dimensional strong-convexity transport lemma}

Let
\begin{equation}
 d\nu_h(x)
 =
 Z_h^{-1}e^{-W(x)+hx}\,dx,
 \label{eq:v5v51-tilted-law}
\end{equation}
where \(W\in C^2(\mathbb R)\), \(W''\ge\lambda>0\), and \(Z_h<\infty\).

\begin{lemma}[One-dimensional Brascamp--Lieb inequality]
\label{lem:v5v51-BL}
For every smooth \(f\) with finite variance under \(\nu_h\),
\begin{equation}
 \operatorname{Var}_{\nu_h}(f)
 \le
 \int\frac{|f'(x)|^2}{W''(x)}\,d\nu_h(x)
 \le
 \frac1\lambda\int|f'(x)|^2\,d\nu_h(x).
 \label{eq:v5v51-BL}
\end{equation}
\end{lemma}

\begin{proof}
The linear tilt does not change the second derivative of the potential.
Let
\[
 L=\partial_x^2-(W'-h)\partial_x
\]
be the symmetric diffusion generator in \(L^2(\nu_h)\).  For a centred
smooth \(f\), solve first on a compact interval and then pass by
approximation to \(-Lu=f\).  Integration by parts and the one-dimensional
Bochner identity give
\begin{align*}
 \int f^2\,d\nu_h
 &=\int f'u'\,d\nu_h,\\
 \int(Lu)^2\,d\nu_h
 &=\int (u'')^2\,d\nu_h
 +\int W''(u')^2\,d\nu_h.
\end{align*}
Weighted Cauchy--Schwarz in the first identity yields
\[
 \int f^2
 \le
 \left(\int\frac{|f'|^2}{W''}\right)^{1/2}
 \left(\int W''|u'|^2\right)^{1/2}.
\]
The second identity bounds the last factor by
\((\int f^2)^{1/2}\).  Divide by it and approximate general \(f\).
The lower bound \(W''\ge\lambda\) gives the second inequality.
\end{proof}

\begin{theorem}[Linear-tilt Wasserstein Lipschitz bound]
\label{thm:v5v51-tilt-W1}
For all \(h,h'\in\mathbb R\),
\begin{equation}
 W_1(\nu_h,\nu_{h'})
 \le\frac{|h-h'|}{\lambda}.
 \label{eq:v5v51-tilt-W1}
\end{equation}
\end{theorem}

\begin{proof}
Let \(f\) be smooth and one-Lipschitz.  Differentiation under the integral
gives
\[
 \frac d{dh}\int f\,d\nu_h
 =
 \operatorname{Cov}_{\nu_h}(f,X).
\]
By Cauchy--Schwarz and
Lemma~\ref{lem:v5v51-BL},
\[
 |\operatorname{Cov}_{\nu_h}(f,X)|
 \le
 \sqrt{\operatorname{Var}(f)\operatorname{Var}(X)}
 \le\frac1\lambda,
\]
because \(|f'|\le1\) and \(X'=1\).  Integrate from \(h\) to \(h'\), then
take the supremum over one-Lipschitz \(f\) in the Kantorovich dual formula.
General Lipschitz \(f\) follow by smooth approximation.
\end{proof}

\subsection{The scalar conditional influence matrix}

Let \(G=(\Lambda,E)\) be a finite graph with degrees \(z_i\).  Consider the
energy
\begin{equation}
 S(\phi)
 =
 \sum_{i\in\Lambda}U_i(\phi_i)
 +\frac\kappa2\sum_{\{i,j\}\in E}(\phi_i-\phi_j)^2,
 \qquad \kappa\ge0.
 \label{eq:v5v51-lattice-action}
\end{equation}
Assume
\begin{equation}
 U_i''(x)\ge\mu_i
 \quad\text{for all }x,
 \qquad
 \lambda_i:=\mu_i+\kappa z_i>0.
 \label{eq:v5v51-site-curvature}
\end{equation}
The conditional law at \(i\), given the boundary values
\(\phi_{\partial i}\), has potential
\begin{equation}
 W_i(x)
 =
 U_i(x)+\frac{\kappa z_i}{2}x^2
 -\kappa x\sum_{j\sim i}\phi_j,
 \label{eq:v5v51-conditional-potential}
\end{equation}
up to a boundary-dependent constant.

\begin{theorem}[Exact strong-convexity influence card]
\label{thm:v5v51-site-influence}
Equip each scalar site with \(d_i(x,y)=|x-y|\).  Then the conditional
kernels satisfy
\begin{equation}
 W_1\bigl(
 \Gamma_i(\cdot\mid\phi),
 \Gamma_i(\cdot\mid\psi)
 \bigr)
 \le
 \sum_{j\sim i}
 \frac{\kappa}{\mu_i+\kappa z_i}
 |\phi_j-\psi_j|.
 \label{eq:v5v51-site-influence}
\end{equation}
Thus one may take
\begin{equation}
 c_{ij}
 =
 \begin{cases}
 \displaystyle\frac{\kappa}{\mu_i+\kappa z_i},
 &j\sim i,\\
 0,&j\not\sim i.
 \end{cases}
 \label{eq:v5v51-cij}
\end{equation}
\end{theorem}

\begin{proof}
The non-tilted part
\(U_i(x)+\kappa z_ix^2/2\) has second derivative at least
\(\lambda_i\).  The boundary enters only through the linear tilt
\[
 h_i(\phi)=\kappa\sum_{j\sim i}\phi_j.
\]
Apply Theorem~\ref{thm:v5v51-tilt-W1} and use
\[
 |h_i(\phi)-h_i(\psi)|
 \le\kappa\sum_{j\sim i}|\phi_j-\psi_j|.
\]
\end{proof}

\begin{corollary}[Positive, critical, and negative curvature]
\label{cor:v5v51-three-regimes}
The unweighted row sum is
\begin{equation}
 \theta_i
 =
 \sum_jc_{ij}
 =
 \frac{\kappa z_i}{\mu_i+\kappa z_i}.
 \label{eq:v5v51-row-sum}
\end{equation}
Hence:
\begin{enumerate}
 \item if \(\inf_i\mu_i>0\), then
 \(\sup_i\theta_i<1\) on graphs of uniformly bounded degree;
 \item if \(\mu_i=0\), then \(\theta_i=1\);
 \item if \(-\kappa z_i<\mu_i<0\), the conditional law is still strongly
 convex but \(\theta_i>1\), so the Dobrushin row test fails.
\end{enumerate}
\end{corollary}

\begin{proof}
All statements follow from
\eqref{eq:v5v51-row-sum}.
\end{proof}

\begin{assessment}[Higgs interpretation]
\label{ass:v5v51-Higgs}
For
\[
 U(\phi)=\lambda_H\phi^4+\frac{m_H^2}{2}\phi^2,
 \qquad\lambda_H>0,
\]
one has \(U''(\phi)=12\lambda_H\phi^2+m_H^2\).  The symmetric
positive-mass regime \(m_H^2>0\) satisfies the strict row test.  A broken
bare potential with \(m_H^2<0\) remains coercive but is not uniformly
convex, and the one-site Dobrushin argument does not cover it.  Block
convexification, phase conditioning, or a different metric is required.
\end{assessment}

\begin{firewall}[Coercivity is weaker than influence contraction]
\label{fw:v5v51-coercive-not-convex}
A positive quartic makes the density integrable as in V49 even when its
quadratic curvature is negative near the origin.  Integrability does not
imply \(\mu_i>0\), and therefore does not imply a Dobrushin row sum below
one.
\end{firewall}

\subsection{The continuum lattice scaling}

Let the spatial dimension be \(D\), the lattice spacing be \(a>0\), and
the nearest-neighbour degree be \(z=2D\).  The standard discretization of
\[
 \int\left(\frac12|\nabla\phi|^2+\frac{m^2}{2}\phi^2\right)\,dx
\]
has
\begin{equation}
 \kappa_a=a^{D-2},
 \qquad
 \mu_a=m^2a^D.
 \label{eq:v5v51-continuum-coefficients}
\end{equation}

\begin{theorem}[Microscopic contraction gap vanishes]
\label{thm:v5v51-gap-vanishes}
For the coefficients \eqref{eq:v5v51-continuum-coefficients}, the
Dobrushin row sum is
\begin{equation}
 \theta_a
 =
 \frac{z}{z+m^2a^2}
 =
 1-\frac{m^2}{z}a^2+O(a^4).
 \label{eq:v5v51-theta-a}
\end{equation}
Thus no cutoff-independent constant \(\theta<1\) bounds
\(\theta_a\) as \(a\downarrow0\).
\end{theorem}

\begin{proof}
Insert \(\kappa_a,\mu_a\) into
\eqref{eq:v5v51-row-sum} and cancel \(a^{D-2}\).  Expand
\((1+m^2a^2/z)^{-1}\).
\end{proof}

\begin{firewall}[A fixed-site RDIC is not continuum uniform]
\label{fw:v5v51-fixed-site}
The fixed-lattice strict inequality in
Corollary~\ref{cor:v5v51-three-regimes} cannot be inserted unchanged into
V50's cofinal certificate.  Its gap collapses at the canonical
diffusive rate \(a^2\).
\end{firewall}

\subsection{Diffusive acceleration}

\begin{theorem}[Finite contraction after diffusive time]
\label{thm:v5v51-diffusive-contraction}
Let \(T>0\) and let \(n(a)\) be integers with
\begin{equation}
 a^2n(a)\longrightarrow T.
 \label{eq:v5v51-diffusive-time}
\end{equation}
Then
\begin{equation}
 \theta_a^{\,n(a)}
 \longrightarrow
 \exp\!\left(-\frac{m^2T}{z}\right).
 \label{eq:v5v51-diffusive-limit}
\end{equation}
\end{theorem}

\begin{proof}
From \eqref{eq:v5v51-theta-a},
\[
 \log\theta_a
 =
 -\frac{m^2}{z}a^2+O(a^4).
\]
Multiply by \(n(a)\).  The leading term tends to
\(-m^2T/z\), while \(n(a)O(a^4)=O(a^2)\to0\).  Exponentiate.
\end{proof}

\begin{assessment}[Correct shell unit]
\label{ass:v5v51-shell-unit}
A continuum shell kernel should correspond to a fixed positive amount of
physical diffusion or block heat flow, not to one microscopic
single-site update.  Theorem~\ref{thm:v5v51-diffusive-contraction} gives a
nontrivial mass-dependent factor for that unit.  Implementing it with
overlapping gauge, Higgs, fermion, and metric blocks remains a separate
construction.
\end{assessment}

\subsection{Directional Green decay survives the limit}

For the translation-invariant quadratic model on \(\mathbb Z^D\), let
\[
 c_a=\frac1{z+m^2a^2},
 \qquad
 (C_au)_i=c_a\sum_{|e|=1}u_{i+e}.
\]
For \(h\in\mathbb R^D\), conjugation by \(u_i\mapsto e^{h\cdot i}u_i\)
has row sum
\begin{equation}
 \Theta_a(h)
 =
 \frac{2\sum_{\ell=1}^D\cosh h_\ell}
 {z+m^2a^2}.
 \label{eq:v5v51-directional-row}
\end{equation}

\begin{theorem}[Physical exponential Green bound]
\label{thm:v5v51-directional-Green}
Fix \(k\in\mathbb R^D\) with \(|k|<m\), and set \(h=ak\).  For all
sufficiently small \(a\), \(\Theta_a(h)<1\), and
\begin{equation}
 (I-C_a)^{-1}_{ij}
 \le
 \frac{\exp(-h\cdot(j-i))}
 {1-\Theta_a(h)}
 \label{eq:v5v51-directional-Green}
\end{equation}
whenever \(h\cdot(j-i)\ge0\).  Moreover,
\begin{equation}
 1-\Theta_a(ak)
 =
 \frac{m^2-|k|^2}{z}a^2+O(a^4).
 \label{eq:v5v51-directional-gap}
\end{equation}
Thus the decay exponent \(h\cdot(j-i)=k\cdot(a(j-i))\) is measured in
physical distance, while the prefactor has the expected resolvent scale
\(a^{-2}\).
\end{theorem}

\begin{proof}
Taylor expansion gives
\[
 2\sum_{\ell=1}^D\cosh(ak_\ell)
 =
 z+a^2|k|^2+O(a^4).
\]
Division by \(z+m^2a^2\) proves
\eqref{eq:v5v51-directional-gap} and strictness for \(|k|<m\).
The conjugated nonnegative matrix has \(\ell^\infty\) norm
\(\Theta_a(h)\).  Its Neumann series is bounded by
\((1-\Theta_a(h))^{-1}\).  Undoing the conjugation gives
\eqref{eq:v5v51-directional-Green}.
\end{proof}

\begin{assessment}[Why the absolute-distance row was too crude]
\label{ass:v5v51-directional}
The V50 weight \(e^{m d(i,j)}\) charges every nearest-neighbour step as
outward.  Directional conjugation pairs opposite steps and produces
\(\cosh h_\ell=1+h_\ell^2/2+O(h_\ell^4)\); the linear cost cancels.
That cancellation is essential for the physical massive continuum decay.
\end{assessment}

\subsection{Consequences for SM--Einstein sectors}

\begin{proposition}[Massive scalar fibre certificate]
\label{prop:v5v51-massive-scalar}
For a scalar sector satisfying a uniform positive conditional Hessian
bound and standard quadratic spatial coupling, the exact influence
coefficients are given by
\eqref{eq:v5v51-cij}.  A fixed-physical-time block update yields the
nontrivial contraction factor
\eqref{eq:v5v51-diffusive-limit}, and the linearized boundary response has
the physical directional decay
\eqref{eq:v5v51-directional-Green}.
\end{proposition}

\begin{proof}
Combine Theorems~\ref{thm:v5v51-site-influence},
\ref{thm:v5v51-diffusive-contraction}, and
\ref{thm:v5v51-directional-Green}.
\end{proof}

\begin{firewall}[Massless physical fields cannot be forced into this fibre]
\label{fw:v5v51-massless}
When \(m=0\), the row sum in
\eqref{eq:v5v51-theta-a} is exactly one and no positive directional
exponent satisfies \eqref{eq:v5v51-directional-gap}.  Photons, any
unproved Yang--Mills massless regime, and physical graviton directions
must remain in the retained/base or stratified long-range sector until a
separate gap or gauge-invariant decay theorem is proved.
\end{firewall}

\begin{firewall}[The conformal metric sign fails the convex test]
\label{fw:v5v51-conformal}
The Euclidean conformal metric direction identified in V8 has the wrong
quadratic sign before a valid physical contour or constraint quotient is
constructed.  It does not satisfy
\eqref{eq:v5v51-site-curvature}, so the scalar strong-convexity proof
cannot be transferred to it.
\end{firewall}

\begin{definition}[Convex diffusive influence certificate]
\label{def:v5v51-CDIC}
A CDIC for an unbounded matter fibre consists of:
\begin{enumerate}
 \item a cutoff-uniform conditional Hessian lower bound in the physical
 block metric;
 \item the resulting complete Wasserstein influence matrix;
 \item a shell update representing a fixed amount of physical block
 diffusion, with a nontrivial cofinal contraction factor;
 \item directional or spectral Green bounds at fixed physical distance;
 \item Lyapunov control of nonuniformly convex tails and complete
 cumulant/Ward assembly of cross-sector perturbations;
 \item exclusion or separate treatment of massless and wrong-sign
 physical directions.
\end{enumerate}
\end{definition}

\subsection{V51 status}

\begin{theorem}[V51 conclusion]
\label{thm:v5v51-status}
Strong conditional convexity gives an explicit Wasserstein influence
matrix for an unbounded scalar fibre.  Positive on-site curvature makes
the fixed-lattice row strict, while critical or negative curvature does
not.  Under continuum scaling, the microscopic gap vanishes as \(a^2\);
after diffusive acceleration it converges to a finite mass-dependent
factor, and directional conjugation retains physical exponential decay.
This proves a CDIC for the scalar quadratic model, not for the interacting
Higgs, gauge, fermion, or Einstein sectors.
\end{theorem}

\begin{proof}
Use Corollary~\ref{cor:v5v51-three-regimes},
Theorems~\ref{thm:v5v51-gap-vanishes},
\ref{thm:v5v51-diffusive-contraction}, and
\ref{thm:v5v51-directional-Green}.
\end{proof}

V52 will add a nonquadratic but uniformly convex on-site interaction and
a small cross-sector Hessian.  It will derive a block Schur influence
matrix, determine the exact dominance condition preserving the diffusive
mass gap, and isolate the first place where fermionic determinants or
Einstein constraint reduction can violate positivity.

\clearpage
\section{V52: block strong convexity and physical-mass-scale dominance}
\label{sec:v5v52}

V51 proved an exact conditional influence estimate for one scalar species
and found that the microscopic contraction gap is only of order \(a^2\).
This note adds finitely many coupled species and a nonquadratic uniformly
convex on-site interaction.  The central point is a scale distinction.
Cross-sector Hessians must be small relative to the physical on-site mass
matrix, whose lattice coefficient is \(a^D\).  Smallness only relative to
the full conditional Hessian, dominated by the gradient term
\(a^{D-2}\), does not preserve the continuum mass gap.

The weighted Schur theorem of V10 is used rather than reproved.  Its
application gives a positive block mass matrix and hence a vector-valued
Wasserstein influence coefficient.  Integrating out a sector subtracts a
conditional covariance from the retained Hessian, as computed in V36;
that subtraction must obey the same physical-mass-scale dominance.
Fermionic determinants are favourable only when their Hessian is proved
positive on the chosen real physical slice.

\subsection{Multivariate linear-tilt transport}

Let \(M\) be a positive-definite matrix on \(\mathbb R^r\), and let
\begin{equation}
 d\nu_b(x)
 =
 Z_b^{-1}e^{-W(x)+\langle b,x\rangle}\,dx,
 \qquad
 \nabla^2W(x)\ge M.
 \label{eq:v5v52-vector-tilt}
\end{equation}
Equip \(\mathbb R^r\) with the distance
\begin{equation}
 d_M(x,y)=|M^{1/2}(x-y)|.
 \label{eq:v5v52-M-distance}
\end{equation}

\begin{theorem}[Matrix strong-convexity transport]
\label{thm:v5v52-matrix-transport}
For all \(b,b'\in\mathbb R^r\),
\begin{equation}
 W_{2,M}(\nu_b,\nu_{b'})
 \le
 |M^{-1/2}(b-b')|,
 \label{eq:v5v52-matrix-transport}
\end{equation}
where \(W_{2,M}\) is the quadratic Wasserstein distance induced by
\eqref{eq:v5v52-M-distance}.  The same bound holds for \(W_{1,M}\).
\end{theorem}

\begin{proof}
Set \(y=M^{1/2}x\) and
\(\widetilde W(y)=W(M^{-1/2}y)\).  Then
\(\nabla^2\widetilde W\ge I\), while the tilt is
\(\widetilde b=M^{-1/2}b\).
Run the two overdamped Langevin diffusions with tilts
\(\widetilde b,\widetilde b'\) using the same Brownian motion.  Their
difference \(D_t\) obeys
\[
 \frac d{dt}|D_t|
 \le-|D_t|+|\widetilde b-\widetilde b'|
\]
whenever \(D_t\ne0\), by strong monotonicity of
\(\nabla\widetilde W\).  At stationarity this synchronous coupling gives
\[
 \bigl(\mathbb E|D|^2\bigr)^{1/2}
 \le|\widetilde b-\widetilde b'|.
\]
Transform back to \(x\).  Since \(W_1\le W_2\), the last assertion
follows.  Standard truncation of \(W\) and passage to the stationary
limit justify the coupling when its derivatives are unbounded.
\end{proof}

\subsection{Vector lattice influence}

Let \(\phi_i\in\mathbb R^r\), let \(K>0\) be a kinetic matrix, and consider
\begin{equation}
 S_a(\phi)
 =
 \sum_i U_{i,a}(\phi_i)
 +\frac{a^{D-2}}2
 \sum_{\{i,j\}\in E}
 \langle\phi_i-\phi_j,K(\phi_i-\phi_j)\rangle.
 \label{eq:v5v52-vector-action}
\end{equation}
Assume
\begin{equation}
 \nabla^2U_{i,a}(x)\ge a^D M_i^2
 \label{eq:v5v52-onsite-mass}
\end{equation}
for symmetric matrices \(M_i^2\).  At a site of degree \(z_i\), define
\begin{equation}
 H_{i,a}
 =
 z_i a^{D-2}K+a^DM_i^2.
 \label{eq:v5v52-conditional-Hessian}
\end{equation}

\begin{theorem}[Complete vector influence card]
\label{thm:v5v52-vector-influence}
If \(H_{i,a}>0\), then in the conditional metric induced by \(H_{i,a}\),
the influence of a neighbour \(j\sim i\) is bounded by the operator
\begin{equation}
 \mathcal C_{ij,a}
 =
 H_{i,a}^{-1/2}a^{D-2}K\,R_{j,a}^{-1/2},
 \label{eq:v5v52-vector-cij}
\end{equation}
when the boundary block \(j\) is measured in the metric \(R_{j,a}>0\).
In particular,
\begin{equation}
 W_{1,H_{i,a}}\bigl(
 \Gamma_i(\cdot\mid\phi),
 \Gamma_i(\cdot\mid\psi)
 \bigr)
 \le
 \sum_{j\sim i}
 \|\mathcal C_{ij,a}\|\,
 |\phi_j-\psi_j|_{R_{j,a}}.
 \label{eq:v5v52-vector-influence}
\end{equation}
\end{theorem}

\begin{proof}
The boundary values enter the site action through the linear tilt
\[
 b_i(\phi)=a^{D-2}K\sum_{j\sim i}\phi_j.
\]
Apply Theorem~\ref{thm:v5v52-matrix-transport} with
\(M=H_{i,a}\), then insert
\(\phi_j-\psi_j=R_{j,a}^{-1/2}
 R_{j,a}^{1/2}(\phi_j-\psi_j)\) and take operator norms.
\end{proof}

For a translation-invariant system, use the kinetic metric \(K\) at every
boundary site and set
\begin{equation}
 \mathcal M=K^{-1/2}M^2K^{-1/2}.
 \label{eq:v5v52-kinetic-mass}
\end{equation}

\begin{theorem}[Vector microscopic row and diffusive limit]
\label{thm:v5v52-vector-row}
Suppose \(\mathcal M\ge m_*^2I\) with \(m_*^2>0\).  The exact Gaussian
conditional-mean derivative has per-neighbour matrix
\begin{equation}
 C_a=(zI+a^2\mathcal M)^{-1},
 \label{eq:v5v52-exact-vector-C}
\end{equation}
in kinetic coordinates.  Hence
\begin{equation}
 \|zC_a\|
 \le
 \frac{z}{z+a^2m_*^2}
 =
 1-\frac{m_*^2}{z}a^2+O(a^4).
 \label{eq:v5v52-vector-row}
\end{equation}
For \(a^2n(a)\to T\),
\begin{equation}
 \|(zC_a)^{n(a)}\|
 \le
 \left(\frac{z}{z+a^2m_*^2}\right)^{n(a)}
 \longrightarrow e^{-m_*^2T/z}.
 \label{eq:v5v52-vector-diffusive}
\end{equation}
\end{theorem}

\begin{proof}
Factor
\[
 H_a
 =
 a^{D-2}K^{1/2}(zI+a^2\mathcal M)K^{1/2}.
\]
The derivative of the conditional mean with respect to one neighbour is
\(H_a^{-1}a^{D-2}K\), which is similar in kinetic coordinates to
\eqref{eq:v5v52-exact-vector-C}.  Functional calculus and
\(\mathcal M\ge m_*^2I\) give
\eqref{eq:v5v52-vector-row}.  The last limit is the scalar calculation of
V51 applied to the smallest spectral value.
\end{proof}

\begin{corollary}[Uniformly convex nonquadratic on-site terms]
\label{cor:v5v52-nonquadratic}
If
\[
 \nabla^2U_{i,a}(x)\ge a^DM^2
\]
uniformly in \(x,i,a\), then the same upper influence bound
\eqref{eq:v5v52-vector-row} holds for the nonquadratic conditional
kernels.
\end{corollary}

\begin{proof}
Theorem~\ref{thm:v5v52-matrix-transport} uses only the Hessian lower
bound.  Replacing the exact quadratic Hessian by a larger
field-dependent one can only improve the stated upper card.
\end{proof}

\subsection{Cross-sector Schur dominance}

Split the physical species space into two blocks, for example Higgs/matter
and a sliced gauge/metric block.  Let the physical on-site mass matrix be
\begin{equation}
 M_{\mathrm{phys}}^2
 =
 \begin{pmatrix}
  A & B\\
  B^* & D
 \end{pmatrix},
 \qquad A>0,\quad D>0.
 \label{eq:v5v52-mass-block}
\end{equation}
Put
\begin{equation}
 q=\|A^{-1/2}BD^{-1/2}\|.
 \label{eq:v5v52-cross-q}
\end{equation}

\begin{theorem}[Physical block-mass lower bound]
\label{thm:v5v52-block-mass}
If \(q<1\), then
\begin{equation}
 M_{\mathrm{phys}}^2
 \ge
 (1-q)\begin{pmatrix}A&0\\0&D\end{pmatrix}.
 \label{eq:v5v52-block-mass-lower}
\end{equation}
Consequently, if
\begin{equation}
 K^{-1/2}
 \begin{pmatrix}A&0\\0&D\end{pmatrix}
 K^{-1/2}
 \ge m_0^2I,
 \label{eq:v5v52-diagonal-mass}
\end{equation}
then Theorem~\ref{thm:v5v52-vector-row} holds with
\begin{equation}
 m_*^2=(1-q)m_0^2.
 \label{eq:v5v52-effective-mass}
\end{equation}
\end{theorem}

\begin{proof}
Apply Theorem~\ref{thm:v5v10-weighted-Schur} to the diagonal operator
\(\operatorname{diag}(A,D)\) and the off-diagonal operator with entries
\(B,B^*\).  This gives
\eqref{eq:v5v52-block-mass-lower}.  Conjugate by \(K^{-1/2}\), use
\eqref{eq:v5v52-diagonal-mass}, and apply
Theorem~\ref{thm:v5v52-vector-row}.
\end{proof}

\begin{assessment}[What the Schur number measures]
\label{ass:v5v52-Schur-number}
The dimensionless quantity \(q\) compares the cross-Hessian with both
physical diagonal mass blocks.  It is not a comparison with the
nearest-neighbour kinetic Hessian.  This is the two-sided weighting
required by V10 and is the scale that survives the continuum limit.
\end{assessment}

\subsection{The physical-mass-scale necessity}

Suppose a scalar or smallest block eigenvalue has effective on-site
coefficient
\begin{equation}
 \mu_{\mathrm{eff},a}
 =
 m^2a^D-\Delta_a.
 \label{eq:v5v52-effective-onsite}
\end{equation}

\begin{theorem}[Continuum mass dominance criterion]
\label{thm:v5v52-mass-dominance}
For the standard kinetic coefficient \(\kappa_a=a^{D-2}\), the conditional
row is
\begin{equation}
 \theta_a
 =
 \frac{z}{z+a^2m^2-a^{2-D}\Delta_a}.
 \label{eq:v5v52-corrected-row}
\end{equation}
If
\begin{equation}
 a^{-D}\Delta_a\longrightarrow\delta
 \quad\text{with }\delta<m^2,
 \label{eq:v5v52-physical-correction}
\end{equation}
then
\begin{equation}
 \theta_a
 =
 1-\frac{m^2-\delta}{z}a^2+o(a^2),
 \label{eq:v5v52-renormalized-gap}
\end{equation}
and diffusive acceleration yields
\(\exp(-(m^2-\delta)T/z)\).
By contrast, a bound
\begin{equation}
 |\Delta_a|\le\varepsilon a^{D-2}
 \label{eq:v5v52-kinetic-small}
\end{equation}
with fixed \(\varepsilon>0\) does not preserve the mass gap: it allows
\(a^{2-D}\Delta_a=O(1)\), which is larger than \(a^2m^2\).
\end{theorem}

\begin{proof}
Insert \eqref{eq:v5v52-effective-onsite} into
\[
 \theta_a
 =
 \frac{z\kappa_a}{z\kappa_a+\mu_{\mathrm{eff},a}}
\]
and divide by \(a^{D-2}\), proving
\eqref{eq:v5v52-corrected-row}.  Under
\eqref{eq:v5v52-physical-correction}, its added denominator is
\((m^2-\delta)a^2+o(a^2)\); expansion gives
\eqref{eq:v5v52-renormalized-gap} and the V51 logarithmic limit.
Condition \eqref{eq:v5v52-kinetic-small} controls
\(\Delta_a/a^{D-2}\), not \(\Delta_a/a^D\), so it permits an order-one
change in the denominator after kinetic rescaling.
\end{proof}

\begin{firewall}[Conditional-Hessian smallness can be two powers too weak]
\label{fw:v5v52-two-powers}
The conditional Hessian is of order \(a^{D-2}\), while the physical mass
term is of order \(a^D\).  A perturbation small by a fixed percentage of
the former can overwhelm the latter as \(a\to0\).  Every determinant,
constraint, and cross-sector correction must therefore be renormalized
and estimated at the \(a^D\) mass scale before claiming a continuum
influence gap.
\end{firewall}

\subsection{Integrating out a sector}

Let \(p\) denote retained fields and \(q\) integrated fields.  At fixed
reference covariance, V36 gives
\begin{equation}
 D^2V_{\mathrm{red}}(p)
 =
 \left\langle D_{pp}^2V\right\rangle_{p,V}
 -\operatorname{Cov}_{p,V}(D_pV,D_pV).
 \label{eq:v5v52-reduced-Hessian}
\end{equation}
Write the two positive forms as \(A_p\) and \(R_p\), respectively.

\begin{theorem}[Variance-debit preservation of mass]
\label{thm:v5v52-variance-debit}
Assume
\begin{equation}
 A_p\ge a^DM_0^2
 \label{eq:v5v52-direct-mass}
\end{equation}
and, as quadratic forms,
\begin{equation}
 0\le R_p\le rA_p,
 \qquad 0\le r<1.
 \label{eq:v5v52-variance-relative}
\end{equation}
Then
\begin{equation}
 D^2V_{\mathrm{red}}(p)
 \ge(1-r)A_p
 \ge a^D(1-r)M_0^2.
 \label{eq:v5v52-reduced-mass}
\end{equation}
This retained mass can be inserted into
Theorem~\ref{thm:v5v52-vector-row}.
\end{theorem}

\begin{proof}
Subtract \eqref{eq:v5v52-variance-relative} from \(A_p\) in
\eqref{eq:v5v52-reduced-Hessian}.  The second inequality is
\eqref{eq:v5v52-direct-mass}.
\end{proof}

\begin{firewall}[The relative form bound must be proved at physical scale]
\label{fw:v5v52-variance-proof}
The covariance term in \eqref{eq:v5v52-reduced-Hessian} is nonnegative
and can cancel the direct curvature.  A rooted covariance estimate which
is merely finite, or small relative to the kinetic block, does not imply
\eqref{eq:v5v52-variance-relative}.  The complete conditional source
tensor must supply the strict relative factor \(r<1\).
\end{firewall}

\subsection{Fermionic determinant sign}

Let \(F(x)\) be an invertible finite-dimensional fermion matrix and put
\begin{equation}
 \Gamma_F(x)=-\log\det F(x)
 \label{eq:v5v52-fermion-effective}
\end{equation}
on a branch where this expression is real and differentiable.

\begin{lemma}[Exact determinant Hessian]
\label{lem:v5v52-det-Hessian}
For directions \(h,k\),
\begin{align}
 D^2\Gamma_F[h,k]
 &=
 \operatorname{Tr}\!\left(
 F^{-1}DF[k]F^{-1}DF[h]\right)
 -
 \operatorname{Tr}\!\left(F^{-1}D^2F[h,k]\right).
 \label{eq:v5v52-det-Hessian}
\end{align}
\end{lemma}

\begin{proof}
Differentiate
\(D\log\det F[h]=\operatorname{Tr}(F^{-1}DF[h])\) and use
\(D(F^{-1})[k]=-F^{-1}DF[k]F^{-1}\).
\end{proof}

\begin{corollary}[Positive affine Hermitian case]
\label{cor:v5v52-positive-det}
If \(F(x)>0\) is Hermitian and affine in \(x\), then
\begin{equation}
 D^2\Gamma_F[h,h]
 =
 \operatorname{Tr}\!\left(
 (F^{-1/2}DF[h]F^{-1/2})^2\right)\ge0.
 \label{eq:v5v52-positive-det}
\end{equation}
\end{corollary}

\begin{proof}
The second term in \eqref{eq:v5v52-det-Hessian} vanishes.  Cyclicity of
the trace gives the displayed square.
\end{proof}

\begin{firewall}[The physical fermion determinant need not satisfy the
positive case]
\label{fw:v5v52-fermion-sign}
Chiral gauge operators, complex phases, zero crossings, gauge fixing, and
nonaffine field dependence can invalidate the hypotheses of
Corollary~\ref{cor:v5v52-positive-det}.  In those cases the determinant
Hessian must enter the signed physical mass block and satisfy
Theorem~\ref{thm:v5v52-mass-dominance}; it cannot be declared convex from
the formal trace square alone.
\end{firewall}

\subsection{Block convex diffusive certificate}

\begin{definition}[Block convex diffusive certificate]
\label{def:v5v52-BCDC}
A BCDC consists of:
\begin{enumerate}
 \item a positive physical kinetic matrix \(K\) after gauge and constraint
 reduction;
 \item a renormalized physical mass matrix with a uniform kinetic-normalized
 lower bound \(m_*^2>0\) on the declared massive fibre;
 \item two-sided Schur control of every cross-sector Hessian at the
 physical mass scale;
 \item the variance-debit estimate
 \eqref{eq:v5v52-variance-relative} for every integrated sector;
 \item a sign-resolved determinant Hessian and local counterterm
 allocation;
 \item a fixed-physical-time block update and complete conditional
 influence assembly;
 \item separate retained treatment of massless and wrong-sign directions.
\end{enumerate}
\end{definition}

\subsection{V52 status}

\begin{theorem}[V52 conclusion]
\label{thm:v5v52-status}
Uniform block strong convexity yields a vector-valued conditional
influence estimate and a diffusive contraction governed by the smallest
kinetic-normalized physical mass.  V10 Schur control preserves that mass
under small cross-Hessians, while V36 reduction preserves it only when the
subtracted conditional covariance has strict relative size below one.
All corrections must be controlled at order \(a^D\); smallness relative
to the order-\(a^{D-2}\) conditional kinetic Hessian is insufficient by
two lattice powers.  A BCDC is proved for the stated convex block model,
but not for the chiral fermion determinant, broken Higgs phase, or
Einstein constraint sector.
\end{theorem}

\begin{proof}
Use Theorems~\ref{thm:v5v52-vector-row},
\ref{thm:v5v52-block-mass},
\ref{thm:v5v52-mass-dominance}, and
\ref{thm:v5v52-variance-debit}.
\end{proof}

V53 will attack the variance debit
\eqref{eq:v5v52-variance-relative}.  It will derive a Brascamp--Lieb
upper bound for the conditional covariance of cross-gradients, express
the result as a Schur complement of the joint Hessian, and determine
whether convex integration can preserve the retained physical mass
without an independent small-coupling assumption.

\clearpage
\section{V53: conditional Brascamp--Lieb and the reduced Hessian Schur complement}
\label{sec:v5v53}

V52 isolated the variance debit in the Hessian obtained after integrating
out a sector.  This note controls that debit for a positive bosonic
conditional measure.  Conditional Brascamp--Lieb bounds the covariance of
the retained-field gradient by the inverse integrated-field Hessian.  When
inserted into the exact V36 reduced-Hessian formula, the result is the
pointwise Schur complement of the joint action Hessian.

This gives a quantitative form of convexity preservation under bosonic
marginalization.  An independent small-coupling hypothesis is unnecessary
if the full joint Hessian already has a uniform Schur mass gap.  The
argument does not apply to a signed or complex fermion/ghost measure, to a
wrong-sign Einstein contour, or to a retained-field-dependent reference
measure until its extra covariance-derivative terms are included.

\subsection{Conditional matrix Brascamp--Lieb}

Let \(p\in P\simeq\mathbb R^r\) be retained and
\(q\in Q\simeq\mathbb R^s\) be integrated.  Let
\begin{equation}
 d\nu_p(q)
 =
 Z(p)^{-1}e^{-V(p,q)}\,dq,
 \label{eq:v5v53-conditional-law}
\end{equation}
where \(V\in C^2\), \(Z(p)<\infty\), and
\begin{equation}
 D_{qq}^2V(p,q)>0.
 \label{eq:v5v53-q-convexity}
\end{equation}

\begin{theorem}[Conditional Brascamp--Lieb inequality]
\label{thm:v5v53-conditional-BL}
For every smooth scalar \(f:Q\to\mathbb R\) with finite variance,
\begin{equation}
 \operatorname{Var}_{\nu_p}(f)
 \le
 \int_Q
 \left\langle
 D_qf,
 (D_{qq}^2V)^{-1}D_qf
 \right\rangle
 d\nu_p.
 \label{eq:v5v53-conditional-BL}
\end{equation}
\end{theorem}

\begin{proof}
Fix \(p\) and write
\[
 L=\Delta_q-\langle D_qV,D_q\rangle.
\]
For centred smooth \(f\), solve \(-Lu=f\) first on bounded convex domains
with reflecting boundary and pass to the full space by coercive
approximation.  Integration by parts gives
\[
 \int f^2\,d\nu_p=\int\langle D_qf,D_qu\rangle\,d\nu_p.
\]
The Bochner identity is
\[
 \int(Lu)^2\,d\nu_p
 =
 \int\|D_q^2u\|_{\mathrm{HS}}^2\,d\nu_p
 +
 \int\langle D_qu,(D_{qq}^2V)D_qu\rangle\,d\nu_p.
\]
Apply Cauchy--Schwarz with the positive matrices
\((D_{qq}^2V)^{-1}\) and \(D_{qq}^2V\) to the first display.  The second
factor is at most \((\int f^2d\nu_p)^{1/2}\) by the Bochner identity.
Divide by it.  Approximation proves the asserted domain extension.
\end{proof}

\subsection{The reduced Hessian}

Define
\begin{equation}
 V_{\mathrm{red}}(p)
 =
 -\log\int_Qe^{-V(p,q)}\,dq.
 \label{eq:v5v53-reduced-potential}
\end{equation}
At a fixed reference measure, V36 proved
\begin{equation}
 D^2V_{\mathrm{red}}(p)[h,h]
 =
 \left\langle D_{pp}^2V[h,h]\right\rangle_{\nu_p}
 -
 \operatorname{Var}_{\nu_p}(D_pV[h]).
 \label{eq:v5v53-V36-Hessian}
\end{equation}
Introduce the pointwise blocks
\begin{equation}
 A=D_{pp}^2V,
 \qquad
 B=D_{qp}^2V,
 \qquad
 D=D_{qq}^2V.
 \label{eq:v5v53-Hessian-blocks}
\end{equation}

\begin{theorem}[Reduced Schur-complement lower bound]
\label{thm:v5v53-reduced-Schur}
Under \eqref{eq:v5v53-q-convexity},
\begin{equation}
 D^2V_{\mathrm{red}}(p)
 \ge
 \int_Q
 \left(A-B^*D^{-1}B\right)(p,q)\,d\nu_p(q)
 \label{eq:v5v53-reduced-Schur}
\end{equation}
as quadratic forms on \(P\).
\end{theorem}

\begin{proof}
For fixed \(h\in P\), apply
Theorem~\ref{thm:v5v53-conditional-BL} to
\[
 f_h(q)=D_pV(p,q)[h].
\]
Its \(q\)-gradient is
\[
 D_qf_h=B h.
\]
Therefore
\[
 \operatorname{Var}_{\nu_p}(D_pV[h])
 \le
 \int\langle Bh,D^{-1}Bh\rangle\,d\nu_p.
\]
Insert this into \eqref{eq:v5v53-V36-Hessian}.
\end{proof}

\begin{corollary}[Quantitative convex marginalization]
\label{cor:v5v53-convex-marginal}
If a fixed positive form \(M\) on \(P\) satisfies
\begin{equation}
 A-B^*D^{-1}B\ge M
 \quad\text{for every }(p,q),
 \label{eq:v5v53-pointwise-Schur-gap}
\end{equation}
then
\begin{equation}
 D^2V_{\mathrm{red}}(p)\ge M
 \quad\text{for every }p.
 \label{eq:v5v53-reduced-gap}
\end{equation}
\end{corollary}

\begin{proof}
Average \eqref{eq:v5v53-pointwise-Schur-gap} in
\eqref{eq:v5v53-reduced-Schur}.
\end{proof}

\begin{corollary}[Joint Hessian formulation]
\label{cor:v5v53-joint-Hessian}
Suppose
\begin{equation}
 \begin{pmatrix}
  A&B^*\\
  B&D
 \end{pmatrix}
 \ge
 \begin{pmatrix}
  M&0\\
  0&0
 \end{pmatrix},
 \qquad D>0.
 \label{eq:v5v53-joint-gap}
\end{equation}
Then \eqref{eq:v5v53-reduced-gap} holds.
\end{corollary}

\begin{proof}
Test \eqref{eq:v5v53-joint-gap} on
\((h,-D^{-1}Bh)\).  The resulting quadratic form is
\[
 \langle h,(A-B^*D^{-1}B)h\rangle
 \ge\langle h,Mh\rangle.
\]
Apply Corollary~\ref{cor:v5v53-convex-marginal}.
\end{proof}

\begin{assessment}[No separate perturbative smallness is needed]
\label{ass:v5v53-no-small-coupling}
The cross block \(B\) may be large in its unweighted norm.  What matters
is the positive joint Schur gap
\eqref{eq:v5v53-pointwise-Schur-gap}.  This is the nonlinear conditional
version of V10's two-sided weighted Schur test.
\end{assessment}

\subsection{Gaussian saturation}

\begin{proposition}[The Schur bound is sharp]
\label{prop:v5v53-Gaussian-sharp}
Let
\begin{equation}
 V(p,q)
 =
 \frac12\langle p,Ap\rangle
 +\langle q,Bp\rangle
 +\frac12\langle q,Dq\rangle,
 \qquad D>0.
 \label{eq:v5v53-Gaussian-joint}
\end{equation}
Then, up to an additive constant,
\begin{equation}
 V_{\mathrm{red}}(p)
 =
 \frac12
 \langle p,(A-B^*D^{-1}B)p\rangle.
 \label{eq:v5v53-Gaussian-reduced}
\end{equation}
Equality holds in
\eqref{eq:v5v53-reduced-Schur}.
\end{proposition}

\begin{proof}
Complete the square:
\[
 \frac12\langle q,Dq\rangle+\langle q,Bp\rangle
 =
 \frac12\langle q+D^{-1}Bp,D(q+D^{-1}Bp)\rangle
 -
 \frac12\langle p,B^*D^{-1}Bp\rangle.
\]
The shifted Gaussian integral is independent of \(p\), yielding
\eqref{eq:v5v53-Gaussian-reduced}.  The Gaussian Brascamp--Lieb
inequality is an equality for linear functions.
\end{proof}

\begin{firewall}[Direct retained curvature is not the reduced mass]
\label{fw:v5v53-direct-not-reduced}
The positive block \(A\) alone does not survive integration.  The exact
Gaussian model subtracts \(B^*D^{-1}B\), and no estimate can generally
replace the Schur complement by \(A\).
\end{firewall}

\subsection{Relative variance debit}

The V52 variance condition can now be verified from pointwise Hessian
data.

\begin{theorem}[Pointwise debit implies the V52 relative bound]
\label{thm:v5v53-relative-debit}
Assume \(A(p,q)\ge A_0>0\), \(D(p,q)>0\), and
\begin{equation}
 \left\|
 A(p,q)^{-1/2}B(p,q)^*D(p,q)^{-1}
 B(p,q)A(p,q)^{-1/2}
 \right\|
 \le r<1
 \label{eq:v5v53-relative-cross}
\end{equation}
uniformly.  Then
\begin{equation}
 D^2V_{\mathrm{red}}(p)
 \ge
 (1-r)\int A(p,q)\,d\nu_p(q)
 \ge(1-r)A_0.
 \label{eq:v5v53-relative-reduced}
\end{equation}
\end{theorem}

\begin{proof}
Condition \eqref{eq:v5v53-relative-cross} is equivalent to
\[
 B^*D^{-1}B\le rA.
\]
Insert this in
Theorem~\ref{thm:v5v53-reduced-Schur} and average.  The lower bound by
\(A_0\) is preserved under averaging.
\end{proof}

\begin{corollary}[Physical lattice mass survives convex integration]
\label{cor:v5v53-physical-mass}
If
\begin{equation}
 A_0\ge a^DM_0^2
 \label{eq:v5v53-A0-mass}
\end{equation}
and \eqref{eq:v5v53-relative-cross} holds with cutoff-independent
\(r<1\), then the retained mass matrix is at least
\begin{equation}
 a^D(1-r)M_0^2.
 \label{eq:v5v53-retained-physical-mass}
\end{equation}
The V52 diffusive influence theorem applies with the corresponding
kinetic-normalized smallest eigenvalue.
\end{corollary}

\begin{proof}
Combine \eqref{eq:v5v53-relative-reduced} and
\eqref{eq:v5v53-A0-mass}, then apply
Theorem~\ref{thm:v5v52-vector-row}.
\end{proof}

\begin{firewall}[The cross ratio must be uniform in the field]
\label{fw:v5v53-uniform-field}
An estimate of \eqref{eq:v5v53-relative-cross} only near the minimum of
the action does not control the full conditional covariance.  Large-field
regions enter the Brascamp--Lieb average.  Uniform convexity, a Lyapunov
localization plus a tail estimate, or a phase-restricted theorem is
required.
\end{firewall}

\subsection{A bounded nonconvex perturbation}

The globally convex model can be enlarged by a bounded oscillatory
perturbation at a quantitative cost.

\begin{lemma}[Coarse Holley--Stroock comparison]
\label{lem:v5v53-HS}
Let \(\mu_0\) satisfy
\begin{equation}
 \operatorname{Var}_{\mu_0}(f)
 \le C_0\int|\nabla f|^2\,d\mu_0.
 \label{eq:v5v53-Poincare0}
\end{equation}
Let
\[
 d\mu=Z^{-1}e^{-R}\,d\mu_0
\]
with \(\operatorname{osc}(R)\le\omega<\infty\).  Then
\begin{equation}
 \operatorname{Var}_{\mu}(f)
 \le
 e^{2\omega}C_0
 \int|\nabla f|^2\,d\mu.
 \label{eq:v5v53-HS}
\end{equation}
\end{lemma}

\begin{proof}
The normalized density ratio \(r=d\mu/d\mu_0\) satisfies
\[
 e^{-\omega}\le r\le e^\omega.
\]
Using the \(\mu_0\)-mean as a trial constant,
\begin{align*}
 \operatorname{Var}_{\mu}(f)
 &\le
 \int|f-\mu_0f|^2\,d\mu\\
 &\le
 e^\omega\operatorname{Var}_{\mu_0}(f)\\
 &\le
 e^\omega C_0\int|\nabla f|^2\,d\mu_0\\
 &\le
 e^{2\omega}C_0\int|\nabla f|^2\,d\mu.
\end{align*}
\end{proof}

\begin{proposition}[Perturbed variance debit]
\label{prop:v5v53-perturbed-debit}
If the conditional convex reference has Hessian \(D_0>0\), the actual
conditional measure differs by a perturbation of oscillation at most
\(\omega\), and
\[
 \|D_q(D_pV[h])\|_{D_0^{-1}}^2
 \le r\,A[h,h]
\]
pointwise, then its variance debit is at most
\begin{equation}
 e^{2\omega}r\,\langle A[h,h]\rangle.
 \label{eq:v5v53-perturbed-debit}
\end{equation}
The retained mass stays positive when
\begin{equation}
 e^{2\omega}r<1.
 \label{eq:v5v53-perturbed-smallness}
\end{equation}
\end{proposition}

\begin{proof}
Apply Lemma~\ref{lem:v5v53-HS} to
\(f=D_pV[h]\), using the Brascamp--Lieb Poincaré constant of the convex
reference in the \(D_0\) metric.  Insert the resulting variance estimate
into \eqref{eq:v5v53-V36-Hessian}.
\end{proof}

\begin{assessment}[Broken phases remain outside the bounded perturbation]
\label{ass:v5v53-broken}
A negative quadratic perturbation of a positive quartic has unbounded
oscillation on \(\mathbb R^s\).  Lemma~\ref{lem:v5v53-HS} therefore does
not prove a broken-Higgs conditional gap.  One needs a phase-localized
block measure, a two-scale Lyapunov/Poincaré theorem, or convexification by
the block boundary conditions.
\end{assessment}

\subsection{Signed and moving-reference obstructions}

\begin{firewall}[Fermion and ghost integration is not a positive marginal]
\label{fw:v5v53-signed}
The proof of Theorem~\ref{thm:v5v53-conditional-BL} uses a positive
probability density and variance positivity.  A complex chiral determinant
or a signed ghost representation does not satisfy those hypotheses.
After exact determinant and BRST assembly, one must either obtain a
positive physical density or estimate its Hessian directly in the Schur
mass block.
\end{firewall}

\begin{firewall}[A moving reference adds terms]
\label{fw:v5v53-moving-reference}
If the conditional Gaussian covariance depends on \(p\), the derivative
formula contains the trace insertion derived in V36.  Then
\eqref{eq:v5v53-V36-Hessian} is incomplete unless the reference derivative
is absorbed into \(V\) relative to a fixed dominating measure or retained
as an additional complete block.  The fixed-reference Schur theorem cannot
be quoted without that conversion.
\end{firewall}

\begin{firewall}[Convexity still does not imply OS positivity]
\label{fw:v5v53-convex-not-OS}
The reduced Hessian bound controls Euclidean coercivity and conditional
mixing.  As V10 proved, reflection positivity is a separate sign condition
on the time-reflected inverse covariance and interacting observable
algebra.
\end{firewall}

\subsection{Conditional Schur convexity certificate}

\begin{definition}[Conditional Schur convexity certificate]
\label{def:v5v53-CSCC}
A CSCC consists of:
\begin{enumerate}
 \item a positive conditional bosonic density relative to a fixed
 dominating measure;
 \item strict integrated-field Hessian positivity \(D>0\);
 \item a pointwise physical-mass Schur gap
 \(A-B^*D^{-1}B\ge a^DM_*^2\);
 \item cutoff-uniform control of unbounded tails or a bounded perturbation
 satisfying \eqref{eq:v5v53-perturbed-smallness};
 \item explicit treatment of determinant, ghost, and moving-reference
 terms outside the positive variance argument;
 \item insertion of the retained mass into the BCDC and RDIC.
\end{enumerate}
\end{definition}

\subsection{V53 status}

\begin{theorem}[V53 conclusion]
\label{thm:v5v53-status}
For a positive conditionally log-concave bosonic sector, the covariance
debit in the reduced Hessian is bounded by
\(B^*D^{-1}B\), and a pointwise joint Schur mass gap passes unchanged to
the retained potential.  The estimate is exact for a joint Gaussian and
preserves the physical \(a^D\) mass scale under the uniform relative
condition \eqref{eq:v5v53-relative-cross}.  The result does not cover
broken nonconvex phases, signed or complex fermionic/ghost weights,
moving-reference trace terms, or the wrong-sign Einstein conformal
sector.
\end{theorem}

\begin{proof}
Use Theorem~\ref{thm:v5v53-reduced-Schur},
Corollary~\ref{cor:v5v53-convex-marginal},
Proposition~\ref{prop:v5v53-Gaussian-sharp}, and
Corollary~\ref{cor:v5v53-physical-mass}.
\end{proof}

V54 will address the first excluded case without leaving positive
measures: a symmetric double-well scalar conditioned to one phase.  It
will quantify the local convexity radius, the exponentially small phase
leakage, and the block size needed for a positive effective Hessian, then
test whether those constants can remain uniform near the Higgs crossover.

\clearpage
\section{V54: phase-restricted convexity and Higgs radial leakage}
\label{sec:v5v54}

V53 propagated a physical mass gap through a positive conditionally convex
bosonic marginal.  A broken-symmetry polynomial is coercive but not
globally convex, so that theorem does not apply directly.  This note
separates a convex phase basin from its complement.  Within the basin,
the strong-convexity Wasserstein estimate remains valid.  The full
conditional law differs from the restricted law by exactly its leakage
probability, which is exponentially small when a block action has a
positive barrier.

For the gauge-invariant Higgs radius
\(\rho=|H|\), the barrier and convexity constants are explicit.  They
yield a controlled phase-restricted fibre away from the Higgs crossover.
The barrier vanishes as the symmetry-breaking radius tends to zero, so
the estimate is not uniform through the crossover.  The phase or radial
stratum must then be retained in the V39 stratified base chart.

\subsection{A general energy-gap leakage lemma}

Let \((X,\mu)\) be a probability space, let \(F:X\to\mathbb R\), and let
\begin{equation}
 d\nu_\beta
 =
 Z_\beta^{-1}e^{-\beta F}\,d\mu,
 \qquad \beta>0.
 \label{eq:v5v54-Gibbs-beta}
\end{equation}
Let \(A\subset X\) be the desired basin and \(B=X\setminus A\).

\begin{lemma}[Energy-gap leakage]
\label{lem:v5v54-leakage}
Suppose there are \(A_0\subset A\), \(m\in\mathbb R\),
\(\Delta>\varepsilon\ge0\), with \(\mu(A_0)>0\), such that
\begin{align}
 F&\le m+\varepsilon
 &&\text{on }A_0,\label{eq:v5v54-good-window}\\
 F&\ge m+\Delta
 &&\text{on }B.\label{eq:v5v54-bad-gap}
\end{align}
Then
\begin{equation}
 \nu_\beta(B)
 \le
 \frac{\mu(B)}{\mu(A_0)}
 e^{-\beta(\Delta-\varepsilon)}.
 \label{eq:v5v54-leakage}
\end{equation}
\end{lemma}

\begin{proof}
The numerator of \(\nu_\beta(B)\) is at most
\(\mu(B)e^{-\beta(m+\Delta)}\), while
\[
 Z_\beta
 \ge
 \int_{A_0}e^{-\beta F}\,d\mu
 \ge
 \mu(A_0)e^{-\beta(m+\varepsilon)}.
\]
Divide the two estimates.
\end{proof}

\begin{corollary}[Cofinal leakage criterion]
\label{cor:v5v54-cofinal-leakage}
For a sequence of blocks, if
\begin{equation}
 \beta_j(\Delta_j-\varepsilon_j)
 -
 \log\frac{\mu_j(B_j)}{\mu_j(A_{0,j})}
 \longrightarrow+\infty,
 \label{eq:v5v54-cofinal-condition}
\end{equation}
then the phase leakage tends to zero.
\end{corollary}

\begin{proof}
Take logarithms in
\eqref{eq:v5v54-leakage}.
\end{proof}

\begin{firewall}[A local minimum is not an energy gap]
\label{fw:v5v54-local-minimum}
The condition requires the entire complement \(B\) to lie above the
chosen phase window.  A second minimum of equal or lower energy makes
\(\Delta\le\varepsilon\) and invalidates the estimate, even if the desired
minimum has a positive local Hessian.
\end{firewall}

\subsection{Convex transport inside a phase basin}

Let \(A\subset\mathbb R^r\) be closed and convex.  For a boundary parameter
\(b\), define
\begin{equation}
 d\nu_b^A(x)
 =
 (Z_b^A)^{-1}
 \mathbf 1_A(x)e^{-W(x)+\langle b,x\rangle}\,dx.
 \label{eq:v5v54-restricted-law}
\end{equation}

\begin{theorem}[Restricted strong-convexity transport]
\label{thm:v5v54-restricted-transport}
If \(\nabla^2W\ge M>0\) on \(A\), then
\begin{equation}
 W_{2,M}(\nu_b^A,\nu_{b'}^A)
 \le
 |M^{-1/2}(b-b')|.
 \label{eq:v5v54-restricted-transport}
\end{equation}
\end{theorem}

\begin{proof}
Apply the synchronous Langevin coupling from
Theorem~\ref{thm:v5v52-matrix-transport} with normal reflection at the
boundary of \(A\).  Convexity of \(A\) makes the two reflection terms
nonexpansive:
\[
 \langle X-Y,n_A(X)\,dL_X-n_A(Y)\,dL_Y\rangle\ge0
\]
with the sign convention in which they are subtracted from the
difference equation.  Strong monotonicity of \(\nabla W\) then gives the
same differential inequality as in the unconstrained proof.  Approximate
a nonsmooth convex boundary by smooth convex domains and pass to the
limit.
\end{proof}

Let \(\nu_b\) be the unrestricted law and define its leakage
\begin{equation}
 \epsilon_b=\nu_b(A^c).
 \label{eq:v5v54-epsilon-b}
\end{equation}
Let \(d\le1\) be a bounded version of the phase metric.

\begin{lemma}[Restriction comparison]
\label{lem:v5v54-restriction-TV}
The normalized restriction \(\nu_b^A=\nu_b(\,\cdot\,\mid A)\) satisfies
\begin{equation}
 \|\nu_b-\nu_b^A\|_{\mathrm{TV}}=\epsilon_b.
 \label{eq:v5v54-restriction-TV}
\end{equation}
Consequently,
\begin{equation}
 W_d(\nu_b,\nu_{b'})
 \le
 W_d(\nu_b^A,\nu_{b'}^A)
 +\epsilon_b+\epsilon_{b'}.
 \label{eq:v5v54-full-vs-restricted}
\end{equation}
\end{lemma}

\begin{proof}
The density of \(\nu_b^A\) relative to \(\nu_b\) is
\(\mathbf1_A/(1-\epsilon_b)\).  Direct integration of its absolute
difference from one gives \(2\epsilon_b\), hence total variation
\(\epsilon_b\).  Since \(d\le1\), \(W_d\le\|\cdot\|_{\mathrm{TV}}\).
Apply the triangle inequality through the two restricted laws.
\end{proof}

\begin{theorem}[Affine phase-fibre contraction]
\label{thm:v5v54-affine-contraction}
Suppose a restricted block update contracts a discrepancy by
\(\theta<1\), and both endpoint conditional laws have leakage at most
\(\epsilon\).  Then the unrestricted update obeys
\begin{equation}
 d'\le\theta d+2\epsilon.
 \label{eq:v5v54-affine-contraction}
\end{equation}
After \(n\) iterations,
\begin{equation}
 d_n
 \le
 \theta^nd_0+
 \frac{2\epsilon(1-\theta^n)}{1-\theta}.
 \label{eq:v5v54-affine-iterate}
\end{equation}
\end{theorem}

\begin{proof}
Equation~\eqref{eq:v5v54-affine-contraction} is
Lemma~\ref{lem:v5v54-restriction-TV} followed by the restricted
contraction.  Iterate the affine recurrence and sum the geometric series.
\end{proof}

\begin{firewall}[Exponentially small is not zero]
\label{fw:v5v54-small-not-zero}
At fixed block size, leakage leaves the nonzero error floor
\(2\epsilon/(1-\theta)\).  Exact uniqueness in a selected phase requires
cofinal leakage tending to zero fast enough, or inclusion of the phase
label in the retained state.
\end{firewall}

\subsection{The gauge-invariant Higgs radius}

Consider the radial polynomial
\begin{equation}
 U(\rho)
 =
 \frac{\lambda}{4}(\rho^2-v^2)^2,
 \qquad
 \rho\ge0,
 \qquad
 \lambda>0,\quad v>0.
 \label{eq:v5v54-Higgs-radial}
\end{equation}
Its first two derivatives are
\begin{equation}
 U'(\rho)=\lambda\rho(\rho^2-v^2),
 \qquad
 U''(\rho)=\lambda(3\rho^2-v^2).
 \label{eq:v5v54-Higgs-derivatives}
\end{equation}
Choose
\begin{equation}
 \frac v{\sqrt3}<r_0<v,
 \qquad
 A=[r_0,\infty).
 \label{eq:v5v54-r0}
\end{equation}

\begin{theorem}[Explicit radial convexity and barrier]
\label{thm:v5v54-radial-barrier}
On \(A\),
\begin{equation}
 U''(\rho)\ge
 \mu_A:=\lambda(3r_0^2-v^2)>0.
 \label{eq:v5v54-radial-convexity}
\end{equation}
On the complement \(B=[0,r_0)\),
\begin{equation}
 U(\rho)\ge
 \Delta_A:=
 U(r_0)
 =
 \frac{\lambda}{4}(v^2-r_0^2)^2>0,
 \label{eq:v5v54-radial-gap}
\end{equation}
while \(\min_AU=U(v)=0\).
\end{theorem}

\begin{proof}
The second derivative in
\eqref{eq:v5v54-Higgs-derivatives} is increasing for \(\rho\ge0\), so
its minimum on \(A\) is its value at \(r_0\), proving
\eqref{eq:v5v54-radial-convexity}.  The function \(U\) decreases on
\([0,v]\); hence its minimum on \([0,r_0]\) is \(U(r_0)\).
The minimum \(v\) lies in \(A\) and has value zero.
\end{proof}

Let \(\mu_{\mathrm{rad}}\) be any finite-cutoff radial probability
reference assigning positive mass to every interval about \(v\).  For
\(\delta>0\) small, let
\begin{equation}
 A_0=[v-\delta,v+\delta]\subset A,
 \qquad
 \varepsilon_\delta=\sup_{\rho\in A_0}U(\rho).
 \label{eq:v5v54-radial-window}
\end{equation}

\begin{corollary}[Radial phase leakage]
\label{cor:v5v54-radial-leakage}
If \(\varepsilon_\delta<\Delta_A\), then the block radial law
\[
 d\nu_\beta(\rho)
 \propto e^{-\beta U(\rho)}\,d\mu_{\mathrm{rad}}(\rho)
\]
satisfies
\begin{equation}
 \nu_\beta([0,r_0))
 \le
 \frac{\mu_{\mathrm{rad}}([0,r_0))}
 {\mu_{\mathrm{rad}}(A_0)}
 e^{-\beta(\Delta_A-\varepsilon_\delta)}.
 \label{eq:v5v54-radial-leakage}
\end{equation}
\end{corollary}

\begin{proof}
Apply Lemma~\ref{lem:v5v54-leakage} with \(m=0\).
\end{proof}

\begin{assessment}[Radial Jacobian]
\label{ass:v5v54-radial-Jacobian}
For an \(N\)-component Higgs field, polar coordinates contribute
\(\rho^{N-1}\).  Keeping it in the reference
\(\mu_{\mathrm{rad}}\) is exact.  If it is moved into the action, the term
\(-(N-1)\log\rho/\beta\) has positive second derivative on
\(\rho>0\) and therefore improves the restricted radial convexity.
\end{assessment}

\subsection{Adding spatial coupling}

At one radial block, a quadratic boundary coupling adds
\[
 \frac{\kappa z}{2}\rho^2
 -\kappa\rho\sum_{j\sim i}\rho_j
\]
to the conditional potential.

\begin{proposition}[Restricted radial influence]
\label{prop:v5v54-radial-influence}
On the basin \(A=[r_0,\infty)\), the conditional curvature is at least
\begin{equation}
 \lambda_A=\mu_A+\kappa z.
 \label{eq:v5v54-radial-total-curvature}
\end{equation}
The restricted conditional influence of each neighbour is at most
\begin{equation}
 c_{ij}^A=\frac{\kappa}{\mu_A+\kappa z},
 \label{eq:v5v54-radial-cij}
\end{equation}
and the restricted row sum is
\begin{equation}
 \theta_A=\frac{\kappa z}{\mu_A+\kappa z}<1.
 \label{eq:v5v54-radial-row}
\end{equation}
\end{proposition}

\begin{proof}
Use \eqref{eq:v5v54-radial-convexity} and apply
Theorem~\ref{thm:v5v54-restricted-transport} exactly as in the scalar
site proof of V51.
\end{proof}

\begin{corollary}[Full radial affine row]
\label{cor:v5v54-full-radial-row}
If all conditional radial laws in the compared boundary class have
leakage at most \(\epsilon_\beta\), one restricted update satisfies
\begin{equation}
 d'\le\theta_A d+2\epsilon_\beta.
 \label{eq:v5v54-full-radial-row}
\end{equation}
\end{corollary}

\begin{proof}
Combine Proposition~\ref{prop:v5v54-radial-influence} and
Theorem~\ref{thm:v5v54-affine-contraction}.
\end{proof}

\begin{firewall}[Boundary coupling can move the barrier]
\label{fw:v5v54-boundary-barrier}
The linear boundary term changes the energy difference between \(A\) and
its complement.  Corollary~\ref{cor:v5v54-radial-leakage} applies to the
conditional law only after the worst permitted boundary field has been
included in \(\Delta_A\).  Unbounded adverse boundary values require a
Lyapunov-weighted estimate rather than a uniform leakage constant.
\end{firewall}

\subsection{Failure at the Higgs crossover}

Write \(r_0=cv\) with
\begin{equation}
 \frac1{\sqrt3}<c<1.
 \label{eq:v5v54-c}
\end{equation}
Then
\begin{align}
 \mu_A&=\lambda v^2(3c^2-1),
 \label{eq:v5v54-mu-scaling}\\
 \Delta_A&=\frac{\lambda v^4}{4}(1-c^2)^2.
 \label{eq:v5v54-Delta-scaling}
\end{align}

\begin{theorem}[No uniform radial-basin certificate at \(v=0\)]
\label{thm:v5v54-crossover-failure}
For fixed \(\lambda,c\), both the restricted curvature and the barrier
vanish as \(v\to0\):
\begin{equation}
 \mu_A=O(v^2),
 \qquad
 \Delta_A=O(v^4).
 \label{eq:v5v54-crossover-rates}
\end{equation}
Consequently, neither a row gap bounded away from zero nor a fixed-block
exponential leakage estimate follows uniformly through the crossover.
\end{theorem}

\begin{proof}
Equations~\eqref{eq:v5v54-mu-scaling} and
\eqref{eq:v5v54-Delta-scaling} give the rates.  The row gap is
\[
 1-\theta_A
 =
 \frac{\mu_A}{\mu_A+\kappa z},
\]
which tends to zero when \(\kappa z\) does not vanish faster than
\(\mu_A\).  The leakage exponent in
\eqref{eq:v5v54-radial-leakage} contains
\(\beta\Delta_A\), which is not uniformly large at fixed block size.
\end{proof}

\begin{assessment}[Necessary block growth near crossover]
\label{ass:v5v54-block-growth}
The leakage criterion requires at least
\[
 \beta_j\lambda_jv_j^4\longrightarrow\infty
\]
up to the reference-volume term in
\eqref{eq:v5v54-cofinal-condition}.  Such growing blocks can conflict
with locality and with the fixed shell atlas.  At or near \(v=0\), the
radial mode should instead be retained as a crossover coordinate.
\end{assessment}

\subsection{Relation to the physical Higgs and Einstein sectors}

\begin{firewall}[A scalar sign is not the Higgs gauge quotient]
\label{fw:v5v54-Higgs-gauge}
The physical Higgs phase is not selected by choosing one Cartesian sign.
Gauge rotations move points on the vacuum orbit.  The radial argument is
gauge invariant, but angular gauge, Goldstone, ghost, and gauge-boson
blocks must still be assembled through the BRST/constraint quotient before
the radial influence coefficient is physical.
\end{firewall}

\begin{firewall}[The Einstein conformal mode has no analogous positive
basin yet]
\label{fw:v5v54-Einstein}
V8 found no cutoff-uniform coercive straight contour for the conformal
metric direction.  V54 assumes a positive real radial measure and cannot
be transferred to that direction without a proved physical contour or
constraint elimination.
\end{firewall}

\begin{definition}[Phase-restricted radial certificate]
\label{def:v5v54-PRRC}
A PRRC consists of:
\begin{enumerate}
 \item a gauge-invariant convex radial basin with curvature
 \(\mu_A>0\);
 \item a conditional energy gap including all allowed boundary fields;
 \item a cofinal leakage estimate satisfying
 \eqref{eq:v5v54-cofinal-condition};
 \item a restricted diffusive influence row and its coupling to angular
 and gauge blocks;
 \item retention of the phase/radial label whenever leakage does not
 vanish;
 \item a separate crossover chart where \(\mu_A\) or \(\Delta_A\)
 degenerates.
\end{enumerate}
\end{definition}

\subsection{V54 status}

\begin{theorem}[V54 conclusion]
\label{thm:v5v54-status}
A positive-energy phase basin gives a rigorously controlled approximation
to a nonconvex conditional fibre: strong convexity contracts the
restricted law, and the unrestricted error is exactly bounded by phase
leakage.  For the gauge-invariant Higgs radius, the convexity and barrier
are \eqref{eq:v5v54-radial-convexity} and
\eqref{eq:v5v54-radial-gap}.  Both degenerate at the Higgs crossover, so
the PRRC is valid only away from that stratum unless the block size grows
fast enough.  No full Higgs--gauge or Einstein PRRC is yet proved.
\end{theorem}

\begin{proof}
Use Lemma~\ref{lem:v5v54-leakage},
Theorem~\ref{thm:v5v54-restricted-transport},
Proposition~\ref{prop:v5v54-radial-influence}, and
Theorem~\ref{thm:v5v54-crossover-failure}.
\end{proof}

V55 will perform the compulsory YM/NS audit and then build the explicit
stratified handoff between the convex massive Higgs radial fibre and the
retained crossover coordinate.  It will determine the overlap estimates
needed to prevent double counting of radial fluctuations and to preserve
the V49 relative-density metric across the handoff.

\clearpage
\section{V55: companion audit and exact stratified radial handoff}
\label{sec:v5v55}

V54 constructed a convex massive Higgs radial basin and proved that both
its curvature and its barrier degenerate at the crossover.  This note
performs the compulsory companion audit and gives an exact handoff between
two density charts: one retains the radial variable, while the other
treats it as part of a massive conditional fibre.  Both charts are
factorizations of the same joint density.  Their overlap transition is
therefore an isometry in the V49 weighted \(L^1\) metric and cannot double
count a fluctuation.

Collapsing the massive radial conditional is a separate approximation.
It is justified only when the conditional fibre diameter tends to zero.
The V54 leakage and V50/V51 contraction estimates give the explicit error
\(2\epsilon/(1-\theta)\).  Near the crossover this error need not vanish,
so the radial variable remains in the retained stratified state.

\subsection{Mandatory V55 companion audit}

The newest YM file inspected was
\begin{equation}
 \texttt{Renewal\_Yang\_Mills\_Mass\_Gap\_Research\_Notes\_V62.tex},
 \qquad
 \mathrm{SHA256}=
 \texttt{C9642B20E450F9FBB7383558803A97962793475BD425AF5C0FC56B9FC5479516}.
 \label{eq:v5v55-YM-hash}
\end{equation}
The NS file remains
\begin{equation}
 \texttt{Renewal\_NS\_Stability\_Research\_Notes\_V31.tex},
 \qquad
 \mathrm{SHA256}=
 \texttt{F1121B62238AD5491BB7CF03635EA9934942EDF573ED8FAAA9EE79E1D38A6696}.
 \label{eq:v5v55-NS-hash}
\end{equation}
No companion file was modified.

The relevant YM advances after V57 are:
\begin{enumerate}
 \item V58 closes the output-dominant discrete quintic branch;
 \item V59 closes the branching-dominant branch, proves
 \(\operatorname{PGSEL}_5\) on every selected fibre, and completes the
 global quintic marked-tree atlas;
 \item V60 tracks the all-order factorial-exponential constant ledger;
 \item V61 proves an all-order centred displacement cumulant estimate and
 an all-order hot-star local tree card at the required arity scale;
 \item V62 proves, from an exact \(SU(2)\) Bessel lower bound, that the
 proposed uniform \(m!K^m\) private Casimir-moment estimate is false.
\end{enumerate}
The remaining YM alternatives are to re-fibre high-degree private
terminals into hot local stars, or to change the arity/representation
Gevrey scale.  The all-order MTA, WUFC, continuum measure, OS
reconstruction, and mass gap remain unproved.  The NS endpoint is
unchanged and supplies no new radial crossover estimate.

\begin{assessment}[Audit transfer]
\label{ass:v5v55-audit-transfer}
The exact YM constants do not transfer to the Higgs sector.  Its logical
lesson does: when an all-order moment bound for a collapsed terminal is
false, one must not keep using that terminal representation.  The
high-degree variable should be retained as an active local coordinate or
placed in a sharper grade.  The V55 radial crossover chart implements the
first option exactly at the density level.
\end{assessment}

\begin{firewall}[Quintic closure is not an SM all-order theorem]
\label{fw:v5v55-quintic-scope}
The complete YM quintic atlas and hot-star card concern the declared
Wilson representation norm.  They justify an assembly and re-fibration
strategy, not a numerical SM influence coefficient or a Higgs moment
bound.
\end{firewall}

\subsection{One joint density, two augmented factorizations}

Let \(P\) be the retained nonradial base, \(R\) the Higgs radial variable,
and \(Q\) the remaining conditional fibre.  Let
\begin{equation}
 \mu=\mu_P\otimes\mu_R\otimes\mu_Q
 \label{eq:v5v55-reference}
\end{equation}
and let \(r(p,\rho,q)\ge0\) be a normalized density.

The \emph{crossover chart} retains \((p,\rho)\):
\begin{align}
 m_C(p,\rho)
 &=
 \int_Qr(p,\rho,q)\,d\mu_Q(q),
 \label{eq:v5v55-mC}\\
 k_C(q\mid p,\rho)
 &=
 \frac{r(p,\rho,q)}{m_C(p,\rho)}.
 \label{eq:v5v55-kC}
\end{align}
The \emph{massive chart} retains only \(p\):
\begin{align}
 m_M(p)
 &=
 \int_R\int_Qr(p,\rho,q)\,d\mu_Q(q)d\mu_R(\rho),
 \label{eq:v5v55-mM}\\
 k_M(\rho,q\mid p)
 &=
 \frac{r(p,\rho,q)}{m_M(p)}.
 \label{eq:v5v55-kM}
\end{align}
As in V49, conditional densities on zero-marginal sets may be chosen
arbitrarily.

\begin{theorem}[Exact radial re-fibration]
\label{thm:v5v55-refibration}
On positive-marginal sets, the two charts are related by
\begin{align}
 m_M(p)
 &=
 \int_Rm_C(p,\rho)\,d\mu_R(\rho),
 \label{eq:v5v55-C-to-M-marginal}\\
 k_M(\rho,q\mid p)
 &=
 \frac{m_C(p,\rho)}{m_M(p)}
 k_C(q\mid p,\rho),
 \label{eq:v5v55-C-to-M-conditional}\\
 m_C(p,\rho)
 &=
 m_M(p)\int_Qk_M(\rho,q\mid p)\,d\mu_Q(q),
 \label{eq:v5v55-M-to-C-marginal}\\
 k_C(q\mid p,\rho)
 &=
 \frac{k_M(\rho,q\mid p)}
 {\int_Qk_M(\rho,q'\mid p)\,d\mu_Q(q')}.
 \label{eq:v5v55-M-to-C-conditional}
\end{align}
Each pair reconstructs exactly the same joint density:
\begin{equation}
 r=m_Ck_C=m_Mk_M.
 \label{eq:v5v55-same-density}
\end{equation}
\end{theorem}

\begin{proof}
Integrate \eqref{eq:v5v55-mC} over \(\rho\) to obtain
\eqref{eq:v5v55-C-to-M-marginal}.  Divide
\(r=m_Ck_C\) by \(m_M\) to obtain
\eqref{eq:v5v55-C-to-M-conditional}.  Conversely, integrate
\(r=m_Mk_M\) over \(q\) to obtain
\eqref{eq:v5v55-M-to-C-marginal}, and divide by the result to obtain
\eqref{eq:v5v55-M-to-C-conditional}.  Equation
\eqref{eq:v5v55-same-density} is the defining factorization in either
chart.
\end{proof}

\begin{corollary}[No radial double counting]
\label{cor:v5v55-no-double-count}
The radial variable occurs once in either chart: as a retained coordinate
in \(m_C\), or as a conditional coordinate in \(k_M\).  Applying
Theorem~\ref{thm:v5v55-refibration} changes its role without multiplying or
adding a second radial measure.
\end{corollary}

\begin{proof}
Both products in \eqref{eq:v5v55-same-density} equal the same
Radon--Nikodym derivative \(r\).
\end{proof}

\subsection{The common weighted metric}

Let \(\Omega(p,\rho,q)\ge1\) be a physical Lyapunov weight.  For either
augmented chart \(\mathcal X=(m,k)\), define its reconstruction
\(\mathscr R\mathcal X=mk\) and the distance
\begin{equation}
 d_\Omega(\mathcal X,\mathcal Y)
 =
 \int\Omega
 |\mathscr R\mathcal X-\mathscr R\mathcal Y|\,d\mu.
 \label{eq:v5v55-chart-metric}
\end{equation}

\begin{theorem}[Isometric overlap transition]
\label{thm:v5v55-overlap-isometry}
The transition
Theorem~\ref{thm:v5v55-refibration} is an isometry between the massive and
crossover augmented charts:
\begin{equation}
 d_\Omega(\mathcal X_M,\mathcal Y_M)
 =
 d_\Omega(\mathcal X_C,\mathcal Y_C)
 \label{eq:v5v55-overlap-isometry}
\end{equation}
whenever the paired charts represent the same two physical densities.
\end{theorem}

\begin{proof}
Both sides equal
\[
 \int\Omega|r-s|\,d\mu
\]
by \eqref{eq:v5v55-same-density}.
\end{proof}

\begin{corollary}[Closed overlap graph]
\label{cor:v5v55-closed-graph}
In the metric completions defined by
\eqref{eq:v5v55-chart-metric}, the graph of the overlap transition is
closed.
\end{corollary}

\begin{proof}
An isometry is continuous.  If
\((\mathcal X_{M,n},\mathcal X_{C,n})\) lies in the graph and converges to
\((\mathcal X_M,\mathcal X_C)\), their reconstructed densities have the
same limit in weighted \(L^1\).  Hence
\(\mathcal X_M,\mathcal X_C\) represent the same density and remain in
the graph.
\end{proof}

\begin{assessment}[Compatibility with the V39 quotient]
\label{ass:v5v55-V39}
The closed graph supplies the topological input needed to identify the
two descriptions in the V39 compact stratified quotient, provided the
cutoff chart sets themselves are compact or tight in the chosen topology.
The identification introduces no metric discontinuity.
\end{assessment}

\subsection{When the radial conditional may be collapsed}

Fix a canonical massive conditional density
\(\bar k_M(\rho,q\mid p)\), for example the unique limit of the
phase-restricted conditional update.  Define the collapse and section
\begin{align}
 \mathscr C(m_M,k_M)&=m_M,
 \label{eq:v5v55-collapse}\\
 \mathscr I(m_M)&=(m_M,\bar k_M).
 \label{eq:v5v55-section}
\end{align}

\begin{lemma}[Exact collapse error]
\label{lem:v5v55-collapse-error}
For the unweighted density metric,
\begin{equation}
 d_1\bigl(
 (m_M,k_M),\mathscr I\mathscr C(m_M,k_M)
 \bigr)
 =
 \int_Pm_M(p)
 \|k_M(\cdot\mid p)-\bar k_M(\cdot\mid p)\|_{L^1(\mu_R\otimes\mu_Q)}
 d\mu_P(p).
 \label{eq:v5v55-collapse-error}
\end{equation}
The same identity holds with a conditional physical weight inside the
fibre norm.
\end{lemma}

\begin{proof}
The two reconstructed densities are \(m_Mk_M\) and
\(m_M\bar k_M\).  Factor out the nonnegative marginal \(m_M\) and use
Fubini.
\end{proof}

Suppose one complete phase-fibre shell update has contraction factor
\(\theta_j<1\) and leakage error at most \(2\epsilon_j\), as in V54.
Then the diameter of its invariant conditional set is controlled by
\begin{equation}
 \delta_j=\frac{2\epsilon_j}{1-\theta_j}.
 \label{eq:v5v55-fibre-diameter}
\end{equation}

\begin{theorem}[Cofinal radial collapse]
\label{thm:v5v55-cofinal-collapse}
Assume every admissible massive conditional at scale \(j\) lies within
\(\delta_j\) in fibre \(L^1\) of \(\bar k_{M,j}\), uniformly in the
retained base, and
\begin{equation}
 \delta_j\longrightarrow0.
 \label{eq:v5v55-delta-zero}
\end{equation}
Then collapse followed by the canonical section converges to the identity
in the joint density metric:
\begin{equation}
 \sup_{\mathcal X_{M,j}}
 d_1(\mathcal X_{M,j},
 \mathscr I_j\mathscr C_j\mathcal X_{M,j})
 \longrightarrow0.
 \label{eq:v5v55-collapse-identity}
\end{equation}
Every bounded physical observable has the same cofinal expectation in the
augmented and collapsed massive descriptions.
\end{theorem}

\begin{proof}
Lemma~\ref{lem:v5v55-collapse-error} bounds the joint error by
\[
 \int_Pm_M(p)\delta_j\,d\mu_P(p)=\delta_j.
\]
This proves \eqref{eq:v5v55-collapse-identity}.  The expectation
difference of an observable \(A\) is at most
\(\|A\|_\infty\delta_j\) by the \(L^1\)-\(L^\infty\) duality.
\end{proof}

\begin{corollary}[Leakage--gap collapse criterion]
\label{cor:v5v55-leakage-gap}
The radial conditional may be collapsed if
\begin{equation}
 \frac{\epsilon_j}{1-\theta_j}\longrightarrow0.
 \label{eq:v5v55-leakage-gap}
\end{equation}
\end{corollary}

\begin{proof}
Use \eqref{eq:v5v55-fibre-diameter} in
Theorem~\ref{thm:v5v55-cofinal-collapse}.
\end{proof}

\begin{firewall}[Leakage alone does not prove collapse]
\label{fw:v5v55-leakage-alone}
Even when \(\epsilon_j\to0\), a simultaneously vanishing contraction gap
can make \(\epsilon_j/(1-\theta_j)\) nonzero or divergent.  The barrier and
the conditional mixing rate must be compared in the same cofinal scaling.
\end{firewall}

\subsection{The massive and crossover strata}

For the V54 radial basin, let
\begin{align}
 g_j&=1-\theta_{A,j}
 =
 \frac{\mu_{A,j}}{\mu_{A,j}+\kappa_jz},
 \label{eq:v5v55-gap-gj}\\
 \epsilon_j&\le
 C_j e^{-\beta_j(\Delta_{A,j}-\varepsilon_j)}.
 \label{eq:v5v55-epsilon-j}
\end{align}
Define the massive admissible set by
\begin{equation}
 \mathfrak M
 =
 \left\{
 \frac{\epsilon_j}{g_j}\to0
 \right\},
 \label{eq:v5v55-massive-stratum}
\end{equation}
and the crossover set \(\mathfrak C\) as the complementary scales on which
the radial coordinate is retained.

\begin{theorem}[Explicit radial handoff condition]
\label{thm:v5v55-handoff-condition}
For \(r_{0,j}=cv_j\), \(1/\sqrt3<c<1\), a sufficient massive-collapse
condition is
\begin{equation}
 C_j
 \frac{\mu_{A,j}+\kappa_jz}{\mu_{A,j}}
 \exp\!\left[
 -\beta_j\left(
 \frac{\lambda_jv_j^4}{4}(1-c^2)^2-\varepsilon_j
 \right)
 \right]
 \longrightarrow0,
 \label{eq:v5v55-explicit-handoff}
\end{equation}
where
\begin{equation}
 \mu_{A,j}=\lambda_jv_j^2(3c^2-1).
 \label{eq:v5v55-muA}
\end{equation}
When \eqref{eq:v5v55-explicit-handoff} fails, V54 does not justify
collapsing the radial conditional.
\end{theorem}

\begin{proof}
Insert the V54 barrier and curvature formulas into
\(\epsilon_j/g_j\).  The displayed expression is exactly the resulting
upper bound.  Its convergence to zero implies
\eqref{eq:v5v55-leakage-gap}; failure only says the present estimate does
not prove collapse.
\end{proof}

\begin{assessment}[Crossover retention]
\label{ass:v5v55-crossover-retention}
As \(v_j\to0\), the barrier exponent is of order
\(\beta_j\lambda_jv_j^4\), while the prefactor contains a factor
comparable to \(1+\kappa_jz/(\lambda_jv_j^2)\).  Unless block growth beats
both degeneracies, the radial law belongs in \(m_C(p,\rho)\), not in a
collapsed conditional.
\end{assessment}

\subsection{Shell maps descend through the handoff}

Let \(\mathscr S\) be a shell Markov map on joint densities.  Define its
representatives in either chart by reconstruction, application of
\(\mathscr S\), and refactorization.

\begin{theorem}[Exact intertwining of chart shell maps]
\label{thm:v5v55-shell-intertwining}
Let \(\mathscr T_{MC}\) denote the overlap transition from the massive to
the crossover augmented chart.  Then
\begin{equation}
 \mathscr T_{MC}\mathscr S_M
 =
 \mathscr S_C\mathscr T_{MC}
 \label{eq:v5v55-shell-intertwining}
\end{equation}
wherever both sides are defined.
\end{theorem}

\begin{proof}
Each route reconstructs the same input density \(r\), applies the same
physical Markov map \(\mathscr S\), and factorizes the same output density
\(\mathscr Sr\).  Theorem~\ref{thm:v5v55-refibration} makes the two output
factorizations identical on the overlap.
\end{proof}

\begin{corollary}[Descent to the stratified quotient]
\label{cor:v5v55-quotient-descent}
If the massive and crossover chart sets form a closed equivalence relation
as in Corollary~\ref{cor:v5v55-closed-graph}, then the shell map descends
to their quotient.
\end{corollary}

\begin{proof}
Equation~\eqref{eq:v5v55-shell-intertwining} says equivalent
representatives have equivalent images.  Therefore the value of the
quotient map is independent of the representative.
\end{proof}

\begin{firewall}[Collapsed dynamics is only asymptotically intertwined]
\label{fw:v5v55-collapsed-intertwining}
The exact intertwining theorem uses the full conditional \(k_M\).
Replacing it by \(\bar k_M\) introduces the error
\eqref{eq:v5v55-collapse-error}.  The collapsed shell map is a valid
cofinal representative only under
\eqref{eq:v5v55-delta-zero}.
\end{firewall}

\subsection{Stratified radial handoff certificate}

\begin{definition}[Stratified radial handoff certificate]
\label{def:v5v55-SRHC}
An SRHC consists of:
\begin{enumerate}
 \item a common positive joint relative density on \(P\times R\times Q\);
 \item the exact augmented factorizations
 \eqref{eq:v5v55-mC}--\eqref{eq:v5v55-kM};
 \item the common weighted metric
 \eqref{eq:v5v55-chart-metric} and closed overlap graph;
 \item a massive region satisfying
 \(\epsilon_j/(1-\theta_j)\to0\);
 \item a crossover region retaining the radial marginal;
 \item exact shell intertwining before collapse and quantified asymptotic
 intertwining after collapse;
 \item gauge/BRST completion of the radial--angular split and compatibility
 with the V39 stratified ancestry.
\end{enumerate}
\end{definition}

\subsection{V55 status}

\begin{theorem}[V55 conclusion]
\label{thm:v5v55-status}
The massive and crossover Higgs radial descriptions are exact augmented
factorizations of one joint density, and their overlap transition is an
isometry which intertwines the physical shell map.  A massive radial
conditional may be collapsed only when its leakage divided by its
contraction gap tends to zero; formula
\eqref{eq:v5v55-explicit-handoff} gives a sufficient condition.  If it
fails near the crossover, the radial coordinate remains in the retained
stratum.  This proves the abstract density-level SRHC and its scalar
radial criterion, but not the gauge/BRST-completed Higgs handoff or its
coupling to the Einstein sector.
\end{theorem}

\begin{proof}
Use Theorems~\ref{thm:v5v55-refibration},
\ref{thm:v5v55-overlap-isometry},
\ref{thm:v5v55-cofinal-collapse},
\ref{thm:v5v55-handoff-condition}, and
\ref{thm:v5v55-shell-intertwining}.
\end{proof}

V56 will gauge-complete the radial decomposition at finite cutoff.  It
will separate the Higgs modulus, gauge orbit, stabilizer, and Faddeev--
Popov factors, identify the singular orbit at zero modulus, and determine
which parts define a positive conditional density suitable for the SRHC.

\clearpage
\section{V56: finite-cutoff electroweak radial quotient and orbit strata}
\label{sec:v5v56}

V55 constructed an exact density-level handoff for an abstract Higgs radial
variable.  This note derives that variable from one electroweak Higgs
doublet at finite cutoff.  Away from zero modulus, the compact gauge orbit
can be separated exactly: Lebesgue measure contributes the positive radial
Jacobian \(\rho^3\), and unitary-gauge Faddeev--Popov volume has the same
cubic scaling.  At \(\rho=0\), the stabilizer jumps and the orbit quotient
is singular.  This is precisely the crossover stratum which V55 must
retain.

The Higgs kinetic term also gives the exact mass split on the nonzero
stratum.  The charged \(W\) modes and neutral \(Z\) mode acquire masses
proportional to \(\rho\), while the photon remains massless.  Thus only the
three broken gauge directions can enter the massive conditional fibre.
The photon, Higgs zeros, chiral determinant phase, and Einstein variables
remain outside the convex radial certificate.

\subsection{One-doublet orbit decomposition}

Identify a Higgs doublet \(H\in\mathbb C^2\) with \(\mathbb R^4\), and set
\begin{equation}
 \rho=|H|\ge0.
 \label{eq:v5v56-rho}
\end{equation}
For \(\rho>0\), write \(H=\rho u\) with \(u\in S^3\).

\begin{theorem}[Exact polar disintegration]
\label{thm:v5v56-polar}
Lebesgue measure on \(\mathbb C^2\) disintegrates as
\begin{equation}
 d^4H=\rho^3\,d\rho\,d\sigma_{S^3}(u).
 \label{eq:v5v56-polar}
\end{equation}
Consequently, for every nonnegative gauge-invariant function
\(F(H)=f(|H|)\),
\begin{equation}
 \int_{\mathbb C^2}F(H)\,d^4H
 =
 |S^3|\int_0^\infty f(\rho)\rho^3\,d\rho.
 \label{eq:v5v56-radial-integral}
\end{equation}
\end{theorem}

\begin{proof}
This is the Euclidean polar-coordinate formula in four real dimensions.
The volume element is \(\rho^{4-1}d\rho\,d\sigma_{S^3}\).  A
gauge-invariant function is constant on the angular sphere, giving
\eqref{eq:v5v56-radial-integral}.
\end{proof}

Let
\[
 G_{\mathrm{EW}}=SU(2)_L\times U(1)_Y
\]
act on the doublet with hypercharge \(Y=1/2\).

\begin{theorem}[Orbit and stabilizer strata]
\label{thm:v5v56-orbit-strata}
For every \(\rho>0\), the electroweak orbit through
\begin{equation}
 H_\rho=\frac1{\sqrt2}\binom{0}{\sqrt2\rho}
 =\binom0\rho
 \label{eq:v5v56-Hrho}
\end{equation}
is the sphere of doublets of modulus \(\rho\).  Its stabilizer is a copy of
\(U(1)_{\mathrm{em}}\), and
\begin{equation}
 G_{\mathrm{EW}}/U(1)_{\mathrm{em}}\simeq S^3.
 \label{eq:v5v56-orbit}
\end{equation}
At \(H=0\), the stabilizer is all of \(G_{\mathrm{EW}}\).  Hence the orbit
dimension drops from three to zero at \(\rho=0\).
\end{theorem}

\begin{proof}
\(SU(2)\) acts transitively on unit vectors of \(\mathbb C^2\), so every
nonzero doublet can be sent to \(H_\rho\).  The combined generator which
annihilates \(H_\rho\) is
\[
 Q=T^3+Y,
\]
and it generates \(U(1)_{\mathrm{em}}\).  The quotient has dimension
\(4-1=3\) and is the orbit \(S^3\).  Every group element fixes the zero
vector, proving the last assertion.
\end{proof}

\begin{corollary}[The quotient is stratified at zero modulus]
\label{cor:v5v56-stratified-zero}
The gauge-orbit quotient has a principal radial stratum \(\rho>0\) with
stabilizer \(U(1)_{\mathrm{em}}\) and a singular stratum \(\rho=0\) with
stabilizer \(G_{\mathrm{EW}}\).  No single smooth unitary-gauge chart
extends through \(\rho=0\).
\end{corollary}

\begin{proof}
Orbit type is locally constant on a smooth principal stratum.  The
stabilizer dimension jumps by three at zero, so a principal-orbit chart
cannot remain a local diffeomorphism there.
\end{proof}

\subsection{The unitary-gauge Jacobian}

Choose three gauge conditions which set the three angular components of
the doublet to zero and require the remaining component to be positive.
Let \(\chi(H)\in\mathbb R^3\) denote these conditions, and let
\(\xi\in\mathbb R^3\) parametrize the broken generators transverse to
\(U(1)_{\mathrm{em}}\).

\begin{theorem}[Cubic Faddeev--Popov scaling]
\label{thm:v5v56-FP}
At \(H_\rho\) with \(\rho>0\),
\begin{equation}
 \left|
 \det\frac{\partial\chi(e^\xi H_\rho)}{\partial\xi}
 \right|_{\xi=0}
 =
 C_{\mathrm{FP}}\rho^3
 \label{eq:v5v56-FP}
\end{equation}
for a convention-dependent constant \(C_{\mathrm{FP}}>0\).  The
determinant is positive after fixing the orientation of the broken
generators and vanishes cubically at \(\rho=0\).
\end{theorem}

\begin{proof}
Each infinitesimal broken generator acts linearly on \(H_\rho\), so each
of the three transverse tangent vectors is \(\rho\) times its value at
\(\rho=1\).  The derivative matrix of the three independent gauge
conditions therefore equals \(\rho\) times a fixed invertible
\(3\times3\) matrix.  Its determinant is
\(C_{\mathrm{FP}}\rho^3\).  Invertibility follows because the three
broken generators span the tangent space to the principal orbit.
\end{proof}

\begin{assessment}[Polar and Faddeev--Popov factors agree]
\label{ass:v5v56-Jacobian-agreement}
The cubic factor in \eqref{eq:v5v56-FP} is the orbit-volume Jacobian in
\eqref{eq:v5v56-polar}.  Introducing unitary-gauge ghosts or integrating
the compact orbit are two coordinate realizations of the same finite-
cutoff volume factor on the principal stratum.
\end{assessment}

\begin{firewall}[The zero is a chart singularity]
\label{fw:v5v56-FP-zero}
The vanishing determinant at \(\rho=0\) does not remove that physical
configuration.  It says the unitary-gauge angular coordinates cease to be
independent there.  Deleting the zero-modulus stratum would invalidate the
V55 crossover handoff.
\end{firewall}

\subsection{Exact positive bosonic quotient measure}

Let a finite lattice have Higgs variables \(H_i\), compact electroweak
link variables \(U_e\), and a real gauge-invariant bosonic action
\(S_B(H,U)\) bounded below strongly enough for integrability.  Its
unnormalized measure is
\begin{equation}
 e^{-S_B(H,U)}
 \prod_i d^4H_i
 \prod_e d\mathrm{Haar}(U_e).
 \label{eq:v5v56-bosonic-measure}
\end{equation}

\begin{theorem}[Positive gauge-orbit pushforward]
\label{thm:v5v56-positive-pushforward}
The pushforward of \eqref{eq:v5v56-bosonic-measure} to the algebra of
gauge-invariant variables is a positive measure.  On the open set
\(\rho_i>0\) for all sites, it admits local coordinates consisting of:
\begin{enumerate}
 \item the moduli \(\rho_i\), with factor
 \(\prod_i\rho_i^3d\rho_i\);
 \item gauge-transformed link variables;
 \item the residual \(U(1)_{\mathrm{em}}\) gauge quotient.
\end{enumerate}
No signed Faddeev--Popov factor is needed for gauge-invariant bosonic
expectations on this finite compact lattice.
\end{theorem}

\begin{proof}
A measurable pushforward of a positive measure is positive.  On the
principal stratum, choose at each site a measurable \(SU(2)\) rotation
sending \(H_i/|H_i|\) to the fixed unit vector in
\eqref{eq:v5v56-Hrho}, and transform all incident links by the same gauge
rotation.  Haar invariance leaves the link volume unchanged, while
Theorem~\ref{thm:v5v56-polar} supplies
\(\rho_i^3d\rho_i\).  Rotations preserving the fixed doublet form the
residual electromagnetic group.  Integrating the original compact orbit,
rather than inserting an oriented determinant with changing sign,
establishes positivity.
\end{proof}

\begin{firewall}[Principal-stratum coordinates are not global topology]
\label{fw:v5v56-global-section}
The local or measurable gauge rotations used above do not prove a smooth
global continuum unitary gauge in topologically nontrivial bundles.
Higgs zeros, defects, boundary conditions, and Gribov-type global
identifications must remain in the stratified quotient.
\end{firewall}

\subsection{Gauge-boson mass split}

Let the covariant derivative be
\begin{equation}
 D_\mu
 =
 \partial_\mu
 -ig\frac{\tau^a}{2}W_\mu^a
 -ig'\frac12B_\mu.
 \label{eq:v5v56-covariant-derivative}
\end{equation}
Evaluate the quadratic gauge-field term in
\(|D_\mu H_\rho|^2\) for a constant \(\rho\).

\begin{theorem}[Exact electroweak mass quadratic form]
\label{thm:v5v56-mass-form}
At \(H_\rho=(0,\rho)^T\),
\begin{equation}
 |D_\mu H_\rho|^2
 =
 \frac{\rho^2}{4}
 \left[
 g^2\bigl((W_\mu^1)^2+(W_\mu^2)^2\bigr)
 +(gW_\mu^3-g'B_\mu)^2
 \right].
 \label{eq:v5v56-mass-form}
\end{equation}
With
\begin{align}
 W_\mu^\pm&=\frac{W_\mu^1\mp iW_\mu^2}{\sqrt2},
 \label{eq:v5v56-Wpm}\\
 Z_\mu&=\frac{gW_\mu^3-g'B_\mu}{\sqrt{g^2+g'^2}},
 \label{eq:v5v56-Z}\\
 A_\mu&=\frac{g'W_\mu^3+gB_\mu}{\sqrt{g^2+g'^2}},
 \label{eq:v5v56-photon}
\end{align}
the masses are
\begin{equation}
 m_W^2=\frac{g^2\rho^2}{2},
 \qquad
 m_Z^2=\frac{(g^2+g'^2)\rho^2}{2},
 \qquad
 m_A^2=0
 \label{eq:v5v56-masses-convention}
\end{equation}
under the present choice \(|H_\rho|=\rho\) and the convention that a real
vector mass term is \(m^2V_\mu V_\mu/2\).  If the standard vacuum
parameter is defined by \(H=(0,v/\sqrt2)\), these become
\(m_W^2=g^2v^2/4\) and
\(m_Z^2=(g^2+g'^2)v^2/4\).
\end{theorem}

\begin{proof}
Use
\[
 \frac{\tau^1}{2}H_\rho=\frac{\rho}{2}\binom10,\quad
 \frac{\tau^2}{2}H_\rho=-\frac{i\rho}{2}\binom10,\quad
 \frac{\tau^3}{2}H_\rho=-\frac{\rho}{2}\binom01.
\]
The upper and lower components are orthogonal, giving
\eqref{eq:v5v56-mass-form}.  The charged and neutral changes of basis
diagonalize it.  The neutral \(2\times2\) matrix has eigenvalues zero and
\(g^2+g'^2\), with eigenvectors \(A_\mu\) and \(Z_\mu\).
Comparison with the \(m^2V^2/2\) convention gives
\eqref{eq:v5v56-masses-convention}.
\end{proof}

\begin{corollary}[Massive and retained gauge directions]
\label{cor:v5v56-gauge-strata}
For \(\rho>0\), the \(W^\pm\) and \(Z\) directions may enter a massive
conditional fibre if their interacting Schur mass remains positive.  The
photon must remain in the massless retained sector.  As \(\rho\to0\), all
three massive eigenvalues vanish and the full electroweak gauge block
returns to the crossover stratum.
\end{corollary}

\begin{proof}
This is the spectrum in
\eqref{eq:v5v56-masses-convention}.
\end{proof}

\begin{assessment}[Match to the V55 radial criterion]
\label{ass:v5v56-match-V55}
On a basin \(\rho\ge r_0>0\), the three broken gauge masses have positive
lower bounds proportional to \(r_0^2\).  Their conditional influence can
therefore be combined with the radial curvature by the V52 block Schur
test.  When \(r_0\to0\), the same lower bounds degenerate together with
the V54 radial barrier, requiring the exact V55 re-fibration.
\end{assessment}

\subsection{Fermion masses and the positivity boundary}

A Yukawa term has the schematic form
\begin{equation}
 y_f\overline\psi_LH\psi_R+\mathrm{h.c.}
 \label{eq:v5v56-Yukawa}
\end{equation}
and therefore produces, on the principal unitary-gauge stratum, a mass
linear in \(\rho\).

\begin{proposition}[Radial scaling of Yukawa masses]
\label{prop:v5v56-Yukawa-mass}
For a Yukawa eigenvalue \(y_f\), the corresponding Dirac mass satisfies
\begin{equation}
 m_f(\rho)=y_f\rho
 \label{eq:v5v56-Yukawa-mass}
\end{equation}
under the convention \(H_\rho=(0,\rho)\), or
\(m_f=y_fv/\sqrt2\) under \(H=(0,v/\sqrt2)\).
\end{proposition}

\begin{proof}
Substitution of the constant radial representative into
\eqref{eq:v5v56-Yukawa} leaves the bilinear coefficient \(y_f\rho\).
\end{proof}

\begin{firewall}[The chiral determinant is not part of the positive
bosonic pushforward]
\label{fw:v5v56-chiral-determinant}
Theorem~\ref{thm:v5v56-positive-pushforward} concerns the bosonic
gauge--Higgs measure.  Integrating chiral fermions can introduce a complex
phase or zero determinant.  Anomaly cancellation is necessary for gauge
consistency but does not by itself prove a nonnegative Euclidean
determinant or the V53 conditional Brascamp--Lieb hypotheses.
\end{firewall}

\subsection{The radial chart with its exact Jacobian}

Let \(r_{\mathrm{bos}}(\rho,\mathcal U)\) denote the normalized density of
the gauge-invariant bosonic variables relative to a product reference
which includes
\begin{equation}
 d\mu_R(\rho)=Z_R^{-1}\rho^3e^{-c\rho^2}\,d\rho
 \label{eq:v5v56-radial-reference}
\end{equation}
for any convenient \(c>0\).

\begin{theorem}[Positive electroweak radial density chart]
\label{thm:v5v56-radial-chart}
At finite cutoff, the bosonic gauge--Higgs state defines a positive
relative density with respect to the radial reference
\eqref{eq:v5v56-radial-reference}, compact transformed-link Haar measure,
and the residual electromagnetic quotient.  The V55 massive and crossover
factorizations apply to this density.  The singular set \(\rho=0\) is
retained as an orbit-type stratum even though it has zero single-site
Lebesgue measure.
\end{theorem}

\begin{proof}
Theorem~\ref{thm:v5v56-positive-pushforward} gives positivity and the
factor \(\rho^3d\rho\).  Multiplication and division by
\(e^{-c\rho^2}\) merely choose an equivalent positive finite-cutoff
reference.  The V55 factorization theorems apply to every normalized
positive joint density.  The orbit-type statement is
Corollary~\ref{cor:v5v56-stratified-zero}.
\end{proof}

\begin{firewall}[Measure zero does not imply dynamical irrelevance]
\label{fw:v5v56-zero-measure}
Configurations with an exact isolated zero may have zero product
Lebesgue measure, but neighbourhoods of small modulus control crossover
and topological defects.  The cofinal theory cannot remove the singular
stratum using a finite-cutoff measure-zero statement.
\end{firewall}

\subsection{Electroweak radial quotient certificate}

\begin{definition}[Electroweak radial quotient certificate]
\label{def:v5v56-ERQC}
An ERQC consists of:
\begin{enumerate}
 \item the exact compact-gauge pushforward of the positive bosonic
 finite-cutoff measure;
 \item the radial Jacobian \(\rho^3\) and residual
 \(U(1)_{\mathrm{em}}\) quotient;
 \item the principal and zero-modulus orbit strata;
 \item the mass split \(W^\pm,Z\) versus the photon;
 \item BCDC/CSCC estimates for the principal massive block;
 \item the V55 exact handoff when radial or gauge masses degenerate;
 \item a separate anomaly, determinant-phase, BRST, OS, and Einstein
 completion.
\end{enumerate}
\end{definition}

\subsection{V56 status}

\begin{theorem}[V56 conclusion]
\label{thm:v5v56-status}
The finite-cutoff bosonic electroweak gauge--Higgs measure has an exact
positive radial quotient on its principal orbit stratum, with Jacobian
\(\rho^3\).  Its unitary-gauge determinant vanishes at the singular
zero-modulus stratum.  Three broken gauge directions acquire masses
proportional to \(\rho^2\), while the photon remains massless, giving the
precise massive/retained split required by V55.  This proves the bosonic
finite-cutoff part of an ERQC, not the chiral determinant, BRST/OS
continuum limit, or coupling to Einstein gravity.
\end{theorem}

\begin{proof}
Use Theorems~\ref{thm:v5v56-polar},
\ref{thm:v5v56-orbit-strata},
\ref{thm:v5v56-FP},
\ref{thm:v5v56-positive-pushforward}, and
\ref{thm:v5v56-mass-form}.
\end{proof}

V57 will examine anomaly cancellation and determinant-line descent for
one Standard Model generation.  It will distinguish cancellation of the
local gauge anomaly from positivity of the Euclidean chiral determinant
and identify the exact additional datum needed by the relative-density
chart.

\clearpage
\section{V57: one-generation anomaly cancellation and the determinant-phase gate}
\label{sec:v5v57}

V56 produced a positive finite-cutoff density chart for the bosonic
electroweak gauge--Higgs sector.  Fermions cannot be inserted into that
chart merely by citing Standard Model anomaly cancellation.  This note
computes the perturbative gauge and mixed gravitational anomaly sums for
one generation and checks the \(SU(2)\) parity obstruction.  These
cancellations are necessary for gauge descent of the chiral determinant
line.

They do not make the Euclidean chiral determinant nonnegative.  Factoring
the determinant into modulus and phase gives an exact positive
phase-quenched reference, but reconstruction of physical expectations
requires a cutoff-uniform noncancellation bound for the average phase, or
an independent reflection-positive Grassmann construction.  This
determinant-phase gate is the next missing datum in the V49
relative-density programme.

\subsection{One generation in left-handed notation}

Use the gauge algebra
\[
 \mathfrak{su}(3)_c\oplus\mathfrak{su}(2)_L\oplus\mathfrak u(1)_Y
\]
and write every fermion as a left-handed Weyl field.  One generation is
\begin{align}
 Q&:(\mathbf3,\mathbf2)_{1/6},
 &
 u^c&:(\overline{\mathbf3},\mathbf1)_{-2/3},
 &
 d^c&:(\overline{\mathbf3},\mathbf1)_{1/3},
 \label{eq:v5v57-quarks}\\
 L&:(\mathbf1,\mathbf2)_{-1/2},
 &
 e^c&:(\mathbf1,\mathbf1)_{1}.
 \label{eq:v5v57-leptons}
\end{align}
A sterile \(\nu^c:(\mathbf1,\mathbf1)_0\) changes none of the following
sums.

\begin{theorem}[Pure colour anomaly cancellation]
\label{thm:v5v57-SU3-cubic}
Normalize the cubic colour anomaly by
\(A(\mathbf3)=1\) and
\(A(\overline{\mathbf3})=-1\).  Then
\begin{equation}
 \mathcal A_{SU(3)^3}
 =
 2A(\mathbf3)+A(\overline{\mathbf3})
 +A(\overline{\mathbf3})
 =2-1-1=0.
 \label{eq:v5v57-SU3-cubic}
\end{equation}
\end{theorem}

\begin{proof}
The weak doublet \(Q\) contains two colour triplets.  The conjugate
right-handed quarks are two colour antitriplets.  Sum their cubic indices.
\end{proof}

\begin{theorem}[Mixed nonabelian--hypercharge cancellation]
\label{thm:v5v57-mixed}
With \(T(\mathbf3)=T(\mathbf2)=1/2\),
\begin{align}
 \mathcal A_{SU(3)^2Y}
 &=
 2\left(\frac16\right)\frac12
 +\left(-\frac23\right)\frac12
 +\left(\frac13\right)\frac12
 =0,
 \label{eq:v5v57-SU3Y}\\
 \mathcal A_{SU(2)^2Y}
 &=
 3\left(\frac16\right)\frac12
 +\left(-\frac12\right)\frac12
 =0.
 \label{eq:v5v57-SU2Y}
\end{align}
\end{theorem}

\begin{proof}
For \(SU(3)^2Y\), multiply the quark-doublet term by its two weak
components.  For \(SU(2)^2Y\), multiply the quark-doublet term by three
colours.  The displayed arithmetic gives zero.
\end{proof}

\begin{theorem}[Abelian cubic and gravitational cancellation]
\label{thm:v5v57-abelian}
Including colour and weak multiplicities,
\begin{align}
 \mathcal A_{Y^3}
 &=
 6\left(\frac16\right)^3
 +3\left(-\frac23\right)^3
 +3\left(\frac13\right)^3
 +2\left(-\frac12\right)^3
 +1^3
 =0,
 \label{eq:v5v57-Y3}\\
 \mathcal A_{\mathrm{grav}^2Y}
 &=
 6\left(\frac16\right)
 +3\left(-\frac23\right)
 +3\left(\frac13\right)
 +2\left(-\frac12\right)
 +1
 =0.
 \label{eq:v5v57-gravY}
\end{align}
\end{theorem}

\begin{proof}
For the cubic sum, a common denominator \(36\) gives
\[
 1-32+4-9+36=0.
\]
The linear sum is \(1-2+1-1+1=0\).
\end{proof}

\begin{proposition}[Remaining perturbative nonabelian anomalies]
\label{prop:v5v57-remaining-local}
The perturbative \(SU(2)^3\) anomaly vanishes because the doublet is
pseudoreal and has zero symmetric cubic invariant.  Mixed anomalies with
one nonabelian generator and two hypercharge insertions vanish because the
trace of every nonabelian generator is zero.
\end{proposition}

\begin{proof}
For a pseudoreal representation,
\(\operatorname{tr}(T^{(a}T^bT^{c)})=0\).  For one nonabelian generator,
the hypercharge is scalar on each irreducible multiplet, so the group
factor is \(Y^2\operatorname{tr}T^a=0\).
\end{proof}

\subsection{The \(SU(2)\) parity obstruction}

\begin{theorem}[Even number of weak doublets]
\label{thm:v5v57-Witten-count}
One Standard Model generation contains four left-handed
\(SU(2)\) doublets: three colour copies of \(Q\) and one \(L\).  Hence the
mod-two doublet count is
\begin{equation}
 3+1=4\equiv0\pmod2.
 \label{eq:v5v57-Witten-count}
\end{equation}
The usual \(SU(2)\) global sign obstruction detected by an odd number of
fundamental doublets is absent.
\end{theorem}

\begin{proof}
Colour supplies three independent weak doublets in \(Q\), and \(L\)
supplies the fourth.  Four is even.
\end{proof}

\begin{firewall}[This is not a full global-anomaly classification]
\label{fw:v5v57-global-classification}
The parity count addresses the standard \(SU(2)\) mod-two obstruction on
the assumed spin/gauge structure.  The global gauge group may be a
nontrivial quotient, spacetime may carry boundaries or additional
structures, and determinant-line holonomy may involve other bordism
classes.  Those global cases require a separate classification.
\end{firewall}

\subsection{What local anomaly cancellation proves}

Let \(\mathcal L\to\mathcal A/\mathcal G\) denote the determinant line over
gauge fields modulo gauge transformations, on a regulator where this
bundle is defined.

\begin{proposition}[Local descent consequence]
\label{prop:v5v57-local-descent}
The cancellations
\eqref{eq:v5v57-SU3-cubic}--\eqref{eq:v5v57-gravY}, together with
Proposition~\ref{prop:v5v57-remaining-local}, make the local anomaly
polynomial vanish for one generation.  Therefore the local curvature and
infinitesimal gauge variation of the combined determinant line cancel,
provided all multiplets are regulated in one gauge-covariant scheme.
\end{proposition}

\begin{proof}
The six-dimensional anomaly polynomial is a linear combination of the
group invariants computed in the preceding theorems: cubic nonabelian,
nonabelian-squared times \(Y\), \(Y^3\), and \(Yp_1(T)\).  Every
coefficient is zero.  The descent equations identify these coefficients
with the infinitesimal gauge variation and the gauge-direction curvature
of the determinant line.
\end{proof}

\begin{firewall}[Flat is not canonically trivial]
\label{fw:v5v57-flat-not-trivial}
Vanishing local curvature makes the determinant line flat along gauge
directions.  A flat line can still have nontrivial global holonomy.
Gauge-invariant regulator existence and trivialization of every relevant
global holonomy remain required.
\end{firewall}

\begin{firewall}[Anomaly cancellation is not determinant positivity]
\label{fw:v5v57-anomaly-not-positive}
The anomaly sums constrain gauge variation of the determinant phase.
They do not constrain that phase to be zero.  A gauge-invariant complex
number can have any phase, so gauge descent alone cannot make the
fermionic weight a positive density.
\end{firewall}

\subsection{Exact modulus--phase factorization}

At finite cutoff, suppose the combined regulated fermion factor is a
complex measurable function
\begin{equation}
 \mathcal D_F(\Phi)
 =
 |\mathcal D_F(\Phi)|e^{i\Theta(\Phi)}
 \label{eq:v5v57-det-phase}
\end{equation}
of the bosonic fields \(\Phi\), with the phase chosen arbitrarily on its
zero set.  Let
\begin{equation}
 Z_+
 =
 \int e^{-S_B(\Phi)}|\mathcal D_F(\Phi)|\,d\mu_B(\Phi)
 \label{eq:v5v57-Zplus}
\end{equation}
be finite and positive, and define the phase-quenched probability
\begin{equation}
 d\nu_+(\Phi)
 =
 Z_+^{-1}e^{-S_B(\Phi)}
 |\mathcal D_F(\Phi)|\,d\mu_B(\Phi).
 \label{eq:v5v57-phase-quenched}
\end{equation}

\begin{theorem}[Exact phase reconstruction]
\label{thm:v5v57-phase-reconstruction}
If
\begin{equation}
 \zeta:=\mathbb E_{\nu_+}e^{i\Theta}\ne0,
 \label{eq:v5v57-zeta}
\end{equation}
then every integrable physical observable satisfies
\begin{equation}
 \langle\mathcal O\rangle_{\mathrm{phys}}
 =
 \frac{
 \mathbb E_{\nu_+}[\mathcal Oe^{i\Theta}]
 }{
 \mathbb E_{\nu_+}e^{i\Theta}
 }.
 \label{eq:v5v57-phase-reconstruction}
\end{equation}
\end{theorem}

\begin{proof}
The physical numerator is
\[
 Z_+\mathbb E_{\nu_+}[\mathcal Oe^{i\Theta}],
\]
and its partition function is
\(Z_+\mathbb E_{\nu_+}e^{i\Theta}\).  Divide them.
\end{proof}

\begin{theorem}[Phase-denominator stability]
\label{thm:v5v57-phase-stability}
Suppose two cutoff states have phase-quenched laws \(\nu_+,\widetilde\nu_+\),
phases \(\Theta,\widetilde\Theta\), and
\begin{equation}
 |\zeta|,|\widetilde\zeta|\ge c>0.
 \label{eq:v5v57-phase-lower}
\end{equation}
For \(\|\mathcal O\|_\infty\le1\),
\begin{align}
 \left|
 \langle\mathcal O\rangle_{\mathrm{phys}}
 -
 \langle\mathcal O\rangle_{\widetilde{\mathrm{phys}}}
 \right|
 \le&
 2(c^{-1}+c^{-2})
 \|\nu_+-\widetilde\nu_+\|_{\mathrm{TV}}
 \notag\\
 &+
 \left(\frac1c+\frac1{c^2}\right)
 \|e^{i\Theta}-e^{i\widetilde\Theta}\|_
 {L^1(\widetilde\nu_+)}.
 \label{eq:v5v57-phase-stability}
\end{align}
\end{theorem}

\begin{proof}
Write the two ratios as \(N/\zeta\) and
\(\widetilde N/\widetilde\zeta\).  Then
\[
 \left|\frac N\zeta-\frac{\widetilde N}{\widetilde\zeta}\right|
 \le
 \frac{|N-\widetilde N|}{c}
 +
 \frac{|\widetilde N|\,|\zeta-\widetilde\zeta|}{c^2}.
\]
Here \(|\widetilde N|\le1\).  Add and subtract the integrals with the
same phase to obtain
\begin{align*}
 |N-\widetilde N|
 &\le
 2\|\nu_+-\widetilde\nu_+\|_{\mathrm{TV}}
 +
 \|e^{i\Theta}-e^{i\widetilde\Theta}\|_{L^1(\widetilde\nu_+)},\\
 |\zeta-\widetilde\zeta|
 &\le
 2\|\nu_+-\widetilde\nu_+\|_{\mathrm{TV}}
 +
 \|e^{i\Theta}-e^{i\widetilde\Theta}\|_{L^1(\widetilde\nu_+)}.
\end{align*}
Insert these estimates in the ratio bound.
\end{proof}

\begin{firewall}[Average-phase lower bound is substantive]
\label{fw:v5v57-sign-problem}
The modulus--phase identity is algebraic, but a continuum density chart
requires a cutoff-uniform \(c>0\) in
\eqref{eq:v5v57-phase-lower}.  If the average phase tends to zero, the
ratio is unstable and the positive phase-quenched measure does not provide
the physical continuum state.
\end{firewall}

\subsection{Zeros and the Higgs crossover}

\begin{proposition}[Yukawa zeros meet the radial singular stratum]
\label{prop:v5v57-Yukawa-zero}
For a fermion with Yukawa mass \(m_f(\rho)=y_f\rho\), the mass vanishes at
\(\rho=0\).  Therefore the determinant can acquire zero modes on the same
orbit-type stratum where the V56 unitary-gauge Jacobian vanishes.
\end{proposition}

\begin{proof}
This is Proposition~\ref{prop:v5v56-Yukawa-mass} evaluated at
\(\rho=0\).
\end{proof}

\begin{firewall}[The phase is undefined at a determinant zero]
\label{fw:v5v57-phase-zero}
The value assigned to \(e^{i\Theta}\) on
\(\mathcal D_F=0\) does not affect the phase-quenched measure because its
density contains \(|\mathcal D_F|\).  Continuity and topological winding
around the zero can still obstruct a global phase section and must be
included in the determinant-line analysis.
\end{firewall}

\subsection{Alternatives to determinant positivity}

\begin{assessment}[Grassmann route]
\label{ass:v5v57-Grassmann}
One may retain fermions as Grassmann variables and seek reflection
positivity directly for the fermionic observable algebra, rather than
integrating them into a positive bosonic density.  That route avoids an
average-phase denominator but requires a super-Banach or algebraic shell
chart, compatible Berezin integration, gauge descent, and OS
reconstruction.  None is supplied by the bosonic V49 metric.
\end{assessment}

\begin{assessment}[Vectorlike comparison]
\label{ass:v5v57-vectorlike}
In a vectorlike theory, antiunitary or
\(\gamma_5\)-Hermiticity symmetries can pair eigenvalues and sometimes
make the determinant real or nonnegative.  The chiral Standard Model does
not inherit that conclusion without an explicit regulator-level pairing
for the full multiplet.
\end{assessment}

\subsection{Fermionic determinant-line certificate}

\begin{definition}[Fermionic determinant-line certificate]
\label{def:v5v57-FDLC}
An FDLC consists of:
\begin{enumerate}
 \item one gauge-covariant regulator for the complete anomaly-cancelling
 chiral multiplet;
 \item vanishing local anomaly coefficients, as proved above;
 \item trivial global determinant-line holonomy for the chosen global
 gauge group and spacetime structure;
 \item control of determinant zeros across the Higgs radial strata;
 \item either a nonnegative physical fermion factor, a cutoff-uniform
 average-phase lower bound with phase stability, or a reflection-positive
 Grassmann shell construction;
 \item compatibility with the V55 handoff, BRST identities, local
 counterterms, and Einstein coupling.
\end{enumerate}
\end{definition}

\subsection{V57 status}

\begin{theorem}[V57 conclusion]
\label{thm:v5v57-status}
For one Standard Model generation, all perturbative gauge and mixed
gravitational anomaly coefficients vanish, and the usual mod-two
\(SU(2)\) doublet obstruction is absent.  These results establish the
local algebraic anomaly cancellation needed for determinant-line descent.
They do not prove global triviality or positivity.  The exact
phase-quenched representation reconstructs physical expectations only
when its average phase is nonzero, and a stable continuum chart requires
a cutoff-uniform lower bound.  No full FDLC is yet proved.
\end{theorem}

\begin{proof}
Use Theorems~\ref{thm:v5v57-SU3-cubic},
\ref{thm:v5v57-mixed},
\ref{thm:v5v57-abelian},
\ref{thm:v5v57-Witten-count}, and
\ref{thm:v5v57-phase-reconstruction}, with the limitations in
Firewalls~\ref{fw:v5v57-global-classification} and
\ref{fw:v5v57-sign-problem}.
\end{proof}

V58 will avoid assuming determinant positivity and construct the finite-
cutoff Grassmann conditional chart.  It will prove the exact Berezin
Schur-complement formula for a fermion block, identify the norm needed for
conditional source insertions, and test compatibility with the bosonic
relative-density metric.

\clearpage
\section{V58: the finite-cutoff Grassmann conditional chart}
\label{sec:v5v58}

V57 isolated two alternatives for chiral fermions: control the determinant
phase inside a positive bosonic density, or retain the fermions as
Grassmann variables.  This note constructs the second alternative at
finite cutoff.  A weighted coefficient norm makes the finite Grassmann
algebra a Banach algebra, and Berezin differentiation and integration have
explicit norm costs.  For a quadratic fermion action, integrating a
fermion shell gives the exact Schur complement and determinant factor.

If the shell block is singular, the formula is invalid.  Its zero modes
must be re-fibred into the retained Grassmann sector, exactly as the Higgs
radial mode was retained at the V55 crossover.  Combining a positive
bosonic relative density with Grassmann-valued coefficients gives a
Banach module, not a probability space.  Its norm controls algebraic shell
operations but does not by itself prove fermionic reflection positivity.

\subsection{A weighted finite Grassmann algebra}

Let \(\vartheta_1,\ldots,\vartheta_N\) be anticommuting generators.  For
an ordered subset \(I=\{i_1<\cdots<i_k\}\), put
\(\vartheta_I=\vartheta_{i_1}\cdots\vartheta_{i_k}\).  Every element is
\begin{equation}
 F=\sum_{I\subset\{1,\ldots,N\}}F_I\vartheta_I.
 \label{eq:v5v58-Grassmann-expansion}
\end{equation}
For \(r>0\), define
\begin{equation}
 \|F\|_r
 =
 \sum_I |F_I|r^{|I|}.
 \label{eq:v5v58-Grassmann-norm}
\end{equation}

\begin{theorem}[Grassmann Banach algebra]
\label{thm:v5v58-Grassmann-algebra}
For all \(F,G\),
\begin{equation}
 \|FG\|_r\le\|F\|_r\|G\|_r.
 \label{eq:v5v58-Grassmann-product}
\end{equation}
Consequently,
\begin{equation}
 \|e^F\|_r\le e^{\|F\|_r}.
 \label{eq:v5v58-Grassmann-exponential}
\end{equation}
\end{theorem}

\begin{proof}
The product \(\vartheta_I\vartheta_J\) is zero when
\(I\cap J\ne\varnothing\), and otherwise equals a sign times
\(\vartheta_{I\cup J}\), with
\[
 r^{|I\cup J|}=r^{|I|}r^{|J|}.
\]
Take absolute values after collecting coefficients and sum.  This gives
\eqref{eq:v5v58-Grassmann-product}.  Apply it to every term of the finite
exponential series and sum the scalar majorant.
\end{proof}

\begin{theorem}[Derivative and Berezin-integration bounds]
\label{thm:v5v58-calculus-bounds}
For every generator,
\begin{equation}
 \|\partial_{\vartheta_i}F\|_r
 \le r^{-1}\|F\|_r,
 \qquad
 \|\vartheta_iF\|_r\le r\|F\|_r.
 \label{eq:v5v58-derivative-multiply}
\end{equation}
If \(J\) is an ordered set of \(k\) integrated generators, then
\begin{equation}
 \left\|
 \int d\vartheta_J\,F
 \right\|_r
 \le r^{-k}\|F\|_r.
 \label{eq:v5v58-Berezin-bound}
\end{equation}
\end{theorem}

\begin{proof}
Differentiation deletes one generator and multiplication adds one when
the result is nonzero, giving the two weight ratios.  Berezin integration
extracts the coefficient containing every generator in \(J\) and deletes
those \(k\) generators.  Its output weight is therefore \(r^{-k}\) times
the corresponding input weight.
\end{proof}

\begin{assessment}[Finite cutoff is essential]
\label{ass:v5v58-finite}
The algebra has \(2^N\) coefficients and every exponential terminates.
The continuum problem is to choose a compatible family of radii,
localizations, and normalized Gaussian contractions whose constants do
not grow with \(N\).
\end{assessment}

\subsection{Quadratic Berezin Schur complement}

Split complex Grassmann pairs into retained variables
\((\bar\psi,\psi)\) and shell variables
\((\bar\chi,\chi)\).  Let
\begin{equation}
 \mathcal D
 =
 \begin{pmatrix}
  A&B\\
  C&E
 \end{pmatrix},
 \qquad E\ \text{invertible},
 \label{eq:v5v58-D-block}
\end{equation}
and consider
\begin{equation}
 S_F
 =
 \bar\psi A\psi+\bar\psi B\chi
 +\bar\chi C\psi+\bar\chi E\chi.
 \label{eq:v5v58-fermion-action}
\end{equation}

\begin{theorem}[Exact fermion-shell Schur formula]
\label{thm:v5v58-Berezin-Schur}
With a consistent ordered Berezin convention,
\begin{equation}
 \int d\bar\chi\,d\chi\ e^{-S_F}
 =
 \det(E)\,
 \exp\!\left[
 -\bar\psi(A-BE^{-1}C)\psi
 \right].
 \label{eq:v5v58-Berezin-Schur}
\end{equation}
In particular,
\begin{equation}
 \det\mathcal D
 =
 \det E\,
 \det(A-BE^{-1}C).
 \label{eq:v5v58-det-Schur}
\end{equation}
\end{theorem}

\begin{proof}
Make the Berezin translations
\begin{equation}
 \chi'=\chi+E^{-1}C\psi,
 \qquad
 \bar\chi'=\bar\chi+\bar\psi BE^{-1}.
 \label{eq:v5v58-Grassmann-shift}
\end{equation}
Berezin measure is translation invariant.  Direct expansion gives
\[
 S_F
 =
 \bar\chi'E\chi'
 +\bar\psi(A-BE^{-1}C)\psi.
\]
The shell Gaussian integral is \(\det E\), proving
\eqref{eq:v5v58-Berezin-Schur}.  Integrating the retained variables as
well gives \eqref{eq:v5v58-det-Schur}.
\end{proof}

\begin{corollary}[Associativity of exact fermion elimination]
\label{cor:v5v58-associativity}
If two disjoint shell blocks are successively invertible, integrating them
in either order gives the Schur complement of their union and the same
total determinant.
\end{corollary}

\begin{proof}
Both iterated Berezin integrals equal the integral over the union by
finite-dimensional Fubini for coefficient extraction.  Applying
Theorem~\ref{thm:v5v58-Berezin-Schur} in either order yields the same
result.
\end{proof}

\begin{firewall}[No intermediate inverse unless the block is exact]
\label{fw:v5v58-no-formal-inverse}
The factor \(E^{-1}\) is legal because it is the inverse of the exact
quadratic shell block being integrated.  It is not the final RG or Newton
inverse of V42--V45.  Replacing \(E\) by a diagonal approximation without
controlling the remainder would invalidate
\eqref{eq:v5v58-Berezin-Schur}.
\end{firewall}

\subsection{Conditional Grassmann sources}

Add shell sources \(\bar\eta,\eta\) through
\begin{equation}
 S_{F,\eta}
 =
 S_F-\bar\eta\chi-\bar\chi\eta.
 \label{eq:v5v58-source-action}
\end{equation}

\begin{theorem}[Exact conditional source polynomial]
\label{thm:v5v58-source}
Normalized by its value at zero source, the shell integral is
\begin{align}
 \mathcal K_F(\bar\eta,\eta;\bar\psi,\psi)
 &=
 \exp\!\left[
 \bar\eta E^{-1}\eta
 -\bar\eta E^{-1}C\psi
 -\bar\psi BE^{-1}\eta
 \right].
 \label{eq:v5v58-source}
\end{align}
Its finite Grassmann derivatives at zero give every conditional shell
correlator and every retained--shell cross contraction.
\end{theorem}

\begin{proof}
Complete the same square as in
\eqref{eq:v5v58-Grassmann-shift}, now with
\(C\psi-\eta\) and \(\bar\psi B-\bar\eta\).  The shell integral equals
\[
 \det E\,
 e^{-\bar\psi(A-BE^{-1}C)\psi}
 e^{\bar\eta E^{-1}\eta-\bar\eta E^{-1}C\psi
 -\bar\psi BE^{-1}\eta}.
\]
Divide by its zero-source value.  Grassmann differentiation extracts the
finite source coefficients.
\end{proof}

\begin{corollary}[Finite algebraic injectivity]
\label{cor:v5v58-source-injective}
The reduced Schur action together with
\(\mathcal K_F\) determines the matrices
\(E^{-1}\), \(E^{-1}C\), \(BE^{-1}\), and
\(A-BE^{-1}C\).  If \(E\) is known and invertible, these data reconstruct
the full block matrix \(\mathcal D\).
\end{corollary}

\begin{proof}
Read the four matrices from the coefficients of
\(\bar\eta\eta\), \(\bar\eta\psi\), \(\bar\psi\eta\), and
\(\bar\psi\psi\).  Then
\[
 C=E(E^{-1}C),\qquad
 B=(BE^{-1})E,\qquad
 A=(A-BE^{-1}C)+BE^{-1}C.
\]
\end{proof}

\subsection{The zero-mode re-fibration}

\begin{theorem}[Invertibility is necessary]
\label{thm:v5v58-invertibility-necessary}
If \(E\) has a nonzero kernel vector, neither
\eqref{eq:v5v58-Berezin-Schur} nor
\eqref{eq:v5v58-source} is defined as written.  Moreover, the normalized
zero-source shell Gaussian integral vanishes when \(\det E=0\).
\end{theorem}

\begin{proof}
\(E^{-1}\) does not exist, and
\[
 \int d\bar\chi\,d\chi\ e^{-\bar\chi E\chi}=\det E=0.
\]
Division by the zero-source integral is therefore impossible.
\end{proof}

Let \(P_0\) be the spectral projection onto a declared small-eigenvalue
subspace of \(E\), and \(P_1=I-P_0\).  Move
\((P_0\bar\chi,P_0\chi)\) into the retained variables and integrate only
the \(P_1\) block.

\begin{theorem}[Exact fermion zero-mode handoff]
\label{thm:v5v58-zero-handoff}
If
\begin{equation}
 E_1=P_1EP_1
 \label{eq:v5v58-E1}
\end{equation}
is invertible on \(P_1\), then the enlarged retained/shell decomposition
obeys Theorem~\ref{thm:v5v58-Berezin-Schur} with \(E_1\).  No fermion mode
is discarded: the \(P_0\) variables remain explicit retained generators.
\end{theorem}

\begin{proof}
Reorder the finite Grassmann generators so that the \(P_0\) variables are
adjacent to \(\psi\) and the \(P_1\) variables are the shell block.  The
action retains the form \eqref{eq:v5v58-fermion-action} with lower-right
block \(E_1\).  Apply
Theorem~\ref{thm:v5v58-Berezin-Schur}.
\end{proof}

\begin{firewall}[A small singular value cannot be hidden in the shell]
\label{fw:v5v58-small-eigenvalue}
Even before an exact zero, the source norm contains \(E^{-1}\).  A
cutoff-uniform chart requires a proved lower singular-value bound on the
integrated block.  Modes crossing that threshold must be retained by
Theorem~\ref{thm:v5v58-zero-handoff}.
\end{firewall}

\begin{assessment}[Match to the Higgs radial strata]
\label{ass:v5v58-Higgs-match}
Yukawa masses are proportional to the Higgs radius by V56.  As
\(\rho\to0\), the corresponding low fermion modes naturally move from the
massive Grassmann shell into the retained crossover algebra.  This is the
fermionic counterpart of the exact V55 radial re-fibration.
\end{assessment}

\subsection{A Grassmann module over the bosonic density chart}

Let \((X,\nu_B)\) be a positive bosonic probability space and let
\(\mathfrak G_{N,r}\) be the finite Grassmann algebra above.  For a
measurable Grassmann-valued function \(F(x)\), define
\begin{equation}
 \|F\|_{1,\Omega;r}
 =
 \int_X\Omega(x)\|F(x)\|_r\,d\nu_B(x),
 \qquad \Omega\ge1.
 \label{eq:v5v58-module-norm}
\end{equation}

\begin{theorem}[Boson--Grassmann Banach module]
\label{thm:v5v58-module}
The completion of finite measurable functions in
\eqref{eq:v5v58-module-norm} is a Banach module over bounded
Grassmann-valued functions.  In particular,
\begin{equation}
 \|AF\|_{1,\Omega;r}
 \le
 \|A\|_{\infty;r}\|F\|_{1,\Omega;r}.
 \label{eq:v5v58-module-product}
\end{equation}
\end{theorem}

\begin{proof}
Completeness is the Bochner \(L^1\) theorem for the finite-dimensional
Banach space \(\mathfrak G_{N,r}\).  Pointwise
Theorem~\ref{thm:v5v58-Grassmann-algebra} gives
\(\|A(x)F(x)\|_r\le\|A(x)\|_r\|F(x)\|_r\); integrate.
\end{proof}

\begin{theorem}[Bounded super-kernel pushforward]
\label{thm:v5v58-super-kernel}
Let \(K(x,dy)\) be a bosonic Markov kernel satisfying the V49 Lyapunov
condition, and let \(T_{x,y}\) be linear maps between finite Grassmann
algebras with
\begin{equation}
 \|T_{x,y}\|_{r\to r}\le M
 \label{eq:v5v58-T-bound}
\end{equation}
uniformly.  Then
\begin{equation}
 (\mathscr KF)(y)
 =
 \int T_{x,y}F(x)\,K^\leftarrow(y,dx)
 \label{eq:v5v58-super-push}
\end{equation}
obeys the corresponding weighted \(L^1\) bound
\begin{equation}
 \|\mathscr KF\|_{1,\Omega_{\mathrm{out}};r}
 \le M\|F\|_{1,\Omega_{\mathrm{in}};r},
 \label{eq:v5v58-super-bound}
\end{equation}
whenever \(K^\leftarrow\) denotes the disintegration of the joint
pushforward measure.
\end{theorem}

\begin{proof}
Apply the triangle inequality and
\eqref{eq:v5v58-T-bound} inside the conditional integral.  Jensen,
Fubini, and the Lyapunov domination of the output weight by the input
weight give \eqref{eq:v5v58-super-bound}.
\end{proof}

\begin{firewall}[This norm is not a probability metric]
\label{fw:v5v58-not-probability}
Grassmann coefficients are neither ordered nor nonnegative.  The module
norm permits bounded algebraic evolution over a positive bosonic
reference, but total-variation coupling and Dobrushin positivity arguments
do not apply to its fermionic coefficients.
\end{firewall}

\subsection{Cutoff-uniformity requirements}

For an integrated block of \(n\) complex fermion pairs, the raw Berezin
bound \eqref{eq:v5v58-Berezin-bound} costs \(r^{-2n}\).  The determinant
and Wick contractions can compensate this only after normalization.

\begin{firewall}[Raw top-form extraction is volume dependent]
\label{fw:v5v58-volume}
A bound containing \(r^{-2n}\) with \(n\) proportional to cutoff volume
is not a continuum estimate.  The shell chart must factor into local
blocks, normalize each Gaussian integral, and prove a rooted cluster or
determinant estimate before summing blocks.
\end{firewall}

\begin{firewall}[Schur complement may be nonlocal]
\label{fw:v5v58-nonlocal-Schur}
Even when \(E^{-1}\) exists, the term \(BE^{-1}C\) can connect distant
retained variables.  A continuum local action requires exponential or
rooted decay of \(E^{-1}\), compatible local Taylor subtraction, and
uniform control near every re-fibration threshold.
\end{firewall}

\begin{firewall}[Algebra norm does not prove fermionic OS positivity]
\label{fw:v5v58-OS}
Fermionic reflection positivity uses a time reflection, order reversal,
and a graded conjugation on the observable algebra.  Submultiplicativity
of \(\|\cdot\|_r\) has no sign content and cannot establish the OS
inequality.
\end{firewall}

\subsection{Grassmann conditional chart certificate}

\begin{definition}[Grassmann conditional chart certificate]
\label{def:v5v58-GCCC}
A GCCC consists of:
\begin{enumerate}
 \item a gauge-covariant finite-cutoff Grassmann algebra for the complete
 anomaly-cancelling multiplet;
 \item an exact retained/shell block matrix with a uniform inverse on the
 integrated block;
 \item the Schur and source identities
 \eqref{eq:v5v58-Berezin-Schur} and
 \eqref{eq:v5v58-source};
 \item exact re-fibration of every small or zero mode;
 \item a local normalized module norm with cutoff-independent rooted
 bounds for Berezin integration and the Schur kernel;
 \item determinant-line global descent, BRST compatibility, and
 fermionic OS positivity;
 \item compatibility with the V55 radial and V39 continuum strata.
\end{enumerate}
\end{definition}

\subsection{V58 status}

\begin{theorem}[V58 conclusion]
\label{thm:v5v58-status}
At finite cutoff, the Grassmann coefficient norm is submultiplicative,
Berezin operations have explicit bounds, and an invertible quadratic
fermion shell integrates to the exact Schur complement with an injective
conditional source polynomial.  Small and zero modes can be re-fibred
exactly into the retained algebra.  This constructs the algebraic core of
a GCCC over the positive bosonic density chart, but does not provide
cutoff-uniform locality, determinant-line global descent, or fermionic
reflection positivity.
\end{theorem}

\begin{proof}
Use Theorems~\ref{thm:v5v58-Grassmann-algebra},
\ref{thm:v5v58-calculus-bounds},
\ref{thm:v5v58-Berezin-Schur},
\ref{thm:v5v58-source},
\ref{thm:v5v58-zero-handoff}, and
\ref{thm:v5v58-module}.
\end{proof}

V59 will attack locality of the fermionic Schur complement.  It will prove
a Combes--Thomas bound for a gapped finite-range shell block \(E\), insert
that bound into \(BE^{-1}C\), and determine how the decay degenerates as a
Yukawa eigenvalue approaches the retained threshold.

\clearpage
\section{V59: Combes--Thomas locality of the fermionic Schur complement}
\label{sec:v5v59}

V58 showed that a quadratic fermion shell produces the exact retained
operator \(A-BE^{-1}C\).  This note proves locality of that correction
when the integrated block \(E\) has a singular-value gap and finite-range
hopping.  The argument does not require self-adjointness: exponential
conjugation changes \(E\) by a small operator, and the inverse follows
from a Neumann estimate.

The resulting Schur kernel decays exponentially when the cross blocks
\(B,C\) are local.  Its decay length and norm diverge when the fermion gap
closes.  Since Yukawa masses vanish with the Higgs radius, this gives the
quantitative trigger for the V58 zero-mode handoff and the V55 crossover
chart.

\subsection{Exponentially weighted hopping norm}

Let \((\Lambda,d)\) be a finite metric graph and
\[
 \mathscr H=\bigoplus_{x\in\Lambda}\mathbb C^{n_x}.
\]
Write an operator \(E\) in blocks \(E_{xy}\).  For \(\mu\ge0\), define
\begin{align}
 J_\mu^{\mathrm r}
 &=
 \sup_x\sum_y
 \bigl(e^{\mu d(x,y)}-1\bigr)\|E_{xy}\|,
 \label{eq:v5v59-Jrow}\\
 J_\mu^{\mathrm c}
 &=
 \sup_y\sum_x
 \bigl(e^{\mu d(x,y)}-1\bigr)\|E_{xy}\|,
 \label{eq:v5v59-Jcol}\\
 J_\mu&=\sqrt{J_\mu^{\mathrm r}J_\mu^{\mathrm c}}.
 \label{eq:v5v59-Jmu}
\end{align}

\begin{lemma}[Weighted conjugation perturbation]
\label{lem:v5v59-conjugation}
Fix \(y_0\in\Lambda\) and let
\[
 (W_{\mu,y_0}f)(x)=e^{\mu d(x,y_0)}f(x).
\]
Then
\begin{equation}
 \|W_{\mu,y_0}EW_{\mu,y_0}^{-1}-E\|
 \le J_\mu.
 \label{eq:v5v59-conjugation}
\end{equation}
\end{lemma}

\begin{proof}
The \(x,y\) block of the difference is
\[
 \left(
 e^{\mu(d(x,y_0)-d(y,y_0))}-1
 \right)E_{xy}.
\]
The triangle inequality gives
\[
 \left|
 e^{\mu(d(x,y_0)-d(y,y_0))}-1
 \right|
 \le e^{\mu d(x,y)}-1.
\]
The block Schur test bounds the operator norm by the geometric mean of the
maximum row and column sums, which is \(J_\mu\).
\end{proof}

\subsection{A nonnormal Combes--Thomas theorem}

Assume
\begin{equation}
 \|E^{-1}\|\le m^{-1}
 \label{eq:v5v59-singular-gap}
\end{equation}
for some \(m>0\).  This is a lower singular-value bound, not an eigenvalue
gap.

\begin{theorem}[Singular-gap Combes--Thomas estimate]
\label{thm:v5v59-CT}
If
\begin{equation}
 J_\mu<m,
 \label{eq:v5v59-CT-condition}
\end{equation}
then
\begin{equation}
 \|(E^{-1})_{xy}\|
 \le
 \frac{e^{-\mu d(x,y)}}{m-J_\mu}.
 \label{eq:v5v59-CT}
\end{equation}
\end{theorem}

\begin{proof}
Put
\[
 E_{\mu,y}=W_{\mu,y}EW_{\mu,y}^{-1}=E+K_{\mu,y}.
\]
By Lemma~\ref{lem:v5v59-conjugation},
\(\|K_{\mu,y}\|\le J_\mu\).  Hence
\[
 E_{\mu,y}^{-1}
 =
 (I+E^{-1}K_{\mu,y})^{-1}E^{-1}
\]
and the Neumann series gives
\[
 \|E_{\mu,y}^{-1}\|
 \le\frac1{m-J_\mu}.
\]
Its \(x,y\) block is
\[
 (E_{\mu,y}^{-1})_{xy}
 =
 e^{\mu d(x,y)}(E^{-1})_{xy},
\]
because the weight equals one at \(y\).  Rearrangement proves
\eqref{eq:v5v59-CT}.
\end{proof}

\begin{corollary}[Finite-range explicit exponent]
\label{cor:v5v59-finite-range}
Suppose \(E_{xy}=0\) for \(d(x,y)>R\), and
\begin{equation}
 J_0^*
 =
 \max\left\{
 \sup_x\sum_{y\ne x}\|E_{xy}\|,
 \sup_y\sum_{x\ne y}\|E_{xy}\|
 \right\}<\infty.
 \label{eq:v5v59-J0}
\end{equation}
If
\begin{equation}
 0<\mu\le
 \frac1R\log\left(1+\frac{m}{2J_0^*}\right),
 \label{eq:v5v59-mu-choice}
\end{equation}
then
\begin{equation}
 \|(E^{-1})_{xy}\|
 \le\frac2m e^{-\mu d(x,y)}.
 \label{eq:v5v59-CT-explicit}
\end{equation}
\end{corollary}

\begin{proof}
Finite range gives
\[
 J_\mu^{\mathrm r},J_\mu^{\mathrm c}
 \le(e^{\mu R}-1)J_0^*\le\frac m2.
\]
Thus \(J_\mu\le m/2\), and
Theorem~\ref{thm:v5v59-CT} applies.
\end{proof}

\begin{firewall}[Eigenvalues do not control a nonnormal inverse]
\label{fw:v5v59-nonnormal}
For a nonnormal chiral block, eigenvalues bounded away from zero do not
imply \eqref{eq:v5v59-singular-gap}.  The smallest singular value or a
direct inverse norm is required.
\end{firewall}

\subsection{Locality of the Schur correction}

Let retained sites be indexed by \(\Lambda_P\) and shell sites by
\(\Lambda_Q\), inside a common metric graph.  Suppose
\begin{align}
 B_{xu}&=0 &&\text{if }d(x,u)>R_B,
 \label{eq:v5v59-B-range}\\
 C_{vy}&=0 &&\text{if }d(v,y)>R_C,
 \label{eq:v5v59-C-range}
\end{align}
and define
\begin{equation}
 b=\sup_x\sum_u\|B_{xu}\|,
 \qquad
 c=\sup_y\sum_v\|C_{vy}\|.
 \label{eq:v5v59-bc}
\end{equation}

\begin{theorem}[Exponential Schur-kernel locality]
\label{thm:v5v59-Schur-locality}
Under Theorem~\ref{thm:v5v59-CT},
\begin{equation}
 \|(BE^{-1}C)_{xy}\|
 \le
 \frac{bc\,e^{\mu(R_B+R_C)}}{m-J_\mu}
 e^{-\mu d(x,y)}.
 \label{eq:v5v59-Schur-locality}
\end{equation}
\end{theorem}

\begin{proof}
Expand
\[
 (BE^{-1}C)_{xy}
 =
 \sum_{u,v}B_{xu}(E^{-1})_{uv}C_{vy}.
\]
On the support of a term,
\[
 d(u,v)\ge d(x,y)-R_B-R_C.
\]
Apply \eqref{eq:v5v59-CT}, then sum the \(B\) row and \(C\) column
norms using \eqref{eq:v5v59-bc}.
\end{proof}

\begin{corollary}[Rooted summability]
\label{cor:v5v59-rooted}
If the graph has volume growth
\begin{equation}
 \#\{y:d(x,y)=n\}\le C_Ve^{\sigma n}
 \label{eq:v5v59-volume-growth}
\end{equation}
with \(\sigma<\mu\), then every Schur row is absolutely summable and
\begin{equation}
 \sup_x\sum_y\|(BE^{-1}C)_{xy}\|<\infty.
 \label{eq:v5v59-rooted-sum}
\end{equation}
\end{corollary}

\begin{proof}
Insert \eqref{eq:v5v59-Schur-locality}, group sites by distance, and sum
the geometric series with ratio \(e^{-(\mu-\sigma)}\).
\end{proof}

\begin{firewall}[Decay must beat graph growth]
\label{fw:v5v59-growth}
Pointwise exponential decay is not a rooted operator bound on a graph
whose shell multiplicity grows at the same or a faster exponential rate.
The strict inequality \(\mu>\sigma\), or an appropriate polynomial-growth
analogue, is part of the locality certificate.
\end{firewall}

\subsection{Yukawa degeneration and the handoff threshold}

Assume a vectorlike comparison block has the form
\begin{equation}
 E(\rho)=D_0+y\rho,
 \label{eq:v5v59-Yukawa-block}
\end{equation}
where \(D_0^*=-D_0\), \(y\rho\) is a real scalar mass, and \(y\rho\ne0\).

\begin{lemma}[Vectorlike singular gap]
\label{lem:v5v59-vectorlike-gap}
For every \(\psi\),
\begin{equation}
 \|E(\rho)\psi\|^2
 =
 \|D_0\psi\|^2+y^2\rho^2\|\psi\|^2.
 \label{eq:v5v59-vectorlike-gap}
\end{equation}
Hence
\begin{equation}
 \|E(\rho)^{-1}\|\le\frac1{|y|\rho}.
 \label{eq:v5v59-vectorlike-inverse}
\end{equation}
\end{lemma}

\begin{proof}
The cross term is
\[
 2y\rho\,\operatorname{Re}\langle D_0\psi,\psi\rangle=0
\]
because \(D_0\) is anti-Hermitian.  The lower bound and inverse estimate
follow.
\end{proof}

\begin{corollary}[Radial locality scale]
\label{cor:v5v59-radial-locality}
In the vectorlike comparison, Corollary~\ref{cor:v5v59-finite-range}
permits
\begin{equation}
 \mu(\rho)
 \le
 \frac1R
 \log\left(1+\frac{|y|\rho}{2J_0^*}\right).
 \label{eq:v5v59-radial-mu}
\end{equation}
As \(\rho\downarrow0\),
\begin{equation}
 \mu(\rho)
 =
 \frac{|y|\rho}{2RJ_0^*}+O(\rho^2),
 \label{eq:v5v59-radial-mu-small}
\end{equation}
while the inverse prefactor grows as \(\rho^{-1}\).
\end{corollary}

\begin{proof}
Use \(m=|y|\rho\) in
\eqref{eq:v5v59-mu-choice} and expand the logarithm.
\end{proof}

\begin{theorem}[Quantitative fermion re-fibration trigger]
\label{thm:v5v59-refibration-trigger}
Fix required constants \(m_{\min}>0\) and
\(\mu_{\min}>0\).  A mode may remain in the local integrated fermion shell
only where
\begin{equation}
 s_{\min}(E(\rho))\ge m_{\min},
 \qquad
 J_{\mu_{\min}}<m_{\min}.
 \label{eq:v5v59-shell-threshold}
\end{equation}
When either inequality fails, V58 requires that mode to be moved to the
retained Grassmann sector.  In the vectorlike Yukawa comparison, the first
failure occurs no later than
\begin{equation}
 \rho<\frac{m_{\min}}{|y|}.
 \label{eq:v5v59-rho-threshold}
\end{equation}
\end{theorem}

\begin{proof}
The two inequalities are exactly the hypotheses needed for a uniform
inverse and the desired Combes--Thomas exponent.  If they fail, the V58
Schur and source norms lose their declared constants.  The last statement
uses Lemma~\ref{lem:v5v59-vectorlike-gap}.
\end{proof}

\begin{assessment}[Chiral use]
\label{ass:v5v59-chiral-use}
Lemma~\ref{lem:v5v59-vectorlike-gap} is not asserted for the regulated
chiral Standard Model block.  In that block,
\eqref{eq:v5v59-shell-threshold} must be verified directly after anomaly
and determinant-line assembly.  The threshold formulation remains valid
because it uses singular values rather than vectorlike spectral pairing.
\end{assessment}

\subsection{Moving projectors}

Let \(P_0(\rho)\) be a Riesz projector onto singular values below a
threshold, defined through the self-adjoint operator \(E(\rho)^*E(\rho)\).

\begin{proposition}[Stable retained projector]
\label{prop:v5v59-projector}
Suppose the selected part of the spectrum of \(E^*E\) is separated from
the remainder by a contour gap \(\gamma>0\), and
\[
 \|\partial_\rho(E^*E)\|\le L.
\]
Then the Riesz projector is differentiable and
\begin{equation}
 \|\partial_\rho P_0(\rho)\|
 \le
 \frac{C_\Gamma L}{\gamma^2},
 \label{eq:v5v59-projector-bound}
\end{equation}
where \(C_\Gamma\) depends on the length of the fixed contour.
\end{proposition}

\begin{proof}
Write
\[
 P_0=\frac1{2\pi i}\oint_\Gamma
 (z-E^*E)^{-1}\,dz.
\]
Differentiate under the integral.  Each resolvent has norm at most
\(\gamma^{-1}\), giving the stated bound after integration.
\end{proof}

\begin{firewall}[Threshold crossing is a stratum]
\label{fw:v5v59-threshold-crossing}
At a singular-value crossing, \(\gamma\to0\) and
\eqref{eq:v5v59-projector-bound} diverges.  A single smooth
retained/shell splitting cannot pass through the crossing; the V39/V55
stratified quotient must identify the adjacent projector charts.
\end{firewall}

\subsection{Fermionic Schur locality certificate}

\begin{definition}[Fermionic Schur locality certificate]
\label{def:v5v59-FSLC}
An FSLC consists of:
\begin{enumerate}
 \item a gauge-covariant finite-range integrated block \(E\);
 \item a uniform lower singular-value bound and weighted hopping bound
 \(J_\mu<m\);
 \item local cross blocks \(B,C\) satisfying
 \eqref{eq:v5v59-B-range}--\eqref{eq:v5v59-bc};
 \item rooted graph growth strictly below the decay exponent;
 \item exact re-fibration at every singular-value or projector-gap
 threshold;
 \item compatibility with determinant-line descent, the GCCC, radial
 handoff, and fermionic OS structure.
\end{enumerate}
\end{definition}

\subsection{V59 status}

\begin{theorem}[V59 conclusion]
\label{thm:v5v59-status}
A finite-range fermion shell block with a uniform singular-value gap has
an exponentially local inverse even when it is nonnormal, and its exact
Schur correction \(BE^{-1}C\) is exponentially local.  The bounds
degenerate as the Yukawa mass or singular-value gap closes, giving the
explicit re-fibration trigger
\eqref{eq:v5v59-shell-threshold}.  This proves the analytic core of an
FSLC for a gapped finite matrix block, not the required gap, global
determinant-line descent, or OS positivity for the regulated chiral
Standard Model.
\end{theorem}

\begin{proof}
Use Theorems~\ref{thm:v5v59-CT},
\ref{thm:v5v59-Schur-locality}, and
\ref{thm:v5v59-refibration-trigger}.
\end{proof}

V60 will perform the compulsory YM/NS audit.  It will then combine the
bosonic radial, massive gauge, retained photon, and fermionic projector
strata into one finite-cutoff electroweak shell atlas, with exact overlap
maps and a ledger of the still-missing positivity and continuum-uniformity
conditions.

\clearpage
\section{V60: finite-cutoff electroweak shell atlas and exact common refinements}
\label{chap:v5v60-ew-atlas}
\label{sec:v5v60}

\section{Compulsory V60 companion audit}

Before making the next Standard Model reduction, we inspected the newest
live Yang--Mills and Navier--Stokes notes in the shared directory.  The
files inspected were
\begin{align*}
 &\texttt{Renewal\_Yang\_Mills\_Mass\_Gap\_Research\_Notes\_V66.tex},
 &&\operatorname{SHA256}=
 \texttt{58665C5F01FA47B69F2D9553245E0A93}\notag\\
 &&&\hspace{32mm}\texttt{A5CF1CE9EB30721EFA7436C752C7E430},\\
 &\texttt{Renewal\_NS\_Stability\_Research\_Notes\_V31.tex},
 &&\operatorname{SHA256}=
 \texttt{F1121B62238AD5491BB7CF03635EA993}\notag\\
 &&&\hspace{32mm}\texttt{4942EDF573ED8FAAA9EE79E1D38A6696}.
\end{align*}
The files named V65 and V66 in the Yang--Mills sequence are byte-for-byte
identical.  Thus V66 is the newest filename, but its mathematical endpoint
is the V65 chapter.  No companion file was modified.

The relevant Yang--Mills result is negative as well as constructive.
Pruefer-type exact summation removes a combinatorial multiplicity, but
rowwise first moments do not control the remaining all-thermal anisotropy.
The V65 countermodel shows that a sharp all-order Wilson/Bessel tail, of
the schematic form
\begin{equation}
 \frac{I_n(\alpha)}{I_1(\alpha)}
 \lesssim
 \exp\{-c\alpha h(n/\alpha)\},
 \qquad
 h(t)=t\operatorname{arsinh}t-\sqrt{1+t^2}+1,
 \label{eq:v5v60-YM-tail}
\end{equation}
is genuinely new analytic input.  The Navier--Stokes endpoint remains V31
and adds no new change to the source/concentration interface used here.

The transfer to the present programme is precise: exact shell summation
does not create the smallness required for a uniform continuum limit, and
a finite list of radial or fermionic moments cannot replace an all-order
tail.  We therefore glue shell descriptions only by exact common
refinement.  Any mode for which the declared analytic estimate fails is
retained as a coordinate rather than hidden in an uncontrolled remainder.

\section{Finite-cutoff field space and quantitative cells}

Fix a finite lattice or spectral cutoff.  Let \(X\) be the resulting
finite-dimensional bosonic configuration space and let \(H_F\) be the
finite-dimensional one-particle fermion space.  Gauge fixing and all
finite-dimensional Jacobians are regarded as part of the bosonic density.
The photon subspace is denoted by \(H_\gamma\); it is retained in every
cell.

\begin{definition}[Electroweak cell datum]
\label{def:v5v60-cell}
An electroweak cell datum
\(
 \mathfrak C=(U_\mathfrak C,R^B_\mathfrak C,
 S^B_\mathfrak C,R^F_\mathfrak C,S^F_\mathfrak C)
\)
consists of:
\begin{enumerate}
 \item a Borel set \(U_\mathfrak C\subset X\);
 \item a bosonic retained/shell splitting
 \(
 X=R^B_\mathfrak C\times S^B_\mathfrak C
 \)
 in local coordinates, with \(H_\gamma\subset R^B_\mathfrak C\);
 \item an orthogonal fermion splitting
 \(
 H_F=R^F_\mathfrak C\oplus S^F_\mathfrak C;
 \)
 \item a positive conditional bosonic kernel on
 \(S^B_\mathfrak C\) and a fixed Berezin orientation on
 \(S^F_\mathfrak C\);
 \item invertibility of the fermion block compressed to
 \(S^F_\mathfrak C\).
\end{enumerate}
The radial variable is placed in \(S^B_\mathfrak C\) only in a massive
cell and in \(R^B_\mathfrak C\) in a crossover cell.
\end{definition}

The preceding notes produced two quantitative gates.  We now put them in
one cell definition.  Let
\begin{equation}
 q_{\rm rad}(x):=
 \frac{\varepsilon_{\rm rad}(x)}{1-\vartheta_{\rm rad}(x)},
 \qquad \vartheta_{\rm rad}(x)<1,
 \label{eq:v5v60-radial-quality}
\end{equation}
and, for a proposed integrated fermion block \(E_\mathfrak C(x)\), let
\begin{equation}
 q_F(x;m,\mu,\kappa):=
 \left\{
 \begin{array}{l}
 s_{\min}(E_\mathfrak C(x))\ge m,\\
 J_\mu(E_\mathfrak C(x))\le(1-\kappa)m
 \end{array}
 \right. ,
 \qquad 0<\kappa<1.
 \label{eq:v5v60-fermion-quality}
\end{equation}

\begin{definition}[Massive and crossover cells]
\label{def:v5v60-massive-crossover}
For \(r,\delta_{\rm rad},m,\mu>0\) and \(0<\kappa<1\), the massive-safe
cell is the set on which
\begin{equation}
 \rho\ge r,
 \qquad q_{\rm rad}\le\delta_{\rm rad},
 \qquad q_F(x;m,\mu,\kappa)\text{ holds}.
 \label{eq:v5v60-safe-cell}
\end{equation}
There the certified radial and massive \(W/Z\) variables may be integrated,
as may the certified fermion shell.  On the complementary crossover cells,
the failed radial variable, every failed massive vector block, and every
fermion singular band violating \eqref{eq:v5v60-fermion-quality} are moved
to the retained sector.  The photon remains retained on both kinds of
cell.
\end{definition}

This definition is algorithmic at finite cutoff: there are finitely many
singular values and finitely many candidate blocks.  It makes no assertion
that the number of cells or their constants are uniform in the cutoff.

\begin{lemma}[Quantitative locality on a safe cell]
\label{lem:v5v60-safe-locality}
On a cell satisfying \eqref{eq:v5v60-fermion-quality}, the integrated
fermion inverse obeys
\begin{equation}
 \|(E_\mathfrak C^{-1})_{xy}\|
 \le \frac{1}{\kappa m}e^{-\mu d(x,y)}.
 \label{eq:v5v60-safe-inverse}
\end{equation}
If the cross blocks \(B,C\) have range at most \(R_0\) and norms bounded by
\(b,c\), then
\begin{equation}
 \|(BE_\mathfrak C^{-1}C)_{xy}\|
 \le \frac{bc}{\kappa m}
 e^{2\mu R_0}e^{-\mu d(x,y)}.
 \label{eq:v5v60-safe-Schur}
\end{equation}
\end{lemma}

\begin{proof}
The V59 weighted conjugation proof gives the inverse prefactor
\((m-J_\mu)^{-1}\).  The second inequality in
\eqref{eq:v5v60-fermion-quality} implies
\(m-J_\mu\ge\kappa m\), proving
\eqref{eq:v5v60-safe-inverse}.  Inserting the two finite-range cross blocks
can change each endpoint by at most \(R_0\), which gives
\eqref{eq:v5v60-safe-Schur}.
\end{proof}

\section{The exact refinement calculus}

The word ``atlas'' requires care.  Integration loses information, so a
coarser reduced density cannot in general be inverted to recover a finer
one.  The exact object is a projective atlas: two reductions are compared
through a common finer retained sector, and all routes from that refinement
to either coarser description commute.

\subsection{Bosonic disintegration}

Let \(Y=Y_0\times Y_1\times Y_2\), and suppose a finite positive measure
has the disintegration
\begin{equation}
 d\nu(y_0,y_1,y_2)
 =d\nu_0(y_0)
 K_1(y_0,dy_1)K_2(y_0,y_1,dy_2).
 \label{eq:v5v60-disintegration}
\end{equation}
Write \(\mathcal I_j\) for conditional integration in \(Y_j\).

\begin{lemma}[Bosonic refinement identity]
\label{lem:v5v60-bosonic-Fubini}
For every \(f\in L^1(\nu)\),
\begin{equation}
 \mathcal I_1\mathcal I_2 f
 =\mathcal I_{1,2}f
 \quad\text{in }L^1(\nu_0).
 \label{eq:v5v60-bosonic-Fubini}
\end{equation}
The same identity holds after restricting to a measurable atlas cell and
renormalising only when the normalising factor is recorded as part of the
retained density.
\end{lemma}

\begin{proof}
The first claim is Tonelli's theorem applied to (|f|), followed by
Fubini's theorem.  For a restriction \(1_Uf\), the unnormalised identity
is unchanged.  Conditional normalisation divides by a retained-coordinate
dependent partition function; recording that factor is therefore exactly
what restores the identity.
\end{proof}

The radial factorisation of V55 is an instance of
Lemma~\ref{lem:v5v60-bosonic-Fubini}: the density-valued state, rather than
only a collection of moments, is the overlap coordinate.

\subsection{Berezin refinement}

Let \(V=R\oplus T\oplus S\) be an oriented finite-dimensional fermion
space.  For a Grassmann polynomial \(F\in\Lambda V^*\), denote Berezin
integration over a summand by \(\mathcal B_R,\mathcal B_T,\mathcal B_S\).
We fix the orientation convention
\(
 \operatorname{or}(V)=
 \operatorname{or}(R)\wedge\operatorname{or}(T)\wedge
 \operatorname{or}(S)
\).

\begin{lemma}[Berezin refinement identity]
\label{lem:v5v60-Berezin-Fubini}
With the displayed orientation convention,
\begin{equation}
 \mathcal B_T\mathcal B_S F
 =\mathcal B_{T\oplus S}F.
 \label{eq:v5v60-Berezin-Fubini}
\end{equation}
If another ordering is used, the two sides differ only by the determinant
sign of the ordering permutation.
\end{lemma}

\begin{proof}
Expand \(F\) in the ordered exterior basis.  Both sides vanish on every
monomial that does not contain the top exterior monomial of \(T\oplus S\).
On a monomial that does contain it, both sides extract the same coefficient.
Reordering the extraction variables multiplies the top exterior monomial
by the determinant sign of the permutation.
\end{proof}

\begin{proposition}[Associativity of fermionic Schur elimination]
\label{prop:v5v60-Schur-associativity}
Suppose the quadratic fermion matrix is written in blocks for
\(R\oplus T\oplus S\), and suppose the blocks eliminated along either
route are invertible.  Eliminating \(S\) and then \(T\) gives the same
effective quadratic form on \(R\) and the same determinant factor as
eliminating \(T\oplus S\) at once.
\end{proposition}

\begin{proof}
Apply Lemma~\ref{lem:v5v60-Berezin-Fubini} to
\(
 F=\exp(-\bar\psi D\psi)
\).
The finite Grassmann Gaussian identity identifies each side with its Schur
complement times its determinant factor.  Equality as Grassmann
polynomials proves both assertions.  Equivalently, it is block Gaussian
elimination together with the block determinant identity.
\end{proof}

The proposition includes source polynomials: insert arbitrary retained
Grassmann sources before applying the lemma.  Consequently the exact V58
source polynomial, rather than only the determinant, is the datum that
survives a change of retained/shell splitting.

\subsection{Common refinement of two cells}

\begin{definition}[Admissible common refinement]
\label{def:v5v60-common-refinement}
Two cells \(\mathfrak C_1,\mathfrak C_2\) have an admissible common
refinement \(\mathfrak R\) on \(U_1\cap U_2\) if:
\begin{enumerate}
 \item \(R^B_{\mathfrak R}\) contains both bosonic retained sectors and
 \(R^F_{\mathfrak R}\) contains both fermionic retained sectors;
 \item the shell of \(\mathfrak R\) is the common complement and every
 block integrated on a map \(\mathfrak R\to\mathfrak C_i\) satisfies the
 required bosonic and fermionic invertibility bounds;
 \item the conditional kernels and Berezin orientations are restrictions
 of one full-cutoff density and one determinant-line orientation.
\end{enumerate}
\end{definition}

At finite cutoff a canonical candidate is obtained by retaining the span
of the two retained spectral bands, the radial coordinate if either cell
retains it, and every massive vector band retained by either cell.  If the
common complement fails a declared bound, its failed band is retained as
well.  Since the dimension is finite, this enlargement terminates, in the
worst case at the unreduced full theory.

\begin{theorem}[Exact common-refinement square]
\label{thm:v5v60-refinement-square}
Let \(F\) be an integrable finite-cutoff superdensity and let
\(\mathfrak C_1,\mathfrak C_2\) admit a common refinement
\(\mathfrak R\).  Let \(\Pi_{\mathfrak R\mathfrak C_i}\) denote exact
integration of the additional variables retained in
\(\mathfrak R\) but shelled in \(\mathfrak C_i\).  Then
\begin{equation}
 \Pi_{\mathfrak R\mathfrak C_i}
 \bigl(\mathcal I_{\mathfrak R}F\bigr)
 =\mathcal I_{\mathfrak C_i}F,
 \qquad i=1,2.
 \label{eq:v5v60-refinement-square}
\end{equation}
Moreover, if a third cell \(\mathfrak C_3\) is coarser than
\(\mathfrak C_i\), then
\begin{equation}
 \Pi_{\mathfrak C_i\mathfrak C_3}
 \Pi_{\mathfrak R\mathfrak C_i}
 =\Pi_{\mathfrak R\mathfrak C_3}.
 \label{eq:v5v60-projective-cocycle}
\end{equation}
\end{theorem}

\begin{proof}
Split the bosonic variables into those shelled already in
\(\mathfrak R\), those integrated on the map to
\(\mathfrak C_i\), and those retained by \(\mathfrak C_i\).
Lemma~\ref{lem:v5v60-bosonic-Fubini} proves the bosonic identity.  Make the
analogous threefold split of the Grassmann generators.
Lemma~\ref{lem:v5v60-Berezin-Fubini} proves the fermionic identity, with
the determinant-line orientation removing the possible ordering sign.
Their tensor product gives \eqref{eq:v5v60-refinement-square}.  Applying
the same two Fubini identities to four nested groups of variables gives
\eqref{eq:v5v60-projective-cocycle}.
\end{proof}

\begin{firewall}[No inverse transition is asserted]
\label{fw:v5v60-no-inverse}
The maps in Theorem~\ref{thm:v5v60-refinement-square} run from a finer
retained description to a coarser one.  They need not be injective.
Calling them ordinary invertible coordinate changes would be false.  The
common refinement and cocycle identities are the exact replacement.
\end{firewall}

\section{Normed overlap control}

Let
\(
 F(x,\theta)=\sum_I F_I(x)\theta_I
\)
be a superdensity on a cell.  For \(r>0\), set
\begin{equation}
 \|F\|_{\Omega,r}
 :=\int_\Omega\sum_I |F_I(x)|r^{|I|}\,d\nu(x).
 \label{eq:v5v60-superdensity-norm}
\end{equation}
For two representatives on the same refinement, define
\begin{equation}
 d_{\Omega,r}(F,G):=\|F-G\|_{\Omega,r}.
 \label{eq:v5v60-superdensity-distance}
\end{equation}

\begin{lemma}[Bosonic and Grassmann contraction]
\label{lem:v5v60-integration-contraction}
Conditional integration against a probability kernel is a contraction in
the coefficientwise \(L^1\) part of
\eqref{eq:v5v60-superdensity-norm}.  Berezin integration over \(k\)
generators obeys
\begin{equation}
 \|\mathcal B_SF\|_{\Omega,r}
 \le r^{-k}\|F\|_{\Omega,r}.
 \label{eq:v5v60-Berezin-bound}
\end{equation}
\end{lemma}

\begin{proof}
The first statement is Jensen's inequality followed by Tonelli's theorem.
Berezin integration discards all coefficients except those containing the
top monomial of \(S\), and lowers their degrees by \(k\).  Each surviving
weight \(r^{|I|-k}\) equals \(r^{-k}r^{|I|}\), proving the second claim.
\end{proof}

\begin{lemma}[Finite basis-change bound]
\label{lem:v5v60-basis-change}
Let the Grassmann generators transform by \(\theta'=U\theta\), and let
\(L=\max_j\sum_i|U_{ij}|\).  The induced exterior-algebra map satisfies
\begin{equation}
 \|\Lambda U(F)\|_{\Omega,r}
 \le \|F\|_{\Omega,Lr}.
 \label{eq:v5v60-basis-change}
\end{equation}
If both \(U\) and \(U^{-1}\) have column sums bounded by \(L\), the two
coordinate norms are equivalent with the corresponding radius changes.
\end{lemma}

\begin{proof}
Each transformed generator has coefficient norm at most \(Lr\).  The V58
submultiplicative Grassmann norm bounds a transformed degree-\(k\) monomial
by \((Lr)^k\).  Sum the coefficient bounds over all monomials and integrate
over \(\Omega\).
\end{proof}

The dependence on \(L\), on the number of integrated generators in
\eqref{eq:v5v60-Berezin-bound}, and on conditional normalising factors is
harmless at a fixed cutoff.  Uniform control of exactly these quantities is
part of the continuum problem and is not supplied by finite-dimensional
norm equivalence.

\section{No double counting and symmetry descent}

\begin{proposition}[Each physical mode occurs once]
\label{prop:v5v60-no-double-counting}
Within every cell, every cutoff mode belongs to exactly one of the retained
or shell sectors.  On an overlap, the common-refinement construction does
not add a second copy of a mode.  It moves a disputed mode to the retained
sector and derives each coarser description by one exact integration.
\end{proposition}

\begin{proof}
The bosonic coordinate splitting is a product decomposition and the
fermionic splitting is a direct sum, so retained and shell sectors are
disjoint and exhaustive.  The common refinement takes the span of disputed
retained subspaces inside the original cutoff space; it is not an external
direct sum of copies.  Equations~\eqref{eq:v5v60-refinement-square} and
\eqref{eq:v5v60-projective-cocycle} then integrate each additional variable
exactly once.
\end{proof}

\begin{proposition}[Equivariance of exact shell maps]
\label{prop:v5v60-equivariance}
Let a compact residual gauge group \(G\) preserve a cell, its retained and
shell splittings, its conditional bosonic kernel, and its determinant-line
orientation.  Then
\begin{equation}
 \mathcal I_\mathfrak C(g\cdot F)
 =g\cdot\mathcal I_\mathfrak C(F),
 \qquad g\in G.
 \label{eq:v5v60-equivariance}
\end{equation}
The refinement maps therefore descend to the finite-cutoff quotient.
\end{proposition}

\begin{proof}
Change variables by \(g\) in the finite-dimensional bosonic integral.  The
kernel invariance removes its Jacobian.  The Berezin change of variables
contributes the determinant-line action, which is trivial under the stated
orientation hypothesis.  The retained variables carry the remaining
action, giving \eqref{eq:v5v60-equivariance}.  Apply the same argument to
the nested integrations defining a refinement map.
\end{proof}

The orientation hypothesis is deliberately explicit.  The one-generation
anomaly sums proved in V57 are necessary consistency checks, but they do
not by themselves construct the global regulated chiral determinant line
or trivialise its phase.

\section{Why finitely many moments cannot close the atlas}

The companion audit can be stated as a general obstruction independent of
the detailed Yang--Mills countermodel.

\begin{lemma}[Finite constraints do not determine a new observable]
\label{lem:v5v60-finite-moment-obstruction}
Let \(\Omega=\{1,\ldots,N\}\), let
\(L_0,L_1,\ldots,L_k,T\) be linear functionals on
\(\mathbb R^N\), where \(L_0\) is total mass.  If
\(
 T\notin\operatorname{span}\{L_0,\ldots,L_k\}
\)
and \(N>k+1\), then there are two strictly positive probability vectors
\(p_+,p_-\) such that
\begin{equation}
 L_j(p_+)=L_j(p_-),\quad 0\le j\le k,
 \qquad T(p_+)\ne T(p_-).
 \label{eq:v5v60-moment-counterpair}
\end{equation}
\end{lemma}

\begin{proof}
Linear algebra gives a vector \(v\) in the common kernel of
\(L_0,\ldots,L_k\) with \(T(v)\ne0\).  Let \(p_0\) be the uniform
probability vector.  For sufficiently small \(\epsilon>0\), both
\(p_\pm=p_0\pm\epsilon v\) are strictly positive.  They have mass one and
the same listed constraints, while
\(
 T(p_+)-T(p_-)=2\epsilon T(v)\ne0
\).
\end{proof}

Thus a finite set of low radial moments or low Grassmann coefficients
cannot determine an independent all-order Wilson or determinant
observable.  Closure requires either the full density/source polynomial,
as retained in the present atlas, or a proven tail estimate such as
\eqref{eq:v5v60-YM-tail}.

\section{Finite-cutoff atlas theorem and continuum ledger}

\begin{theorem}[Exact finite-cutoff electroweak shell atlas]
\label{thm:v5v60-ew-atlas}
Assume at a fixed finite cutoff:
\begin{enumerate}
 \item the bosonic density admits the conditional disintegrations of the
 radial and massive-vector shell certificates;
 \item every massive-safe cell satisfies the radial gate
 \eqref{eq:v5v60-safe-cell} and the previously stated block convexity and
 Schur-complement certificates;
 \item every integrated fermion block satisfies
 \eqref{eq:v5v60-fermion-quality};
 \item shell orientations are induced by one finite-cutoff determinant
 line and the residual symmetry hypotheses of
 Proposition~\ref{prop:v5v60-equivariance} hold;
 \item whenever a hypothesis fails, the failed block is moved to the
 retained sector.
\end{enumerate}
Then the massive and crossover cells form an exact finite projective atlas
for the cutoff electroweak superdensity.  Its transition data are the
common-refinement maps of Theorem~\ref{thm:v5v60-refinement-square}; they
satisfy the cocycle law \eqref{eq:v5v60-projective-cocycle}, do not double
count modes, retain the photon on every cell, and have exponentially local
fermionic Schur corrections on every massive-safe cell with the explicit
bounds \eqref{eq:v5v60-safe-inverse}--\eqref{eq:v5v60-safe-Schur}.
\end{theorem}

\begin{proof}
Finite dimensionality and the rule in item 5 produce an admissible cell
after finitely many enlargements of the retained sector.  The radial and
massive-vector certificates give the bosonic conditional kernels.  V58
gives exact Grassmann shell integration, while
Lemma~\ref{lem:v5v60-safe-locality} supplies locality in safe cells.
Theorem~\ref{thm:v5v60-refinement-square} gives exact overlap maps and
their cocycle.  Proposition~\ref{prop:v5v60-no-double-counting} proves the
mode ledger, and Proposition~\ref{prop:v5v60-equivariance} gives descent.
The photon is retained by Definition~\ref{def:v5v60-cell}.
\end{proof}

\begin{firewall}[What the atlas theorem does not prove]
\label{fw:v5v60-continuum-firewall}
Theorem~\ref{thm:v5v60-ew-atlas} is algebraically exact at one finite
cutoff.  It does not prove cutoff-uniform cell number, cutoff-uniform
transition norms, a global chiral determinant-line trivialisation, a lower
bound on the determinant phase, reflection positivity, removal of gauge
fixing, control of the massless photon, or coupling to a dynamical Einstein
metric.  Consequently it is not a construction of the Standard
Model--Einstein continuum.
\end{firewall}

For the next stage, the unresolved conditions are recorded in dependency
order:
\begin{enumerate}
 \item a cutoff-uniform determinant-line and phase chart compatible with
 the anomaly cancellation identities;
 \item cutoff-uniform locality and norm bounds for all common-refinement
 maps, including singular-band crossings;
 \item a massless photon base theory with scale-uniform charged insertions;
 \item reflection-positive reconstruction of the coupled boson--fermion
 atlas;
 \item a conserved, renormalised stress tensor whose metric response can
 enter the Einstein source equation;
 \item joint removal of lattice, volume, gauge, and gravitational
 regulators.
\end{enumerate}

\subsection{V60 status}

\begin{theorem}[V60 conclusion]
\label{thm:v5v60-status}
The radial massive/crossover charts, massive \(W/Z\) blocks, retained
photon, and singular-value fermion charts admit one exact finite-cutoff
projective shell atlas under the stated certificates.  Common refinement,
rather than an inverse to shell integration, supplies the rigorous overlap
law.  This closes the finite-dimensional bookkeeping and locality problem
for certified cells, while exposing determinant-line uniformity and the
massless photon sector as the next upstream bottlenecks.
\end{theorem}

\begin{proof}
Combine Theorem~\ref{thm:v5v60-ew-atlas} with the continuum firewall and
the obstruction Lemma~\ref{lem:v5v60-finite-moment-obstruction}.
\end{proof}

V61 will go upstream to the first bottleneck: it will formulate the
finite-cutoff chiral determinant line over this atlas, prove its exact
transition cocycle, and isolate the additional estimates needed to make
that cocycle uniform and reflection compatible.

\clearpage
\section{V61: determinant-line inheritance and the massless photon base}
\label{sec:v5v61}

V60 named the determinant line and the photon as the next two
dependencies.  A check against the frozen manuscripts changes that
ordering.  The first-series V142 and V145 notes already construct the
exact determinant-functor shell maps, prove their fusion cocycle, and
obtain a unit parallel phase section for the complete three-generation
Standard Model on the faithful gauge-group branch, under their stated
smooth-family and boundary-cobordism hypotheses.  Fourth-series V74 then
fits a soft singular packet into that same line without assigning it an
independent anomaly phase.  Those results remain inherited and must not be
reproved as if they were absent.

The one-generation computation in V57 is a local diagnostic.  Its
firewall about flat lines does not revoke the stronger inherited
three-generation faithful-branch theorem.  What is new after V58--V60 is
the projective radial/Grassmann atlas.  We first prove that its Schur maps
inherit the old determinant line.  We then attack the genuinely open
massless photon base, including its harmonic zero modes and the failure of
a canonical photon line at the zero-Higgs stratum.

\subsection{Transfer of the inherited determinant line}

At one cutoff, write a chiral operator in retained/shell blocks as
\begin{equation}
 D=
 \begin{pmatrix}
 A&B\\ C&E
 \end{pmatrix}:
 R^+\oplus S^+\longrightarrow R^-\oplus S^-,
 \qquad E:S^+\longrightarrow S^-.
 \label{eq:v5v61-chiral-block}
\end{equation}
Assume \(E\) is invertible and set
\begin{equation}
 D_{\mathrm{eff}}:=A-BE^{-1}C.
 \label{eq:v5v61-effective-Dirac}
\end{equation}

\begin{proposition}[Inherited-line Schur isomorphism]
\label{prop:v5v61-line-Schur}
There is a canonical determinant-line isomorphism
\begin{equation}
 \lambda(D)\cong\lambda(D_{\mathrm{eff}})
 \label{eq:v5v61-line-isomorphism}
\end{equation}
after the acyclic shell complex \(E:S^+\to S^-\) is trivialised by its
canonical torsion.  In local frames the scalar representative of this
isomorphism is the exact factor
\begin{equation}
 \det D=\det E\,\det D_{\mathrm{eff}}
 \label{eq:v5v61-block-determinant}
\end{equation}
whenever both sides are ordinary square determinants.  If
\(D_{\mathrm{eff}}\) has a kernel, that zero remains a zero of the
retained determinant-line section.
\end{proposition}

\begin{proof}
Multiply \(D\) on the left and right by
\begin{equation}
 L=
 \begin{pmatrix}1&-BE^{-1}\\0&1\end{pmatrix},
 \qquad
 R=
 \begin{pmatrix}1&0\\-E^{-1}C&1\end{pmatrix}.
 \label{eq:v5v61-eliminators}
\end{equation}
Both are block triangular with determinant one, and
\begin{equation}
 LDR=
 \begin{pmatrix}D_{\mathrm{eff}}&0\\0&E\end{pmatrix}.
 \label{eq:v5v61-diagonalisation}
\end{equation}
The determinant functor is invariant under the two isomorphisms and is
multiplicative under direct sum.  The invertible shell complex is acyclic,
so its canonical torsion identifies its determinant line with the scalar
line.  This gives \eqref{eq:v5v61-line-isomorphism} and, in frames,
\eqref{eq:v5v61-block-determinant}.  No inverse of
\(D_{\mathrm{eff}}\) was used, so its kernel is retained.
\end{proof}

\begin{theorem}[Compatibility with the V60 refinement atlas]
\label{thm:v5v61-line-refinement}
Suppose every V60 fermion shell block is a block of the complete regulated
three-generation chiral operator and every Berezin orientation is induced
from the inherited faithful-branch determinant line.  Then:
\begin{enumerate}
 \item every V60 shell map carries the isomorphism
 \eqref{eq:v5v61-line-isomorphism};
 \item on a common refinement, eliminating shells in either admissible
 order gives the same determinant-line map, including its Koszul sign;
 \item re-fibring a singular band into the retained sector introduces no
 new anomaly phase.
\end{enumerate}
\end{theorem}

\begin{proof}
The first assertion is Proposition~\ref{prop:v5v61-line-Schur} applied to
the shell block.  For the second, V60 Berezin--Fubini says the two iterated
integrals are the same exterior-coefficient extraction.  Equivalently,
successive block eliminations and one combined elimination have the same
unit-determinant triangular factorisation.  The determinant functor's
ordered exterior products supply the same Koszul sign on both routes.
For the third, refine until the singular band is retained.  The remaining
shell is invertible, so the preceding argument applies, while the
crossing line stays inside the original determinant line.  No second
line or phase convention is introduced.
\end{proof}

\begin{assessment}[Disposition of the V60 phase item]
\label{ass:v5v61-phase-disposition}
Under the inherited smooth-family, faithful-quotient, and boundary
hypotheses, the determinant-line phase is not a new V61 obstruction.
What remains analytic is convergence of the renormalised modulus and
compatibility with reflection positivity.  If those inherited hypotheses
fail on a new stratum, the failure must be exhibited there; it cannot be
inferred merely from the one-generation V57 diagnostic.
\end{assessment}

\subsection{Finite cochain model of the photon quotient}

Let
\begin{equation}
 C^0\xrightarrow{d_0}C^1\xrightarrow{d_1}C^2,
 \qquad d_1d_0=0,
 \label{eq:v5v61-cochain-complex}
\end{equation}
be the real finite-dimensional cellular cochain complex of a connected
finite lattice, with fixed inner products and adjoints.  Define
\begin{equation}
 \mathcal H^1:=\ker d_1\cap\ker d_0^*,
 \qquad
 \mathcal T^1:=\operatorname{im}d_1^*.
 \label{eq:v5v61-harmonic-transverse}
\end{equation}

\begin{theorem}[Finite Hodge splitting]
\label{thm:v5v61-Hodge}
There is an orthogonal decomposition
\begin{equation}
 C^1
 =
 \operatorname{im}d_0
 \oplus\mathcal H^1
 \oplus\mathcal T^1.
 \label{eq:v5v61-Hodge}
\end{equation}
On \(\mathcal T^1\), the Maxwell operator
\begin{equation}
 M:=d_1^*d_1
 \label{eq:v5v61-Maxwell-operator}
\end{equation}
is strictly positive and hence invertible.
\end{theorem}

\begin{proof}
Finite-dimensional orthogonal decomposition gives
\[
 C^1=\operatorname{im}d_0\oplus\ker d_0^*.
\]
Inside \(\ker d_0^*\), the orthogonal complement of
\(\operatorname{im}d_1^*\) is
\(\ker d_1\cap\ker d_0^*=\mathcal H^1\), proving
\eqref{eq:v5v61-Hodge}.  If \(a=d_1^*b\in\mathcal T^1\) and
\(d_1a=0\), then
\[
 \|a\|^2=\langle d_1^*b,a\rangle
 =\langle b,d_1a\rangle=0.
\]
Thus \(d_1\) is injective on \(\mathcal T^1\), and
\(\langle a,Ma\rangle=\|d_1a\|^2>0\) for nonzero \(a\) there.
\end{proof}

Gauge transformations act by \(A\mapsto A+d_0\lambda\).  The Coulomb
slice is
\begin{equation}
 \ker d_0^*=\mathcal H^1\oplus\mathcal T^1.
 \label{eq:v5v61-Coulomb-slice}
\end{equation}
The Maxwell action with coupling \(e>0\) is
\begin{equation}
 S_{\mathrm{M}}(A)=\frac1{2e^2}\|d_1A\|^2.
 \label{eq:v5v61-Maxwell-action}
\end{equation}
It is identically zero on \(\mathcal H^1\).  Therefore harmonic modes
cannot be included in a normalised noncompact Gaussian.

\begin{definition}[Finite photon base]
\label{def:v5v61-photon-base}
The finite photon base is the pair \((h,\gamma_e)\), where
\(h\in\mathcal H^1\) is retained and \(\gamma_e\) is the centred Gaussian
probability on \(\mathcal T^1\) with covariance
\begin{equation}
 C_e=e^2M^{-1}.
 \label{eq:v5v61-photon-covariance}
\end{equation}
If the gauge field is compact, \(h\) is instead a point of the harmonic
torus and carries its normalised Haar measure.
\end{definition}

\begin{proposition}[Exact photon characteristic functional]
\label{prop:v5v61-photon-characteristic}
Let \(j\in C^1\) satisfy \(d_0^*j=0\), and write
\(j=j_H+j_T\) according to
\eqref{eq:v5v61-Coulomb-slice}.  At fixed harmonic coordinate \(h\),
\begin{equation}
 \int_{\mathcal T^1}
 e^{i\langle j,h+a\rangle}\,d\gamma_e(a)
 =
 \exp\left(
 i\langle j_H,h\rangle
 -\frac{e^2}{2}\langle j_T,M^{-1}j_T\rangle
 \right).
 \label{eq:v5v61-characteristic}
\end{equation}
For a compact harmonic torus, subsequent Haar integration is zero for
every nontrivial harmonic character and is one for the trivial character.
\end{proposition}

\begin{proof}
The exact component of \(j\) vanishes because
\(\langle j,d_0\lambda\rangle=\langle d_0^*j,\lambda\rangle=0\).
The harmonic and transverse subspaces are orthogonal.  The standard
finite-dimensional Gaussian characteristic formula on
\(\mathcal T^1\) gives the second exponential in
\eqref{eq:v5v61-characteristic}.  Haar orthogonality of characters proves
the last statement.
\end{proof}

\begin{firewall}[Deleting the harmonic kernel changes the theory]
\label{fw:v5v61-harmonic-kernel}
Replacing \(M^{-1}\) by an inverse on all of \(\ker d_0^*\), or silently
setting \(h=0\), discards physical toron/flux sectors.  The inverse exists
only on \(\mathcal T^1\).  The harmonic coordinate must be retained,
compactly integrated, or fixed by an explicitly declared boundary
sector.
\end{firewall}

\subsection{Gauss compatibility of charged insertions}

Let a matter insertion \(\mathcal M_q\) have infinitesimal gauge
transformation
\begin{equation}
 \mathcal M_q\longmapsto
 e^{i\langle q,\lambda\rangle}\mathcal M_q,
 \qquad q\in C^0.
 \label{eq:v5v61-matter-charge}
\end{equation}
Dress it by a photon current \(j\in C^1\).

\begin{theorem}[Exact finite Gauss condition]
\label{thm:v5v61-Gauss-condition}
The dressed insertion
\begin{equation}
 \mathcal O_{q,j}
 :=
 \mathcal M_q e^{i\langle j,A\rangle}
 \label{eq:v5v61-dressed-insertion}
\end{equation}
is invariant under every finite-cutoff gauge transformation connected to
the identity if and only if
\begin{equation}
 q+d_0^*j=0.
 \label{eq:v5v61-Gauss-condition}
\end{equation}
On a connected closed lattice this implies charge neutrality,
\begin{equation}
 \langle q,\mathbf1\rangle=0.
 \label{eq:v5v61-neutrality}
\end{equation}
\end{theorem}

\begin{proof}
Under \(A\mapsto A+d_0\lambda\), the photon factor acquires
\[
 \exp\{i\langle j,d_0\lambda\rangle\}
 =
 \exp\{i\langle d_0^*j,\lambda\rangle\}.
\]
Multiplication by \eqref{eq:v5v61-matter-charge} is invariant for all
\(\lambda\) exactly when \eqref{eq:v5v61-Gauss-condition} holds.  Pair
that equation with the constant cochain.  Since
\(d_0\mathbf1=0\),
\[
 \langle d_0^*j,\mathbf1\rangle
 =\langle j,d_0\mathbf1\rangle=0,
\]
which proves \eqref{eq:v5v61-neutrality}.
\end{proof}

\begin{corollary}[No isolated charged observable on the closed quotient]
\label{cor:v5v61-no-single-charge}
A single nonzero local charge cannot define a gauge-invariant observable
of the closed finite-volume photon quotient.  It must be paired with
opposite charge, attached to a boundary flux sector, or represented in a
different charged superselection construction.
\end{corollary}

\begin{proof}
A single charge violates \eqref{eq:v5v61-neutrality}.
\end{proof}

\subsection{A thermodynamic infrared bound for local neutral currents}

The exact zero-mode quotient removes only the algebraic divergence.  We
now check whether the remaining massless covariance is uniformly finite
as the periodic volume grows.  Consider the unit-spacing \(d\)-dimensional
periodic lattice of side \(L\).  Its nonzero Fourier momenta lie in
\((2\pi/L)\mathbb Z^d\cap[-\pi,\pi)^d\), and
\begin{equation}
 \lambda(k)=4\sum_{\nu=1}^d\sin^2(k_\nu/2).
 \label{eq:v5v61-lattice-symbol}
\end{equation}
For a one-cochain with zero harmonic component, define
\begin{equation}
 Q_L(j):=
 L^{-d}\sum_{k\ne0}
 \frac{|\widehat{P_Tj}(k)|^2}{\lambda(k)}.
 \label{eq:v5v61-Q}
\end{equation}

\begin{theorem}[Volume-uniform local-current bound]
\label{thm:v5v61-IR-bound}
Suppose the support of \(j\) has diameter less than \(L/2\), choose a lift
of that support to \(\mathbb Z^d\), and assume each component has zero
mean,
\begin{equation}
 \sum_xj_\mu(x)=0.
 \label{eq:v5v61-zero-mean-current}
\end{equation}
For a base point \(x_0\) in the lifted support, set
\begin{equation}
 M_1(j):=
 \sum_{x,\mu}|j_\mu(x)|\,|x-x_0|.
 \label{eq:v5v61-first-moment}
\end{equation}
Then
\begin{equation}
 Q_L(j)
 \le C_d\bigl(M_1(j)^2+\|j\|_1^2\bigr),
 \label{eq:v5v61-IR-bound}
\end{equation}
where \(C_d\) is independent of \(L\).
\end{theorem}

\begin{proof}
Condition \eqref{eq:v5v61-zero-mean-current} gives
\begin{equation}
 \widehat j_\mu(k)
 =
 \sum_xj_\mu(x)
 \left(e^{-ik\cdot(x-x_0)}-1\right),
 \label{eq:v5v61-Fourier-cancellation}
\end{equation}
so
\begin{equation}
 |\widehat j(k)|\le |k|M_1(j).
 \label{eq:v5v61-Fourier-linear}
\end{equation}
Orthogonal transverse projection cannot increase the Fourier norm.  On
\(|k|\le1\), the elementary sine bound gives
\(\lambda(k)\ge c_d|k|^2\), hence every summand ratio is at most
\(c_d^{-1}M_1(j)^2\).  On \(|k|>1\), compactness of the Brillouin zone
away from zero gives \(\lambda(k)\ge c'_d>0\), while
\(|\widehat j(k)|\le\|j\|_1\).  The normalised number
\(L^{-d}\#\{k\}\) is one.  Summing the two regions proves
\eqref{eq:v5v61-IR-bound}.
\end{proof}

\begin{corollary}[No thermodynamic photon divergence for a fixed local
neutral current]
\label{cor:v5v61-thermodynamic}
For a family obtained by embedding one fixed, bounded-support,
zero-harmonic current into larger periodic boxes, the free-photon
characteristic functional \eqref{eq:v5v61-characteristic} has a
volume-uniform exponent.
\end{corollary}

\begin{proof}
The quantities on the right of \eqref{eq:v5v61-IR-bound} are fixed under
the embeddings.
\end{proof}

\begin{firewall}[This is not an ultraviolet or charged-particle theorem]
\label{fw:v5v61-not-UV}
Theorem~\ref{thm:v5v61-IR-bound} keeps the lattice spacing fixed.  Thin
Wilson lines and pointlike dressings can have ultraviolet self-energy
divergences when the spacing is removed.  A dressing whose size grows with
particle separation can also make \(M_1(j)\) grow.  Neither effect is
controlled by the volume-uniform estimate.
\end{firewall}

\subsection{The photon label at the Higgs crossover}

On the principal Higgs stratum, V56 defines the neutral combinations
\begin{equation}
 Z_\mu=\frac{gW_\mu^3-g'B_\mu}{\sqrt{g^2+g'^2}},
 \qquad
 A_\mu=\frac{g'W_\mu^3+gB_\mu}{\sqrt{g^2+g'^2}}.
 \label{eq:v5v61-AZ}
\end{equation}

\begin{proposition}[No canonical photon projector at zero modulus]
\label{prop:v5v61-no-photon-projector}
For \(\rho>0\), the neutral mass matrix has a one-dimensional kernel
spanned by the photon in \eqref{eq:v5v61-AZ}, separated from the \(Z\)
eigenvalue by a gap proportional to
\(\rho^2(g^2+g'^2)\).  At \(\rho=0\), the complete electroweak mass matrix
vanishes and its kernel dimension jumps to four.  Hence the rank-one
photon projector is not characterised by the mass spectrum at the
zero-Higgs stratum, and no gauge-invariant unitary-gauge photon line
extends canonically through that stratum.
\end{proposition}

\begin{proof}
The V56 mass quadratic form is
\[
 \frac{\rho^2}{4}
 \left[g^2\bigl((W^1)^2+(W^2)^2\bigr)
 +(gW^3-g'B)^2\right].
\]
For \(\rho>0\), its neutral block has eigenvalues zero and
\(\rho^2(g^2+g'^2)/4\), up to the declared real-vector convention, with
the eigenvectors in \eqref{eq:v5v61-AZ}.  At \(\rho=0\), all four gauge
directions have zero mass.  Moreover the unitary-gauge Higgs direction
used to identify \(T^3+Y\) is undefined there.  A fixed matrix
continuation of \eqref{eq:v5v61-AZ} is therefore a chart choice, not a
spectrally or gauge-invariantly distinguished line.
\end{proof}

\begin{corollary}[Correction to the V60 retained-sector wording]
\label{cor:v5v61-V60-correction}
On principal massive cells, the photon is retained as the residual
\(U(1)_{\mathrm{em}}\) field.  On crossover cells meeting \(\rho=0\), the
retained sector must contain the full electroweak band whose mass gap
collapses.  The photon label on that larger band is only an overlap
coordinate.
\end{corollary}

\begin{proof}
Apply Proposition~\ref{prop:v5v61-no-photon-projector} and the V60 rule
that every band failing its gap certificate is re-fibred into the retained
sector.
\end{proof}

\subsection{Photon-base certificate and V61 status}

\begin{definition}[Photon-base certificate]
\label{def:v5v61-PBC}
A photon-base certificate consists of:
\begin{enumerate}
 \item the exact Hodge quotient
 \eqref{eq:v5v61-Hodge}, with harmonic modes retained or compactly
 integrated;
 \item the transverse covariance \eqref{eq:v5v61-photon-covariance};
 \item the Gauss compatibility condition
 \eqref{eq:v5v61-Gauss-condition} for every charged insertion;
 \item volume-uniform bounds of the form
 \eqref{eq:v5v61-IR-bound} for the chosen neutral or dressed observable
 class;
 \item ultraviolet renormalisation of thin dressings and charged
 insertions;
 \item reflection positivity on the gauge-invariant photon/matter
 algebra;
 \item a stratified handoff from \(U(1)_{\mathrm{em}}\) to the full
 electroweak retained band at Higgs zeros;
 \item a cutoff trajectory for the abelian coupling consistent with the
 inherited hypercharge ultraviolet obstruction.
\end{enumerate}
\end{definition}

\begin{theorem}[V61 conclusion]
\label{thm:v5v61-status}
The new V60 Grassmann/refinement atlas inherits the already constructed
three-generation faithful-branch determinant line and does not create an
independent phase.  At finite cutoff, the free electromagnetic quotient
has an exact Hodge-factorised base: gauge directions are removed,
harmonic modes are retained, and the coexact modes carry a nondegenerate
Gaussian with characteristic functional
\eqref{eq:v5v61-characteristic}.  Gauge-invariant charged insertions obey
the exact Gauss condition \eqref{eq:v5v61-Gauss-condition}, and fixed local
zero-harmonic currents have the volume-uniform infrared bound
\eqref{eq:v5v61-IR-bound}.  At the zero-Higgs crossover there is no
canonical photon line, so the whole gapless electroweak band must be
retained.  This proves the finite free and thermodynamic-infrared parts of
a photon-base certificate, not its ultraviolet, interacting,
reflection-positive, or hypercharge-continuum parts.
\end{theorem}

\begin{proof}
Use Theorem~\ref{thm:v5v61-line-refinement},
Theorem~\ref{thm:v5v61-Hodge},
Proposition~\ref{prop:v5v61-photon-characteristic},
Theorem~\ref{thm:v5v61-Gauss-condition},
Theorem~\ref{thm:v5v61-IR-bound}, and
Corollary~\ref{cor:v5v61-V60-correction}.
\end{proof}

V62 will attack the next missing photon-base item: reflection positivity
of the finite free Maxwell quotient on a gauge-invariant half-space
algebra, with harmonic sectors kept explicit.  If a chosen charged
dressing crosses the reflection plane or destroys the positive
factorisation, it will be rejected or moved into a boundary-sector
coordinate rather than justified by gauge fixing.

\clearpage
\section{V62: reflection positivity of the free photon quotient}
\label{sec:v5v62}

V61 constructed the finite Hodge quotient and isolated the harmonic
kernel.  Reflection positivity cannot be read from positivity of the
four-dimensional inverse Maxwell operator alone: the gauge projection can
be nonlocal in Euclidean time, and proper-time shell differences need not
be reflection positive separately.  We therefore use the physical
Hamiltonian quotient.  This produces two transverse oscillators for every
nonzero spatial momentum and a separate compact quantum rotor for each
harmonic spatial direction.

The scalar Gaussian argument of V11 and the transverse-traceless
oscillator argument of V13 are inherited.  The new work is to identify
the electromagnetic physical variables, include their harmonic sector,
and prove that gauge-invariant photon observables and admissible
positive-time dressings lie in the resulting OS algebra.

\subsection{Spatial Coulomb quotient}

Let \(\Sigma\) be a finite connected periodic spatial lattice.  Its real
cochain complex begins
\begin{equation}
 C^0(\Sigma)\xrightarrow{d_\Sigma}C^1(\Sigma)
 \xrightarrow{d_{\Sigma,1}}C^2(\Sigma).
 \label{eq:v5v62-spatial-complex}
\end{equation}
Put
\begin{equation}
 \mathcal H_\Sigma^1
 :=
 \ker d_{\Sigma,1}\cap\ker d_\Sigma^*,
 \qquad
 \mathcal T_\Sigma^1
 :=
 \operatorname{im}d_{\Sigma,1}^*.
 \label{eq:v5v62-spatial-spaces}
\end{equation}
The V61 Hodge proof, applied on each time slice, gives
\begin{equation}
 C^1(\Sigma)
 =
 \operatorname{im}d_\Sigma
 \oplus\mathcal H_\Sigma^1
 \oplus\mathcal T_\Sigma^1.
 \label{eq:v5v62-spatial-Hodge}
\end{equation}

Work first in continuous Euclidean time with this finite spatial cutoff.
Write
\begin{equation}
 A(t)=d_\Sigma\phi(t)+h(t)+a_T(t)
 \label{eq:v5v62-time-slice-split}
\end{equation}
according to \eqref{eq:v5v62-spatial-Hodge}.  The exact term is gauge,
\(h(t)\in\mathcal H_\Sigma^1\), and
\(a_T(t)\in\mathcal T_\Sigma^1\).

\begin{proposition}[Physical free-photon quadratic form]
\label{prop:v5v62-physical-form}
After imposing Gauss law and quotienting the exact spatial gauge
directions, the free Maxwell quadratic form is
\begin{equation}
 S_{\mathrm{phys}}
 =
 \frac1{2e^2}\int_{\mathbb R}
 \left(
 \|\dot a_T(t)\|^2
 \langle a_T(t),M_\Sigma a_T(t)\rangle
 \|\dot h(t)\|^2
 \right)\,dt,
 \label{eq:v5v62-physical-action}
\end{equation}
where
\begin{equation}
 M_\Sigma
 :=
 d_{\Sigma,1}^*d_{\Sigma,1}
 \quad\hbox{on }\mathcal T_\Sigma^1
 \label{eq:v5v62-spatial-Maxwell}
\end{equation}
is strictly positive.
\end{proposition}

\begin{proof}
Choose temporal gauge locally in time, or equivalently impose the
Hamiltonian Gauss constraint before quotienting.  Then the electric term
is \(\|\dot A\|^2/(2e^2)\), while the magnetic term is
\(\|d_{\Sigma,1}A\|^2/(2e^2)\).  Orthogonality in
\eqref{eq:v5v62-spatial-Hodge} removes the exact coordinate and separates
\(h\) from \(a_T\).  The harmonic coordinate has no magnetic term.
Strict positivity of \(M_\Sigma\) on \(\mathcal T_\Sigma^1\) is
Theorem~\ref{thm:v5v61-Hodge} on the spatial complex.
\end{proof}

\begin{firewall}[Temporal gauge is not the proof]
\label{fw:v5v62-temporal-gauge}
Proposition~\ref{prop:v5v62-physical-form} is a statement about the
Gauss-law quotient.  A formal choice \(A_0=0\) by itself leaves residual
time-independent gauge transformations and does not dispose of harmonic
coordinates.  The Hodge quotient and the separate \(h\) sector are
essential.
\end{firewall}

\subsection{Nonzero transverse modes}

Choose an orthonormal eigenbasis
\(\{u_\alpha\}_{\alpha=1}^{N_T}\) of \(M_\Sigma\),
\begin{equation}
 M_\Sigma u_\alpha=\omega_\alpha^2u_\alpha,
 \qquad \omega_\alpha>0,
 \label{eq:v5v62-transverse-spectrum}
\end{equation}
and expand \(a_T(t)=\sum_\alpha q_\alpha(t)u_\alpha\).  The vacuum
covariance is
\begin{equation}
 \mathbb E[q_\alpha(t)q_\beta(s)]
 =
 \delta_{\alpha\beta}
 \frac{e^2}{2\omega_\alpha}
 e^{-\omega_\alpha|t-s|}.
 \label{eq:v5v62-mode-covariance}
\end{equation}
Let \(\theta t=-t\), and let the transverse spatial potential be even
under time reflection.

\begin{theorem}[Transverse photon reflection positivity]
\label{thm:v5v62-transverse-OS}
The centred Gaussian functional with covariance
\eqref{eq:v5v62-mode-covariance} is reflection positive on the cylinder
algebra generated by \(a_T(f)\) with time support in \(t>0\).
\end{theorem}

\begin{proof}
For positive-time test coefficients \(f_\alpha(t)\), the reflected
covariance form is
\begin{align}
 \langle\theta f,C_Tf\rangle
 &=
 \sum_{\alpha=1}^{N_T}
 \frac{e^2}{2\omega_\alpha}
 \int_0^\infty\!\!\int_0^\infty
 \overline{f_\alpha(t)}
 e^{-\omega_\alpha(t+s)}
 f_\alpha(s)\,dt\,ds
 \notag\\
 &=
 \sum_{\alpha=1}^{N_T}
 \frac{e^2}{2\omega_\alpha}
 \left|
 \int_0^\infty e^{-\omega_\alpha t}f_\alpha(t)\,dt
 \right|^2
 \ge0.
 \label{eq:v5v62-reflected-form}
\end{align}
The Gaussian exponential-vector argument used in V11 then proves
positivity for all finite cylinder polynomials.
\end{proof}

\begin{corollary}[Thermodynamic masslessness is compatible with OS
positivity]
\label{cor:v5v62-massless-OS}
The lower transverse frequency may tend to zero as the spatial volume
grows without changing the sign in
\eqref{eq:v5v62-reflected-form}.  Reflection positivity therefore
survives every pointwise limit of transverse Schwinger matrices for which
the entries converge.
\end{corollary}

\begin{proof}
Every term in \eqref{eq:v5v62-reflected-form} is nonnegative for every
\(\omega_\alpha>0\).  The inherited closure of positive semidefinite
finite matrices under entrywise limits gives the second assertion.
\end{proof}

\subsection{The harmonic rotor}

For a compact \(U(1)\) gauge field,
\(\mathcal H_\Sigma^1\) exponentiates to a torus
\(\mathbb T_H\).  The harmonic part of
\eqref{eq:v5v62-physical-action} is a quantum rotor.  On
\(\mathscr H_H=L^2(\mathbb T_H,dh)\), let
\begin{equation}
 H_H=-\frac{e^2}{2}\Delta_{\mathbb T_H}\ge0,
 \qquad
 \Omega_H(h)=1.
 \label{eq:v5v62-rotor-Hamiltonian}
\end{equation}
The precise volume factors can be absorbed into the flat metric defining
\(\Delta_{\mathbb T_H}\); their positivity is all that is used below.

\begin{lemma}[Positive semigroup reflection lemma]
\label{lem:v5v62-semigroup-OS}
Let \(H\ge0\) be self-adjoint on a Hilbert space, let
\(\Omega\in\ker H\), and let \(A_1,\ldots,A_n\) be bounded multiplication
operators closed under adjoint.  The Euclidean time-ordered functions
formed from \(e^{-tH}\), \(\Omega\), and these operators are reflection
positive.
\end{lemma}

\begin{proof}
For a monomial supported at ordered positive times
\(0<t_1<\cdots<t_n\), associate the vector
\begin{equation}
 \Psi=
 e^{-t_1H}A_1e^{-(t_2-t_1)H}\cdots
 A_ne^{-(T-t_n)H}\Omega
 \label{eq:v5v62-transfer-vector}
\end{equation}
with any common \(T\ge t_n\).  Reflection reverses the order and replaces
each \(A_i\) by \(A_i^*\).  The reflected product of two monomials is
therefore the Hilbert inner product of their associated vectors, after a
common harmless terminal semigroup factor.  A finite linear combination
has OS form equal to the squared norm of the corresponding linear
combination of vectors.
\end{proof}

\begin{theorem}[Harmonic-sector reflection positivity]
\label{thm:v5v62-harmonic-OS}
The compact harmonic photon sector with Hamiltonian
\eqref{eq:v5v62-rotor-Hamiltonian} is reflection positive on bounded
cylinder functions of \(h(t)\).  A positive mixture of electric-flux
superselection sectors is also reflection positive.
\end{theorem}

\begin{proof}
Apply Lemma~\ref{lem:v5v62-semigroup-OS} to \(H_H\), the constant ground
state, and trigonometric multiplication operators on \(\mathbb T_H\).
They generate a dense cylinder algebra.  For a positive mixture, each OS
matrix is a positive weighted sum of positive matrices.
\end{proof}

\begin{assessment}[Noncompact harmonic coordinates]
\label{ass:v5v62-noncompact-harmonic}
For a noncompact linear photon field, the zero-frequency harmonic
coordinate has no normalisable translation-invariant vacuum.  One must
fix a harmonic sector, compactify it, or use a non-vacuum representation.
The transverse proof does not manufacture a probability measure for that
coordinate.
\end{assessment}

\subsection{Gauge-invariant field strengths}

Let \(\mathcal A_{T,+}\) be the positive-time algebra generated by
transverse coordinates and let \(\mathcal A_{F,+}\) be the subalgebra
generated by smeared electric and magnetic field strengths whose test
supports have positive distance from the reflection plane.  Time
reflection acts with the usual parity: magnetic components are even and
components with one time index are odd.

\begin{proposition}[OS descent to the field-strength algebra]
\label{prop:v5v62-field-strength-OS}
The free-photon Schwinger functional is reflection positive on
\(\mathcal A_{F,+}\).
\end{proposition}

\begin{proof}
After the Gauss quotient, the nonharmonic part of every smeared field
strength in \(\mathcal A_{F,+}\) is a real linear combination of \(a_T\)
and its time derivative, with the induced reflection parity.  The harmonic
electric part is a cylinder observable of the rotor sector.  Integration by parts
moves the time derivative to a test function without crossing \(t=0\)
because the test support is separated from the plane.  Hence the nonharmonic part of
\(\mathcal A_{F,+}\) embeds as a reflection-stable subalgebra of the
transverse oscillator algebra.  Restricting a positive quadratic form to
a subalgebra preserves positivity.  Tensoring with the positive harmonic
rotor sector and using
Theorem~\ref{thm:v5v62-harmonic-OS} preserves it again.
\end{proof}

\begin{firewall}[The four-dimensional Hodge projector is not
half-space local]
\label{fw:v5v62-four-dimensional-projector}
The covariant projector
\begin{equation}
 P_{\mathrm{cov}}
 =
 1-d_0(d_0^*d_0)^{-1}d_0^*
 \label{eq:v5v62-covariant-projector}
\end{equation}
uses a four-dimensional inverse Laplacian.  Applied to a current supported
at positive time, it generally has tails at negative time.  Therefore
reflection positivity cannot be obtained by projecting an arbitrary
positive-time covariant potential test function with
\eqref{eq:v5v62-covariant-projector}.  The spatial Gauss quotient used
above preserves time support and avoids this error.
\end{firewall}

\subsection{Dressed insertions}

For a time-slice matter charge \(q(t)\), V61 requires a spatial current
\(j(t)\) with
\begin{equation}
 q(t)+d_\Sigma^*j(t)=0.
 \label{eq:v5v62-time-slice-Gauss}
\end{equation}
Let \(P_{\Sigma,T}\) be the spatial transverse projector.  Unlike
\eqref{eq:v5v62-covariant-projector}, it acts independently at each time.

\begin{theorem}[Sufficient OS support condition for a photon dressing]
\label{thm:v5v62-dressing-OS}
Suppose a finite family of dressed insertions satisfies:
\begin{enumerate}
 \item the exact Gauss condition
 \eqref{eq:v5v62-time-slice-Gauss};
 \item all matter endpoints and all dressing currents are supported in
 \(t\ge\epsilon>0\);
 \item the matter factor, before electromagnetic interaction is turned
 on, has a reflection-positive positive-time algebra;
 \item the harmonic character of each complete insertion is included in
 the compact rotor algebra.
\end{enumerate}
Then the product free matter--photon functional is reflection positive on
the algebra generated by these insertions.
\end{theorem}

\begin{proof}
Because \(P_{\Sigma,T}\) is spatial, it leaves the time support in item 2
unchanged.  The transverse photon factor is therefore an element of
\(\mathcal A_{T,+}\), and Theorem~\ref{thm:v5v62-transverse-OS} applies.
The harmonic character is covered by
Theorem~\ref{thm:v5v62-harmonic-OS}.  The tensor product of the photon OS
Gram matrix and the matter OS Gram matrix is positive by the Schur product
theorem.  Finite linear combinations complete the claim.
\end{proof}

\begin{firewall}[Interaction is not a tensor product]
\label{fw:v5v62-interaction}
Theorem~\ref{thm:v5v62-dressing-OS} concerns a product free functional.
After the covariant derivative couples matter and photons, positivity must
come from a common transfer operator or a reflection-factorised
interaction.  Multiplying separately proved free inequalities does not
establish the interacting Standard Model.
\end{firewall}

\begin{firewall}[A crossing dressing leaves the positive-time algebra]
\label{fw:v5v62-crossing-dressing}
A Wilson path or Coulomb dressing that crosses the reflection plane is not
an element of the declared positive-time algebra.  Its gauge invariance
does not imply OS positivity.  It must be split with an explicit boundary
state, replaced by a same-half-space dressing, or treated in a larger
transfer-matrix construction.
\end{firewall}

\subsection{Shell compatibility}

\begin{proposition}[Complete covariance, not isolated shells]
\label{prop:v5v62-complete-covariance}
Let \(C_T\) be the covariance
\eqref{eq:v5v62-mode-covariance} and write it algebraically as a sum of
Wilsonian covariances \(C_T=\sum_n C_n\).  Positivity of each \(C_n\) as an
ordinary covariance does not imply that each \(C_n\) is reflection
positive.  The V60 projective atlas preserves photon OS positivity only if
the assembled characteristic functional equals that of \(C_T\), or if a
separate reflected transfer factorisation is proved shell by shell.
\end{proposition}

\begin{proof}
Reflection positivity tests the kernel
\(C_n(-t,s)\), not the ordinary quadratic form \(C_n(t,s)\).
A positive difference of spectral cutoffs can fail the former test, as
already recorded in the V11 shell firewall.  Exact V60 refinement and
Fubini recover the complete covariance when all shells are assembled; then
Theorems~\ref{thm:v5v62-transverse-OS} and
\ref{thm:v5v62-harmonic-OS} apply.  Without exact assembly, an independent
reflected factorisation is necessary.
\end{proof}

\subsection{Free photon OS certificate and V62 status}

\begin{definition}[Free photon OS certificate]
\label{def:v5v62-FPOSC}
A free photon OS certificate consists of:
\begin{enumerate}
 \item the spatial Hodge/Gauss quotient
 \eqref{eq:v5v62-spatial-Hodge};
 \item the positive transverse oscillator covariance
 \eqref{eq:v5v62-mode-covariance};
 \item a declared compact, fixed, or superselected harmonic sector;
 \item a positive-time gauge-invariant field-strength algebra;
 \item for charged observables, the support and Gauss conditions of
 Theorem~\ref{thm:v5v62-dressing-OS};
 \item exact shell reassembly or a shellwise transfer factorisation;
 \item an interacting extension proved from one common transfer operator.
\end{enumerate}
\end{definition}

\begin{theorem}[V62 conclusion]
\label{thm:v5v62-status}
At finite spatial cutoff, the physical free photon quotient is reflection
positive.  Its nonzero transverse modes are positive oscillators, and its
compact harmonic modes form positive quantum rotors.  Reflection
positivity descends to the gauge-invariant field-strength algebra and to
free dressed matter insertions satisfying the exact Gauss and
positive-time support conditions.  A covariant four-dimensional Hodge
projection, an isolated Wilsonian covariance band, or a dressing crossing
the reflection plane does not supply this conclusion.  The interacting
electroweak transfer operator and ultraviolet removal remain open.
\end{theorem}

\begin{proof}
Combine Theorems~\ref{thm:v5v62-transverse-OS} and
\ref{thm:v5v62-harmonic-OS},
Proposition~\ref{prop:v5v62-field-strength-OS}, and
Theorem~\ref{thm:v5v62-dressing-OS}, subject to
Firewalls~\ref{fw:v5v62-four-dimensional-projector} and
\ref{fw:v5v62-interaction}.
\end{proof}

V63 will move from the free photon quotient to the interacting finite
lattice.  It will derive a sufficient reflection factorisation for a
compact \(U(1)\) Wilson action coupled to gauge-covariant matter hopping,
identify which chiral-fermion step is still missing from a positive
transfer kernel, and keep the Higgs-zero electroweak crossover as a
separate retained stratum.

\clearpage
\section{V63: interacting compact electroweak bosons and the chiral OS gate}
\label{sec:v5v63}

V62 proved reflection positivity after solving the free photon constraint.
For an interacting compact lattice there is a more robust route: avoid
gauge fixing and expand every cross-plane Boltzmann factor in positive
group characters.  The compact \(U(1)\) Wilson weight has exactly such an
expansion, and the gauge-covariant Higgs hopping factor has a positive
rank-one expansion even when the Higgs modulus fluctuates and vanishes.

This closes the finite-cutoff bosonic electromagnetic/Higgs OS step.  It
does not close the chiral theory.  A globally consistent determinant-line
phase need not factor across a reflection plane.  We identify the exact
reflection-coboundary or Grassmann transfer condition which would close
that remaining step.

\subsection{The reflection-positive cone}

Let a finite reflection-symmetric lattice decompose into positive and
negative halves \(\Lambda_+,\Lambda_-\), with a possible fixed reflection
slice \(\Lambda_0\).  Let \(\Theta\) be the antilinear involution on
observables which reflects the lattice, complex conjugates scalar
coefficients, reverses link orientations, and reverses the order of
Grassmann monomials when fermions are present.  Denote the positive-half
algebra by \(\mathcal A_+\).

\begin{definition}[Reflection-positive cone]
\label{def:v5v63-RP-cone}
The algebraic reflection-positive cone is
\begin{equation}
 \mathcal P_\Theta
 :=
 \left\{
 \sum_{\alpha=1}^N
 c_\alpha(\Theta F_\alpha)F_\alpha:
 N<\infty,\ c_\alpha\ge0,\ F_\alpha\in\mathcal A_+
 \right\}.
 \label{eq:v5v63-RP-cone}
\end{equation}
Its \(L^1\)-closure with respect to the uncoupled product measure is
denoted by \(\overline{\mathcal P_\Theta}^{\,L^1}\).
\end{definition}

\begin{lemma}[Cone multiplication]
\label{lem:v5v63-cone-product}
If \(X,Y\in\mathcal P_\Theta\) are bosonic, then
\(XY\in\mathcal P_\Theta\).  The same holds in the \(L^1\)-closed cone
whenever the products converge in \(L^1\).
\end{lemma}

\begin{proof}
Write
\[
 X=\sum_\alpha c_\alpha(\Theta F_\alpha)F_\alpha,
 \qquad
 Y=\sum_\beta d_\beta(\Theta G_\beta)G_\beta.
\]
Bosonic commutativity and multiplicativity of \(\Theta\) give
\[
 XY
 =
 \sum_{\alpha,\beta}c_\alpha d_\beta
 \Theta(F_\alpha G_\beta)(F_\alpha G_\beta),
\]
which is in the cone.  Approximation and \(L^1\) continuity give the
closure statement.
\end{proof}

\begin{theorem}[Reflection factorisation criterion]
\label{thm:v5v63-factorisation}
Let \(d\mu_0\) be a reflection-invariant product measure which factorises
between the two open halves conditional on \(\Lambda_0\).  Suppose
\begin{equation}
 e^{-S}
 =
 (\Theta e^{-S_+})e^{-S_+}W_0W_\times,
 \label{eq:v5v63-action-factorisation}
\end{equation}
where \(W_0\ge0\) depends only on fixed-slice variables and
\(W_\times\in\overline{\mathcal P_\Theta}^{\,L^1}\).  Then the normalised
functional
\begin{equation}
 \mathbb E_\mu[F]
 :=
 Z^{-1}\int F e^{-S}\,d\mu_0
 \label{eq:v5v63-interacting-functional}
\end{equation}
is reflection positive, provided \(0<Z<\infty\).
\end{theorem}

\begin{proof}
First take
\(W_\times=\sum_\alpha c_\alpha(\Theta G_\alpha)G_\alpha\).
For \(F=\sum_i z_iF_i\) with \(F_i\in\mathcal A_+\), conditional Fubini
factorisation gives
\begin{align}
 \int(\Theta F)F e^{-S}\,d\mu_0
 &=
 \int_{\Lambda_0}W_0
 \sum_\alpha c_\alpha
 \left|
 \int_{\Lambda_+}
 F e^{-S_+}G_\alpha\,d\mu_{0,+|\Lambda_0}
 \right|^2
 d\mu_{0,0}
 \ge0.
 \label{eq:v5v63-factorisation-square}
\end{align}
An \(L^1\) approximation proves the closed-cone case for bounded cylinder
observables; truncation gives the declared integrable domain.
Normalisation by positive \(Z\) preserves the sign.
\end{proof}

\subsection{Compact \(U(1)\) Wilson plaquettes}

For \(z=e^{i\vartheta}\in U(1)\), the Wilson plaquette weight is
\begin{equation}
 w_\beta(z):=\exp\{\beta\operatorname{Re}z\},
 \qquad \beta\ge0.
 \label{eq:v5v63-Wilson-weight}
\end{equation}

\begin{lemma}[Positive Wilson character expansion]
\label{lem:v5v63-Wilson-character}
The exact Fourier expansion is
\begin{equation}
 w_\beta(e^{i\vartheta})
 =
 \sum_{n\in\mathbb Z}I_{|n|}(\beta)e^{in\vartheta},
 \qquad I_{|n|}(\beta)\ge0,
 \label{eq:v5v63-Wilson-character}
\end{equation}
and the series converges absolutely and uniformly.
\end{lemma}

\begin{proof}
The coefficient is
\[
 I_n(\beta)
 =
 \frac1{2\pi}\int_{-\pi}^{\pi}
 e^{\beta\cos\vartheta}e^{-in\vartheta}\,d\vartheta.
\]
For \(n\ge0\), its power series is
\begin{equation}
 I_n(\beta)
 =
 \sum_{k=0}^\infty
 \frac{(\beta/2)^{2k+n}}{k!(k+n)!},
 \label{eq:v5v63-Bessel-series}
\end{equation}
which is nonnegative.  Symmetry gives \(I_{-n}=I_n\).
The sum of the absolute coefficients equals
\(e^\beta\), proving absolute uniform convergence.
\end{proof}

Use link reflection through a plane chosen so that every crossing
plaquette holonomy can be written
\begin{equation}
 U_p=X_p\,\Theta X_p
 \label{eq:v5v63-crossing-plaquette}
\end{equation}
after the standard orientation convention; \(X_p\) is the ordered
positive-half plaquette path, including any fixed-plane half links.

\begin{proposition}[Crossing Wilson weight lies in the cone]
\label{prop:v5v63-Wilson-cone}
For every crossing plaquette,
\begin{equation}
 w_\beta(U_p)
 =
 \sum_{n\in\mathbb Z}
 I_{|n|}(\beta)
 \bigl(\Theta X_p^n\bigr)X_p^n
 \in\overline{\mathcal P_\Theta}^{\,L^1}.
 \label{eq:v5v63-Wilson-cone}
\end{equation}
\end{proposition}

\begin{proof}
Insert \eqref{eq:v5v63-crossing-plaquette} into
\eqref{eq:v5v63-Wilson-character}.  Each coefficient is nonnegative by
Lemma~\ref{lem:v5v63-Wilson-character}; absolute convergence permits the
\(L^1\) limit.
\end{proof}

\subsection{Radial Higgs hopping across the plane}

Write the Higgs field in a local polar representative as
\(\Phi_x=\rho_x u_x\), where \(\rho_x\ge0\) and \(u_x\) is a unit vector
in the appropriate compact representation orbit.  Across a time-like
link \(\ell=(x,\theta x)\), the gauge-covariant hopping weight has the
form
\begin{equation}
 W_{\kappa,\ell}
 =
 \exp\left\{
 \kappa\rho_x\rho_{\theta x}
 \operatorname{Re}
 \langle u_x,U_\ell u_{\theta x}\rangle
 \right\},
 \qquad \kappa\ge0.
 \label{eq:v5v63-Higgs-hopping}
\end{equation}

\begin{lemma}[Positive power expansion of a crossing hop]
\label{lem:v5v63-Higgs-expansion}
With the reflection convention
\[
 Z_\ell:=\rho_x\langle v,U_{\ell,+}u_x\rangle
\]
for a fixed-plane representation basis \(v\), the hopping factor is an
absolutely convergent sum of terms
\begin{equation}
 c_{mn}\,
 \Theta\!\left(Z_\ell^m\overline{Z_\ell}^{\,n}\right)
 \left(Z_\ell^m\overline{Z_\ell}^{\,n}\right),
 \qquad
 c_{mn}=\frac{(\kappa/2)^{m+n}}{m!\,n!}\ge0,
 \label{eq:v5v63-Higgs-rank-one}
\end{equation}
summed also over the finite representation indices introduced by
contracting the half-link.  Hence
\(W_{\kappa,\ell}\in
\overline{\mathcal P_\Theta}^{\,L^1}\) whenever the radial potential
makes the series integrable.
\end{lemma}

\begin{proof}
Write
\[
 e^{\kappa\operatorname{Re}z}
 =
 \exp\{(\kappa/2)z\}
 \exp\{(\kappa/2)\bar z\}
 =
 \sum_{m,n\ge0}
 \frac{(\kappa/2)^{m+n}}{m!\,n!}z^m\bar z^n.
\]
For a reflected link, the representation contraction \(z\) is the finite
sum of a positive-half half-link amplitude times its reflected amplitude.
Expand each tensor power and group equal multi-indices.  Every resulting
term has the form in \eqref{eq:v5v63-Higgs-rank-one} with a nonnegative
coefficient.  Absolute convergence follows pointwise from the
exponential series and in \(L^1\) from the stated radial integrability.
\end{proof}

\begin{assessment}[Higgs zeros are harmless for the cone]
\label{ass:v5v63-Higgs-zero}
At \(\rho_x=0\), every nonconstant term in
\eqref{eq:v5v63-Higgs-rank-one} vanishes and the hopping weight equals
one.  Thus the orbit-type singularity found in V56 does not create a sign
in the bosonic reflection factorisation.  It still changes the retained
electroweak band and cannot be deleted from the atlas.
\end{assessment}

\subsection{Finite bosonic electroweak reflection positivity}

Consider a finite compact electroweak lattice whose bosonic action
contains:
\begin{enumerate}
 \item Wilson plaquette weights with nonnegative couplings;
 \item gauge-covariant Higgs nearest-neighbour hoppings with
 nonnegative couplings;
 \item real on-site Higgs potentials bounded below and sufficiently
 coercive for radial integrability;
 \item reflection-symmetric local counterterms lying wholly in one half
 or on the fixed slice.
\end{enumerate}

\begin{theorem}[Finite compact gauge--Higgs OS positivity]
\label{thm:v5v63-bosonic-OS}
The finite-cutoff bosonic gauge--Higgs functional just described is
reflection positive on the gauge-invariant positive-half cylinder
algebra.  Compact gauge integration includes harmonic and toron sectors;
no gauge fixing or deletion of zero modes is required.
\end{theorem}

\begin{proof}
Plaquettes and hopping terms wholly in one open half enter \(S_+\) and
\(\Theta S_+\).  Fixed-slice potentials give a nonnegative \(W_0\).
Every crossing Wilson factor lies in the closed cone by
Proposition~\ref{prop:v5v63-Wilson-cone}, and every crossing Higgs factor
does so by Lemma~\ref{lem:v5v63-Higgs-expansion}.  Cone multiplication,
Lemma~\ref{lem:v5v63-cone-product}, puts their product in the cone.
Haar link measure, radial measure, and orbit measure have the conditional
product factorisation required by
Theorem~\ref{thm:v5v63-factorisation}.  That theorem proves the claim.
Harmonic link holonomies are integrated by the same compact Haar measure,
so they have not been removed.
\end{proof}

\begin{corollary}[Agreement with the free photon quotient]
\label{cor:v5v63-free-limit}
In a weak-field principal-stratum expansion, the nonzero transverse and
harmonic small-field limits of
Theorem~\ref{thm:v5v63-bosonic-OS} agree with the oscillator and rotor
sectors of V62, whenever that Gaussian limit exists.
\end{corollary}

\begin{proof}
The quadratic expansion of the compact Wilson action gives the Maxwell
quadratic form of Proposition~\ref{prop:v5v62-physical-form}; compact
harmonic holonomies become the toron rotor coordinates.  Reflection
positivity is preserved under convergent finite Schwinger matrices, as in
Corollary~\ref{cor:v5v62-massless-OS}.
\end{proof}

\subsection{A precise phase factorisation which preserves OS positivity}

The inherited determinant line supplies a global phase ancestry, but OS
positivity is a half-space statement.  The following exact lemma
separates the two requirements.

\begin{lemma}[Reflection-coboundary multiplier]
\label{lem:v5v63-reflection-coboundary}
Let \(\mathbb E_0\) be reflection positive on \(\mathcal A_+\).  If
\(G\in\mathcal A_+\) and
\begin{equation}
 W_G=(\Theta G)G
 \label{eq:v5v63-coboundary-weight}
\end{equation}
is integrable with
\(0<\mathbb E_0[W_G]<\infty\), then
\begin{equation}
 \mathbb E_G[F]
 :=
 \frac{\mathbb E_0[W_GF]}{\mathbb E_0[W_G]}
 \label{eq:v5v63-coboundary-functional}
\end{equation}
is reflection positive.
\end{lemma}

\begin{proof}
For \(F\in\mathcal A_+\),
\[
 \mathbb E_0[(\Theta F)FW_G]
 =
 \mathbb E_0[\Theta(FG)(FG)]
 \ge0.
\]
Divide by the positive denominator.
\end{proof}

\begin{corollary}[Half-space phase is sufficient]
\label{cor:v5v63-half-phase}
If a determinant factor relative to an OS-positive modulus has the form
\begin{equation}
 e^{i\Theta_{\mathrm{det}}}
 =
 \Theta(e^{i\phi_+})e^{i\phi_+},
 \qquad e^{i\phi_+}\in\mathcal A_+,
 \label{eq:v5v63-half-phase}
\end{equation}
then it preserves reflection positivity, provided the normalising
denominator is positive.
\end{corollary}

\begin{proof}
Take \(G=e^{i\phi_+}\) in
Lemma~\ref{lem:v5v63-reflection-coboundary}.
\end{proof}

\begin{firewall}[Global triviality is not half-space factorisation]
\label{fw:v5v63-global-not-half}
A nowhere-zero parallel section of the determinant line fixes the global
phase consistently.  It does not imply
\eqref{eq:v5v63-half-phase}.  The latter requires locality with respect to
one reflection cut and compatibility of the eta/boundary contribution
with that cut.  The inherited faithful-branch phase theorem therefore
cannot be cited as a proof of chiral OS positivity.
\end{firewall}

\subsection{The chiral Grassmann transfer gate}

Let \(\mathcal G_+\) be the positive-half Grassmann algebra, with
\(\Theta\) reversing monomial order.  Before integrating fermions, write
the cross-plane fermion factor as a finite Grassmann polynomial
\(W_F\).

\begin{definition}[Chiral reflection factorisation gate]
\label{def:v5v63-CRFG}
The chiral reflection factorisation gate holds when, in the complete
anomaly-cancelling multiplet and one common regulator,
\begin{equation}
 W_F
 =
 \sum_{\alpha=1}^N
 c_\alpha(\Theta G_\alpha)G_\alpha,
 \qquad
 c_\alpha\ge0,\quad G_\alpha\in\mathcal G_+,
 \label{eq:v5v63-fermion-cone}
\end{equation}
with the Berezin orientation induced by the inherited determinant line.
An equivalent transfer formulation is allowed if it gives the same
nonnegative reflected Gram matrices.
\end{definition}

\begin{proposition}[Exact Schur elimination preserves a proved gate only
when it respects the cut]
\label{prop:v5v63-Schur-OS-gate}
Suppose \eqref{eq:v5v63-fermion-cone} holds and a fermion shell is a direct
sum of a positive-half shell and its reflected negative-half shell, with
no shell variable supported on both sides.  Exact Berezin integration of
that shell preserves the reflected Gram inequalities.  A Schur block
whose inverse connects both sides of the reflection plane does not satisfy
this conclusion without an additional factorisation.
\end{proposition}

\begin{proof}
In the first case, Berezin--Fubini integrates each
\((\Theta G_\alpha)G_\alpha\) as the reflected product of the corresponding
positive-half shell integral, with the common determinant-line
orientation.  The coefficients remain nonnegative.  In the second case,
the Schur term contains the nonlocal kernel \(BE^{-1}C\) across the cut.
Algebraic equality of determinants gives no representation of that kernel
in the cone, so the first argument cannot be applied.
\end{proof}

\begin{firewall}[Anomaly cancellation does not prove the chiral gate]
\label{fw:v5v63-anomaly-not-CRFG}
Cancellation of local and global anomalies controls gauge covariance of
the complete fermion line.  It does not prove nonnegative coefficients in
\eqref{eq:v5v63-fermion-cone}, a positive transfer operator, or the
reflection-coboundary phase
\eqref{eq:v5v63-half-phase}.
\end{firewall}

\subsection{V63 status}

\begin{theorem}[V63 conclusion]
\label{thm:v5v63-status}
The finite compact \(U(1)\) Wilson action coupled to a radially fluctuating
gauge-covariant Higgs field is reflection positive under the stated
nearest-neighbour and nonnegative-coupling hypotheses.  The proof uses
positive character and rank-one cross-plane expansions and includes
harmonic compact sectors and Higgs zeros.  The complete chiral Standard
Model requires, in addition, either the reflection-coboundary determinant
factorisation \eqref{eq:v5v63-half-phase} or the complete Grassmann
transfer gate \eqref{eq:v5v63-fermion-cone}.  Neither follows from anomaly
cancellation or from an exact Schur determinant identity.
\end{theorem}

\begin{proof}
The bosonic statement is
Theorem~\ref{thm:v5v63-bosonic-OS}.  The two sufficient chiral routes are
Corollary~\ref{cor:v5v63-half-phase} and
Definition~\ref{def:v5v63-CRFG}; their non-implications are
Firewalls~\ref{fw:v5v63-global-not-half} and
\ref{fw:v5v63-anomaly-not-CRFG}.
\end{proof}

V64 will attack the Grassmann route locally at the reflection plane.  It
will put the finite chiral hopping matrix in temporal blocks, derive the
coefficient matrix of its cross-plane Berezin polynomial, and determine
whether Standard Model chirality permits a nonnegative factorisation or
forces a doubled/vectorlike transfer space.  A failed principal minor will
be recorded as an obstruction rather than hidden in the determinant
phase.

\clearpage
\section{V64: the chiral cross-plane matrix and Schur positivity}
\label{sec:v5v64}

V63 left a finite algebraic question at the reflection plane.  This note
answers that question for a quadratic hopping polynomial.  After the
negative-half generators are identified with reflected positive-half
generators, the coefficient matrix must be Hermitian positive
semidefinite.  This condition is both necessary and sufficient for the
exponential cross factor to lie in the Grassmann reflection cone.

The principal Weyl temporal symbol passes the test when the reflection
identifies the right-handed codomain with the reflected left-handed
domain.  A same-chirality compression is the wrong test and kills the
temporal symbol.  Shell elimination preserves the test only for a
reflection-adjoint block matrix; the general nonnormal Schur complement
of V59 does not do so automatically.

\subsection{A finite Grassmann cone theorem}

Let \(\xi_1,\ldots,\xi_N\) be positive-half Grassmann generators.  Their
reflections \(\Theta\xi_i\) are independent negative-half generators.
The antilinear involution reverses monomial order:
\begin{equation}
 \Theta(\xi_{i_1}\cdots\xi_{i_k})
 =
 (\Theta\xi_{i_k})\cdots(\Theta\xi_{i_1}).
 \label{eq:v5v64-Theta-order}
\end{equation}
For an \(N\times N\) complex matrix \(M\), define the even cross factor
\begin{equation}
 W_M
 :=
 \exp\left\{
 \sum_{i,j=1}^N
 (\Theta\xi_i)M_{ij}\xi_j
 \right\}.
 \label{eq:v5v64-WM}
\end{equation}

\begin{theorem}[Grassmann cross-matrix criterion]
\label{thm:v5v64-cross-matrix}
The polynomial \(W_M\) has a finite reflection-cone expansion
\begin{equation}
 W_M
 =
 \sum_\alpha c_\alpha
 (\Theta F_\alpha)F_\alpha,
 \qquad c_\alpha\ge0,
 \label{eq:v5v64-cone-expansion}
\end{equation}
with \(F_\alpha\) in the positive Grassmann algebra if and only if
\begin{equation}
 M=M^*\ge0.
 \label{eq:v5v64-M-positive}
\end{equation}
\end{theorem}

\begin{proof}
Assume first that \(M\ge0\).  Choose a unitary \(U\) with
\(M=U^*\operatorname{diag}(\lambda_1,\ldots,\lambda_N)U\),
\(\lambda_a\ge0\), and set
\(\eta_a=\sum_jU_{aj}\xi_j\).  Then
\begin{equation}
 W_M
 =
 \prod_{a=1}^N
 \left(1+\lambda_a(\Theta\eta_a)\eta_a\right).
 \label{eq:v5v64-product}
\end{equation}
The even nilpotent factors commute.  Expanding the product and using the
order convention \eqref{eq:v5v64-Theta-order} gives
\begin{equation}
 W_M
 =
 \sum_{S\subset\{1,\ldots,N\}}
 \left(\prod_{a\in S}\lambda_a\right)
 \Theta\!\left(\prod_{a\in S}^{\nearrow}\eta_a\right)
 \left(\prod_{a\in S}^{\nearrow}\eta_a\right),
 \label{eq:v5v64-subset-expansion}
\end{equation}
which is \eqref{eq:v5v64-cone-expansion}.

Conversely, inspect the component of bidegree \((1,1)\), one reflected
and one unreflected generator.  In \(W_M\) its coefficient matrix is
\(M\).  If
\[
 F_\alpha
 =
 a_\alpha+\sum_jv_{\alpha j}\xi_j+\text{higher degree},
\]
the bidegree-\((1,1)\) coefficient of
\((\Theta F_\alpha)F_\alpha\) is
\(v_\alpha^*v_\alpha\).  Constant terms paired with higher terms have
bidegrees \((0,k)\) or \((k,0)\) and do not contribute.  Hence
\[
 M=\sum_\alpha c_\alpha v_\alpha^*v_\alpha,
\]
which is Hermitian positive semidefinite.
\end{proof}

\begin{corollary}[Principal-minor obstruction]
\label{cor:v5v64-principal-minor}
If \(M\) is non-Hermitian or has a negative principal minor, then the
cross factor \(W_M\) cannot satisfy the chiral reflection factorisation
gate of V63.
\end{corollary}

\begin{proof}
A Hermitian matrix is positive semidefinite if and only if all of its
principal minors are nonnegative.  Apply
Theorem~\ref{thm:v5v64-cross-matrix}.
\end{proof}

\begin{assessment}[Basis covariance]
\label{ass:v5v64-basis-covariance}
If the positive generators change by \(\xi'=V\xi\), the cross matrix
changes by congruence \(M'=(V^{-1})^*MV^{-1}\).  Positivity is preserved.
A unitary gauge half-link can therefore be absorbed into the
positive-half frame when the reflection maps the other half-link to its
adjoint.  Similarity without the adjoint is not enough.
\end{assessment}

\subsection{The principal Weyl temporal crossing}

Let \(\gamma_0=\gamma_0^*\), \(\gamma_0^2=1\), and let
\(\gamma_5=\gamma_5^*\) anticommute with \(\gamma_0\).  Put
\begin{equation}
 P_L=\frac{1-\gamma_5}{2},
 \qquad
 P_R=\frac{1+\gamma_5}{2}.
 \label{eq:v5v64-chiral-projectors}
\end{equation}
Then
\begin{equation}
 \gamma_0P_L=P_R\gamma_0,
 \qquad
 P_L\gamma_0P_L=0.
 \label{eq:v5v64-chirality-flip}
\end{equation}
The chiral temporal symbol is the map
\begin{equation}
 B_0:=P_R\gamma_0P_L:
 P_LH\longrightarrow P_RH.
 \label{eq:v5v64-Weyl-temporal}
\end{equation}

\begin{definition}[Chiral reflection identification]
\label{def:v5v64-reflection-identification}
The one-particle reflection identification is
\begin{equation}
 J_0:=P_L\gamma_0P_R:
 P_RH\longrightarrow P_LH.
 \label{eq:v5v64-J0}
\end{equation}
It identifies the negative-half right-handed codomain with the reflected
dual of the positive-half left-handed domain.
\end{definition}

\begin{proposition}[The principal Weyl block is positive]
\label{prop:v5v64-Weyl-positive}
On \(P_LH\),
\begin{equation}
 J_0B_0=P_L.
 \label{eq:v5v64-Weyl-Gram}
\end{equation}
Thus the reflected cross matrix of the principal temporal Weyl hopping is
the identity on its chiral domain and satisfies
\eqref{eq:v5v64-M-positive}.
\end{proposition}

\begin{proof}
Using \(\gamma_0P_R\gamma_0=P_L\),
\[
 J_0B_0
 =
 P_L\gamma_0P_R\gamma_0P_L
 =
 P_L.
\]
Restricted to \(P_LH\), this is the identity and is positive.
\end{proof}

\begin{firewall}[Same-chirality compression is the wrong pairing]
\label{fw:v5v64-same-chirality}
Equation \eqref{eq:v5v64-chirality-flip} gives
\(P_L\gamma_0P_L=0\).  Therefore placing both the domain and reflected
codomain in the same \(P_L\) subspace erases the first-order temporal
symbol.  The OS pairing must use the chiral complex
\(P_LH\to P_RH\), consistent with the inherited determinant line.
\end{firewall}

Let \(V\) be a unitary internal parallel transport on a crossing link,
commuting with the spin projectors.  If reflection sends the negative
half-link to \(V^*\), use
\begin{equation}
 J_V:=V^*J_0.
 \label{eq:v5v64-JV}
\end{equation}

\begin{corollary}[Gauge transport does not change the principal spectrum]
\label{cor:v5v64-gauge-positive}
For the transported temporal block \(B_V=VB_0\),
\begin{equation}
 J_VB_V=P_L.
 \label{eq:v5v64-gauge-Gram}
\end{equation}
Hence a reflection-compatible unitary half-link preserves the local
principal cross-plane cone.
\end{corollary}

\begin{proof}
Internal \(V\) commutes with \(J_0\), and \(V^*V=1\).  Apply
Proposition~\ref{prop:v5v64-Weyl-positive}.
\end{proof}

\begin{firewall}[A principal-symbol check is not a regulator theorem]
\label{fw:v5v64-principal-only}
Proposition~\ref{prop:v5v64-Weyl-positive} tests one nearest-plane
principal hopping block.  A gauge-covariant chiral regulator can contain
Wilson terms, Ginsparg--Wilson projectors depending on the gauge field,
longer-range temporal couplings, boundary eta data, and zero-mode
insertions.  Their complete cross matrix must pass
\eqref{eq:v5v64-M-positive}; the principal result does not imply that it
does.
\end{firewall}

\subsection{Reflection-adjoint Schur elimination}

Split the positive generators into retained and shell groups.  Let the
complete reflected cross matrix be
\begin{equation}
 \mathbb M=
 \begin{pmatrix}
 A&B\\ B^*&E
 \end{pmatrix},
 \qquad E>0.
 \label{eq:v5v64-positive-block}
\end{equation}

\begin{theorem}[Positive Schur inheritance]
\label{thm:v5v64-positive-Schur}
If \(\mathbb M\ge0\), then the retained Schur matrix
\begin{equation}
 M_{\mathrm{ret}}
 :=
 A-BE^{-1}B^*
 \label{eq:v5v64-positive-Schur}
\end{equation}
is positive semidefinite.  Consequently exact shell elimination preserves
the quadratic cross-plane Grassmann cone when the shell splitting respects
the reflection-adjoint block structure.
\end{theorem}

\begin{proof}
For every retained vector \(x\), choose
\(y=-E^{-1}B^*x\).  Then
\begin{align}
 0
 &\le
 \left\langle
 \binom{x}{y},
 \mathbb M\binom{x}{y}
 \right\rangle
 =
 \langle x,(A-BE^{-1}B^*)x\rangle.
 \label{eq:v5v64-Schur-square}
\end{align}
Thus \(M_{\mathrm{ret}}\ge0\).  Theorem~\ref{thm:v5v64-cross-matrix}
converts this matrix statement to the cone statement.
\end{proof}

\begin{corollary}[Quantitative retained lower bound]
\label{cor:v5v64-Schur-lower}
If \(A\ge aI\), \(E\ge eI\), and \(\|B\|\le b\), then
\begin{equation}
 M_{\mathrm{ret}}
 \ge
 \left(a-\frac{b^2}{e}\right)I.
 \label{eq:v5v64-Schur-lower}
\end{equation}
In particular, strict cross-plane positivity survives when \(b^2<ae\).
\end{corollary}

\begin{proof}
Since \(E^{-1}\le e^{-1}I\),
\[
 BE^{-1}B^*\le e^{-1}BB^*\le b^2e^{-1}I.
\]
Subtract from \(A\).
\end{proof}

\begin{firewall}[Nonnormal locality does not imply OS positivity]
\label{fw:v5v64-nonnormal}
V59 permits a general Schur term \(A-BE^{-1}C\) and proves its locality
from singular values even when \(E\) is nonnormal.  Theorem
\ref{thm:v5v64-positive-Schur} additionally requires
\(C=B^*\), \(E=E^*>0\), after the reflection identification.  Without
that structure, the Schur matrix can be non-Hermitian or negative even
when \(E^{-1}\) is exponentially local.
\end{firewall}

\begin{example}[Local inverse with failed reflection sign]
\label{ex:v5v64-negative-Schur}
Take scalar blocks
\begin{equation}
 A=1,\qquad E=1,\qquad B=2,\qquad C=2.
 \label{eq:v5v64-negative-example}
\end{equation}
The shell inverse is perfectly local and has norm one, but
\begin{equation}
 A-BE^{-1}C=-3<0.
 \label{eq:v5v64-negative-result}
\end{equation}
Thus locality and invertibility alone do not preserve the reflection
cone.
\end{example}

\subsection{Multiple temporal ranges}

A regulator with temporal range \(R>1\) couples several positive and
negative time layers.  Collect all generators in the first \(R\) layers
into one boundary vector \(\Xi\).  Reflection turns all cross-plane
quadratic terms into
\begin{equation}
 \sum_{a,b}(\Theta\Xi_a)\mathbb M^{(R)}_{ab}\Xi_b.
 \label{eq:v5v64-range-R-matrix}
\end{equation}

\begin{proposition}[Finite-range boundary test]
\label{prop:v5v64-range-R}
For a finite-range quadratic regulator, its complete exponential
cross-plane factor belongs to the Grassmann reflection cone if and only if
\begin{equation}
 \mathbb M^{(R)}=(\mathbb M^{(R)})^*\ge0.
 \label{eq:v5v64-range-R-test}
\end{equation}
This is a finite principal-minor test at every finite cutoff.
\end{proposition}

\begin{proof}
The enlarged boundary vector is still finite.  Apply
Theorem~\ref{thm:v5v64-cross-matrix} with \(M=\mathbb M^{(R)}\).
\end{proof}

\begin{assessment}[Nonlocal chiral kernels]
\label{ass:v5v64-nonlocal}
For an exponentially decaying but infinite-range temporal kernel, truncate
at range \(R\) and apply
Proposition~\ref{prop:v5v64-range-R}.  Passing \(R\to\infty\) requires
both convergence of the Grassmann functional and a uniform nonnegative
boundary form.  Small operator-norm tails do not repair a negative
principal minor already present at finite \(R\).
\end{assessment}

\subsection{Reflection-compatible fermion shell certificate}

\begin{definition}[Reflection-compatible fermion shell certificate]
\label{def:v5v64-RCFSC}
An RCFSC consists of:
\begin{enumerate}
 \item the complete three-generation regulated chiral complex with the
 inherited determinant-line orientation;
 \item an antilinear reflection pairing its right-handed codomain with the
 reflected dual of its left-handed domain;
 \item a complete cross-plane matrix satisfying
 \eqref{eq:v5v64-range-R-test}, or a convergent infinite-range analogue;
 \item shell splittings with the reflection-adjoint form
 \eqref{eq:v5v64-positive-block};
 \item strict Schur margins such as
 \eqref{eq:v5v64-Schur-lower} wherever shell elimination is used;
 \item exact retention of every zero mode or band at which these margins
 fail;
 \item compatibility with gauge averaging, Higgs crossover strata, and
 the V60 common-refinement cocycle.
\end{enumerate}
\end{definition}

\begin{theorem}[V64 conclusion]
\label{thm:v5v64-status}
A quadratic fermion cross-plane exponential lies in the Grassmann
reflection cone exactly when its reflected coefficient matrix is
Hermitian positive semidefinite.  The principal Weyl temporal symbol
passes this test with the correct left-domain/right-codomain reflection
pairing, including reflection-compatible unitary gauge transport.
Reflection-adjoint shell elimination preserves the cone through the
positive Schur complement, whereas the general nonnormal local Schur
block of V59 does not.  This proves the algebraic and principal-symbol
parts of an RCFSC, not the complete regulator-level Standard Model
certificate.
\end{theorem}

\begin{proof}
Use Theorem~\ref{thm:v5v64-cross-matrix},
Proposition~\ref{prop:v5v64-Weyl-positive},
Corollary~\ref{cor:v5v64-gauge-positive}, and
Theorem~\ref{thm:v5v64-positive-Schur}, with
Firewall~\ref{fw:v5v64-principal-only}.
\end{proof}

V65 will perform the compulsory YM/NS audit.  It will then apply the
finite principal-minor test to the simplest complete temporal regulator
block available in the inherited Standard Model construction.  If that
block is not specified strongly enough to compute
\(\mathbb M^{(R)}\), V65 will record the missing regulator datum and move
upstream to a regulator-independent reflection-transfer criterion rather
than assume positivity.

\clearpage
\section{V65: companion audit and the heavy-determinant cancellation gate}
\label{sec:v5v65}

\subsection{Compulsory V65 companion audit}

The newest live companion files inspected before this note were
\begin{align}
 &\texttt{Renewal\_Yang\_Mills\_Mass\_Gap\_Research\_Notes\_V78.tex},
 \notag\\
 &\qquad
 \operatorname{SHA256}=
 \texttt{870ECB0BDD16CB287B2780AEDFD2A880}
 \texttt{AA176525AEC3FBC2BABDDB34A7ADE962},
 \label{eq:v5v65-YM-audit}\\
 &\texttt{Renewal\_NS\_Stability\_Research\_Notes\_V31.tex},
 \notag\\
 &\qquad
 \operatorname{SHA256}=
 \texttt{F1121B62238AD5491BB7CF03635EA993}
 \texttt{4942EDF573ED8FAAA9EE79E1D38A6696}.
 \label{eq:v5v65-NS-audit}
\end{align}
The YM file had \(1\,150\,236\) bytes.  The NS file had
\(887\,537\) bytes and is byte-identical to the V60 audit.  Neither
companion file was modified.

YM V78 proves a full one-link \(SU(2)\) Wilson Poincare inequality with
spectral gap of order \(\alpha\), an intrinsic covariance/resolvent
identity, and the first differentiated-resolvent equation.  Its next
obstruction is all-order iteration without Bell or factorial growth.
The note explicitly does not prove reflection positivity, the continuum
Yang--Mills measure, or a mass gap.

The useful transfer is structural.  A positive one-link spectral gap can
later control centered bosonic boundary dependence, but it cannot be
inserted into a chiral determinant quotient as if it were a fermionic
transfer estimate.  As in the YM all-order problem, exact cancellation of
a multiplicity does not create the analytic positivity needed after the
cancellation.

\subsection{Correction from the frozen Standard Model archive}

The search required by V64 found an earlier result that changes the next
step.  Frozen third-series V76 already constructs a finite,
nearest-neighbour five-dimensional Wilson domain-wall operator.  It proves
exact gauge covariance and reflection positivity of the complete
finite two-wall Grassmann block.  Its cross-plane temporal hopping uses
the positive projectors \(P_\pm=(1\pm\gamma_4)/2\), which is an explicit
instance of the V64 cross-matrix criterion.

That inherited theorem retains the mirror wall and all heavy modes.  Its
own firewall leaves unproved reflection positivity after division by a
Pauli--Villars determinant, passage to the overlap determinant ratio, or
discarding the remote wall.  Thus V65 does not repeat the two-wall
calculation.  It examines the cancellation operation itself.

\subsection{Finite reflection kernels}

Let \(X=\{1,\ldots,N\}\) be a finite set of boundary configurations.  A
reflection-positive transfer kernel is a Hermitian matrix
\begin{equation}
 K=(K_{ij})_{i,j\in X}\ge0.
 \label{eq:v5v65-positive-kernel}
\end{equation}
For matrices \(A,B\) of the same size, write \(A\circ B\) for their
entrywise product.  If every entry of \(H\) is nonzero, write
\begin{equation}
 H^{\circ-1}_{ij}:=\frac1{H_{ij}}.
 \label{eq:v5v65-Hadamard-inverse}
\end{equation}

\begin{theorem}[Exact quotient-kernel criterion]
\label{thm:v5v65-quotient-criterion}
Let \(H\) have no zero entries.  The entrywise quotient map
\begin{equation}
 \mathcal Q_H(K):=
 K\circ H^{\circ-1}
 \label{eq:v5v65-quotient-map}
\end{equation}
sends every positive semidefinite matrix \(K\) to a positive semidefinite
matrix if and only if
\begin{equation}
 H^{\circ-1}\ge0.
 \label{eq:v5v65-reciprocal-positive}
\end{equation}
\end{theorem}

\begin{proof}
If \eqref{eq:v5v65-reciprocal-positive} holds, the Schur product theorem
gives
\(K\circ H^{\circ-1}\ge0\) for every \(K\ge0\).
Conversely, apply \(\mathcal Q_H\) to the all-ones matrix
\(\mathbf1\mathbf1^*\ge0\).  Its image is exactly
\(H^{\circ-1}\), proving necessity.
\end{proof}

\begin{theorem}[Universal positive cancellation is a coboundary]
\label{thm:v5v65-coboundary-necessity}
Suppose \(H=H^*\ge0\), every entry of \(H\) is nonzero, and the quotient
map \(\mathcal Q_H\) preserves every positive semidefinite kernel.  Then
\(H\) has rank one:
\begin{equation}
 H_{ij}=\overline{g_i}g_j
 \label{eq:v5v65-rank-one}
\end{equation}
for nonzero numbers \(g_i\).  Conversely, every matrix of the form
\eqref{eq:v5v65-rank-one} has this universal cancellation property.
\end{theorem}

\begin{proof}
By Theorem~\ref{thm:v5v65-quotient-criterion},
\(R:=H^{\circ-1}\ge0\).  Positivity of every \(2\times2\) principal
minor of \(H\) gives
\begin{equation}
 |H_{ij}|^2\le H_{ii}H_{jj}.
 \label{eq:v5v65-H-CS}
\end{equation}
The same inequality for \(R\) gives
\[
 |H_{ij}|^{-2}
 \le (H_{ii}H_{jj})^{-1},
\]
which is the reverse inequality to
\eqref{eq:v5v65-H-CS}.  Equality therefore holds for every pair.  View
\(H\) as a Gram matrix \(H_{ij}=\langle v_i,v_j\rangle\).  Equality in
Cauchy--Schwarz makes every pair \(v_i,v_j\) collinear.  Since no entry is
zero, all \(v_i\) lie on one nonzero line, so \(H\) has rank one and can
be written as \eqref{eq:v5v65-rank-one}.

Conversely, if \eqref{eq:v5v65-rank-one} holds, then
\[
 H^{\circ-1}_{ij}
 =
 \overline{g_i^{-1}}g_j^{-1}
\]
is positive semidefinite.  Apply
Theorem~\ref{thm:v5v65-quotient-criterion}.
\end{proof}

\begin{corollary}[Reflection-coboundary interpretation]
\label{cor:v5v65-reflection-coboundary}
The rank-one kernel \eqref{eq:v5v65-rank-one} is the sampled form of a
reflection coboundary
\begin{equation}
 H(x_-,x_+)
 =
 (\Theta G)(x_-)G(x_+).
 \label{eq:v5v65-kernel-coboundary}
\end{equation}
Division by it is the congruence
\[
 K\longmapsto
 \operatorname{diag}(g)^{-*}
 K\operatorname{diag}(g)^{-1},
\]
and hence preserves positivity in both directions.
\end{corollary}

\begin{proof}
Sampling \eqref{eq:v5v65-kernel-coboundary} gives
\eqref{eq:v5v65-rank-one}.  The displayed congruence has entries
\(K_{ij}/(\overline{g_i}g_j)\), and congruence by an invertible matrix
preserves positive semidefiniteness.
\end{proof}

\subsection{Why two positive determinants are not enough}

\begin{example}[Positive numerator and denominator, negative quotient]
\label{ex:v5v65-positive-ratio-fails}
Let \(0<a<1\), and take
\begin{equation}
 K=
 \begin{pmatrix}1&1\\1&1\end{pmatrix},
 \qquad
 H=
 \begin{pmatrix}1&a\\a&1\end{pmatrix}.
 \label{eq:v5v65-two-positive}
\end{equation}
Both matrices are positive semidefinite, but
\begin{equation}
 K\circ H^{\circ-1}
 =
 \begin{pmatrix}1&a^{-1}\\a^{-1}&1\end{pmatrix}
 \label{eq:v5v65-negative-quotient}
\end{equation}
has eigenvalue \(1-a^{-1}<0\).
\end{example}

\begin{proof}
The eigenvalues of \(H\) are \(1\pm a>0\), and \(K\) has eigenvalues
\(2,0\).  The quotient eigenvalues are \(1\pm a^{-1}\).
\end{proof}

\begin{corollary}[Strict transfer mixing obstructs universal division]
\label{cor:v5v65-strict-mixing}
If a positive kernel \(H\) has two boundary configurations satisfying
\begin{equation}
 |H_{ij}|^2<H_{ii}H_{jj},
 \label{eq:v5v65-strict-CS}
\end{equation}
then entrywise division by \(H\) fails to preserve some
reflection-positive kernel.
\end{corollary}

\begin{proof}
The \(2\times2\) principal minor of \(H^{\circ-1}\) has determinant
\[
 \frac1{H_{ii}H_{jj}}-\frac1{|H_{ij}|^2}<0.
\]
Use Theorem~\ref{thm:v5v65-quotient-criterion}.
\end{proof}

\begin{assessment}[Universal versus ratio-specific positivity]
\label{ass:v5v65-universal-specific}
Theorems~\ref{thm:v5v65-quotient-criterion} and
\ref{thm:v5v65-coboundary-necessity} concern cancellation that is safe
for every physical kernel.  A particular correlated numerator \(K\) may
still have \(K\circ H^{\circ-1}\ge0\) even when \(H\) is not rank one.
That weaker statement requires direct positivity of every principal
minor of the actual quotient kernel.
\end{assessment}

\subsection{Application to the domain-wall/overlap ratio}

At finite wall separation, denote the sampled reflected kernel of the
complete two-wall theory by \(K_{\mathrm{DW}}\).  The inherited V76 theorem
gives
\begin{equation}
 K_{\mathrm{DW}}\ge0.
 \label{eq:v5v65-DW-positive}
\end{equation}
Let \(H_{\mathrm{PV}}\) denote the boundary kernel whose determinant
implements the heavy Pauli--Villars cancellation, and let
\begin{equation}
 K_{\mathrm{ratio}}
 :=
 K_{\mathrm{DW}}\circ H_{\mathrm{PV}}^{\circ-1}
 \label{eq:v5v65-ratio-kernel}
\end{equation}
whenever this sampled scalar quotient is defined.

\begin{definition}[Heavy-determinant cancellation certificate]
\label{def:v5v65-HDCC}
An HDCC is either:
\begin{enumerate}
 \item a nonvanishing positive-half functional \(G_{\mathrm{PV}}\) for
 which the complete heavy kernel is the reflection coboundary
 \begin{equation}
 H_{\mathrm{PV}}=(\Theta G_{\mathrm{PV}})G_{\mathrm{PV}};
 \label{eq:v5v65-PV-coboundary}
 \end{equation}
 or
 \item a direct proof that every finite sampled matrix of the actual
 quotient kernel \(K_{\mathrm{ratio}}\) is positive semidefinite, compatible
 with refinement and the determinant-line orientation.
\end{enumerate}
\end{definition}

\begin{theorem}[What is required to pass from V76 to a chiral quotient]
\label{thm:v5v65-DW-to-chiral}
The inherited positivity
\eqref{eq:v5v65-DW-positive} descends through Pauli--Villars cancellation
if an HDCC holds.  Positivity of \(H_{\mathrm{PV}}\) by itself is
insufficient.  If the cancellation is required to preserve arbitrary
reflection-positive physical boundary kernels, the coboundary branch
\eqref{eq:v5v65-PV-coboundary} is necessary as well as sufficient.
\end{theorem}

\begin{proof}
The coboundary branch acts by congruence, so
Corollary~\ref{cor:v5v65-reflection-coboundary} proves positivity.  The
ratio-specific branch is the definition of reflected Gram positivity.
Example~\ref{ex:v5v65-positive-ratio-fails} proves that separate
positivity is insufficient.  Universal necessity is
Theorem~\ref{thm:v5v65-coboundary-necessity}.
\end{proof}

\begin{firewall}[Inverse determinants are not positive bosonic Gaussians
for free]
\label{fw:v5v65-inverse-determinant}
Writing
\((\det D_{\mathrm{PV}})^{-1}\) as a commuting spinor integral requires a
quadratic form with a convergent contour.  Replacing it by
\((\det D_{\mathrm{PV}}^*D_{\mathrm{PV}})^{-1}\) supplies a positive
Gaussian only for the modulus squared and changes the chiral phase.
Neither representation proves an HDCC for the intended determinant ratio.
\end{firewall}

\begin{firewall}[Marginalisation is not cancellation]
\label{fw:v5v65-marginal-not-cancel}
Exact Berezin integration of a heavy shell retains its determinant factor
and can preserve OS positivity under the V64 reflection-adjoint
hypotheses.  Dividing by that determinant removes closed heavy loops and
is a different operation.  The positivity of a marginal cannot be cited
as positivity of the determinant ratio.
\end{firewall}

\begin{firewall}[A common boundary-kernel representation is required]
\label{fw:v5v65-common-kernel}
The entrywise formula \eqref{eq:v5v65-ratio-kernel} applies only after
the numerator and cancelling factor have been represented on the same
reflection boundary, with all purely positive- and negative-half factors
stripped consistently.  A determinant ratio defined first on complete
four-dimensional gauge configurations need not become the quotient of two
separately integrated transfer kernels.  Constructing the common boundary
representation, or retaining the Pauli--Villars fields in one joint
transfer functional, is part of the HDCC.
\end{firewall}

\subsection{A finite computable test}

For boundary configurations \(x_1,\ldots,x_N\), form
\begin{equation}
 \mathsf K^{(N)}_{ij}
 =
 \frac{K_{\mathrm{DW}}(x_i,x_j)}
 {H_{\mathrm{PV}}(x_i,x_j)}.
 \label{eq:v5v65-sampled-ratio}
\end{equation}

\begin{proposition}[Sampled witness alternative]
\label{prop:v5v65-witness}
At each finite cutoff and wall separation, exactly one of the following
holds:
\begin{enumerate}
 \item every finite sample \(\mathsf K^{(N)}\) is positive semidefinite,
 so the scalar boundary ratio is a positive-type kernel;
 \item some finite sample has a negative principal minor, which is a
 rigorous obstruction to reflection positivity of that ratio.
\end{enumerate}
If the boundary configuration space is finite, samples of size at most
its cardinality decide the alternative.
\end{proposition}

\begin{proof}
A kernel is of positive type precisely when all of its finite Gram
matrices are positive semidefinite.  A Hermitian matrix fails positivity
precisely when some principal minor is negative, or equivalently when it
has a negative eigenvalue.  The finite-space statement is immediate.
\end{proof}

\begin{assessment}[Missing inherited datum]
\label{ass:v5v65-missing-datum}
The frozen notes specify the local Wilson domain-wall operator and prove
the complete two-wall cone, but they do not provide the full
background-dependent cross-plane kernel
\(H_{\mathrm{PV}}(x_-,x_+)\) or prove
\eqref{eq:v5v65-PV-coboundary}.  Therefore V64's finite matrix test cannot
yet be evaluated on the chiral determinant ratio.  The missing object is
the reflected heavy-cancellation kernel, not another anomaly sum.
\end{assessment}

\subsection{V65 status}

\begin{theorem}[V65 conclusion]
\label{thm:v5v65-status}
The mandatory audit records YM V78 and unchanged NS V31.  The inherited
Standard Model archive already proves reflection positivity of the
complete finite two-wall Wilson domain-wall block.  Passing to the chiral
overlap determinant ratio requires a new heavy-determinant cancellation
certificate.  Entrywise division by a positive reflection kernel is not
positivity preserving in general; universal safe cancellation is
equivalent to a rank-one reflection coboundary.  The existing notes do not
supply that Pauli--Villars boundary kernel or a ratio-specific
principal-minor proof, so chiral OS positivity remains open at this exact
gate.
\end{theorem}

\begin{proof}
Use the audit, Theorems~\ref{thm:v5v65-quotient-criterion},
\ref{thm:v5v65-coboundary-necessity}, and
\ref{thm:v5v65-DW-to-chiral}, together with
Assessment~\ref{ass:v5v65-missing-datum}.
\end{proof}

V66 will go upstream to the local operator ratio.  It will derive the
finite-dimensional Pauli--Villars cancellation as a block Gaussian
identity, distinguish determinant, modulus, and phase factors, and test
whether the heavy ratio is a reflection coboundary.  If a nontrivial
connected boundary kernel remains, it will formulate the correlated
numerator/denominator Gram condition rather than divide two positivity
proofs.

\clearpage
\section{V66: an exact local Pauli--Villars supertransfer representation}
\label{sec:v5v66}

The inherited V77 determinant identity already performs the common-bulk
algebra:
\[
 \frac{\det D_{\mathrm{DW},N}(m_f)}
 {\det D_{\mathrm{DW},N}(1)}
 =
 \det D_N(m_f).
\]
Repeating that Schur calculation would not advance the programme.  Its
proof is an identity of determinants, not a proof that the quotient has a
positive reflected transfer kernel.  V65 showed why positivity cannot be
obtained by dividing two separately positive kernels.

This note constructs a different exact finite representation.  The
inverse Pauli--Villars determinant is written as one conjugate Grassmann
determinant divided by the positive square
\(D_{\mathrm{PV}}^*D_{\mathrm{PV}}\).  The representation preserves the
full complex phase and is local with twice the original hopping range.
It converts the open quotient question into explicit Grassmann and
commuting-field cross-matrix tests.

\subsection{Why the naive commuting-spinor integral can diverge}

For a complex \(n\times n\) matrix \(D\), the formal identity
\begin{equation}
 \int_{\mathbb C^n}
 e^{-\phi^*D\phi}\,\frac{d^{2n}\phi}{\pi^n}
 \stackrel{\mathrm{formal}}{=}
 \frac1{\det D}
 \label{eq:v5v66-naive-boson}
\end{equation}
is valid on the standard real contour only under a convergence
hypothesis, for example
\begin{equation}
 \frac{D+D^*}{2}>0.
 \label{eq:v5v66-accretive}
\end{equation}

\begin{lemma}[Standard-contour convergence]
\label{lem:v5v66-standard-contour}
If the Hermitian part of \(D\) has a negative direction, the integral in
\eqref{eq:v5v66-naive-boson} is not absolutely convergent.  If
\eqref{eq:v5v66-accretive} holds, it is absolutely convergent and equals
\((\det D)^{-1}\), with the determinant branch obtained by continuous
deformation from a positive Hermitian matrix.
\end{lemma}

\begin{proof}
If \(\operatorname{Re}\langle v,Dv\rangle<0\), restrict the integrand to
the ray \(\phi=tv\).  Its modulus grows as
\(\exp\{c t^2\}\), so absolute convergence fails.  Under
\eqref{eq:v5v66-accretive}, the modulus is bounded by a strictly
decaying Gaussian.  Schur triangularisation, or analytic continuation
from the positive Hermitian cone inside the accretive cone, gives the
determinant formula.
\end{proof}

For the inherited domain-wall operator,
\begin{equation}
 D_{\mathrm{DW}}
 =
 \sum_A\gamma_A\nabla_A
 \sum_A\frac{r a_A}{2}\Delta_A
 m(s),
 \label{eq:v5v66-domain-wall}
\end{equation}
the covariant differences are anti-Hermitian and the Wilson Laplacians
are nonnegative.  Hence
\begin{equation}
 \frac{D_{\mathrm{DW}}+D_{\mathrm{DW}}^*}{2}
 =
 \sum_A\frac{r a_A}{2}\Delta_A+m(s).
 \label{eq:v5v66-domain-wall-Hermitian}
\end{equation}

\begin{proposition}[Negative wall plateau defeats the naive integral]
\label{prop:v5v66-negative-plateau}
Suppose the wall mass obeys \(m(s)\le-m_0<0\) on \(K\) consecutive
fifth-direction sites and the other directions admit a covariantly
constant test section on the chosen background.  For \(K\) sufficiently
large compared with \(r/(a_5m_0)\), the Hermitian part
\eqref{eq:v5v66-domain-wall-Hermitian} has a negative direction.
Consequently the standard-contour commuting-spinor representation
\eqref{eq:v5v66-naive-boson} is unavailable for that domain-wall block.
\end{proposition}

\begin{proof}
Choose a fifth-direction test sequence which equals one on the middle
\(K-2\) sites of the plateau, vanishes outside the plateau, and changes
only across the two endpoint links.  Extend it by the covariantly constant
section in the other directions.  Its mass form is at most
\(-m_0(K-2)\) times the constant fibre norm.  The fifth Wilson form is
\[
 \frac{r a_5}{2}\langle v,\Delta_5v\rangle
 =
 \frac{r}{2a_5}\sum_s|v(s+1)-v(s)|^2,
\]
which is bounded by a constant independent of \(K\) for this test
sequence.  All other Wilson forms vanish.  After division by
\(\|v\|^2\), the positive contribution tends to zero and the mass
contribution tends to at most \(-m_0\).  It is therefore negative for
large \(K\).  Lemma~\ref{lem:v5v66-standard-contour} applies.
\end{proof}

\begin{assessment}[Contour rotation is additional OS data]
\label{ass:v5v66-contour}
A complex contour may define \((\det D)^{-1}\) even when
\eqref{eq:v5v66-accretive} fails.  To use it in the continuum programme
one must prove convergence, absence of crossed singularities, reflection
invariance of the cycle, and equality of all reflected matrix elements.
The determinant identity alone supplies none of these facts.
\end{assessment}

\subsection{Exact Hermitianised inverse determinant}

Let \(D\) now be invertible.  Introduce an independent Grassmann pair
\((\bar\chi,\chi)\) and a complex commuting field \(\phi\).

\begin{theorem}[Exact phase-preserving inverse determinant]
\label{thm:v5v66-hermitianised-inverse}
With the standard finite-dimensional normalisations,
\begin{align}
 \int d\bar\chi\,d\chi\,
 e^{-\bar\chi D^*\chi}
 &=
 \det D^*,
 \label{eq:v5v66-conjugate-Grassmann}\\
 \int_{\mathbb C^n}
 e^{-\phi^*D^*D\phi}\,
 \frac{d^{2n}\phi}{\pi^n}
 &=
 \frac1{\det(D^*D)}.
 \label{eq:v5v66-positive-pseudofermion}
\end{align}
Their product is exactly
\begin{equation}
 \frac{\det D^*}{\det(D^*D)}
 =
 \frac1{\det D}.
 \label{eq:v5v66-exact-inverse}
\end{equation}
No determinant phase is discarded.
\end{theorem}

\begin{proof}
The first formula is the finite Grassmann Gaussian identity.  The matrix
\(D^*D\) is strictly positive, so diagonalising it proves the commuting
Gaussian formula.  Finally
\[
 \det(D^*D)=\det D^*\det D,
\]
and invertibility permits cancellation of \(\det D^*\).
\end{proof}

\begin{corollary}[Exact local representation of the finite ratio]
\label{cor:v5v66-ratio-representation}
For invertible \(D_{\mathrm{PV}}\),
\begin{equation}
 \frac{\det D_{\mathrm{phys}}}{\det D_{\mathrm{PV}}}
 =
 \int
 e^{-\bar\psi D_{\mathrm{phys}}\psi
    -\bar\chi D_{\mathrm{PV}}^*\chi
    -\phi^*D_{\mathrm{PV}}^*D_{\mathrm{PV}}\phi}
 \,d\bar\psi\,d\psi\,d\bar\chi\,d\chi\,
 \frac{d\phi}{\pi^n}.
 \label{eq:v5v66-supertransfer-integral}
\end{equation}
If both Dirac matrices have hopping range \(R\), the two Grassmann
actions have range \(R\) and the commuting action has range at most
\(2R\).  All three actions are gauge covariant when the original
operators are.
\end{corollary}

\begin{proof}
Integrate the \(\psi\) variables and use
Theorem~\ref{thm:v5v66-hermitianised-inverse} for the other two fields.
Products of two range-\(R\) matrices have range at most \(2R\).
Gauge covariance
\(D(U^g)=gD(U)g^{-1}\) implies the same identity for \(D^*\) and
\(D^*D\); transform all auxiliary fields by \(g\).
\end{proof}

\begin{firewall}[The Pauli--Villars block must be invertible]
\label{fw:v5v66-PV-invertible}
If \(D_{\mathrm{PV}}\) has a zero mode, neither side of
\eqref{eq:v5v66-exact-inverse} is defined.  Such a mode cannot be removed
by formal cancellation; it must be excluded by a proved heavy gap or
re-fibred into an explicitly retained regulator sector.
\end{firewall}

\subsection{The commuting cross-matrix criterion}

Let \(z=(z_1,\ldots,z_N)\) be complex positive-half commuting variables,
with reflected variables \(\Theta z_i\).  For a matrix \(M\), define the
entire cross factor
\begin{equation}
 \mathcal W_M(z)
 :=
 \exp\left\{
 \sum_{i,j}(\Theta z_i)M_{ij}z_j
 \right\}.
 \label{eq:v5v66-bosonic-cross}
\end{equation}

\begin{theorem}[Commuting cross-matrix criterion]
\label{thm:v5v66-bosonic-cross}
The factor \(\mathcal W_M\) has an absolutely convergent expansion
\begin{equation}
 \mathcal W_M
 =
 \sum_\alpha c_\alpha
 (\Theta F_\alpha)F_\alpha,
 \qquad c_\alpha\ge0,
 \label{eq:v5v66-bosonic-cone}
\end{equation}
on every compact set if and only if \(M=M^*\ge0\).
\end{theorem}

\begin{proof}
If \(M=U^*\operatorname{diag}(\lambda_a)U\), set \(w=Uz\).  For
\(\lambda_a\ge0\),
\begin{align}
 \mathcal W_M
 &=
 \prod_a
 \exp\{\lambda_a(\Theta w_a)w_a\}
 \notag\\
 &=
 \sum_{(k_a)\in\mathbb N^N}
 \left(\prod_a\frac{\lambda_a^{k_a}}{k_a!}\right)
 \Theta\!\left(\prod_aw_a^{k_a}\right)
 \left(\prod_aw_a^{k_a}\right).
 \label{eq:v5v66-bosonic-multiindex}
\end{align}
This is an absolutely convergent cone expansion on compact sets.
Conversely, the coefficient of bidegree \((1,1)\) in any expansion
\eqref{eq:v5v66-bosonic-cone} is a positive sum of matrices
\(v_\alpha^*v_\alpha\).  For \(\mathcal W_M\), that coefficient is \(M\).
\end{proof}

\subsection{The finite supertransfer test}

Reflect the range-\(2R\) strip about the time plane and collect all
positive-side boundary generators into:
\begin{equation}
 \Xi_\psi,\qquad \Xi_\chi,\qquad Z_\phi.
 \label{eq:v5v66-boundary-variables}
\end{equation}
Let the reflected cross matrices of the three quadratic actions in
\eqref{eq:v5v66-supertransfer-integral} be
\begin{equation}
 M_\psi,\qquad M_\chi,\qquad M_\phi.
 \label{eq:v5v66-three-cross-matrices}
\end{equation}
The reflection identifications include spin chirality, link half
transport, and the determinant-line/Berezin orientation.

\begin{definition}[Pauli--Villars supertransfer certificate]
\label{def:v5v66-PVSTC}
A PVSTC consists of:
\begin{enumerate}
 \item an invertible gauge-covariant \(D_{\mathrm{PV}}\);
 \item the exact representation
 \eqref{eq:v5v66-supertransfer-integral};
 \item reflection covariance of all within-half actions and convergence
 of the commuting Gaussian;
 \item the finite matrix inequalities
 \begin{equation}
 M_\psi\ge0,\qquad M_\chi\ge0,\qquad M_\phi\ge0;
 \label{eq:v5v66-PVSTC-matrices}
 \end{equation}
 \item compatibility of these inequalities with Higgs-zero
 re-fibration, gauge averaging, and the V60 common refinements;
 \item bounds uniform in volume, lattice spacing, and fifth-direction
 wall separation.
\end{enumerate}
\end{definition}

\begin{theorem}[Finite PV supertransfer sufficiency]
\label{thm:v5v66-PVSTC-sufficient}
At a fixed finite cutoff, items 1--4 of a PVSTC imply reflection
positivity of the exact determinant-ratio functional
\eqref{eq:v5v66-supertransfer-integral}.
\end{theorem}

\begin{proof}
The V64 Grassmann cross-matrix theorem places the \(\psi\) and \(\chi\)
cross exponentials in their reflection cones.  Theorem
\ref{thm:v5v66-bosonic-cross} does the same for \(\phi\).  Their product
lies in the graded product cone; the reversed Grassmann orientation gives
the standard reflected square.  The within-half actions are reflections
of one another, and the positive matrix \(D_{\mathrm{PV}}^*D_{\mathrm{PV}}\)
makes the full commuting integral finite.  Apply the reflection
factorisation proof of V63 to the joint field algebra.
\end{proof}

\begin{firewall}[Global positivity of \(D^*D\) is not the cross test]
\label{fw:v5v66-global-vs-cross}
The operator \(D_{\mathrm{PV}}^*D_{\mathrm{PV}}\) is positive on the full
lattice, but its reflected cross matrix \(M_\phi\) is a different object.
Range-two paths can cross the reflection plane with signs and spin/link
contractions.  The inequality \(M_\phi\ge0\) must be calculated; it does
not follow from \(D_{\mathrm{PV}}^*D_{\mathrm{PV}}\ge0\).
\end{firewall}

\begin{firewall}[Three cones must use one reflection]
\label{fw:v5v66-one-reflection}
Separate changes of variables may make the three matrices in
\eqref{eq:v5v66-three-cross-matrices} look positive under different
reflections.  That is insufficient.  One antilinear involution, one gauge
half-link convention, and one boundary orientation must make all three
inequalities in \eqref{eq:v5v66-PVSTC-matrices} hold simultaneously.
\end{firewall}

\subsection{Relation to the inherited Schur identity}

\begin{proposition}[Equality with the overlap determinant]
\label{prop:v5v66-overlap-equality}
For the finite Shamir domain-wall matrices of inherited V77, take
\[
 D_{\mathrm{phys}}=D_{\mathrm{DW},N}(m_f),
 \qquad
 D_{\mathrm{PV}}=D_{\mathrm{DW},N}(1).
\]
Whenever \(D_{\mathrm{PV}}\) is invertible, the value of
\eqref{eq:v5v66-supertransfer-integral} is exactly
\begin{equation}
 \det D_N(m_f).
 \label{eq:v5v66-overlap-value}
\end{equation}
Thus a PVSTC would prove reflection positivity of the same finite
determinant ratio that converges to the overlap operator; it would not
define a different regulator.
\end{proposition}

\begin{proof}
Corollary~\ref{cor:v5v66-ratio-representation} evaluates the joint
integral as the domain-wall determinant ratio.  The inherited V77 finite
fifth-direction determinant identity evaluates that ratio as
\(\det D_N(m_f)\).
\end{proof}

\begin{assessment}[The exact remaining computation]
\label{ass:v5v66-remaining-computation}
Inherited V76 proves \(M_\psi\ge0\) for the complete physical two-wall
Grassmann block.  Reflection covariance gives a closely related test for
\(D_{\mathrm{PV}}^*\), but no inherited theorem checks the complete joint
orientation.  Most importantly, no note computes the range-two
pseudofermion matrix \(M_\phi\).  This finite matrix, not the determinant
identity, is the next pass-or-witness calculation.
\end{assessment}

\subsection{V66 status}

\begin{theorem}[V66 conclusion]
\label{thm:v5v66-status}
The inverse Pauli--Villars determinant has an exact finite local
representation using a conjugate Grassmann regulator and a positive
pseudofermion square.  It preserves the determinant phase, is gauge
covariant, and increases hopping range by at most a factor of two.  A
simultaneous nonnegative cross-matrix test for the physical Grassmann,
conjugate Grassmann, and pseudofermion fields is sufficient for
reflection positivity of the exact finite overlap ratio.  The naive
commuting-spinor integral can diverge on a negative wall-mass plateau, and
global positivity of \(D_{\mathrm{PV}}^*D_{\mathrm{PV}}\) does not prove
the range-two cross-matrix condition.  The latter remains the explicit
open calculation.
\end{theorem}

\begin{proof}
Use Proposition~\ref{prop:v5v66-negative-plateau},
Theorem~\ref{thm:v5v66-hermitianised-inverse},
Corollary~\ref{cor:v5v66-ratio-representation}, and
Theorem~\ref{thm:v5v66-PVSTC-sufficient}, with
Firewall~\ref{fw:v5v66-global-vs-cross}.
\end{proof}

V67 will compute the pseudofermion cross matrix for a nearest-neighbour
reflection-covariant Wilson operator in temporal gauge.  It will separate
one-step and two-step paths in \(D^*D\), diagonalise the boundary-layer
matrix, and either prove \(M_\phi\ge0\) in a declared Wilson-parameter
window or exhibit a negative principal minor.

\clearpage
\section{V67: the pseudofermion cross matrix at Wilson parameter one}
\label{sec:v5v67}

V66 reduced the finite determinant-ratio OS problem to the cross matrix of
\(D_{\mathrm{PV}}^*D_{\mathrm{PV}}\).  For a general Wilson parameter this
square has temporal range two.  At the standard temporal value \(r=1\),
the two-step coefficient vanishes exactly because the forward and
backward spin projectors are orthogonal.  The remaining cross matrix is
the Hermitian within-slice Wilson operator.  This gives an explicit
positive parameter window rather than another abstract certificate.

\subsection{Temporal block form}

Let the time spacing be \(a_t>0\), let
\(\gamma_0=\gamma_0^*\), \(\gamma_0^2=1\), and put
\begin{equation}
 P_\pm=\frac{1\pm\gamma_0}{2}.
 \label{eq:v5v67-temporal-projectors}
\end{equation}
In temporal gauge, a nearest-neighbour Wilson operator with temporal
Wilson parameter one can be written
\begin{equation}
 (D\psi)_t
 =
 A_t\psi_t
 C_+\psi_{t+1}
 C_-\psi_{t-1},
 \qquad
 C_+=-\frac1{a_t}P_-,
 \quad
 C_-=-\frac1{a_t}P_+.
 \label{eq:v5v67-temporal-D}
\end{equation}
All spatial, fifth-direction, mass, Higgs, and boundary-wall terms are
contained in the finite slice operator \(A_t\).

On a reflection-symmetric pair of slices adjacent to the reflection
plane, assume the reflection identification converts both slice
operators to one matrix
\begin{equation}
 A=H+K,
 \qquad
 H=H^*,\quad [H,\gamma_0]=0,
 \qquad
 K^*=-K,\quad \{K,\gamma_0\}=0.
 \label{eq:v5v67-HK-split}
\end{equation}
For the minimal Wilson domain-wall operator, \(H\) contains the temporal
Wilson diagonal, all non-temporal Wilson Laplacians, and the real wall
mass; \(K\) contains the non-temporal Clifford differences.

\begin{lemma}[Projector cancellation identities]
\label{lem:v5v67-projector-cancellation}
Under \eqref{eq:v5v67-HK-split},
\begin{align}
 C_-^*C_+&=0,
 &
 C_+^*C_-&=0,
 \label{eq:v5v67-two-step-zero}\\
 A^*P_-+P_+A&=H,
 &
 A^*P_++P_-A&=H.
 \label{eq:v5v67-nearest-identity}
\end{align}
\end{lemma}

\begin{proof}
The first line follows from \(P_+P_-=P_-P_+=0\).  Since \(H\) commutes
with \(\gamma_0\), it commutes with \(P_\pm\).  Since \(K\) anticommutes
with \(\gamma_0\),
\[
 KP_-=P_+K,\qquad KP_+=P_-K.
\]
Using \(A^*=H-K\),
\[
 A^*P_-+P_+A
 =
 HP_- -KP_-+P_+H+P_+K
 =H.
\]
The second identity is analogous.
\end{proof}

\subsection{The exact square across the reflection plane}

Put \(Q=D^*D\).  Away from temporal boundaries, its two-step forward
coefficient is
\begin{equation}
 Q_{t,t+2}=C_-^*C_+,
 \label{eq:v5v67-Q-two-step}
\end{equation}
and, for two reflection-adjacent slices with common \(A\), its one-step
coefficient is
\begin{equation}
 Q_{0,1}
 =
 A^*C_+ + C_-^*A.
 \label{eq:v5v67-Q-one-step}
\end{equation}

\begin{theorem}[Exact pseudofermion boundary matrix]
\label{thm:v5v67-pseudofermion-matrix}
Under \eqref{eq:v5v67-temporal-D}--\eqref{eq:v5v67-HK-split},
\begin{equation}
 Q_{0,2}=0,
 \qquad
 Q_{0,1}=-\frac1{a_t}H.
 \label{eq:v5v67-Q-boundary}
\end{equation}
Consequently the cross-plane part of the pseudofermion Boltzmann factor
\(\exp(-\phi^*Q\phi)\) is
\begin{equation}
 \exp\left\{
 \frac1{a_t}
 \left[
 (\Theta\phi)H\phi+\Theta\bigl((\Theta\phi)H\phi\bigr)
 \right]
 \right\},
 \label{eq:v5v67-pseudofermion-cross}
\end{equation}
with reflected cross matrix
\begin{equation}
 M_\phi=\frac1{a_t}H.
 \label{eq:v5v67-Mphi}
\end{equation}
\end{theorem}

\begin{proof}
Equation \eqref{eq:v5v67-two-step-zero} gives the first statement.
For the nearest coefficient,
\[
 A^*C_+ + C_-^*A
 =
 -\frac1{a_t}(A^*P_-+P_+A)
 =
 -\frac1{a_t}H
\]
by Lemma~\ref{lem:v5v67-projector-cancellation}.  The quadratic action
contains this term and its adjoint.  Changing from \(S\) to the
Boltzmann exponent \(-S\) reverses the sign, giving
\eqref{eq:v5v67-pseudofermion-cross} and
\eqref{eq:v5v67-Mphi}.
\end{proof}

\begin{corollary}[Necessary and sufficient pseudofermion cone test]
\label{cor:v5v67-pseudofermion-test}
Under the hypotheses of
Theorem~\ref{thm:v5v67-pseudofermion-matrix}, the pseudofermion
cross-plane exponential lies in the commuting reflection cone if and only
if
\begin{equation}
 H\ge0.
 \label{eq:v5v67-H-positive}
\end{equation}
\end{corollary}

\begin{proof}
Apply the commuting cross-matrix theorem of V66 to
\eqref{eq:v5v67-Mphi}.
\end{proof}

\begin{firewall}[Global square positivity is weaker]
\label{fw:v5v67-square-weaker}
The identity \(Q=D^*D\ge0\) holds without
\eqref{eq:v5v67-H-positive}.  If \(H\) has a negative eigenvalue, the
cross-plane cone fails even though the full pseudofermion Gaussian
converges.  V67 therefore supplies information not contained in global
positivity of the square.
\end{firewall}

\subsection{A Wilson mass subwindow}

For the minimal reflection-even domain-wall block, suppose
\begin{equation}
 H
 =
 \frac1{a_t}I+W_\perp+m(s),
 \qquad W_\perp\ge0,
 \label{eq:v5v67-minimal-H}
\end{equation}
where \(W_\perp\) is the sum of the spatial and fifth-direction Wilson
Laplacians and the real wall profile satisfies
\begin{equation}
 m(s)\ge-M.
 \label{eq:v5v67-mass-lower}
\end{equation}

\begin{proposition}[Explicit positive subwindow]
\label{prop:v5v67-positive-window}
If
\begin{equation}
 0<a_tM<1,
 \label{eq:v5v67-positive-window}
\end{equation}
then
\begin{equation}
 H\ge\left(\frac1{a_t}-M\right)I>0
 \label{eq:v5v67-H-gap}
\end{equation}
and the pseudofermion cross matrix has margin
\begin{equation}
 M_\phi
 \ge
 \frac{1-a_tM}{a_t^2}I.
 \label{eq:v5v67-Mphi-gap}
\end{equation}
\end{proposition}

\begin{proof}
Drop the nonnegative \(W_\perp\) in
\eqref{eq:v5v67-minimal-H} and use
\eqref{eq:v5v67-mass-lower}.  Divide by \(a_t\) for the second
inequality.
\end{proof}

The inherited one-species Wilson-degree window at temporal parameter one
is
\begin{equation}
 0<a_tM<2.
 \label{eq:v5v67-one-species}
\end{equation}

\begin{corollary}[Nonempty topological/OS overlap window]
\label{cor:v5v67-overlap-window}
For the minimal free block, the interval
\begin{equation}
 0<a_tM<1
 \label{eq:v5v67-overlap-window}
\end{equation}
simultaneously lies in the one-species Wilson phase and satisfies the
strict pseudofermion reflection-cone condition.
\end{corollary}

\begin{proof}
Combine \eqref{eq:v5v67-one-species} with
Proposition~\ref{prop:v5v67-positive-window}.
\end{proof}

\begin{assessment}[Anisotropic spacings]
\label{ass:v5v67-anisotropic}
Only the temporal Wilson diagonal \(a_t^{-1}\) pays the cross-plane
margin.  Non-temporal Wilson terms are nonnegative and may improve it.
Thus the relevant dimensionless wall mass is \(a_tM\), not a value formed
from an unrelated spatial spacing.
\end{assessment}

\subsection{Reflection-even perturbations}

Let \(H=H_0+V\), where \(H_0\) satisfies
\begin{equation}
 H_0\ge h_0I,\qquad h_0>0,
 \label{eq:v5v67-H0}
\end{equation}
and \(V=V^*\) commutes with \(\gamma_0\).

\begin{lemma}[Stability of the cross margin]
\label{lem:v5v67-margin-stability}
If \(\|V\|<h_0\), then
\begin{equation}
 M_\phi
 \ge
 \frac{h_0-\|V\|}{a_t}I>0.
 \label{eq:v5v67-margin-stability}
\end{equation}
\end{lemma}

\begin{proof}
For every \(v\),
\[
 \langle v,Hv\rangle
 \ge(h_0-\|V\|)\|v\|^2.
\]
Use \eqref{eq:v5v67-Mphi}.
\end{proof}

\begin{corollary}[Small reflection-even gauge/Higgs patch]
\label{cor:v5v67-small-patch}
If the gauge and Higgs perturbations preserve the structural split
\eqref{eq:v5v67-HK-split} and change its Hermitian part by operator norm
less than the free margin in \eqref{eq:v5v67-H-gap}, then the
pseudofermion cross-plane cone persists on that patch.
\end{corollary}

\begin{proof}
Apply Lemma~\ref{lem:v5v67-margin-stability}.
\end{proof}

\begin{firewall}[Plaquette admissibility is not this norm bound]
\label{fw:v5v67-plaquette-not-H}
A small plaquette controls curvature but need not place every link or
Higgs background in the operator-norm ball of
Corollary~\ref{cor:v5v67-small-patch}; torons and topology can remain.
One needs a gauge-covariant estimate for the slice operator \(H\), or a
finite atlas of patches carrying the same positive margin.
\end{firewall}

\subsection{The conjugate Grassmann regulator}

The V66 representation also contains
\(\bar\chi D_{\mathrm{PV}}^*\chi\).

\begin{lemma}[Adjoint Wilson block retains the Grassmann cone]
\label{lem:v5v67-adjoint-Grassmann}
If the nearest-neighbour Wilson action of \(D\) has the reflected
Grassmann cone with cross coefficient \(\kappa P_\pm\),
\(\kappa\ge0\), then the action of \(D^*\), with positive and negative
reflection generators interchanged, has the same nonnegative coefficient
list.
\end{lemma}

\begin{proof}
Taking the adjoint exchanges forward and backward hopping and exchanges
\(P_+\) with \(P_-\).  The reflection involution already makes the same
exchange.  Relabel the orthonormal bases of the two projector ranges in
the inherited V76 cone expansion.  The scalar coefficient \(\kappa\)
is unchanged.
\end{proof}

\begin{theorem}[Finite minimal PV supertransfer window]
\label{thm:v5v67-finite-PVST}
Assume:
\begin{enumerate}
 \item the finite physical and Pauli--Villars domain-wall operators have
 the reflection-covariant nearest-neighbour form
 \eqref{eq:v5v67-temporal-D};
 \item their common slice operator satisfies
 \eqref{eq:v5v67-HK-split} and \(H>0\);
 \item the Pauli--Villars operator is invertible;
 \item all within-half terms and Berezin orientations obey the inherited
 V76 reflection convention.
\end{enumerate}
Then the exact joint representation
\eqref{eq:v5v66-supertransfer-integral} is reflection positive at that
finite cutoff.  In the minimal free wall model,
\(0<a_tM<1\) is a sufficient nonempty one-species window.
\end{theorem}

\begin{proof}
The inherited V76 cone gives \(M_\psi\ge0\).
Lemma~\ref{lem:v5v67-adjoint-Grassmann} gives \(M_\chi\ge0\).
Theorem~\ref{thm:v5v67-pseudofermion-matrix} and \(H>0\) give
\(M_\phi>0\).  These are the three cross inequalities in the V66
Pauli--Villars supertransfer certificate, so its sufficiency theorem
proves reflection positivity.  The free parameter statement is
Corollary~\ref{cor:v5v67-overlap-window}.
\end{proof}

\begin{firewall}[Boundary-wall and Yukawa terms must enter \(H\) honestly]
\label{fw:v5v67-extra-slice-terms}
The Shamir boundary mass, flavour mixing, and Higgs--Yukawa matrices may
add Hermitian slice terms not displayed in
\eqref{eq:v5v67-minimal-H}.  Theorem~\ref{thm:v5v67-finite-PVST}
continues to apply only after their contribution to
\(\lambda_{\min}(H)\) is included.  The scalar bound
\(0<a_tM<1\) cannot be quoted if those terms can lower the slice
spectrum.
\end{firewall}

\begin{firewall}[Finite window is not a cofinal theorem]
\label{fw:v5v67-not-cofinal}
Theorem~\ref{thm:v5v67-finite-PVST} is at one finite cutoff and on a
reflection-even patch.  A continuum result needs a cutoff-independent
lower margin, compatibility with all gauge/topological patches, control
as wall separation grows, and convergence of every reflected Standard
Model cylinder matrix.
\end{firewall}

\subsection{V67 status}

\begin{theorem}[V67 conclusion]
\label{thm:v5v67-status}
For temporal Wilson parameter one, the range-two pseudofermion crossings
in \(D^*D\) vanish exactly and the remaining reflected cross matrix is
\(H/a_t\), where \(H\) is the Hermitian reflection-even slice operator.
Thus \(H\ge0\) is the exact commuting-field cone test.  In the minimal
free domain-wall model, \(0<a_tM<1\) is a nonempty subwindow of the
one-species phase in which the exact phase-preserving Pauli--Villars
supertransfer representation is reflection positive.  The result extends
to reflection-even perturbations smaller than the cross margin, but no
uniform atlas or cutoff-removal estimate is yet proved.
\end{theorem}

\begin{proof}
Use Theorem~\ref{thm:v5v67-pseudofermion-matrix},
Proposition~\ref{prop:v5v67-positive-window},
Lemma~\ref{lem:v5v67-adjoint-Grassmann}, and
Theorem~\ref{thm:v5v67-finite-PVST}.
\end{proof}

V68 will attempt the uniform patch extension.  It will express
\(\lambda_{\min}(H)\) in terms of gauge-covariant non-temporal Wilson
forms, wall/boundary mass matrices, and Higgs--Yukawa singular values.
Where a uniform lower bound fails, it will re-fibre the offending slice
band into the V60 retained sector and test the resulting positive Schur
margin rather than discard the background.

\clearpage
\section{V68: gauge-uniform slice margins and Yukawa polar strata}
\label{sec:v5v68}

V67 expressed the pseudofermion reflection matrix as \(M_\phi=H/a_t\).
The remaining issue is to estimate the Hermitian slice operator \(H\)
without imposing an artificial small-link gauge.  For the minimal Wilson
operator, every non-temporal gauge link occurs in a nonnegative covariant
Laplacian.  Hence the lower bound is uniform over all link configurations,
including torons and nontrivial representatives.

Yukawa terms require a more careful statement.  On a massive Higgs
stratum, polar decomposition permits the reflected right-handed frame to
be aligned with the left-handed mass map, making its cross contribution
the nonnegative operator \(|Y|\).  At a Yukawa kernel or Higgs zero, that
polar frame ceases to be invertible; the kernel must remain in the
retained Grassmann sector.  Only Hermitian slice terms not absorbed by
these positive structures debit the Wilson margin.

\subsection{Gauge-covariant Wilson forms}

Let \(T_j^+\) be a unitary covariant shift in a non-temporal direction and
\(T_j^-=(T_j^+)^*\).  Define
\begin{equation}
 \Delta_j
 :=
 \frac{2-T_j^+-T_j^-}{a_j^2}.
 \label{eq:v5v68-covariant-Laplacian}
\end{equation}

\begin{lemma}[Link-uniform positivity]
\label{lem:v5v68-link-positive}
For every gauge-link configuration and every vector \(u\),
\begin{equation}
 \langle u,\Delta_ju\rangle
 =
 \frac1{a_j^2}\|(1-T_j^+)u\|^2
 \ge0.
 \label{eq:v5v68-link-positive}
\end{equation}
The estimate is independent of volume, plaquette size, toron holonomy,
and gauge representative.
\end{lemma}

\begin{proof}
Unitarity gives
\[
 (1-T_j^+)^*(1-T_j^+)
 =
 (1-T_j^-)(1-T_j^+)
 =
 2-T_j^+-T_j^-.
\]
Take the quadratic form.  No curvature or gauge choice enters.
\end{proof}

Let
\begin{equation}
 W_\perp
 :=
 \sum_{j\ne0}\frac{r_ja_j}{2}\Delta_j,
 \qquad r_j\ge0.
 \label{eq:v5v68-Wperp}
\end{equation}

\begin{corollary}[Gauge-uniform Wilson lower bound]
\label{cor:v5v68-Wperp}
For all unitary non-temporal links,
\begin{equation}
 W_\perp\ge0.
 \label{eq:v5v68-Wperp-positive}
\end{equation}
\end{corollary}

\begin{proof}
Sum Lemma~\ref{lem:v5v68-link-positive} with nonnegative coefficients.
\end{proof}

\begin{assessment}[Where admissibility is still needed]
\label{ass:v5v68-admissibility}
The lower bound \eqref{eq:v5v68-Wperp-positive} does not require small
plaquettes.  Plaquette admissibility remains relevant to locality of the
overlap sign function, heat-kernel comparison, topology of the chiral
projector, and continuum approximation.  It is not needed for the bare
pseudofermion cross sign.
\end{assessment}

\subsection{A general slice ledger}

Write the reflection-even Hermitian slice operator as
\begin{equation}
 H
 =
 \frac1{a_t}I
 W_\perp
 +M_{\mathrm{wall}}
 +M_{\mathrm{Y}}
 +B_{\mathrm{bd}}
 +B_{\mathrm{imp}},
 \label{eq:v5v68-H-ledger}
\end{equation}
where:
\begin{enumerate}
 \item \(M_{\mathrm{wall}}=M_{\mathrm{wall}}^*\) is the wall profile;
 \item \(M_{\mathrm{Y}}=M_{\mathrm{Y}}^*\) is the Yukawa contribution
 after the declared reflection-frame identification;
 \item \(B_{\mathrm{bd}}\) contains Shamir boundary-wall couplings not
 already included in \(M_{\mathrm{wall}}\);
 \item \(B_{\mathrm{imp}}\) contains clover, improvement, or other
 Hermitian slice terms.
\end{enumerate}
Assume
\begin{align}
 M_{\mathrm{wall}}&\ge-MI,
 \label{eq:v5v68-wall-lower}\\
 B_{\mathrm{bd}}+B_{\mathrm{imp}}&\ge-bI.
 \label{eq:v5v68-debit-lower}
\end{align}

\begin{theorem}[Gauge-uniform slice margin]
\label{thm:v5v68-slice-margin}
If \(M_{\mathrm{Y}}\ge0\), then for every non-temporal gauge-link
configuration
\begin{align}
 H
 &\ge
 \left(\frac1{a_t}-M-b\right)I,
 \label{eq:v5v68-H-lower}\\
 M_\phi=\frac{H}{a_t}
 &\ge
 \frac{1-a_t(M+b)}{a_t^2}I.
 \label{eq:v5v68-Mphi-lower}
\end{align}
In particular, for a fixed reserve \(0<\kappa<1\), the condition
\begin{equation}
 a_t(M+b)\le1-\kappa
 \label{eq:v5v68-reserve}
\end{equation}
gives the uniform cross margin
\begin{equation}
 M_\phi\ge\frac{\kappa}{a_t^2}I.
 \label{eq:v5v68-cross-margin}
\end{equation}
\end{theorem}

\begin{proof}
Use Corollary~\ref{cor:v5v68-Wperp}, the assumed nonnegativity of
\(M_{\mathrm{Y}}\), and
\eqref{eq:v5v68-wall-lower}--\eqref{eq:v5v68-debit-lower} in
\eqref{eq:v5v68-H-ledger}.  Divide by \(a_t\) and insert
\eqref{eq:v5v68-reserve}.
\end{proof}

\begin{corollary}[Volume and wall-separation independence]
\label{cor:v5v68-volume-wall}
If the constants \(M,b,\kappa\) in
\eqref{eq:v5v68-reserve} are independent of spatial volume and
fifth-direction wall separation, then the cross margin
\eqref{eq:v5v68-cross-margin} is independent of those quantities.
\end{corollary}

\begin{proof}
No other constant occurs in
\eqref{eq:v5v68-cross-margin}.
\end{proof}

\subsection{Polar alignment of a Yukawa map}

Let \(Y:R\to L\) be a finite-dimensional Yukawa mass map between
right- and left-handed retained fibres.  Its polar decomposition is
\begin{equation}
 Y=U|Y|,
 \qquad
 |Y|=(Y^*Y)^{1/2},
 \label{eq:v5v68-polar-Yukawa}
\end{equation}
where \(U\) is a partial isometry with initial space
\((\ker Y)^\perp\) and final space \(\overline{\operatorname{ran}Y}\).

\begin{lemma}[Positive massive Yukawa frame]
\label{lem:v5v68-positive-Yukawa}
On a stratum where \(Y\) is square and
\begin{equation}
 s_{\min}(Y)\ge y_0>0,
 \label{eq:v5v68-Yukawa-gap}
\end{equation}
the polar factor \(U\) is unitary.  Identifying the reflected left fibre
with the right fibre by \(U^*\) changes the Yukawa cross coefficient to
\begin{equation}
 U^*Y=|Y|\ge y_0I.
 \label{eq:v5v68-Yukawa-positive}
\end{equation}
This unitary change of reflection frame preserves the Grassmann and
commuting reflection cones.
\end{lemma}

\begin{proof}
The singular-value lower bound makes \(Y\) invertible, so its polar
partial isometry is unitary.  Equation
\eqref{eq:v5v68-Yukawa-positive} is the polar identity.  Reflection-cone
coefficient matrices transform by unitary congruence, which preserves
positive semidefiniteness.
\end{proof}

\begin{proposition}[Kernel-refined polar decomposition]
\label{prop:v5v68-kernel-refinement}
For an arbitrary square or rectangular \(Y\), there are orthogonal
splittings
\begin{equation}
 R=\ker Y\oplus R_m,
 \qquad
 L=\ker Y^*\oplus L_m
 \label{eq:v5v68-Yukawa-splitting}
\end{equation}
such that
\[
 Y|_{R_m}:R_m\longrightarrow L_m
\]
has a unitary polar factor and positive modulus on every singular-value
cell separated from zero.  The kernel and cokernel dimensions are exact
retained data and are not included in the massive Yukawa shell.
\end{proposition}

\begin{proof}
Take \(R_m=(\ker Y)^\perp\) and
\(L_m=\overline{\operatorname{ran}Y}=(\ker Y^*)^\perp\).
The polar partial isometry restricts to a unitary
\(R_m\to L_m\).  On any finite-cutoff cell whose nonzero singular values
are bounded below, the restricted modulus is strictly positive.  The two
kernels are precisely the cohomological source and target defects and
must remain in the retained determinant line.
\end{proof}

\begin{corollary}[Higgs radial scaling]
\label{cor:v5v68-radial-Yukawa}
If \(Y(\rho)=\rho Y_0\), then on \(\rho\ge r>0\),
\begin{equation}
 |Y(\rho)|=\rho|Y_0|,
 \qquad
 s_{\min}(Y(\rho)|_{R_m})
 \ge r\,s_{\min}(Y_0|_{R_m}).
 \label{eq:v5v68-radial-Yukawa}
\end{equation}
At \(\rho=0\), the massive subspace joins the retained Higgs-zero
crossover stratum.
\end{corollary}

\begin{proof}
Functional calculus gives
\[
 |Y(\rho)|
 =
 (\rho^2Y_0^*Y_0)^{1/2}
 =
 \rho|Y_0|
\]
for \(\rho\ge0\).  The singular-value estimate follows.  At zero the
map vanishes.
\end{proof}

\begin{firewall}[Polar frames are local atlas data]
\label{fw:v5v68-polar-local}
The unitary \(U\) can wind over gauge--Higgs configuration space and is
undefined when singular values cross zero.  Lemma
\ref{lem:v5v68-positive-Yukawa} is a cellwise reflection-frame
construction.  Its overlap determinants must use the inherited
determinant line, and its zero crossings must use the V60 common
refinement.
\end{firewall}

\subsection{Unaligned Yukawa debit}

If a chosen reflection convention cannot absorb the polar factor, retain
the Hermitian Yukawa contribution as an explicit debit.  Set
\begin{equation}
 y_-(\Phi)
 :=
 \max\{0,-\lambda_{\min}(M_{\mathrm{Y}}(\Phi))\}.
 \label{eq:v5v68-Yukawa-debit}
\end{equation}

\begin{proposition}[General unaligned margin]
\label{prop:v5v68-unaligned-margin}
Without assuming \(M_{\mathrm{Y}}\ge0\),
\begin{equation}
 M_\phi
 \ge
 \frac{1-a_t(M+b+y_-(\Phi))}{a_t^2}I.
 \label{eq:v5v68-unaligned-margin}
\end{equation}
Thus the exact safe cell with reserve \(\kappa\) is
\begin{equation}
 a_t(M+b+y_-(\Phi))\le1-\kappa.
 \label{eq:v5v68-unaligned-cell}
\end{equation}
\end{proposition}

\begin{proof}
By definition,
\(M_{\mathrm{Y}}\ge-y_-(\Phi)I\).  Repeat the proof of
Theorem~\ref{thm:v5v68-slice-margin}.
\end{proof}

\begin{assessment}[Why a norm bound is only a fallback]
\label{ass:v5v68-polar-vs-norm}
The crude estimate \(y_-\le\|M_{\mathrm{Y}}\|\) can incorrectly make a
large positive fermion mass look dangerous.  Polar alignment uses the
sign and proves nonnegativity on the massive stratum.  The norm debit is
appropriate only for terms whose reflection orientation has not been
resolved.
\end{assessment}

\subsection{Spectral re-fibration of a failed slice margin}

For a fixed cell, let
\begin{equation}
 P_{\mathrm{bad}}
 :=
 1_{(-\infty,h_*)}(H),
 \qquad h_*>0,
 \label{eq:v5v68-bad-projector}
\end{equation}
and \(P_{\mathrm{safe}}=1-P_{\mathrm{bad}}\).

\begin{proposition}[Positive conditional pseudofermion shell]
\label{prop:v5v68-positive-conditional-shell}
On the \(P_{\mathrm{safe}}\) sector,
\begin{equation}
 P_{\mathrm{safe}}M_\phi P_{\mathrm{safe}}
 \ge\frac{h_*}{a_t}P_{\mathrm{safe}}.
 \label{eq:v5v68-safe-shell}
\end{equation}
Because \(P_{\mathrm{bad}}\) and \(P_{\mathrm{safe}}\) are spectral
projectors of \(H\), the cross matrix has no mixed block between them.
Thus the safe pseudofermion shell has a reflection-cone factorisation
conditional on retaining the bad band.
\end{proposition}

\begin{proof}
Spectral calculus gives
\(P_{\mathrm{safe}}HP_{\mathrm{safe}}\ge h_*P_{\mathrm{safe}}\)
and
\(P_{\mathrm{bad}}HP_{\mathrm{safe}}=0\).  Divide by \(a_t\) and apply
the commuting cross-matrix criterion of V66 to the safe block.
\end{proof}

\begin{firewall}[Conditional positivity is not full OS positivity]
\label{fw:v5v68-conditional-only}
Moving \(P_{\mathrm{bad}}\) to the retained sector prevents an invalid
Gaussian shell integration, but it does not make the retained negative
cross block positive.  A full OS theorem must eventually give that band a
positive transfer representation, show that its physical contribution
vanishes, or exclude it by a proved uniform margin.  Re-fibration is
exact bookkeeping, not a sign cure.
\end{firewall}

\begin{proposition}[Projector stability away from a crossing]
\label{prop:v5v68-projector-stability}
Suppose the spectrum of \(H(x)\) stays a distance \(\gamma>0\) from
\(h_*\) on a parameter cell and
\(\|\partial H\|\le L\).  Then
\begin{equation}
 \|\partial P_{\mathrm{bad}}\|
 \le C_\Gamma\frac{L}{\gamma^2}.
 \label{eq:v5v68-projector-stability}
\end{equation}
At a threshold crossing, the adjacent cells glue through their common
refinement with the crossing band retained.
\end{proposition}

\begin{proof}
Use the Riesz formula for
\(1_{(-\infty,h_*)}(H)\) on a contour in the spectral gap and
differentiate.  The two resolvents each cost at most \(\gamma^{-1}\).
The crossing statement is the V59/V60 re-fibration rule.
\end{proof}

\subsection{Uniform slice-margin certificate}

\begin{definition}[Uniform slice-margin certificate]
\label{def:v5v68-USMC}
A USMC consists of:
\begin{enumerate}
 \item temporal Wilson parameter one and the structural split of V67;
 \item unitary non-temporal links and nonnegative Wilson parameters;
 \item bounds \eqref{eq:v5v68-wall-lower} and
 \eqref{eq:v5v68-debit-lower};
 \item cellwise polar Yukawa frames, with kernels retained;
 \item a cutoff-independent reserve \(\kappa>0\) in
 \eqref{eq:v5v68-reserve}, or an explicit retained bad band;
 \item overlap maps compatible with the determinant line and V60
 refinement atlas;
 \item a separate full OS disposition for every retained bad band.
\end{enumerate}
\end{definition}

\begin{theorem}[V68 conclusion]
\label{thm:v5v68-status}
For the minimal temporal Wilson block, non-temporal gauge fields do not
debit the pseudofermion reflection margin: their covariant Wilson
Laplacians are nonnegative for every unitary link configuration.  On
massive Higgs strata, polar alignment turns an invertible Yukawa map into
its nonnegative modulus while retaining all kernels and crossings.
Consequently the exact margin is
\(\kappa/a_t^2\) under
\eqref{eq:v5v68-reserve}, uniformly in volume, wall separation, torons,
and gauge representative when the printed slice debits are uniform.
Failed spectral bands admit an exact positive conditional shell, but
their retained cross sign remains an open OS obligation.
\end{theorem}

\begin{proof}
Use Lemma~\ref{lem:v5v68-link-positive},
Theorem~\ref{thm:v5v68-slice-margin},
Lemma~\ref{lem:v5v68-positive-Yukawa}, and
Proposition~\ref{prop:v5v68-positive-conditional-shell}, with
Firewall~\ref{fw:v5v68-conditional-only}.
\end{proof}

V69 will address cofinal reflected convergence.  It will prove a
quantitative stability theorem for finite reflected cylinder matrices
under changes of the physical, conjugate-regulator, and pseudofermion
cross kernels, using the uniform reserve from V68.  This will separate
entrywise convergence, which preserves OS positivity, from determinant
or operator-norm convergence that does not control the relevant
cylinder observables.

\clearpage
\section{V69: quantitative convergence of reflected supertransfer matrices}
\label{sec:v5v69}

Finite reflection positivity is useful for the continuum only when the
same physical cylinder observables have convergent reflected matrix
elements.  Determinant convergence, convergence of an unreflected
quadratic form, or a pointwise limit of cross matrices is not by itself
that statement.  V69 supplies a quantitative stability estimate for the
joint physical-fermion, conjugate-regulator, and pseudofermion functional
of V66--V68.

The inherited notes already record that positive semidefinite matrices
are closed under entrywise limits.  The new step here is to give norms
which imply those entrywise limits and to display their dimension costs.
The pseudofermion estimate shows that operator-norm convergence alone can
lose a factor \(\sqrt{\operatorname{rank}}\); a Hilbert--Schmidt or
localized trace debit is required in a cofinal limit.

\subsection{An abstract normalized superfunctional estimate}

Let \((\mathfrak B,\|\cdot\|)\) be a finite-dimensional unital Banach
superalgebra with
\begin{equation}
 \|FG\|\le\|F\|\,\|G\|,
 \qquad
 \|\Theta F\|=\|F\|.
 \label{eq:v5v69-Banach-superalgebra}
\end{equation}
The V58 weighted Grassmann coefficient norm tensored with a bosonic
\(L^1\) norm is an example.  Let
\(\tau:\mathfrak B\to\mathbb C\) be the unnormalised integration
functional, with
\begin{equation}
 |\tau(F)|\le C_\tau\|F\|.
 \label{eq:v5v69-tau-bound}
\end{equation}
For a weight \(W\), define
\begin{equation}
 Z_W:=\tau(W),
 \qquad
 \mathcal S_W(F):=\frac{\tau(WF)}{Z_W}.
 \label{eq:v5v69-normalized-functional}
\end{equation}

\begin{lemma}[Normalized weight stability]
\label{lem:v5v69-weight-stability}
Suppose
\begin{equation}
 |Z_W|,|Z_{\widetilde W}|\ge z>0,
 \qquad
 \|W\|,\|\widetilde W\|\le M.
 \label{eq:v5v69-normalizer-reserve}
\end{equation}
Then
\begin{equation}
 |\mathcal S_W(F)-\mathcal S_{\widetilde W}(F)|
 \le
 \left(
 \frac{C_\tau}{z}
 \frac{C_\tau^2M}{z^2}
 \right)
 \|F\|\,\|W-\widetilde W\|.
 \label{eq:v5v69-weight-stability}
\end{equation}
\end{lemma}

\begin{proof}
Write
\[
 \frac{\tau(WF)}{Z_W}
 -
 \frac{\tau(\widetilde WF)}{Z_{\widetilde W}}
 =
 \frac{\tau((W-\widetilde W)F)}{Z_W}
 +
 \tau(\widetilde WF)
 \frac{Z_{\widetilde W}-Z_W}{Z_WZ_{\widetilde W}}.
\]
Use \eqref{eq:v5v69-Banach-superalgebra},
\eqref{eq:v5v69-tau-bound},
\[
 |Z_W-Z_{\widetilde W}|
 \le C_\tau\|W-\widetilde W\|,
\]
and \(\|\widetilde W\|\le M\).
\end{proof}

\begin{corollary}[Reflected entry stability]
\label{cor:v5v69-reflected-entry}
If \(\|F_i\|\le R\), then
\begin{equation}
 \left|
 \mathcal S_W((\Theta F_i)F_j)
 -
 \mathcal S_{\widetilde W}((\Theta F_i)F_j)
 \right|
 \le C_{z,M,\tau}R^2\|W-\widetilde W\|,
 \label{eq:v5v69-entry-bound}
\end{equation}
where
\[
 C_{z,M,\tau}
 =
 \frac{C_\tau}{z}+\frac{C_\tau^2M}{z^2}.
\]
\end{corollary}

\begin{proof}
By \eqref{eq:v5v69-Banach-superalgebra},
\(\|(\Theta F_i)F_j\|\le R^2\).  Apply
Lemma~\ref{lem:v5v69-weight-stability}.
\end{proof}

\subsection{Finite reflected eigenvalue stability}

For \(F_1,\ldots,F_N\), let
\begin{equation}
 G_W=(\mathcal S_W((\Theta F_i)F_j))_{i,j=1}^N.
 \label{eq:v5v69-Gram}
\end{equation}

\begin{lemma}[Gram eigenvalue perturbation]
\label{lem:v5v69-Gram-perturbation}
If
\[
 \max_{i,j}|(G_W-G_{\widetilde W})_{ij}|\le\varepsilon,
\]
then
\begin{equation}
 \|G_W-G_{\widetilde W}\|_{\mathrm{op}}
 \le N\varepsilon
 \label{eq:v5v69-Gram-op}
\end{equation}
and
\begin{equation}
 \lambda_{\min}(G_{\widetilde W})
 \ge
 \lambda_{\min}(G_W)-N\varepsilon.
 \label{eq:v5v69-Gram-eigen}
\end{equation}
\end{lemma}

\begin{proof}
The maximum row sum and maximum column sum are at most \(N\varepsilon\),
so the spectral norm is at most \(N\varepsilon\).  Weyl's eigenvalue
inequality gives \eqref{eq:v5v69-Gram-eigen}.
\end{proof}

\begin{corollary}[Strict finite OS reserve]
\label{cor:v5v69-strict-OS}
If \(G_W\ge\delta I_N\) and the right side of
\eqref{eq:v5v69-entry-bound} is less than \(\delta/N\), then
\(G_{\widetilde W}>0\).
\end{corollary}

\begin{proof}
Use Lemma~\ref{lem:v5v69-Gram-perturbation}.
\end{proof}

\subsection{Grassmann exponential differences}

\begin{lemma}[Banach-algebra exponential bound]
\label{lem:v5v69-exponential-bound}
For \(A,B\) in a unital Banach algebra,
\begin{equation}
 \|e^A-e^B\|
 \le
 e^{\max(\|A\|,\|B\|)}\|A-B\|.
 \label{eq:v5v69-exponential-bound}
\end{equation}
\end{lemma}

\begin{proof}
The Duhamel identity is
\[
 e^A-e^B
 =
 \int_0^1
 e^{(1-s)A}(A-B)e^{sB}\,ds.
\]
Submultiplicativity gives an integrand norm at most
\[
 e^{(1-s)\|A\|+s\|B\|}\|A-B\|
 \le e^{\max(\|A\|,\|B\|)}\|A-B\|.
\]
\end{proof}

Let
\[
 A_D=-\bar\psi D\psi
\]
be a finite Grassmann quadratic form, measured in the V58 weighted
coefficient norm.

\begin{corollary}[Grassmann Dirac stability]
\label{cor:v5v69-Grassmann-stability}
If the quadratic-form map obeys
\begin{equation}
 \|A_D-A_{\widetilde D}\|
 \le c_r\|D-\widetilde D\|_{\mathrm{mat}},
 \label{eq:v5v69-quadratic-map}
\end{equation}
then
\begin{equation}
 \|e^{A_D}-e^{A_{\widetilde D}}\|
 \le
 c_r e^{\max(\|A_D\|,\|A_{\widetilde D}\|)}
 \|D-\widetilde D\|_{\mathrm{mat}}.
 \label{eq:v5v69-Grassmann-stability}
\end{equation}
\end{corollary}

\begin{proof}
Apply Lemma~\ref{lem:v5v69-exponential-bound}.
\end{proof}

For the physical and conjugate-regulator Grassmann weights, a telescoping
product therefore reduces their contribution to the two corresponding
Dirac-matrix differences.  The exponential prefactor must be controlled
in a rooted or local norm; a raw global coefficient norm can grow with
the number of lattice generators.

\subsection{Normalized pseudofermion stability}

For a positive Hermitian \(n\times n\) matrix \(Q\), let
\begin{equation}
 d\gamma_Q(\phi)
 =
 \det Q\,e^{-\phi^*Q\phi}
 \frac{d^{2n}\phi}{\pi^n}.
 \label{eq:v5v69-complex-Gaussian}
\end{equation}
This is a probability measure with covariance \(Q^{-1}\).

\begin{lemma}[Relative Gaussian entropy]
\label{lem:v5v69-Gaussian-entropy}
For \(Q,\widetilde Q>0\), set
\begin{equation}
 X=Q^{-1/2}(\widetilde Q-Q)Q^{-1/2}.
 \label{eq:v5v69-relative-X}
\end{equation}
If \(\|X\|_{\mathrm{op}}\le\eta<1\), then
\begin{align}
 D_{\mathrm{KL}}(\gamma_Q\|\gamma_{\widetilde Q})
 &=
 \operatorname{tr}\bigl(X-\log(1+X)\bigr)
 \label{eq:v5v69-KL-exact}\\
 &\le
 \frac1{2(1-\eta)}
 \|X\|_{\mathrm{HS}}^2.
 \label{eq:v5v69-KL-bound}
\end{align}
\end{lemma}

\begin{proof}
Direct Gaussian integration gives
\[
 D_{\mathrm{KL}}(\gamma_Q\|\gamma_{\widetilde Q})
 =
 \operatorname{tr}(\widetilde QQ^{-1}-I)
 -
 \log\det(\widetilde QQ^{-1}).
\]
Similarity with \(1+X\) gives
\eqref{eq:v5v69-KL-exact}.  For every eigenvalue
\(x\in[-\eta,\eta]\),
\[
 x-\log(1+x)
 =
 \int_0^x\frac{t}{1+t}\,dt
 \le\frac{x^2}{2(1-\eta)}.
\]
Sum over the eigenvalues.
\end{proof}

\begin{corollary}[Pseudofermion total-variation bound]
\label{cor:v5v69-pseudofermion-TV}
Under the hypotheses of
Lemma~\ref{lem:v5v69-Gaussian-entropy},
\begin{equation}
 \|\gamma_Q-\gamma_{\widetilde Q}\|_{\mathrm{TV}}
 \le
 \frac{\|X\|_{\mathrm{HS}}}
 {2\sqrt{1-\eta}}.
 \label{eq:v5v69-pseudofermion-TV}
\end{equation}
\end{corollary}

\begin{proof}
Pinsker's inequality gives
\[
 \|\gamma_Q-\gamma_{\widetilde Q}\|_{\mathrm{TV}}
 \le\sqrt{\frac12D_{\mathrm{KL}}
 (\gamma_Q\|\gamma_{\widetilde Q})}.
\]
Insert \eqref{eq:v5v69-KL-bound}.
\end{proof}

\begin{firewall}[Operator norm alone pays rank]
\label{fw:v5v69-operator-rank}
The estimate
\[
 \|X\|_{\mathrm{HS}}
 \le\sqrt n\,\|X\|_{\mathrm{op}}
\]
shows that an operator-norm error tending to zero need not give a
cutoff-uniform pseudofermion measure error when the rank grows.  A
Hilbert--Schmidt, trace-local, or rooted summability estimate is required.
The V68 cross margin controls positivity but does not by itself pay this
rank.
\end{firewall}

\subsection{A combined finite stability theorem}

Let \(W_q\) denote the joint V66 supertransfer weight at cutoff \(q\),
expressed over a common refinement of the bosonic and Grassmann variables.
Let its three quadratic matrices be
\[
 D_{\mathrm{phys},q},\qquad
 D_{\mathrm{PV},q}^*,\qquad
 Q_q=D_{\mathrm{PV},q}^*D_{\mathrm{PV},q}.
\]

\begin{theorem}[Finite reflected supertransfer stability]
\label{thm:v5v69-supertransfer-stability}
Fix \(N\) positive-time observables \(F_1,\ldots,F_N\) of norm at most
\(R\).  Suppose two common-refinement cutoffs \(q,\widetilde q\) satisfy:
\begin{enumerate}
 \item the normalizer and weight reserves
 \eqref{eq:v5v69-normalizer-reserve};
 \item the Grassmann quadratic norms are bounded by \(L\);
 \item their physical and conjugate-regulator Dirac differences have
 matrix-norm sum at most \(\delta_F\);
 \item for the pseudofermion matrices,
 \(\|X\|_{\mathrm{op}}\le\eta<1\) and
 \(\|X\|_{\mathrm{HS}}\le\delta_B\);
 \item all remaining bosonic factors differ in total variation by at most
 \(\delta_0\).
\end{enumerate}
Then every reflected entry differs by at most
\begin{equation}
 C R^2
 \left(e^L\delta_F+\delta_B+\delta_0\right),
 \label{eq:v5v69-combined-entry}
\end{equation}
where \(C\) depends on the displayed normalizer reserves, the fixed
Grassmann norm map, and \(\eta\), but not on the chosen \(F_i\).
Consequently
\begin{equation}
 \lambda_{\min}(G_{\widetilde q})
 \ge
 \lambda_{\min}(G_q)
 -
 NC R^2
 \left(e^L\delta_F+\delta_B+\delta_0\right).
 \label{eq:v5v69-combined-eigen}
\end{equation}
\end{theorem}

\begin{proof}
Use a telescoping product for the two Grassmann exponentials and apply
Corollary~\ref{cor:v5v69-Grassmann-stability}.  Couple the two
pseudofermion measures maximally; the expectation of a bounded integrand
changes by at most twice their total variation, which is controlled by
Corollary~\ref{cor:v5v69-pseudofermion-TV}.  Treat the remaining bosonic
law in the same way.  These estimates give a common-algebra weight error
bounded by the parenthesis in
\eqref{eq:v5v69-combined-entry}.  Apply
Corollary~\ref{cor:v5v69-reflected-entry} and then
Lemma~\ref{lem:v5v69-Gram-perturbation}.
\end{proof}

\subsection{Cofinal OS passage}

\begin{definition}[Reflected supertransfer convergence certificate]
\label{def:v5v69-RSCC}
An RSCC consists of:
\begin{enumerate}
 \item common-refinement representatives for a dense gauge-invariant
 positive-time cylinder algebra;
 \item a cutoff-uniform positive normalizer reserve;
 \item local or rooted Grassmann bounds replacing any volume-growing raw
 exponential factor;
 \item relative pseudofermion errors tending to zero in Hilbert--Schmidt
 or a stronger summable topology;
 \item convergence in total variation, or direct cylinder-entry
 convergence, for the remaining bosonic factors;
 \item the V68 cross margin and exact retention of every failed band.
\end{enumerate}
\end{definition}

\begin{theorem}[Conditional cofinal OS theorem]
\label{thm:v5v69-cofinal-OS}
Let \(q_n\) be a cofinal sequence of finite supertransfer functionals,
each reflection positive.  If an RSCC holds and the errors in
\eqref{eq:v5v69-combined-entry} tend to zero for every fixed finite
cylinder family, then the limiting functional is reflection positive on
the limiting cylinder algebra.
\end{theorem}

\begin{proof}
Theorem~\ref{thm:v5v69-supertransfer-stability} makes every entry of every
fixed reflected Gram matrix Cauchy.  Its limit exists.  Each finite-cutoff
matrix is positive semidefinite, so for every coefficient vector \(c\),
\[
 c^*G_\infty c
 =
 \lim_{n\to\infty}c^*G_{q_n}c
 \ge0.
\]
Thus every limiting reflected Gram matrix is positive semidefinite.
\end{proof}

\begin{firewall}[Determinants do not control cylinder entries]
\label{fw:v5v69-determinant-not-entry}
Convergence of
\(\det D_{\mathrm{phys},q}/\det D_{\mathrm{PV},q}\) at zero source controls
only the partition function.  Reflected matrices contain arbitrary
positive-time source insertions.  Their convergence requires the marked
matrix, Gaussian, and normalization bounds in an RSCC or an equivalent
direct proof.
\end{firewall}

\begin{firewall}[The raw global Grassmann norm is not cofinal]
\label{fw:v5v69-global-Grassmann}
The finite V58 coefficient norm proves algebraic continuity, but its
exponential bound can grow with the number of generators.  V69 does not
claim that this raw norm is cutoff uniform.  A rooted connected norm and
the inherited cluster/shell estimates must supply the local version of
item 3 in an RSCC.
\end{firewall}

\subsection{V69 status}

\begin{theorem}[V69 conclusion]
\label{thm:v5v69-status}
Finite reflected cylinder matrices of the exact Pauli--Villars
supertransfer are quantitatively stable under common-refinement changes
of their Grassmann Dirac matrices, pseudofermion square, and bosonic law.
The pseudofermion change is controlled by the relative
Hilbert--Schmidt norm through an exact Gaussian entropy formula, while the
Grassmann change is controlled by a Banach-algebra Duhamel estimate.
These bounds give a conditional cofinal OS theorem.  They also expose two
unproved uniform inputs: a rooted Grassmann norm and a summable relative
pseudofermion error.  Operator-norm or determinant convergence alone does
not supply either input.
\end{theorem}

\begin{proof}
Use Theorem~\ref{thm:v5v69-supertransfer-stability} and
Theorem~\ref{thm:v5v69-cofinal-OS}, with
Firewalls~\ref{fw:v5v69-operator-rank} and
\ref{fw:v5v69-global-Grassmann}.
\end{proof}

V70 will perform the compulsory companion audit.  It will then attack the
pseudofermion summability input by writing the relative square
\(Q_{q'}-Q_q\) as marked local shell terms and testing whether the V59
exponential locality plus the inherited rooted smoother yields a
cutoff-summable Hilbert--Schmidt debit.  If the volume trace diverges, the
observable-local relative entropy will replace the global Gaussian
distance.

\clearpage
\section{V70: companion audit and observable-local pseudofermion convergence}
\label{sec:v5v70}

\subsection{Compulsory V70 companion audit}

The newest live companion files inspected for this note were
\begin{align}
 &\texttt{Renewal\_Yang\_Mills\_Mass\_Gap\_Research\_Notes\_V83.tex},
 \notag\\
 &\quad
 \operatorname{SHA256}=
 \texttt{5DD4C5C871EAAD632803A7DB1E0BEEB5}
 \texttt{A2780B28CC8F36CCC692D18DA8654658},
 \label{eq:v5v70-YM-audit}\\
 &\texttt{Renewal\_NS\_Stability\_Research\_Notes\_V31.tex},
 \notag\\
 &\quad
 \operatorname{SHA256}=
 \texttt{F1121B62238AD5491BB7CF03635EA993}
 \texttt{4942EDF573ED8FAAA9EE79E1D38A6696}.
 \label{eq:v5v70-NS-audit}
\end{align}
The YM file had \(1\,238\,627\) bytes.  The NS file had
\(887\,537\) bytes and is unchanged from the V65 audit.  Neither file was
modified.

YM V79--V83 develops the intrinsic Wilson resolvent, nested gradient
energy cards, and explicit ternary and quartic connected estimates.  V83
proves the all-thermal local tree card through arity four.  Its own
firewall leaves arbitrary-order Gaussian-defect summation, continuum
measure construction, reflection positivity, and the mass gap open.

The transfer to the present note is a restriction, not an imported
estimate.  The pseudofermion field is exactly Gaussian, so its complete
source dependence is determined by one covariance and needs no
finite-arity extrapolation.  The remaining interacting gauge/Grassmann
sector still requires its own all-order rooted estimates; the V83
arity-four result cannot close it.

\subsection{Why the global Hilbert--Schmidt distance is too strong}

Let \(\Lambda_L\) be a periodic lattice of volume \(V_L\), and let
\(X_L\) be a translation-covariant finite-range matrix with a nonzero
kernel independent of \(L\).

\begin{lemma}[Extensive Hilbert--Schmidt obstruction]
\label{lem:v5v70-extensive-HS}
If
\begin{equation}
 c_X:=
 \sum_y\|X_L(0,y)\|_{\mathrm{HS}}^2>0
 \label{eq:v5v70-row-HS}
\end{equation}
is independent of \(L\), then
\begin{equation}
 \|X_L\|_{\mathrm{HS}}^2
 =
 V_Lc_X.
 \label{eq:v5v70-extensive-HS}
\end{equation}
Thus the global Gaussian bound of V69 grows like
\(\sqrt{V_L}\) for a fixed nonzero local perturbation.
\end{lemma}

\begin{proof}
Translation covariance makes every row have the same squared
Hilbert--Schmidt sum \(c_X\).  Sum over the \(V_L\) rows.
\end{proof}

\begin{firewall}[Thermodynamic failure of a metric is not failure of
local convergence]
\label{fw:v5v70-global-metric}
Equation \eqref{eq:v5v70-extensive-HS} says that total variation between
two complete infinite-volume Gaussian measures can be maximal even when
all fixed local marginals converge.  The continuum programme needs
convergence of each fixed cylinder matrix, so the covariance must be
tested after localization to its observable support.
\end{firewall}

\subsection{A weighted kernel algebra}

On a graph with distance \(d\), define
\begin{equation}
 \|A\|_{\mu,1}
 :=
 \max\left\{
 \sup_x\sum_y e^{\mu d(x,y)}\|A_{xy}\|,
 \sup_y\sum_x e^{\mu d(x,y)}\|A_{xy}\|
 \right\}.
 \label{eq:v5v70-weighted-kernel}
\end{equation}

\begin{lemma}[Weighted Schur algebra]
\label{lem:v5v70-weighted-algebra}
For \(\mu\ge0\),
\begin{equation}
 \|AB\|_{\mu,1}
 \le
 \|A\|_{\mu,1}\|B\|_{\mu,1}.
 \label{eq:v5v70-weighted-product}
\end{equation}
\end{lemma}

\begin{proof}
The triangle inequality
\(d(x,y)\le d(x,z)+d(z,y)\) gives
\[
 e^{\mu d(x,y)}
 \|(AB)_{xy}\|
 \le
 \sum_z
 e^{\mu d(x,z)}\|A_{xz}\|
 e^{\mu d(z,y)}\|B_{zy}\|.
\]
Sum first in \(y\), then take the supremum in \(x\).  The column estimate
is identical.
\end{proof}

Let \(Q,\widetilde Q>0\) be two pseudofermion squares on a common
refinement and set
\begin{equation}
 V:=\widetilde Q-Q,
 \qquad
 C:=Q^{-1},
 \quad
 \widetilde C:=\widetilde Q^{-1}.
 \label{eq:v5v70-QCV}
\end{equation}

\begin{theorem}[Weighted local covariance stability]
\label{thm:v5v70-local-covariance}
Suppose
\begin{equation}
 \|C\|_{\mu,1}\le K,
 \qquad
 \|\widetilde C\|_{\mu,1}\le K,
 \qquad
 \|V\|_{\mu,1}\le\varepsilon.
 \label{eq:v5v70-weighted-hypotheses}
\end{equation}
Then
\begin{equation}
 \|C-\widetilde C\|_{\mu,1}
 \le K^2\varepsilon.
 \label{eq:v5v70-local-covariance}
\end{equation}
For a finite observable set \(A\) with projection \(P_A\),
\begin{equation}
 \|P_A(C-\widetilde C)P_A\|_{\mathrm{HS}}
 \le
 \sqrt{\dim P_A}\,K^2\varepsilon.
 \label{eq:v5v70-local-HS}
\end{equation}
The constants are independent of total volume.
\end{theorem}

\begin{proof}
The resolvent identity gives
\[
 C-\widetilde C
 =
 C(\widetilde Q-Q)\widetilde C
 =
 CV\widetilde C.
\]
Apply Lemma~\ref{lem:v5v70-weighted-algebra} twice.  The unweighted
operator norm of any compression is bounded by the weighted Schur norm.
For a matrix of rank at most \(\dim P_A\), its Hilbert--Schmidt norm is at
most the square root of that rank times its operator norm.
\end{proof}

\begin{corollary}[Input from the V59/V68 gaps]
\label{cor:v5v70-gap-input}
If \(D_{\mathrm{PV}}\) and its common-refinement comparison have
volume-uniform singular gaps and finite-range hopping satisfying the V59
weighted Combes--Thomas condition, then their inverse kernels have a
finite bound in \(\|\cdot\|_{\mu',1}\) for every exponent
\(\mu'\) below both the Combes--Thomas exponent and the graph-growth
rate.  If the corresponding square difference \(V\) tends to zero in the
same weighted norm, \eqref{eq:v5v70-local-HS} tends to zero for every
fixed \(A\).
\end{corollary}

\begin{proof}
V59 gives pointwise exponential decay.  Summing it against
\(e^{\mu'd(x,y)}\) is finite when the remaining decay exceeds graph
growth.  Apply Theorem~\ref{thm:v5v70-local-covariance}.
\end{proof}

\subsection{Local Gaussian cylinder entries}

Use the realification of the complex pseudofermion field.  Its centered
Gaussian covariance is a positive matrix \(\mathcal C\) obtained from
\(C=Q^{-1}\).  For a real source \(j\) supported in \(A\), define the Weyl
cylinder observable
\begin{equation}
 W(j):=e^{i\langle j,\phi\rangle}.
 \label{eq:v5v70-Weyl}
\end{equation}

\begin{lemma}[Exact local characteristic stability]
\label{lem:v5v70-characteristic-stability}
For two centered real Gaussians with covariances
\(\mathcal C,\widetilde{\mathcal C}\),
\begin{equation}
 \left|
 \mathbb E_{\mathcal C}W(j)
 -
 \mathbb E_{\widetilde{\mathcal C}}W(j)
 \right|
 \le
 \frac12\|j\|^2
 \|P_A(\mathcal C-\widetilde{\mathcal C})P_A\|_{\mathrm{op}}.
 \label{eq:v5v70-characteristic-stability}
\end{equation}
\end{lemma}

\begin{proof}
The characteristic functions are
\[
 \exp\left(-\frac12\langle j,\mathcal Cj\rangle\right),
 \qquad
 \exp\left(-\frac12\langle j,\widetilde{\mathcal C}j\rangle\right).
\]
Both exponents are nonpositive.  The mean-value theorem and
\(|e^{-x}-e^{-y}|\le|x-y|\) for \(x,y\ge0\) give the result.
\end{proof}

\begin{theorem}[Volume-uniform local pseudofermion entries]
\label{thm:v5v70-local-Gaussian}
Under \eqref{eq:v5v70-weighted-hypotheses}, every fixed finite product of
pseudofermion Weyl observables supported in \(A\) changes by at most
\begin{equation}
 C_JK^2\varepsilon,
 \label{eq:v5v70-local-Weyl}
\end{equation}
where \(C_J\) depends only on the finite source family, not on total
volume.  The same conclusion holds for its reflected products when
\(A\cup\theta A\) is used in place of \(A\).
\end{theorem}

\begin{proof}
A product of Weyl observables is one Weyl observable with source equal to
the signed sum of the individual sources.  Apply
Lemma~\ref{lem:v5v70-characteristic-stability} and
Theorem~\ref{thm:v5v70-local-covariance}.  Reflection only enlarges the
fixed support to \(A\cup\theta A\).
\end{proof}

\begin{corollary}[Local polynomial cylinders]
\label{cor:v5v70-local-polynomial}
If the source norms remain in a fixed bounded set, differentiation of the
Gaussian characteristic functions at the origin gives convergence of
every fixed-degree local polynomial cylinder and its reflected matrix
entries, with constants determined by its degree and support.
\end{corollary}

\begin{proof}
Finite-dimensional Gaussian characteristic functions are entire.
Differentiate their explicit quadratic exponent.  Each derivative is a
finite Wick polynomial in the compressed covariance entries, which
converge by Theorem~\ref{thm:v5v70-local-covariance}.
\end{proof}

\subsection{Distant shell perturbations}

Suppose \(V\) is supported on \(S\times S\), and
\(d(A,S)\ge R\).  Assume pointwise inverse bounds
\begin{equation}
 \|C_{xy}\|,\|\widetilde C_{xy}\|
 \le K_0e^{-\mu d(x,y)}
 \label{eq:v5v70-pointwise-inverse}
\end{equation}
and a uniform unweighted row sum
\(\sup_x\sum_y\|V_{xy}\|\le v_0\).

\begin{proposition}[Off-support exponential debit]
\label{prop:v5v70-distant-shell}
For every \(\mu'<\mu\) above the graph-growth threshold,
\begin{equation}
 \|P_A(C-\widetilde C)P_A\|_{\mathrm{op}}
 \le
 C_{\mu,\mu'}K_0^2v_0e^{-2\mu'R}.
 \label{eq:v5v70-distant-shell}
\end{equation}
\end{proposition}

\begin{proof}
Insert the resolvent identity
\[
 (C-\widetilde C)_{xy}
 =
 \sum_{z,w}C_{xz}V_{zw}\widetilde C_{wy}.
\]
For \(x,y\in A\), every nonzero term has \(z,w\in S\).
The two inverse legs together travel from \(A\) to \(S\) and back, paying
at least \(2R\).  Spend the difference \(\mu-\mu'\) to sum graph growth
and use the row bound on \(V\).
\end{proof}

\begin{assessment}[Two valid local routes]
\label{ass:v5v70-two-routes}
An everywhere-supported adjacent-cutoff perturbation should be controlled
by the weighted smallness
\(\|V\|_{\mu,1}\to0\).  A shell perturbation with order-one local size but
moving away from the observable should use
\eqref{eq:v5v70-distant-shell}.  Combining the two mechanisms is valid by
splitting \(V\) before taking a norm.
\end{assessment}

\subsection{Observable-local RSCC}

\begin{definition}[Observable-local pseudofermion certificate]
\label{def:v5v70-OLPC}
An OLPC consists of:
\begin{enumerate}
 \item the V68 uniform cross margin and a uniform singular gap for the
 Pauli--Villars operator;
 \item common-refinement inverse kernels satisfying weighted exponential
 bounds below graph growth;
 \item for every fixed cylinder support, a decomposition of the relative
 square into a weighted-small local part and an exponentially distant
 part;
 \item convergence of the compressed covariance on the support and its
 reflection;
 \item compatibility with the physical and conjugate-regulator
 Grassmann rooted estimates.
\end{enumerate}
\end{definition}

\begin{theorem}[Observable-local replacement for the V69 global debit]
\label{thm:v5v70-local-RSCC}
In the conditional cofinal OS theorem of V69, the global relative
Hilbert--Schmidt hypothesis for pseudofermions may be replaced by an OLPC.
Then every fixed pseudofermion cylinder block has convergent reflected
entries uniformly in total volume.
\end{theorem}

\begin{proof}
Theorem~\ref{thm:v5v70-local-Gaussian} controls the weighted-small part,
and Proposition~\ref{prop:v5v70-distant-shell} controls the distant part.
Their sum makes every fixed local Gaussian reflected entry Cauchy.  The
remaining factors are handled by the other hypotheses of the V69 RSCC,
so its entrywise limiting proof applies unchanged.
\end{proof}

\begin{firewall}[The relative square still needs a cutoff estimate]
\label{fw:v5v70-relative-square}
V70 proves that a weighted or distant-shell estimate for
\(Q_{q'}-Q_q\) is sufficient.  It does not prove that the actual
domain-wall/overlap refinement has either property after rescaling.
That requires an explicit comparison map between adjacent lattice
operators and ownership of their local counterterms.
\end{firewall}

\begin{firewall}[Gaussian exactness does not close the interacting
cumulants]
\label{fw:v5v70-Gaussian-only}
The all-source pseudofermion result follows from a quadratic covariance.
Gauge, Higgs, and physical-fermion interactions generate connected
cumulants.  The YM V83 arity-four card and the present Gaussian theorem do
not prove their arbitrary-order rooted summation.
\end{firewall}

\subsection{V70 status}

\begin{theorem}[V70 conclusion]
\label{thm:v5v70-status}
The mandatory audit records YM V83 and unchanged NS V31.  A global
Hilbert--Schmidt comparison of pseudofermion covariances is extensive and
is unnecessarily strong for local OS reconstruction.  Weighted resolvent
algebra and the exact resolvent identity give a volume-uniform compressed
covariance bound for each fixed observable support; exponentially distant
shell changes gain two inverse-kernel legs.  These estimates replace the
global pseudofermion debit in the V69 cofinal OS theorem, conditional on an
explicit weighted adjacent-cutoff comparison.  That comparison and the
all-order interacting rooted bounds remain open.
\end{theorem}

\begin{proof}
Use Lemma~\ref{lem:v5v70-extensive-HS},
Theorem~\ref{thm:v5v70-local-covariance},
Proposition~\ref{prop:v5v70-distant-shell}, and
Theorem~\ref{thm:v5v70-local-RSCC}, with the two final firewalls.
\end{proof}

V71 will construct the adjacent-cutoff comparison for the minimal
Wilson domain-wall square.  It will embed the coarse lattice into the
fine lattice by a gauge-covariant block map, expand
\(Q_{\mathrm{fine}}-\iota Q_{\mathrm{coarse}}\iota^*\) into local symbol
and counterterm pieces, and determine which part is weighted-small after
physical rescaling.  If a relevant local term remains order one, it will
be assigned to its running coupling rather than misclassified as an
irrelevant covariance error.


\clearpage
\section{V71: Feshbach comparison of the blocked Pauli--Villars covariance}
\label{sec:v5v71}

\subsection{Why the inherited block map is not yet an OLPC}

V70 reduced pseudofermion convergence to an observable-local covariance
comparison.  The third-series V83 construction already supplies the
required gauge-covariant isometry from a coarse Wilson band into the fine
Hilbert space, together with a local massive spectator complement.  We use
that construction as an inherited input and do not repeat it here.

A direct estimate of
\(Q_f-\iota Q_c\iota^*\) on the whole fine space is generally the wrong
next step.  The fine operator acts nontrivially on the wavelet complement,
whereas the comparison operator carries a deliberately chosen spectator
mass.  Their difference can therefore remain of order one even when all
observables under consideration lie in the retained band.  This note
eliminates the high block exactly before taking a norm.  The resulting
formula isolates two and only two low-energy debits: the Schur defect of
the effective Dirac operator and a positive mixing covariance.

All statements below are finite-dimensional at fixed cutoff and volume.
Their constants are formulated in the weighted Schur algebra of
V70 and hence can be uniform in volume.

\subsection{The covariant low--high decomposition}

Let \(\mathcal H_c\) be the coarse Pauli--Villars spinor space and let
\(\mathcal H_f\) be the fine space.  On a fixed small-link chart, take the
third-series V83 gauge-covariant block isometry
\begin{equation}
 \iota_U:\mathcal H_c\longrightarrow\mathcal H_f,
 \qquad \iota_U^*\iota_U=I_c,
 \qquad P_U=\iota_U\iota_U^*.
 \label{eq:v5v71-isometry}
\end{equation}
Choose its finite-range covariant wavelet completion
\begin{equation}
 \mathscr U_U:\mathcal H_c\oplus\mathcal H_h
 \longrightarrow\mathcal H_f,
 \qquad
 \mathscr U_U(j,0)=\iota_Uj.
 \label{eq:v5v71-wavelet-unitary}
\end{equation}
Here \(\mathcal H_h\) is the orthogonal high-mode space.  A change of
wavelet frame acts by a unitary only on \(\mathcal H_h\); all compressed
identities below are therefore independent of that choice.

After the canonical lattice-field rescaling, write the dimensionless fine
Pauli--Villars operator in this decomposition as
\begin{equation}
 \widehat D_f
 :=\mathscr U_U^*(a_fD_f)\mathscr U_U
 =
 \begin{pmatrix}
  A&K\\ L&E
 \end{pmatrix}.
 \label{eq:v5v71-Dirac-blocks}
\end{equation}
The hats matter: a raw Wilson operator has size \(a_f^{-1}\), so its
unscaled difference cannot tend to zero merely from refinement.  Assume
that the high block \(E\) is invertible and define
\begin{equation}
 S:=A-KE^{-1}L,
 \qquad M:=KE^{-1}.
 \label{eq:v5v71-SM}
\end{equation}
The coarse comparison operator, after the corresponding canonical
rescaling and transport to the same block graph, is denoted
\(\widehat D_c\).

\begin{lemma}[Exact block inverse]
\label{lem:v5v71-block-inverse}
If \(E\) and \(S\) are invertible, then \(\widehat D_f\) is invertible and
\begin{equation}
 \widehat D_f^{-1}
 =
 \begin{pmatrix}
  S^{-1}&-S^{-1}KE^{-1}\\
  -E^{-1}LS^{-1}&
  E^{-1}+E^{-1}LS^{-1}KE^{-1}
 \end{pmatrix}.
 \label{eq:v5v71-block-inverse}
\end{equation}
Conversely, if \(E\) and \(\widehat D_f\) are invertible, then \(S\) is
invertible.
\end{lemma}

\begin{proof}
The exact triangular factorization
\begin{equation}
 \begin{pmatrix}I&-KE^{-1}\\0&I\end{pmatrix}
 \begin{pmatrix}A&K\\L&E\end{pmatrix}
 \begin{pmatrix}I&0\\-E^{-1}L&I\end{pmatrix}
 =\begin{pmatrix}S&0\\0&E\end{pmatrix}
 \label{eq:v5v71-triangular-factorization}
\end{equation}
proves both invertibility assertions.  Multiplying the inverse triangular
factors gives \eqref{eq:v5v71-block-inverse}.
\end{proof}

This is a Dirac Schur complement, but its use here differs from the
inherited Hamiltonian gap test: we now compute the retained covariance of
the positive Pauli--Villars square exactly.

\subsection{Exact retained covariance}

Put
\begin{equation}
 \widehat Q_f=\widehat D_f^*\widehat D_f,
 \qquad
 \widehat C_f=\widehat Q_f^{-1}
 =\widehat D_f^{-1}\widehat D_f^{-*},
 \qquad
 \widehat C_c
 =\widehat D_c^{-1}\widehat D_c^{-*}.
 \label{eq:v5v71-covariances}
\end{equation}
Let \(P_c\) denote projection onto the first summand of
\(\mathcal H_c\oplus\mathcal H_h\).

\begin{theorem}[Feshbach formula for the retained pseudofermion covariance]
\label{thm:v5v71-retained-covariance}
Under the hypotheses of Lemma~\ref{lem:v5v71-block-inverse},
\begin{equation}
 P_c\widehat C_fP_c
 =S^{-1}(I+MM^*)S^{-*}.
 \label{eq:v5v71-retained-covariance}
\end{equation}
In particular,
\begin{equation}
 P_c\widehat C_fP_c-S^{-1}S^{-*}
 =S^{-1}MM^*S^{-*}\ge0.
 \label{eq:v5v71-positive-mixing}
\end{equation}
The correction vanishes exactly if \(K E^{-1}=0\).
\end{theorem}

\begin{proof}
The first block row of \eqref{eq:v5v71-block-inverse} is
\begin{equation}
 \bigl(S^{-1},-S^{-1}M\bigr).
\end{equation}
Multiplying this row by its adjoint gives
\(S^{-1}S^{-*}+S^{-1}MM^*S^{-*}\), proving
\eqref{eq:v5v71-retained-covariance}.  The second term is
\((S^{-1}M)(S^{-1}M)^*\), so it is positive and vanishes exactly when
\(S^{-1}M=0\).  Invertibility of \(S\) makes this equivalent to \(M=0\).
\end{proof}

\begin{assessment}[The missing datum in an effective Dirac comparison]
\label{ass:v5v71-mixing-datum}
Control of the effective operator \(S-\widehat D_c\) alone does not compare
pseudofermion covariances.  The positive factor \(MM^*\) records the norm
of the high-mode component created when a retained source is propagated by
\(\widehat D_f^{-1}\).  It is a separate datum and cannot cancel against
the Schur defect by a sign estimate.
\end{assessment}

\subsection{A volume-uniform weighted estimate}

Use the block graph whose vertices are coarse cells and whose internal
indices contain the retained and wavelet channels.  Equip all kernels with
the symmetric weighted Schur norm \(\|\cdot\|_{\mu,1}\) of
\eqref{eq:v5v70-weighted-kernel}.  It obeys
\begin{equation}
 \|T^*\|_{\mu,1}=\|T\|_{\mu,1}
 \label{eq:v5v71-adjoint-norm}
\end{equation}
and is submultiplicative by
Lemma~\ref{lem:v5v70-weighted-algebra}.

\begin{theorem}[Two-debit adjacent-cutoff covariance bound]
\label{thm:v5v71-two-debit}
Suppose, for some \(K_0,\varepsilon_S,\varepsilon_M<\infty\),
\begin{align}
 \|S^{-1}\|_{\mu,1},\ 
 \|\widehat D_c^{-1}\|_{\mu,1}&\le K_0,
 \label{eq:v5v71-inverse-bounds}\\
 \|S-\widehat D_c\|_{\mu,1}&\le\varepsilon_S,
 \label{eq:v5v71-Schur-defect}\\
 \|KE^{-1}\|_{\mu,1}&\le\varepsilon_M.
 \label{eq:v5v71-mixing-bound}
\end{align}
Then
\begin{equation}
 \left\|P_c\widehat C_fP_c-\widehat C_c\right\|_{\mu,1}
 \le
 2K_0^3\varepsilon_S+K_0^2\varepsilon_M^2.
 \label{eq:v5v71-two-debit-bound}
\end{equation}
Consequently every fixed retained pseudofermion Weyl cylinder, and every
fixed reflected product for a reflection-compatible block map, changes by
at most a constant times the right-hand side of
\eqref{eq:v5v71-two-debit-bound}.
\end{theorem}

\begin{proof}
The resolvent identity at Dirac level gives
\begin{equation}
 S^{-1}-\widehat D_c^{-1}
 =S^{-1}(\widehat D_c-S)\widehat D_c^{-1},
 \label{eq:v5v71-Dirac-resolvent}
\end{equation}
and hence
\begin{equation}
 \|S^{-1}-\widehat D_c^{-1}\|_{\mu,1}
 \le K_0^2\varepsilon_S.
 \label{eq:v5v71-inverse-difference}
\end{equation}
Writing \(R=S^{-1}\) and \(T=\widehat D_c^{-1}\),
\begin{equation}
 RR^*-TT^*=(R-T)R^*+T(R^*-T^*).
\end{equation}
Equations \eqref{eq:v5v71-adjoint-norm} and
\eqref{eq:v5v71-inverse-difference} therefore bound this part by
\(2K_0^3\varepsilon_S\).  The positive mixing term in
\eqref{eq:v5v71-retained-covariance} is bounded by
\begin{equation}
 \|S^{-1}MM^*S^{-*}\|_{\mu,1}
 \le K_0^2\varepsilon_M^2.
\end{equation}
This proves \eqref{eq:v5v71-two-debit-bound}.  Compression to a fixed
support is bounded by the weighted norm, and the characteristic-function
argument of Lemma~\ref{lem:v5v70-characteristic-stability} proves the
cylinder assertion.  If the block partition, weights, paths, and wavelet
frame intertwine time reflection, then
\(\mathscr U_{\theta U}\theta=\theta\mathscr U_U\); the same estimate
applies on the union of a support and its reflection.
\end{proof}

\begin{corollary}[A sufficient one-step OLPC input]
\label{cor:v5v71-OLPC-input}
On a common small-link patch, assume the constants in
\eqref{eq:v5v71-inverse-bounds} are uniform in volume and cutoff, while
\begin{equation}
 \varepsilon_S(a_f)\longrightarrow0,
 \qquad
 \varepsilon_M(a_f)\longrightarrow0.
 \label{eq:v5v71-two-limits}
\end{equation}
Then the retained pseudofermion cylinders satisfy the local covariance
part of the V70 OLPC under one adjacent-cutoff step.
\end{corollary}

\begin{proof}
The right-hand side of \eqref{eq:v5v71-two-debit-bound} tends to zero.
Apply Theorem~\ref{thm:v5v70-local-Gaussian} after identifying a coarse
source \(j\) with the fine covariant source \(\iota_Uj\).
\end{proof}

The conclusion is pointwise in the gauge and Higgs background and uniform
on the declared patch.  Because \(\iota_U\) depends on the dynamical gauge
field, a blocked cylinder is itself a gauge-dependent local observable.
The full functional-integral comparison must carry that dependence; it
cannot replace \(\iota_U\) by its zero-field value after applying the
theorem.

\subsection{Locality of the two debits}

The first debit contains the exact high-mode self-energy
\begin{equation}
 \Sigma:=KE^{-1}L,
 \qquad S=A-\Sigma.
 \label{eq:v5v71-self-energy}
\end{equation}
Finite range of \(K,L\) and a weighted inverse estimate for \(E^{-1}\)
give
\begin{equation}
 \|\Sigma\|_{\mu,1}
 \le
 \|K\|_{\mu,1}\|E^{-1}\|_{\mu,1}\|L\|_{\mu,1},
 \qquad
 \|M\|_{\mu,1}
 \le
 \|K\|_{\mu,1}\|E^{-1}\|_{\mu,1}.
 \label{eq:v5v71-locality-debits}
\end{equation}
These bounds prove locality, but they do not prove smallness: in
dimensionless variables all three factors may have nonzero cutoff limits.

\begin{proposition}[Exact normalized-square identity and its limitation]
\label{prop:v5v71-square-identity}
Let \(\widehat D_{
m blk}\) be any block-diagonal comparison and write
\(\widehat D_f=\widehat D_{
m blk}+\mathcal E\).  Then
\begin{equation}
 \widehat Q_f-\widehat Q_{\rm blk}
 =\widehat D_{\rm blk}^*\mathcal E
  +\mathcal E^*\widehat D_{\rm blk}
  +\mathcal E^*\mathcal E,
 \label{eq:v5v71-square-expansion}
\end{equation}
and hence
\begin{equation}
 \|\widehat Q_f-\widehat Q_{\rm blk}\|_{\mu,1}
 \le
 \bigl(2\|\widehat D_{\rm blk}\|_{\mu,1}
       +\|\mathcal E\|_{\mu,1}\bigr)
 \|\mathcal E\|_{\mu,1}.
 \label{eq:v5v71-square-bound}
\end{equation}
For the unscaled operators the same identity carries the overall factor
\(a_f^{-2}\).  Thus a claim of square smallness must specify both the
canonical field rescaling and a norm in which \(\mathcal E\to0\).
\end{proposition}

\begin{proof}
Expand
\((\widehat D_{\rm blk}+\mathcal E)^*
 (\widehat D_{\rm blk}+\mathcal E)\) and apply the weighted algebra
inequality.  Since \(D=a_f^{-1}\widehat D\), its square is
\(a_f^{-2}\widehat Q\).
\end{proof}

Theorem~\ref{thm:v5v71-two-debit} is stronger for retained observables:
it permits an order-one high block, provided its inverse is local and the
effective defect and retained-to-high propagation become small.

\subsection{Counterterm ownership}

Let \(\mathfrak R\) be the finite-dimensional space of local Dirac kernels
allowed by the exact lattice symmetries at the chosen engineering orders.
Choose a bounded local jet projection
\begin{equation}
 \Pi_{\mathfrak R}:\mathfrak A_{\mu,1}\longrightarrow\mathfrak R
 \label{eq:v5v71-relevant-projection}
\end{equation}
on the weighted kernel algebra and decompose
\begin{equation}
 S-\widehat D_c
 =R_{
m rel}+R_{
m irr},
 \qquad
 R_{
m rel}=\Pi_{\mathfrak R}(S-\widehat D_c).
 \label{eq:v5v71-relevant-split}
\end{equation}
This is bookkeeping only after \(\mathfrak R\), its normalization
conditions, and the projection have been fixed.  It becomes an estimate
when the coarse running parameters are chosen to absorb \(R_{
m rel}\).

\begin{proposition}[Renormalized Schur defect]
\label{prop:v5v71-renormalized-defect}
Suppose the coarse family \(\widehat D_c(g)\) has local coordinates
\(g=(g_1,\ldots,g_N)\) whose derivatives span \(\mathfrak R\) at the
matching point, with inverse coordinate condition number at most
\(K_{\rm match}\).  If
\begin{equation}
 \|R_{\rm rel}\|_{\mu,1}
 <\rho/K_{\rm match}
 \label{eq:v5v71-matching-radius}
\end{equation}
inside a coordinate ball of radius \(\rho\), then there is a unique local
parameter shift \(\delta g\) satisfying the declared normalization
conditions, and
\begin{equation}
 \|\delta g\|\le
 2K_{\rm match}\|R_{\rm rel}\|_{\mu,1}
 \label{eq:v5v71-parameter-shift}
\end{equation}
after shrinking the ball so that the nonlinear coordinate remainder has
Lipschitz constant at most \((2K_{\rm match})^{-1}\).  With
\(\widehat D_c\) evaluated at \(g+\delta g\), the remaining Schur defect is
bounded by the irrelevant remainder plus a quadratic matching error.
\end{proposition}

\begin{proof}
Apply the contraction mapping theorem to the finite-dimensional
normalization map.  Its derivative on \(\mathfrak R\) has a right inverse
of norm at most \(K_{\rm match}\).  Subtract the linear solution and use
the stated Lipschitz bound for the derivative; the resulting map is a
contraction on the ball of radius
\(2K_{\rm match}\|R_{\rm rel}\|_{\mu,1}\).  Taylor's formula then leaves
\(R_{\rm irr}\) plus a remainder quadratic in \(\delta g\).
\end{proof}

The proposition assigns, but does not calculate, the mass, kinetic,
Yukawa, and gauge-background coefficients in \(R_{\rm rel}\).  Their
values must be obtained from the actual Wilson/domain-wall kernel.  An
order-one relevant coefficient is a running-coupling shift; inserting it
into \(\varepsilon_S\) would falsely label it an irrelevant error.

\subsection{Firewalls and next estimate}

\begin{firewall}[A spectator gap gives locality, not decoupling]
\label{fw:v5v71-gap-not-small}
A uniform lower bound on the singular values of \(E\) bounds \(E^{-1}\)
and, with finite range, gives exponential locality.  It does not imply
\(\|KE^{-1}\|_{\mu,1}\to0\).  The mixing numerator must carry a scale gain,
or the retained observable must be redefined by a local whitening map.
\end{firewall}

\begin{firewall}[Pointwise Gaussian comparison is not gauge integration]
\label{fw:v5v71-pointwise-only}
Theorem~\ref{thm:v5v71-two-debit} compares conditional pseudofermion
Gaussians on a common background.  It does not compare the interacting
gauge--Higgs laws, prove arbitrary-order rooted cumulant bounds, or extend
the small-link chart across all topological sectors.
\end{firewall}

\begin{theorem}[V71 conclusion]
\label{thm:v5v71-status}
The inherited covariant block isometry yields an exact adjacent-cutoff
formula for retained Pauli--Villars covariance.  Its difference from the
coarse covariance is controlled by the effective Dirac Schur defect
\(S-\widehat D_c\) and the positive mixing amplitude \(KE^{-1}\), with the
volume-uniform bound \eqref{eq:v5v71-two-debit-bound}.  This avoids the
false requirement that the entire high spectator square approach the fine
square.  Relevant local terms must be absorbed by explicit normalization
conditions before the Schur defect can be tested for an irrelevant gain.
Neither required gain has yet been proved for the concrete Wilson block.
\end{theorem}

\begin{proof}
Combine Theorem~\ref{thm:v5v71-retained-covariance},
Theorem~\ref{thm:v5v71-two-debit},
Proposition~\ref{prop:v5v71-square-identity}, and
Proposition~\ref{prop:v5v71-renormalized-defect}, subject to the two
firewalls.
\end{proof}

V72 will test the mixing debit for an explicit free Wilson block.  It will
show whether a piecewise-constant retained vector has a scale gain, and, if
not, will identify the moment or energy-adapted correction needed to make
\(KE^{-1}\) small without changing the gauge-covariant topology.


\clearpage
\section{V72: the Haar mixing obstruction and exact covariance whitening}
\label{sec:v5v72}

\subsection{A free calculation that rules out uniform mixing smallness}

V71 isolated the retained-to-high propagation
\(M=KE^{-1}\) as a sufficient smallness debit.  Before attempting to
estimate it in the interacting theory, one must test the debit on the
free Wilson symbol.  The following one-dimensional slice is already
decisive.  It occurs inside every product block in four dimensions by
setting the transverse momenta to zero.

Let \(T\) be fine-lattice translation by one site and consider the
dimensionless Wilson operator
\begin{equation}
 d=m_0+\frac12\gamma(T-T^*)
   +\frac r2(2I-T-T^*),
 \qquad \gamma^*=\gamma,
 \qquad \gamma^2=I.
 \label{eq:v5v72-Wilson-slice}
\end{equation}
Group sites \(2n,2n+1\) into cells and introduce the Haar channels
\begin{equation}
 c_n=\frac{\psi_{2n}+\psi_{2n+1}}{\sqrt2},
 \qquad
 h_n=\frac{\psi_{2n}-\psi_{2n+1}}{\sqrt2}.
 \label{eq:v5v72-Haar-channels}
\end{equation}
For coarse Bloch momentum \(p\) put \(z=e^{ip}\).  In the even--odd
basis,
\begin{equation}
 T(p)=\begin{pmatrix}0&1\\z&0\end{pmatrix}.
 \label{eq:v5v72-translation}
\end{equation}

\begin{lemma}[Exact Haar translation symbol]
\label{lem:v5v72-Haar-translation}
In the ordered \((c,h)\) basis,
\begin{equation}
 \widetilde T(p)
 =\frac12
 \begin{pmatrix}
  1+z&z-1\\
  1-z&-(1+z)
 \end{pmatrix}.
 \label{eq:v5v72-Haar-translation-symbol}
\end{equation}
Consequently the low-to-high block \(K(p)\) and high block \(E(p)\) of
\(d(p)\) are
\begin{align}
 K(p)&=-\gamma\sin^2\frac p2-\frac{ir}{2}\sin p,
 \label{eq:v5v72-K-symbol}\\
 E(p)&=m_0-\frac{i\gamma}{2}\sin p
       +\frac r2(3+\cos p).
 \label{eq:v5v72-E-symbol}
\end{align}
\end{lemma}

\begin{proof}
The normalized low and high column vectors are
\(2^{-1/2}(1,1)^T\) and \(2^{-1/2}(1,-1)^T\).  Taking their matrix
elements against \eqref{eq:v5v72-translation} gives
\eqref{eq:v5v72-Haar-translation-symbol}.  Its low--high entries in
\(T-T^*\) and \(T+T^*\) are respectively
\begin{equation}
 \cos p-1,
 \qquad i\sin p.
\end{equation}
Its high--high entries are \(-i\sin p\) and \(-(1+\cos p)\).
Substitution in \eqref{eq:v5v72-Wilson-slice} proves
\eqref{eq:v5v72-K-symbol}--\eqref{eq:v5v72-E-symbol}.
\end{proof}

\begin{theorem}[No cutoff gain for the piecewise-constant block]
\label{thm:v5v72-no-Haar-gain}
Assume \(m_0+r\ne0\).  At the edge \(p=\pi\) of the reduced Brillouin
zone,
\begin{equation}
 K(\pi)=-\gamma,
 \qquad E(\pi)=m_0+r,
 \qquad
 \|K(\pi)E(\pi)^{-1}\|=|m_0+r|^{-1}.
 \label{eq:v5v72-edge-mixing}
\end{equation}
Hence for every refinement level,
\begin{equation}
 \|KE^{-1}\|_{\mathrm{op}}
 \ge |m_0+r|^{-1}.
 \label{eq:v5v72-mixing-lower}
\end{equation}
The same lower bound holds for the zero-field four-dimensional product
Haar block on the corresponding one-axis momentum slice.
\end{theorem}

\begin{proof}
Set \(p=\pi\) in
\eqref{eq:v5v72-K-symbol}--\eqref{eq:v5v72-E-symbol} and use
\(\|\gamma\|=1\).  The dimensionless reduced symbol is identical at every
refinement level, so the value does not acquire a power of the lattice
spacing.  In four dimensions, take all transverse translations to have
Bloch value one and choose the wavelet channel in the displayed direction;
the same two-channel restriction is present.
\end{proof}

Thus the sufficient condition
\(\varepsilon_M(a_f)\to0\) in
Corollary~\ref{cor:v5v71-OLPC-input} is too strong for this natural block.
The failure is an alias-edge effect, not a failure of locality or of the
high gap.  At small \(p\), one instead has
\begin{equation}
 K(p)=-\frac{ir}{2}p-\frac14\gamma p^2+O(p^3),
 \label{eq:v5v72-small-p-mixing}
\end{equation}
but a weighted kernel norm sees the whole reduced zone and cannot use this
pointwise infrared gain.

\subsection{Exact integration of the mixing covariance}

Return to the covariant block decomposition of V71.  Define
\begin{equation}
 Z:=I+MM^*=I+KE^{-1}E^{-*}K^*.
 \label{eq:v5v72-Z}
\end{equation}
It is strictly positive and satisfies \(Z\ge I\).  Let
\begin{equation}
 D_{\rm cov}:=Z^{-1/2}S,
 \label{eq:v5v72-Dcov}
\end{equation}
where the positive square root is used.

\begin{theorem}[Exact whitening of the retained covariance]
\label{thm:v5v72-whitening}
Under the invertibility hypotheses of V71,
\begin{equation}
 P_c\widehat C_fP_c
 =D_{\rm cov}^{-1}D_{\rm cov}^{-*},
 \label{eq:v5v72-whitened-covariance}
\end{equation}
and its inverse is
\begin{equation}
 (P_c\widehat C_fP_c)^{-1}
 =D_{\rm cov}^*D_{\rm cov}
 =S^*Z^{-1}S.
 \label{eq:v5v72-marginal-precision}
\end{equation}
Therefore the marginal of the normalized fine pseudofermion Gaussian on
the retained variable is exactly the Gaussian with action
\begin{equation}
 \|D_{\rm cov}\phi_c\|^2.
 \label{eq:v5v72-marginal-action}
\end{equation}
No smallness assumption on \(M\) is required.
\end{theorem}

\begin{proof}
Since \(D_{\rm cov}^{-1}=S^{-1}Z^{1/2}\),
\begin{equation}
 D_{\rm cov}^{-1}D_{\rm cov}^{-*}
 =S^{-1}ZS^{-*}.
\end{equation}
Theorem~\ref{thm:v5v71-retained-covariance} identifies the right-hand side
with \(P_c\widehat C_fP_c\), proving
\eqref{eq:v5v72-whitened-covariance}.  Inversion gives
\eqref{eq:v5v72-marginal-precision}.  A finite-dimensional centered
Gaussian marginal is determined by its compressed covariance, proving
\eqref{eq:v5v72-marginal-action} after normalization.
\end{proof}

\begin{corollary}[Exact positive-square Schur complement]
\label{cor:v5v72-square-Schur}
Let \(\widehat Q_f\) be written in low--high blocks.  Its Schur complement
after Gaussian integration of the high pseudofermion variable equals
\begin{equation}
 \widehat Q_f/\widehat Q_{f,hh}
 =S^*Z^{-1}S.
 \label{eq:v5v72-square-Schur}
\end{equation}
This identity remains valid although the Dirac operator is nonnormal.
\end{corollary}

\begin{proof}
The inverse of a positive block matrix has upper-left block equal to the
inverse of its high-block Schur complement.  Apply this fact to
\(\widehat Q_f\), then use
\eqref{eq:v5v72-marginal-precision}.  The proof never commutes \(S\),
\(M\), or their adjoints.
\end{proof}

This identifies the right adjacent-cutoff object.  Comparing the raw
square or demanding \(M\to0\) discards an exact local wavefunction factor;
the marginal Dirac kernel \(D_{\rm cov}\) retains it.

\subsection{Locality and covariance of the whitening map}

The square root must be controlled before it is used in a local continuum
argument.  The next statement makes the needed hypotheses explicit.

\begin{proposition}[Local covariant square root]
\label{prop:v5v72-local-square-root}
Suppose \(M(U)\) is gauge covariant and
\begin{equation}
 \|M(U)\|_{\mu,1}\le m_*,
 \label{eq:v5v72-M-local}
\end{equation}
uniformly in volume and on a common field patch.  Assume the weighted
holomorphic functional calculus from V59 at every exponent
\(0<\mu'<\mu\).  Then
\(Z(U)^{\pm1/2}\) are gauge covariant, have kernels exponentially local
with exponent \(\mu'\), and obey volume-uniform weighted bounds depending
only on \(m_*\), \(\mu-\mu'\), and the graph-growth constants.  If the
block data intertwine time reflection, so do \(Z^{\pm1/2}\) and
\(D_{\rm cov}\).
\end{proposition}

\begin{proof}
The weighted algebra gives
\(\|MM^*\|_{\mu,1}\le m_*^2\), while
\(\operatorname{spec}Z\subset[1,1+m_*^2]\) in operator norm.  The
functions \(z^{1/2}\) and \(z^{-1/2}\) are holomorphic on a neighbourhood
of this compact interval disjoint from the nonpositive real axis.
Weighted holomorphic functional calculus at the slightly smaller exponent
therefore places both functions of \(Z\) in the same local kernel algebra
with uniform bounds.  Functional calculus commutes with unitary gauge
conjugation and with the reflection intertwiner.  Equation
\eqref{eq:v5v72-Dcov} completes the claim.
\end{proof}

\begin{firewall}[Locality requires a uniform functional-calculus margin]
\label{fw:v5v72-functional-margin}
Finite-volume existence of \(Z^{-1/2}\) is automatic because \(Z\ge I\).
Uniform exponential locality additionally requires uniform weighted bounds
for \(M\) and a common functional-calculus exponent.  A volume-dependent
wavelet frame or a collapsing high-block gap would not suffice.
\end{firewall}

\subsection{The improved adjacent-cutoff estimate}

\begin{theorem}[One-debit retained covariance comparison]
\label{thm:v5v72-one-debit}
Assume \(D_{\rm cov}\) and \(\widehat D_c\) are invertible and, for some
\(K_1,\varepsilon_{\rm cov}<\infty\),
\begin{equation}
 \|D_{\rm cov}^{-1}\|_{\mu',1},\ 
 \|\widehat D_c^{-1}\|_{\mu',1}\le K_1,
 \qquad
 \|D_{\rm cov}-\widehat D_c\|_{\mu',1}
 \le\varepsilon_{\rm cov}.
 \label{eq:v5v72-one-debit-hypotheses}
\end{equation}
Then
\begin{equation}
 \|P_c\widehat C_fP_c-\widehat C_c\|_{\mu',1}
 \le2K_1^3\varepsilon_{\rm cov}.
 \label{eq:v5v72-one-debit-bound}
\end{equation}
This conclusion permits \(\|M\|_{\mu,1}\) to stay of order one.
\end{theorem}

\begin{proof}
Use \eqref{eq:v5v72-whitened-covariance} and the resolvent identity
\begin{equation}
 D_{\rm cov}^{-1}-\widehat D_c^{-1}
 =D_{\rm cov}^{-1}(\widehat D_c-D_{\rm cov})
  \widehat D_c^{-1}.
\end{equation}
Its weighted norm is at most \(K_1^2\varepsilon_{\rm cov}\).  Expanding
the difference of the two inverse squares as in the proof of
Theorem~\ref{thm:v5v71-two-debit} yields
\eqref{eq:v5v72-one-debit-bound}.
\end{proof}

\begin{corollary}[Whitened OLPC input]
\label{cor:v5v72-whitened-OLPC}
If the hypotheses of Proposition~\ref{prop:v5v72-local-square-root} and
Theorem~\ref{thm:v5v72-one-debit} hold uniformly along adjacent cutoffs and
\begin{equation}
 \varepsilon_{\rm cov}(a_f)\longrightarrow0,
 \label{eq:v5v72-cov-limit}
\end{equation}
then the retained pseudofermion sector satisfies the local Gaussian part of
the V70 OLPC.  The same holds for reflected cylinders under a
reflection-compatible block system.
\end{corollary}

\begin{proof}
Equation \eqref{eq:v5v72-one-debit-bound} gives convergence of compressed
covariances.  Apply Theorem~\ref{thm:v5v70-local-Gaussian} and its reflected
version.
\end{proof}

\subsection{Ownership of the whitening factor}

At small fields and momenta, the local kernel \(Z^{-1/2}\) contributes a
wavefunction normalization and higher-derivative terms.  Its zero-momentum
symmetry-allowed component belongs to the relevant/marginal normalization
space \(\mathfrak R\) of V71.  It must be included when matching
\(D_{\rm cov}\) to \(\widehat D_c\).  The remaining momentum-dependent
part can be called irrelevant only after a Taylor remainder bound in the
chosen weighted norm.

\begin{firewall}[Infrared Taylor gain is not a weighted global gain]
\label{fw:v5v72-IR-not-weighted}
Equation \eqref{eq:v5v72-small-p-mixing} is a pointwise small-momentum
expansion.  It does not prove
\(\|D_{\rm cov}-\widehat D_c\|_{\mu',1}\to0\), because that norm includes
alias-edge modes.  A proof must either match the full reduced-zone symbol,
use a scale-dependent smooth low-pass block, or localize the observable in
both position and spectral scale with a separately controlled error.
\end{firewall}

\begin{theorem}[V72 conclusion]
\label{thm:v5v72-status}
The piecewise-constant covariant block has an order-one retained/high
mixing amplitude already for the free Wilson operator, so the two-debit
criterion of V71 cannot close by requiring \(KE^{-1}\to0\).  Exact
Gaussian elimination instead produces the local positive factor
\(Z=I+KE^{-1}E^{-*}K^*\) and the whitened marginal Dirac kernel
\(D_{\rm cov}=Z^{-1/2}S\).  Retained covariance convergence is reduced to
the single local comparison
\(D_{\rm cov}-\widehat D_c\to0\), subject to uniform high-gap and
functional-calculus bounds.  That comparison remains unproved for the
full reduced Brillouin zone.
\end{theorem}

\begin{proof}
Use Theorem~\ref{thm:v5v72-no-Haar-gain},
Theorem~\ref{thm:v5v72-whitening},
Proposition~\ref{prop:v5v72-local-square-root}, and
Theorem~\ref{thm:v5v72-one-debit}, with the two firewalls.
\end{proof}

V73 will decide whether a finite-range coarse Wilson stencil can match the
free whitened symbol uniformly on the reduced Brillouin zone.  It will
separate exact representability from approximation, quantify the best
finite-range symbol error, and identify the range growth required for that
error to vanish along refinement.


\clearpage
\section{V73: perfect-stencil representation of the whitened coarse kernel}
\label{sec:v5v73}

\subsection{The representability problem}

V72 replaced the false small-mixing requirement by the exact marginal
kernel
\begin{equation}
 D_{\rm cov}(U)=Z(U)^{-1/2}S(U).
 \label{eq:v5v73-Dcov-recall}
\end{equation}
To close the pseudofermion part of the OLPC, one must decide what is meant
by the next coarse operator.  Restricting it a priori to the original
nearest-neighbour Wilson stencil is not innocuous: Schur elimination and
the square root in \eqref{eq:v5v73-Dcov-recall} generally create infinitely
many exponentially decaying hoppings.  This note proves that such a kernel
has an exact exponentially local perfect-stencil representation, while a
fixed finite hopping range cannot in general reproduce it uniformly.

Let \(k\in\mathbb T^4_{\rm red}\) be reduced-zone momentum, including the
finite block-channel indices.  Write
\begin{equation}
 d_*(k):=D_{\rm cov}(I;k).
 \label{eq:v5v73-target-symbol}
\end{equation}
All masses and dimensionless Wilson parameters are kept inside a compact
safe set.

\subsection{Analyticity of the free whitened symbol}

For a matrix Laurent polynomial \(F(k)\), define its holomorphic adjoint
by
\begin{equation}
 F^{\sharp}(k):=F(-k)^*,
 \label{eq:v5v73-sharp}
\end{equation}
where the star conjugates coefficients, not the complex variable.  On real
momentum this equals \(F(k)^*\) for the translation conventions used here.

\begin{lemma}[A common complex strip]
\label{lem:v5v73-complex-strip}
Suppose the free high block \(E(k)\) has
\begin{equation}
 \inf_{k\in\mathbb T^4_{\rm red}}s_{\min}(E(k))\ge e_0>0.
 \label{eq:v5v73-high-gap}
\end{equation}
Then there is \(\rho_0>0\), depending only on the finite stencil,
\(e_0\), and the compact parameter set, such that \(E(k)^{-1}\) is
holomorphic for \(|\operatorname{Im}k_j|<\rho_0\).  After possibly
shrinking \(\rho_0\),
\begin{equation}
 Z(k)=I+M(k)M^{\sharp}(k)
 \label{eq:v5v73-complex-Z}
\end{equation}
has spectrum in a compact subset of
\(\mathbb C\setminus(-\infty,0]\) throughout that strip.  Hence the
principal \(Z(k)^{-1/2}\) and \(d_*(k)\) are holomorphic there and uniformly
bounded.
\end{lemma}

\begin{proof}
The determinant of the Laurent-polynomial matrix \(E(k)\) is continuous
and nonzero on the compact real torus by
\eqref{eq:v5v73-high-gap}.  Uniform continuity of its complex extension
and the resolvent identity preserve invertibility on a sufficiently thin
strip.  On the real torus, \(Z=I+MM^*\ge I\).  Its spectrum therefore has
positive distance from the square-root cut.  The complex matrix
\eqref{eq:v5v73-complex-Z} depends holomorphically on \(k\); compactness and
spectral upper semicontinuity preserve a positive distance from the cut on
a thinner strip.  Holomorphic matrix functional calculus gives the
principal inverse square root.  Products with the holomorphic Schur
complement preserve holomorphy and boundedness.
\end{proof}

\subsection{Fourier decay and finite-range approximation}

Expand
\begin{equation}
 d_*(k)=\sum_{n\in\mathbb Z^4}d_ne^{in\cdot k}.
 \label{eq:v5v73-Fourier-series}
\end{equation}

\begin{theorem}[Exponential perfect-stencil coefficients]
\label{thm:v5v73-Fourier-decay}
For every \(0<\rho<\rho_0\), there is \(C_\rho<\infty\) such that
\begin{equation}
 \|d_n\|\le C_\rho e^{-\rho|n|_1}.
 \label{eq:v5v73-Fourier-decay}
\end{equation}
If
\begin{equation}
 d_*^{(R)}(k)=\sum_{|n|_1\le R}d_ne^{in\cdot k},
 \label{eq:v5v73-truncation}
\end{equation}
then for \(0\le\mu<\rho\) and every
\(0<\sigma<\rho-\mu\),
\begin{equation}
 \sum_{|n|_1>R}e^{\mu|n|_1}\|d_n\|
 \le C_{\rho,\mu,\sigma}e^{-\sigma R}.
 \label{eq:v5v73-weighted-tail}
\end{equation}
The constants are independent of lattice volume.
\end{theorem}

\begin{proof}
Shift the \(j\)-th Fourier contour to
\(\operatorname{Im}k_j=-\rho\operatorname{sgn}n_j\).  Cauchy's theorem
and the strip supremum give \eqref{eq:v5v73-Fourier-decay}.  The number of
lattice points with \(|n|_1=N\) is bounded by a cubic polynomial in \(N\).
Thus the weighted tail is bounded by
\begin{equation}
 C_\rho\sum_{N>R}C(1+N)^3e^{-(\rho-\mu)N}.
\end{equation}
Absorb the polynomial into
\(e^{((\rho-\mu)-\sigma)N}\) to obtain
\eqref{eq:v5v73-weighted-tail}.
\end{proof}

Lattice symmetries are preserved by truncating with a symmetry-invariant
norm such as \(|n|_1\).  If a finite symmetry group is represented with a
different fundamental cell, average the truncated kernel over that group;
the averaging is a norm contraction and does not enlarge the maximal
range beyond the largest orbit radius.

\begin{proposition}[Fixed-range obstruction]
\label{prop:v5v73-fixed-range-obstruction}
Let \(\mathcal P_R\) be the finite-dimensional space of matrix
trigonometric polynomials supported in \(|n|_1\le R\).  If
\(d_*\notin\mathcal P_R\), then
\begin{equation}
 \delta_R:=\inf_{p\in\mathcal P_R}
 \|d_*-p\|_{C(\mathbb T^4_{\rm red})}>0.
 \label{eq:v5v73-positive-best-error}
\end{equation}
Consequently a self-similar refinement whose dimensionless target symbol
is \(d_*\) cannot have uniform symbol error tending to zero while its
coarse hopping range remains bounded by \(R\).
\end{proposition}

\begin{proof}
Every finite-dimensional subspace of a normed space is closed.  Hence a
point outside \(\mathcal P_R\) has strictly positive distance from it.
Self-similar dimensionless refinement repeats the same approximation
problem at every level, so the lower bound is unchanged.
\end{proof}

The hypothesis is typical rather than automatic.  For example, if
\(\alpha>0\), the periodic function
\begin{equation}
 f(p)=\bigl(1+\alpha(1-\cos p)\bigr)^{-1/2}
 \label{eq:v5v73-square-root-example}
\end{equation}
is not a trigonometric polynomial.  Indeed, if a Laurent polynomial
\(g(z)\) equalled \(f\), then
\(g(z)^2[1+\alpha-\alpha(z+z^{-1})/2]=1\).  The bracket is a nonconstant
Laurent polynomial and hence not a unit in
\(\mathbb C[z,z^{-1}]\), whose only units are \(cz^m\); contradiction.
The whitening square root creates precisely this algebraic mechanism in
generic scalar channels.

\begin{corollary}[Logarithmic range schedule]
\label{cor:v5v73-log-range}
Let \(a\downarrow0\) and choose
\begin{equation}
 R(a)=\left\lceil\frac{s}{\sigma}
 \log\frac{a_0}{a}\right\rceil,
 \qquad s>0.
 \label{eq:v5v73-range-schedule}
\end{equation}
Then the weighted truncation error is
\begin{equation}
 O\bigl((a/a_0)^s\bigr),
 \label{eq:v5v73-power-error}
\end{equation}
while its physical hopping radius is
\begin{equation}
 aR(a)=O\bigl(a\log(a_0/a)\bigr)\longrightarrow0.
 \label{eq:v5v73-physical-radius}
\end{equation}
\end{corollary}

\begin{proof}
Insert \eqref{eq:v5v73-range-schedule} in
\eqref{eq:v5v73-weighted-tail}.  The physical-radius limit is elementary.
\end{proof}

\subsection{Exact gauge-covariant exponential stencil}

Finite truncation is not necessary for locality.  For each displacement
\(n\), fix a translation-covariant family of lattice paths of length at
most \(c_p|n|_1+c_0\), and let \(T_U(x+n,x)\) denote parallel transport
along the selected path.  Define
\begin{equation}
 (D_{\rm perf}(U)\psi)(x)
 :=\sum_{n\in\mathbb Z^4}d_n
 T_U(x,x-n)\psi(x-n),
 \label{eq:v5v73-perfect-covariantization}
\end{equation}
with the internal block channels transported in the inherited covariant
frame.  Symmetry-related paths are averaged when required.

\begin{theorem}[Covariant perfect-stencil bound]
\label{thm:v5v73-perfect-stencil}
The series \eqref{eq:v5v73-perfect-covariantization} converges in every
weighted Schur norm with \(\mu<\rho\), is gauge covariant, and satisfies
\begin{equation}
 D_{\rm perf}(I;k)=d_*(k).
 \label{eq:v5v73-perfect-free}
\end{equation}
If \(\|U_e-I\|\le e^\eta-1\) for every link and
\begin{equation}
 \mu+c_p\eta<\rho,
 \label{eq:v5v73-chart-condition}
\end{equation}
then
\begin{equation}
 \|D_{\rm perf}(U)-D_{\rm perf}(I)\|_{\mu,1}
 \le C\eta,
 \label{eq:v5v73-perfect-Lipschitz}
\end{equation}
where \(C\) is uniform in volume and for \(\eta\) in a compact subinterval
of \eqref{eq:v5v73-chart-condition}.
\end{theorem}

\begin{proof}
Equations \eqref{eq:v5v73-Fourier-decay} and the path-length bound make the
weighted row and column sums absolutely convergent.  Parallel transport
transforms only at its endpoints, proving gauge covariance, and it equals
the identity at the free field, proving
\eqref{eq:v5v73-perfect-free}.  A product of \(\ell\) links differs from
the identity by at most \(e^{\ell\eta}-1\).  The mean-value bound
\begin{equation}
 e^{\ell\eta}-1\le\ell\eta e^{\ell\eta}
\end{equation}
reduces \eqref{eq:v5v73-perfect-Lipschitz} to the convergent sum
\begin{equation}
 C\eta\sum_n(1+|n|_1)
 e^{-(\rho-\mu-c_p\eta)|n|_1}.
\end{equation}
\end{proof}

\begin{proposition}[Conditional free-to-background comparison]
\label{prop:v5v73-background-comparison}
Assume the exact marginal kernel \(D_{\rm cov}(U)\) has a volume-uniform
weighted first link derivative on the same chart:
\begin{equation}
 \sup_{0\le t\le1}
 \left\|\frac d{dt}D_{\rm cov}(U(t))\right\|_{\mu,1}
 \le C_{\rm cov}\eta
 \label{eq:v5v73-Dcov-derivative}
\end{equation}
along \(U(t)=\exp(t\log U)\).  Then
\begin{equation}
 \|D_{\rm cov}(U)-D_{\rm perf}(U)\|_{\mu,1}
 \le(C_{\rm cov}+C)\eta.
 \label{eq:v5v73-background-defect}
\end{equation}
In particular, on a canonically gauge-fixed smooth-field family with
\(\eta(a)=O(a^\alpha)\), \(\alpha>0\), the whitened OLPC debit tends to
zero.
\end{proposition}

\begin{proof}
At \(U=I\), the two kernels agree by definition and
\eqref{eq:v5v73-perfect-free}.  Integrate
\eqref{eq:v5v73-Dcov-derivative} along the radial link path and combine it
with \eqref{eq:v5v73-perfect-Lipschitz}.  The last assertion follows from
Theorem~\ref{thm:v5v72-one-debit} once the inverse bounds are uniform.
\end{proof}

The derivative hypothesis is plausible from the V59 joint cluster bounds
and the local functional calculus of V72, but it has not been verified for
the complete domain-wall kernel with its fifth-direction parameters.

\subsection{The link-chart and reflection costs}

For the truncated range \(R\), direct covariantization changes a length
\(R\) path by \(e^{R\eta}-1\).  Thus a fixed perturbative gap estimate of
the V84 type requires
\begin{equation}
 R\eta\le c_*.
 \label{eq:v5v73-range-chart-cost}
\end{equation}
With the logarithmic schedule
\eqref{eq:v5v73-range-schedule}, a cutoff-independent link chart cannot be
deduced from that estimate.  Smooth continuum fields instead have
\(\eta=O(a)\), for which
\(R(a)\eta(a)=O(a\log(a_0/a))\to0\).

\begin{firewall}[A perfect stencil is not automatically an OS action]
\label{fw:v5v73-perfect-not-RP}
Exponential locality and gauge covariance do not imply reflection
positivity.  The perfect stencil contains arbitrarily long temporal
hoppings, so the nearest-neighbour cross-plane matrix test of V67 does not
apply term by term.  The exact marginal has a stronger possible route---
reflection positivity may descend from the fine measure---but that descent
must use a block map which preserves the positive-time algebra.
\end{firewall}

\begin{firewall}[The smooth-link family is not the full gauge measure]
\label{fw:v5v73-smooth-link-only}
The condition \(\eta(a)=O(a^\alpha)\) describes a controlled continuum
background after a compatible gauge choice.  It is not a concentration
theorem for the interacting lattice gauge law and does not cover large
gauge representatives or nontrivial topological charts.
\end{firewall}

\begin{theorem}[V73 conclusion]
\label{thm:v5v73-status}
The free whitened marginal kernel is analytic on a uniform complex momentum
strip whenever the high block remains gapped.  It therefore has an exact
gauge-covariant exponentially local perfect stencil and finite-range
approximants with exponentially small weighted error.  A bounded hopping
range cannot approximate a genuinely nonpolynomial self-similar symbol;
logarithmic range growth gives an arbitrary power error while its physical
radius still vanishes.  On a smooth-link family, a uniform link-derivative
bound makes the perfect-stencil OLPC debit tend to zero.  Reflection
positivity of the resulting coarse description has not yet been inherited.
\end{theorem}

\begin{proof}
Combine Lemma~\ref{lem:v5v73-complex-strip},
Theorem~\ref{thm:v5v73-Fourier-decay},
Proposition~\ref{prop:v5v73-fixed-range-obstruction},
Theorem~\ref{thm:v5v73-perfect-stencil}, and
Proposition~\ref{prop:v5v73-background-comparison}, subject to the two
firewalls.
\end{proof}

V74 will take the upstream measure-theoretic route suggested by the first
firewall.  It will prove exact reflection-positivity descent under a
reflection-compatible, time-one-sided covariant block map and determine
whether the whitened marginal may be iterated without re-establishing a
finite temporal cross-matrix factorization at every scale.


\clearpage
\section{V74: reflection-positive descent through exact block marginals}
\label{sec:v5v74}

\subsection{The upstream measure-theoretic route}

V73 showed that the exact whitened marginal kernel is exponentially local
but generally has arbitrarily long temporal hoppings.  Reapplying the V67
nearest-neighbour cross-matrix calculation to such a kernel would require a
new factorization at every range.  This is unnecessary when the kernel is
kept as the exact marginal of a reflection-positive fine measure.
Reflection positivity is inherited by pullback of observables, provided
the block map preserves the positive-time algebra.

We formulate this first without Gaussian assumptions.  Let
\(\mathcal A_f\) and \(\mathcal A_c\) be unital graded observable algebras,
with antilinear reflections \(\Theta_f,\Theta_c\) and declared
positive-time subalgebras \(\mathcal A_{f,+},\mathcal A_{c,+}\).  A state
\(\omega_f\) is reflection positive when
\begin{equation}
 \omega_f\bigl((\Theta_fF)F\bigr)\ge0
 \qquad(F\in\mathcal A_{f,+}),
 \label{eq:v5v74-fine-RP}
\end{equation}
with the standard graded product convention in the fermionic case.

\begin{definition}[Reflection-admissible block pullback]
\label{def:v5v74-reflection-block}
A unital grading-preserving algebra map
\(\mathcal B:\mathcal A_c\to\mathcal A_f\) is
reflection-admissible if
\begin{equation}
 \mathcal B(\mathcal A_{c,+})\subset\mathcal A_{f,+},
 \qquad
 \mathcal B\Theta_c=\Theta_f\mathcal B.
 \label{eq:v5v74-reflection-admissible}
\end{equation}
The exact coarse state is the pullback
\begin{equation}
 \omega_c:=\omega_f\circ\mathcal B.
 \label{eq:v5v74-coarse-state}
\end{equation}
\end{definition}

\begin{theorem}[Reflection positivity descends exactly]
\label{thm:v5v74-RP-descent}
If \(\omega_f\) is reflection positive and \(\mathcal B\) is
reflection-admissible, then \(\omega_c\) is reflection positive.
No locality or finite-range assumption on a coarse action is needed.
\end{theorem}

\begin{proof}
For \(F\in\mathcal A_{c,+}\), multiplicativity and
\eqref{eq:v5v74-reflection-admissible} give
\begin{equation}
 \omega_c((\Theta_cF)F)
 =\omega_f\bigl((\Theta_f\mathcal BF)(\mathcal BF)\bigr)\ge0,
\end{equation}
because \(\mathcal BF\in\mathcal A_{f,+}\).
\end{proof}

\subsection{A reflection-compatible covariant block geometry}

Take reflection in the link plane between fine time slices \(0\) and
\(1\).  Choose the block side \(q\) and time-cell boundaries so that no
block intersects both half-lattices.  Pair every positive block \(C_b\)
with \(\theta C_b=C_{\theta b}\).  Choose weights and internal paths with
\begin{equation}
 w_{\theta x}=\overline{w_x},
 \qquad
 T_{\theta U}(\theta x,\theta b)
 =\Theta T_U(x,b)\Theta^{-1}.
 \label{eq:v5v74-reflected-block-data}
\end{equation}
The second relation is obtained by choosing the path in the reflected
block to be the reflected path, with link orientation reversed when
required.

\begin{lemma}[One-sided block intertwining]
\label{lem:v5v74-one-sided-block}
For the V83 block isometry constructed from
\eqref{eq:v5v74-reflected-block-data},
\begin{equation}
 \iota_{\theta U}\Theta_c=\Theta_f\iota_U.
 \label{eq:v5v74-isometry-reflection}
\end{equation}
Moreover, a coarse variable in the positive half depends only on fine
fields and links in the positive half.  Therefore the observable pullback
induced by
\begin{equation}
 \phi_c=\iota_U^*\phi_f,
 \qquad U_c=r(U),
 \label{eq:v5v74-block-variables}
\end{equation}
is reflection-admissible.
\end{lemma}

\begin{proof}
Insert \eqref{eq:v5v74-reflected-block-data} into the block-map formula
\((\iota_U\phi)(x)=w_xT_U(x,b)\phi(b)\).  This proves
\eqref{eq:v5v74-isometry-reflection}.  A positive coarse cell and every
path used inside it lie entirely in the positive fine half by construction.
Thus a positive-time coarse cylinder pulls back to a positive-time fine
cylinder.  Products, grading, and reflection are preserved, which is
Definition~\ref{def:v5v74-reflection-block}.
\end{proof}

\begin{firewall}[Blocks may not straddle the reflection plane]
\label{fw:v5v74-no-straddling}
A symmetric average over a cell containing sites on both sides of the
reflection plane intertwines reflection but fails the first condition in
\eqref{eq:v5v74-reflection-admissible}: a positive coarse observable then
depends on negative-time fine variables.  The descent theorem does not
apply.  Alignment of the time partition is a structural requirement, not
a boundary error.
\end{firewall}

\subsection{The exact pseudofermion marginal}

At fixed gauge--Higgs background, let the normalized complex fine Gaussian
have positive precision \(Q_f\):
\begin{equation}
 d\nu_f(\phi_f)
 =\frac{\det Q_f}{\pi^{N_f}}
 e^{-\langle\phi_f,Q_f\phi_f\rangle},d\phi_f.
 \label{eq:v5v74-normalized-fine-Gaussian}
\end{equation}
Use the unitary wavelet coordinates
\(\phi_f=\mathscr U_U(\phi_c,\phi_h)\).  Write the transformed precision as
\begin{equation}
 \mathscr U_U^*Q_f\mathscr U_U
 =\begin{pmatrix}Q_{cc}&Q_{ch}\\Q_{hc}&Q_{hh}\end{pmatrix}.
 \label{eq:v5v74-Q-blocks}
\end{equation}

\begin{proposition}[Normalized Gaussian marginal and determinant ledger]
\label{prop:v5v74-Gaussian-marginal}
The pushforward of \(\nu_f\) under
\(\phi_c=\iota_U^*\phi_f\) is the normalized Gaussian
\begin{equation}
 d\nu_c(\phi_c)
 =\frac{\det Q_{\rm eff}}{\pi^{N_c}}
 e^{-\langle\phi_c,Q_{\rm eff}\phi_c\rangle},d\phi_c,
 \qquad
 Q_{\rm eff}=Q_{cc}-Q_{ch}Q_{hh}^{-1}Q_{hc}.
 \label{eq:v5v74-effective-Gaussian}
\end{equation}
For the Dirac block of V71,
\begin{equation}
 Q_{\rm eff}=S^*Z^{-1}S=D_{\rm cov}^*D_{\rm cov}.
 \label{eq:v5v74-Qeff-Dcov}
\end{equation}
For an unnormalized Gaussian integral, eliminating \(\phi_h\) additionally
produces the background-dependent factor
\begin{equation}
 \frac{\pi^{N_h}}{\det Q_{hh}}.
 \label{eq:v5v74-high-determinant}
\end{equation}
\end{proposition}

\begin{proof}
Complete the square:
\begin{align}
 \langle\phi,Q_f\phi\rangle
 &=\left\langle
 \phi_h+Q_{hh}^{-1}Q_{hc}\phi_c,
 Q_{hh}(\phi_h+Q_{hh}^{-1}Q_{hc}\phi_c)
 \right\rangle
 +\langle\phi_c,Q_{\rm eff}\phi_c\rangle.
 \label{eq:v5v74-complete-square}
\end{align}
Integration of the translated high Gaussian gives
\eqref{eq:v5v74-high-determinant}.  The determinant identity
\begin{equation}
 \det Q_f=\det Q_{hh}\det Q_{\rm eff}
 \label{eq:v5v74-Schur-determinant}
\end{equation}
cancels that factor in the normalized density, proving
\eqref{eq:v5v74-effective-Gaussian}.  Corollary
\ref{cor:v5v72-square-Schur} proves \eqref{eq:v5v74-Qeff-Dcov}.
\end{proof}

\begin{theorem}[Reflection positivity of the whitened Gaussian]
\label{thm:v5v74-whitened-RP}
Assume the joint or conditional fine pseudofermion state is reflection
positive on its declared patch, and use the one-sided block geometry of
Lemma~\ref{lem:v5v74-one-sided-block}.  Then the exact whitened coarse
Gaussian with precision \(D_{\rm cov}^*D_{\rm cov}\) is reflection
positive.  Arbitrarily long temporal couplings in that precision do not
alter the conclusion.
\end{theorem}

\begin{proof}
Proposition~\ref{prop:v5v74-Gaussian-marginal} identifies the whitened
Gaussian with the exact pushforward of the fine Gaussian.  Its observable
pullback is reflection-admissible by
Lemma~\ref{lem:v5v74-one-sided-block}.  Apply
Theorem~\ref{thm:v5v74-RP-descent}.
\end{proof}

For the full joint measure, the block is gauge dependent.  This causes no
Jacobian anomaly: with \(U\) held fixed, \(\mathscr U_U\) is unitary in
the pseudofermion variables, and the triangular change
\((U,\phi_f)\mapsto(U,\phi_c,\phi_h)\) has bosonic fibre Jacobian one.
Nevertheless, the factor \eqref{eq:v5v74-high-determinant} is a genuine
induced gauge--Higgs weight in an unnormalized functional integral.

\begin{firewall}[The high determinant may not be discarded]
\label{fw:v5v74-determinant-ledger}
Reflection positivity descends for the exact pushforward measure.  If one
keeps only \(e^{-\|D_{\rm cov}\phi_c\|^2}\) in an unnormalized
supertransfer and drops \((\det Q_{hh})^{-1}\), one has changed the
background law.  That determinant must remain in the Wilsonian action or
cancel by an explicitly proved regulator identity.
\end{firewall}

\subsection{Iteration and the reflection-positive limit}

Let \((\mathcal A_n,\Theta_n,\mathcal A_{n,+},\omega_n)\) be a sequence of
cutoff algebras and states, and let
\(\mathcal B_n:\mathcal A_n\to\mathcal A_{n+1}\) be the pullback from
level \(n\) to the finer level \(n+1\).

\begin{theorem}[Finite-depth exact RG preserves reflection positivity]
\label{thm:v5v74-iterated-RP}
Suppose every \(\mathcal B_n\) is reflection-admissible and
\begin{equation}
 \omega_n=\omega_{n+1}\circ\mathcal B_n.
 \label{eq:v5v74-projective-consistency}
\end{equation}
If one fine state \(\omega_N\) is reflection positive, then all exact
marginals \(\omega_n\), \(n<N\), are reflection positive.  If the family
is consistent for all \(n\), it defines a reflection-positive state on the
algebraic inductive limit of the cylinder algebras.
\end{theorem}

\begin{proof}
Repeated application of Theorem~\ref{thm:v5v74-RP-descent} proves the first
claim.  In the inductive limit, represent a cylinder by \(F\in\mathcal A_n\)
and define \(\omega([F])=\omega_n(F)\).  Consistency makes this independent
of the representative.  Every positive-time element has a representative
at a finite level, where its reflected square has nonnegative expectation;
hence the limit state is reflection positive.
\end{proof}

\begin{proposition}[Closedness under local entrywise limits]
\label{prop:v5v74-RP-closed}
Let \(\omega^{(j)}\) be reflection-positive states on a common cylinder
algebra.  If \(\omega^{(j)}(F)\to\omega(F)\) for every cylinder \(F\),
then \(\omega\) is reflection positive.
\end{proposition}

\begin{proof}
For fixed \(F\in\mathcal A_+\), each real number
\(\omega^{(j)}((\Theta F)F)\) is nonnegative.  Its limit is
\(\omega((\Theta F)F)\), which is therefore nonnegative.
\end{proof}

This proposition joins V74 to the V69--V70 entrywise convergence route.
Exact marginal states preserve positivity at every finite depth; a
convergent cofinal family preserves it in the limit without requiring a
uniform finite temporal range.

\subsection{What has and has not been closed}

\begin{assessment}[Pseudofermion reflection ledger]
\label{ass:v5v74-pseudofermion-ledger}
For the pseudofermion factor, the V67 positive fine-lattice gate, the V72
exact marginal identity, a one-sided covariant block, and the V74 descent
theorem form a closed finite-depth reflection argument.  Local convergence
then uses the whitened OLPC estimate rather than a new cross-plane
factorization of the perfect action.
\end{assessment}

\begin{firewall}[Exact marginal recursion is not yet the SM recursion]
\label{fw:v5v74-not-full-SM}
The joint Standard Model step also integrates physical chiral Grassmann
fields, gauge and Higgs shells, ghosts or gauge-fixing data, and the
background-dependent high determinant.  V74 does not prove that their
combined induced action lies in a cutoff-uniform contractive polymer
domain.  Nor does it supply the arbitrary-order rooted cumulant theorem
left open in V70.
\end{firewall}

\begin{firewall}[Continuum symmetries need separate convergence]
\label{fw:v5v74-no-continuum-symmetry}
An algebraic reflection-positive inductive-limit state is not by itself an
Osterwalder--Schrader continuum theory.  Euclidean covariance, regularity,
clustering, the spectral condition after reconstruction, and the Einstein
diffeomorphism constraints still require the corresponding cofinal
estimates.
\end{firewall}

\begin{theorem}[V74 conclusion]
\label{thm:v5v74-status}
Reflection positivity of the exact whitened pseudofermion marginal follows
directly from reflection positivity of the fine state when the covariant
block partition is time-one-sided and reflection compatible.  The result
is stable under arbitrary finite iteration and under local entrywise
limits, even though the effective precision has long exponentially
decaying temporal hoppings.  The high Gaussian determinant is part of the
induced background action and cannot be omitted.  The unresolved problem
has therefore moved from pseudofermion cross-plane signs to uniform control
of the complete induced SM--gravity shell action.
\end{theorem}

\begin{proof}
Combine Theorem~\ref{thm:v5v74-RP-descent},
Proposition~\ref{prop:v5v74-Gaussian-marginal},
Theorem~\ref{thm:v5v74-whitened-RP},
Theorem~\ref{thm:v5v74-iterated-RP}, and
Proposition~\ref{prop:v5v74-RP-closed}, subject to the determinant ledger
and the final two firewalls.
\end{proof}

V75 is the next compulsory YM/NS audit.  It will use the newest companion
files to choose the next attack on the induced high determinant and the
all-order rooted shell action.  If the companion programmes do not supply
an all-arity transport theorem, V75 will formulate the exact determinant
derivative expansion and identify the first unowned cumulant rather than
assuming a convergent polymer recursion.


\clearpage
\section{V75: companion audit and rooted trace-log control of the induced high determinant}
\label{sec:v5v75}

\subsection{Compulsory YM/NS audit}

The companion files were read passively before the present argument.  The
newest named Yang--Mills file at the time of inspection was
\begin{equation}
 \begin{split}
 &\texttt{Renewal\_Yang\_Mills\_Mass\_Gap\_Research\_Notes\_V88.tex},\\
 &1\,307\,549\text{ bytes},\qquad
 \operatorname{SHA256}=\\
 &\texttt{6B6795D5940C05088C75D1A7AE1DF34991982F5C106DC26BECBB459223963D3E}.
 \end{split}
 \label{eq:v5v75-YM-fingerprint}
\end{equation}
The shared folder also contained its V87 predecessor and frozen YM
archives; none was altered.  V88 proves an exact subset recursion for
fully symmetrized centered resolvent chains and shows that a single SPCC
would imply the all-arity Pruefer tree card.  The SPCC remains open even
at its first two-label column estimate.  Its decisive methodological
lesson is that permutation cancellation must occur before taking norms;
ordinary triangle iteration loses every required incidence power.

The Navier--Stokes file remained
\begin{equation}
 \begin{split}
 &\texttt{Renewal\_NS\_Stability\_Research\_Notes\_V31.tex},\\
 &887\,537\text{ bytes},\qquad
 \operatorname{SHA256}=\\
 &\texttt{F1121B62238AD5491BB7CF03635EA9934942EDF573ED8FAAA9EE79E1D38A6696}.
 \end{split}
 \label{eq:v5v75-NS-fingerprint}
\end{equation}
Its relevant conclusion is unchanged: a probability mark can be exact on
every dyadic annulus, while removing the mark costs the missing critical
first moment.  The Duhamel lift pays that moment only in the critical
source norm; estimating a weaker source simply restores the missing scale
weight elsewhere.  Thus a global total must be rooted at the physical
entry before it is summed.

These lessons meet at the determinant left by V74.  A global determinant
is extensive, and its differentiated permutation chains have factorial
multiplicity.  We therefore root the exact trace logarithm at one lattice
cell and resum each closed loop before taking absolute values.  This gives
an all-order result for the Gaussian high determinant without asserting
the unproved YM SPCC for non-Gaussian Wilson cumulants.

\subsection{Relative high precision and the trace logarithm}

Let \(Q_0>0\) be a local reference precision on the high wavelet sector,
and write the actual high precision as
\begin{equation}
 Q_h=Q_0+V=Q_0(I+X),
 \qquad X:=Q_0^{-1}V.
 \label{eq:v5v75-relative-high}
\end{equation}
Although \(X\) need not be self-adjoint, it is similar to the self-adjoint
matrix \(Q_0^{-1/2}VQ_0^{-1/2}\).  Work on the block graph and assume
\begin{equation}
 \|X\|_{\mu,1}\le\kappa<1.
 \label{eq:v5v75-kappa}
\end{equation}
This is a relative high-shell hypothesis, not a claim that the full
fine-to-coarse action is small.

\begin{lemma}[Exact determinant ratio]
\label{lem:v5v75-determinant-series}
Under \eqref{eq:v5v75-kappa}, the principal logarithm satisfies
\begin{equation}
 \log\frac{\det Q_h}{\det Q_0}
 =\operatorname{Tr}\log(I+X)
 =\sum_{m\ge1}\frac{(-1)^{m+1}}m\operatorname{Tr}X^m,
 \label{eq:v5v75-trace-log}
\end{equation}
and the series is absolutely convergent at every finite volume.
\end{lemma}

\begin{proof}
The weighted Schur norm dominates operator norm, so
\(\|X\|_{\rm op}<1\).  The power series for \(\log(I+X)\) converges in
operator norm.  Taking the finite-dimensional trace term by term is valid.
The determinant identity follows from
\(Q_h=Q_0(I+X)\) and the principal logarithm, whose spectrum stays in the
disk centered at one of radius \(\kappa\), disjoint from the logarithm cut.
\end{proof}

\subsection{Rooted closed-loop density}

Let the internal fibre dimension at a block be at most \(d_*\).  Define
the trace-log density rooted at block \(x\) by
\begin{equation}
 \ell_x(X)
 :=\sum_{m\ge1}\frac{(-1)^{m+1}}m
 \operatorname{tr}_{\rm int}(X^m)_{xx}.
 \label{eq:v5v75-rooted-density}
\end{equation}

\begin{theorem}[Volume-uniform all-order rooted trace-log bound]
\label{thm:v5v75-rooted-trace-log}
Under \eqref{eq:v5v75-kappa}, the series
\eqref{eq:v5v75-rooted-density} converges absolutely and
\begin{equation}
 \sup_x|\ell_x(X)|
 \le d_*\sum_{m\ge1}\frac{\kappa^m}{m}
 =d_*\log\frac1{1-\kappa}.
 \label{eq:v5v75-rooted-bound}
\end{equation}
Moreover,
\begin{equation}
 \operatorname{Tr}\log(I+X)=\sum_x\ell_x(X).
 \label{eq:v5v75-density-sum}
\end{equation}
The bound is independent of total volume; the final sum is extensive only
because it deliberately sums one uniformly controlled density over all
roots.
\end{theorem}

\begin{proof}
Expand the diagonal block as a sum over closed paths:
\begin{equation}
 (X^m)_{xx}
 =\sum_{x_1,\ldots,x_{m-1}}
 X_{xx_1}X_{x_1x_2}\cdots X_{x_{m-1}x}.
 \label{eq:v5v75-closed-paths}
\end{equation}
The row part of the weighted Schur algebra estimate, with the exponential
weights discarded only after multiplication, bounds the sum of norms in
\eqref{eq:v5v75-closed-paths} by \(\kappa^m\).  The internal trace is at
most \(d_*\) times the block norm.  Summation over \(m\) proves
\eqref{eq:v5v75-rooted-bound}.  Absolute convergence permits summation
over the finite set of roots, which reconstructs the trace in
\eqref{eq:v5v75-trace-log}.
\end{proof}

Every term in \(\ell_x\) is a closed high-mode propagation loop.  Under a
gauge transformation, its diagonal block is conjugated at \(x\), so the
internal trace is gauge invariant.  Thus the rooted density is locally
gauge invariant without fixing a wavelet basis.

\begin{assessment}[The NS first-moment analogue]
\label{ass:v5v75-root-before-sum}
The determinant itself corresponds to the extensive sum in
\eqref{eq:v5v75-density-sum}.  The local continuum datum is
\(\ell_x\), whose bound \eqref{eq:v5v75-rooted-bound} is obtained before
that sum.  Estimating \(|\operatorname{Tr}\log(I+X)|\) first would insert a
volume factor and then offer no legitimate operation which removes it.
This is the determinant analogue of the NS instruction to retain the
physical mark until its exact first moment is assigned.
\end{assessment}

\subsection{Stability and distant-shell decay}

\begin{theorem}[Local Lipschitz stability of the resummed determinant]
\label{thm:v5v75-trace-log-stability}
If
\begin{equation}
 \|X\|_{\mu,1},\ \|Y\|_{\mu,1}\le\kappa<1,
 \qquad
 \|X-Y\|_{\mu,1}\le\delta,
 \label{eq:v5v75-two-kernels}
\end{equation}
then
\begin{equation}
 \sup_x|\ell_x(X)-\ell_x(Y)|
 \le\frac{d_*\delta}{1-\kappa}.
 \label{eq:v5v75-density-Lipschitz}
\end{equation}
\end{theorem}

\begin{proof}
Use the noncommutative telescoping identity
\begin{equation}
 X^m-Y^m
 =\sum_{j=0}^{m-1}X^j(X-Y)Y^{m-1-j}.
 \label{eq:v5v75-power-telescope}
\end{equation}
The weighted algebra bounds its norm by
\(m\kappa^{m-1}\delta\).  After division by \(m\), summation of the
geometric series and the internal trace bound give
\eqref{eq:v5v75-density-Lipschitz}.
\end{proof}

\begin{proposition}[Two-leg decay for a remote determinant change]
\label{prop:v5v75-remote-trace-log}
Let \(A\) be a fixed set of roots and suppose \(X-Y\) is supported on
\(S\times S\), with \(d(A,S)\ge R\).  Assume pointwise exponential kernel
bounds compatible with \eqref{eq:v5v75-two-kernels}.  For every exponent
\(\mu'<\mu\) above the graph-growth threshold, there is a volume-uniform
constant \(C\) such that
\begin{equation}
 \sup_{x\in A}|\ell_x(X)-\ell_x(Y)|
 \le
 C\frac{\delta}{(1-\kappa)^2}e^{-2\mu'R}.
 \label{eq:v5v75-remote-density}
\end{equation}
\end{proposition}

\begin{proof}
Insert \eqref{eq:v5v75-power-telescope} into the rooted diagonal.  Every
closed path contributing to a difference contains one insertion supported
in \(S\).  Starting at \(x\in A\), the path must reach \(S\) before that
insertion and return to \(x\) afterwards.  These two legs have total length
at least \(2R\).  Spend \(\mu-\mu'\) to sum graph growth on the two legs.
The remaining portions on either side of the insertion give two geometric
sums in \(\kappa\), bounded by \((1-\kappa)^{-2}\).  The insertion gives
\(\delta\), proving \eqref{eq:v5v75-remote-density}.
\end{proof}

\subsection{Symmetry-first derivative identity}

The exact resummation above is preferable when
\eqref{eq:v5v75-kappa} holds.  For comparison with the YM audit, its mixed
Taylor coefficients admit an exact symmetry-first formula.  Let
\begin{equation}
 Q(t)=Q_0+\sum_{i=1}^mt_iV_i,
 \qquad R_0=Q_0^{-1}.
 \label{eq:v5v75-affine-precision}
\end{equation}

\begin{theorem}[Symmetrized cyclic determinant derivative]
\label{thm:v5v75-cyclic-derivative}
For distinct formal sources \(t_1,\ldots,t_m\),
\begin{equation}
 \left.\partial_{t_1}\cdots\partial_{t_m}
 \log\det Q(t)\right|_{t=0}
 =\frac{(-1)^{m-1}}m
 \sum_{\pi\in S_m}
 \operatorname{Tr}
 \bigl(R_0V_{\pi(1)}\cdots R_0V_{\pi(m)}\bigr).
 \label{eq:v5v75-cyclic-derivative}
\end{equation}
The permutation sum must be assembled before a norm is taken.
\end{theorem}

\begin{proof}
Expand
\begin{equation}
 \log\det Q(t)-\log\det Q_0
 =\sum_{r\ge1}\frac{(-1)^{r+1}}r
 \operatorname{Tr}\left(\sum_it_iR_0V_i\right)^r.
\end{equation}
Only \(r=m\) contains the monomial \(t_1\cdots t_m\).  Its coefficient is
the sum over all permutations displayed in
\eqref{eq:v5v75-cyclic-derivative}; differentiating the monomial introduces
no further factor.
\end{proof}

Cyclicity makes many summands identical, but applying the triangle
inequality to the individual words reinstates factorial growth.  The
rooted trace-log theorem avoids that loss by summing the complete power
\(X^m\) with its \(1/m\) coefficient.  This does not prove the YM SPCC,
whose non-Gaussian projected tails lack a single small operator power.

\subsection{Insertion into the exact block marginal}

The unnormalized high Gaussian factor from V74 is
\((\det Q_h)^{-1}\).  Relative to the reference high Gaussian it contributes
the induced action
\begin{equation}
 \Delta S_h
 =\log\frac{\det Q_h}{\det Q_0}
 =\sum_x\ell_x(X).
 \label{eq:v5v75-induced-action}
\end{equation}

\begin{corollary}[A determinant-sector all-order local certificate]
\label{cor:v5v75-determinant-certificate}
Suppose along the cutoff ancestry:
\begin{enumerate}
 \item \(Q_0^{-1}\) has a common weighted locality bound;
 \item the locally matched relative high defect satisfies
 \(\|Q_0^{-1}(Q_h-Q_0)\|_{\mu,1}\le\kappa<1\);
 \item adjacent-cutoff defects are weighted-small near each fixed root or
 are supported at a distance tending to infinity.
\end{enumerate}
Then the local densities of \(\Delta S_h\) are uniformly bounded and form
a Cauchy family at every fixed root.  All powers of the high determinant
are included; no fixed-arity truncation is used.
\end{corollary}

\begin{proof}
Uniform boundedness is Theorem~\ref{thm:v5v75-rooted-trace-log}.
Weighted-small changes are controlled by
Theorem~\ref{thm:v5v75-trace-log-stability}; remote changes are controlled
by Proposition~\ref{prop:v5v75-remote-trace-log}.  Apply the assumed
Cauchy conditions.
\end{proof}

\begin{firewall}[The relative high defect has not been made small]
\label{fw:v5v75-kappa-open}
The certificate is conditional on \(\kappa<1\).  A high spectral gap gives
a bound on \(Q_0^{-1}\), but it does not make
\(Q_0^{-1}(Q_h-Q_0)\) small.  The reference must absorb the order-one local
high symbol and its relevant gauge--Higgs response before the remaining
relative defect is tested.
\end{firewall}

\begin{firewall}[Local determinant control is not the full shell card]
\label{fw:v5v75-not-full-shell}
Corollary~\ref{cor:v5v75-determinant-certificate} closes the all-order
trace-log series of one Gaussian high block under a small relative
operator hypothesis.  It does not prove the YM SPCC, the physical
Grassmann rooted bounds, or convergence of mixed gauge--Higgs--gravity
cumulants.  Products between the induced density and the other shell
factors still require the measure-first assembly of V49 and the RSCC of
V69.
\end{firewall}

\begin{theorem}[V75 conclusion]
\label{thm:v5v75-status}
The compulsory audit records YM V88 and unchanged NS V31.  Their shared
instruction---assemble the physical symmetry or root before taking a
norm---leads to a volume-uniform diagonal trace-loop density for the high
pseudofermion determinant.  If the locally matched relative high precision
has weighted norm \(\kappa<1\), the exact trace logarithm converges to all
orders, is locally Lipschitz, and gains two exponential legs for remote
changes.  This supplies a determinant-sector local certificate while
leaving the quantitative construction of the matched reference and the
non-Gaussian all-order shell card open.
\end{theorem}

\begin{proof}
Use Theorem~\ref{thm:v5v75-rooted-trace-log},
Theorem~\ref{thm:v5v75-trace-log-stability},
Proposition~\ref{prop:v5v75-remote-trace-log}, and
Corollary~\ref{cor:v5v75-determinant-certificate}, subject to the two
firewalls.
\end{proof}

V76 will construct the optimal local reference for the free high Haar
block.  It will compute the relative precision after matching the full
zero-field high symbol, determine whether the gauge-linear remainder has
weighted norm below one on a cutoff-uniform chart, and assign any
order-one gauge--Higgs jets to their induced local couplings before the
trace-log certificate is invoked.


\clearpage
\section{V76: jet-matched high-sector references and a quadratic relative defect}
\label{sec:v5v76}

\subsection{Covariant local coordinates for the background patch}

V75 reduced the induced high determinant to the relative operator
\(X=Q_0^{-1}(Q_h-Q_0)\) and required
\(\|X\|_{\mu,1}<1\).  A spectral gap alone does not provide this
smallness.  The reference must first contain the order-one free high symbol
and the nonvanishing local background jets.

Work in one contractible small-field chart of the V60 projective atlas.
After its local gauge pinning, let \(z\) collect based link holonomies,
Higgs fluctuations in a polar chart, and any declared metric/tetrad
background coordinates.  Equip these coordinates with a local norm
\(\|z\|_{\mathcal Z}\).  Transition maps act equivariantly, so a kernel
constructed tensorially from the jets below descends to the unpinned chart.

Let
\begin{equation}
 Q_h:\overline B_{\mathcal Z}(0,\eta_*)
 \longrightarrow\mathfrak A_{\mu,1}
 \label{eq:v5v76-Qh-map}
\end{equation}
be the high pseudofermion precision in the weighted kernel algebra.  Assume
it is twice continuously Fr\'echet differentiable and self-adjoint on real
backgrounds, with
\begin{align}
 Q_h(z)&\ge g_h I,
 \label{eq:v5v76-high-positive}\\
 \sup_{\|z\|\le\eta_*}
 \|D^2Q_h(z)\|_{\mathcal Z^2\to\mathfrak A_{\mu,1}}
 &\le M_2.
 \label{eq:v5v76-second-derivative}
\end{align}
The V59 cluster calculus is intended to supply
\eqref{eq:v5v76-second-derivative}; it is kept explicit here.

\subsection{The unique first-jet match}

Define the affine jet reference
\begin{equation}
 Q_0^{[1]}(z):=Q_h(0)+DQ_h(0)[z].
 \label{eq:v5v76-affine-reference}
\end{equation}
The zeroth term is the full free high symbol, represented by its
exponentially local perfect stencil as in V73.  The linear term contains
the complete gauge, Higgs, and geometric one-background vertices in the
same covariant chart.

\begin{theorem}[Quadratic jet remainder]
\label{thm:v5v76-quadratic-remainder}
Under \eqref{eq:v5v76-second-derivative},
\begin{equation}
 \|Q_h(z)-Q_0^{[1]}(z)\|_{\mu,1}
 \le\frac{M_2}{2}\|z\|_{\mathcal Z}^2.
 \label{eq:v5v76-quadratic-remainder}
\end{equation}
If
\begin{equation}
 \frac{M_2}{2}\eta_*^2<g_h,
 \label{eq:v5v76-reference-positive-radius}
\end{equation}
then \(Q_0^{[1]}(z)>0\) throughout the chart and
\begin{equation}
 \|(Q_0^{[1]}(z))^{-1}\|_{\rm op}
 \le\left(g_h-\frac{M_2}{2}\eta_*^2\right)^{-1}.
 \label{eq:v5v76-reference-op-inverse}
\end{equation}
\end{theorem}

\begin{proof}
The Banach-space Taylor formula with integral remainder gives
\begin{equation}
 Q_h(z)-Q_h(0)-DQ_h(0)[z]
 =\int_0^1(1-t)D^2Q_h(tz)[z,z],dt,
 \label{eq:v5v76-Taylor-integral}
\end{equation}
which proves \eqref{eq:v5v76-quadratic-remainder}.  The operator norm is
bounded by the weighted norm.  Hence
\(\|Q_h(z)-Q_0^{[1]}(z)\|_{\rm op}\le M_2\eta_*^2/2\), and Weyl's
spectral inequality with \eqref{eq:v5v76-high-positive} proves positivity
and \eqref{eq:v5v76-reference-op-inverse}.
\end{proof}

\begin{proposition}[Optimality among affine references]
\label{prop:v5v76-affine-optimality}
Let \(\widetilde Q_0(z)=Q_h(0)+L[z]\) for a bounded linear map \(L\).
Then
\begin{equation}
 \|Q_h(z)-\widetilde Q_0(z)\|_{\mu,1}=o(\|z\|_{\mathcal Z})
 \quad(z\to0)
 \label{eq:v5v76-little-o}
\end{equation}
if and only if \(L=DQ_h(0)\).  Thus
\eqref{eq:v5v76-affine-reference} is the unique affine reference with a
quadratic relative defect.
\end{proposition}

\begin{proof}
Fr\'echet differentiability gives
\begin{equation}
 Q_h(z)-\widetilde Q_0(z)
 =(DQ_h(0)-L)[z]+o(\|z\|_{\mathcal Z}).
\end{equation}
If \eqref{eq:v5v76-little-o} holds, set \(z=tv\), divide by \(|t|\), and
let \(t\to0\); then \((DQ_h(0)-L)[v]=0\) for every \(v\).  The converse is
the definition of differentiability, with the quadratic bound supplied by
Theorem~\ref{thm:v5v76-quadratic-remainder}.
\end{proof}

\subsection{Weighted inverse control and the V75 smallness gate}

Operator positivity does not by itself bound the weighted inverse.  Assume
the V59 Combes--Thomas hypotheses for \(Q_0^{[1]}(z)\), uniformly on the
smaller chart \(\|z\|\le\eta_1\), so that for some
\(0<\mu'<\mu\),
\begin{equation}
 \sup_{\|z\|\le\eta_1}
 \|(Q_0^{[1]}(z))^{-1}\|_{\mu',1}\le K_h.
 \label{eq:v5v76-weighted-inverse}
\end{equation}

\begin{theorem}[Quadratic relative high defect]
\label{thm:v5v76-relative-defect}
For \(\|z\|_{\mathcal Z}\le\eta_1\), set
\begin{equation}
 X^{[1]}(z)
 =(Q_0^{[1]}(z))^{-1}
  (Q_h(z)-Q_0^{[1]}(z)).
 \label{eq:v5v76-X1}
\end{equation}
Then
\begin{equation}
 \|X^{[1]}(z)\|_{\mu',1}
 \le\frac{K_hM_2}{2}\|z\|_{\mathcal Z}^2.
 \label{eq:v5v76-X1-bound}
\end{equation}
In particular, the V75 trace-log certificate holds uniformly on every
chart radius satisfying
\begin{equation}
 \frac{K_hM_2}{2}\eta_1^2\le\kappa<1.
 \label{eq:v5v76-kappa-radius}
\end{equation}
\end{theorem}

\begin{proof}
Apply the weighted algebra inequality to
\eqref{eq:v5v76-X1}, then use
\eqref{eq:v5v76-quadratic-remainder} and
\eqref{eq:v5v76-weighted-inverse}.  Condition
\eqref{eq:v5v76-kappa-radius} is precisely
\eqref{eq:v5v75-kappa}.
\end{proof}

\begin{corollary}[Quadratic bound for the residual determinant density]
\label{cor:v5v76-residual-density}
Let \(\ell_x^{[1]}(z)=\ell_x(X^{[1]}(z))\) be the relative determinant
density of V75.  Under \eqref{eq:v5v76-kappa-radius},
\begin{equation}
 \sup_x|\ell_x^{[1]}(z)|
 \le d_*\log\!\left(
 \frac1{1-K_hM_2\|z\|_{\mathcal Z}^2/2}
 \right)
 =O(\|z\|_{\mathcal Z}^2).
 \label{eq:v5v76-density-quadratic}
\end{equation}
\end{corollary}

\begin{proof}
Insert \eqref{eq:v5v76-X1-bound} into
Theorem~\ref{thm:v5v75-rooted-trace-log}.
\end{proof}

\subsection{Higher jets and scale gain}

If \(Q_h\) has \(p+1\) uniformly bounded derivatives, define
\begin{equation}
 Q_0^{[p]}(z)
 =\sum_{j=0}^p\frac1{j!}D^jQ_h(0)[z^{\otimes j}].
 \label{eq:v5v76-p-reference}
\end{equation}

\begin{proposition}[Order-(p+1) relative defect]
\label{prop:v5v76-higher-jet}
If
\begin{equation}
 \sup_{\|z\|\le\eta_1}
 \|D^{p+1}Q_h(z)\|le M_{p+1},
 \qquad
 \|(Q_0^{[p]}(z))^{-1}\|_{\mu',1}\le K_p,
 \label{eq:v5v76-higher-hypotheses}
\end{equation}
then
\begin{equation}
 \left\|(Q_0^{[p]})^{-1}(Q_h-Q_0^{[p]})
 \right\|_{\mu',1}
 \le\frac{K_pM_{p+1}}{(p+1)!}
 \|z\|_{\mathcal Z}^{p+1}.
 \label{eq:v5v76-higher-defect}
\end{equation}
\end{proposition}

\begin{proof}
Use the order-\(p\) Taylor formula with integral remainder and the weighted
algebra inequality.
\end{proof}

For a smooth continuum family with \(\|z(a)\|=O(a^\alpha)\), the residual
relative determinant begins at
\(O(a^{\alpha(p+1)})\).  Increasing \(p\) is legitimate only when the
corresponding background vertices belong to the declared local action
space and their derivative constants do not destroy the factorial gain.

\subsection{Ownership and exact covariance}

The jet reference is equivariant because differentiation of the identity
\begin{equation}
 Q_h(g\cdot z)=G(g)Q_h(z)G(g)^{-1}
 \label{eq:v5v76-Q-covariance}
\end{equation}
produces equivariant multilinear jets in the pinned chart.  Its scalar
determinant contributions are therefore gauge invariant after the complete
SM representation sum.  The zeroth and first jets have explicit owners:
\begin{enumerate}
 \item \(Q_h(0)\) is the full free high perfect symbol, not an on-site
 spectator approximation;
 \item \(DQ_h(0)\) contains the induced one-background vertices and must be
 carried by the regulator reference or assigned to the matching gauge,
 Higgs, Yukawa, and geometric couplings.
\end{enumerate}

\begin{firewall}[Reference matching moves an order-one determinant]
\label{fw:v5v76-reference-determinant}
The identity
\begin{equation}
 \det Q_h=\det Q_0^{[1]}\det(I+X^{[1]})
 \label{eq:v5v76-determinant-split}
\end{equation}
makes the residual trace logarithm small, but
\(\det Q_0^{[1]}\) remains.  It must be retained as a reference Gaussian,
cancelled by the conjugate regulator through an exact block identity, or
entered in the induced action.  Declaring only the second factor to be
small does not remove the first.
\end{firewall}

\begin{firewall}[Coordinate Taylor order is not engineering irrelevance]
\label{fw:v5v76-Taylor-not-RG}
A quadratic remainder in the amplitude of a smooth background is not
automatically irrelevant under renormalization.  Quadratic gauge and Higgs
terms include marginal kinetic and relevant mass operators.  Their
engineering projection and normalization conditions must still be applied
to the rooted density coefficients.
\end{firewall}

\begin{theorem}[V76 conclusion]
\label{thm:v5v76-status}
Matching the full free high precision and its first covariant background
jet is the unique affine reference choice that turns the high relative
defect from linear to quadratic order.  Uniform second-derivative and
weighted inverse bounds give the explicit V75 smallness constant
\(K_hM_2\eta_1^2/2<1\), and the residual rooted determinant density is
quadratic on that chart.  The matched reference determinant and the
relevant parts of its higher background jets remain to be assigned in the
phase-preserving supertransfer ledger.
\end{theorem}

\begin{proof}
Use Theorem~\ref{thm:v5v76-quadratic-remainder},
Proposition~\ref{prop:v5v76-affine-optimality},
Theorem~\ref{thm:v5v76-relative-defect}, and
Corollary~\ref{cor:v5v76-residual-density}, with the two firewalls.
\end{proof}

V77 will perform that exact ledger.  It will block the conjugate
Grassmann determinant and the positive pseudofermion square simultaneously,
track the whitening determinant \(\det Z\), and prove which high and low
factors cancel before any absolute estimate is taken.


\clearpage
\section{V77: exact block factorization of the phase-preserving supertransfer}
\label{sec:v5v77}

\subsection{Why the determinants must be blocked together}

V66 represented the inverse determinant without discarding its phase:
\begin{equation}
 \frac1{\det D}
 =\frac{\det D^*}{\det(D^*D)}.
 \label{eq:v5v77-supertransfer}
\end{equation}
V72 then integrated the high pseudofermion variable and produced the
positive whitening factor \(Z\).  Treating that factor as an independent
counterterm obscures an exact cancellation.  This note blocks the
conjugate Grassmann determinant and the positive pseudofermion square in
the same low--high decomposition before any norm is taken.

Use the V71 notation
\begin{equation}
 D=\begin{pmatrix}A&K\\L&E\end{pmatrix},
 \qquad
 S=A-KE^{-1}L,
 \qquad
 M=KE^{-1},
 \qquad
 Z=I+MM^*.
 \label{eq:v5v77-block-data}
\end{equation}
Assume \(D\), \(E\), and \(S\) are invertible.  Put \(Q=D^*D\) and
write \(Q_{hh}\) for its high--high block.

\subsection{The three determinant identities}

\begin{lemma}[Dirac and high-square factors]
\label{lem:v5v77-three-factors}
One has
\begin{align}
 \det D&=\det E\det S,
 \label{eq:v5v77-det-D}\\
 Q_{hh}&=K^*K+E^*E,
 \label{eq:v5v77-Qhh}\\
 \det Q_{hh}&=|\det E|^2\det Z.
 \label{eq:v5v77-det-Qhh}
\end{align}
Moreover, the positive-square Schur complement is
\begin{equation}
 Q_{\rm eff}=Q/Q_{hh}=S^*Z^{-1}S,
 \qquad
 \det Q_{\rm eff}=\frac{|\det S|^2}{\det Z}.
 \label{eq:v5v77-Qeff-determinant}
\end{equation}
\end{lemma}

\begin{proof}
The triangular factorization
\eqref{eq:v5v71-triangular-factorization} gives
\eqref{eq:v5v77-det-D}.  Direct multiplication gives
\eqref{eq:v5v77-Qhh}.  Factor it as
\begin{equation}
 Q_{hh}
 =E^*\left(I+E^{-*}K^*KE^{-1}\right)E.
 \label{eq:v5v77-Qhh-factor}
\end{equation}
Sylvester's identity \(\det(I+AB)=\det(I+BA)\) gives
\begin{equation}
 \det\left(I+E^{-*}K^*KE^{-1}\right)
 =\det(I+KE^{-1}E^{-*}K^*)=\det Z,
\end{equation}
which proves \eqref{eq:v5v77-det-Qhh}.  The operator identity in
\eqref{eq:v5v77-Qeff-determinant} is
Corollary~\ref{cor:v5v72-square-Schur}; taking determinants proves its
second part.
\end{proof}

\begin{corollary}[Exact cancellation inside the positive square]
\label{cor:v5v77-square-cancellation}
The whitening determinant cancels between the high block and the retained
Schur precision:
\begin{equation}
 \det Q_{hh}\det Q_{\rm eff}
 =|\det E|^2|\det S|^2
 =|\det D|^2.
 \label{eq:v5v77-square-cancellation}
\end{equation}
\end{corollary}

\begin{proof}
Multiply \eqref{eq:v5v77-det-Qhh} and
\eqref{eq:v5v77-Qeff-determinant}, then use
\eqref{eq:v5v77-det-D}.
\end{proof}

\subsection{Natural and symmetric allocations}

High Grassmann elimination of \(D^*\) gives \(\det E^*\), while the
retained conjugate operator is \(S^*\).  Combining these with the two
pseudofermion Gaussian factors yields the natural allocation
\begin{align}
 \mathcal H_{\rm nat}
 &:=\frac{\det E^*}{\det Q_{hh}}
 =\frac1{\det E\det Z},
 \label{eq:v5v77-high-natural}\\
 \mathcal L_{\rm nat}
 &:=\frac{\det S^*}{\det Q_{\rm eff}}
 =\frac{\det Z}{\det S}.
 \label{eq:v5v77-low-natural}
\end{align}

\begin{theorem}[Exact phase-preserving block product]
\label{thm:v5v77-block-supertransfer}
The high and low natural factors satisfy
\begin{equation}
 \mathcal H_{\rm nat}\mathcal L_{\rm nat}
 =\frac1{\det E\det S}
 =\frac1{\det D}.
 \label{eq:v5v77-natural-product}
\end{equation}
In particular, every power of \(\det Z\) cancels exactly in the combined
product.
\end{theorem}

\begin{proof}
Substitute \eqref{eq:v5v77-high-natural} and
\eqref{eq:v5v77-low-natural}; then use
\eqref{eq:v5v77-det-D}.
\end{proof}

The whitened marginal Dirac operator is
\begin{equation}
 D_{\rm cov}=Z^{-1/2}S,
 \qquad
 Q_{\rm eff}=D_{\rm cov}^*D_{\rm cov}.
 \label{eq:v5v77-Dcov}
\end{equation}
It is often preferable to pair its conjugate determinant directly with
its square.  Define the symmetric allocation
\begin{align}
 \mathcal H_{1/2}
 &:=\frac1{\det E\det Z^{1/2}},
 \label{eq:v5v77-high-half}\\
 \mathcal L_{1/2}
 &:=\frac{\det D_{\rm cov}^*}
          {\det(D_{\rm cov}^*D_{\rm cov})}
 =\frac1{\det D_{\rm cov}}
 =\frac{\det Z^{1/2}}{\det S}.
 \label{eq:v5v77-low-half}
\end{align}

\begin{proposition}[Local half-density redistribution]
\label{prop:v5v77-half-density}
The natural and symmetric allocations have the same product:
\begin{equation}
 \mathcal H_{1/2}\mathcal L_{1/2}
 =\mathcal H_{\rm nat}\mathcal L_{\rm nat}
 =\frac1{\det D}.
 \label{eq:v5v77-half-product}
\end{equation}
The redistribution is realized in the retained conjugate Grassmann
integral by the change of variable
\begin{equation}
 \chi=Z^{-1/2}\chi',
 \qquad
 d\chi=(\det Z^{1/2}),d\chi',
 \label{eq:v5v77-Grassmann-change}
\end{equation}
which changes the retained operator from \(S^*\) to
\(S^*Z^{-1/2}=D_{\rm cov}^*\) and leaves its integral equal to
\(\det S^*\).
\end{proposition}

\begin{proof}
The first identity follows immediately from
\eqref{eq:v5v77-high-half}--\eqref{eq:v5v77-low-half}.  For Grassmann
variables, a linear change \(\chi=A\chi'\) changes the Berezin measure by
\((\det A)^{-1}\).  With \(A=Z^{-1/2}\), the Jacobian is
\(\det Z^{1/2}\), and
\begin{equation}
 (\det Z^{1/2})\det(S^*Z^{-1/2})=\det S^*.
\end{equation}
\end{proof}

\subsection{Phase and locality ledger}

\begin{proposition}[The whitening factor carries no phase]
\label{prop:v5v77-Z-no-phase}
Since \(Z=I+MM^*>0\),
\begin{equation}
 \det Z>0,
 \qquad
 \arg\det D_{\rm cov}=\arg\det S.
 \label{eq:v5v77-phase-ledger}
\end{equation}
Thus the phase of the inverse determinant splits only between
\(\det E\) and \(\det S\); whitening neither creates nor cancels an
anomalous phase.
\end{proposition}

\begin{proof}
All eigenvalues of \(Z\) are positive, so its principal square root has
positive determinant.  Equation \eqref{eq:v5v77-Dcov} gives
\(\det D_{\rm cov}=(\det Z)^{-1/2}\det S\), proving the phase identity.
\end{proof}

If \(M\) is uniformly local, V72 makes \(Z^{\pm1/2}\) uniformly local in
a slightly weaker exponent.  Consequently both allocations use local
operators.  Their determinant densities must nevertheless be assembled in
the exact product before applying absolute estimates: separately bounding
\(\log\det Z\) in the high and low factors would pay the same local loops
twice and erase the algebraic cancellation.

\begin{corollary}[Reference determinants must be paired simultaneously]
\label{cor:v5v77-reference-pairing}
Let \(E_0\), \(Z_0\), and \(S_0\) be matched references.  The relative
supertransfer factor is
\begin{equation}
 \frac{\det D_0}{\det D}
 =\frac{\det E_0}{\det E}
  \frac{\det S_0}{\det S},
 \label{eq:v5v77-relative-supertransfer}
\end{equation}
with no \(Z\)-ratio.  Equivalently, in the half-density allocation the two
opposite \(Z^{1/2}\)-ratios cancel before the rooted trace-log norm is
taken.
\end{corollary}

\begin{proof}
Apply \eqref{eq:v5v77-det-D} to \(D\) and \(D_0\).  The equivalent
statement follows from Proposition~\ref{prop:v5v77-half-density} on both
backgrounds.
\end{proof}

This refines the V76 reference firewall.  The positive reference precision
does not disappear, but its whitening part has an exact owner on both sides
of the block split.  The genuinely remaining high reference is the complex
Dirac factor \(\det E\), while the retained recursion is governed by
\(D_{\rm cov}\) or, equivalently, \(S\) together with its positive local
normalization.

\subsection{Reflection and anomaly firewalls}

\begin{firewall}[A local Grassmann change need not be time-one-sided]
\label{fw:v5v77-half-density-RP}
Although \(Z^{-1/2}\) is exponentially local, it can connect the two sides
of the reflection plane.  The variable change
\eqref{eq:v5v77-Grassmann-change} is therefore an algebraic determinant
redistribution, not by itself a reflection-admissible observable block.
Reflection positivity should continue to be inherited from the exact
one-sided marginal as in V74.
\end{firewall}

\begin{firewall}[The high Dirac phase remains]
\label{fw:v5v77-high-phase}
The factor \((\det E)^{-1}\) contains the phase of the eliminated high
Dirac block.  Positivity of \(Q_{hh}\), locality of \(E^{-1}\), and
positivity of \(Z\) do not trivialize that phase.  Its complete Standard
Model representation sum must satisfy the inherited local anomaly and
large-holonomy gates before it can be absorbed into a gauge-invariant
induced action.
\end{firewall}

\begin{theorem}[V77 conclusion]
\label{thm:v5v77-status}
Simultaneous low--high factorization of the conjugate Grassmann determinant
and positive pseudofermion square gives an exact phase-preserving product.
The whitening determinant \(\det Z\) cancels between high and low sectors;
it may also be split symmetrically by a local half-density Grassmann
change.  Because \(Z>0\), this redistribution carries no determinant
phase.  The unresolved induced factor is the high complex Dirac determinant
\((\det E)^{-1}\), whose local modulus can use the V75--V76 reference
method but whose phase still requires the complete SM anomaly assembly.
\end{theorem}

\begin{proof}
Use Lemma~\ref{lem:v5v77-three-factors},
Theorem~\ref{thm:v5v77-block-supertransfer},
Proposition~\ref{prop:v5v77-half-density},
Proposition~\ref{prop:v5v77-Z-no-phase}, and
Corollary~\ref{cor:v5v77-reference-pairing}, subject to the two firewalls.
\end{proof}

V78 will isolate the high Dirac phase by a polar decomposition of \(E\).
It will separate the positive modulus, which admits a rooted trace-log
certificate, from the unitary phase connection, and test whether the
complete Standard Model representation sum cancels its local first
variation before any branch of the logarithm is chosen.


\clearpage
\section{V78: a rooted local branch for the high Dirac phase}
\label{sec:v5v78}

\subsection{Polar separation and the local branch problem}

V77 proved that the whitening determinant is positive and cancels exactly
between the high and retained sectors.  The genuinely complex eliminated
factor is \((\det E)^{-1}\).  For an invertible high block, its right polar
decomposition is
\begin{equation}
 E=UH,
 \qquad H=(E^*E)^{1/2}>0,
 \qquad U=EH^{-1}\text{ unitary}.
 \label{eq:v5v78-polar}
\end{equation}
Thus
\begin{equation}
 \log|\det E|=\operatorname{Tr}\log H
 =\frac12\operatorname{Tr}\log(E^*E),
 \qquad
 \frac{\det E}{|\det E|}=\det U.
 \label{eq:v5v78-modulus-phase}
\end{equation}
The positive modulus fits V75 directly.  The phase requires a branch of
the logarithm.  A relative operator ball supplies such a branch and, more
strongly, gives the same rooted loop expansion for both parts.

Let \(E_0\) be an invertible local reference high Dirac block and define
\begin{equation}
 R:=E_0^{-1}(E-E_0),
 \qquad E=E_0(I+R).
 \label{eq:v5v78-relative-E}
\end{equation}
Assume
\begin{equation}
 \|R\|_{\mu,1}\le\kappa_E<1.
 \label{eq:v5v78-relative-ball}
\end{equation}

\subsection{The complex rooted trace logarithm}

Define at every block root \(x\)
\begin{equation}
 \lambda_x(R)
 :=\sum_{m\ge1}\frac{(-1)^{m+1}}m
 \operatorname{tr}_{\rm int}(R^m)_{xx}.
 \label{eq:v5v78-complex-density}
\end{equation}

\begin{theorem}[Local modulus and phase certificate]
\label{thm:v5v78-complex-trace-log}
Under \eqref{eq:v5v78-relative-ball},
\begin{equation}
 \log\frac{\det E}{\det E_0}
 =\operatorname{Tr}\log(I+R)
 =\sum_x\lambda_x(R)
 \label{eq:v5v78-logdet-ratio}
\end{equation}
on the principal relative branch, and
\begin{equation}
 \sup_x|\lambda_x(R)|
 \le d_*\log\frac1{1-\kappa_E}.
 \label{eq:v5v78-complex-bound}
\end{equation}
Consequently
\begin{align}
 \log\left|\frac{\det E}{\det E_0}\right|
 &=\sum_x\operatorname{Re}\lambda_x(R),
 \label{eq:v5v78-rooted-modulus}\\
 \arg\frac{\det E}{\det E_0}
 &=\sum_x\operatorname{Im}\lambda_x(R)
 \quad\pmod{2\pi}.
 \label{eq:v5v78-rooted-phase}
\end{align}
Each rooted density is volume-uniform and gauge invariant whenever
\(E\) and \(E_0\) transform by the same left and right gauge frames.
\end{theorem}

\begin{proof}
The operator norm of \(R\) is below one, so the principal power series for
\(\log(I+R)\) converges.  The determinant identity follows from
\eqref{eq:v5v78-relative-E}.  The closed-path proof of
Theorem~\ref{thm:v5v75-rooted-trace-log} does not use self-adjointness and
applies verbatim to \(R\), proving \eqref{eq:v5v78-complex-bound} and the
rooted sum.  Real and imaginary parts yield
\eqref{eq:v5v78-rooted-modulus}--\eqref{eq:v5v78-rooted-phase}.  Under a
common pair of gauge frames,
\(R\) is conjugated on its high domain; every diagonal closed-loop trace is
therefore invariant.
\end{proof}

\begin{theorem}[Complex local stability]
\label{thm:v5v78-complex-stability}
If \(R_1,R_2\) satisfy
\begin{equation}
 \|R_1\|_{\mu,1},\ \|R_2\|_{\mu,1}\le\kappa_E<1,
 \qquad
 \|R_1-R_2\|_{\mu,1}\le\delta_E,
 \label{eq:v5v78-two-relative-E}
\end{equation}
then
\begin{equation}
 \sup_x|\lambda_x(R_1)-\lambda_x(R_2)|
 \le\frac{d_*\delta_E}{1-\kappa_E}.
 \label{eq:v5v78-complex-stability}
\end{equation}
The same bound holds separately for the rooted modulus and phase
increments.  A remote relative change gains the two exponential legs of
Proposition~\ref{prop:v5v75-remote-trace-log}.
\end{theorem}

\begin{proof}
Use the power telescope \eqref{eq:v5v75-power-telescope}.  Its proof and
the remote closed-loop proof require only submultiplicativity, not
positivity.  Taking real or imaginary parts cannot increase the absolute
bound.
\end{proof}

\subsection{The infinitesimal phase connection}

For a differentiable family of invertible high blocks, define
\begin{equation}
 \mathcal A_E
 :=\operatorname{Im}\operatorname{Tr}(E^{-1}dE).
 \label{eq:v5v78-phase-connection}
\end{equation}

\begin{proposition}[Local exactness and rooted connection density]
\label{prop:v5v78-phase-connection}
On the relative ball \eqref{eq:v5v78-relative-ball},
\begin{equation}
 \mathcal A_E-\mathcal A_{E_0}
 =d\operatorname{Im}\operatorname{Tr}\log(I+R).
 \label{eq:v5v78-relative-connection}
\end{equation}
For a tangent \(\dot E\), the absolute connection has the rooted density
\begin{equation}
 a_x(\dot E)
 :=\operatorname{Im}\operatorname{tr}_{\rm int}
 \sum_y(E^{-1})_{xy}\dot E_{yx},
 \qquad
 \mathcal A_E(\dot E)=\sum_xa_x(\dot E),
 \label{eq:v5v78-rooted-connection}
\end{equation}
and
\begin{equation}
 \sup_x|a_x(\dot E)|
 \le d_*\|E^{-1}\|_{\mu,1}\|\dot E\|_{\mu,1}.
 \label{eq:v5v78-connection-bound}
\end{equation}
\end{proposition}

\begin{proof}
Jacobi's formula gives
\(d\log\det E=\operatorname{Tr}(E^{-1}dE)\).  Apply it to
\(E=E_0(I+R)\) on the principal relative branch and take imaginary parts,
proving \eqref{eq:v5v78-relative-connection}.  Expanding the trace in
blocks gives \eqref{eq:v5v78-rooted-connection}.  The weighted row--column
estimate and the internal trace inequality give
\eqref{eq:v5v78-connection-bound}.
\end{proof}

The one-form \(\operatorname{Tr}(E^{-1}dE)\) is closed wherever \(E\) is
invertible.  Indeed, the Maurer--Cartan equation gives
\begin{equation}
 d\operatorname{Tr}(E^{-1}dE)
 =-\operatorname{Tr}(E^{-1}dE\wedge E^{-1}dE)=0,
 \label{eq:v5v78-closed-connection}
\end{equation}
because the coefficient of each two-form is the trace of a commutator.
It is locally exact but can have integral periods around noncontractible
loops in the invertible-matrix space.

\begin{firewall}[A local branch does not remove winding]
\label{fw:v5v78-local-not-global}
The ball \(\|R\|<1\) is contractible and fixes a principal relative phase.
A large field loop may leave that ball while \(E\) remains invertible, and
\(\det E\) may wind around zero.  The closed connection can then have
nonzero period \(2\pi\mathbb Z\).  Weighted local estimates on one chart do
not prove global determinant-line triviality.
\end{firewall}

\subsection{Jet matching and complete-SM assembly}

Apply the V76 construction directly to the complex high Dirac map
\(E(z)\).  Let
\begin{equation}
 E_0^{[1]}(z)=E(0)+DE(0)[z].
 \label{eq:v5v78-E-jet}
\end{equation}
If \(E_0^{[1]}(z)^{-1}\) has weighted norm at most \(K_E\) and
\(\sup\|D^2E\|\le M_E\), Taylor's theorem gives
\begin{equation}
 \left\|(E_0^{[1]})^{-1}(E-E_0^{[1]})\right\|_{\mu',1}
 \le\frac{K_EM_E}{2}\|z\|_{\mathcal Z}^2.
 \label{eq:v5v78-E-quadratic}
\end{equation}
Thus both the residual rooted modulus and the residual rooted phase begin
quadratically on a sufficiently small uniform chart.

For Standard Model representations \(f\), orientations \(\epsilon_f\),
and multiplicities \(n_f\), the phase must be assembled as
\begin{equation}
 \lambda_x^{\rm SM}
 :=\sum_fn_f\epsilon_f
 \operatorname{Im}\lambda_{f,x}(R_f)
 \label{eq:v5v78-SM-phase-sum}
\end{equation}
before applying a local norm.  The zero-field gauge character cancellation
proved in the inherited anomaly ledger removes the corresponding vertical
constant coefficient.  Equation \eqref{eq:v5v78-SM-phase-sum} names the
higher local coefficients which still require the complete anomaly
cancellation and counterterm identities.

\begin{firewall}[Finite-matrix closedness is not anomaly cancellation]
\label{fw:v5v78-closed-not-anomaly}
Equation \eqref{eq:v5v78-closed-connection} states that the ordinary
determinant phase is locally a scalar logarithm on an invertible chart.
It does not prove that the chiral determinant-line connection is gauge
basic, that its complete SM vertical variation vanishes, or that phases
from different projective charts have trivial holonomy.  Those are the
local and global anomaly gates retained from V60--V61.
\end{firewall}

\begin{theorem}[V78 conclusion]
\label{thm:v5v78-status}
Relative norm \(\|E_0^{-1}(E-E_0)\|_{\mu,1}<1\) gives a single all-order
rooted trace-log certificate for both the modulus and the phase of the
high Dirac determinant.  The rooted phase connection is volume-uniform and
locally exact, and first-jet matching makes its residual quadratic in the
background amplitude.  The construction fixes only a contractible local
branch; complete Standard Model anomaly assembly and global phase holonomy
remain open.
\end{theorem}

\begin{proof}
Use Theorem~\ref{thm:v5v78-complex-trace-log},
Theorem~\ref{thm:v5v78-complex-stability},
Proposition~\ref{prop:v5v78-phase-connection}, and
\eqref{eq:v5v78-E-quadratic}, subject to the two firewalls.
\end{proof}

V79 will combine the exact V74 marginal state, the V77 supertransfer
factorization, and the V75/V78 rooted determinant densities into one
finite-depth Wilsonian block identity.  It will state the precise normed
inputs needed for projective consistency of the joint gauge--Higgs,
Grassmann, and pseudofermion cylinders without claiming the still-open
all-order non-Gaussian shell estimate.


\clearpage
\section{V79: a joint finite-depth Wilsonian block identity}
\label{sec:v5v79}

\subsection{Exact states, block pullbacks, and why a second theorem is needed}

V74 proves reflection positivity for an exact marginal, V77 factors the
phase-preserving regulator supertransfer, and V75/V78 control the rooted
local densities of its high determinants on a relative operator ball.
These statements must now be assembled with the gauge--Higgs and physical
Grassmann variables in one normalized block step.

At cutoff level \(n\), let \(\mathcal X_n\) denote the complete finite
configuration space, including all bosonic fields and the graded physical,
conjugate-regulator, and ghost variables.  Let \(\mathcal A_n\) be its
cylinder algebra.  A reflection-compatible gauge-covariant block map
\(b_n:\mathcal X_{n+1}\to\mathcal X_n\) induces the algebra pullback
\begin{equation}
 (\mathcal B_nF)(X)=F(b_nX).
 \label{eq:v5v79-block-pullback}
\end{equation}
For a normalized fine state \(\omega_{n+1}\), its exact marginal is
\begin{equation}
 \omega_n^{\rm ex}(F)
 :=\omega_{n+1}(\mathcal B_nF).
 \label{eq:v5v79-exact-state}
\end{equation}
This identity is tautological at one fixed finite cutoff.  The continuum
problem is harder because independently regularized actions at adjacent
cutoffs must be shown to approach these exact marginals in a common local
topology.

\subsection{Exact one-step density identity}

Write the fine variables after the covariant wavelet change as
\(X=(x,y)\), where \(x\) is retained and \(y\) is the complete shell.  The
bosonic fibre Jacobian of the unitary pseudofermion wavelet map is one.  All
other changes of variables carry their Berezinian explicitly.  If
\(d\Lambda_{n+1}(x,y)\) is a fixed product reference measure, define
\begin{equation}
 e^{-S_n^{\rm ex}(x)}
 :=\int e^{-S_{n+1}(x,y)}\,d\Lambda_{n+1}(y\mid x).
 \label{eq:v5v79-exact-effective-action}
\end{equation}
The equality is an identity of finite-dimensional densities before
normalization.

For each Pauli--Villars species \(f\), use the symmetric allocation of V77:
\begin{align}
 \mathcal H_{n,f}(x)
 &:=\frac{1}{\det E_{n,f}(x)\det Z_{n,f}(x)^{1/2}},
 \label{eq:v5v79-high-factor}\\
 \mathcal L_{n,f}(x)
 &:=\frac{\det D_{{\rm cov},n,f}(x)^*}
 {\det(D_{{\rm cov},n,f}(x)^*D_{{\rm cov},n,f}(x))}.
 \label{eq:v5v79-low-factor}
\end{align}

\begin{theorem}[Joint one-step block identity]
\label{thm:v5v79-joint-block}
Assume all finite-dimensional shell integrals exist and every high block
\(E_{n,f}\) is invertible.  After simultaneous high-shell integration,
the exact effective density can be written
\begin{equation}
 e^{-S_n^{\rm ex}(x)}
 =\mathcal W_n^{\rm bos}(x)
  \mathcal W_n^{\rm phys}(x)
  \prod_f\mathcal H_{n,f}(x)
  \;e^{-S_n^{\rm ret}(x)},
 \label{eq:v5v79-joint-density}
\end{equation}
where:
\begin{enumerate}
 \item \(\mathcal W_n^{\rm bos}\) is the exact conditional gauge--Higgs,
 gravitational, gauge-fixing, and ordinary bosonic shell integral;
 \item \(\mathcal W_n^{\rm phys}\) is the exact physical-Grassmann and
 ghost shell Berezin factor, with its determinant-line phase retained;
 \item the retained regulator action in \(S_n^{\rm ret}\) uses
 \(D_{{\rm cov},n,f}^*D_{{\rm cov},n,f}\) and the matched conjugate
 Grassmann operator \(D_{{\rm cov},n,f}^*\);
 \item for every regulator species,
 \begin{equation}
  \mathcal H_{n,f}\mathcal L_{n,f}
  =\frac1{\det D_{n+1,f}}.
  \label{eq:v5v79-species-identity}
 \end{equation}
\end{enumerate}
No \(\det Z_{n,f}\) remains in the product of the high and retained
regulator factors.
\end{theorem}

\begin{proof}
Disintegrate the bosonic reference measure and perform the shell integrals
in \eqref{eq:v5v79-exact-effective-action}.  The physical Grassmann and
ghost integrations define \(\mathcal W_n^{\rm phys}\) without taking an
absolute value.  Proposition~\ref{prop:v5v74-Gaussian-marginal} identifies
the retained pseudofermion precision, while
Proposition~\ref{prop:v5v77-half-density} supplies the conjugate Grassmann
half-density change and its Berezinian.  Theorem
\ref{thm:v5v77-block-supertransfer} proves
\eqref{eq:v5v79-species-identity} and the cancellation of \(Z_{n,f}\).
Multiplication of the remaining exact factors yields
\eqref{eq:v5v79-joint-density}.
\end{proof}

\begin{assessment}[Order of assembly]
\label{ass:v5v79-order}
Equation \eqref{eq:v5v79-joint-density} is to be assembled in the displayed
order as a complete representation and field-sector sum.  Taking the
modulus of \(\mathcal W_n^{\rm phys}\) or \(\mathcal H_{n,f}\) separately
would discard the chiral phase cancellation.  Estimating the two powers of
\(\det Z_{n,f}^{1/2}\) separately would pay a determinant which cancels
algebraically.
\end{assessment}

\subsection{Normalized conditional operators}

Let \(\mathsf T_n\) denote the exact normalized conditional expectation
from fine observables to retained observables.  Thus
\begin{equation}
 (\mathsf T_nG)(x)
 =\int G(x,y)\,d\kappa_n(y\mid x),
 \label{eq:v5v79-conditional-operator}
\end{equation}
where \(\kappa_n\) is the normalized conditional law derived from
\eqref{eq:v5v79-joint-density}.  On bounded bosonic cylinders,
\begin{equation}
 \|\mathsf T_nG\|_\infty\le\|G\|_\infty.
 \label{eq:v5v79-Markov-contraction}
\end{equation}
In the graded sector, use the V58 weighted Grassmann coefficient norm and
the V69 joint tensor norm; their one-step operator constants are recorded
as \(c_n\ge1\).

Let \(\widetilde{\mathsf T}_n\) be the transfer constructed from the
declared cutoff-\(n\) local action after relevant-parameter matching.  For
a cylinder class \(\mathfrak C_n(A)\) supported in \(A\), define the local
one-step defect
\begin{equation}
 \delta_n(A)
 :=\sup_{0\ne F\in\mathfrak C_n(A)}
 \frac{\| (\mathsf T_n-\widetilde{\mathsf T}_n)F\|_n}
 {\|F\|_{n+1}}.
 \label{eq:v5v79-local-defect}
\end{equation}

\begin{lemma}[Product telescope for nonidentical transfers]
\label{lem:v5v79-transfer-telescope}
For bounded operators \(T_j,\widetilde T_j\),
\begin{equation}
 T_m\cdots T_{N-1}-\widetilde T_m\cdots\widetilde T_{N-1}
 =\sum_{j=m}^{N-1}
 T_m\cdots T_{j-1}(T_j-\widetilde T_j)
 \widetilde T_{j+1}\cdots\widetilde T_{N-1}.
 \label{eq:v5v79-product-telescope}
\end{equation}
Empty products are identities.
\end{lemma}

\begin{proof}
Insert and subtract products in which the factors are replaced one at a
time, starting at the left.  Consecutive intermediate products cancel.
\end{proof}

\begin{theorem}[Summable one-step defects imply local projective
consistency]
\label{thm:v5v79-summable-consistency}
Fix a coarse cylinder \(F\) and let \(A_j\) be the support obtained after
pullback to level \(j\).  Suppose the exact and model transfers have
operator bounds whose accumulated distortion satisfies
\begin{equation}
 \sup_{N>j}
 \left(\prod_{i=m}^{j-1}\|\mathsf T_i\|\right)
 \left(\prod_{i=j+1}^{N-1}\|\widetilde{\mathsf T}_i\|\right)
 \le C_F
 \label{eq:v5v79-distortion-bound}
\end{equation}
and
\begin{equation}
 \sum_{j=m}^{\infty}\delta_j(A_j)<\infty.
 \label{eq:v5v79-summable-defects}
\end{equation}
Then the exact and locally modeled finite-depth expectations differ by at
most
\begin{equation}
 C_F\|F\|\sum_{j=m}^{N-1}\delta_j(A_j),
 \label{eq:v5v79-telescope-bound}
\end{equation}
and the modeled cylinder expectations form a Cauchy family whenever the
exact ancestry is projectively consistent.
\end{theorem}

\begin{proof}
Apply Lemma~\ref{lem:v5v79-transfer-telescope} to the pullback of \(F\).
Bound each summand using \eqref{eq:v5v79-local-defect} and
\eqref{eq:v5v79-distortion-bound}; summation gives
\eqref{eq:v5v79-telescope-bound}.  The tail of
\eqref{eq:v5v79-summable-defects} tends to zero.  Exact consistency removes
the difference between two exact representatives of the same cylinder, so
the triangle inequality makes the modeled expectations Cauchy.
\end{proof}

\subsection{The joint Wilsonian block certificate}

\begin{definition}[Joint Wilsonian block identity certificate]
\label{def:v5v79-JWBIC}
A JWBIC consists of the following volume-uniform data on compatible field
charts:
\begin{enumerate}
 \item reflection-admissible gauge-covariant block pullbacks and exact
 normalized disintegrations;
 \item the joint density identity \eqref{eq:v5v79-joint-density}, with all
 representation sums taken before local norms;
 \item V68 cross-plane margins and V72/V74 exact pseudofermion marginal
 locality and reflection descent;
 \item V75/V76 relative positive high-block bounds and V78 relative complex
 high-Dirac phase bounds, including their rooted local densities;
 \item the V58/V69 rooted physical and conjugate Grassmann coefficient
 estimates in the same support and provenance grades;
 \item a measure-first gauge--Higgs--gravity conditional estimate of V49
 type, including the induced high determinant rather than treating it as
 an external scalar;
 \item explicit projection of every relevant or marginal local jet to its
 running coupling and Ward owner;
 \item the distortion bound \eqref{eq:v5v79-distortion-bound} and summable
 local defects \eqref{eq:v5v79-summable-defects} for every fixed cylinder;
 \item compatible chart overlaps, global determinant-line holonomy data,
 and cofinal tightness or compact ancestry for the retained fields.
\end{enumerate}
\end{definition}

\begin{theorem}[Conditional joint local continuum state]
\label{thm:v5v79-conditional-state}
If a cofinal cutoff family carries a JWBIC, then every fixed joint cylinder
has a cutoff-independent limiting expectation.  The limits define a state
on the algebraic cylinder limit.  If the finite-cutoff states are
reflection positive and the block maps are reflection-admissible, the
limit state is reflection positive.  Exact gauge Ward identities which are
preserved by every matched transfer pass to the limit.
\end{theorem}

\begin{proof}
Theorem~\ref{thm:v5v79-summable-consistency} gives the Cauchy property and
independence of the cofinal representative.  Linearity and normalization
pass to limits.  Positivity of ordinary squares and reflected squares is
closed under entrywise convergence by
Proposition~\ref{prop:v5v74-RP-closed}.  A Ward identity is a finite linear
relation among cylinder expectations; if it holds at every level and the
terms converge, it holds in the limit.
\end{proof}

\begin{firewall}[The certificate has not been acquired]
\label{fw:v5v79-JWBIC-open}
V79 proves that the listed data are sufficient and shows their exact
finite-depth algebraic compatibility.  The physical-Grassmann all-order
rooted estimate, the non-Gaussian gauge--Higgs--gravity shell estimate,
summability of their joint defects, chart holonomy, and cofinal compactness
have not all been proved.  The JWBIC is therefore not yet a construction of
the SM--Einstein continuum state.
\end{firewall}

\begin{firewall}[Cylinder convergence is not full OS reconstruction]
\label{fw:v5v79-not-full-OS}
Even a JWBIC supplies only the limiting state and reflection positivity on
the declared cylinder algebra.  Continuum Euclidean covariance,
regularity, clustering, nontriviality, reconstruction of the Hamiltonian,
and the Einstein constraint algebra require the separate inherited gates.
\end{firewall}

\begin{theorem}[V79 conclusion]
\label{thm:v5v79-status}
The exact gauge--Higgs, physical Grassmann, conjugate-regulator, and
pseudofermion shell factors now form one finite-depth Wilsonian density
identity.  A product-telescope theorem reduces projective consistency of
independently regularized cutoff actions to summable observable-local
one-step defects with controlled norm distortion.  The JWBIC records the
precise compatible inputs which would yield a reflection-positive joint
limiting cylinder state.  Several non-Gaussian and global inputs remain
open, so no continuum construction is claimed.
\end{theorem}

\begin{proof}
Use Theorem~\ref{thm:v5v79-joint-block},
Lemma~\ref{lem:v5v79-transfer-telescope},
Theorem~\ref{thm:v5v79-summable-consistency}, and
Theorem~\ref{thm:v5v79-conditional-state}, with the two firewalls.
\end{proof}

V80 is the next compulsory YM/NS audit.  It will compare the newest
all-arity YM acquisition and NS stopping-time acquisition with the two
unproved analytic entries of the JWBIC, then attack the first joint defect
which can be sharpened without assuming either companion programme's final
theorem.


\clearpage
\section{V80: companion audit and negative-to-strong interpolation of a joint shell defect}
\label{sec:v5v80}

\subsection{Compulsory companion audit}

The newest active Yang--Mills note at the time of inspection was
\begin{equation}
 \begin{split}
 &\texttt{Renewal\_Yang\_Mills\_Mass\_Gap\_Research\_Notes\_V90.tex},\\
 &1\,335\,805\text{ bytes},\qquad
 \operatorname{SHA256}=\\
 &\texttt{1EFAEC15BD87B8E293C48821E50B3095E9274F8EDE8919DE786AE04475CB8397}.
 \end{split}
 \label{eq:v5v80-YM-fingerprint}
\end{equation}
It proves that the projected Wilson product subtracts the exact
finite-parameter covariance, proves the two-label matrix certificate in
the high-thermal sector, and reduces the low-thermal sector to an assembled
negative-norm bound plus a separate high-spectral-tail bound.  It explicitly
forbids reversing the spectral-gap inequality; the low-thermal certificate,
higher subsets, and the all-order card remain open.

The active Navier--Stokes note was unchanged:
\begin{equation}
 \begin{split}
 &\texttt{Renewal\_NS\_Stability\_Research\_Notes\_V31.tex},\\
 &887\,537\text{ bytes},\qquad
 \operatorname{SHA256}=\\
 &\texttt{F1121B62238AD5491BB7CF03635EA9934942EDF573ED8FAAA9EE79E1D38A6696}.
 \end{split}
 \label{eq:v5v80-NS-fingerprint}
\end{equation}
Its relevant warning remains that a global estimate cannot be unmarked for
free: the exact physical root must be retained until the missing first
moment is paid.  Both files were read passively and were not modified.

For the V79 JWBIC, the corresponding object is a normalized, root-local
shell density defect.  Unlike the YM Pruefer target, projective consistency
does not require a particular sharp incidence power: any positive dyadic
cutoff gain is summable.  This permits an interpolation route which keeps
the negative-norm cancellation and pays the high tail through one fixed
positive shell derivative.

\subsection{A shell Hilbert scale}

Let \((\Omega,\nu)\) be one normalized conditional shell probability space
at fixed retained background.  Let \(\mathscr L\ge0\) be its self-adjoint
Markov generator in \(L^2(\nu)\), with constants as its kernel.  On the
mean-zero subspace define, for \(s>0\),
\begin{equation}
 \|h\|_{-1}:=\|\mathscr L^{-1/2}h\|_2,
 \qquad
 \|h\|_{s}:=\|\mathscr L^{s/2}h\|_2.
 \label{eq:v5v80-Hilbert-scale}
\end{equation}
The generator may be a conditional Langevin/Witten generator rather than a
free Ornstein--Uhlenbeck operator; only self-adjoint functional calculus is
used below.

\begin{theorem}[Exact negative-to-strong interpolation]
\label{thm:v5v80-interpolation}
For every mean-zero \(h\) in the domains appearing in
\eqref{eq:v5v80-Hilbert-scale},
\begin{equation}
 \boxed{\qquad
 \|h\|_2
 \le
 \|h\|_{-1}^{\,s/(s+1)}
 \|h\|_{s}^{\,1/(s+1)}.
 \qquad}
 \label{eq:v5v80-interpolation}
\end{equation}
The same inequality holds in the column norm for finite-dimensional
operator-valued \(h\), with constants independent of the target dimension.
\end{theorem}

\begin{proof}
Let \(d\mu_h(\lambda)\) be the spectral measure of \(h\) for
\(\mathscr L\) on the positive spectrum.  H\"older's inequality with
exponents \((s+1)/s\) and \(s+1\) gives
\begin{align}
 \|h\|_2^2
 &=\int
 (\lambda^{-1})^{s/(s+1)}
 (\lambda^s)^{1/(s+1)}\,d\mu_h(\lambda)\\
 &\le
 \left(\int\lambda^{-1}\,d\mu_h\right)^{s/(s+1)}
 \left(\int\lambda^s\,d\mu_h\right)^{1/(s+1)}.
\end{align}
Taking square roots proves \eqref{eq:v5v80-interpolation}.  For an
operator-valued function, apply the scalar Hilbert-space proof to \(hv\)
for each unit input vector \(v\), bound both factors by their column suprema,
and then take the supremum over \(v\).
\end{proof}

\begin{corollary}[Spectral split without reversing the gap]
\label{cor:v5v80-spectral-split}
For every \(\Lambda>0\),
\begin{equation}
 \|h\|_2
 \le
 \Lambda^{1/2}\|h\|_{-1}
 +\Lambda^{-s/2}\|h\|_s.
 \label{eq:v5v80-spectral-split}
\end{equation}
Optimizing \(\Lambda\) recovers
\eqref{eq:v5v80-interpolation} up to a constant depending only on \(s\).
\end{corollary}

\begin{proof}
On the spectral interval \((0,\Lambda]\), one has
\(1\le\Lambda\lambda^{-1}\); on \((\Lambda,\infty)\), one has
\(1\le\Lambda^{-s}\lambda^s\).  Apply these inequalities to the two
orthogonal spectral pieces and use the triangle inequality.  Balancing the
two displayed terms gives the interpolation exponents.
\end{proof}

The theorem strengthens the raw split for the present purpose: one need not
name a sharp spectral cutoff once a single positive shell derivative is
available.  It does not obtain that derivative from a spectral gap.

\subsection{Normalized conditional density defects}

Let \(\widetilde\kappa_n(dy\mid x)\) be a reference conditional shell law
at retained background \(x\), and let the exact conditional law have
normalized density
\begin{equation}
 \kappa_n(dy\mid x)
 =r_n(x,y)\widetilde\kappa_n(dy\mid x),
 \qquad
 r_n=1+h_n,
 \qquad
 \mathbb E_{\widetilde\kappa_n}h_n=0.
 \label{eq:v5v80-density-defect}
\end{equation}
The centering is exact because both conditional laws have mass one.  It is
the measure-level analogue of subtracting the exact Wilson covariance in
YM V90.

\begin{lemma}[Strong density norm controls the transfer defect]
\label{lem:v5v80-density-transfer}
For every square-integrable shell observable \(F(x,\cdot)\),
\begin{equation}
 \left|
 \mathbb E_{\kappa_n}F-
 \mathbb E_{\widetilde\kappa_n}F
 \right|
 \le
 \|F\|_{L^2(\widetilde\kappa_n)}
 \|h_n\|_{L^2(\widetilde\kappa_n)}.
 \label{eq:v5v80-transfer-L2}
\end{equation}
In particular, the local transfer defect of V79 is bounded by
\(\|h_n\|_2\) on a unit \(L^2\) cylinder class.
\end{lemma}

\begin{proof}
Equation \eqref{eq:v5v80-density-defect} makes the difference equal to
\(\mathbb E_{\widetilde\kappa_n}(Fh_n)\).  Apply Cauchy--Schwarz.
\end{proof}

\begin{theorem}[Negative-norm shell certificate]
\label{thm:v5v80-negative-certificate}
Let \(a_n=q^{-n}a_0\) with \(q>1\).  Suppose, uniformly in volume,
retained background, and every root in a fixed cylinder support,
\begin{align}
 \|h_n\|_{-1}&\le A a_n^\alpha,
 \qquad \alpha>0,
 \label{eq:v5v80-negative-gain}\\
 \|h_n\|_s&\le B
 \label{eq:v5v80-positive-bound}
\end{align}
for one fixed \(s>0\).  Then
\begin{equation}
 \|h_n\|_2
 \le A^{s/(s+1)}B^{1/(s+1)}
 a_n^{\alpha s/(s+1)},
 \label{eq:v5v80-strong-gain}
\end{equation}
and
\begin{equation}
 \sum_{n\ge0}\|h_n\|_2<\infty.
 \label{eq:v5v80-summable-L2}
\end{equation}
Consequently these shell defects satisfy the summability entry of the V79
JWBIC for unit local \(L^2\) cylinders, provided the support-pullback and
norm-distortion constants remain bounded.
\end{theorem}

\begin{proof}
Apply Theorem~\ref{thm:v5v80-interpolation} to
\eqref{eq:v5v80-negative-gain}--\eqref{eq:v5v80-positive-bound}.  Since
\(a_n^{\alpha s/(s+1)}=a_0^{\alpha s/(s+1)}
q^{-n\alpha s/(s+1)}\), the series is geometric.  Lemma
\ref{lem:v5v80-density-transfer} and
Theorem~\ref{thm:v5v79-summable-consistency} give the JWBIC assertion.
\end{proof}

\begin{assessment}[Why a nonsharp gain can suffice here]
\label{ass:v5v80-any-gain}
The YM SPCC seeks a sharp representation-incidence power and therefore
cannot replace its high spectral tail by an arbitrary interpolation loss.
The SM projective-consistency telescope asks for a summable sequence.  On a
geometric cutoff scale, every positive exponent in
\eqref{eq:v5v80-strong-gain} is summable.  The weaker target permits the
negative-norm cancellation to be converted into useful progress without
claiming the YM low-thermal pair estimate.
\end{assessment}

\subsection{Rooted and joint versions}

The defect \(h_n\) in Theorem~\ref{thm:v5v80-negative-certificate} must be
the complete normalized local density ratio after the SM representation
sum and the V77 determinant cancellations.  If it is replaced by a global
density on the whole lattice, both \(A\) and \(B\) can grow with volume and
the theorem gives no local certificate.

Let \(P_Ah_n\) denote the conditional density component connected to a
fixed retained support \(A\), using the rooted polymer projection of the
V58/V69 norm.  A usable joint estimate is
\begin{equation}
 \|P_Ah_n\|_{-1}\le A_Aa_n^\alpha,
 \qquad
 \|P_Ah_n\|_s\le B_A,
 \label{eq:v5v80-rooted-joint}
\end{equation}
with \(A_A,B_A\) independent of total volume.  The NS audit explains why
the projection must precede the norm: summing over all possible roots and
then trying to divide by volume loses the one-use ancestry.

\begin{definition}[Negative-interpolation local certificate]
\label{def:v5v80-NILC}
A NILC for one joint shell consists of:
\begin{enumerate}
 \item an exact normalized density ratio \(1+h_n\) after all algebraic
 determinant and representation cancellations;
 \item a self-adjoint conditional shell generator with a common functional
 calculus on the mean-zero subspace;
 \item rooted bounds \eqref{eq:v5v80-rooted-joint} for some fixed
 \(s,\alpha>0\);
 \item compatibility of the generator domain with the physical Grassmann
 coefficient norm and the bosonic Lyapunov weight;
 \item summable support-pullback and transfer-distortion constants.
\end{enumerate}
\end{definition}

\begin{firewall}[Positive regularity is an independent payment]
\label{fw:v5v80-positive-payment}
Neither centering nor a spectral gap proves
\eqref{eq:v5v80-positive-bound}.  The bound differentiates the complete
normalized density in shell directions.  For gauge, Higgs, and gravity
fields it must control the interaction, Jacobian, determinant, and chart
terms together.  Estimating them separately may destroy the cancellation
which produced \eqref{eq:v5v80-negative-gain}.
\end{firewall}

\begin{firewall}[An \(L^2\) density chart has a finite domain]
\label{fw:v5v80-L2-domain}
Stable real-contour interactions need not have a uniformly square-integrable
density ratio relative to a poorly chosen Gaussian reference.  The
reference and Lyapunov weight must be chosen so that \(h_n\in L^2\) and in
the domain of \(\mathscr L^{s/2}\).  V49's \(L^1\) contraction alone does
not imply this stronger property.
\end{firewall}

\begin{theorem}[V80 conclusion]
\label{thm:v5v80-status}
The compulsory audit records YM V90 and unchanged NS V31.  Their
negative-norm and rooted-first-moment lessons yield an exact Hilbert-scale
interpolation theorem for a normalized joint shell density defect.  A
rooted negative norm with any positive cutoff power, combined with one
uniform positive shell derivative, gives a summable strong transfer error
and therefore supplies the local-defect entry of the V79 JWBIC.  Neither of
those two rooted estimates has yet been proved for the complete
SM--Einstein shell density.
\end{theorem}

\begin{proof}
Use Theorem~\ref{thm:v5v80-interpolation},
Lemma~\ref{lem:v5v80-density-transfer}, and
Theorem~\ref{thm:v5v80-negative-certificate}, subject to the two
firewalls.
\end{proof}

V81 will seek the first NILC line from an exact Poisson/Stein equation for
the normalized conditional density.  It will identify the source whose
divergence represents \(h_n\), prove the resulting negative-norm duality,
and determine which complete SM Ward cancellation can supply a positive
power of the cutoff before any sectorwise absolute value is taken.


\clearpage
\section{V81: a Poisson--Stein bound for the normalized shell density}
\label{sec:v5v81}

\subsection{The exact negative norm}

V80 showed that a positive cutoff gain in a rooted negative norm is enough
for summable local transfer errors, provided one fixed positive shell
derivative is uniformly bounded.  We now identify a rigorous source for
that negative norm in a positive conditional Gibbs sector.

Let \(\nu\) be a normalized reference shell measure and let
\(d\) be a closed shell gradient with self-adjoint generator
\begin{equation}
 \mathscr L=d^*d\ge0.
 \label{eq:v5v81-generator}
\end{equation}
Assume constants are the kernel of \(\mathscr L\).  For mean-zero
\(h\in\operatorname{Dom}(\mathscr L^{-1/2})\), define
\begin{equation}
 \|h\|_{-1,\nu}:=\|\mathscr L^{-1/2}h\|_{L^2(\nu)}.
 \label{eq:v5v81-negative-norm}
\end{equation}

\begin{lemma}[Poisson duality]
\label{lem:v5v81-Poisson-duality}
One has
\begin{equation}
 \|h\|_{-1,\nu}
 =\sup\left\{
 |\mathbb E_\nu(Fh)|:
 F\in\operatorname{Dom}(d),\ 
 \|dF\|_{L^2(\nu)}\le1
 \right\}.
 \label{eq:v5v81-negative-dual}
\end{equation}
If \(u=\mathscr L^{-1}h\), then
\begin{equation}
 h=d^*du,
 \qquad
 \mathbb E_\nu(Fh)=\langle dF,du\rangle_{L^2(\nu)},
 \qquad
 \|du\|_2=\|h\|_{-1,\nu}.
 \label{eq:v5v81-Poisson-current}
\end{equation}
\end{lemma}

\begin{proof}
The spectral theorem gives
\(u=\mathscr L^{-1}h\) on the mean-zero subspace and
\begin{equation}
 \|du\|_2^2
 =\langle u,\mathscr Lu\rangle
 =\langle h,\mathscr L^{-1}h\rangle
 =\|h\|_{-1,\nu}^2.
\end{equation}
The definition of the adjoint gives the middle identity in
\eqref{eq:v5v81-Poisson-current}; Cauchy--Schwarz proves the upper bound in
\eqref{eq:v5v81-negative-dual}.  Equality follows by taking
\(F=u/\|du\|_2\), with the zero case immediate.
\end{proof}

\subsection{Gibbs interpolation}

Let \(W\) be a real residual shell interaction after the relevant and
marginal reference terms have been matched.  Define
\begin{equation}
 Z_t=\mathbb E_\nu e^{-tW},
 \qquad
 r_t=Z_t^{-1}e^{-tW},
 \qquad
 d\nu_t=r_t\,d\nu,
 \qquad 0\le t\le1.
 \label{eq:v5v81-Gibbs-path}
\end{equation}
The final normalized density defect is
\begin{equation}
 h=r_1-1,
 \qquad \mathbb E_\nu h=0.
 \label{eq:v5v81-final-defect}
\end{equation}

\begin{lemma}[Exact covariance derivative]
\label{lem:v5v81-covariance-derivative}
For every integrable \(F\) for which differentiation is justified,
\begin{equation}
 \frac d{dt}\mathbb E_{\nu_t}F
 =-\operatorname{Cov}_{\nu_t}(F,W).
 \label{eq:v5v81-covariance-derivative}
\end{equation}
Consequently
\begin{equation}
 \mathbb E_\nu(Fh)
 =-\int_0^1\operatorname{Cov}_{\nu_t}(F,W)\,dt.
 \label{eq:v5v81-exact-Gibbs-difference}
\end{equation}
\end{lemma}

\begin{proof}
Differentiate \(Z_t^{-1}\mathbb E_\nu(Fe^{-tW})\).  Since
\(Z_t'=-Z_t\mathbb E_{\nu_t}W\), the two terms combine to
\(-\mathbb E_{\nu_t}(FW)+mathbb E_{\nu_t}F\mathbb E_{\nu_t}W\).
Integrate from zero to one.
\end{proof}

\begin{theorem}[Poisson--Stein Gibbs bound]
\label{thm:v5v81-Gibbs-negative}
Assume every \(\nu_t\) satisfies the common Poincar\'e inequality
\begin{equation}
 \operatorname{Var}_{\nu_t}(G)
 \le C_P\mathbb E_{\nu_t}|dG|^2
 \label{eq:v5v81-Poincare}
\end{equation}
and that the interpolating densities obey
\begin{equation}
 0\le r_t\le M
 \quad\nu\text{-almost everywhere},
 \qquad 0\le t\le1.
 \label{eq:v5v81-density-domination}
\end{equation}
Then
\begin{equation}
 \boxed{\qquad
 \|r_1-1\|_{-1,\nu}
 \le
 C_P\sqrt M
 \int_0^1
 \left(\mathbb E_{\nu_t}|dW|^2\right)^{1/2}dt.
 \qquad}
 \label{eq:v5v81-Gibbs-negative}
\end{equation}
\end{theorem}

\begin{proof}
Polarization of \eqref{eq:v5v81-Poincare} and Cauchy--Schwarz give
\begin{equation}
 |\operatorname{Cov}_{\nu_t}(F,W)|
 \le C_P
 \left(\mathbb E_{\nu_t}|dF|^2\right)^{1/2}
 \left(\mathbb E_{\nu_t}|dW|^2\right)^{1/2}.
 \label{eq:v5v81-covariance-Poincare}
\end{equation}
By \eqref{eq:v5v81-density-domination}, the first factor is at most
\(\sqrt M\|dF\|_{L^2(\nu)}\).  Insert this in
\eqref{eq:v5v81-exact-Gibbs-difference}, integrate in \(t\), and take the
supremum from Lemma~\ref{lem:v5v81-Poisson-duality}.
\end{proof}

\begin{corollary}[A cutoff gain from the assembled residual gradient]
\label{cor:v5v81-cutoff-gain}
If, uniformly in cutoff, volume, and retained background,
\begin{equation}
 C_P\le C_*,
 \qquad M\le M_*,
 \qquad
 \sup_{0\le t\le1}
 \left(\mathbb E_{\nu_t}|dW_n|^2\right)^{1/2}
 \le C_Wa_n^\alpha,
 \quad\alpha>0,
 \label{eq:v5v81-gradient-gain}
\end{equation}
then
\begin{equation}
 \|h_n\|_{-1,\nu}
 \le C_*\sqrt{M_*}C_Wa_n^\alpha.
 \label{eq:v5v81-negative-gain}
\end{equation}
Together with a uniform \(s\)-derivative bound from V80, this gives the
summable strong defect
\begin{equation}
 \|h_n\|_2
 \le C a_n^{\alpha s/(s+1)}.
 \label{eq:v5v81-strong-gain}
\end{equation}
\end{corollary}

\begin{proof}
Apply Theorem~\ref{thm:v5v81-Gibbs-negative}, then
Theorem~\ref{thm:v5v80-negative-certificate}.
\end{proof}

\subsection{Root-local conditional form}

For a fixed retained cylinder support \(A\), let \(W_{n,A}\) be the
connected shell interaction whose polymers meet the pulled-back support
\(A_n\).  The remaining vacuum polymers cancel between numerator and
normalization.  Let \(d_{n,A}\) differentiate only shell variables in the
rooted dependency neighbourhood and let \(\mathscr L_{n,A}=d_{n,A}^*
d_{n,A}\).

\begin{proposition}[Rooted transfer bound]
\label{prop:v5v81-rooted-transfer}
Suppose the conditional Gibbs interpolation for \(W_{n,A}\) satisfies
\eqref{eq:v5v81-Poincare}--\eqref{eq:v5v81-density-domination} with
constants \(C_A,M_A\) independent of total volume.  Then the local transfer
defect obeys
\begin{equation}
 \delta_n(A)
 \le
 C_A\sqrt{M_A}
 \int_0^1
 \|d_{n,A}W_{n,A}\|_{L^2(\nu_{n,A,t})}\,dt
 \label{eq:v5v81-rooted-transfer}
\end{equation}
after the V80 interpolation payment when a positive derivative is needed.
Vacuum volume does not appear.
\end{proposition}

\begin{proof}
Apply Theorem~\ref{thm:v5v81-Gibbs-negative} to the rooted conditional
measure.  Lemma~\ref{lem:v5v80-density-transfer} converts its strong norm
to the transfer defect.  Polymers disjoint from the observable cancel in
the normalized conditional expectation, so the constants depend on the
rooted neighbourhood rather than the total number of lattice cells.
\end{proof}

This proposition gives the exact location at which the NS first-moment
lesson enters: the connected projection meeting \(A_n\) is made before
the shell gradient and probability norm are taken.

\subsection{Complete-sector assembly of the gradient}

The residual interaction in \eqref{eq:v5v81-gradient-gain} is not a sum of
independently normed sector remainders.  It is the logarithm of the complete
positive conditional density after:
\begin{enumerate}
 \item the V77 cancellation of whitening determinants;
 \item the complete SM representation sum in the rooted determinant
 densities of V75 and V78;
 \item projection of relevant and marginal local jets to their running
 couplings;
 \item gauge Ward and BRST partners have been placed in the same rooted
 polymer;
 \item the physical Grassmann coefficients have been retained in their
 graded norm rather than replaced by an absolute determinant.
\end{enumerate}
Only after this assembly may one test
\(\|dW_{n,A}\|_2=O(a_n^\alpha)\).  A cancellation of the complete gradient
can yield a positive power even when each displayed sector has an
order-one derivative.

\begin{firewall}[The chiral phase is not a positive Gibbs tilt]
\label{fw:v5v81-phase-not-Gibbs}
Theorem~\ref{thm:v5v81-Gibbs-negative} applies to a real normalized density.
The chiral determinant phase and Grassmann coefficients must remain in the
V58/V69 graded complex norm and use the V78 local phase connection.  Folding
them into \(W\) as a complex Gibbs potential would destroy the probability
and Poincar\'e argument.
\end{firewall}

\begin{firewall}[Uniform Poincar\'e control is a genuine dynamical gate]
\label{fw:v5v81-Poincare-open}
Strict convexity can give \eqref{eq:v5v81-Poincare} for some massive local
shells, but compact gauge orbits, Higgs near-zero strata, gravitational
conformal directions, and phase transitions can remove a uniform gap.
Gauge quotienting and stratified conditional charts must be performed
before a common \(C_P\) is asserted.
\end{firewall}

\begin{firewall}[Pointwise density domination is strong]
\label{fw:v5v81-domination-open}
The bound \(r_t\le M\) is sufficient and transparent but may fail for
unbounded polynomial interactions.  It may be replaced by a weighted
H\"older comparison between \(L^2(\nu_t)\) and the reference Dirichlet
space, but such a comparison needs a proved Lyapunov or log-Sobolev bound;
normalization alone supplies none.
\end{firewall}

\begin{theorem}[V81 conclusion]
\label{thm:v5v81-status}
For a positive normalized conditional Gibbs sector, the exact covariance
interpolation and a uniform Poincar\'e inequality bound the reference
negative norm of the density defect by the integrated shell gradient of
the complete residual interaction.  A rooted gradient gain
\(O(a_n^\alpha)\) therefore supplies the first NILC line and, with V80's
independent positive derivative, a summable strong transfer defect.  The
uniform Poincar\'e, density-comparison, assembled-gradient, and chiral
graded estimates remain open for the full SM--Einstein shell.
\end{theorem}

\begin{proof}
Use Lemma~\ref{lem:v5v81-Poisson-duality},
Lemma~\ref{lem:v5v81-covariance-derivative},
Theorem~\ref{thm:v5v81-Gibbs-negative},
Corollary~\ref{cor:v5v81-cutoff-gain}, and
Proposition~\ref{prop:v5v81-rooted-transfer}, subject to the three
firewalls.
\end{proof}

V82 will address the second NILC line.  It will differentiate the normalized
Gibbs density exactly, derive a weighted Sobolev bound in terms of finitely
many derivatives of the assembled residual interaction, and determine
whether a stable Higgs quartic and compact gauge shell supply the required
uniform positive regularity without a pointwise density-ratio bound.


\clearpage
\section{V82: the exact Fisher payment for positive shell regularity}
\label{sec:v5v82}

\subsection{One positive derivative is sufficient}

V80 permits any fixed \(s>0\) in the positive shell norm.  Thus there is no
need to differentiate the Gibbs density to all orders.  In the Dirichlet
setting of V81, take \(s=1\).  Let
\begin{equation}
 r=Z^{-1}e^{-W},
 \qquad Z=\mathbb E_\nu e^{-W},
 \qquad h=r-1.
 \label{eq:v5v82-density}
\end{equation}
Here \(Z\) may depend on retained variables but is constant in the shell
directions differentiated by \(d\).

\begin{lemma}[Exact normalized-density derivative]
\label{lem:v5v82-density-derivative}
If \(W\in\operatorname{Dom}(d)\) and the chain rule is valid, then
\begin{equation}
 dh=dr=-r\,dW.
 \label{eq:v5v82-density-derivative}
\end{equation}
Consequently
\begin{equation}
 \|h\|_1^2
 :=\|\mathscr L^{1/2}h\|_2^2
 =\|dh\|_2^2
 =\mathbb E_\nu\bigl[r^2|dW|^2\bigr].
 \label{eq:v5v82-Fisher-identity}
\end{equation}
\end{lemma}

\begin{proof}
The shell derivative annihilates \(Z\), so the chain rule gives
\(dr=-Z^{-1}e^{-W}dW\).  Since \(d1=0\), this proves
\eqref{eq:v5v82-density-derivative}.  The quadratic-form identity
\(\|\mathscr L^{1/2}h\|_2=\|dh\|_2\) follows from
\(\mathscr L=d^*d\), giving \eqref{eq:v5v82-Fisher-identity}.
\end{proof}

Define the rooted order-two Fisher payment
\begin{equation}
 \mathcal I_{2,\nu}(W)
 :=\mathbb E_\nu\bigl[r^2|dW|^2\bigr].
 \label{eq:v5v82-I2}
\end{equation}
Unlike pointwise density domination, it permits an unbounded density ratio
provided its gradient-weighted second moment is finite.

\begin{theorem}[Fisher completion of the NILC]
\label{thm:v5v82-Fisher-NILC}
Suppose a rooted normalized density defect satisfies
\begin{equation}
 \|h_n\|_{-1}\le Aa_n^\alpha,
 \qquad
 \mathcal I_{2,\nu_n}(W_n)\le B^2,
 \qquad \alpha>0,
 \label{eq:v5v82-two-lines}
\end{equation}
uniformly in cutoff and volume.  Then
\begin{equation}
 \|h_n\|_2
 \le\sqrt{AB}\,a_n^{\alpha/2},
 \qquad
 \sum_n\|h_n\|_2<\infty.
 \label{eq:v5v82-half-gain}
\end{equation}
Hence the local transfer defects are summable on unit \(L^2\) cylinder
classes, subject to the V79 support-distortion bound.
\end{theorem}

\begin{proof}
Lemma~\ref{lem:v5v82-density-derivative} gives \(\|h_n\|_1\le B\).
Apply Theorem~\ref{thm:v5v80-interpolation} with \(s=1\):
\begin{equation}
 \|h_n\|_2\le\|h_n\|_{-1}^{1/2}\|h_n\|_1^{1/2}.
\end{equation}
This proves the first inequality.  The dyadic series in
\(a_n^{\alpha/2}\) is geometric.  Lemma
\ref{lem:v5v80-density-transfer} completes the transfer assertion.
\end{proof}

\subsection{A Lyapunov criterion without pointwise density domination}

Let \(V_L\ge0\) be a shell Lyapunov function.  The next criterion is
designed for a stable Higgs polynomial relative to a massive Gaussian or
already stable reference law.

\begin{proposition}[Stable-potential Fisher bound]
\label{prop:v5v82-stable-Fisher}
Assume, uniformly in the retained background,
\begin{align}
 W&\ge cV_L-C_0,
 \qquad c>0,
 \label{eq:v5v82-coercive-W}\\
 |dW|^2&\le C_1(1+V_L)^p,
 \label{eq:v5v82-gradient-growth}\\
 \mathbb E_\nu W&\le C_2.
 \label{eq:v5v82-mean-upper}
\end{align}
Then
\begin{equation}
 \mathcal I_{2,\nu}(W)
 \le
 e^{2C_2+2C_0}C_1
 \sup_{v\ge0}(1+v)^pe^{-2cv}<\infty.
 \label{eq:v5v82-stable-Fisher-bound}
\end{equation}
No upper bound on \(r\) is assumed.
\end{proposition}

\begin{proof}
Jensen's inequality gives
\begin{equation}
 Z=\mathbb E_\nu e^{-W}\ge e^{-\mathbb E_\nu W}\ge e^{-C_2}.
\end{equation}
Therefore
\begin{align}
 \mathcal I_{2,\nu}(W)
 &=Z^{-2}\mathbb E_\nu(e^{-2W}|dW|^2)\\
 &\le e^{2C_2+2C_0}C_1
 \mathbb E_\nu\bigl[(1+V_L)^pe^{-2cV_L}\bigr],
\end{align}
which is bounded by the displayed supremum.
\end{proof}

The criterion is insensitive to the dimension only if
\(C_0,C_1,C_2,c,p\) are root-local and uniform.  Applying it to the total
lattice potential would generally make \(C_0\) and \(C_2\) extensive.

\begin{corollary}[Stable Higgs shell]
\label{cor:v5v82-Higgs}
On a rooted finite shell, suppose the matched real Higgs residual satisfies
\begin{equation}
 W_H(\varphi)
 \ge c_4\|\varphi\|_4^4+c_2\|\varphi\|_2^2-C_0,
 \qquad c_4>0,
 \label{eq:v5v82-Higgs-stability}
\end{equation}
and its shell gradient has polynomial growth of fixed degree, with uniform
local coefficients.  If \(\mathbb E_\nu W_H\) is uniformly bounded above,
then \(\mathcal I_{2,\nu}(W_H)\) is uniformly bounded.
\end{corollary}

\begin{proof}
Take
\(V_L=\|\varphi\|_4^4+\|\varphi\|_2^2\), after absorbing a bounded negative
quadratic part into \(C_0\).  Polynomial growth of \(dW_H\) is bounded by a
power of \(1+V_L\).  Apply
Proposition~\ref{prop:v5v82-stable-Fisher}.
\end{proof}

\subsection{Compact gauge shells}

Let the rooted gauge shell range over a fixed finite product of compact
groups with Haar reference probability \(\nu_G\).

\begin{proposition}[Compact-shell Fisher bound]
\label{prop:v5v82-compact-gauge}
If a real gauge-shell residual \(W_G\) is continuously differentiable and
\begin{equation}
 \operatorname{osc}W_G
 :=\sup W_G-\inf W_G\le O_G,
 \qquad
 \sup|dW_G|\le G_1,
 \label{eq:v5v82-compact-bounds}
\end{equation}
then its normalized density satisfies
\begin{equation}
 \mathcal I_{2,\nu_G}(W_G)
 \le e^{2O_G}G_1^2.
 \label{eq:v5v82-gauge-Fisher}
\end{equation}
\end{proposition}

\begin{proof}
Since
\(e^{-\sup W_G}\le Z\le e^{-\inf W_G}\), one has
\(r\le e^{\operatorname{osc}W_G}\).  Insert this and the gradient bound in
\eqref{eq:v5v82-I2}.
\end{proof}

The dimension-free content again comes from rooting: the number of compact
group factors meeting one fixed local polymer is bounded.  A product over
the whole lattice would make \(O_G\) extensive.

\subsection{Products and complete residuals}

If \(W=\sum_jW_j\), the exact payment is
\begin{equation}
 \mathcal I_{2,\nu}(W)
 =\mathbb E_\nu\left[r^2\left|\sum_jdW_j\right|^2\right].
 \label{eq:v5v82-assembled-Fisher}
\end{equation}
The sum is inside the square.  Ward partners, determinant phase-free
moduli, gauge-fixing terms, and Higgs counterterms must therefore be
assembled before estimating it.  The sectorwise inequality
\(|\sum_jdW_j|^2\le N\sum_j|dW_j|^2\) is valid but can lose both the Ward
cancellation and a growing sector multiplicity.

\begin{proposition}[Bounded finite-sector assembly]
\label{prop:v5v82-finite-assembly}
If the number of rooted real sectors is at most \(N_*<\infty\), each
\(|dW_j|^2\le C_j(1+V_L)^{p_j}\), and the complete potential satisfies the
coercive lower bound \eqref{eq:v5v82-coercive-W}, then
\(\mathcal I_{2,\nu}(W)\) is uniformly bounded, with a constant depending
on \(N_*\), the local coefficient bounds, and the Lyapunov data.
\end{proposition}

\begin{proof}
Use
\(|\sum_jdW_j|^2\le N_*\sum_j|dW_j|^2\), dominate the finite sum by one
power of \(1+V_L\), and apply the proof of
Proposition~\ref{prop:v5v82-stable-Fisher} to the complete \(W\).
\end{proof}

This proposition gives a robust fallback.  A sharper proof should retain
the complete sum in \eqref{eq:v5v82-assembled-Fisher} when cancellation is
needed for uniformity or a cutoff gain.

\subsection{The gravitational boundary of the positive chart}

\begin{firewall}[The Euclidean conformal mode is not covered]
\label{fw:v5v82-conformal-mode}
The Euclidean Einstein--Hilbert action is not bounded below along the
conformal direction before a contour prescription, constraint reduction,
or stabilizing ultraviolet completion is supplied.  Therefore it does not
satisfy \eqref{eq:v5v82-coercive-W}, and
Proposition~\ref{prop:v5v82-stable-Fisher} cannot be applied to the full
metric shell.  Treating that direction as an ordinary positive Gibbs
variable would assume away the conformal-factor problem.
\end{firewall}

\begin{firewall}[Gauge fixing must precede a uniform Poincar\'e claim]
\label{fw:v5v82-gauge-Poincare}
Proposition~\ref{prop:v5v82-compact-gauge} proves the positive derivative
payment on a rooted compact shell.  V81 also needs a uniform conditional
Poincar\'e inequality.  Gauge orbit zero modes and transitions between
strata must first be quotiented or pinned; compactness by itself does not
give the required gap on the physical conditional algebra.
\end{firewall}

\begin{firewall}[The complex sectors retain their own norm]
\label{fw:v5v82-complex-sectors}
The Fisher identity applies to the real positive density after the exact
V77 cancellations.  Chiral phases and physical Grassmann coefficients are
not squared inside \eqref{eq:v5v82-assembled-Fisher}; they remain governed
by the V58/V69 graded estimates and V78 phase connection.
\end{firewall}

\begin{theorem}[V82 conclusion]
\label{thm:v5v82-status}
The positive derivative required by the NILC is exactly the rooted
order-two Fisher payment
\(\mathbb E_\nu[r^2|dW|^2]\).  A uniform bound on that payment, together
with V81's negative-norm gain, gives a summable transfer defect with half
the negative-norm cutoff exponent.  Stable Higgs shells and rooted compact
gauge shells satisfy explicit sufficient Fisher criteria.  The Euclidean
gravitational conformal direction does not belong to this positive Gibbs
chart without an additional contour, reduction, or ultraviolet input.
\end{theorem}

\begin{proof}
Use Lemma~\ref{lem:v5v82-density-derivative},
Theorem~\ref{thm:v5v82-Fisher-NILC},
Proposition~\ref{prop:v5v82-stable-Fisher},
Corollary~\ref{cor:v5v82-Higgs}, and
Proposition~\ref{prop:v5v82-compact-gauge}, subject to the three firewalls.
\end{proof}

V83 will go upstream to the gravitational obstruction.  It will decompose
the quadratic Euclidean metric action into transverse-traceless, gauge,
and conformal components, prove the sign of each principal block, and state
the minimal contour or constraint datum needed before the joint positive
conditional-measure argument can include an Einstein shell.


\clearpage
\section{V83: a nonlinear mean-zero York constraint chart and its massless tail}
\label{sec:v5v83}

\subsection{Inherited sign results and the actual next problem}

The provisional V82 target has already been answered upstream in this
manuscript.  V9 proves that the negative Euclidean conformal Hessian
survives every ordinary diffeomorphism Faddeev--Popov slice.  V13 proves a
different statement: in linearized vacuum ADM gravity on a flat three-torus,
the Hamiltonian constraint removes every nonzero spatial conformal mode,
leaving two positive transverse--traceless oscillators.  We use both results
without repeating their proofs.

The missing bridge is nonlinear and sourced.  One must solve the
Hamiltonian constraint on a local field chart, retain its global zero-mode
condition, and place the resulting inverse-Laplacian tail in a norm
compatible with the Wilsonian ancestry.  This note constructs the
mean-zero constraint chart near flat data and proves why it cannot belong
to the exponentially local massive kernel class used for the PV sector.

\subsection{The projected Lichnerowicz equation}

Let \(\Sigma=\mathbb T^3\) with a fixed smooth reference volume form.  In
the constant-mean-curvature conformal method write
\begin{equation}
 \gamma=\phi^4\widetilde\gamma,
 \qquad
 K=\phi^{-2}\widetilde A+\frac13\tau\gamma,
 \qquad \phi>0,
 \label{eq:v5v83-York-data}
\end{equation}
where \(\widetilde A\) is transverse traceless with respect to the conformal
metric and \(\tau\) is the mean curvature.  With one conventional choice of
matter conformal weights, the Hamiltonian constraint has the schematic
Lichnerowicz form
\begin{equation}
 \mathcal F(\phi,p)
 :=-8\widetilde\Delta\phi
 +R(\widetilde\gamma)\phi
 -|\widetilde A|_{\widetilde\gamma}^2\phi^{-7}
 +\frac23\tau^2\phi^5
 -16\pi G\,\widetilde\rho\,\phi^{q_\rho}=0.
 \label{eq:v5v83-Lichnerowicz}
\end{equation}
Here \(p=(\widetilde\gamma-\delta,\widetilde A,\tau,widetilde\rho)\), and
the fixed exponent \(q_\rho\) changes with the chosen matter scaling but
does not affect the local argument below.  At flat vacuum data
\begin{equation}
 p=0,
 \qquad \phi=1,
 \qquad \mathcal F(1,0)=0.
 \label{eq:v5v83-flat-data}
\end{equation}

Let \(P_0f=|\Sigma|^{-1}\int_\Sigma f\) and \(P_\perp=I-P_0\).  Fix the
conformal-volume coordinate by writing
\begin{equation}
 \phi=1+u,
 \qquad P_0u=0.
 \label{eq:v5v83-mean-zero-u}
\end{equation}
The nonzero-mode constraint is
\begin{equation}
 \mathcal F_\perp(u,p)
 :=P_\perp\mathcal F(1+u,p)=0.
 \label{eq:v5v83-projected-constraint}
\end{equation}

\begin{theorem}[Analytic mean-zero constraint solution]
\label{thm:v5v83-IFT}
Fix \(s>3/2\).  There are neighbourhoods
\(\mathcal U\subset H^{s+2}_0(\Sigma)\) and \(\mathcal P\) of the flat
data in the corresponding \(H^{s+2}\times H^{s+1}\times\mathbb R\times
H^s\) data space such that for every \(p\in\mathcal P\) there is a unique
\begin{equation}
 u=\mathscr C(p)\in\mathcal U
 \label{eq:v5v83-constraint-map}
\end{equation}
satisfying \eqref{eq:v5v83-projected-constraint}.  The map \(\mathscr C\)
is real analytic, \(1+u\ge1/2\) after shrinking the neighbourhood, and
\begin{equation}
 \|\mathscr C(p)\|_{H^{s+2}}
 \le C\|p\|_{\mathcal P}.
 \label{eq:v5v83-constraint-bound}
\end{equation}
Its derivative at the flat point is
\begin{equation}
 D\mathscr C(0)[\dot p]
 =\frac18(-\Delta)^{-1}_{0}
 P_\perp\,D_p\mathcal F(1,0)[\dot p],
 \label{eq:v5v83-linear-constraint-map}
\end{equation}
with the sign fixed by the convention in
\eqref{eq:v5v83-Lichnerowicz}.
\end{theorem}

\begin{proof}
Because \(s>3/2\), \(H^{s+2}(\mathbb T^3)\) is a Banach algebra and embeds
continuously in \(C^0\).  On the open set \(1+u>0\), the powers
\((1+u)^{-7},(1+u)^5,(1+u)^{q_\rho}\) are analytic Nemytskii maps.
The metric inverse, scalar curvature, Laplacian, and pointwise products are
analytic maps between the displayed Sobolev spaces near the flat metric.
Thus \(\mathcal F_\perp\) is analytic.

At \((u,p)=(0,0)\),
\begin{equation}
 D_u\mathcal F_\perp(0,0)=-8\Delta:
 H^{s+2}_0(\Sigma)\longrightarrow H^s_0(\Sigma).
 \label{eq:v5v83-linearization}
\end{equation}
The mean-zero Laplacian is an isomorphism on the torus.  The analytic
implicit-function theorem gives a unique analytic solution map.  Its local
Lipschitz estimate gives \eqref{eq:v5v83-constraint-bound}.  Sobolev
embedding and a smaller neighbourhood give \(\|u\|_\infty\le1/2\).
Differentiating \(\mathcal F_\perp(\mathscr C(p),p)=0\) at zero and
inverting \(-8\Delta\) proves
\eqref{eq:v5v83-linear-constraint-map}.
\end{proof}

\begin{corollary}[The sourced conformal mode is dependent]
\label{cor:v5v83-conformal-dependent}
On this chart, every nonzero spatial conformal mode is an analytic
functional of the conformal metric, TT momentum, mean curvature, and matter
source data.  It is not an independent positive or negative Gaussian shell
coordinate.
\end{corollary}

\begin{proof}
The conformal factor is \(1+\mathscr C(p)\), uniquely determined by
Theorem~\ref{thm:v5v83-IFT} after its mean is fixed.
\end{proof}

\subsection{The unsolved global Hamiltonian constraint}

Solving the projected equation leaves the scalar zero-mode condition
\begin{equation}
 \mathcal Z(p)
 :=P_0\mathcal F(1+\mathscr C(p),p)=0.
 \label{eq:v5v83-global-constraint}
\end{equation}

\begin{proposition}[Projected solution is not a full solution]
\label{prop:v5v83-global-gate}
For \(p\in\mathcal P\), the pair
\((1+\mathscr C(p),p)\) solves the complete Hamiltonian constraint if and
only if \(\mathcal Z(p)=0\).  At leading nontrivial order around flat data,
the global equation contains the balance
\begin{equation}
 -\int_\Sigma|\widetilde A|^2
 +\frac23\tau^2|\Sigma|
 -16\pi G\int_\Sigma\widetilde\rho
 +\text{metric-curvature and higher-order terms}=0.
 \label{eq:v5v83-global-balance}
\end{equation}
The signs follow the convention in \eqref{eq:v5v83-Lichnerowicz}.
\end{proposition}

\begin{proof}
The projected equation says
\(P_\perp\mathcal F=0\), so \(\mathcal F\) is constant on \(\Sigma\).
It vanishes exactly when its mean vanishes.  Integrate
\eqref{eq:v5v83-Lichnerowicz}; the Laplacian term integrates to zero.
Insert \(\phi=1+O(p)\) and retain the first nonzero terms to obtain
\eqref{eq:v5v83-global-balance}.
\end{proof}

\begin{firewall}[The volume and lapse zero modes remain]
\label{fw:v5v83-zero-mode}
The inverse \(( -\Delta)_0^{-1}\) in
\eqref{eq:v5v83-linear-constraint-map} exists only after fixing the mean of
\(\phi\).  The global equation \eqref{eq:v5v83-global-constraint}, torus
moduli, lapse normalization, and cosmological constant sector require a
separate finite-dimensional reduction.  The local implicit-function
theorem does not solve them.
\end{firewall}

\subsection{Constraint energy and the Coulomb norm}

For a mean-zero source \(f\), let \(u=(-\Delta)_0^{-1}f\).  The natural
constraint norm is
\begin{equation}
 \|f\|_{H^{-1}_0}^2
 :=\langle f,(-\Delta)_0^{-1}f\rangle
 =\|\nabla u\|_2^2.
 \label{eq:v5v83-Coulomb-norm}
\end{equation}

\begin{lemma}[Exact Coulomb payment]
\label{lem:v5v83-Coulomb-payment}
For every mean-zero \(f\in H^{-1}\), the unique mean-zero solution of
\(-\Delta u=f\) satisfies
\begin{equation}
 \|\nabla u\|_2=\|f\|_{H^{-1}_0},
 \qquad
 \|u\|_{H^1}\le C_\Sigma\|f\|_{H^{-1}_0}.
 \label{eq:v5v83-Coulomb-bound}
\end{equation}
\end{lemma}

\begin{proof}
Test the weak equation against \(u\) to obtain the first identity.  The
mean-zero Poincar\'e inequality controls \(\|u\|_2\) by
\(\|\nabla u\|_2\), proving the second.
\end{proof}

\begin{proposition}[Nonlinear constraint Lipschitz estimate]
\label{prop:v5v83-nonlinear-Lipschitz}
After shrinking \(\mathcal P\), there is \(C<\infty\) such that
\begin{equation}
 \|\mathscr C(p)-\mathscr C(p')\|_{H^{s+2}}
 \le C\|P_\perp(\mathfrak S(p)-\mathfrak S(p'))\|_{H^s},
 \label{eq:v5v83-nonlinear-Lipschitz}
\end{equation}
where \(\mathfrak S\) denotes the complete non-Laplacian source in
\eqref{eq:v5v83-Lichnerowicz}, evaluated on the corresponding solutions.
In the energy topology, the leading source-to-conformal map is bounded
from \(H^{-1}_0\) to \(H^1_0\), with the exact payment
\eqref{eq:v5v83-Coulomb-norm}.
\end{proposition}

\begin{proof}
Subtract the two projected equations.  The averaged linearization in the
\(u\) variable is a small perturbation of \(-8\Delta\) and is therefore
invertible from \(H^{s+2}_0\) to \(H^s_0\) with a uniform local bound.
Apply that inverse to the source difference.  The energy statement is
Lemma~\ref{lem:v5v83-Coulomb-payment} applied to the leading linearization.
\end{proof}

\subsection{Why exponential locality is the wrong constraint topology}

On large tori or in infinite volume, the kernel of
\(( -\Delta)^{-1}\) in three dimensions has the Coulomb asymptotic
\begin{equation}
 G(x,y)\sim\frac{1}{4\pi|x-y|}.
 \label{eq:v5v83-Coulomb-tail}
\end{equation}

\begin{theorem}[No volume-uniform exponential constraint kernel]
\label{thm:v5v83-no-exponential}
For an exhausting sequence of spatial tori, there are no constants
\(C,\mu>0\), independent of volume, such that the mean-zero Green kernels
satisfy
\begin{equation}
 |G_L(x,y)|\le Ce^{-\mu d_L(x,y)}
 \label{eq:v5v83-false-exponential}
\end{equation}
for all pairs away from the diagonal.  Equivalently, the volume-uniform
weighted Schur norm with any fixed positive exponential weight diverges.
\end{theorem}

\begin{proof}
Choose separations \(1\ll r\ll L\).  The torus Green kernel converges on
such scales to the Euclidean Green kernel, whose magnitude is bounded below
by \(c/r\) along a subsequence.  The proposed bound would give
\(c/r\le Ce^{-\mu r}\), impossible as \(r\to\infty\).  Multiplication by
\(e^{\mu r}\) then makes the corresponding weighted row sums diverge.
\end{proof}

Thus the PV and massive Higgs inverse kernels may use the exponential
weighted algebra of V70, while the solved gravitational constraint must
use a massless Coulomb/Sobolev grade.  Combining them in one maximum norm
would falsely demand exponential decay of the Newtonian tail.

\begin{definition}[Hybrid constrained-gravity shell certificate]
\label{def:v5v83-HCGSC}
An HCGSC consists of:
\begin{enumerate}
 \item the analytic mean-zero constraint map of
 Theorem~\ref{thm:v5v83-IFT};
 \item a solution of the global equation
 \eqref{eq:v5v83-global-constraint} and the finite-dimensional zero modes;
 \item exponential local grades for massive and regulator kernels;
 \item a separate \(H^{-1}_0\to H^1_0\) Coulomb grade for the constraint
 map, with composition rules against matter source norms;
 \item a time-slice-local block map, so solving the spatial constraint does
 not mix positive and negative Euclidean times;
 \item nonlinear coercivity or a canonical transfer construction on the
 reduced variables;
 \item cutoff-uniform chart transitions and a treatment of topology beyond
 the flat York chart.
\end{enumerate}
\end{definition}

\begin{proposition}[Time-slice reduction is reflection-admissible]
\label{prop:v5v83-time-slice}
If the constraint is solved independently on each time slice and the
spatial Green operator commutes with time reflection, then a retained
positive-time observable pulls back to variables on positive-time slices
only.  Hence the constraint substitution satisfies the algebraic
reflection-admissibility condition of V74.
\end{proposition}

\begin{proof}
The value of \(\mathscr C(p)(t)\) depends on the data \(p(t)\) on the same
time slice and on no data at reflected negative time.  Reflection acts by
relabeling \(t\) and commutes with the fixed spatial elliptic solution.
Both conditions in \eqref{eq:v5v74-reflection-admissible} follow.
\end{proof}

\begin{firewall}[Canonical reduction is not yet a Euclidean measure]
\label{fw:v5v83-canonical-Euclidean}
Proposition~\ref{prop:v5v83-time-slice} is an algebraic compatibility
statement.  It does not construct a positive canonical transfer measure,
prove that its Jacobian and lapse determinant match the covariant
Faddeev--Popov formulation, or establish OS positivity for nonlinear
gravity.
\end{firewall}

\begin{firewall}[Fixed-torus ellipticity is not a thermodynamic bound]
\label{fw:v5v83-torus-gap}
The inverse of the mean-zero Laplacian is bounded on a fixed torus because
its first nonzero eigenvalue is positive.  That eigenvalue vanishes as the
spatial volume grows.  Volume-uniform estimates must use the homogeneous
Coulomb norm and neutral-source conditions, not the fixed-volume Poincar\'e
constant.
\end{firewall}

\begin{theorem}[V83 conclusion]
\label{thm:v5v83-status}
Near flat constant-mean-curvature data, the nonlinear mean-zero
Lichnerowicz constraint has a unique analytic solution with a quantitative
Sobolev bound.  The nonzero conformal factor is therefore a dependent
field, extending the linear V13 conclusion on a local nonlinear chart.
The global Hamiltonian constraint remains as a separate scalar equation,
and the solution map has an unavoidable spatial Coulomb tail.  A joint
SM--Einstein shell norm must therefore combine massive exponential grades
with a massless constraint grade; the positive Gibbs chart of V81--V82
cannot simply include an unconstrained real conformal variable.
\end{theorem}

\begin{proof}
Use Theorem~\ref{thm:v5v83-IFT},
Proposition~\ref{prop:v5v83-global-gate},
Lemma~\ref{lem:v5v83-Coulomb-payment},
Proposition~\ref{prop:v5v83-nonlinear-Lipschitz}, and
Theorem~\ref{thm:v5v83-no-exponential}, subject to the three firewalls.
\end{proof}

V84 will construct the first mixed massive--Coulomb product estimates.  It
will prove how an exponentially localized matter or regulator polymer
sources the constraint tail, determine which neutrality moment removes the
infrared zero mode, and test whether the resulting rooted Coulomb debit is
summable under adjacent cutoff blocking.


\clearpage
\section{V84: mixed massive--Coulomb estimates and the constraint-source payer}
\label{sec:v5v84}

\subsection{The source defect, not the full Newtonian field}

V83 proved that solving the Hamiltonian constraint produces a spatial
Coulomb kernel.  The physical long-range gravitational field should not be
forced to become exponentially local.  Adjacent-cutoff consistency instead
requires the \emph{difference} between two matched constraint sources to be
small in a massless norm.  Conservation and moment matching provide the
possible payers.

Work first on \(\mathbb R^3\).  For a compactly supported source \(f\), set
\begin{equation}
 u=G*f,
 \qquad G(x)=\frac1{4\pi|x|},
 \qquad -\Delta u=f.
 \label{eq:v5v84-Newton-potential}
\end{equation}
The homogeneous constraint norm is
\begin{equation}
 \|f\|_{\dot H^{-1}}:=\|\nabla u\|_2.
 \label{eq:v5v84-Hminus1}
\end{equation}

\subsection{Divergence form gives the exact energy payer}

\begin{theorem}[Conserved-current constraint bound]
\label{thm:v5v84-divergence-payer}
Let \(J\in L^2(\mathbb R^3;\mathbb R^3)\) and
\(f=\nabla\!\cdot J\) in distributions.  Then
\begin{equation}
 \|f\|_{\dot H^{-1}}
 \le\|J\|_2.
 \label{eq:v5v84-divergence-bound}
\end{equation}
Equivalently, the mean-zero weak solution satisfies
\begin{equation}
 \|\nabla(-\Delta)^{-1}\nabla\!\cdot J\|_2
 \le\|J\|_2.
 \label{eq:v5v84-Riesz-projection}
\end{equation}
The same estimate holds on a torus after removing the zero mode, with
constant one and no dependence on the torus size.
\end{theorem}

\begin{proof}
In Fourier variables,
\begin{equation}
 \widehat{\nabla u}(k)
 =-\frac{k\,k\cdot\widehat J(k)}{|k|^2}.
\end{equation}
The multiplier \(k\otimes k/|k|^2\) is an orthogonal rank-one projection
and has norm one.  Plancherel proves
\eqref{eq:v5v84-Riesz-projection}.  The torus proof is identical on every
nonzero Fourier mode.
\end{proof}

\begin{corollary}[Summable constraint-source ancestry]
\label{cor:v5v84-summable-source}
If an adjacent-cutoff constraint-source defect has the Ward form
\begin{equation}
 \delta f_n=\nabla\!\cdot J_n,
 \qquad
 \|J_n\|_2\le Ca_n^\beta,
 \qquad \beta>0,
 \label{eq:v5v84-current-gain}
\end{equation}
then
\begin{equation}
 \sum_n\|\delta f_n\|_{\dot H^{-1}}<\infty.
 \label{eq:v5v84-source-sum}
\end{equation}
The corresponding linear constraint solutions are Cauchy in
\(\dot H^1\).
\end{corollary}

\begin{proof}
Apply Theorem~\ref{thm:v5v84-divergence-payer}; the series
\(\sum_na_n^\beta\) is geometric.  The identity
\(\|\nabla(-\Delta)^{-1}\delta f_n\|_2=
\|\delta f_n\|_{\dot H^{-1}}\) gives the last assertion.
\end{proof}

The divergence identity supplies neutrality on a compact slice.  It does
not supply the positive power in \eqref{eq:v5v84-current-gain}; that power
must come from irrelevant-order matching of the complete stress tensor.

\subsection{Composition with an exponentially local matter kernel}

Let \(K\) be a translation-invariant scalar or finite-matrix convolution
operator with kernel satisfying
\begin{equation}
 \sum_xe^{\mu|x|}\|K(x)\|\le C_K
 \label{eq:v5v84-exponential-K}
\end{equation}
on a lattice, or the analogous integral bound in the continuum.  It acts
componentwise on currents.

\begin{theorem}[Massive localization preserves the Coulomb payment]
\label{thm:v5v84-massive-Coulomb}
If \(f=\nabla\!\cdot J\) and \(K\) commutes with translations and spatial
derivatives, then
\begin{equation}
 Kf=\nabla\!\cdot(KJ),
 \qquad
 \|Kf\|_{\dot H^{-1}}
 \le\|K\|_{2\to2}\|J\|_2
 \le C_K\|J\|_2.
 \label{eq:v5v84-mixed-bound}
\end{equation}
Thus insertion of any volume-uniform exponentially local massive or
regulator response does not worsen the massless constraint norm.
\end{theorem}

\begin{proof}
Convolution commutes with derivatives, proving the first identity.  Apply
Theorem~\ref{thm:v5v84-divergence-payer} to \(KJ\).  Young's convolution
inequality gives \(\|K\|_{2\to2}\le\|K\|_1\le C_K\).
\end{proof}

\begin{proposition}[Background commutator ledger]
\label{prop:v5v84-background-commutator}
For a covariant derivative \(\nabla_A\) and a background-dependent local
kernel \(K_A\),
\begin{equation}
 K_A\nabla_A\!\cdot J
 =\nabla_A\!\cdot(K_AJ)-[\nabla_A\!\cdot,K_A]J.
 \label{eq:v5v84-commutator}
\end{equation}
If
\begin{equation}
 \|[\nabla_A,K_A]J\|_{\dot H^{-1}}
 \le\varepsilon_A\|J\|_2,
 \label{eq:v5v84-commutator-bound}
\end{equation}
then
\begin{equation}
 \|K_A\nabla_A\!\cdot J\|_{\dot H^{-1}}
 \le(\|K_A\|_{2\to2}+\varepsilon_A)\|J\|_2.
 \label{eq:v5v84-covariant-mixed}
\end{equation}
\end{proposition}

\begin{proof}
Equation \eqref{eq:v5v84-commutator} is the definition of the commutator.
Apply the divergence estimate to the first term and
\eqref{eq:v5v84-commutator-bound} to the second.
\end{proof}

The commutator is a curvature and background-gradient debit.  Calling a
kernel ``gauge covariant'' does not make it vanish; it must be estimated in
the massless grade or absorbed by the corresponding Ward-complete source.

\subsection{Multipole cancellation and pointwise tails}

Fix a root \(x_0\), let \(f\) be supported in \(B(x_0,r_0)\), and define
\begin{equation}
 M_m(f;x_0):=int|y-x_0|^m|f(y)|\,dy.
 \label{eq:v5v84-moment-size}
\end{equation}
Say that \(f\) has vanishing moments below order \(m\) when
\begin{equation}
 \int(y-x_0)^\alpha f(y)\,dy=0
 \qquad(|\alpha|<m).
 \label{eq:v5v84-moment-cancellation}
\end{equation}

\begin{theorem}[Matched multipole tail]
\label{thm:v5v84-multipole}
If \eqref{eq:v5v84-moment-cancellation} holds for an integer \(m\ge0\),
then for \(|x-x_0|\ge2r_0\),
\begin{equation}
 |\nabla(G*f)(x)|
 \le C_mM_m(f;x_0)|x-x_0|^{-m-2}.
 \label{eq:v5v84-gradient-tail}
\end{equation}
Consequently
\begin{equation}
 \|\nabla(G*f)\|_{L^2(|x-x_0|>R)}
 \le C_mM_m(f;x_0)R^{-m-1/2}
 \qquad(R\ge2r_0).
 \label{eq:v5v84-energy-tail}
\end{equation}
For \(m=0\) no moment is assumed; \(m=1\) is monopole cancellation, and
\(m=2\) additionally cancels the dipole.
\end{theorem}

\begin{proof}
Taylor-expand \(\nabla G(x-y)\) in \(y-x_0\) through order \(m-1\).
Every displayed polynomial term integrates to zero by
\eqref{eq:v5v84-moment-cancellation}.  Derivatives of order \(m+1\) of
\(G\) are bounded by
\(C_m|x-x_0|^{-m-2}\) when \(|y-x_0|\le|x-x_0|/2\).  The integral
remainder proves \eqref{eq:v5v84-gradient-tail}.  Squaring and integrating
in spherical shells gives
\begin{equation}
 C M_m^2\int_R^\infty r^2r^{-2m-4}\,dr
 =C_mM_m^2R^{-2m-1},
\end{equation}
whose square root is \eqref{eq:v5v84-energy-tail}.
\end{proof}

\begin{assessment}[Which moments adjacent cutoffs must match]
\label{ass:v5v84-moment-owners}
The full physical stress source may have a nonzero monopole and therefore a
Newtonian \(1/r\) field.  The \emph{difference} of two cutoff descriptions
should have its monopole assigned to energy/cosmological matching and its
dipole assigned to translation or momentum Ward matching.  After those
owners agree, the residual begins at \(m=2\) and its gradient tail is
\(O(r^{-4})\).  This is a matching condition to prove, not an automatic
property of lattice discretization.
\end{assessment}

\subsection{Preservation of matched moments}

\begin{lemma}[Convolution preserves vanishing lower moments]
\label{lem:v5v84-moment-preservation}
Let \(K\in L^1\) have finite moments through order \(m\).  If \(f\) has
vanishing moments below order \(m\), then \(K*f\) has the same vanishing
moments.  Moreover,
\begin{equation}
 M_m(K*f;x_0)
 \le C_m\left(\sum_{j=0}^mM_j(K;0)\right)M_m(f;x_0),
 \label{eq:v5v84-moment-bound}
\end{equation}
where lower moments of \(f\) are bounded by the support radius times
\(M_m(f)\).
\end{lemma}

\begin{proof}
For \(|\alpha|<m\), expand
\((z+y-x_0)^\alpha\) in the convolution integral.  Every term contains a
moment of \(f\) of order at most \(|\alpha|\), hence vanishes.  For order
\(m\), use
\(|z+y-x_0|^m\le C_m(|z|^m+|y-x_0|^m)\), Tonelli's theorem, and the finite
kernel moments.
\end{proof}

Every exponentially local kernel has all finite moments, so the free
massive response in Theorem~\ref{thm:v5v84-massive-Coulomb} preserves any
finite set of matched multipoles.  In a nontranslation-invariant
background, the corresponding statement becomes a Ward/moment commutator
estimate analogous to \eqref{eq:v5v84-commutator-bound}.

\subsection{A mixed norm and a sufficient adjacent-cutoff gate}

Define the rooted hybrid source norm
\begin{equation}
 \|f\|_{\mathfrak C_m(x_0)}
 :=\|f\|_{\dot H^{-1}}+M_m(f;x_0),
 \label{eq:v5v84-hybrid-source-norm}
\end{equation}
on sources with moments below \(m\) matched to zero.

\begin{theorem}[Mixed massive--constraint summability]
\label{thm:v5v84-mixed-summability}
Let \(\delta f_n\) be rooted adjacent-cutoff source defects satisfying
\begin{equation}
 \delta f_n=\nabla\!\cdot J_n,
 \qquad
 \|J_n\|_2+M_m(\delta f_n;x_0)
 \le Ca_n^\beta,
 \qquad \beta>0,
 \label{eq:v5v84-hybrid-gain}
\end{equation}
and the moment conditions \eqref{eq:v5v84-moment-cancellation}.  Let
\(K_n\) be massive response kernels with uniform exponential moment bounds
and the required Ward commutator bounded by \(C'a_n^\beta\).  Then
\begin{equation}
 \sum_n\|K_n\delta f_n\|_{\dot H^{-1}}<\infty,
 \label{eq:v5v84-mixed-Hminus-sum}
\end{equation}
and the remote constraint-gradient tails are dominated by the summable
majorant
\begin{equation}
 C a_n^\beta R^{-m-1/2}.
 \label{eq:v5v84-mixed-tail-majorant}
\end{equation}
\end{theorem}

\begin{proof}
Theorems~\ref{thm:v5v84-divergence-payer} and
\ref{thm:v5v84-massive-Coulomb}, together with the commutator estimate, give
\(\|K_n\delta f_n\|_{\dot H^{-1}}\le C''a_n^\beta\).  Sum the geometric
series.  Lemma~\ref{lem:v5v84-moment-preservation} controls the output
moment, and Theorem~\ref{thm:v5v84-multipole} gives
\eqref{eq:v5v84-mixed-tail-majorant}.
\end{proof}

\begin{firewall}[Projection to mean zero is not local compensation]
\label{fw:v5v84-mean-projection}
On a torus, \(P_\perp f=f-|\Sigma|^{-1}\int f\) solves the zero-mode
compatibility condition by adding a uniform background.  This operation is
not a rooted local counterterm.  The global Hamiltonian constraint of V83
must own that mean; local multipole cancellation cannot be inferred merely
from applying \(P_\perp\).
\end{firewall}

\begin{firewall}[Stress conservation does not give irrelevant size]
\label{fw:v5v84-conservation-size}
A Ward identity can put a source difference in divergence form and thereby
prove \eqref{eq:v5v84-divergence-bound}.  It does not prove
\(\|J_n\|_2=O(a_n^\beta)\).  The positive cutoff power requires matched
relevant stress coefficients and a quantitative discretization or shell
remainder estimate.
\end{firewall}

\begin{theorem}[V84 conclusion]
\label{thm:v5v84-status}
A divergence-form constraint-source defect is paid exactly by the
\(L^2\) norm of its conserved current, uniformly in spatial volume.
Exponentially local massive response kernels preserve that Coulomb payment
and, in the translation-invariant case, every matched lower multipole.
Monopole and dipole matching improve the remote constraint-gradient tail
by explicit powers.  These results give a sufficient mixed summability
gate, but the required positive cutoff power and global zero-mode owner
remain unproved for the complete SM stress tensor.
\end{theorem}

\begin{proof}
Use Theorem~\ref{thm:v5v84-divergence-payer},
Theorem~\ref{thm:v5v84-massive-Coulomb},
Theorem~\ref{thm:v5v84-multipole},
Lemma~\ref{lem:v5v84-moment-preservation}, and
Theorem~\ref{thm:v5v84-mixed-summability}, subject to the two firewalls.
\end{proof}

V85 is the next compulsory YM/NS audit.  It will then compute the
adjacent-cutoff stress-tensor defect for the free gauge, Higgs, fermion,
and TT-graviton blocks, test its monopole and dipole moments after canonical
matching, and identify the first interacting Ward remainder which lacks a
positive cutoff power.


\clearpage
\section{V85: companion audit and the free adjacent-cutoff stress defect}
\label{sec:v5v85}

\subsection{Compulsory companion audit}

At the time of inspection, the highest-named Yang--Mills file was
\begin{equation}
 \begin{split}
 &\texttt{Renewal\_Yang\_Mills\_Mass\_Gap\_Research\_Notes\_V93.tex},\\
 &1\,365\,654\text{ bytes},\qquad
 \operatorname{SHA256}=\\
 &\texttt{F83A4A67010765B1E8427AFF21C22099F91695190632CCA559AE3B1F51055E71}.
 \end{split}
 \label{eq:v5v85-YM-fingerprint}
\end{equation}
It was byte-identical to the simultaneously present V92 file and contained
no appended V93 mathematical section.  Its actual endpoint is therefore
V92.  That endpoint replaces the logarithm chart by global smooth
quaternion coordinates, closes the complete two-label matrix RSPCC in the
terminal negative norm, and recovers the ternary local tree card by that
route.  It does not prove the three-label RSPCC or the all-order local card.
Neither companion file was altered.

The active Navier--Stokes file remained
\begin{equation}
 \begin{split}
 &\texttt{Renewal\_NS\_Stability\_Research\_Notes\_V31.tex},\\
 &887\,537\text{ bytes},\qquad
 \operatorname{SHA256}=\\
 &\texttt{F1121B62238AD5491BB7CF03635EA9934942EDF573ED8FAAA9EE79E1D38A6696}.
 \end{split}
 \label{eq:v5v85-NS-fingerprint}
\end{equation}
Its root-before-first-moment warning is unchanged.

The audit selects a fixed-arity task.  The gravitational constraint couples
first to the one-graviton stress vertex.  We compute the adjacent-cutoff
free defect and its external-momentum moments before taking a norm.  No
all-order Yang--Mills input is imported.

\subsection{Moment matching as exact Laurent divisibility}

Let \(f:\mathbb Z^d\to\mathbb C^N\) have finite support and Fourier--Laurent
polynomial
\begin{equation}
 \widehat f(z)=\sum_{x\in\mathbb Z^d}f(x)z^x,
 \qquad z^x=\prod_{j=1}^dz_j^{x_j}.
 \label{eq:v5v85-Laurent-transform}
\end{equation}
Let \(\mathfrak m=(z_1-1,\ldots,z_d-1)\) be the maximal ideal at
\(z=(1,\ldots,1)\).

\begin{theorem}[Exact moment--difference correspondence]
\label{thm:v5v85-Laurent-divisibility}
The following statements hold componentwise:
\begin{enumerate}
 \item \(\sum_xf(x)=0\) if and only if
 \begin{equation}
  \widehat f(z)=\sum_i(z_i-1)\widehat J_i(z)
  \label{eq:v5v85-single-difference}
 \end{equation}
 for Laurent polynomials \(\widehat J_i\).
 \item If additionally \(\sum_xx_jf(x)=0\) for every \(j\), then
 \begin{equation}
  \widehat f(z)=\sum_{i,j}(z_i-1)(z_j-1)
  \widehat M_{ij}(z)
  \label{eq:v5v85-double-difference}
 \end{equation}
 for Laurent polynomials \(\widehat M_{ij}\).
\end{enumerate}
Equivalently, a finite-range source with matched monopole is a finite-range
lattice divergence, and one with matched monopole and dipole is a
finite-range double divergence.
\end{theorem}

\begin{proof}
Evaluation at \(z=1\) maps the Laurent polynomial ring onto \(\mathbb C\)
with kernel \(\mathfrak m\).  This proves the first algebraic statement.
Explicitly, telescope each monomial \(z^x-1\) one coordinate at a time;
negative powers obey
\(z_j^{-n}-1=-(z_j-1)(z_j^{-1}+\cdots+z_j^{-n})\).

The quotient \(\mathfrak m/\mathfrak m^2\) has basis represented by
\(z_j-1\).  The coefficient of this class is
\(\partial_{z_j}\widehat f(1)=\sum_xx_jf(x)\).  Hence vanishing of all first
moments places \(\widehat f\) in \(\mathfrak m^2\), which is generated by
the products \((z_i-1)(z_j-1)\).  This gives
\eqref{eq:v5v85-double-difference}.  Inverse Fourier transformation turns
the factors \(z_j-1\) into lattice differences.
\end{proof}

\begin{corollary}[Constraint norm of a double difference]
\label{cor:v5v85-double-Hminus}
Let \(\nabla_{a,j}\) be the physical lattice difference with symbol
\((e^{iak_j}-1)/a\).  If
\begin{equation}
 f=\sum_{i,j}\nabla_{a,i}\nabla_{a,j}M_{ij},
 \label{eq:v5v85-physical-double}
\end{equation}
then
\begin{equation}
 \|f\|_{\dot H^{-1}_a}
 \le C_d\|\nabla_aM\|_2.
 \label{eq:v5v85-double-Hminus}
\end{equation}
\end{corollary}

\begin{proof}
In Fourier variables, divide the symbol in
\eqref{eq:v5v85-physical-double} by the lattice Laplacian modulus
\(|\widehat\nabla_a(k)|\).  One difference remains, and the finite sum over
\(i,j\) costs only a dimension-dependent constant.  Apply Plancherel.
\end{proof}

\subsection{Adjacent free symbols}

For one momentum component define
\begin{equation}
 \widehat p_a(p)=\frac2a\sin\frac{ap}{2},
 \qquad
 \overline p_a(p)=\frac1a\sin(ap),
 \qquad
 w_a(p)=\frac r a(1-\cos ap).
 \label{eq:v5v85-lattice-momenta}
\end{equation}
Let the adjacent coarse spacing be \(a_c=qa\), where \(q>1\) is fixed.

\begin{lemma}[Sharp fixed-momentum symbol defects]
\label{lem:v5v85-symbol-defects}
For \(|p|\le\Lambda\) and \(qa\Lambda\le c<\pi\),
\begin{align}
 \widehat p_{qa}(p)^2-\widehat p_a(p)^2
 &=-\frac{q^2-1}{12}a^2p^4+O_c(a^4p^6),
 \label{eq:v5v85-Laplacian-defect}\\
 \overline p_{qa}(p)-\overline p_a(p)
 &=-\frac{q^2-1}{6}a^2p^3+O_c(a^4p^5),
 \label{eq:v5v85-Dirac-defect}\\
 w_{qa}(p)-w_a(p)
 &=\frac r2(q-1)ap^2+O_c(a^3p^4).
 \label{eq:v5v85-Wilson-defect}
\end{align}
The remainders and their finitely many momentum derivatives are uniform on
the displayed band.
\end{lemma}

\begin{proof}
Insert the Taylor expansions
\(\sin x=x-x^3/6+O_c(x^5)\) and
\(1-\cos x=x^2/2-x^4/24+O_c(x^6)\), first with \(x=ap\) and then with
\(x=qap\).  Subtraction gives the three formulas.  Taylor's theorem on the
compact interval \([-c,c]\) gives uniform differentiated remainders.
\end{proof}

Thus the free scalar, transverse gauge, and TT-graviton Laplacian blocks
have an \(O(a^2)\) adjacent symbol defect on a fixed physical band.  The
symmetric Dirac derivative also has order two, while an unimproved Wilson
term leaves an \(O(a)\) defect.  This order-one irrelevant fermion operator
sets the minimum free gain unless it is Symanzik-improved or included in the
coarse perfect action of V73.

\subsection{One-graviton stress vertices}

For a free quadratic kernel \(K_a(g)\), define its weak one-graviton vertex
at the flat metric by
\begin{equation}
 \Gamma_a^{\mu\nu}(k;p)
 :=\left.
 \frac{\delta K_a(g;p+k,p)}{\delta h_{\mu\nu}(k)}
 \right|_{g=\delta}.
 \label{eq:v5v85-stress-vertex}
\end{equation}
The source defect is
\begin{equation}
 \delta\Gamma_a^{\mu\nu}
 :=\Gamma_{qa}^{\mu\nu}-\Gamma_a^{\mu\nu}
 \label{eq:v5v85-adjacent-stress}
\end{equation}
after canonical field, mass, vacuum-energy, and kinetic normalization have
been matched.

\begin{theorem}[Free band-limited stress-defect orders]
\label{thm:v5v85-free-stress-orders}
Fix a physical momentum band
\(|p|,|k|\le\Lambda\).  For centered symmetric blocks and matched free
normalizations, the scalar, transverse gauge, and TT-graviton vertices
satisfy
\begin{equation}
 \|\delta\Gamma_a(k;p)\|
 \le C_\Lambda a^2,
 \label{eq:v5v85-boson-stress-order}
\end{equation}
while the Wilson fermion vertex satisfies
\begin{equation}
 \|\delta\Gamma_a(k;p)\|
 \le C_\Lambda a
 \label{eq:v5v85-fermion-stress-order}
\end{equation}
without improvement, and \(O(a^2)\) after cancellation of the dimension-five
Wilson stress vertex.  The same bounds hold for a fixed finite number of
external-momentum derivatives.
\end{theorem}

\begin{proof}
The free scalar, gauge-fixed transverse gauge, and TT-graviton kernels are
finite matrix polynomials in the lattice Laplacian symbols and bounded
projectors on the fixed band.  Their metric variations are finite linear
combinations of the same symbols and their momentum derivatives.  Apply
\eqref{eq:v5v85-Laplacian-defect}.  The fermion kernel is a finite matrix
combination of \(\overline p_a\), \(w_a\), masses, and gamma matrices.
Equations \eqref{eq:v5v85-Dirac-defect} and
\eqref{eq:v5v85-Wilson-defect} show that the Wilson term dominates at
order \(a\).  Removing its matched dimension-five vertex leaves the
order-\(a^2\) symmetric-derivative defect.  Uniformity of differentiated
remainders follows from Lemma~\ref{lem:v5v85-symbol-defects}.
\end{proof}

\subsection{Exact monopole and dipole ownership}

Let \(\delta\tau_a(x)\) be the spatial constraint-source defect obtained
from \eqref{eq:v5v85-adjacent-stress} for fixed smooth external fields.  Its
spatial Fourier transform is a Laurent polynomial in the external lattice
momentum because the underlying stencil has finite range.

\begin{proposition}[Matched free stress is a double divergence]
\label{prop:v5v85-stress-double-divergence}
Suppose the normalization conditions impose
\begin{equation}
 \widehat{\delta\tau_a}(0)=0
 \label{eq:v5v85-monopole-match}
\end{equation}
and the block centroid and translation Ward normalization impose
\begin{equation}
 \partial_{k_j}\widehat{\delta\tau_a}(0)=0
 \qquad(j=1,2,3).
 \label{eq:v5v85-dipole-match}
\end{equation}
Then there is a finite-range tensor \(M_{a,ij}\) such that
\begin{equation}
 \delta\tau_a=\sum_{i,j}\nabla_{a,i}\nabla_{a,j}M_{a,ij}.
 \label{eq:v5v85-stress-double}
\end{equation}
On a fixed physical momentum band,
\begin{align}
 \|\nabla_aM_a\|_2
 &\le C_\Lambda a^2\mathcal N_{\rm free}
 &&\text{for scalar, gauge, and TT blocks},
 \label{eq:v5v85-boson-current}\\
 \|\nabla_aM_a\|_2
 &\le C_\Lambda a\mathcal N_{\rm free}
 &&\text{for an unimproved Wilson fermion},
 \label{eq:v5v85-fermion-current}
\end{align}
where \(\mathcal N_{\rm free}\) is a fixed Sobolev product norm of the
band-limited external fields.
\end{proposition}

\begin{proof}
Equations \eqref{eq:v5v85-monopole-match}--
\eqref{eq:v5v85-dipole-match} are precisely the zeroth and first spatial
moments.  Apply Theorem~\ref{thm:v5v85-Laurent-divisibility} to obtain
\eqref{eq:v5v85-stress-double}.  On a compact momentum band, division by
the two external lattice differences is regular because the numerator has
a second-order zero at the origin; Taylor's integral remainder bounds the
quotient and one remaining difference.  The coefficient orders are those
of Theorem~\ref{thm:v5v85-free-stress-orders}.  Standard band-limited
Sobolev product estimates control the bilinear external fields by
\(\mathcal N_{\rm free}\).
\end{proof}

\begin{corollary}[Free constraint-source gains]
\label{cor:v5v85-free-Hminus}
Under the hypotheses of Proposition
\ref{prop:v5v85-stress-double-divergence},
\begin{align}
 \|\delta\tau_a\|_{\dot H^{-1}_a}
 &\le C_\Lambda a^2\mathcal N_{\rm free}
 &&\text{for the free bosonic blocks},
 \label{eq:v5v85-boson-Hminus}\\
 \|\delta\tau_a\|_{\dot H^{-1}_a}
 &\le C_\Lambda a\mathcal N_{\rm free}
 &&\text{for the unimproved free Wilson block}.
 \label{eq:v5v85-fermion-Hminus}
\end{align}
Both gains are summable along a geometric cutoff ancestry.
\end{corollary}

\begin{proof}
Use Corollary~\ref{cor:v5v85-double-Hminus} and
\eqref{eq:v5v85-boson-current}--\eqref{eq:v5v85-fermion-current}.  Positive
powers of a geometric spacing are summable.
\end{proof}

\subsection{The exact scope of the free result}

\begin{firewall}[Fixed physical bands do not control cutoff modes]
\label{fw:v5v85-bandlimited}
The Taylor estimates require \(a\Lambda\ll1\).  A Wilsonian shell contains
momenta of order \(a^{-1}\), where the dimensionless adjacent symbols differ
by order one.  Theorem~\ref{thm:v5v85-free-stress-orders} therefore proves
the retained smooth-field source gain, not a global operator-norm estimate
on the full shell.  High modes must be integrated by the exact marginal and
rooted determinant methods of V74--V78.
\end{firewall}

\begin{firewall}[Moment conditions are normalization equations]
\label{fw:v5v85-moment-not-automatic}
Translation covariance gives a Ward identity but does not by itself make
two cutoff stress tensors have the same total energy or centroid.
Equations \eqref{eq:v5v85-monopole-match}--
\eqref{eq:v5v85-dipole-match} must be imposed and solved as vacuum,
cosmological, kinetic, and block-centering normalization conditions.
\end{firewall}

The first interacting source not covered by the free theorem is the
one-graviton, one-gauge, two-fermion Ward packet, together with its fermion
self-energy and Wilson contact partners.  Its divergence is constrained by
the translation and gauge Ward identities only after those terms are
assembled.  Estimating the vertex alone would discard the cancellation and
need not produce the \(O(a)\) current in
\eqref{eq:v5v84-current-gain}.

\begin{theorem}[V85 conclusion]
\label{thm:v5v85-status}
The compulsory audit records a byte-identical YM V93/V92 endpoint and
unchanged NS V31.  For finite-range lattice sources, exact matching of the
monopole and dipole is equivalent to an exact double-divergence
factorization.  On every fixed physical momentum band, adjacent free
scalar, gauge, and TT-graviton stress defects have an \(O(a^2)\)
constraint-source gain, while an unimproved Wilson fermion has an
\(O(a)\) gain.  These gains are summable but do not cover cutoff-scale
modes or the first interacting gauge--fermion Ward packet.
\end{theorem}

\begin{proof}
Use Theorem~\ref{thm:v5v85-Laurent-divisibility},
Lemma~\ref{lem:v5v85-symbol-defects},
Theorem~\ref{thm:v5v85-free-stress-orders},
Proposition~\ref{prop:v5v85-stress-double-divergence}, and
Corollary~\ref{cor:v5v85-free-Hminus}, subject to the two firewalls.
\end{proof}

V86 will assemble the one-graviton/one-gauge/two-fermion lattice Ward
packet.  It will derive the exact translation Ward identity for the Wilson
Dirac kernel, include the contact and self-energy terms, and test whether
the matched packet has a divergence current with the free \(O(a)\) cutoff
gain before any rooted cumulant estimate is invoked.


\clearpage
\section{V86: the mixed graviton--gauge--fermion Ward packet}
\label{sec:v5v86}

\subsection{Ward functionals before vertex extraction}

V85 identified the first interacting stress source not covered by its free
symbol theorem: the one-graviton, one-gauge, two-fermion packet.  Its
individual vertices are not conserved.  The correct object is obtained by
differentiating the complete translation Ward functional, which retains
the fermion two-point and contact partners.

Let \(\Gamma_a[h,A,\psi,\bar\psi]\) be a finite-cutoff one-particle
irreducible functional, or the bare action at tree level, in a fixed local
gauge and metric chart.  Let \(\mathcal W_{a,\nu}^{\rm tr}\) be the local
translation Ward operator.  We write its exact identity as
\begin{equation}
 \mathcal W_{a,\nu}^{\rm tr}\Gamma_a
 =\mathcal B_{a,\nu},
 \label{eq:v5v86-translation-Ward}
\end{equation}
where \(\mathcal B_{a,\nu}=0\) for an exactly translation-gauged
discretization and otherwise records the local hypercubic translation
breaking.  Exact lattice gauge invariance is
\begin{equation}
 \mathcal W_a^{\rm g}(\omega)\Gamma_a=0.
 \label{eq:v5v86-gauge-Ward}
\end{equation}
The Ward operators act on every field, source, gauge-fixing partner, and
Jacobian in \(\Gamma_a\).

Define the complete mixed translation packet at zero fields by
\begin{equation}
 \mathfrak P_{a,\nu}^{hA\bar\psi\psi}
 :=left.
 \frac{\delta^4}{\delta h\,\delta A\,
  \delta\psi\,\delta\bar\psi}
 \bigl(\mathcal W_{a,\nu}^{\rm tr}\Gamma_a\bigr)
 \right|_0.
 \label{eq:v5v86-complete-packet}
\end{equation}
The notation includes all placements of the Ward derivative on the four
external fields.

\begin{theorem}[Differentiated Ward packet identity]
\label{thm:v5v86-differentiated-Ward}
At every finite cutoff,
\begin{equation}
 \mathfrak P_{a,\nu}^{hA\bar\psi\psi}
 =left.
 \frac{\delta^4\mathcal B_{a,\nu}}
 {\delta h\,\delta A\,\delta\psi\,\delta\bar\psi}
 \right|_0.
 \label{eq:v5v86-packet-breaking}
\end{equation}
The left side is the assembled sum of:
\begin{enumerate}
 \item the divergence of the \(hA\bar\psi\psi\) contact vertex;
 \item the two momentum-shifted \(A\bar\psi\psi\) vertices;
 \item the two momentum-shifted \(h\bar\psi\psi\) vertices;
 \item the fermion two-point functions generated when the translation acts
 on either external fermion;
 \item the term in which translation acts on the external gauge field;
 \item gauge-fixing, ghost, and measure contacts required by the chosen
 Ward realization.
\end{enumerate}
Thus the packet vanishes identically when \(\mathcal B_{a,\nu}=0\), but no
individual item is asserted to vanish.
\end{theorem}

\begin{proof}
Differentiate \eqref{eq:v5v86-translation-Ward} by the four declared
sources and evaluate at zero.  The functional Leibniz rule distributes the
derivatives over every field coefficient in
\(\mathcal W_{a,\nu}^{\rm tr}\), producing precisely the listed contact and
lower-vertex terms.  The right side is the corresponding derivative of the
breaking functional.  No estimate or perturbative expansion is used.
\end{proof}

\begin{assessment}[Self-energy ownership]
\label{ass:v5v86-self-energy}
The fermion two-point terms in the packet are the inverse propagator at tree
level and the full self-energy after quantum shell integration.  Removing
them before taking the divergence converts an exact Ward identity into a
false conservation statement for the contact vertex alone.
\end{assessment}

\subsection{Momentum-space two-point Ward identity}

For the three-point stress vertex, use momenta
\(p\) and \(p+k\) on the fermion legs.  Let \(\Sigma_a(p)\) be the complete
inverse two-point function and let \(\pi_{a,\nu}(p)\) be the periodic
symbol representing the chosen lattice translation on a fermion leg.

\begin{proposition}[Exact stress--self-energy relation]
\label{prop:v5v86-stress-self-energy}
The two-fermion derivative of \eqref{eq:v5v86-translation-Ward} has the
form
\begin{equation}
 \widehat k_{a,\mu}
 \Gamma_{a}^{\mu}{}_{\nu;h\bar\psi\psi}(k;p+k,p)
 =pi_{a,\nu}(p)\Sigma_a(p+k)
  -\pi_{a,\nu}(p+k)\Sigma_a(p)
  +B_{a,\nu}^{\bar\psi\psi}(k;p),
 \label{eq:v5v86-stress-Ward-momentum}
\end{equation}
where \(B_{a,\nu}^{\bar\psi\psi}\) is the corresponding derivative of
\(\mathcal B_{a,\nu}\).  Sign changes caused by an alternative all-incoming
momentum convention affect all three terms together and not the assertion.
\end{proposition}

\begin{proof}
Fourier-transform the local Ward identity after two fermion derivatives.
The divergence acting on the stress insertion gives
\(\widehat k_{a,\mu}\).  Translation acting on each external fermion gives
the two contact terms with its incoming or outgoing translation symbol.
The remaining term is the differentiated breaking functional.
\end{proof}

Differentiating \eqref{eq:v5v86-stress-Ward-momentum} with respect to the
gauge background is only part of
Theorem~\ref{thm:v5v86-differentiated-Ward}: translation of the gauge field
itself supplies the missing gauge-vertex contact.  This is why the complete
functional packet is the stable definition.

\subsection{Exact finite-range divided differences}

There is an algebraic lattice version of the divergence factor which does
not divide by a momentum zero.  Let \(D_a(z)\) be any finite-range free
Dirac Laurent polynomial and let \(w_j=z_j'-z_j\) represent an external
momentum shift.

\begin{lemma}[Finite-range divided-difference vertex]
\label{lem:v5v86-divided-difference}
There exist matrix Laurent polynomials \(V_{a,j}(z',z)\) such that
\begin{equation}
 D_a(z')-D_a(z)
 =\sum_j(z_j'-z_j)V_{a,j}(z',z).
 \label{eq:v5v86-divided-difference}
\end{equation}
Their hopping range is bounded by that of \(D_a\), up to a fixed factor,
and they inherit every linear internal symmetry of \(D_a\).
\end{lemma}

\begin{proof}
For a monomial \(z^n\), change the coordinates from \(z\) to \(z'\) one
at a time.  In one coordinate,
\begin{equation}
 (z')^m-z^m=(z'-z)\sum_{r=0}^{m-1}(z')^{m-1-r}z^r
\end{equation}
for \(m>0\), with the analogous finite identity for negative powers.
Apply this to every monomial of \(D_a\).  Finite range and linear symmetry
inheritance are immediate.
\end{proof}

This lemma supplies a finite contact vertex satisfying the free gauge or
translation difference identity exactly.  A metric vertex obtained by a
different lattice covariantization may differ by a transverse improvement
and by the explicit breaking \(\mathcal B_a\); those pieces must stay in
the packet.

\subsection{Smooth-background Wilson expansion}

For a smooth gauge connection, covariant parallel transport along one link
gives the formal covariant Taylor identities
\begin{align}
 \frac{U_\mu(x)\psi(x+a\hat\mu)-
       U_\mu(x-a\hat\mu)^*\psi(x-a\hat\mu)}{2a}
 &=D_\mu\psi+\frac{a^2}{6}D_\mu^3\psi+O(a^4),
 \label{eq:v5v86-covariant-central}\\
 \frac{2\psi-U_\mu\psi_+-U_\mu^*\psi_-}{a}
 &=-aD_\mu^2\psi-\frac{a^3}{12}D_\mu^4\psi+O(a^5),
 \label{eq:v5v86-covariant-Wilson}
\end{align}
with signs adjusted if the lattice Laplacian is defined oppositely.

\begin{theorem}[Tree-level order of the complete mixed packet]
\label{thm:v5v86-tree-packet-order}
On a fixed smooth momentum/background band, the Wilson fermion action has
the Symanzik form
\begin{equation}
 S_{a,\psi}
 =S_{0,\psi}+a\,\mathcal O_{5,W}+a^2\mathcal R_{6,a},
 \label{eq:v5v86-Symanzik}
\end{equation}
where \(\mathcal O_{5,W}\) is the gauge- and metric-covariant dimension-five
Wilson operator and \(\mathcal R_{6,a}\) has uniformly bounded fixed-arity
vertices on that band.  After differentiating the complete translation
Ward functional,
\begin{equation}
 \mathfrak P_{a,\nu}^{hA\bar\psi\psi}
 =a\,
 \left.
 \frac{\delta^4
  (\mathcal W_{0,\nu}^{\rm tr}\mathcal O_{5,W})}
 {\delta h\,\delta A\,\delta\psi\,\delta\bar\psi}
 \right|_0
 +O(a^2).
 \label{eq:v5v86-packet-Oa}
\end{equation}
If the matched coarse action includes the complete dimension-five Wilson
stress operator and all its contact partners, the residual packet is
\(O(a^2)\).
\end{theorem}

\begin{proof}
Equations \eqref{eq:v5v86-covariant-central}--
\eqref{eq:v5v86-covariant-Wilson} give
\eqref{eq:v5v86-Symanzik}; metric and gauge differentiation preserves the
power of \(a\) on a fixed smooth band.  Apply the linear Ward operator and
four fixed functional derivatives to the expansion.  The continuum term
satisfies the exact continuum translation Ward identity after its contact
terms are included.  The first remaining term is therefore the displayed
dimension-five packet.  Matching that entire local operator, rather than
one selected vertex, removes it and leaves the bounded order-six
remainder.
\end{proof}

\subsection{From the packet to a constraint current}

Let \(\delta\tau_{a}^{hA\bar\psi\psi}\) be the adjacent-cutoff spatial
Hamiltonian-constraint source extracted from the complete packet after
matching its zero-momentum stress coefficient.  The Ward identity gives a
lattice divergence plus the Symanzik breaking.

\begin{corollary}[Fixed-arity interacting current gain]
\label{cor:v5v86-current-gain}
On a fixed smooth band, there is a finite-range current
\(J_a^{hA\bar\psi\psi}\) such that
\begin{equation}
 \delta\tau_a^{hA\bar\psi\psi}
 =\nabla_a\!\cdot J_a^{hA\bar\psi\psi},
 \qquad
 \|J_a^{hA\bar\psi\psi}\|_2
 \le Ca\,\mathcal N_{hA\bar\psi\psi},
 \label{eq:v5v86-current-Oa}
\end{equation}
where \(\mathcal N_{hA\bar\psi\psi}\) is a fixed Sobolev product norm of
the external fields.  With complete dimension-five improvement, the right
side is \(Ca^2\mathcal N_{hA\bar\psi\psi}\).  Hence the corresponding
constraint-source defect is summable in \(\dot H^{-1}\).
\end{corollary}

\begin{proof}
After zero-momentum matching, the finite-range Laurent source has zero
monopole.  Theorem~\ref{thm:v5v85-Laurent-divisibility} gives its divergence
form.  The quotient is bounded on the fixed band by Taylor's integral
remainder, and Theorem~\ref{thm:v5v86-tree-packet-order} supplies the
coefficient.  Apply Theorem~\ref{thm:v5v84-divergence-payer} and sum the
geometric powers of \(a\).
\end{proof}

\begin{firewall}[The theorem is fixed-arity and smooth-band]
\label{fw:v5v86-fixed-arity}
Theorem~\ref{thm:v5v86-tree-packet-order} and
Corollary~\ref{cor:v5v86-current-gain} control one fixed external packet on
smooth backgrounds.  They do not bound cutoff-scale loop momenta, an
arbitrary number of gauge insertions, or the rooted quantum cumulant sum.
Those require the exact high-shell integration and all-order companion
gates.
\end{firewall}

\begin{firewall}[A Ward identity does not bound the assembled current]
\label{fw:v5v86-Ward-not-bound}
At the quantum effective-action level, exact gauge and translation Ward
identities determine the divergence and contact structure.  They do not
give a dimension-free norm bound on the full self-energy/contact packet.
Wavefunction, gauge, mass, and stress counterterms must be matched, and the
remaining rooted current must still be estimated.
\end{firewall}

\begin{firewall}[Discrete translations require a declared Ward
realization]
\label{fw:v5v86-discrete-translation}
A hypercubic lattice has exact discrete translations but not automatic
local continuous translation invariance.  The breaking functional
\(\mathcal B_{a,\nu}\) depends on the chosen lattice metric coupling and
energy--momentum operator.  Setting it to zero without constructing that
realization would assume the desired stress Ward identity.
\end{firewall}

\begin{theorem}[V86 conclusion]
\label{thm:v5v86-status}
The first interacting gravitational source must be treated as the complete
differentiated translation Ward packet containing the
graviton--gauge--fermion contact, gauge and stress vertices, fermion
two-point functions, and measure partners.  For the tree-level Wilson
action on a fixed smooth band, the assembled packet is \(O(a)\), with its
coefficient owned by the dimension-five Wilson operator; complete
improvement raises it to \(O(a^2)\).  This yields a summable fixed-arity
constraint current, while the quantum rooted norm and cutoff-scale modes
remain open.
\end{theorem}

\begin{proof}
Use Theorem~\ref{thm:v5v86-differentiated-Ward},
Proposition~\ref{prop:v5v86-stress-self-energy},
Lemma~\ref{lem:v5v86-divided-difference},
Theorem~\ref{thm:v5v86-tree-packet-order}, and
Corollary~\ref{cor:v5v86-current-gain}, subject to the three firewalls.
\end{proof}

V87 will address the quantum fixed-arity step.  It will express the
renormalized mixed Ward packet as a difference of normalized conditional
expectations, use the V81 Poisson--Stein duality in its natural negative
norm, and identify the positive shell derivative needed to turn the
tree-level \(O(a)\) Symanzik remainder into a summable rooted transfer
error.


\clearpage
\section{V87: a Stein estimate for the quantum fixed-arity Ward packet}
\label{sec:v5v87}

\subsection{Why a fixed packet is cheaper than a full transfer norm}

V80 used positive shell regularity to convert a negative density norm into
an \(L^2\) norm uniformly over an entire cylinder class.  The mixed Ward
packet of V86 is one specified observable.  For such a fixed observable,
the Poisson duality of V81 pairs the negative density defect directly with
one shell derivative of the already assembled packet.  No positive
derivative of the density is needed at this step.

Let \(\widetilde\kappa(dy\mid x)\) be a normalized positive reference shell
law and
\begin{equation}
 d\kappa=(1+h)d\widetilde\kappa,
 \qquad
 \mathbb E_{\widetilde\kappa}h=0.
 \label{eq:v5v87-density}
\end{equation}
Let \(d\) and \(\mathscr L=d^*d\) be the reference shell gradient and
generator.  For any observable \(F\), put
\begin{equation}
 QF:=F-\mathbb E_{\widetilde\kappa}F.
 \label{eq:v5v87-centering}
\end{equation}

\begin{theorem}[Exact one-observable Stein bound]
\label{thm:v5v87-one-observable}
If \(h\) is in the negative Sobolev space and
\(F\in\operatorname{Dom}(d)\), then
\begin{equation}
 \mathbb E_\kappa F-
 \mathbb E_{\widetilde\kappa}F
 =\mathbb E_{\widetilde\kappa}(hQF)
 \label{eq:v5v87-exact-difference}
\end{equation}
and
\begin{equation}
 \boxed{\qquad
 \left|
 \mathbb E_\kappa F-
 \mathbb E_{\widetilde\kappa}F
 \right|
 \le
 \|h\|_{-1,\widetilde\kappa}
 \|dF\|_{L^2(\widetilde\kappa)}.
 \qquad}
 \label{eq:v5v87-Stein-bound}
\end{equation}
The statement holds coefficientwise for a finite Grassmann polynomial and
in the associated column norm for finite-dimensional operator targets.
\end{theorem}

\begin{proof}
Equation \eqref{eq:v5v87-density} gives
\begin{equation}
 \mathbb E_\kappa F-
 \mathbb E_{\widetilde\kappa}F
 =\mathbb E_{\widetilde\kappa}(hF).
\end{equation}
The constant part of \(F\) vanishes because \(h\) is centered, proving
\eqref{eq:v5v87-exact-difference}.  Lemma
\ref{lem:v5v81-Poisson-duality} bounds the pairing by
\(\|h\|_{-1}\|dQF\|_2\), and \(dQF=dF\).  Apply the scalar proof to each
Grassmann coefficient or each target vector before taking the declared
finite coefficient or column norm.
\end{proof}

\begin{assessment}[The derivative acts after assembly]
\label{ass:v5v87-assembled-derivative}
In \eqref{eq:v5v87-Stein-bound}, \(F\) is the complete differentiated Ward
packet of V86, including both fermion-leg contacts, the gauge-field contact,
self-energy terms, and measure partners.  Replacing \(\|dF\|_2\) by the sum
of derivatives of those pieces is valid but can discard the Ward
cancellation which makes the complete packet small.
\end{assessment}

\subsection{Comparison of matched packet observables}

Let \(F_a\) be the cutoff packet observable and \(F_*\) its matched
reference representative on a common block chart.  Write
\begin{equation}
 \Delta F_a:=F_a-F_*.
 \label{eq:v5v87-packet-defect}
\end{equation}

\begin{proposition}[Density and symbol debits separate exactly]
\label{prop:v5v87-two-debits}
With the exact law \(\kappa_a=(1+h_a)\widetilde\kappa_a\),
\begin{align}
 \mathbb E_{\kappa_a}F_a-
 \mathbb E_{\widetilde\kappa_a}F_*
 &=\mathbb E_{\widetilde\kappa_a}\Delta F_a
 +\mathbb E_{\widetilde\kappa_a}(h_aQF_a),
 \label{eq:v5v87-two-debits}\
 \left|
 \mathbb E_{\kappa_a}F_a-
 \mathbb E_{\widetilde\kappa_a}F_*
 \right|
 &\le
 \|\Delta F_a\|_{L^1(\widetilde\kappa_a)}
 +\|h_a\|_{-1}
  \|dF_a\|_{L^2(\widetilde\kappa_a)}.
 \label{eq:v5v87-two-debit-bound}
\end{align}
\end{proposition}

\begin{proof}
Add and subtract \(\mathbb E_{\widetilde\kappa_a}F_a\), then apply
Theorem~\ref{thm:v5v87-one-observable} to the resulting density-change
term.  The triangle inequality gives the displayed bound.
\end{proof}

\begin{theorem}[Summable quantum fixed-packet criterion]
\label{thm:v5v87-summable-packet}
Let \(a_n=q^{-n}a_0\).  Suppose for the rooted complete packet
\begin{align}
 \|\Delta F_{a_n}\|_{L^1}&\le C_0a_n^\beta,
 \label{eq:v5v87-symbol-gain}\\
 \|h_{a_n}\|_{-1}&\le C_1a_n^\alpha,
 \label{eq:v5v87-density-negative}\\
 \|dF_{a_n}\|_2&\le C_2a_n^{-\gamma},
 \label{eq:v5v87-packet-derivative}
\end{align}
with \(\beta>0\) and \(\alpha>\gamma\ge0\).  Then
\begin{equation}
 \left|
 \mathbb E_{\kappa_{a_n}}F_{a_n}-
 \mathbb E_{\widetilde\kappa_{a_n}}F_*
 \right|
 \le C\left(a_n^\beta+a_n^{\alpha-\gamma}\right),
 \label{eq:v5v87-packet-gain}
\end{equation}
and the one-step packet defects are summable.
\end{theorem}

\begin{proof}
Insert \eqref{eq:v5v87-symbol-gain}--
\eqref{eq:v5v87-packet-derivative} in
\eqref{eq:v5v87-two-debit-bound}.  Both exponents are positive, so both
dyadic series converge.
\end{proof}

For the tree-level packet of V86, \(\beta=1\) without dimension-five
improvement and \(\beta=2\) with it.  The quantum problem is therefore
reduced to obtaining \(\alpha>\gamma\) for the normalized density and the
shell derivative of the complete packet.

\subsection{A direct current version}

Let \(F_a\) take values in spatial currents and let the constraint source
be \(\nabla_a\!\cdot F_a\).  The Stein estimate may be taken in the current
norm before the gravitational inverse is applied.

\begin{corollary}[Quantum constraint-current bound]
\label{cor:v5v87-quantum-current}
Under the hypotheses of Theorem~\ref{thm:v5v87-summable-packet},
\begin{equation}
 \left\|
 \mathbb E_{\kappa_{a_n}}F_{a_n}-
 \mathbb E_{\widetilde\kappa_{a_n}}F_*
 \right\|_{L^2_x}
 \le C(a_n^\beta+a_n^{\alpha-\gamma}).
 \label{eq:v5v87-current-bound}
\end{equation}
Consequently the induced constraint-source difference has the same bound
in \(\dot H^{-1}_x\).
\end{corollary}

\begin{proof}
Apply Theorem~\ref{thm:v5v87-one-observable} after pairing the current with
an arbitrary unit \(L^2_x\) test vector, and take the supremum.  Then use
Theorem~\ref{thm:v5v84-divergence-payer}.
\end{proof}

\subsection{Conditional covariance form}

When the reference shell generator has a spectral gap, the negative norm
can also be represented by an exact integrated covariance.  Let
\(P_t=e^{-t\mathscr L}\) on the mean-zero subspace.

\begin{lemma}[Semigroup representation]
\label{lem:v5v87-semigroup}
For centered \(h\),
\begin{equation}
 \|h\|_{-1}^2
 =\int_0^\infty\langle h,P_th\rangle,dt.
 \label{eq:v5v87-semigroup-negative}
\end{equation}
\end{lemma}

\begin{proof}
The spectral theorem gives
\(\int_0^\infty e^{-t\lambda}dt=\lambda^{-1}\) on the positive spectrum.
Insert this identity in
\(\langle h,\mathscr L^{-1}h\rangle\).
\end{proof}

This formula shows that the required density gain is a time-integrated
connected shell correlation.  It may be attacked by mixing or a Ward
current without proving a pointwise density bound.  It does not itself
create decay.

\subsection{Phase and normalization denominator}

The positive law \(\kappa_a\) in this note contains the gauge--Higgs and
positive regulator modulus.  Let \(\Xi_a\) denote the remaining finite
Grassmann polynomial and determinant phase.  A physical normalized
expectation has the form
\begin{equation}
 \langle F\rangle_a
 =\frac{\mathbb E_{\kappa_a}(\Xi_aF)}
 {\mathbb E_{\kappa_a}\Xi_a}.
 \label{eq:v5v87-complex-normalization}
\end{equation}
Theorem~\ref{thm:v5v87-one-observable} applies coefficientwise to
\(\Xi_aF\), but passage to the ratio additionally requires a uniform lower
bound on the normalization denominator and a shell derivative bound on the
assembled product \(\Xi_aF\).

\begin{firewall}[The sign denominator is an independent gate]
\label{fw:v5v87-denominator}
Reflection positivity of the supertransfer and local control of the phase
do not imply
\(|\mathbb E_{\kappa_a}\Xi_a|\ge c>0\).  Severe phase cancellation can make
the denominator small.  No normalized physical expectation may be bounded
from numerator estimates until this denominator or an equivalent
phase-free representation is controlled.
\end{firewall}

\begin{firewall}[The density gain remains conditional]
\label{fw:v5v87-density-open}
V81 identifies sufficient Poincar\'e and residual-gradient inputs for
\eqref{eq:v5v87-density-negative}, but those inputs have not been acquired
for the full joint shell.  V87 reduces the fixed-packet quantum estimate to
that negative norm; it does not prove its positive cutoff exponent.
\end{firewall}

\begin{firewall}[Rooted derivatives may grow]
\label{fw:v5v87-derivative-growth}
Even at fixed external arity, differentiating the packet in an internal
shell direction can cost inverse powers of the lattice spacing.  The
condition \(\alpha>\gamma\) records the exact required surplus.  Assuming
\(\gamma=0\) without a canonical rescaling and derivative count would hide
the principal ultraviolet debit.
\end{firewall}

\begin{theorem}[V87 conclusion]
\label{thm:v5v87-status}
A normalized shell density defect acts on one assembled Ward packet through
an exact Stein pairing.  Its quantum correction is bounded directly by the
negative density norm times one shell derivative of the packet, so the
positive Fisher derivative needed for a uniform transfer class is not
required at fixed arity.  A tree/Symanzik gain \(a^\beta\), a density gain
\(a^\alpha\), and packet derivative cost \(a^{-\gamma}\) give a summable
quantum constraint current whenever \(\beta>0\) and \(\alpha>\gamma\).
The density exponent, complex normalization denominator, and derivative
cost remain open for the complete theory.
\end{theorem}

\begin{proof}
Use Theorem~\ref{thm:v5v87-one-observable},
Proposition~\ref{prop:v5v87-two-debits},
Theorem~\ref{thm:v5v87-summable-packet}, and
Corollary~\ref{cor:v5v87-quantum-current}, subject to the three firewalls.
\end{proof}

V88 will analyze the complex normalization gate in
\eqref{eq:v5v87-complex-normalization}.  It will prove a quantitative
denominator-stability theorem from a rooted phase variance and Grassmann
constant-term estimate, and determine whether local phase control can be
assembled without assuming a volume-uniform global sign bound.


\clearpage
\section{V88: rooted complex normalization and vacuum-phase cancellation}
\label{sec:v5v88}

\subsection{Why a global sign bound is too strong}

V87 wrote a physical expectation as
\begin{equation}
 \langle F\rangle_Xi
 =\frac{\mathbb E_\kappa(\Xi F)}{\mathbb E_\kappa\Xi},
 \label{eq:v5v88-ratio}
\end{equation}
where \(\kappa\) is the positive assembled shell law and \(\Xi\) contains
the chiral phase and finite Grassmann coefficients.  Requiring the global
denominator in \eqref{eq:v5v88-ratio} to stay uniformly away from zero is
unnecessarily strong for local observables.  Independent or sufficiently
clustered vacuum factors cancel between numerator and denominator before a
norm is taken.

\begin{lemma}[Exact cancellation of an independent vacuum factor]
\label{lem:v5v88-independent-cancellation}
Suppose
\(\kappa=\kappa_A\otimes\kappa_B\), \(F=F_A\), and
\(\Xi=\Xi_A\Xi_B\), with both scalar denominators nonzero.  Then
\begin{equation}
 \frac{\mathbb E_\kappa(\Xi F)}{\mathbb E_\kappa\Xi}
 =\frac{\mathbb E_{\kappa_A}(\Xi_AF_A)}
 {\mathbb E_{\kappa_A}\Xi_A}.
 \label{eq:v5v88-independent-cancellation}
\end{equation}
The size of \(|\mathbb E_{\kappa_B}\Xi_B|\) does not enter the local ratio.
\end{lemma}

\begin{proof}
Both numerator and denominator factor.  Cancel the common factor
\(\mathbb E_{\kappa_B}\Xi_B\).
\end{proof}

This elementary identity is the exact model for a rooted polymer
normalization.  Replacing it by separate absolute bounds on the two global
integrals loses the cancellation.

\subsection{Connected normalization identity}

Let \(V\) be a complex random variable for which
\(\mathbb E e^{V+tF}\) is nonzero and analytic near
\((V,t)=(0,0)\) along the chosen interpolation.  Denote joint cumulants by
\(\kappa_m\).

\begin{theorem}[Local ratios are connected expansions]
\label{thm:v5v88-connected-ratio}
Whenever the following series is absolutely convergent,
\begin{equation}
 \boxed{\qquad
 \frac{\mathbb E(Fe^V)}{\mathbb E e^V}
 =\sum_{m\ge0}\frac1{m!}
 \kappa_{m+1}(F,\underbrace{V,\ldots,V}_{m}).
 \qquad}
 \label{eq:v5v88-connected-ratio}
\end{equation}
If \(V=V_A+V_B\), \(F\) is independent of \(V_B\), and the law
factorizes between \(A\) and \(B\), every cumulant containing a \(V_B\)
vanishes.  Thus vacuum components disconnected from \(F\) are absent term
by term.
\end{theorem}

\begin{proof}
Introduce
\begin{equation}
 K(s,t)=\log\mathbb E e^{sV+tF}.
\end{equation}
Then
\begin{equation}
 \left.\partial_tK(1,t)\right|_{t=0}
 =\frac{\mathbb E(Fe^V)}{\mathbb E e^V}.
\end{equation}
The Taylor expansion of \(K\) at the origin has joint cumulants as its
coefficients.  Differentiate it in \(t\), set \(s=1\), and use absolute
convergence to obtain \eqref{eq:v5v88-connected-ratio}.  Mixed cumulants of
independent families vanish, proving the last assertion.
\end{proof}

For a polymer decomposition \(V=\sum_XV_X\), every nonzero cumulant on the
right of \eqref{eq:v5v88-connected-ratio} has a dependency graph connecting
the support of \(F\) to the participating polymers.  This is the rigorous
location of the root.  Proving an all-order bound on those connected terms
is still required.

\subsection{A quantitative rooted phase denominator}

Let \(\Theta_A\) be a real phase supported in the rooted dependency
neighbourhood of a fixed cylinder \(A\), and let \(G_A\) be the remaining
Grassmann or complex correction after its scalar body has been normalized
to one.  Write
\begin{equation}
 \Xi_A=e^{i\Theta_A}(1+G_A).
 \label{eq:v5v88-local-weight}
\end{equation}
The expectation below means the scalar body after the finite Berezin
contractions have been performed, while \(\|G_A\|_{\rm gr}\) denotes the
V58 coefficient norm.

\begin{theorem}[Rooted denominator lower bound]
\label{thm:v5v88-rooted-denominator}
Suppose
\begin{equation}
 \operatorname{Var}_\kappa(\Theta_A)\le\sigma_A^2,
 \qquad
 \mathbb E_\kappa\|G_A\|_{\rm gr}\le\varepsilon_A.
 \label{eq:v5v88-local-smallness}
\end{equation}
Then, with \(\bar\Theta_A=\mathbb E_\kappa\Theta_A\),
\begin{equation}
 \left|mathbb E_\kappa\Xi_A\right|
 \ge1-\frac12\sigma_A^2-\varepsilon_A.
 \label{eq:v5v88-denominator-lower}
\end{equation}
In particular, the rooted normalization is invertible if
\begin{equation}
 \frac12\sigma_A^2+\varepsilon_A<1.
 \label{eq:v5v88-local-invertibility}
\end{equation}
The deterministic mean phase changes only the argument of the denominator
and costs nothing in the bound.
\end{theorem}

\begin{proof}
Multiply the expectation by \(e^{-i\bar\Theta_A}\), which preserves its
modulus.  Since \(\cos x\ge1-x^2/2\),
\begin{equation}
 \operatorname{Re}\mathbb E_
 \kappa e^{i(\Theta_A-\bar\Theta_A)}
 \ge1-\frac12\operatorname{Var}_\kappa(\Theta_A).
\end{equation}
The correction has modulus at most
\begin{equation}
 \left|\mathbb E_\kappa
 e^{i(\Theta_A-\bar\Theta_A)}G_A\right|
 \le\mathbb E_\kappa\|G_A\|_{\rm gr}.
\end{equation}
Subtract this bound from the real part and use
\(|z|\ge\operatorname{Re}z\).
\end{proof}

\begin{corollary}[Stability of a rooted normalized expectation]
\label{cor:v5v88-ratio-stability}
Let \((N,Z)\) and \((\widetilde N,\widetilde Z)\) be two local numerator and
denominator pairs with
\(|Z|,|\widetilde Z|\ge c_A>0\).  Then
\begin{equation}
 \left|\frac NZ-\frac{\widetilde N}{\widetilde Z}\right|
 \le\frac{|N-\widetilde N|}{c_A}
 +\frac{|\widetilde N|}{c_A^2}|Z-\widetilde Z|.
 \label{eq:v5v88-ratio-stability}
\end{equation}
\end{corollary}

\begin{proof}
Write
\(N/Z-\widetilde N/\widetilde Z=(N-\widetilde N)/Z+
\widetilde N(\widetilde Z-Z)/(Z\widetilde Z)\) and take absolute values.
\end{proof}

\subsection{Variance from rooted phase correlations}

Suppose
\begin{equation}
 \Theta_A=\sum_{x\in\mathcal N(A)}\theta_x
 \label{eq:v5v88-phase-sum}
\end{equation}
with a finite or summably weighted rooted neighbourhood.  Center each
\(\theta_x\).

\begin{proposition}[Covariance-sum phase bound]
\label{prop:v5v88-phase-variance}
If
\begin{equation}
 \sup_x\sum_y
 \left|\operatorname{Cov}_\kappa(\theta_x,\theta_y)\right|
 \le C_\theta a^{2\alpha},
 \qquad \alpha>0,
 \label{eq:v5v88-covariance-sum}
\end{equation}
then for a finite rooted neighbourhood
\begin{equation}
 \operatorname{Var}_\kappa(\Theta_A)
 \le C_\theta|\mathcal N(A)|a^{2\alpha}.
 \label{eq:v5v88-phase-variance}
\end{equation}
For fixed \(A\), the phase part of
\eqref{eq:v5v88-local-invertibility} therefore holds at sufficiently fine
cutoff.
\end{proposition}

\begin{proof}
Expand the variance as the double covariance sum.  Sum first in \(y\) and
use \eqref{eq:v5v88-covariance-sum} for every
\(x\in\mathcal N(A)\).
\end{proof}

The V78 first-jet match makes each residual high-Dirac phase density
quadratic in a smooth background amplitude.  Proposition
\ref{prop:v5v88-phase-variance} converts that local amplitude gain into a
denominator bound only if the rooted covariance sum is volume-uniform.
Pointwise small phase coefficients alone do not control their correlations.

\subsection{Grassmann body and nilpotent inversion}

Let \(\mathfrak G_A\) be the finite Grassmann algebra of a rooted cylinder,
with body map \(\mathrm b\) and nilpotent ideal \(\mathfrak n\).

\begin{lemma}[Finite rooted Grassmann inverse]
\label{lem:v5v88-Grassmann-inverse}
An element \(Z=z_0+N\), \(z_0=\mathrm b(Z)\ne0\),
\(N\in\mathfrak n\), is invertible and
\begin{equation}
 Z^{-1}=z_0^{-1}
 \sum_{j=0}^{J_A}(-z_0^{-1}N)^j,
 \label{eq:v5v88-Grassmann-inverse}
\end{equation}
where \(J_A\) is at most the number of Grassmann generators.  If the V58
coefficient norm is submultiplicative and
\(\|N\|_{\rm gr}<|z_0|\), then
\begin{equation}
 \|Z^{-1}\|_{\rm gr}
 \le\frac1{|z_0|-\|N\|_{\rm gr}}.
 \label{eq:v5v88-Grassmann-inverse-bound}
\end{equation}
\end{lemma}

\begin{proof}
Nilpotence makes the geometric series terminate and direct multiplication
gives one.  The norm bound follows by comparison with the infinite scalar
geometric series.
\end{proof}

Thus the true denominator gate is its scalar rooted body, followed by a
finite nilpotent inverse.  A global Grassmann generator count would make
\(J_A\) extensive; the rooted algebra keeps it fixed for each cylinder.

\subsection{A local normalization certificate}

\begin{definition}[Rooted complex normalization certificate]
\label{def:v5v88-RCNC}
An RCNC for a fixed cylinder support \(A\) consists of:
\begin{enumerate}
 \item a connected decomposition which cancels every vacuum component
 disjoint from \(A\), as in
 Theorem~\ref{thm:v5v88-connected-ratio};
 \item a rooted phase covariance sum
 \eqref{eq:v5v88-covariance-sum};
 \item a rooted Grassmann correction with
 \(\mathbb E\|G_A\|_{\rm gr}\le\varepsilon_A\) and
 \(\sigma_A^2/2+\varepsilon_A<1\);
 \item summable adjacent-cutoff differences of the rooted numerators and
 denominators;
 \item compatibility with block pullback, reflection, and determinant-line
 chart transitions.
\end{enumerate}
\end{definition}

\begin{theorem}[RCNC removes the global denominator hypothesis]
\label{thm:v5v88-RCNC}
For a fixed cylinder, an RCNC and the V87 numerator estimates give a
well-defined normalized local expectation with summable adjacent-cutoff
differences.  No volume-uniform lower bound on the complete vacuum partition
function is required.
\end{theorem}

\begin{proof}
The connected expansion cancels vacuum components outside the rooted
dependency cluster.  Theorem~\ref{thm:v5v88-rooted-denominator} and
Lemma~\ref{lem:v5v88-Grassmann-inverse} invert the remaining rooted
denominator.  Corollary~\ref{cor:v5v88-ratio-stability} converts summable
numerator and denominator differences into a summable ratio difference.
\end{proof}

\begin{firewall}[Connected expansion is still an all-order gate]
\label{fw:v5v88-connected-open}
Equation \eqref{eq:v5v88-connected-ratio} is exact under absolute
convergence.  V88 does not prove the required rooted all-order cumulant
bound for the interacting SM--Einstein shell.  Without it, one cannot
delete distant factors merely by drawing a connected graph.
\end{firewall}

\begin{firewall}[Local denominators must agree on overlaps]
\label{fw:v5v88-overlap-phase}
Two rooted determinant-line charts can each satisfy
\eqref{eq:v5v88-local-invertibility} while their phase choices differ by a
nontrivial transition cocycle.  The RCNC does not replace the global
anomaly and holonomy gates of V60--V61 and V78.
\end{firewall}

\begin{theorem}[V88 conclusion]
\label{thm:v5v88-status}
A global sign-denominator bound is unnecessary for a local continuum
ratio: independent vacuum factors cancel exactly, and an absolutely
convergent cumulant expansion contains only components connected to the
observable root.  A rooted phase variance and Grassmann body estimate give
an explicit local denominator lower bound and stable finite nilpotent
inverse.  The remaining obstruction is the all-order connected bound and
compatibility of local phase choices, not the possible exponential
smallness of a distant vacuum factor.
\end{theorem}

\begin{proof}
Use Lemma~\ref{lem:v5v88-independent-cancellation},
Theorem~\ref{thm:v5v88-connected-ratio},
Theorem~\ref{thm:v5v88-rooted-denominator},
Proposition~\ref{prop:v5v88-phase-variance},
Lemma~\ref{lem:v5v88-Grassmann-inverse}, and
Theorem~\ref{thm:v5v88-RCNC}, subject to the two firewalls.
\end{proof}

V89 will derive the rooted phase covariance sum for the high Dirac
trace-loop density.  It will combine the V78 local Lipschitz estimate with a
massive conditional covariance inequality and separate the massless
constraint contribution, whose Coulomb tail cannot satisfy an exponential
cluster sum.


\clearpage
\section{V89: rooted phase covariance in massive and Coulomb grades}
\label{sec:v5v89}

\subsection{Differentiating the rooted trace-loop density}

Let \(R\) be the relative high Dirac defect of V78 with
\(\|R\|_{\mu,1}\le\kappa<1\).  Its rooted complex density is
\begin{equation}
 \lambda_x(R)=
 \sum_{m\ge1}\frac{(-1)^{m+1}}m
 \operatorname{tr}_{\rm int}(R^m)_{xx},
 \qquad
 \theta_x:=\operatorname{Im}\lambda_x(R).
 \label{eq:v5v89-rooted-phase}
\end{equation}
Let \(d_z\) denote differentiation in a shell coordinate rooted near
\(z\).  Assume the derivative insertion has the localized bound
\begin{equation}
 \left\|d_zR\right\|_{\mu,1;x}
 \le L_Ra^\alpha e^{-\nu d(x,z)},
 \qquad \alpha>0,
 \label{eq:v5v89-local-derivative-R}
\end{equation}
where the left side denotes the weighted sum of kernels in which the closed
loop is rooted at \(x\) and contains the displayed derivative insertion.

\begin{lemma}[Rooted phase-gradient bound]
\label{lem:v5v89-phase-gradient}
Under \eqref{eq:v5v89-local-derivative-R},
\begin{equation}
 |d_z\theta_x|
 \le\frac{d_*L_R}{1-\kappa}
 a^\alpha e^{-\nu d(x,z)}.
 \label{eq:v5v89-phase-gradient}
\end{equation}
The same estimate holds for the real rooted modulus density.
\end{lemma}

\begin{proof}
Differentiate the series in \eqref{eq:v5v89-rooted-phase}.  The derivative
of \(R^m\) is the sum of \(m\) words with one \(d_zR\) insertion.  The
factor \(m\) cancels the coefficient \(1/m\).  Weighted multiplication
bounds the remaining \(m-1\) factors by \(\kappa^{m-1}\), while
\eqref{eq:v5v89-local-derivative-R} supplies the rooted insertion.  Sum
\(\sum_{m\ge1}\kappa^{m-1}=(1-\kappa)^{-1}\) and use the internal trace
bound.  Taking real or imaginary parts cannot increase the estimate.
\end{proof}

Thus the quadratic first-jet residual of V78 yields a phase-gradient gain
whenever the background amplitude and its shell derivative have a positive
cutoff power.  The next step turns this gradient locality into a covariance
sum.

\subsection{The one-form covariance resolvent}

Let \(\kappa\) be a positive conditional shell law with scalar generator
\(\mathscr L=d^*d\).  On shell one-forms, let \(\mathscr L_1\) be the
Helffer--Sj\"ostrand operator obtained by commuting \(d\) through
\(\mathscr L\).  Assume it is invertible on the relevant one-form sector.

\begin{lemma}[Exact one-form covariance identity]
\label{lem:v5v89-HS-identity}
For centered smooth scalar observables \(F,G\),
\begin{equation}
 \operatorname{Cov}_\kappa(F,G)
 =\mathbb E_\kappa
 \left\langle dF,\mathscr L_1^{-1}dG\right\rangle.
 \label{eq:v5v89-HS-identity}
\end{equation}
\end{lemma}

\begin{proof}
Let \(u=\mathscr L^{-1}G\).  Integration by parts gives
\(\operatorname{Cov}(F,G)=\langle dF,du\rangle\).  The commutation identity
\(d\mathscr L=\mathscr L_1d\) gives
\(du=\mathscr L_1^{-1}dG\).  Combine the two formulas.
\end{proof}

Assume the one-form inverse has an annealed or pointwise kernel majorant
\begin{equation}
 \left\|mathscr L_1^{-1}(z,w)\right\|
 \le K_1e^{-m_1d(z,w)}
 \label{eq:v5v89-one-form-decay}
\end{equation}
with constants uniform in volume and retained background.

\begin{theorem}[Massive rooted phase covariance]
\label{thm:v5v89-massive-phase-covariance}
Suppose \eqref{eq:v5v89-phase-gradient} and
\eqref{eq:v5v89-one-form-decay} hold, and the shell graph has polynomial
growth.  For every
\begin{equation}
 0<m'<\min\{\nu,m_1\},
 \label{eq:v5v89-decay-exponent}
\end{equation}
there is \(C<\infty\), independent of volume, such that
\begin{equation}
 \left|\operatorname{Cov}_\kappa(\theta_x,\theta_y)\right|
 \le Ca^{2\alpha}e^{-m'd(x,y)},
 \label{eq:v5v89-massive-covariance}
\end{equation}
and
\begin{equation}
 \sup_x\sum_y
 \left|\operatorname{Cov}_\kappa(\theta_x,\theta_y)\right|
 \le C'a^{2\alpha}.
 \label{eq:v5v89-massive-covariance-sum}
\end{equation}
\end{theorem}

\begin{proof}
Insert the two phase gradients in
\eqref{eq:v5v89-HS-identity} and use
\eqref{eq:v5v89-one-form-decay}.  The resulting triple convolution is
bounded by
\begin{equation}
 Ca^{2\alpha}\sum_{z,w}
 e^{-\nu d(x,z)}e^{-m_1d(z,w)}e^{-\nu d(w,y)}.
\end{equation}
Spend the difference between the exponents and \(m'\) to sum polynomial
graph growth, leaving \eqref{eq:v5v89-massive-covariance}.  Summation in
\(y\) gives \eqref{eq:v5v89-massive-covariance-sum}.
\end{proof}

\begin{corollary}[Massive phase denominator]
\label{cor:v5v89-massive-denominator}
For a fixed rooted neighbourhood \(\mathcal N(A)\), the massive phase
\(\Theta_A^{\rm M}=\sum_{x\in\mathcal N(A)}\theta_x\) satisfies
\begin{equation}
 \operatorname{Var}(\Theta_A^{\rm M})
 \le C|\mathcal N(A)|a^{2\alpha}.
 \label{eq:v5v89-massive-phase-variance}
\end{equation}
It therefore meets the phase part of the V88 rooted denominator criterion
at sufficiently fine cutoff.
\end{corollary}

\begin{proof}
Apply Proposition~\ref{prop:v5v88-phase-variance} to
\eqref{eq:v5v89-massive-covariance-sum}.
\end{proof}

\subsection{Constraint substitution changes the decay class}

Suppose the local phase density depends on the solved conformal constraint
through a source \(f\).  Its response to a remote source perturbation is
then composed with the spatial Green operator of V83.  If the source defect
has moments below order \(m\) matched, V84 gives the gradient-response
majorant
\begin{equation}
 L_m(x,z)
 \le C a^\alpha(1+d(x,z))^{-m-2}.
 \label{eq:v5v89-Coulomb-response}
\end{equation}

\begin{lemma}[The three-dimensional summability threshold]
\label{lem:v5v89-summability-threshold}
On a graph with three-dimensional polynomial volume growth,
\begin{equation}
 \sup_x\sum_z(1+d(x,z))^{-m-2}<\infty
 \quad\Longleftrightarrow\quad m>1.
 \label{eq:v5v89-l1-threshold}
\end{equation}
Thus monopole cancellation alone, \(m=1\), gives the logarithmically
non-summable borderline \(r^{-3}\), while monopole and dipole cancellation,
\(m=2\), give the summable tail \(r^{-4}\).
\end{lemma}

\begin{proof}
The number of vertices in a shell of radius \(r\) is comparable to
\(r^2\).  The row sum has the same convergence as
\(\sum_{r\ge1}r^2r^{-m-2}=\sum_{r\ge1}r^{-m}\), which converges exactly
when \(m>1\).
\end{proof}

\begin{theorem}[Coulomb-grade phase covariance]
\label{thm:v5v89-Coulomb-phase}
Assume the phase gradients satisfy the response bound
\eqref{eq:v5v89-Coulomb-response} with \(m\ge2\), and the remaining
conditional one-form covariance kernel has a volume-uniform \(\ell^1\)
row and column bound \(K_C\).  Then
\begin{equation}
 \sup_x\sum_y
 |\operatorname{Cov}(\theta_x^{\rm C},\theta_y^{\rm C})|
 \le C_mK_Ca^{2\alpha}.
 \label{eq:v5v89-Coulomb-covariance-sum}
\end{equation}
If only \(m=1\) is known, this absolute-sum conclusion does not follow from
the power-counting bound.
\end{theorem}

\begin{proof}
The covariance identity is a product of the response kernel on the left,
the one-form covariance kernel, and the response kernel on the right.  The
\(\ell^1\) operator norm of the product is at most the product of their
three \(\ell^1\) norms.  Lemma~\ref{lem:v5v89-summability-threshold} makes
each response norm at most \(C_ma^\alpha\), proving
\eqref{eq:v5v89-Coulomb-covariance-sum}.  At \(m=1\), that response norm has
a logarithmic volume divergence, so the same argument is unavailable.
\end{proof}

\begin{assessment}[The dipole has a second role]
\label{ass:v5v89-dipole-role}
V85 used dipole matching to write the free adjacent-cutoff stress defect as
a double divergence.  V89 shows an independent necessity: in three spatial
dimensions, the resulting \(r^{-4}\) response is the first Coulomb
multipole tail whose absolute row sum is finite.  Monopole matching alone
controls the \(\dot H^{-1}\) energy but does not give the rooted phase
covariance sum required by V88.
\end{assessment}

\subsection{Mixed massive and constraint phase}

Write \(\theta_x=\theta_x^{\rm M}+\theta_x^{\rm C}\).  Cross covariance is
bounded by the same one-form representation with one exponential and one
Coulomb response kernel.

\begin{corollary}[Mixed rooted phase variance]
\label{cor:v5v89-mixed-phase}
Under the hypotheses of Theorems
\ref{thm:v5v89-massive-phase-covariance} and
\ref{thm:v5v89-Coulomb-phase} with \(m\ge2\),
\begin{equation}
 \sup_x\sum_y|\operatorname{Cov}(\theta_x,\theta_y)|
 \le Ca^{2\alpha}.
 \label{eq:v5v89-mixed-phase-sum}
\end{equation}
Hence every fixed rooted phase satisfies the V88 variance criterion.
\end{corollary}

\begin{proof}
The massive--massive and Coulomb--Coulomb sums are the two theorems.  For a
cross term, use the \(\ell^1\) norm of one exponential response, the
assumed one-form covariance bound, and the \(\ell^1\) norm of the
\(m\ge2\) Coulomb response.  Add the four covariance components and apply
Proposition~\ref{prop:v5v88-phase-variance}.
\end{proof}

\subsection{Unproved inputs}

\begin{firewall}[A scalar Poincar\'e gap is not one-form clustering]
\label{fw:v5v89-one-form-open}
The Poincar\'e inequality of V81 bounds the scalar inverse globally.  It does
not prove the spatial kernel estimate
\eqref{eq:v5v89-one-form-decay}.  Strict local convexity or a proved
massive cluster expansion can supply such a one-form resolvent; gauge zero
modes and critical Higgs directions can obstruct it.
\end{firewall}

\begin{firewall}[First-jet amplitude does not automatically control its
shell derivative]
\label{fw:v5v89-phase-derivative-open}
V78 makes the residual phase density quadratic in the background amplitude
on a smooth chart.  Equation \eqref{eq:v5v89-local-derivative-R} also
requires a cutoff-uniform derivative insertion bound.  A shell derivative
can cost an inverse lattice spacing, so its exponent \(\alpha\) must be
computed after canonical rescaling.
\end{firewall}

\begin{firewall}[Dipole matching remains a normalization gate]
\label{fw:v5v89-dipole-open}
The Coulomb phase theorem needs the \(m\ge2\) response.  V85 proves that
free matched stencils have the required double difference on a smooth band;
it does not prove dipole matching for the full interacting stress source or
for cutoff-scale shell modes.
\end{firewall}

\begin{theorem}[V89 conclusion]
\label{thm:v5v89-status}
The rooted high-Dirac trace-loop phase inherits an exponentially local
gradient with the exact geometric factor \((1-\kappa)^{-1}\).  A local
one-form covariance resolvent converts that gradient into the V88
volume-uniform phase covariance sum for massive sectors.  After nonlinear
constraint substitution, the same conclusion requires both monopole and
dipole matching: the \(r^{-4}\) response is summable in three dimensions,
whereas the monopole-only \(r^{-3}\) response is logarithmically
borderline.  The one-form clustering, derivative exponent, and interacting
dipole match remain open.
\end{theorem}

\begin{proof}
Use Lemma~\ref{lem:v5v89-phase-gradient},
Theorem~\ref{thm:v5v89-massive-phase-covariance},
Lemma~\ref{lem:v5v89-summability-threshold},
Theorem~\ref{thm:v5v89-Coulomb-phase}, and
Corollary~\ref{cor:v5v89-mixed-phase}, subject to the three firewalls.
\end{proof}

V90 is the next compulsory YM/NS audit.  It will then attack the one-form
clustering input on the strictly massive-safe gauge--Higgs--PV patch,
using a block Hessian/Schur criterion to separate genuine massive
directions from gauge, Goldstone, and gravitational constraint modes.

\clearpage
\section{V90: a massive-safe one-form resolvent}
\label{sec:v5v90}

\subsection{Compulsory companion audit}

Before changing the SM direction, I inspected the newest companion files in
the shared folder.  The Yang--Mills file was
\texttt{Renewal\_Yang\_Mills\_Mass\_Gap\_Research\_Notes\_V96.tex}, with
SHA--256
\begin{center}
\texttt{D8F972D600529AAC4BF6B8B7206AD2D8EE002A4AFC650252A056BE27722465C9}.
\end{center}
Its new mathematical endpoint is a dimension-free, fixed-
\(p\) gradient--resolvent estimate obtained from a global lower curvature
bound and reverse semigroup Poincar\'e estimate.  It completes the stated
three-label conditional card there.  It does not establish spatial
exponential decay of the resolvent kernel, the four-label/all-order card,
reflection positivity, cutoff removal, or a Yang--Mills mass gap.  The
Navier--Stokes file remained
\texttt{Renewal\_NS\_Stability\_Research\_Notes\_V31.tex}, with SHA--256
\begin{center}
\texttt{F1121B62238AD5491BB7CF03635EA9934942EDF573ED8FAAA9EE79E1D38A6696}.
\end{center}
Its relevant lesson remains the separation of the root mode and the first
moment before using a coercive inverse.

The audit changes the next step only in one precise way.  A scalar or
fixed-\(p\) semigroup gap controls the size of a resolvent, but V89 needs a
spatial row-summable kernel.  Locality and a one-form gap must therefore be
proved together.  This note supplies that implication on a strictly
massive block and identifies the Standard Model directions for which its
hypotheses can hold.

\subsection{The block form and the correct quotient}

Work first at finite lattice volume.  Let the real one-form tangent space
of the conditioned bosonic variables split orthogonally as
\begin{equation}
 \mathcal K=\mathcal K_M\oplus\mathcal K_A.
 \label{eq:v5v90-tangent-split}
\end{equation}
The first summand contains only retained massive physical directions.  The
second contains auxiliary gauge-fixed or otherwise already-coercive
directions.  No photon, gauge zero mode, Goldstone mode before gauge
fixing, or gravitational constraint mode is admitted to either summand.
Those directions will be returned to below.

Let the one-form Witten operator, or the corresponding positive covariance
operator in a real Gibbs chart, have quadratic form
\begin{equation}
 \mathscr L_1=
 \begin{pmatrix}A&B\\ B^*&C\end{pmatrix}
 \quad\hbox{on }\mathcal K_M\oplus\mathcal K_A,
 \label{eq:v5v90-block-operator}
\end{equation}
where \(A\) and \(C\) are positive definite.  Define the normalized mixing
\begin{equation}
 T=A^{-1/2}BC^{-1/2}.
 \label{eq:v5v90-normalized-mixing}
\end{equation}

\begin{lemma}[Normalized Schur coercivity]
\label{lem:v5v90-Schur-coercivity}
If \(\|T\|\le q<1\), then
\begin{equation}
 \mathscr L_1\ge (1-q)(A\oplus C).
 \label{eq:v5v90-block-lower}
\end{equation}
In particular, if \(A\ge g_M I\) and \(C\ge g_A I\), then
\begin{equation}
 \mathscr L_1\ge g I,
 \qquad g=(1-q)\min\{g_M,g_A\}>0.
 \label{eq:v5v90-full-gap}
\end{equation}
The retained Schur complement satisfies
\begin{equation}
 S_M:=A-BC^{-1}B^*ge(1-q^2)A.
 \label{eq:v5v90-Schur-lower}
\end{equation}
\end{lemma}

\begin{proof}
For \(u\in\mathcal K_M\), \(v\in\mathcal K_A\), put
\(x=A^{1/2}u\) and \(y=C^{1/2}v\).  Then
\[
 \langle(u,v),\mathscr L_1(u,v)\rangle
 =\|x\|^2+\|y\|^2+2\operatorname{Re}\langle x,Ty\rangle.
\]
The last term is bounded below by
\(-2q\|x\|\|y\|\ge-q(\|x\|^2+\|y\|^2)\), proving
\eqref{eq:v5v90-block-lower} and then \eqref{eq:v5v90-full-gap}.
Moreover
\[
 BC^{-1}B^*=A^{1/2}TT^*A^{1/2}\le q^2A,
\]
which proves \eqref{eq:v5v90-Schur-lower}.
\end{proof}

\begin{firewall}[A quotient is logically prior to the Schur complement]
\label{fw:v5v90-quotient-first}
If \(C\) has a gauge, Goldstone, photon, or constraint zero mode, then
\(C^{-1}\) in \eqref{eq:v5v90-Schur-lower} is undefined.  Inserting a
pseudoinverse silently changes both the kernel and the compatibility
conditions.  One must first quotient exact gauge directions, fix and pin
the intended auxiliary directions, and leave every physical massless
direction in a separate polynomial grade.
\end{firewall}

\subsection{Gap plus locality gives a local inverse}

Let \(X\) be the set of lattice blocks, with graph distance \(d\), and let
\(P_x\) project onto the finite-dimensional fibre over \(x\).  For an
operator \(K\), write \(K_{xy}=P_xKP_y\).  The following elementary
Combes--Thomas estimate is the missing implication in the massive part of
V89.

\begin{theorem}[Finite-range one-form inverse]
\label{thm:v5v90-CT}
Suppose \(L=L^*\ge gI\), \(g>0\), and
\begin{equation}
 L_{xy}=0\quad\hbox{if }d(x,y)>R,
 \qquad
 J:=\max\left\{
 \sup_x\sum_y\|L_{xy}\|,
 \sup_y\sum_x\|L_{xy}\|
 \right\}<\infty .
 \label{eq:v5v90-range-row}
\end{equation}
If \(\mu>0\) obeys
\begin{equation}
 J\bigl(e^{\mu R}-1\bigr)\le \frac g2,
 \label{eq:v5v90-mu-choice}
\end{equation}
then, uniformly in the finite volume,
\begin{equation}
 \|(L^{-1})_{xy}\|\le \frac2g e^{-\mu d(x,y)}.
 \label{eq:v5v90-inverse-pointwise}
\end{equation}
Consequently, on graphs of uniformly bounded three-dimensional polynomial
growth,
\begin{equation}
 \sup_x\sum_y\|(L^{-1})_{xy}\|
 +\sup_y\sum_x\|(L^{-1})_{xy}\|
 \le \frac{C_{\mu}}g,
 \label{eq:v5v90-inverse-l1}
\end{equation}
where \(C_\mu\) is independent of the volume.
\end{theorem}

\begin{proof}
Fix \(y\) and let \(\rho_y(x)=d(x,y)\).  Multiplication by
\(W_y(x)=e^{\mu\rho_y(x)}\) is bounded in finite volume.  If
\(L_{xz}\ne0\), then \(|\rho_y(x)-\rho_y(z)|\le R\), and hence the two
Schur row and column bounds give
\begin{equation}
 \|W_yLW_y^{-1}-L\|
 \le J(e^{\mu R}-1)\le g/2.
 \label{eq:v5v90-conjugation-error}
\end{equation}
Indeed, the kernel of the difference is
\((e^{\mu(\rho_y(x)-\rho_y(z))}-1)L_{xz}\), and the operator norm is at
most the geometric mean of its maximal row and column sums.  Since
\(\|L^{-1}\|\le g^{-1}\), the resolvent identity and a Neumann series imply
\[
 \|(W_yLW_y^{-1})^{-1}\|\le 2g^{-1}.
\]
But this inverse is \(W_yL^{-1}W_y^{-1}\).  Applying its \((x,y)\) block
to a unit vector in the fibre at \(y\), where \(W_y=1\), proves
\eqref{eq:v5v90-inverse-pointwise}.  Summing the exponential over shells
proves \eqref{eq:v5v90-inverse-l1}; the same argument with rows and columns
interchanged gives the second sum.
\end{proof}

The finite-range assumption is convenient rather than essential.  The
weighted norm below is the form actually needed for exponentially local
fermion-induced terms.

\begin{corollary}[Exponentially local perturbations]
\label{cor:v5v90-exp-local}
Let \(L=L_0+E\), where \(L=L^*\ge gI\), and assume that for some
\(\mu>0\)
\begin{equation}
 \max\left\{
 \sup_x\sum_z(e^{\mu d(x,z)}-1)\|L_{xz}\|,
 \sup_z\sum_x(e^{\mu d(x,z)}-1)\|L_{xz}\|
 \right\}\le g/2.
 \label{eq:v5v90-weighted-commutator}
\end{equation}
Then \eqref{eq:v5v90-inverse-pointwise} and
\eqref{eq:v5v90-inverse-l1} hold with this \(\mu\).  In particular, an
exponentially local perturbation preserves the conclusion whenever its
weighted row and column norms are small enough that the positive gap and
\eqref{eq:v5v90-weighted-commutator} remain valid.
\end{corollary}

\begin{proof}
For the distance weight used in the proof of
Theorem~\ref{thm:v5v90-CT}, the triangle inequality gives
\[
 |e^{\mu(\rho_y(x)-\rho_y(z))}-1|
 \le e^{\mu d(x,z)}-1.
\]
Thus \eqref{eq:v5v90-weighted-commutator} replaces
\eqref{eq:v5v90-conjugation-error}; the rest of the proof is unchanged.
\end{proof}

\begin{proposition}[Local block elimination]
\label{prop:v5v90-local-Schur}
Assume \(A,B,C\) have finite range and volume-uniform row and column
bounds, \(C\ge g_AI\), and the hypotheses of
Theorem~\ref{thm:v5v90-CT} hold for \(C\).  Then
\(BC^{-1}B^*\) is exponentially local.  If in addition
Lemma~\ref{lem:v5v90-Schur-coercivity} holds, the retained Schur operator
\(S_M\) has a volume-uniform gap and exponentially local entries.  Whenever
its weighted conjugation error is smaller than half that gap,
\(S_M^{-1}\) has the pointwise and \(\ell^1\) bounds of
Theorem~\ref{thm:v5v90-CT}.
\end{proposition}

\begin{proof}
The kernel
\[
 (BC^{-1}B^*)_{xy}=\sum_{z,w}B_{xz}(C^{-1})_{zw}B^*_{wy}
\]
has only boundedly many \(z\) near \(x\) and \(w\) near \(y\).  The
exponential estimate for \(C^{-1}\), together with
\(d(z,w)\ge d(x,y)-2R_B\), proves exponential locality and uniform weighted
row and column bounds.  Equation \eqref{eq:v5v90-Schur-lower} supplies the
gap of \(S_M\), and Corollary~\ref{cor:v5v90-exp-local} applies.
\end{proof}

\subsection{Which Standard Model directions are massive-safe?}

Let \(v\) be a nonzero Higgs vacuum and let \(T_a\) denote the
anti-Hermitian infinitesimal gauge generators in the Higgs representation.
Up to the chosen normalization of the gauge fields, the gauge-boson mass
quadratic form is
\begin{equation}
 (M_A^2)_{ab}=g_ag_b\,
 \operatorname{Re}\langle T_av,T_bv\rangle.
 \label{eq:v5v90-gauge-mass-matrix}
\end{equation}

\begin{lemma}[Kernel of the Higgs gauge mass]
\label{lem:v5v90-mass-kernel}
The matrix \(M_A^2\) is positive semidefinite, and its kernel is the Lie
algebra of the stabilizer of \(v\):
\begin{equation}
 \sum_a c_a(M_A^2)_{ab}c_b=0
 \quad\Longleftrightarrow\quad
 \sum_a g_ac_aT_av=0.
 \label{eq:v5v90-stabilizer}
\end{equation}
Thus in the electroweak broken phase it is strictly positive on the
massive \(W^\pm,Z\) subspace and vanishes on the electromagnetic photon
direction.
\end{lemma}

\begin{proof}
The quadratic form in \eqref{eq:v5v90-gauge-mass-matrix} is the squared
norm \(\|\sum_ag_ac_aT_av\|^2\).  It vanishes exactly for infinitesimal
generators fixing \(v\).  The last statement is the defining breaking
pattern \(SU(2)_L\times U(1)_Y\to U(1)_{\rm em}\).
\end{proof}

For the convention
\(V(\phi)=\lambda(|\phi|^2-v^2/2)^2\), the radial Hessian at the minimum is
strictly positive, while tangent directions to the gauge orbit are
Goldstone directions.  An \(R_\xi\) gauge can make the gauge--Goldstone
auxiliary block coercive for \(\xi>0\), after the exact gauge kernel is
removed.  This does not give a photon mass.  The high Pauli--Villars block
has the singular gap already isolated in V68.  These observations give the
following conditional certification.

\begin{theorem}[Massive-safe SM one-form card]
\label{thm:v5v90-massive-card}
Fix a broken-phase chart on which:
\begin{enumerate}
\item the radial Higgs Hessian and the nonzero eigenvalues of
      \(M_A^2\) are bounded below uniformly;
\item gauge fixing and pinning have removed exact gauge zero modes and
      made the retained auxiliary gauge--Goldstone block uniformly
      positive;
\item the Pauli--Villars singular block has its V68 uniform gap;
\item the normalized off-diagonal mixing has norm at most \(q<1\); and
\item the full real one-form operator obeys either the finite-range
      locality bound of Theorem~\ref{thm:v5v90-CT} or the weighted bound
      \eqref{eq:v5v90-weighted-commutator}.
\end{enumerate}
Then the conditional covariance inverse on the radial-Higgs,
\(W^\pm\), \(Z\), gauge-fixed massive auxiliary, and PV directions has an
exponentially decaying kernel and a volume-uniform \(\ell^1\) row and
column bound.  Consequently the massive-sector one-form covariance input
assumed in V89 follows on this chart.
\end{theorem}

\begin{proof}
Items 1--3 give positive diagonal gaps.  Item 4 and
Lemma~\ref{lem:v5v90-Schur-coercivity} give the full one-form gap.  Item 5
and Theorem~\ref{thm:v5v90-CT}, or its exponentially local corollary, give
the inverse-kernel bound.  The one-form Helffer--Sj\"ostrand representation
used in V89 inserts this inverse between the two local gradients, so
\eqref{eq:v5v90-inverse-l1} is exactly its required kernel estimate.
\end{proof}

\subsection{Why the massless directions must stay outside}

\begin{proposition}[A massless Schur block transmits its tail]
\label{prop:v5v90-massless-tail}
Let \(C_0^{-1}\) be a massless Green kernel with a power-law tail, and let
\(B\) be a nonzero finite-range coupling.  Then the formal retained
operator
\begin{equation}
 S_0=A-BC_0^{-1}B^*
 \label{eq:v5v90-massless-Schur}
\end{equation}
is generally only power-law local.  Exponential locality follows only if
the coupling annihilates the massless channel or if additional moment
cancellations convert the transmitted tail to an independently summable
polynomial grade.
\end{proposition}

\begin{proof}
The kernel of the second term is the convolution displayed in the proof of
Proposition~\ref{prop:v5v90-local-Schur}.  Finite-range factors change each
endpoint by only a bounded distance and therefore do not change a generic
power-law decay into exponential decay.  Exact annihilation removes the
term.  If the finite-range factors are differences, their vanishing
moments differentiate the Green kernel and improve its power, precisely as
in V84--V89; this is a polynomial, not exponential, conclusion.
\end{proof}

Thus the electromagnetic photon and the gravitational elliptic constraint
response belong to the massless grade.  The first requires charge
neutrality or Ward differences; the second requires the monopole and
dipole matching isolated in V85 and V89.  The Yang--Mills V96 fixed-
\(p\) estimate is compatible with the massive card, but it cannot replace
the spatial locality hypothesis: an operator can have a bounded inverse
on every fixed \(L^p\) space while retaining a long-range kernel.

\begin{firewall}[Convexity is local in field space]
\label{fw:v5v90-convexity-chart}
The broken-phase gap can close at a Higgs zero, across a phase boundary, or
along a tunnelling direction.  Theorem~\ref{thm:v5v90-massive-card} is a
chartwise conditional theorem; it is not a global convexity theorem for
the gauge--Higgs measure.
\end{firewall}

\begin{firewall}[The effective fermion Hessian may be complex]
\label{fw:v5v90-complex-Hessian}
After integrating chiral fermions, the effective action need not define a
real positive Hessian.  The block theorem applies to a proved real
positive conditional reference measure.  The residual determinant phase
must still be treated by the rooted cancellation mechanism of V88--V89.
\end{firewall}

\begin{firewall}[Uniformity remains a cutoff problem]
\label{fw:v5v90-uniformity-open}
The constants \(g,q,J,R\), or their exponentially weighted analogues, must
be controlled after canonical rescaling and uniformly along the cutoff
chain.  The algebra above proves what those estimates imply; it does not
supply their interacting cutoff-uniform values.
\end{firewall}

\begin{theorem}[V90 conclusion]
\label{thm:v5v90-status}
On a properly quotiented, strictly massive gauge--Higgs--PV chart, a
normalized block-mixing bound \(q<1\) gives a positive one-form gap.  A
finite-range or exponentially weighted locality bound then gives an
explicit exponentially decaying inverse and the volume-uniform
\(\ell^1\) covariance kernel required in V89.  The photon and gravitational
constraint directions cannot be included in this card: their Schur
elimination transmits a power-law Green tail unless Ward or multipole
cancellations improve it.  Global chart coverage, the real-positive
reference measure, interacting dipole matching, and cutoff-uniform
constants remain open.
\end{theorem}

\begin{proof}
Combine Lemma~\ref{lem:v5v90-Schur-coercivity},
Theorem~\ref{thm:v5v90-CT}, Corollary~\ref{cor:v5v90-exp-local},
Proposition~\ref{prop:v5v90-local-Schur}, and
Theorem~\ref{thm:v5v90-massive-card}.  The exclusions and their reason are
Proposition~\ref{prop:v5v90-massless-tail} and the three firewalls.
\end{proof}

V91 will keep the newly certified massive inverse fixed and analyze the
joint photon--constraint polynomial grade.  The immediate target is an
exact convolution estimate showing which neutral Ward differences and
which gravitational moment cancellations make the mixed covariance row
summable without assigning a false mass to either field.

\clearpage
\section{V91: the joint photon--constraint polynomial grade}
\label{sec:v5v91}

V90 proved the one-form covariance estimate needed by V89 on a strictly
massive block.  This note treats the complementary low-frequency question.
The central issue is not merely whether a source is neutral, but how many
moments vanish at each end of a massless Green kernel.

\subsection{Localized rooted sources and moment order}

Work on \(\mathbb Z^3\), or on a periodic box on distance scales below a
fixed fraction of its side length.  A rooted source of order \(r\ge0\) is a
family \(f_x(z)\), supported in \(|z-x|\le R_0\), such that
\begin{equation}
 \sum_z(z-x)^\alpha f_x(z)=0
 \quad\text{for every multi-index }\alpha\text{ with }|\alpha|<r.
 \label{eq:v5v91-moment-order}
\end{equation}
Thus \(r=1\) is charge neutrality and \(r=2\) is neutrality plus vanishing
dipole.  Set
\begin{equation}
 M_r(f):=\sup_x\sum_z|f_x(z)|(1+|z-x|)^r.
 \label{eq:v5v91-moment-norm}
\end{equation}

Let \(G(z,w)\) be a massless covariance kernel satisfying, away from the
diagonal, the discrete difference bounds
\begin{equation}
 |\nabla_z^\alpha\nabla_w^\beta G(z,w)|
 \le C_{\alpha\beta}(1+|z-w|)^{-1-|\alpha|-|\beta|}
 \label{eq:v5v91-Green-differences}
\end{equation}
for the orders used below.  Local diagonal and contact terms will be
absorbed into constants.

\begin{lemma}[Discrete multipole remainder]
\label{lem:v5v91-multipole-remainder}
Let \(f_x\) have order \(r\) and support radius \(R_0\).  If a function
\(H\) has bounded discrete differences through order \(r\) on the
\(R_0\)-neighbourhood of \(x\), then
\begin{equation}
 \left|\sum_z f_x(z)H(z)\right|
 \le C_{r,R_0}M_r(f)
 \max_{|\alpha|=r}\sup_{|u-x|\le R_0}|\nabla^\alpha H(u)|.
 \label{eq:v5v91-Taylor-bound}
\end{equation}
\end{lemma}

\begin{proof}
Choose coordinate lattice paths from \(x\) to every point in the support.
Repeated use of the fundamental difference identity
\(H(u+e_j)-H(u)=\nabla_jH(u)\) gives the discrete Taylor expansion of
\(H(z)\) at \(x\) through degree \(r-1\), plus a sum of order-
\(r\) differences over at most \(C_{r,R_0}(1+|z-x|)^r\) path cells.
Every polynomial term of degree below \(r\) vanishes after multiplication
by \(f_x(z)\) and summation, by \eqref{eq:v5v91-moment-order}.  Taking
absolute values of the remainder proves \eqref{eq:v5v91-Taylor-bound}.
\end{proof}

\begin{theorem}[Two-ended Coulomb cancellation]
\label{thm:v5v91-two-ended}
Let \(f_x\) and \(h_y\) be rooted sources of orders \(r\) and \(s\), with
uniform support radius and finite moment norms.  Define
\begin{equation}
 K_{xy}:=\sum_{z,w}f_x(z)G(z,w)h_y(w).
 \label{eq:v5v91-sandwich}
\end{equation}
Under \eqref{eq:v5v91-Green-differences},
\begin{equation}
 |K_{xy}|\le C M_r(f)M_s(h)
 (1+|x-y|)^{-1-r-s}.
 \label{eq:v5v91-sandwich-decay}
\end{equation}
Consequently,
\begin{equation}
 \sup_x\sum_y|K_{xy}|<\infty
 \quad\text{whenever}\quad r+s>2.
 \label{eq:v5v91-summable-threshold}
\end{equation}
At \(r+s=2\), the displayed absolute bound has a logarithmically divergent
row sum and does not prove the conclusion.
\end{theorem}

\begin{proof}
When \(|x-y|\le4R_0+4\), bounded support and the local covariance bound
give \eqref{eq:v5v91-sandwich-decay} after increasing \(C\).  Otherwise
apply Lemma~\ref{lem:v5v91-multipole-remainder} first in \(z\), with
\(H(z)=G(z,w)\), and then in \(w\).  The support points remain at distance
comparable to \(|x-y|\), and \eqref{eq:v5v91-Green-differences} gives
exactly \((1+|x-y|)^{-1-r-s}\).  A shell of radius \(n\) in three
dimensions contains \(O(n^2)\) sites, so the row sum is dominated by
\(\sum_{n\ge1}n^{1-r-s}\), which converges exactly when \(r+s>2\).
\end{proof}

\begin{corollary}[The three relevant thresholds]
\label{cor:v5v91-three-thresholds}
For a three-dimensional Coulomb covariance:
\begin{enumerate}
\item two merely neutral insertions, \((r,s)=(1,1)\), give the borderline
      bound \((1+|x-y|)^{-3}\);
\item a neutral insertion and a neutral, dipole-free insertion,
      \((r,s)=(1,2)\), give the summable bound
      \((1+|x-y|)^{-4}\); and
\item two neutral, dipole-free insertions, \((r,s)=(2,2)\), give the
      summable bound \((1+|x-y|)^{-5}\).
\end{enumerate}
\end{corollary}

\begin{proof}
Substitute the three pairs into
\eqref{eq:v5v91-sandwich-decay} and use
\eqref{eq:v5v91-summable-threshold}.
\end{proof}

\subsection{Difference factorization and Ward data}

For translation-covariant finite-range stencils, the moments have an
algebraic form.  Let \(\nabla_j^*q(z)=q_j(z)-q_j(z-e_j)\).

\begin{proposition}[First and second difference cards]
\label{prop:v5v91-difference-cards}
Let \(f\) be a finitely supported scalar stencil on \(\mathbb Z^3\).
\begin{enumerate}
\item If \(\sum_zf(z)=0\), there are finitely supported \(q_j\) such that
\begin{equation}
 f=\sum_j\nabla_j^*q_j.
 \label{eq:v5v91-first-difference}
\end{equation}
\item If in addition \(\sum_zz_kf(z)=0\) for every \(k\), there are
finitely supported \(Q_{jk}\) such that
\begin{equation}
 f=\sum_{j,k}\nabla_j^*\nabla_k^*Q_{jk}.
 \label{eq:v5v91-second-difference}
\end{equation}
\end{enumerate}
The coefficient norms of \(q,Q\) are bounded by a constant depending only
on the support radius times the corresponding moment norm of \(f\).
\end{proposition}

\begin{proof}
Translate the support into a fixed box and identify the stencil with the
Laurent polynomial \(F(t)=\sum_zf(z)t^z\).  The first condition is
\(F(1)=0\), so \(F\) belongs to the ideal generated by
\(t_j-1\).  Division one coordinate at a time gives Laurent polynomials
\(Q_j\) with \(F=\sum_j(t_j-1)Q_j\), which is
\eqref{eq:v5v91-first-difference} after a harmless shift.  The first
moments are \(\partial_{t_k}F(1)\).  Their vanishing says that the image of
\(F\) in the quotient by the square of the maximal ideal at \(t=1\)
vanishes; hence \(F\) lies in the square of that ideal and has the second
factorization.  All spaces involved for a fixed support box are
finite-dimensional.  A linear right inverse on the subspace satisfying
the moment constraints gives the stated norm bounds.
\end{proof}

The first card is the exact finite-range expression of a neutral Ward
insertion.  The second is the dipole match used for the constraint source
in V85.  Summation by parts moves these differences onto \(G\), giving an
alternative proof of Theorem~\ref{thm:v5v91-two-ended} for such stencils.

\begin{theorem}[Joint photon--constraint mixed row sum]
\label{thm:v5v91-mixed-row}
Assume the rooted photon derivative is a uniformly local neutral stencil,
so it has the first factorization
\eqref{eq:v5v91-first-difference}.  Assume the substituted gravitational
constraint derivative has vanishing monopole and dipole and therefore the
second factorization \eqref{eq:v5v91-second-difference}.  If their mixed
massless covariance obeys the Green difference bounds
\eqref{eq:v5v91-Green-differences}, then
\begin{equation}
 \sup_x\sum_y
 \left|\operatorname{Cov}
   (\theta_x^{\rm ph},\theta_y^{\rm gr})\right|
 \le C M_1(f^{\rm ph})M_2(f^{\rm gr}).
 \label{eq:v5v91-mixed-covariance-row}
\end{equation}
The same holds with photon and constraint interchanged.
\end{theorem}

\begin{proof}
The covariance is the sandwich \eqref{eq:v5v91-sandwich} with
\((r,s)=(1,2)\), up to local matrix indices and contact terms.  Apply
Theorem~\ref{thm:v5v91-two-ended}; the decay is \(r^{-4}\), whose row sum
is finite.  A contact term has finite range and hence a uniformly finite
row sum.
\end{proof}

\subsection{The photon--photon obstruction and two valid exits}

Charge neutrality at each end supplies \((r,s)=(1,1)\), which is not
enough for the absolute cluster norm used in V88.  This is a limitation of
that norm, rather than a claim that the photon covariance does not exist.
There are two rigorous ways to close the specific row-sum card.

\begin{proposition}[Extra moment at one photon end]
\label{prop:v5v91-extra-photon-moment}
If every rooted photon insertion is neutral and dipole-free, then its
photon--photon covariance obeys
\begin{equation}
 |\operatorname{Cov}(\theta_x^{\rm ph},\theta_y^{\rm ph})|
 \le C(1+|x-y|)^{-5},
 \label{eq:v5v91-photon-photon-r5}
\end{equation}
and has a uniform absolute row sum.  More generally, for a nonsymmetric
pair of observables, it is enough that one insertion have order two and
the other order one.
\end{proposition}

\begin{proof}
Use Theorem~\ref{thm:v5v91-two-ended} with \((2,2)\), or with
\((2,1)\) in the nonsymmetric case.
\end{proof}

\begin{proposition}[Block-neutral aggregation]
\label{prop:v5v91-block-neutral}
Let \(\Theta_B=\sum_{x\in B}a_x\theta_x^{\rm ph}\), where the coefficients
in every fixed-diameter block obey \(\sum_{x\in B}a_x=0\).  If the
individual rooted photon insertions have order one, then the aggregated
insertion has order at least two.  Covariances between such block-neutral
observables therefore satisfy the summable estimate
\eqref{eq:v5v91-photon-photon-r5}.
\end{proposition}

\begin{proof}
Write the aggregate source as \(F_B(z)=\sum_{x\in B}a_xf_x(z)\).  Its
monopole vanishes because each \(f_x\) is neutral.  Translation covariance
gives
\(\sum_z z_jf_x(z)=m_j+x_j\sum_zf_x(z)=m_j\), independent of \(x\).
Hence \(\sum_zz_jF_B(z)=m_j\sum_xa_x=0\).  The aggregate has order two,
and Proposition~\ref{prop:v5v91-extra-photon-moment} applies.
\end{proof}

\begin{firewall}[A Ward identity does not automatically kill the dipole]
\label{fw:v5v91-Ward-dipole}
Gauge invariance normally proves the zero-momentum condition
\(\sum_zf_x(z)=0\).  Its derivative at zero momentum is an additional
identity and can contain contact or seagull terms.  The order-two photon
card must be proved for the complete differentiated observable, not
inferred from charge neutrality alone.
\end{firewall}

\begin{firewall}[Conditional cancellation is weaker than the V88 norm]
\label{fw:v5v91-conditional-sum}
At the \(r^{-3}\) threshold, angular or tensor cancellation may define a
principal-value operator bounded on \(L^p\).  This does not give the
absolute \(\ell^1\) row sum used in the V88 rooted cumulant estimate.
Replacing that norm would require redoing the cumulant argument in a
Calder\'on--Zygmund or neutral polymer space.
\end{firewall}

\begin{firewall}[Interacting gravitational moments remain open]
\label{fw:v5v91-interacting-moments}
V85 proves the double-difference card for the matched free adjacent-cutoff
stress defect on a smooth band.  Theorem~\ref{thm:v5v91-mixed-row} becomes
an interacting result only after monopole and dipole cancellation are
proved for the complete stress insertion, including measure, ghost,
gauge-fixing, and fermion contact terms.
\end{firewall}

\begin{theorem}[V91 conclusion]
\label{thm:v5v91-status}
For a three-dimensional Coulomb kernel, rooted cancellations of orders
\(r\) and \(s\) produce the exact decay threshold
\((1+R)^{-1-r-s}\), and the absolute covariance row sum is finite when
\(r+s>2\).  A neutral photon insertion and a monopole--dipole matched
gravitational insertion therefore have a summable mixed covariance.  Two
merely neutral photon insertions remain logarithmically borderline; the
V88 norm closes if one gains an extra moment at one end, or after a
block-neutral aggregation that creates that moment.  The complete
interacting Ward and stress moment identities are still unproved.
\end{theorem}

\begin{proof}
Use Theorem~\ref{thm:v5v91-two-ended},
Proposition~\ref{prop:v5v91-difference-cards},
Theorem~\ref{thm:v5v91-mixed-row}, and Propositions
\ref{prop:v5v91-extra-photon-moment}--\ref{prop:v5v91-block-neutral}, with
the three firewalls.
\end{proof}

V92 will attack the missing extra photon moment at the level of the
complete differentiated lattice observable.  It will separate exact
background-gauge Ward identities, contact terms, and the low-momentum
Taylor coefficient, and will record precisely which coefficient is a
renormalization condition rather than a consequence of gauge invariance.

\clearpage
\section{V92: the complete photon Ward jet}
\label{sec:v5v92}

V91 reduced the photon--photon absolute cluster bound to one additional
moment.  This note determines exactly what background gauge invariance does
and does not imply about that moment.

\subsection{A finite-support lattice Ward lemma}

Orient the edges of \(\mathbb Z^3\).  For a finitely supported one-cochain
\(f=(f_i(n))\), set
\begin{equation}
 (\delta f)(n)=\sum_i\bigl(f_i(n)-f_i(n-e_i)\bigr).
 \label{eq:v5v92-codifferential}
\end{equation}
Let \(z_i\) denote the Laurent Fourier variables and
\(p_i=z_i-1\).

\begin{lemma}[Compactly supported Ward cochains are curls]
\label{lem:v5v92-Ward-curl}
If \(f\) has finite support and \(\delta f=0\), then there are finitely
supported cochains \(q_{ij}=-q_{ji}\) such that
\begin{equation}
 f_i=\sum_j\nabla_j^*q_{ij}.
 \label{eq:v5v92-curl-factorization}
\end{equation}
In particular,
\begin{equation}
 \sum_nf_i(n)=0
 \quad\text{for every }i.
 \label{eq:v5v92-component-neutrality}
\end{equation}
\end{lemma}

\begin{proof}
Associate to \(f_i\) its Laurent polynomial \(F_i(z)\).  Up to multiplication
by Laurent monomials, \(\delta f=0\) is the syzygy
\begin{equation}
 \sum_ip_iF_i=0.
 \label{eq:v5v92-Laurent-syzygy}
\end{equation}
The sequence \(p_1,p_2,p_3\) is regular in the Laurent polynomial ring.
Its first syzygy module is therefore generated by the Koszul relations
\(p_je_i-p_ie_j\).  For completeness, eliminate \(p_3\): setting
\(p_1=p_2=0\) in \eqref{eq:v5v92-Laurent-syzygy} shows that \(F_3\) is
divisible by \(p_1\) or \(p_2\) after the elementary two-variable division;
subtract the corresponding Koszul relations.  The remainder has
\(F_3=0\) and satisfies \(p_1F_1+p_2F_2=0\).  Coprimality of \(p_1,p_2\)
gives \(F_1=p_2Q\), \(F_2=-p_1Q\).  Reversing the elimination yields
antisymmetric Laurent polynomials \(Q_{ij}\) with
\(F_i=\sum_jp_jQ_{ij}\).  Transforming back gives
\eqref{eq:v5v92-curl-factorization}, up to the shift convention in
\(\nabla^*\).  Evaluating at \(z=(1,1,1)\) proves
\eqref{eq:v5v92-component-neutrality}.
\end{proof}

The algebraic statement is stronger than the formal momentum equation
\(k_if_i(k)=0\): compact support excludes a harmonic constant one-form.
It supplies the order-one photon insertion assumed in V91.

\subsection{The unconstrained antisymmetric first jet}

For a source rooted at \(x\), define its first-moment tensor
\begin{equation}
 M_{ij}(x):=\sum_n(n_j-x_j)f_{x,i}(n).
 \label{eq:v5v92-first-moment}
\end{equation}

\begin{proposition}[Ward identity fixes only the symmetric jet]
\label{prop:v5v92-antisymmetric-jet}
If \(f_x\) is finitely supported and \(\delta f_x=0\), then
\begin{equation}
 M_{ij}(x)+M_{ji}(x)=0.
 \label{eq:v5v92-M-antisymmetric}
\end{equation}
Conversely, every antisymmetric tensor occurs as the first moment of a
finite-support co-closed one-cochain.  Therefore background gauge
invariance alone does not imply the order-two moment condition of V91.
\end{proposition}

\begin{proof}
Expand the Laurent symbol at zero momentum.  By
\eqref{eq:v5v92-component-neutrality},
\begin{equation}
 \widehat f_i(k)=iM_{ij}k_j+O(|k|^2).
 \label{eq:v5v92-symbol-jet}
\end{equation}
The Ward identity is
\(\sum_i(1-e^{-ik_i})\widehat f_i(k)=0\).  Its homogeneous quadratic
part is \(-k_iM_{ij}k_j\), so the symmetric part of \(M\) vanishes.
Conversely, for antisymmetric \(M\), take a single rooted plaquette
two-cochain \(q_{ij}\) with coefficient \(M_{ij}\) and set
\(f_i=\sum_j\nabla_j^*q_{ij}\).  It is co-closed because commuting
differences contract a symmetric operator with an antisymmetric tensor,
and direct summation gives its first moment proportional to \(M\); the
plaquette orientation fixes the harmless sign and normalization.
\end{proof}

The surviving tensor is the coefficient of a local field-strength term.
It is a physical local Taylor coefficient, not a gauge anomaly.

\begin{lemma}[Local field-strength subtraction]
\label{lem:v5v92-field-strength-subtraction}
Let \(f_x\) be the derivative of a local gauge-invariant rooted observable
\(\Theta_x(A)\), and let \(M_{ij}(x)\) be
\eqref{eq:v5v92-first-moment}.  There is a local gauge-invariant linear
combination \(C_x(A)\) of plaquette field strengths, supported a bounded
distance from \(x\), whose derivative has the same first-moment tensor.
Consequently the matched observable
\begin{equation}
 \Theta_x^{(2)}(A):=\Theta_x(A)-C_x(A)
 \label{eq:v5v92-matched-observable}
\end{equation}
has a derivative with zero monopole and zero dipole.
\end{lemma}

\begin{proof}
The converse construction in
Proposition~\ref{prop:v5v92-antisymmetric-jet} is precisely the derivative
of an oriented plaquette field strength.  Sum those plaquette observables
with coefficients chosen to reproduce \(M_{ij}(x)\).  Both derivatives
are co-closed, their component sums vanish, and their first moments agree.
Their difference therefore satisfies the order-two conditions.
\end{proof}

\begin{corollary}[A symmetry route to the second moment]
\label{cor:v5v92-symmetry-route}
Suppose the rooted observable and its reference background are invariant
under the full spatial cubic group including reflections, and carry no
external vector or pseudovector index.  Then \(M_{ij}(x)=0\), so its photon
derivative already has order two.
\end{corollary}

\begin{proof}
In three dimensions an antisymmetric rank-two tensor is dual to a
pseudovector.  No nonzero pseudovector is invariant under the full cubic
group including reflections.  Proposition~\ref{prop:v5v92-antisymmetric-jet}
then forces \(M=0\).
\end{proof}

On a curved, anisotropic, boundary, or rooted background the symmetry
hypothesis need not hold.  In that case
\eqref{eq:v5v92-matched-observable} is a genuine matching condition.

\subsection{The complete differentiated observable}

Let the regulated rooted functional have the schematic form
\begin{equation}
 \Theta_x(A)=
 \operatorname{Im}\lambda_x(D_{\rm high}(A))
 +\Theta_x^{\rm PV}(A)+\Theta_x^{\rm meas}(A)
 +\Theta_x^{\rm ct}(A).
 \label{eq:v5v92-complete-root}
\end{equation}
Here the four terms denote the high-Dirac trace-loop root, regulator,
measure/Jacobian, and local counterterm contributions.  Differentiation
must be performed before any term is discarded:
\begin{equation}
 f_{x,e}=d_{A_e}\Theta_x
 =f^{\rm high}_{x,e}+f^{\rm PV}_{x,e}
  +f^{\rm meas}_{x,e}+f^{\rm ct}_{x,e}.
 \label{eq:v5v92-complete-insertion}
\end{equation}

\begin{theorem}[Complete Ward-to-moment card]
\label{thm:v5v92-complete-card}
Assume that the finite-cutoff, anomaly-cancelled background gauge identity
holds for the complete rooted functional:
\begin{equation}
 \Theta_x(A+d\omega)=\Theta_x(A)
 \quad\text{for every compactly supported }\omega,
 \label{eq:v5v92-background-Ward}
\end{equation}
and that \(f_x\) in \eqref{eq:v5v92-complete-insertion} has uniformly
bounded support, or an exponentially local approximation with a uniformly
summable tail.  Then:
\begin{enumerate}
\item \(\delta f_x=0\), and the finite-range part has the exact curl
      factorization \eqref{eq:v5v92-curl-factorization};
\item its monopole vanishes component by component;
\item its first moment is antisymmetric and is the coefficient of a local
      field-strength insertion;
\item after either the symmetry condition of
      Corollary~\ref{cor:v5v92-symmetry-route} or the local matching
      \eqref{eq:v5v92-matched-observable}, the photon source has order two
      and its photon--photon covariance has the summable \(r^{-5}\) bound
      of V91.
\end{enumerate}
\end{theorem}

\begin{proof}
Differentiate \eqref{eq:v5v92-background-Ward} at \(\omega=0\):
\[
 0=\sum_ef_{x,e}(d\omega)_e
  =\sum_n(\delta f_x)(n)\omega(n).
\]
Arbitrariness of \(\omega\) gives \(\delta f_x=0\).  The first three
claims follow from Lemma~\ref{lem:v5v92-Ward-curl} and
Proposition~\ref{prop:v5v92-antisymmetric-jet}; exponentially small tails
are controlled by the same weighted sums and do not alter the stated
uniform estimates when their moments are included.  The fourth claim is
Lemma~\ref{lem:v5v92-field-strength-subtraction} or
Corollary~\ref{cor:v5v92-symmetry-route}, followed by
Proposition~\ref{prop:v5v91-extra-photon-moment}.
\end{proof}

\subsection{Anomaly and contact-term firewalls}

\begin{firewall}[Anomaly cancellation must hold for the regulator]
\label{fw:v5v92-regulated-anomaly}
The continuum Standard Model hypercharges cancel the perturbative gauge
anomaly after summing a complete generation.  This representation-theory
fact does not by itself prove
\eqref{eq:v5v92-background-Ward} for a chosen finite lattice regulator.
A Ward defect
\(\Theta_x(A+d\omega)-\Theta_x(A)=\mathcal A_x(A,\omega)\)
produces \(\delta f_x=d_A\mathcal A_x\), destroying the exact curl card
unless it is cancelled by the measure and local counterterms.
\end{firewall}

\begin{firewall}[The trace-loop term is not the complete current]
\label{fw:v5v92-complete-current}
Differentiating only \(\operatorname{Im}\lambda_x\) can omit regulator,
measure, seagull, and counterterm contacts.  Ward and moment identities
apply to the sum \eqref{eq:v5v92-complete-insertion}.  A cancellation among
its terms cannot be assigned to any one summand without a separate proof.
\end{firewall}

\begin{firewall}[Matching changes a local observable]
\label{fw:v5v92-matching-meaning}
The subtraction \eqref{eq:v5v92-matched-observable} is legitimate only
when the renormalization prescription permits the corresponding local
field-strength operator and fixes its coefficient.  It is not a free
algebraic deletion of a physical term.  If the target observable includes
that coefficient, its photon--photon covariance remains at the V91
borderline and must be estimated in a weaker neutral norm.
\end{firewall}

\begin{theorem}[V92 conclusion]
\label{thm:v5v92-status}
For a complete finite-cutoff photon insertion, exact background gauge
invariance implies a compactly supported curl factorization and
componentwise neutrality.  It leaves one antisymmetric first-moment
tensor, the coefficient of a local field-strength term.  Full cubic
reflection symmetry forces that tensor to vanish; otherwise an explicit
renormalization condition may match it.  Only after one of those two steps
does the rooted photon insertion acquire the second moment needed for the
V91 absolute photon--photon row sum.  Regulated anomaly cancellation and
the complete contact packet remain necessary inputs.
\end{theorem}

\begin{proof}
Combine Lemma~\ref{lem:v5v92-Ward-curl},
Proposition~\ref{prop:v5v92-antisymmetric-jet},
Lemma~\ref{lem:v5v92-field-strength-subtraction}, and
Theorem~\ref{thm:v5v92-complete-card}, subject to the three firewalls.
\end{proof}

V93 will turn the finite-support Ward lemma into an exponentially local
statement suited to the high-Dirac loop series.  The target is a controlled
finite-range truncation that repairs its small Ward defect locally, while
preserving exponential tails and the first two moments uniformly in the
cutoff.

\clearpage
\section{V93: gauge-covariant truncation of the rooted loop series}
\label{sec:v5v93}

The target announced in V92 was a local repair of the Ward defect produced
by truncation.  Going one step upstream removes the defect: truncate only
objects whose covariance law is exact.  A distance truncation of a
gauge-covariant kernel is still gauge covariant, and every closed diagonal
loop trace is then gauge invariant term by term.

\subsection{Covariant range and loop-order cutoffs}

Let \(R(A)\) be a matrix kernel on lattice sites and internal indices with
the Abelian background covariance law
\begin{equation}
 R(A+d\omega)_{xy}=U_\omega(x)R(A)_{xy}U_\omega(y)^{-1}.
 \label{eq:v5v93-kernel-covariance}
\end{equation}
For \(H\ge1\), define the range truncation
\begin{equation}
 R^{[H]}_{xy}:={\bf1}_{\{d(x,y)\le H\}}R_{xy}.
 \label{eq:v5v93-range-truncation}
\end{equation}
It obeys the same covariance law.  For \(M\ge1\), define
\begin{equation}
 \lambda_x^{M,H}(A):=
 \sum_{m=1}^{M}\frac{(-1)^{m+1}}m
 \operatorname{tr}_{\rm int}\bigl((R^{[H]}(A))^m\bigr)_{xx},
 \qquad
 \theta_x^{M,H}:=\operatorname{Im}\lambda_x^{M,H}.
 \label{eq:v5v93-double-truncation}
\end{equation}

\begin{lemma}[Exact Ward identity at every truncation]
\label{lem:v5v93-exact-truncated-Ward}
For every \(M,H\),
\begin{equation}
 \lambda_x^{M,H}(A+d\omega)=\lambda_x^{M,H}(A),
 \qquad
 \delta d_A\theta_x^{M,H}=0.
 \label{eq:v5v93-truncated-Ward}
\end{equation}
If the dependence of \(R^{[H]}_{uv}\) on the background is confined to a
fixed neighbourhood of paths of length at most \(H\), then
\(d_A\theta_x^{M,H}\) has support in a ball of radius \(C(MH+H)\) about
\(x\).
\end{lemma}

\begin{proof}
Matrix multiplication telescopes the endpoint transformations in
\eqref{eq:v5v93-kernel-covariance}:
\[
 ((R^{[H]})^m(A+d\omega))_{xx}
 =U_\omega(x)((R^{[H]})^m(A))_{xx}U_\omega(x)^{-1}.
\]
The internal trace is invariant under conjugation, proving the first
identity term by term.  Differentiation in \(\omega\) and lattice summation
by parts prove the second.  A closed product with \(m\) jumps of length at
most \(H\) stays within distance \(mH\) of \(x\); the additional fixed
background-dependence neighbourhood gives the stated radius.
\end{proof}

Thus a cutoff of individual gauge potentials or of an open path expansion
is unnecessary and potentially harmful.  The range projector
\eqref{eq:v5v93-range-truncation} commutes with endpoint gauge covariance.

\subsection{Derivative convergence with an explicit geometric tail}

Use the weighted Schur norm
\begin{equation}
 \|K\|_{\mu,1}:=
 \max\left\{
 \sup_x\sum_y e^{\mu d(x,y)}\|K_{xy}\|,
 \sup_y\sum_x e^{\mu d(x,y)}\|K_{xy}\|
 \right\}.
 \label{eq:v5v93-weighted-Schur}
\end{equation}
It is submultiplicative.  Let \(d_zR\) denote a background derivative
rooted at \(z\), measured in the corresponding rooted weighted norm of
V89.

\begin{lemma}[Range-tail estimate]
\label{lem:v5v93-range-tail}
If \(\|R\|_{\mu,1}\le\kappa<1\), then for
\(0<\mu'<\mu\),
\begin{equation}
 \|R-R^{[H]}\|_{\mu',1}
 \le e^{-(\mu-\mu')H}\kappa.
 \label{eq:v5v93-range-tail}
\end{equation}
The same statement holds for \(d_zR\) whenever its \(\mu\)-weighted
rooted norm is finite.
\end{lemma}

\begin{proof}
On \(d(x,y)>H\),
\(e^{\mu'd(x,y)}\le e^{-(\mu-\mu')H}e^{\mu d(x,y)}\).
Apply this pointwise in each row and column sum.  Differentiation commutes
with the fixed distance projector.
\end{proof}

\begin{theorem}[Differentiated loop-tail bound]
\label{thm:v5v93-loop-tail}
Assume
\begin{equation}
 \|R\|_{\mu,1}\le\kappa<1,
 \qquad
 \|d_zR\|_{\mu,1;x}le L a^\alpha e^{-\nu d(x,z)}.
 \label{eq:v5v93-R-hypotheses}
\end{equation}
For fixed \(H\), the omitted loop orders satisfy
\begin{equation}
 \left\|d_z\lambda_x(R^{[H]})
       -d_z\lambda_x^{M,H}\right\|
 \le \frac{L a^\alpha}{1-\kappa}\,
      \kappa^M e^{-\nu d(x,z)}.
 \label{eq:v5v93-order-tail}
\end{equation}
For \(0<\mu'<\mu\), the difference between the full and range-truncated
differentiated logarithms obeys
\begin{equation}
 \left\|d_z\lambda_x(R)-d_z\lambda_x(R^{[H]})\right\|
 \le C_{\kappa,\mu,\mu'}La^\alpha
 e^{-cH}e^{-\nu' d(x,z)},
 \label{eq:v5v93-spatial-tail}
\end{equation}
for any fixed \(\nu'<\min\{\nu,\mu'\}\), with
\(c>0\) depending only on the sacrificed exponential margins.  The
constants are volume independent.
\end{theorem}

\begin{proof}
Differentiate a power:
\begin{equation}
 d(R^m)=\sum_{\ell=0}^{m-1}R^\ell(dR)R^{m-1-\ell}.
 \label{eq:v5v93-power-derivative}
\end{equation}
Submultiplicativity bounds its norm by
\(m\kappa^{m-1}\|dR\|\).  The factor \(m\) cancels the \(1/m\) in the
logarithm, and summing \(m>M\) gives
\(\kappa^M/(1-\kappa)\), proving
\eqref{eq:v5v93-order-tail}.

For the range error, subtract the two versions of
\eqref{eq:v5v93-power-derivative}.  Telescope each product so that every
term contains either \(R-R^{[H]}\) or
\(dR-dR^{[H]}\).  Lemma~\ref{lem:v5v93-range-tail}, the geometric series
in \(m\), and the gap \(1-\kappa\) give an \(e^{-cH}\) bound in a weaker
weighted norm.  The rooted convolution inequality spends the remaining
exponential margin to retain \(e^{-\nu'd(x,z)}\).  This proves
\eqref{eq:v5v93-spatial-tail}.
\end{proof}

\subsection{Passing the first two moments to the limit}

For a rooted edge source \(f_x(e)\), write
\begin{equation}
 \|f\|_{(r)}:=\sup_x\sum_e
 (1+d(x,e))^r|f_x(e)|.
 \label{eq:v5v93-polynomial-moment-norm}
\end{equation}
Every exponential rooted bound controls this norm for fixed \(r\).

\begin{proposition}[Moment-stable Ward limit]
\label{prop:v5v93-moment-limit}
Choose \(M=M(H)\to\infty\).  Under
\eqref{eq:v5v93-R-hypotheses}, the finite-support sources
\begin{equation}
 f_x^{M,H}:=d_A\theta_x^{M,H}
 \label{eq:v5v93-truncated-source}
\end{equation}
converge to \(f_x=d_A\theta_x\) in \(\|\cdot\|_{(r)}\) for every fixed
\(r\), provided \(\kappa^{M(H)}(1+M(H)H)^r\to0\).  In particular:
\begin{enumerate}
\item the limiting source is co-closed and componentwise neutral;
\item its first-moment tensor is the limit of the antisymmetric tensors of
      the truncations; and
\item if a symmetry or a local field-strength matching condition makes
      every truncated first moment vanish, the full source has order two.
\end{enumerate}
\end{proposition}

\begin{proof}
Theorem~\ref{thm:v5v93-loop-tail} gives exponential convergence of the
range part and a loop-order error supported within radius
\(C(MH+H)\).  Multiplying the latter by its largest polynomial weight gives
the condition in the statement; exponential convergence controls all
polynomial moments of the former.  Hence convergence holds in
\(\|\cdot\|_{(r)}\).

The codifferential and the zeroth and first moment functionals are
continuous in \(\|\cdot\|_{(1)}\).  Lemma
\ref{lem:v5v93-exact-truncated-Ward} and V92 therefore pass co-closedness,
neutrality, antisymmetry, and any imposed zero first moment to the limit.
\end{proof}

One admissible diagonal choice is \(M(H)=H\): exponential decay of
\(\kappa^H\) dominates every power of \(H\).  This gives a single sequence
of exactly gauge-invariant finite-support observables converging with all
fixed moments.

\begin{corollary}[Exponentially local V91 estimate]
\label{cor:v5v93-exponential-V91}
The two-ended cancellation theorem of V91 remains valid for exponentially
localized rooted sources whose moments through orders \(r-1\) and
\(s-1\) vanish and whose corresponding polynomial moment norms are finite.
In particular, the order-two photon source obtained from
Proposition~\ref{prop:v5v93-moment-limit} has a volume-uniform summable
photon--photon covariance row.
\end{corollary}

\begin{proof}
Approximate both sources by the finite-support covariant truncations.
The V91 estimate is uniform in the truncation because its constant uses
only the fixed moment norms.  Convergence in those norms permits passage to
the limit by dominated convergence.  Equivalently, split each source into
a fixed ball and an exponentially small tail and apply the same Taylor
remainder proof directly.
\end{proof}

\subsection{What this construction requires from the regulator}

\begin{firewall}[Covariant kernels are an input]
\label{fw:v5v93-covariant-kernel}
Equation \eqref{eq:v5v93-kernel-covariance} must hold for the regulated
high operator and for every measure or PV contribution entering the
complete V92 root.  A gauge-breaking chiral regulator cannot be repaired
merely by the distance projector; its finite-cutoff Ward defect must first
be cancelled by a proved local counterterm construction.
\end{firewall}

\begin{firewall}[Uniform derivative locality is not implied by
\(\|R\|<1\)]
\label{fw:v5v93-derivative-locality}
The norm gap \(\kappa<1\) controls loop order.  It does not supply the
rooted derivative estimate in \eqref{eq:v5v93-R-hypotheses}.  The latter
can lose a lattice spacing under shell differentiation and remains one of
the exponent-accounting inputs isolated in V89.
\end{firewall}

\begin{firewall}[Matching must commute with the limit]
\label{fw:v5v93-matching-limit}
If the local field-strength coefficient depends on \(M,H\), it must
converge in the chosen renormalization scheme.  Setting each coefficient
to zero with unrelated cutoff-dependent counterterms would not define a
single continuum observable.  Proposition~\ref{prop:v5v93-moment-limit}
requires a coherent symmetry or matching prescription.
\end{firewall}

\begin{theorem}[V93 conclusion]
\label{thm:v5v93-status}
Gauge-covariant range truncation followed by closed-loop-order truncation
produces finite-support rooted observables satisfying the photon Ward
identity exactly at every stage.  The differentiated logarithmic tail is
bounded by \(\kappa^M/(1-\kappa)\), and exponential range locality controls
the spatial truncation.  A diagonal limit preserves co-closedness and the
first two moments, extending the V91 summable photon covariance card to
the exponentially local high-Dirac loop series.  Regulated covariance,
the derivative exponent, and a coherent field-strength matching
prescription remain required.
\end{theorem}

\begin{proof}
Use Lemma~\ref{lem:v5v93-exact-truncated-Ward},
Theorem~\ref{thm:v5v93-loop-tail},
Proposition~\ref{prop:v5v93-moment-limit}, and
Corollary~\ref{cor:v5v93-exponential-V91}, subject to the three firewalls.
\end{proof}

V94 will compute the shell-derivative exponent in canonical lattice units
for the high Wilson--Dirac defect.  It will distinguish a derivative with
respect to a dimensionless link angle from a derivative with respect to a
canonically normalized continuum gauge field, so the remaining
\(a^{-1}\) risk in V89 is decided rather than hidden in notation.

\clearpage
\section{V94: canonical shell derivatives and invariant power counting}
\label{sec:v5v94}

The derivative exponent left open in V89 cannot be assigned before the
field coordinate is specified.  A dimensionless link angle, a sampled
continuum gauge potential, and an \(L^2\)-normalized cell coordinate differ
by powers of the lattice spacing.  This note computes those powers and
shows that the covariance sandwich, rather than an isolated gradient, is
coordinate invariant.

\subsection{Three gauge-field coordinates}

In \(d\) Euclidean dimensions, write a link in a bounded exponential chart
as
\begin{equation}
 U_e=\exp(agA_e)=\exp(\vartheta_e),
 \qquad \vartheta_e:=agA_e.
 \label{eq:v5v94-link-coordinates}
\end{equation}
Here Lie-algebra generators and a harmless factor of \(i\) are suppressed.
Define the cell \(L^2\) coordinate
\begin{equation}
 q_e:=a^{d/2}A_e.
 \label{eq:v5v94-L2-coordinate}
\end{equation}
For an Abelian link, or uniformly on a bounded non-Abelian exponential
chart, the chain rule gives
\begin{equation}
 d_{A_e}=ag\,d_{\vartheta_e},
 \qquad
 d_{q_e}=g a^{1-d/2}d_{\vartheta_e}.
 \label{eq:v5v94-chain-rule}
\end{equation}
For the non-Abelian exponential, the differential is
\begin{equation}
 d\exp_\vartheta[X]
 =\int_0^1e^{(1-s)\vartheta}Xe^{s\vartheta}\,ds,
 \label{eq:v5v94-Duhamel-exponential}
\end{equation}
so the same powers of \(a\) hold with chart-uniform constants.

\begin{lemma}[Wilson--Dirac link derivative]
\label{lem:v5v94-Wilson-link-derivative}
Let \(\mathcal D_W=aD_W\) be the dimensionless Wilson--Dirac operator.
On a bounded link chart,
\begin{equation}
 \|d_{\vartheta_e}\mathcal D_W\|\le C,
 \qquad
 \|d_{A_e}\mathcal D_W\|\le Cga,
 \qquad
 \|d_{q_e}\mathcal D_W\|\le Cg a^{1-d/2}.
 \label{eq:v5v94-Wilson-derivatives}
\end{equation}
For the dimensionful operator \(D_W=a^{-1}\mathcal D_W\), every displayed
power is lower by one.
\end{lemma}

\begin{proof}
Every hopping term of \(\mathcal D_W\) contains a bounded spin projector
times one link matrix and has coefficient independent of \(a\).  Its
derivative with respect to the link angle is therefore uniformly bounded
by \eqref{eq:v5v94-Duhamel-exponential}.  Apply
\eqref{eq:v5v94-chain-rule}.  Multiplication by \(a^{-1}\) gives the last
statement.
\end{proof}

In the physical case \(d=4\), an \(L^2\)-cell derivative costs one power:
\(d_q=ga^{-1}d_\vartheta\).  This is the precise version of the risk noted
in V89.

\subsection{The coordinate-invariant covariance sandwich}

Let \(x\) be any finite-dimensional bosonic coordinate, \(S\) invertible,
and \(y=Sx\).  Let \(H_x\) be a positive one-form Hessian and
\(f_x=d_xF\).

\begin{proposition}[Invariance under field normalization]
\label{prop:v5v94-sandwich-invariance}
Under \(y=Sx\),
\begin{equation}
 f_y=S^{-T}f_x,
 \qquad
 H_y=S^{-T}H_xS^{-1},
 \qquad
 H_y^{-1}=SH_x^{-1}S^T,
 \label{eq:v5v94-transform-laws}
\end{equation}
and hence
\begin{equation}
 \langle f_y,H_y^{-1}f_y\rangle
 =\langle f_x,H_x^{-1}f_x\rangle.
 \label{eq:v5v94-invariant-sandwich}
\end{equation}
The bilinear version with two observables is invariant as well.
\end{proposition}

\begin{proof}
The first identity is the chain rule.  Substitute
\(x=S^{-1}y\) into the quadratic form to obtain the second, invert it for
the third, and multiply the three factors.  The matrices \(S^{-1}\) and
\(S\) cancel pairwise.
\end{proof}

\begin{corollary}[No standalone derivative exponent]
\label{cor:v5v94-no-standalone-alpha}
Suppose \(q=c(a)\vartheta\) is a scalar rescaling.  If
\(\|d_qF\|=|c|^{-1}\|d_\vartheta F\|\), then
\(H_q^{-1}=|c|^2H_\vartheta^{-1}\).  Thus the apparent loss in the
gradient is exactly offset in the covariance sandwich.  A claim of a
positive exponent for \(dF\) is meaningful only together with the norm and
the scaling of the one-form inverse used to measure it.
\end{corollary}

\begin{proof}
Take \(S=cI\) in
Proposition~\ref{prop:v5v94-sandwich-invariance}.
\end{proof}

This does not make power counting vacuous.  V90 assumes a covariance
inverse bounded in a specified coordinate and weighted spatial norm.  Once
that coordinate is fixed, the derivative exponent must beat the stated
bound.

\subsection{The adjacent-cutoff defect exponent}

Let \(R_a(\vartheta)\) be the dimensionless high-Dirac relative defect in
the V78--V93 logarithm.  Assume a matched smooth-band estimate
\begin{equation}
 \|R_a\|_{\mu,1}\le C a^\beta,
 \qquad
 \|d_{\vartheta_e}R_a\|_{\mu,1;x}
 \le C a^\beta e^{-\nu d(x,e)}.
 \label{eq:v5v94-beta-hypothesis}
\end{equation}
The second estimate says that differentiating the dimensionless link does
not destroy the matched symbol order.  This is an assumption about the
complete adjacent-cutoff defect, not a consequence of the first estimate.

\begin{theorem}[Canonical derivative exponent]
\label{thm:v5v94-canonical-exponent}
Under \eqref{eq:v5v94-beta-hypothesis}, and for \(a\) small enough that
\(Ca^\beta<1\), the rooted logarithmic phase satisfies
\begin{align}
 |d_{\vartheta_e}\theta_x|&\le C'a^\beta e^{-\nu d(x,e)},
 \label{eq:v5v94-theta-angle}\\
 |d_{A_e}\theta_x|&\le C'g a^{\beta+1}e^{-\nu d(x,e)},
 \label{eq:v5v94-theta-A}\\
 |d_{q_e}\theta_x|&\le C'g a^{\beta+1-d/2}e^{-\nu d(x,e)}.
 \label{eq:v5v94-theta-q}
\end{align}
In four dimensions the \(L^2\)-cell exponent is
\begin{equation}
 \alpha_q=\beta-1.
 \label{eq:v5v94-four-dimensional-exponent}
\end{equation}
\end{theorem}

\begin{proof}
The differentiated trace-log identity used in V89 is
\[
 d\lambda_x=\operatorname{tr}_{\rm int}
 \bigl((I+R_a)^{-1}dR_a\bigr)_{xx}.
\]
The Neumann inverse has norm at most
\((1-Ca^\beta)^{-1}\), which is absorbed into \(C'\).  This proves
\eqref{eq:v5v94-theta-angle}.  Equations
\eqref{eq:v5v94-theta-A}--\eqref{eq:v5v94-theta-q} follow from
\eqref{eq:v5v94-chain-rule}; set \(d=4\) for
\eqref{eq:v5v94-four-dimensional-exponent}.
\end{proof}

\begin{corollary}[Conditional improvement threshold]
\label{cor:v5v94-improvement-threshold}
If the one-form inverse in the \(q\) coordinate has the volume-uniform
massive bound of V90, then in four dimensions the massive rooted phase
variance obtained from this card is
\begin{equation}
 O(a^{2(\beta-1)}).
 \label{eq:v5v94-phase-variance-beta}
\end{equation}
Thus \(\beta>1\) yields a vanishing phase variance.  A merely first-order
defect \(\beta=1\) gives only an \(O(1)\) bound by this argument, whereas a
fully first-order-improved defect with \(\beta=2\) gives \(O(a^2)\).
\end{corollary}

\begin{proof}
Insert \eqref{eq:v5v94-theta-q} twice in the one-form covariance sandwich
and use the uniform inverse row bound.  The two gradients give the stated
power.  The three cases are immediate.
\end{proof}

The conclusion aligns with the free symbol audit in V86: an unimproved
Wilson defect is generically first order, while cancellation of the full
dimension-five packet can make the smooth-band defect second order.  The
statement here is conditional because the interacting derivative estimate
in \eqref{eq:v5v94-beta-hypothesis} still has to be proved.

\subsection{Shell averages and the massless grade}

Let a normalized shell coordinate be
\(Q_B=|B|^{-1/2}\sum_{e\in B}q_e\).  Its directional derivative is
\begin{equation}
 d_{Q_B}F=|B|^{-1/2}\sum_{e\in B}d_{q_e}F.
 \label{eq:v5v94-shell-derivative}
\end{equation}
For a source localized near one root, Cauchy--Schwarz or absolute summation
can improve the pointwise cost as the block grows.  For a coherent
zero-momentum source it need not.  The Ward moment conditions of V91--V93
are therefore still required in the photon and constraint grades; a change
of normalization cannot turn a physical massless pole into an
exponentially local inverse.

\begin{firewall}[The derivative estimate must be differentiated]
\label{fw:v5v94-differentiated-symbol}
The amplitude bound \(R_a=O(a^\beta)\) does not imply
\(d_\vartheta R_a=O(a^\beta)\).  Differentiation can expose an unmatched
contact or seagull term.  Equation \eqref{eq:v5v94-beta-hypothesis} must be
proved from the complete differentiated translation and gauge Ward packet
identified in V86 and V92.
\end{firewall}

\begin{firewall}[Smooth-band order is not a cutoff-shell order]
\label{fw:v5v94-smooth-versus-shell}
The values \(\beta=1,2\) quoted from the free Symanzik analysis concern a
fixed physical momentum band.  At momenta of order \(a^{-1}\), the same
Taylor expansion is unavailable.  A multiscale proof must combine the
smooth-band gain with high-shell resolvent decay rather than extend
\eqref{eq:v5v94-beta-hypothesis} by assertion.
\end{firewall}

\begin{firewall}[Coordinate invariance requires transforming the norm]
\label{fw:v5v94-norm-transform}
Using the \(q\)-gradient with the \(\vartheta\)-coordinate inverse, or vice
versa, violates \eqref{eq:v5v94-transform-laws} and creates a spurious
power of \(a\).  Every cutoff exponent must be attached to a complete
gradient--covariance--gradient expression in one coordinate convention.
\end{firewall}

\begin{theorem}[V94 conclusion]
\label{thm:v5v94-status}
For a dimensionless Wilson--Dirac operator, a link-angle derivative costs
no lattice power, a sampled continuum-field derivative gains one power,
and a four-dimensional \(L^2\)-cell derivative costs one power.  These
changes cancel exactly against the inverse-Hessian transformation in the
full covariance sandwich.  If the inverse is fixed and uniform in the
\(L^2\)-cell coordinate, a matched defect of order \(a^\beta\) gives the
phase-gradient exponent \(\beta-1\); the present route therefore needs
\(\beta>1\), with full first-order improvement giving the useful value
\(\beta=2\).  The differentiated interacting estimate and the cutoff-shell
regime remain open.
\end{theorem}

\begin{proof}
Use Lemma~\ref{lem:v5v94-Wilson-link-derivative},
Proposition~\ref{prop:v5v94-sandwich-invariance},
Theorem~\ref{thm:v5v94-canonical-exponent}, and
Corollary~\ref{cor:v5v94-improvement-threshold}, subject to the three
firewalls.
\end{proof}

V95 is the next compulsory YM/NS audit.  It will then split the
differentiated defect into a smooth physical band and a cutoff shell and
derive a summable two-regime estimate, using symbol improvement only where
the Taylor expansion is valid and resolvent suppression on the remaining
shells.

\clearpage
\section{V95: companion audit and a two-regime cutoff sum}
\label{sec:v5v95}

\subsection{Compulsory companion audit}

The active Yang--Mills file inspected for this checkpoint was
\texttt{Renewal\_Yang\_Mills\_Mass\_Gap\_Research\_Notes\_V99.tex}, with
\(1\,465\,401\) bytes and SHA--256
\begin{center}
\texttt{52F5EF6046E0E63FDBBF94EB9C2851F26E49B2AF9C60D1CDC1F82CBAB6444B9B}.
\end{center}
Its new endpoint is a complete four-label projected-tail estimate and the
resulting arity-five all-thermal local tree card.  Its first unknown
projected subset has size five.  The notes explicitly show why a
fixed-degree H\"older ladder does not iterate to arbitrary arity.  The
files ending in \texttt{\_f1}, \texttt{\_f2}, and \texttt{\_f3} are frozen
companion archives and were not treated as active files.

The Navier--Stokes file remained
\texttt{Renewal\_NS\_Stability\_Research\_Notes\_V31.tex}, with
\(887\,537\) bytes and SHA--256
\begin{center}
\texttt{F1121B62238AD5491BB7CF03635EA9934942EDF573ED8FAAA9EE79E1D38A6696}.
\end{center}
Its relevant endpoint is exact critical first-moment bookkeeping for a
dyadic heat-formation ladder.  It also proves that charging the same remote
mode to every lower tail recreates the missing critical weight; the valid
organization assigns each swept band to one parent use.

Two consequences guide the SM programme.  First, the smooth/cutoff split
below must be made with disjoint shell ownership rather than repeatedly
charging nested high tails.  Second, a two-point covariance estimate does
not close the all-order rooted cumulant expansion of V88: the Yang--Mills
audit presently supplies a model of the required tree card only through
arity five.

\subsection{An abstract smooth-band/cutoff-shell certificate}

Let \(a_n=a_0\rho^n\), \(0<\rho<1\), be the cutoff sequence.  Let \(T_n\)
be the differentiated adjacent-cutoff defect in one fixed field
coordinate, including the complete contact packet.  Choose a crossover
\begin{equation}
 \Lambda_n=a_n^{-\eta},\qquad \eta>0,
 \label{eq:v5v95-crossover}
\end{equation}
and decompose with spectral projectors
\(P_n=P_{\le\Lambda_n}\), \(Q_n=I-P_n\).

\begin{theorem}[Two-regime summability]
\label{thm:v5v95-two-regime}
Assume there are \(b>0\), \(r,s>0\), \(h\ge0\), and a volume-uniform
constant \(C\) such that
\begin{align}
 \|T_nP_n\|_{\mathcal X}
 &\le C a_n^b\Lambda_n^r,
 \label{eq:v5v95-low-bound}\\
 \|T_nQ_n\|_{\mathcal X}
 &\le C a_n^{-h}\Lambda_n^{-s}.
 \label{eq:v5v95-high-bound}
\end{align}
Here \(\mathcal X\) is the rooted weighted norm paired with the one-form
inverse.  If
\begin{equation}
 sb>rh,
 \label{eq:v5v95-feasibility}
\end{equation}
then one may choose
\begin{equation}
 \frac hs<\eta<\frac br
 \label{eq:v5v95-eta-window}
\end{equation}
and obtain
\begin{equation}
 \sum_{n\ge0}\|T_n\|_{\mathcal X}<\infty.
 \label{eq:v5v95-summable-T}
\end{equation}
The balanced choice
\begin{equation}
 \eta_*:=\frac{b+h}{r+s}
 \label{eq:v5v95-optimal-eta}
\end{equation}
gives both regimes the exponent
\begin{equation}
 \varepsilon_*:=\frac{sb-rh}{r+s}>0,
 \qquad
 \|T_n\|_{\mathcal X}\le 2C a_n^{\varepsilon_*}.
 \label{eq:v5v95-balanced-exponent}
\end{equation}
\end{theorem}

\begin{proof}
Substituting \eqref{eq:v5v95-crossover} into the two bounds gives
\begin{equation}
 \|T_n\|_{\mathcal X}
 \le C\bigl(a_n^{b-\eta r}+a_n^{\eta s-h}\bigr).
 \label{eq:v5v95-two-powers}
\end{equation}
Both exponents are positive exactly under
\eqref{eq:v5v95-eta-window}; that interval is nonempty exactly when
\eqref{eq:v5v95-feasibility} holds.  Equating the two exponents gives
\eqref{eq:v5v95-optimal-eta} and
\eqref{eq:v5v95-balanced-exponent}.  Since \(a_n\) is geometric, every
positive power of \(a_n\) is summable.
\end{proof}

For the four-dimensional \(L^2\)-cell convention of V94, a fully improved
smooth-band derivative has \(b=\beta-1=1\).  If the high-shell estimate
has no power loss, \(h=0\), then every \(s>0\) and finite Taylor growth
order \(r\) satisfy the feasibility condition.  The balanced gain is
\(s/(r+s)\).  By contrast, an unimproved \(\beta=1\) gives \(b=0\) and
cannot satisfy \eqref{eq:v5v95-feasibility} when \(h\ge0\).

\subsection{Why nested high tails must be disjointified}

Let \((\Pi_j)_{j\ge0}\) be mutually orthogonal physical momentum annuli.
The nested tail \(Q_n\) is a sum of many \(\Pi_j\)'s, so summing a bound for
every \(Q_n\) can charge a remote annulus repeatedly.  Let
\(\mathcal A_n\) be disjoint sets of annular indices assigned to cutoff
step \(n\), and put
\begin{equation}
 \widetilde Q_n:=\sum_{j\in\mathcal A_n}\Pi_j,
 \qquad
 \mathcal A_n\cap\mathcal A_m=\varnothing\quad(n\ne m).
 \label{eq:v5v95-owned-shells}
\end{equation}

\begin{lemma}[One-use shell accounting]
\label{lem:v5v95-one-use}
For nonnegative annular costs \(c_{n,j}\),
\begin{equation}
 \sum_n\sum_{j\in\mathcal A_n}c_{n,j}
 =\sum_j {\bf1}_{\{j\in\cup_n\mathcal A_n\}}c_{n(j),j},
 \label{eq:v5v95-one-use-identity}
\end{equation}
where \(n(j)\) is the unique owner.  Hence a bound
\(c_{n(j),j}\le b_j\) costs \(\sum_jb_j\) with no tail-incidence factor.
If instead the same \(b_j\) is charged to every nested tail containing
\(j\), its coefficient is the number of such tails and may diverge.
\end{lemma}

\begin{proof}
Equation \eqref{eq:v5v95-one-use-identity} is a rearrangement of a sum over
disjoint index sets.  The two conclusions follow by comparison.  No
orthogonality estimate is needed because the statement concerns the
absolute shell ledger itself.
\end{proof}

\begin{proposition}[Owned-shell version of the high estimate]
\label{prop:v5v95-owned-high}
Suppose the high-frequency proof yields annular estimates
\begin{equation}
 \|T_n\Pi_j\|_{\mathcal X}
 \le C a_n^{-h}2^{-sj}
 \label{eq:v5v95-annular-high}
\end{equation}
for the unique cutoff step \(n=n(j)\) at which annulus \(j\) is removed or
transferred.  If \(2^j\ge\Lambda_{n(j)}\), then the sum over owned high
annuli obeys the high-regime power
\begin{equation}
 \sum_{j\in\mathcal A_n}\|T_n\Pi_j\|_{\mathcal X}
 \le C'a_n^{-h}\Lambda_n^{-s}.
 \label{eq:v5v95-owned-high-bound}
\end{equation}
There is no additional factor proportional to the number of coarser
cutoffs.
\end{proposition}

\begin{proof}
Sum the geometric series \(\sum_{2^j\ge\Lambda_n}2^{-sj}\), which is at
most \(C_s\Lambda_n^{-s}\).  Lemma~\ref{lem:v5v95-one-use} ensures that no
annulus reappears at another cutoff step.
\end{proof}

This is the direct analogue of the valid high-parent accounting in the NS
audit.  It does not produce the annular estimate
\eqref{eq:v5v95-annular-high}; it prevents a false multiplicity after that
estimate is proved.

\subsection{Consequences for the rooted phase}

Let \(f_n=d\theta_n\) be the rooted phase-gradient increment at cutoff step
\(n\).  Suppose \(\|f_n\|_{\mathcal X}\le C a_n^\varepsilon\) with
\(\varepsilon>0\), as supplied by Theorem~\ref{thm:v5v95-two-regime} once
its two inputs are proved.

\begin{theorem}[Cauchy tail in the one-form covariance norm]
\label{thm:v5v95-phase-Cauchy}
Assume the massive one-form inverse has the V90 uniform bound, and the
massless grades satisfy the V91--V93 moment bounds in the same norm.  Then
for \(N>M\),
\begin{equation}
 \left\|\sum_{n=M}^{N}f_n\right\|_{H^{-1}_{\rm one}}
 \le C'\sum_{n=M}^{N}a_n^\varepsilon
 \le C''a_M^\varepsilon.
 \label{eq:v5v95-phase-Cauchy}
\end{equation}
In particular, the phase-gradient increments converge in the covariance
norm and their tail variance is \(O(a_M^{2\varepsilon})\).
\end{theorem}

\begin{proof}
The massive inverse estimate and the massless two-ended cancellation
estimates bound the covariance norm of each increment by its
\(\mathcal X\)-norm.  Apply the triangle inequality and sum the geometric
sequence.  Squaring gives the variance bound.
\end{proof}

\begin{firewall}[The high-shell estimate is the live analytic input]
\label{fw:v5v95-high-input}
Neither V94 nor Theorem~\ref{thm:v5v95-two-regime} proves
\eqref{eq:v5v95-high-bound}.  It requires a differentiated interacting
resolvent estimate on momenta where the Symanzik Taylor expansion is
invalid.  The exponent pair \((s,h)\) must be obtained without summing the
same high annulus at every cutoff.
\end{firewall}

\begin{firewall}[Second moments do not close all cumulants]
\label{fw:v5v95-all-order}
Theorem~\ref{thm:v5v95-phase-Cauchy} controls a covariance norm.  The V88
connected ratio expansion requires summable rooted cumulants at every
arity.  The audited Yang--Mills programme has proved its analogous local
tree card only through arity five; no all-order conclusion may be imported
from that fixed-arity result.
\end{firewall}

\begin{firewall}[Interacting improvement is still conditional]
\label{fw:v5v95-interacting-improvement}
The favorable smooth exponent \(b=1\) uses the differentiated
\(\beta=2\) estimate, including self-energy, vertex, measure, ghost, and
counterterm contributions.  The free improved Wilson symbol in V86 does
not establish this complete interacting bound.
\end{firewall}

\begin{theorem}[V95 conclusion]
\label{thm:v5v95-status}
A smooth-band gain \(a_n^b\Lambda_n^r\) and an owned high-shell gain
\(a_n^{-h}\Lambda_n^{-s}\) produce a summable adjacent-cutoff derivative
exactly when \(sb>rh\).  The optimal crossover yields the positive exponent
\((sb-rh)/(r+s)\).  Disjoint shell ownership prevents the nested-tail
multiplicity identified by the NS audit.  With the V90--V93 covariance
cards this gives a Cauchy rooted phase at second order.  The interacting
high-shell estimate, full improved Ward packet, and all-order connected
cumulant card remain open.
\end{theorem}

\begin{proof}
Use Theorem~\ref{thm:v5v95-two-regime},
Lemma~\ref{lem:v5v95-one-use},
Proposition~\ref{prop:v5v95-owned-high}, and
Theorem~\ref{thm:v5v95-phase-Cauchy}, subject to the three firewalls.
\end{proof}

V96 will attack the live high-shell input by deriving a differentiated
resolvent identity with one owned annulus.  The target is to expose the
available high-momentum power \(s\), the derivative loss \(h\), and the
commutator term caused by a noncommuting spectral cutoff, before any sum
over scales is taken.

\clearpage
\section{V96: one owned high shell in the differentiated resolvent}
\label{sec:v5v96}

V95 isolated the high-shell estimate as the live input.  This note derives
the exact operator identities behind it.  The useful gain is simple: a
high background momentum entering a local fermion vertex forces at least
one adjacent fermion line to be high.  A first-order Dirac resolvent then
contributes one inverse power of that momentum.  Background spectral
mixing and vertex leakage are kept explicit.

\subsection{High-output resolvent bounds}

Let \(D=D_0+V\), where \(D_0\) is translation invariant and invertible on
the regulated block, \(G=D^{-1}\), and \(G_0=D_0^{-1}\).  Let \(\Pi_K\)
be a spectral annulus or high projector commuting with \(D_0\).  Assume
\begin{equation}
 \|G_0\Pi_K\|+\|\Pi_KG_0\|\le C_0K^{-\sigma},
 \qquad \sigma>0.
 \label{eq:v5v96-free-high-resolvent}
\end{equation}

\begin{lemma}[Interacting high-output identity]
\label{lem:v5v96-high-output}
The exact identities
\begin{align}
 \Pi_KG&=\Pi_KG_0(I-VG),
 \label{eq:v5v96-left-high}\\
 G\Pi_K&=(I-GV)G_0\Pi_K
 \label{eq:v5v96-right-high}
\end{align}
imply
\begin{equation}
 \|\Pi_KG\|+\|G\Pi_K\|
 \le C_0K^{-\sigma}\bigl(2+2\|V\|\|G\|\bigr).
 \label{eq:v5v96-interacting-high}
\end{equation}
\end{lemma}

\begin{proof}
The two resolvent identities are
\(G=G_0-G_0VG\) and \(G=G_0-GVG_0\).  Multiply the first on the left and
the second on the right by \(\Pi_K\), then use commutation with \(G_0\).
Taking norms and applying \eqref{eq:v5v96-free-high-resolvent} proves the
bound.
\end{proof}

No commutator with \(V\) is needed when the high projector remains on the
external side.  If it is moved through the interacting inverse, the exact
price is as follows.

\begin{lemma}[Projector commutator identity]
\label{lem:v5v96-projector-commutator}
If \([\Pi_K,D_0]=0\), then
\begin{equation}
 [\Pi_K,G]=-G[\Pi_K,V]G.
 \label{eq:v5v96-resolvent-commutator}
\end{equation}
Consequently, for \(W=dD\),
\begin{equation}
 \Pi_K(dG)
 =-G\Pi_KWG+G[\Pi_K,V]GWG.
 \label{eq:v5v96-differentiated-commutator}
\end{equation}
\end{lemma}

\begin{proof}
Differentiate \(DG=I\) with the derivation \([\Pi_K,\cdot]\):
\([\Pi_K,D]G+D[\Pi_K,G]=0\).  Multiply by \(G\) on the left and use
\([\Pi_K,D]=[\Pi_K,V]\).  Also \(dG=-GWG\), and
\(\Pi_KG=G\Pi_K-G[\Pi_K,V]G\).  Substitution gives
\eqref{eq:v5v96-differentiated-commutator}.
\end{proof}

Equation \eqref{eq:v5v96-differentiated-commutator} shows why commuting a
shell projector through a rough background is not free.  The external-side
identity of Lemma~\ref{lem:v5v96-high-output} is preferable whenever the
shell organization permits it.

\subsection{A high-transfer vertex}

Let \(h_K\) be a background variation supported in a Fourier annulus
\(K/2\le|k|\le2K\).  Write
\begin{equation}
 W_K:=dD[h_K].
 \label{eq:v5v96-WK}
\end{equation}
Let \(P_K\) project fermion momentum to \(|p|<K/8\), and
\(Q_K=I-P_K\).

\begin{lemma}[Exact free momentum transfer]
\label{lem:v5v96-free-transfer}
For a translation-covariant local vertex linear in \(h_K\),
\begin{equation}
 P_KW_KP_K=0,
 \qquad
 W_K=Q_KW_K+P_KW_KQ_K.
 \label{eq:v5v96-transfer-decomposition}
\end{equation}
\end{lemma}

\begin{proof}
In Fourier space the vertex conserves momentum modulo the lattice
Brillouin lattice and changes fermion momentum by a momentum in the support
of \(h_K\).  Two momenta of size below \(K/8\) differ by less than \(K/4\),
so they cannot be connected to the annulus \([K/2,2K]\), provided that
annulus lies below the no-aliasing cutoff.  Thus \(P_KW_KP_K=0\).  Expand
\(I=P_K+Q_K\) on the left and right of \(W_K\), and combine the two terms
with a left \(Q_K\) to obtain the second identity.
\end{proof}

For a background-dependent link vertex or a smooth partition, introduce
the low--low leakage
\begin{equation}
 E_K:=P_KW_KP_K,
 \qquad
 W_K=Q_KW_K+P_KW_KQ_K+E_K.
 \label{eq:v5v96-leakage-decomposition}
\end{equation}

\begin{theorem}[One-high-line differentiated resolvent]
\label{thm:v5v96-one-high-line}
Assume
\begin{equation}
 \|G\|\le M,
 \qquad
 \|GQ_K\|+\|Q_KG\|\le H_K,
 \label{eq:v5v96-G-bounds}
\end{equation}
and let \(W_K\) have the decomposition
\eqref{eq:v5v96-leakage-decomposition}.  Then
\begin{equation}
 \|dG[h_K]\|
 \le 2MH_K\|W_K\|+M^2\|E_K\|.
 \label{eq:v5v96-dG-high-bound}
\end{equation}
If
\begin{equation}
 H_K\le C_Ha^{-u}K^{-\sigma},
 \quad
 \|W_K\|\le C_Wa^{-w}\|h_K\|,
 \quad
 \|E_K\|\le C_Ea^{-v}K^{-\tau}\|h_K\|,
 \label{eq:v5v96-exponent-inputs}
\end{equation}
then
\begin{equation}
 \|dG[h_K]\|
 \le C\left(a^{-u-w}K^{-\sigma}
             +a^{-v}K^{-\tau}\right)\|h_K\|.
 \label{eq:v5v96-exponent-output}
\end{equation}
For an exact translation-covariant vertex \(E_K=0\); a first-order Dirac
resolvent has \(\sigma=1\).
\end{theorem}

\begin{proof}
Use \(dG[h_K]=-GW_KG\) and
\eqref{eq:v5v96-leakage-decomposition}:
\[
 GW_KG=GQ_KW_KG+GP_KW_KQ_KG+GE_KG.
\]
The first two terms have one high resolvent next to the vertex and are
bounded by \(MH_K\|W_K\|\) each; the last is bounded by
\(M^2\|E_K\|\).  This proves \eqref{eq:v5v96-dG-high-bound}.  Substitute
\eqref{eq:v5v96-exponent-inputs}.  The last assertion follows from
Lemma~\ref{lem:v5v96-free-transfer} and the order-one symbol of \(D_0\).
\end{proof}

\begin{corollary}[Identification of the V95 exponents]
\label{cor:v5v96-V95-exponents}
Under Theorem~\ref{thm:v5v96-one-high-line}, the V95 high-shell bound holds
with
\begin{equation}
 s=\min\{\sigma,\tau\},
 \qquad
 h=\max\{u+w,v\},
 \label{eq:v5v96-sh-pair}
\end{equation}
after enlarging constants and interpreting a missing leakage term as
\(\tau=+\infty\).  In the exact translation-covariant case with a
first-order high resolvent, \(s=1\) and \(h=u+w\).
\end{corollary}

\begin{proof}
Set \(K=\Lambda_n\) in
\eqref{eq:v5v96-exponent-output}.  Each term is bounded by a common
constant times \(a^{-h}K^{-s}\).  This is
\eqref{eq:v5v95-high-bound}.
\end{proof}

\subsection{A quantitative commutator card for smooth backgrounds}

The commutator term in
\eqref{eq:v5v96-differentiated-commutator} can also be controlled when
\(V\) varies slowly relative to the shell.

\begin{proposition}[First commutator gain]
\label{prop:v5v96-first-commutator}
Let \(\Pi_K=\chi(K^{-1}|D_0|)\) be a smooth Fourier multiplier whose kernel
obeys
\begin{equation}
 \sup_x\sum_y |x-y|\,\|\Pi_K(x,y)\|\le C_\chi K^{-1}.
 \label{eq:v5v96-projector-first-moment}
\end{equation}
If \(V\) is multiplication by a matrix field with discrete Lipschitz
constant \(L_V\), then
\begin{equation}
 \|[\Pi_K,V]\|\le C_\chi K^{-1}L_V.
 \label{eq:v5v96-commutator-gain}
\end{equation}
\end{proposition}

\begin{proof}
The commutator kernel is
\(\Pi_K(x,y)(V(y)-V(x))\).  The Lipschitz bound gives
\(\|V(y)-V(x)\|\le L_V|x-y|\).  Apply the row and column Schur test and
\eqref{eq:v5v96-projector-first-moment}.
\end{proof}

Combining this with Lemma~\ref{lem:v5v96-projector-commutator} gives a
second route to one inverse power of \(K\), at the cost of the background
Lipschitz norm and two additional global resolvent factors.  It is useful
on a smooth low-field chart but is not uniform on arbitrary cutoff-scale
backgrounds.

\subsection{Owned annuli and trace-log products}

The differentiated rooted logarithm contains bounded products of
\((I+R)^{-1}\) and \(dR[h_K]\).  The inverse factor has norm at most
\((1-\kappa)^{-1}\) by V89.  Therefore a high-shell estimate for
\(dR[h_K]\) passes to the rooted phase with only that geometric factor.

\begin{proposition}[High-shell inheritance by the rooted phase]
\label{prop:v5v96-phase-inheritance}
Assume \(\|R\|\le\kappa<1\), and that on the owned annulus \(K\)
\begin{equation}
 \|dR[h_K]\|_{\mathcal X}
 \le C a^{-h}K^{-s}\|h_K\|.
 \label{eq:v5v96-dR-high}
\end{equation}
Then
\begin{equation}
 |d\theta_x[h_K]|
 \le \frac{C}{1-\kappa}a^{-h}K^{-s}\|h_K\|
 \label{eq:v5v96-phase-high}
\end{equation}
in the corresponding rooted norm.  Summing only the disjoint annuli owned
by that cutoff step gives the high bound of V95 with no incidence loss.
\end{proposition}

\begin{proof}
Use
\(d\lambda_x=\operatorname{tr}_{\rm int}((I+R)^{-1}dR)_{xx}\), the
Neumann bound on \((I+R)^{-1}\), and the rooted product inequality.  Then
apply Proposition~\ref{prop:v5v95-owned-high}.
\end{proof}

\begin{firewall}[Rough-background leakage is not controlled]
\label{fw:v5v96-rough-leakage}
For a cutoff-scale gauge background, \(P_KW_KP_K\) need not be small and
\([\Pi_K,V]\) need not gain \(K^{-1}\).  The exact free transfer lemma
does not establish the interacting estimate on such configurations.  A
small-field decomposition, a covariant spectral calculus, or a
probabilistic suppression estimate is still required.
\end{firewall}

\begin{firewall}[Global inverse factors must be uniformly bounded]
\label{fw:v5v96-global-inverse}
The constants \(M=\|G\|\) and \(1+\|V\|\|G\|\) enter the high-line bound.
A high free symbol alone does not control near-zero interacting modes.
Those modes must be separated into the low block or removed by the massive
and constraint decomposition of V90.
\end{firewall}

\begin{firewall}[One high line gives only one power]
\label{fw:v5v96-one-power}
The momentum-transfer identity forces at least one adjacent high fermion
line and therefore one first-order resolvent gain.  A second power requires
an additional cancellation, two-sided high transfer, or a differentiated
improved symbol.  It must not be inferred by squaring a one-sided bound.
\end{firewall}

\begin{theorem}[V96 conclusion]
\label{thm:v5v96-status}
A high background annulus in a translation-covariant local Dirac vertex
forces one adjacent fermion line into a high band.  The exact
differentiated resolvent identity then gives one inverse-momentum gain,
\(s=1\), modulo explicit low--low vertex leakage.  Moving a spectral
projector through an interacting inverse produces the exact commutator
\(-G[\Pi,V]G\); a Lipschitz background gives it the same one-power gain.
The estimate passes through the trace logarithm and, with one-use annular
ownership, supplies the form of the V95 high-shell input.  Uniform control
of rough-background leakage and global inverse factors remains open.
\end{theorem}

\begin{proof}
Use Lemmas \ref{lem:v5v96-high-output}--
\ref{lem:v5v96-free-transfer},
Theorem~\ref{thm:v5v96-one-high-line},
Corollary~\ref{cor:v5v96-V95-exponents}, and
Propositions \ref{prop:v5v96-first-commutator} and
\ref{prop:v5v96-phase-inheritance}, subject to the three firewalls.
\end{proof}

V97 will attack the rough-background leakage by a good-field/bad-field
decomposition tied to the positive gauge action.  The target is a
volume-uniform probability estimate strong enough to sum bad blocks while
retaining the V96 high-line gain on every good block.

\clearpage
\section{V97: good curvature blocks and probabilistic leakage}
\label{sec:v5v97}

V96 proves the high-line gain on backgrounds for which the differentiated
vertex has negligible low--low leakage, or for which the background
commutator is controlled.  A positive gauge action controls plaquette
curvature rather than link coordinates.  The correct route is therefore:
define good blocks by curvature, choose a local gauge on each good block,
and pay the complementary blocks with an explicit probability estimate.

\subsection{From small curvature to a smooth local gauge}

Let \(B_L\) be a simply connected lattice cube of side \(L\) in lattice
units, and let \(U_p\) be its plaquette holonomies.  Write
\(\operatorname{dist}(U_p,1)\) for a fixed bi-invariant distance on the
compact gauge group.

\begin{lemma}[Local lattice gauge from small plaquettes]
\label{lem:v5v97-local-gauge}
There are constants \(c_G,C_G>0\) such that if
\begin{equation}
 \max_{p\subset B_L}\operatorname{dist}(U_p,1)
 \le\delta,
 \qquad C_GL^2\delta<c_G,
 \label{eq:v5v97-small-curvature}
\end{equation}
and the block contains no nontrivial topological flux, then a gauge exists
on the concentric cube \(B_{L/2}\) in which
\begin{equation}
 \max_{e\subset B_{L/2}}\operatorname{dist}(U_e,1)
 \le C_GL\delta,
 \qquad
 \max_{e\sim e'}\operatorname{dist}(U_e,U_{e'})
 \le C_GL\delta.
 \label{eq:v5v97-link-smoothness}
\end{equation}
For an Abelian group, the conclusion holds with no small-commutator loss
once a flux sector is fixed.
\end{lemma}

\begin{proof}
Choose a spanning tree rooted at the centre and gauge all tree links to the
identity.  Every remaining link closes a loop with tree paths.  Fill that
loop by a coordinate lattice surface.  In the Abelian case its link angle
is exactly the signed sum of the plaquette angles on the surface.  A
shortest filling inside \(B_L\) uses at most \(C L\) plaquettes for the
links in the smaller cube after the standard radial tree is chosen, giving
the first bound; comparing adjacent fillings gives the second.

For a compact non-Abelian group, order the plaquettes in the filling.
Products of their logarithms differ from the sum by commutators bounded by
\(C L^2\delta^2\).  Condition \eqref{eq:v5v97-small-curvature} keeps the
product in one exponential chart and absorbs this error into
\(C_GL\delta\).  A nontrivial flux would obstruct a single-valued small
logarithm, which is why it is excluded.
\end{proof}

\begin{corollary}[Good-block V96 commutator]
\label{cor:v5v97-good-commutator}
On a block satisfying \eqref{eq:v5v97-small-curvature}, every local link
coefficient \(V(U)\) with bounded first derivative has a gauge in which
its discrete Lipschitz constant is at most \(C L\delta\).  Hence the V96
smooth spectral commutator obeys
\begin{equation}
 \|[\Pi_K,V]\|\le C K^{-1}L\delta.
 \label{eq:v5v97-good-commutator}
\end{equation}
\end{corollary}

\begin{proof}
Apply the mean-value estimate to
\eqref{eq:v5v97-link-smoothness}, then use
Proposition~\ref{prop:v5v96-first-commutator}.
\end{proof}

The gauge used here is local.  Gauge-invariant rooted observables agree on
overlaps, while gauge-fixed covariance kernels require a separately
compatible partition or transition estimate.

\subsection{A moment-to-high-shell interpolation}

Let \(Z_B\ge0\) measure the roughness of a block after its admissible local
gauge is chosen.  For example, on the good chart one may take
\(Z_B=L^{-1}\operatorname{Lip}V\).  Let \(X_K\) denote the high-shell
leakage contribution to the differentiated rooted observable on that
block.

\begin{theorem}[Good/bad high-shell interpolation]
\label{thm:v5v97-good-bad-interpolation}
Assume for some \(m>0\), \(p>1\), and nonnegative cutoff losses
\(\kappa,h,\chi\) that
\begin{align}
 \|Z_B\|_{L^m}&\le C_Za^{-\kappa},
 \label{eq:v5v97-Z-moment}\\
 |X_K|&\le C_Ga^{-h}K^{-1}Z_B
       \quad\text{on }\{Z_B\le a^{-\kappa}K^\xi\},
 \label{eq:v5v97-good-X}\\
 \|X_K\|_{L^p}&\le C_Ba^{-\chi}
       \quad\text{without restriction}.
 \label{eq:v5v97-bad-envelope}
\end{align}
Then, with \(p'=p/(p-1)\),
\begin{equation}
 \mathbb E|X_K|
 \le C\left(
 a^{-h-\kappa}K^{-1+\xi}
 +a^{-\chi}K^{-m\xi/p'}
 \right).
 \label{eq:v5v97-interpolated-high}
\end{equation}
The balanced choice
\begin{equation}
 \xi_*:=\frac1{1+m/p'},
 \qquad
 s_*:=\frac{m/p'}{1+m/p'}>0
 \label{eq:v5v97-balanced-s}
\end{equation}
gives the V95 high-shell form
\begin{equation}
 \mathbb E|X_K|\le C a^{-h_*}K^{-s_*},
 \qquad h_*:=\max\{h+\kappa,\chi\}.
 \label{eq:v5v97-effective-high}
\end{equation}
\end{theorem}

\begin{proof}
On the good event, \eqref{eq:v5v97-good-X} and the threshold give the first
term in \eqref{eq:v5v97-interpolated-high}.  Markov's inequality and
\eqref{eq:v5v97-Z-moment} give
\begin{equation}
 \mathbb P\{Z_B>a^{-\kappa}K^\xi\}
 \le C_Z^mK^{-m\xi}.
 \label{eq:v5v97-bad-probability}
\end{equation}
H\"older's inequality with \eqref{eq:v5v97-bad-envelope} bounds the bad
expectation by
\(C a^{-\chi}K^{-m\xi/p'}\).  Equating
\(1-\xi=m\xi/p'\) gives \eqref{eq:v5v97-balanced-s}; absorbing the larger
cutoff loss gives \eqref{eq:v5v97-effective-high}.
\end{proof}

\begin{corollary}[Compatibility with the V95 crossover]
\label{cor:v5v97-V95-compatibility}
Let \(b,r\) be the V95 smooth-band exponents.  Under
Theorem~\ref{thm:v5v97-good-bad-interpolation}, the adjacent-cutoff sum is
finite if
\begin{equation}
 s_*b>rh_*.
 \label{eq:v5v97-feasibility}
\end{equation}
Its balanced cutoff exponent is
\begin{equation}
 \varepsilon=\frac{s_*b-rh_*}{r+s_*}>0.
 \label{eq:v5v97-resulting-epsilon}
\end{equation}
\end{corollary}

\begin{proof}
Insert \((s,h)=(s_*,h_*)\) into
Theorem~\ref{thm:v5v95-two-regime}.
\end{proof}

This criterion shows exactly what probability has to buy.  A finite
roughness moment always gives some positive high-momentum power, but any
cutoff loss in that moment or in the bad-field envelope must still be
beaten by the improved smooth exponent.

\subsection{What the positive gauge action can certify}

Let the block gauge action be
\begin{equation}
 S_B(U)=\sum_{p\subset B}s(U_p),
 \qquad
 s(U)\ge c_0\operatorname{dist}(U,1)^2
 \quad\text{near }1.
 \label{eq:v5v97-block-action}
\end{equation}

\begin{lemma}[Action exponential moment implies a bad-block tail]
\label{lem:v5v97-action-tail}
Suppose for some \(t>0\) and \(C_t<\infty\), uniformly in volume and in
the admissible conditioning outside \(B\),
\begin{equation}
 \mathbb E\left[e^{t\beta S_B}\mid\mathcal F_{B^c}\right]
 \le e^{C_t|B|}.
 \label{eq:v5v97-action-exp-moment}
\end{equation}
Then for \(u>0\),
\begin{equation}
 \mathbb P\{S_B\ge u|B|\mid\mathcal F_{B^c}\}
 \le \exp\bigl[-(t\beta u-C_t)|B|\bigr].
 \label{eq:v5v97-action-Chernoff}
\end{equation}
If a block is declared bad whenever a fixed fraction of its plaquettes has
curvature at least \(\delta\), its probability obeys the same form with
\(u=c\delta^2\).
\end{lemma}

\begin{proof}
Apply conditional Markov inequality to \(e^{t\beta S_B}\).  On the stated
bad event, \eqref{eq:v5v97-block-action} gives
\(S_B\ge c\delta^2|B|\).
\end{proof}

This lemma is a certificate, not a derivation of
\eqref{eq:v5v97-action-exp-moment}.  Fermion determinants, Higgs couplings,
gauge fixing, and gravity can alter the conditional density.  A lower
bound on the bare plaquette action does not by itself bound the normalized
conditional exponential moment because its partition-function denominator
must also be controlled.

\subsection{Bad polymers and all-order use}

For the all-order rooted expansion, one-point rarity is insufficient.
Let \(\eta_B\) be the bad-block indicator.

\begin{proposition}[A sufficient bad-polymer domination]
\label{prop:v5v97-bad-polymer}
Assume that for every finite family \(\mathcal S\) of mutually separated
blocks,
\begin{equation}
 \mathbb E\prod_{B\in\mathcal S}\eta_B\le p_0^{|\mathcal S|}.
 \label{eq:v5v97-product-domination}
\end{equation}
If the number of connected block animals of size \(m\) containing a fixed
root is at most \(C_{\rm an}^m\), then
\begin{equation}
 \sum_{\substack{\mathcal A\ni B_0\\\mathcal A\ {
m connected}}}
 w^{|\mathcal A|}\,
 \mathbb E\prod_{B\in\mathcal A}\eta_B
 \le
 \sum_{m\ge1}(C_{\rm an}wp_0)^m<\infty
 \label{eq:v5v97-polymer-sum}
\end{equation}
whenever \(C_{\rm an}wp_0<1\), after enlarging blocks by a fixed coloring
factor to enforce separation.
\end{proposition}

\begin{proof}
Color the bounded-degree block graph with finitely many colors.  Every
animal contains a color class of proportional size; applying
\eqref{eq:v5v97-product-domination} to that class only changes \(p_0\) by a
fixed fractional power, which can be renamed.  Count rooted animals and
sum the resulting geometric series.
\end{proof}

The proposition supplies the type of arity-uniform estimate missing from
the V95 Yang--Mills audit, but its product-domination hypothesis has not
been proved for the coupled SM--Einstein measure.

\begin{firewall}[Curvature does not globally bound links]
\label{fw:v5v97-curvature-links}
Small plaquettes imply smooth link representatives only locally, after a
gauge choice and within a controlled flux sector.  Holonomy, topology, and
overlap transitions prevent promotion of
Lemma~\ref{lem:v5v97-local-gauge} to a global link bound without further
work.
\end{firewall}

\begin{firewall}[A bare action is not a probability estimate]
\label{fw:v5v97-action-probability}
The positivity of \(S_B\) does not prove
\eqref{eq:v5v97-action-exp-moment} for the normalized interacting measure.
The conditional partition ratio and the modulus of the fermion determinant
must be bounded uniformly.  The residual complex phase remains in the V88
rooted cancellation problem.
\end{firewall}

\begin{firewall}[One-point moments are not polymer domination]
\label{fw:v5v97-moment-versus-polymer}
The moment hypothesis \eqref{eq:v5v97-Z-moment} suffices for the averaged
two-point high-shell estimate.  It does not imply
\eqref{eq:v5v97-product-domination}; long-range correlations can make many
bad blocks occur together.  All-order cumulants require the stronger
conditional or cluster estimate.
\end{firewall}

\begin{theorem}[V97 conclusion]
\label{thm:v5v97-status}
Small plaquette curvature on a topologically trivial block yields a local
gauge with a quantitative link Lipschitz bound, and hence the V96
one-power commutator gain.  If block roughness has an \(L^m\) bound and the
unrestricted leakage has an \(L^p\) envelope, a good/bad split produces
the explicit high-shell exponent
\(s_*=(m/p')/(1+m/p')\).  It closes the V95 cutoff sum when
\(s_*b>rh_*\).  An all-order bad-polymer sum would follow from uniform
product domination.  The required interacting exponential moment,
topological patching, and product domination remain open.
\end{theorem}

\begin{proof}
Use Lemma~\ref{lem:v5v97-local-gauge},
Theorem~\ref{thm:v5v97-good-bad-interpolation},
Corollary~\ref{cor:v5v97-V95-compatibility},
Lemma~\ref{lem:v5v97-action-tail}, and
Proposition~\ref{prop:v5v97-bad-polymer}, subject to the three firewalls.
\end{proof}

V98 will go upstream to the normalized conditional block measure.  It will
derive a denominator-free comparison inequality for positive reference
measures and identify the precise determinant-modulus bound needed to
transfer the plaquette exponential moment through the Higgs, PV, and
fermion factors.

\clearpage
\section{V98: normalized block comparison without a raw denominator}
\label{sec:v5v98}

The V97 bad-block estimate must hold under the normalized interacting
conditional measure.  Writing an unnormalized weight and then assuming its
partition-function denominator is bounded below would recreate the problem
removed from the complex phase in V88.  For a positive reference measure,
Jensen's inequality gives a cleaner comparison in terms of the centered
logarithm of the reweighting factor.

\subsection{The centered-log comparison}

Let \(\nu_B\) be a probability measure on the fields of a block \(B\), with
the exterior configuration fixed.  Let \(L_B\) be real and define
\begin{equation}
 d\mu_B=\frac{e^{L_B}}{\mathbb E_{\nu_B}e^{L_B}}\,d\nu_B.
 \label{eq:v5v98-positive-reweighting}
\end{equation}

\begin{theorem}[Denominator-free density comparison]
\label{thm:v5v98-centered-log}
For every \(q>1\),
\begin{equation}
 \left\|\frac{d\mu_B}{d\nu_B}\right\|_{L^q(\nu_B)}
 \le
 \left(
 \mathbb E_{\nu_B}e^{q(L_B-\mathbb E_{\nu_B}L_B)}
 \right)^{1/q}.
 \label{eq:v5v98-density-Lq}
\end{equation}
Consequently, if
\begin{equation}
 \log\mathbb E_{\nu_B}
 e^{q(L_B-\mathbb E_{\nu_B}L_B)}
 \le q C_L|B|,
 \label{eq:v5v98-centered-mgf}
\end{equation}
then, with \(p=q/(q-1)\), every nonnegative \(F\) satisfies
\begin{equation}
 \mathbb E_{\mu_B}F
 \le e^{C_L|B|}\|F\|_{L^p(\nu_B)},
 \label{eq:v5v98-positive-comparison}
\end{equation}
and every event \(E\) satisfies
\begin{equation}
 \mu_B(E)\le e^{C_L|B|}\nu_B(E)^{1/p}.
 \label{eq:v5v98-event-comparison}
\end{equation}
\end{theorem}

\begin{proof}
Jensen's inequality gives
\(\mathbb E_{\nu_B}e^{L_B}\ge
e^{\mathbb E_{\nu_B}L_B}\).  Therefore
\[
 \int\left(\frac{e^{L_B}}
 {\mathbb E_{\nu_B}e^{L_B}}\right)^q d\nu_B
 \le\mathbb E_{\nu_B}e^{q(L_B-\mathbb E_{\nu_B}L_B)}.
\]
This proves \eqref{eq:v5v98-density-Lq}.  Under
\eqref{eq:v5v98-centered-mgf}, the density norm is at most
\(e^{C_L|B|}\).  H\"older's inequality proves
\eqref{eq:v5v98-positive-comparison}; take \(F={\bf1}_E\) for
\eqref{eq:v5v98-event-comparison}.
\end{proof}

The normalization denominator has disappeared into a centered positive
moment.  This is useful only because \(L_B\) is real.  The complex
determinant phase remains in the rooted expansion of V88--V96.

\begin{corollary}[Oscillation comparison]
\label{cor:v5v98-oscillation}
If
\begin{equation}
 \operatorname*{ess\,sup}_{\nu_B}L_B-
 \operatorname*{ess\,inf}_{\nu_B}L_B\le\Omega_B,
 \label{eq:v5v98-log-oscillation}
\end{equation}
then
\begin{equation}
 \mathbb E_{\mu_B}F\le e^{\Omega_B}
 \mathbb E_{\nu_B}F
 \label{eq:v5v98-oscillation-comparison}
\end{equation}
for every nonnegative \(F\).
\end{corollary}

\begin{proof}
Pointwise,
\(e^{L_B}\le e^{\operatorname*{ess\,sup}L_B}\), while
\(\mathbb E e^{L_B}\ge e^{\operatorname*{ess\,inf}L_B}\).  Their ratio is
at most \(e^{\Omega_B}\).
\end{proof}

The centered-moment theorem can be much sharper than the oscillation bound
when Higgs or auxiliary fields are unbounded.

\subsection{Transfer of the plaquette tail}

Assume the positive gauge reference measure has a conditional block
estimate
\begin{equation}
 \nu_B\{S_B\ge u|B|\}
 \le \exp[-I_\beta(u)|B|]
 \label{eq:v5v98-reference-tail}
\end{equation}
uniformly in admissible exterior data.

\begin{proposition}[Interacting bad-block tail]
\label{prop:v5v98-transferred-tail}
Under \eqref{eq:v5v98-centered-mgf} and
\eqref{eq:v5v98-reference-tail},
\begin{equation}
 \mu_B\{S_B\ge u|B|\}
 \le
 \exp\left[-\left(\frac{I_\beta(u)}p-C_L\right)|B|\right].
 \label{eq:v5v98-interacting-tail}
\end{equation}
Thus the interacting conditional block is exponentially rare whenever
\(I_\beta(u)>pC_L\).
\end{proposition}

\begin{proof}
Apply \eqref{eq:v5v98-event-comparison} to the event in
\eqref{eq:v5v98-reference-tail}.
\end{proof}

There is an equivalent moment formulation.  If for some \(t>0\)
\begin{equation}
 \mathbb E_{\nu_B}e^{pt\beta S_B}le e^{C_G|B|},
 \label{eq:v5v98-reference-action-mgf}
\end{equation}
then Theorem~\ref{thm:v5v98-centered-log} gives
\begin{equation}
 \mathbb E_{\mu_B}e^{t\beta S_B}
 \le e^{(C_L+C_G/p)|B|},
 \label{eq:v5v98-transferred-action-mgf}
\end{equation}
which is the V97 conditional exponential-moment certificate.

\subsection{From a conditional tail to product domination}

\begin{theorem}[Sequential conditional domination]
\label{thm:v5v98-sequential-domination}
Let \((\eta_B)\) be bad-block indicators.  Suppose that for every block and
almost every configuration of all other blocks,
\begin{equation}
 \mathbb E[\eta_B\mid\mathcal F_{B^c}]\le p_0.
 \label{eq:v5v98-uniform-conditional-bad}
\end{equation}
Then for every finite family \(\mathcal S\),
\begin{equation}
 \mathbb E\prod_{B\in\mathcal S}\eta_B\le p_0^{|\mathcal S|}.
 \label{eq:v5v98-product-domination}
\end{equation}
Consequently the V97 bad-polymer sum converges when
\(C_{\rm an}wp_0<1\).
\end{theorem}

\begin{proof}
Choose one block \(B_1\in\mathcal S\).  The product of all other indicators
is \(\mathcal F_{B_1^c}\)-measurable, so the tower property and
\eqref{eq:v5v98-uniform-conditional-bad} give
\[
 \mathbb E\prod_{B\in\mathcal S}\eta_B
 \le p_0\mathbb E\prod_{B\in\mathcal S\setminus\{B_1\}}\eta_B.
\]
Iterate.  The last assertion is
Proposition~\ref{prop:v5v97-bad-polymer}.
\end{proof}

This theorem is stronger than needed because conditioning on every other
block may be incompatible with boundary plaquettes.  A finite-color or
buffered version suffices, with a fixed change in \(p_0\), but the same
uniform conditional comparison remains the central requirement.

\subsection{The determinant-modulus input}

Put the modulus of the regulated fermionic and PV factors into the positive
reference weight:
\begin{equation}
 L_B=L_B^{\rm Higgs}+L_B^{\rm PV}
     +\log|\det D_B|+L_B^{\rm ct}.
 \label{eq:v5v98-L-decomposition}
\end{equation}
The following finite-dimensional estimate identifies a sufficient
fermionic condition.

\begin{lemma}[Log-determinant Lipschitz bound]
\label{lem:v5v98-logdet-Lipschitz}
Let \(D(t)\), \(0\le t\le1\), be a differentiable path of complex matrices
such that
\begin{equation}
 \sigma_{\min}(D(t))\ge m_D>0.
 \label{eq:v5v98-singular-gap}
\end{equation}
Then
\begin{equation}
 \left|\log|\det D(1)|-\log|\det D(0)|\right|
 \le m_D^{-1}\int_0^1\|D'(t)\|_1\,dt,
 \label{eq:v5v98-logdet-path}
\end{equation}
where \(\|\cdot\|_1\) is the trace norm.
\end{lemma}

\begin{proof}
Invertibility and Jacobi's formula give
\[
 \frac d{dt}\log|\det D(t)|
 =\operatorname{Re}\operatorname{Tr}(D(t)^{-1}D'(t)).
\]
Trace duality bounds the absolute value by
\(\|D(t)^{-1}\|\|D'(t)\|_1\le m_D^{-1}\|D'(t)\|_1\).
Integrate.
\end{proof}

\begin{corollary}[Extensive determinant oscillation]
\label{cor:v5v98-determinant-oscillation}
Suppose any two admissible block gauge--Higgs configurations can be joined
by a path satisfying \eqref{eq:v5v98-singular-gap}, and
\begin{equation}
 \int_0^1\|D'_B(t)\|_1dt\le C_Dm_D|B|.
 \label{eq:v5v98-trace-path-extensive}
\end{equation}
Then \(\log|\det D_B|\) has oscillation at most \(C_D|B|\).  If the other
terms in \eqref{eq:v5v98-L-decomposition} have either extensive
oscillation or the centered moment \eqref{eq:v5v98-centered-mgf}, the full
positive reweighting satisfies the comparison theorem with an extensive
constant \(C_L\).
\end{corollary}

\begin{proof}
Lemma~\ref{lem:v5v98-logdet-Lipschitz} proves the determinant statement.
Oscillation bounds add.  Centered exponential moments combine by H\"older's
inequality, with a fixed enlargement of their exponents and constants.
\end{proof}

For high PV modes, V68 supplies the type of singular gap required here.
For light chiral fermions, a blockwise gap can fail because of physical
near-zero modes.  Those modes must be retained rather than hidden inside
an oscillation constant.

\subsection{Positivity and zero-mode firewalls}

\begin{firewall}[The phase is not part of \(L_B\)]
\label{fw:v5v98-phase-excluded}
The centered-log proof uses order and Jensen's inequality.  It applies to
\(|\det D|\), positive PV factors, and a real positive Higgs/reference
weight.  It does not apply to the complex phase of a chiral determinant;
that phase remains governed by the connected rooted cancellation problem
of V88--V97.
\end{firewall}

\begin{firewall}[Light singular values obstruct oscillation]
\label{fw:v5v98-light-zero-modes}
If \(\det D_B\) vanishes, \(\log|\det D_B|=-\infty\) and the oscillation
card fails.  An integrable determinant zero can still be compatible with
the centered moment theorem, but that requires a direct small-singular-
value distribution estimate.  No uniform light-fermion gap is asserted.
\end{firewall}

\begin{firewall}[The Euclidean conformal mode is not a positive weight]
\label{fw:v5v98-conformal-mode}
The unrotated Euclidean gravitational conformal direction is not a
positive Gibbs factor, as proved earlier in the programme.  It cannot be
inserted into \(L_B\) in \eqref{eq:v5v98-positive-reweighting}.  The block
comparison applies only after a legitimate positive physical/constraint
chart or a separately justified contour construction.
\end{firewall}

\begin{firewall}[The reference tail remains an input]
\label{fw:v5v98-reference-tail}
The comparison theorem transfers \eqref{eq:v5v98-reference-tail}; it does
not prove it.  Boundary frustration, flux sectors, and gauge constraints
must be included in a uniform conditional plaquette estimate for the
chosen positive reference measure.
\end{firewall}

\begin{theorem}[V98 conclusion]
\label{thm:v5v98-status}
A positive interacting block measure can be compared with a gauge
reference measure without a raw partition-function lower bound: a centered
exponential moment of the real log weight controls the normalized density
in \(L^q\).  It transfers exponential plaquette tails, and a uniform
conditional tail implies the all-family product domination required by
V97.  A uniform singular-value path gives a sufficient extensive bound for
the determinant modulus.  Light determinant zeros, the complex phase, the
Euclidean conformal mode, and the reference conditional plaquette tail
remain outside this certificate.
\end{theorem}

\begin{proof}
Use Theorem~\ref{thm:v5v98-centered-log},
Proposition~\ref{prop:v5v98-transferred-tail},
Theorem~\ref{thm:v5v98-sequential-domination},
Lemma~\ref{lem:v5v98-logdet-Lipschitz}, and
Corollary~\ref{cor:v5v98-determinant-oscillation}, subject to the four
firewalls.
\end{proof}

V99 will isolate light fermionic modes by an exact singular-value split.
The target is a Schur/Feshbach determinant factorization in which the high
block satisfies the V98 modulus comparison, while the finite low block is
retained in the observable and its determinant zeros are never divided
out.

\clearpage
\section{V99: exact retention of the light fermion block}
\label{sec:v5v99}

V98 identified a uniform singular-value gap as a sufficient condition for
controlling the determinant modulus.  Physical light fermions do not have
such a gap.  The appropriate operation is not to regularize their zeros
away, but to integrate only a complementary high block whose gap is
proved.  The low determinant is then retained as a finite Grassmann
factor.

\subsection{The Feshbach determinant identity}

Let the finite-cutoff fermion space split as
\begin{equation}
 \mathcal H_F=\mathcal H_L\oplus\mathcal H_H,
 \qquad P+Q=I,
 \label{eq:v5v99-fermion-split}
\end{equation}
and write
\begin{equation}
 D=
 \begin{pmatrix}A&B\\ C&E\end{pmatrix}
 =
 \begin{pmatrix}PDP&PDQ\\ QDP&QDQ\end{pmatrix}.
 \label{eq:v5v99-Dirac-blocks}
\end{equation}
Assume only that \(E\) is invertible, and define
\begin{equation}
 S:=A-BE^{-1}C.
 \label{eq:v5v99-low-Schur}
\end{equation}

\begin{theorem}[Exact high/low determinant split]
\label{thm:v5v99-determinant-split}
Under the preceding hypothesis,
\begin{equation}
 \det D=\det E\,\det S.
 \label{eq:v5v99-determinant-factorization}
\end{equation}
Moreover, \(D\) is singular if and only if \(S\) is singular.  Thus every
determinant zero is retained in the low factor as long as the high gap does
not close.
\end{theorem}

\begin{proof}
Multiply \(D\) on the left by the determinant-one block matrix
\[
 \begin{pmatrix}I&-BE^{-1}\\0&I\end{pmatrix}.
\]
The result is lower triangular with diagonal blocks \(S,E\), proving
\eqref{eq:v5v99-determinant-factorization}.  Since \(E\) is invertible,
the product vanishes exactly when \(\det S=0\).
\end{proof}

\begin{proposition}[Exact low-mode equation]
\label{prop:v5v99-low-equation}
The equation \(D(u,v)=0\) is equivalent to
\begin{equation}
 Su=0,
 \qquad
 v=-E^{-1}Cu.
 \label{eq:v5v99-kernel-lift}
\end{equation}
Hence \(u\mapsto(u,-E^{-1}Cu)\) is an isomorphism from
\(\ker S\) to \(\ker D\).
\end{proposition}

\begin{proof}
The high component gives \(Cu+Ev=0\), which has the unique solution in
\eqref{eq:v5v99-kernel-lift}.  Substitution into the low component gives
\(Su=0\).  The converse follows by reversing the calculation.
\end{proof}

\subsection{Grassmann integration without dividing by the low determinant}

Write the Grassmann variables as \(\psi=(\psi_L,\psi_H)\) and similarly for
\(\bar\psi\).

\begin{theorem}[Exact high Grassmann integration]
\label{thm:v5v99-Grassmann-Feshbach}
With a consistent finite-dimensional Berezin convention,
\begin{equation}
 \int d\bar\psi_Hd\psi_H\,
 e^{-\bar\psi D\psi}
 =\det E\,
 e^{-\bar\psi_LS\psi_L}.
 \label{eq:v5v99-Grassmann-split}
\end{equation}
The right-hand side is a finite polynomial in the entries of \(S\) and in
the retained Grassmann generators even when \(S\) is singular.
\end{theorem}

\begin{proof}
Complete the Grassmann square by the unit-Berezinian translations
\begin{equation}
 \psi_H'=\psi_H+E^{-1}C\psi_L,
 \qquad
 \bar\psi_H'=\bar\psi_H+\bar\psi_LBE^{-1}.
 \label{eq:v5v99-Grassmann-shift}
\end{equation}
Then
\[
 \bar\psi D\psi
 =\bar\psi_H'E\psi_H'+\bar\psi_LS\psi_L.
\]
The Gaussian Berezin integral over the primed high variables is
\(\det E\), proving \eqref{eq:v5v99-Grassmann-split}.  Exponentials in
finitely many Grassmann generators terminate, so no inverse of \(S\) is
used.
\end{proof}

This identity implements the V88 nilpotent inverse principle at the
spectral split: the light factor remains algebraic until its external legs
or observables are specified.

\subsection{Derivative and locality of the low effective operator}

\begin{lemma}[Differentiated Schur complement]
\label{lem:v5v99-dS}
For every parameter derivative \(d\),
\begin{equation}
 dS=dA-(dB)E^{-1}C-BE^{-1}(dC)
     +BE^{-1}(dE)E^{-1}C.
 \label{eq:v5v99-dS-formula}
\end{equation}
If \(\|E^{-1}\|\le m_H^{-1}\), then
\begin{align}
 \|dS\|\le{}&\|dA\|
 +m_H^{-1}\bigl(\|dB\|\|C\|+\|B\|\|dC\|\bigr)
 \notag\\
 &+m_H^{-2}\|B\|\|dE\|\|C\|.
 \label{eq:v5v99-dS-bound}
\end{align}
No inverse of \(S\) appears.
\end{lemma}

\begin{proof}
Differentiate \eqref{eq:v5v99-low-Schur} and use
\(d(E^{-1})=-E^{-1}(dE)E^{-1}\).  The norm bound follows by
submultiplicativity.
\end{proof}

\begin{proposition}[Locality after high-block elimination]
\label{prop:v5v99-Schur-locality}
Suppose \(A,B,C,E\) have finite range or uniform exponential locality, and
the high block \(E\) satisfies the gap and weighted locality hypotheses of
V90.  Then \(E^{-1}\) and the Schur complement \(S\) are exponentially
local, with volume-uniform weighted row and column bounds.  The same holds
for \(dS\) if the differentiated blocks obey those bounds.
\end{proposition}

\begin{proof}
The V90 Combes--Thomas estimate gives the required bound for \(E^{-1}\).
The weighted Schur norm is a Banach algebra.  Apply it to
\(BE^{-1}C\) and to every product in
\eqref{eq:v5v99-dS-formula}.
\end{proof}

The retained light theory is therefore non-ultralocal but exponentially
local.  Its exact zero modes and long-distance propagation reside in
\(S\), rather than in an uncontrolled inverse hidden in the high factor.

\subsection{The high modulus in the positive comparison measure}

Assume along every admissible block interpolation
\begin{equation}
 \sigma_{\min}(E(t))\ge m_H>0.
 \label{eq:v5v99-high-gap}
\end{equation}

\begin{theorem}[High determinant comparison card]
\label{thm:v5v99-high-modulus-card}
If
\begin{equation}
 \int_0^1\|E'(t)\|_1dt\le C_E m_H|B|,
 \label{eq:v5v99-high-trace-path}
\end{equation}
then \(\log|\det E|\) has oscillation at most \(C_E|B|\) and may be
included in the positive log weight \(L_B\) of V98.  The remaining factor
\(e^{-\bar\psi_LS\psi_L}\) is retained in the finite Grassmann algebra and
is not used in the normalization denominator.
\end{theorem}

\begin{proof}
Apply Lemma~\ref{lem:v5v98-logdet-Lipschitz} to \(E(t)\).  Then use
Corollary~\ref{cor:v5v98-oscillation}.  The last statement is exactly the
factorization \eqref{eq:v5v99-Grassmann-split}.
\end{proof}

\begin{corollary}[All determinant zeros remain visible]
\label{cor:v5v99-zeros-visible}
Under \eqref{eq:v5v99-high-gap}, no zero or sign/phase singularity of
\(\det D\) is absorbed into the positive high-block comparison density.
It occurs in \(\det S\), equivalently in the terminating low Grassmann
factor.
\end{corollary}

\begin{proof}
Use Theorem~\ref{thm:v5v99-determinant-split} and the fact that
\(\det E\ne0\).
\end{proof}

\subsection{Gauge covariance of the split}

\begin{proposition}[Covariant projection card]
\label{prop:v5v99-covariant-projection}
Suppose the projectors transform under a background gauge transformation
as
\begin{equation}
 P(A^g)=U_gP(A)U_g^{-1},
 \qquad Q(A^g)=U_gQ(A)U_g^{-1},
 \label{eq:v5v99-projector-covariance}
\end{equation}
and \(D(A^g)=U_gD(A)U_g^{-1}\).  Then \(E\) and \(S\) transform by
conjugation on their respective subspaces, and
\(\det E\det S\) is gauge invariant.
\end{proposition}

\begin{proof}
Insert \eqref{eq:v5v99-projector-covariance} into each block of
\eqref{eq:v5v99-Dirac-blocks}.  Inversion preserves conjugation, so the
same is true of \(S=A-BE^{-1}C\).  Determinants are invariant under the
induced changes of basis.
\end{proof}

\begin{firewall}[A moving spectral projector has connection terms]
\label{fw:v5v99-moving-projector}
If \(P(A)\) is defined by a field-dependent spectral threshold, then
differentation produces \(dP\) terms in every block.  At an eigenvalue
crossing the rank can change and \(dP\) need not exist.  The formula
\eqref{eq:v5v99-dS-formula} applies directly only to a fixed smooth bundle
of low subspaces; a moving split requires its Berry connection and a
separate crossing atlas.
\end{firewall}

\begin{firewall}[Low rank need not be continuum-uniform]
\label{fw:v5v99-low-rank}
The retained Grassmann algebra is finite at finite cutoff and volume, but
its number of generators generally grows with volume and as the cutoff is
removed.  Nilpotence at one cutoff does not give an all-order,
volume-uniform cumulant bound.  A local multiscale ownership rule and tree
estimate are still required.
\end{firewall}

\begin{firewall}[The high gap must be stable along paths]
\label{fw:v5v99-path-gap}
A gap at one configuration does not imply
\eqref{eq:v5v99-high-gap} along all comparison paths.  Spectral flow can
move a light state through the chosen threshold.  Such a state must be
transferred to the retained block before applying the V98 oscillation
estimate.
\end{firewall}

\begin{firewall}[Chiral phase separation remains necessary]
\label{fw:v5v99-chiral-phase}
The positive factor is \(|\det E|\).  Its phase, the low Schur phase, and
the measure counterterms must satisfy the complete regulated Ward and
anomaly-cancellation identities of V92.  The high/low determinant identity
does not prove phase cancellation.
\end{firewall}

\begin{theorem}[V99 conclusion]
\label{thm:v5v99-status}
An exact Feshbach split integrates only a uniformly invertible high fermion
block.  Its determinant modulus satisfies the V98 positive comparison card
under an extensive trace-path estimate.  The low Schur complement remains
an exponentially local finite Grassmann factor; every light determinant
zero is retained there and no inverse of the low operator is introduced.
Gauge covariance holds for a covariant projector bundle.  Moving spectral
crossings, growth of the retained rank, the high path gap, and the chiral
phase remain open.
\end{theorem}

\begin{proof}
Use Theorems \ref{thm:v5v99-determinant-split},
\ref{thm:v5v99-Grassmann-Feshbach}, and
\ref{thm:v5v99-high-modulus-card}, Lemma~\ref{lem:v5v99-dS}, and
Propositions \ref{prop:v5v99-Schur-locality} and
\ref{prop:v5v99-covariant-projection}, subject to the four firewalls.
\end{proof}

V100 is the next compulsory YM/NS audit and the fifth-series boundary.  It
will audit the companion endpoints, construct a smooth covariant projector
atlas with explicit crossing transfers, and then freeze this volume as
\texttt{\_f5}.  The continuation will restart at V1 with context, a compact
ledger of the major proved results, and the live SM--Einstein bottlenecks.

\clearpage
\section{V100: covariant spectral charts and exact crossing transfer}
\label{sec:v5v100}

\subsection{Compulsory companion audit}

The active Yang--Mills continuation has restarted after its fourth frozen
volume.  The file inspected here was
\texttt{Renewal\_Yang\_Mills\_Mass\_Gap\_Research\_Notes\_V3.tex}, with
\(51\,011\) bytes and SHA--256
\begin{center}
\texttt{0194BE468080C1D7960762F1AA0E125985CB9B7B7BAC4C6544B15C22D454B566}.
\end{center}
It proves a connected Gaussian BKAR forest formula, dimension-free complete
derivative rows for representation exponentials, and an all-order Gaussian
tangent tree card with exponential constants.  It thereby removes the
fixed-arity obstruction on the Gaussian tangent model.  The remaining
Yang--Mills acquisition is a uniform all-derivative transport from that
Gaussian law to the Wilson law; its stated UWTC is not yet proved.

The Navier--Stokes file remained
\texttt{Renewal\_NS\_Stability\_Research\_Notes\_V31.tex}, with
\(887\,537\) bytes and SHA--256
\begin{center}
\texttt{F1121B62238AD5491BB7CF03635EA9934942EDF573ED8FAAA9EE79E1D38A6696}.
\end{center}
It still requires a stopping-time-correlated signed formation estimate and
continues to forbid repeated charging of a remote shell.

The YM result supplies the correct architecture for the V88 all-order
rooted phase problem: first prove a forest formula for a positive Gaussian
massive reference, then transport or interpolate to the true local block
law with constants exponential in arity.  It does not prove that transport
for the gauge--Higgs--fermion--gravity measure.  The present note first
finishes the spectral crossing problem left by V99.

\subsection{Gauge-covariant singular subspaces}

Let \(D(A):\mathcal H_R\to\mathcal H_L\) be the finite-cutoff chiral Dirac
map and define
\begin{equation}
 H_R=D^*D,
 \qquad
 H_L=DD^*.
 \label{eq:v5v100-singular-Hamiltonians}
\end{equation}
Fix a threshold \(\Lambda>0\).  On a field-space chart where neither
\(H_R\) nor \(H_L\) has spectrum in
\((\Lambda^2-\delta,\Lambda^2+\delta)\), let
\begin{equation}
 P_R={\bf1}_{[0,\Lambda^2)}(H_R),
 \qquad
 P_L={\bf1}_{[0,\Lambda^2)}(H_L),
 \qquad Q_{R,L}=I-P_{R,L}.
 \label{eq:v5v100-singular-projectors}
\end{equation}

\begin{lemma}[Intertwining of singular projectors]
\label{lem:v5v100-intertwining}
The functional calculus gives
\begin{equation}
 DP_R=P_LD,
 \qquad
 DQ_R=Q_LD.
 \label{eq:v5v100-intertwining}
\end{equation}
The high map
\begin{equation}
 E:=Q_LDQ_R:Q_R\mathcal H_R\longrightarrow Q_L\mathcal H_L
 \label{eq:v5v100-high-map}
\end{equation}
has all singular values at least \(\Lambda\) and therefore
\begin{equation}
 \|E^{-1}\|\le\Lambda^{-1}.
 \label{eq:v5v100-high-inverse}
\end{equation}
\end{lemma}

\begin{proof}
For polynomials, \(D(D^*D)^m=(DD^*)^mD\).  Continuous functional calculus
and spectral approximation extend this identity to the characteristic
functions separated from the spectrum, proving
\eqref{eq:v5v100-intertwining}.  On the range of \(Q_R\),
\(D^*D\ge\Lambda^2\), hence \(\|Du\|\ge\Lambda\|u\|\).  The corresponding
left and right high spaces have equal dimension after zero modes are placed
in the low spaces, so \(E\) is bijective and
\eqref{eq:v5v100-high-inverse} follows.
\end{proof}

\begin{proposition}[Smooth covariant projector chart]
\label{prop:v5v100-Riesz-chart}
On the gap chart, \(P_R\) and \(P_L\) are smooth functions of the
background.  For a parameter derivative \(d\),
\begin{equation}
 dP_R=\frac1{2\pi i}\oint_\Gamma
 (z-H_R)^{-1}(dH_R)(z-H_R)^{-1}\,dz,
 \label{eq:v5v100-dP-Riesz}
\end{equation}
and similarly for \(P_L\), where \(\Gamma\) separates the low and high
spectra.  In particular,
\begin{equation}
 \|dP_R\|+\|dP_L\|
 \le C_\Gamma\delta^{-1}
 \bigl(\|dH_R\|+\|dH_L\|\bigr).
 \label{eq:v5v100-dP-bound}
\end{equation}
If \(D(A^g)=U_L(g)D(A)U_R(g)^{-1}\), the projectors transform by the
corresponding conjugations and the high/low bundles are gauge covariant.
\end{proposition}

\begin{proof}
The spectral projector is the Riesz integral of the resolvent.  Differentiate
under the contour and use
\(d(z-H)^{-1}=(z-H)^{-1}(dH)(z-H)^{-1}\), with the sign fixed by
differentiating \((z-H)(z-H)^{-1}=I\).  The contour can be chosen a fixed
fraction of \(\delta\) from the spectrum, giving
\eqref{eq:v5v100-dP-bound} after absorbing its length into \(C_\Gamma\).
Gauge covariance of \(H_R,H_L\) passes through functional calculus.
\end{proof}

Within such a chart, V99 applies with a smooth covariant projector bundle.
The derivative of every block includes the \(dP\) terms prescribed by
\eqref{eq:v5v100-dP-Riesz}; none may be omitted.

\subsection{Exact transfer at a spectral crossing}

At a threshold crossing, enlarge the retained space by a finite crossing
cluster \(\mathcal H_C\).  Write the fermion space as
\begin{equation}
 \mathcal H_F=\mathcal H_L\oplus\mathcal H_C\oplus\mathcal H_H.
 \label{eq:v5v100-three-way-split}
\end{equation}
Assume only that the terminal high block \(E=D_{HH}\) is invertible.  After
integrating it, define the exact operator on \(L\oplus C\):
\begin{equation}
 \widetilde D_{LC}:=
 D_{LC,LC}-D_{LC,H}E^{-1}D_{H,LC}
 =\begin{pmatrix}\widetilde A&\widetilde B\\
                  \widetilde C&\widetilde F\end{pmatrix}.
 \label{eq:v5v100-LC-Schur}
\end{equation}

\begin{theorem}[Associative crossing transfer]
\label{thm:v5v100-crossing-transfer}
The exact Grassmann identity is
\begin{equation}
 \int d\bar\psi_Hd\psi_H\,e^{-\bar\psi D\psi}
 =\det E\,e^{-\bar\psi_{LC}\widetilde D_{LC}\psi_{LC}}.
 \label{eq:v5v100-retain-crossing}
\end{equation}
It is valid even when \(\widetilde F\) is singular.  On the overlap where
\(\widetilde F\) is invertible, one may also integrate the crossing cluster:
\begin{align}
 \det D
 &=\det E\,\det\widetilde F\,
   \det\left(\widetilde A-widetilde B
          \widetilde F^{-1}\widetilde C\right),
 \label{eq:v5v100-sequential-determinant}\\
 \int d\bar\psi_Cd\psi_Cd\bar\psi_Hd\psi_H\,
 e^{-\bar\psi D\psi}
 &=\det E\det\widetilde F\,
 e^{-\bar\psi_L(\widetilde A-widetilde B
 \widetilde F^{-1}\widetilde C)\psi_L}.
 \label{eq:v5v100-sequential-Grassmann}
\end{align}
Thus moving \(\mathcal H_C\) from high to retained status changes the
factorization but not the full determinant or Grassmann functional.  At
\(\det\widetilde F=0\), formula \eqref{eq:v5v100-retain-crossing} remains
regular and no inverse of the crossing block is taken.
\end{theorem}

\begin{proof}
Apply Theorem~\ref{thm:v5v99-Grassmann-Feshbach} first to the split
\((L\oplus C)\oplus H\), proving
\eqref{eq:v5v100-retain-crossing}.  When \(\widetilde F\) is invertible,
apply the same theorem to \(L\oplus C\), which proves
\eqref{eq:v5v100-sequential-Grassmann}; setting all retained Grassmann
variables to be integrated gives
\eqref{eq:v5v100-sequential-determinant}.  The first identity is a finite
polynomial in \(\widetilde D_{LC}\), so it remains valid when
\(\widetilde F\) is singular.
\end{proof}

\begin{corollary}[Crossing-safe atlas]
\label{cor:v5v100-crossing-atlas}
At every finite cutoff and volume, field space admits local spectral charts
with thresholds avoiding the finite spectrum.  On chart overlaps, the
transition between two choices is a finite sequence of the exact crossing
transfers in Theorem~\ref{thm:v5v100-crossing-transfer}.  Hence all modes
approaching a high-block gap closure can be moved to the retained
Grassmann sector before the V98 modulus comparison is applied.
\end{corollary}

\begin{proof}
At a fixed configuration the finite spectrum has open complementary
intervals; continuity preserves a smaller gap in a neighbourhood.  Two
thresholds differ by finitely many eigenvalue clusters.  Transfer those
clusters one at a time with
Theorem~\ref{thm:v5v100-crossing-transfer}.
\end{proof}

\subsection{Determinant-line phase on overlaps}

The modulus of the high determinant is independent of bases in the high
domain and codomain.  Its complex phase belongs to the determinant line of
the corresponding bundles.

\begin{proposition}[Overlap transition for the high phase]
\label{prop:v5v100-determinant-line}
Let two smooth orthonormal frames of the same high right and left bundles
be related by unitaries \(U_R,U_L\).  Their high matrices satisfy
\begin{equation}
 E'=U_L^*EU_R,
 \qquad
 \det E'=\overline{\det U_L}\,\det U_R\,\det E.
 \label{eq:v5v100-phase-transition}
\end{equation}
Therefore \(|\det E|\) is a scalar positive density, while its phase is a
section with the displayed transition law.  A global phase requires a
trivialization whose cocycle is compatible with the regulated anomaly
cancellation of V92.
\end{proposition}

\begin{proof}
The matrix transformation follows from changing orthonormal bases in the
domain and codomain.  Take determinants and use
\(\det(U_L^*)=\overline{\det U_L}\).
\end{proof}

\subsection{The all-order direction after the audit}

The YM BKAR result can be imported only as an abstract positive-reference
lemma.  Let \(\gamma\) be a Gaussian massive block law and let
\(F_i\) be rooted local functions.  Its connected cumulant is a sum over
trees, with one covariance contraction per edge.  If complete derivative
rows obey \(\|\nabla^dF_i\|\le d!K^dw_i^d\), the tree factorial inequality
gives an \(r!K_1^r\) bound.  V90 supplies the spatially summable massive
edge kernel needed for the SM version.

\begin{proposition}[Conditional SM Gaussian tree card]
\label{prop:v5v100-SM-Gaussian-tree}
On a finite massive-safe chart satisfying V90, suppose the positive
reference block law is Gaussian with covariance kernel \(C\), and the
complete rooted insertions satisfy volume-uniform derivative-row bounds
\begin{equation}
 \|\nabla^dF_x\|_{\rm row}\le d!K_0^dw_x^d.
 \label{eq:v5v100-SM-derivative-row}
\end{equation}
Then their connected cumulants satisfy an all-order rooted tree bound
\begin{equation}
 |\kappa_r(F_{x_1},\ldots,F_{x_r})|
 \le r!K_1^r
 \sum_{T\in\mathfrak T_r}
 \prod_{\{i,j\}\in E(T)}|C(x_i,x_j)|
 \prod_i w_{x_i}^{d_i(T)}.
 \label{eq:v5v100-SM-tree-card}
\end{equation}
If \(C\) has a uniform \(\ell^1\) row bound and the activities are small
enough, the rooted spatial tree sum is finite at every arity with an
exponential arity constant.
\end{proposition}

\begin{proof}
Apply the connected Gaussian BKAR formula audited above.  Each tree edge
contracts one derivative index with \(C\).  Complete row--column
Cauchy--Schwarz and \eqref{eq:v5v100-SM-derivative-row} give the product of
vertex weights and factorials.  Repeated leaf removal gives
\(\prod_i d_i(T)!\le(r-1)!\), which is absorbed into \(r!K_1^r\).
For the spatial sum, remove leaves successively and use the \(\ell^1\) row
bound of \(C\) at each edge.
\end{proof}

\begin{firewall}[The coupled reference is not Gaussian]
\label{fw:v5v100-not-Gaussian}
The positive gauge--Higgs--PV modulus measure of V98 is not the Gaussian
law in Proposition~\ref{prop:v5v100-SM-Gaussian-tree}.  An all-derivative
transport, interpolation, or cluster expansion with exponential arity
constants is still needed.  The YM audit identifies the right target but
does not prove it for the SM measure.
\end{firewall}

\begin{firewall}[The spectral atlas is not uniformly finite]
\label{fw:v5v100-atlas-uniformity}
Corollary~\ref{cor:v5v100-crossing-atlas} is finite-cutoff and
finite-volume.  As volume grows, gaps between singular values can shrink
and the number of crossing charts can grow.  A multiscale ownership rule
must bound the total retained rank and transition complexity.
\end{firewall}

\begin{firewall}[Determinant-line triviality remains an anomaly problem]
\label{fw:v5v100-determinant-line}
Exact Standard Model representation-theory anomaly cancellation does not
automatically trivialize the determinant line of a particular lattice
regulator over the complete gauge-field space.  The required local
counterterms, overlap cocycle, and possible global anomaly must be checked
for the chosen construction.
\end{firewall}

\begin{theorem}[V100 conclusion]
\label{thm:v5v100-status}
Away from a singular-value threshold crossing, the light and high fermion
spaces form smooth gauge-covariant Riesz bundles with explicit derivative
bounds, and the high map has inverse norm at most \(\Lambda^{-1}\).  At a
crossing, associative Grassmann integration transfers the crossing cluster
to the retained sector without dividing by its determinant; the
factorizations agree exactly on overlaps.  The high modulus is a scalar
positive factor, while its phase obeys a determinant-line transition law.
The audited Gaussian BKAR theorem supplies an all-order tree card on the
massive Gaussian reference.  Uniform spectral-atlas complexity,
determinant-line trivialization, and transport of the all-order tree card
to the coupled positive measure remain open.
\end{theorem}

\begin{proof}
Use Lemma~\ref{lem:v5v100-intertwining},
Proposition~\ref{prop:v5v100-Riesz-chart},
Theorem~\ref{thm:v5v100-crossing-transfer},
Corollary~\ref{cor:v5v100-crossing-atlas}, and Propositions
\ref{prop:v5v100-determinant-line} and
\ref{prop:v5v100-SM-Gaussian-tree}, subject to the three firewalls.
\end{proof}

This is the V100 boundary of the fifth SM volume.  The complete file will
now be frozen as \texttt{Renewal\_SM\_Einstein\_Continuum\_Research\_Notes\_V100\_f5.tex}.
The active continuation will restart at V1 with a compact context and
proved-result ledger.  Its first new target is an interacting all-order
transport card from the V100 massive Gaussian reference to the positive
gauge--Higgs--PV modulus measure, with determinant crossings retained by
the atlas above.

\end{document}
